% Parte V del sucesor: mismos propietarios rectores, títulos capitulares
% aprobados y dependencia causal explícita respecto del continuo conjunto.
\makeatletter
\@ifundefined{HMTScientificOwnerSubordinated}{%
  % Entorno editorial para incorporar cuerpos científicos completos sin
% crear capitulos nuevos. El grupo conserva intactos los archivos de origen:
% solo traduce su jerarquia visible un nivel hacia abajo durante el input.
% La procedencia material de cada uso se registra en comentarios del destino
% y en gestion/RESTAURACION_16_CLAVES_SIN_DESTINO_20260824.tsv.
\newenvironment{HMTScientificOwnerSubordinated}{%
  \begingroup
  \let\HMTScientificOwnerSection\section
  \let\HMTScientificOwnerSubsection\subsection
  \let\HMTScientificOwnerSubsubsection\subsubsection
  \let\HMTScientificOwnerParagraph\paragraph
  \let\HMTScientificOwnerSubparagraph\subparagraph
  \let\chapter\HMTScientificOwnerSection
  \let\section\HMTScientificOwnerSubsection
  \let\subsection\HMTScientificOwnerSubsubsection
  \let\subsubsection\HMTScientificOwnerParagraph
  \let\paragraph\HMTScientificOwnerSubparagraph
}{%
  \par
  \endgroup
}
%
}{}
\makeatother

\part{Réalisations terminales de la structure discrète du continu : \texorpdfstring{\(\zeta\)}{zeta}, Navier--Stokes, Yang--Mills, Hodge et Birch--Swinnerton--Dyer}
\label{part:clausuras-terminales-definitiva}

\HMTFlushCurrentChapter
\chapter{Théorème de factorisation conjointe des \(5\) réalisations terminales à travers la structure discrète du continu}
\label{chap:terminal-arquitectura-natural-comun}
La Partie V reçoit le continu unique après sa construction conjointe et
sa réalisation par la mécanique dimensionnelle. Aucune réalisation terminale
ne sélectionne, ne génère ni ne remplace rétrospectivement les constructions
consubstantielles du continu. La séquence d'exposition conserve, dans cet ordre,
la note de lecture, le raccord causal, l'objet commun et la méthode de
démonstration.
Le théorème conjoint établit la factorisation simultanée, la non-
permutabilité des lecteurs et la conservation de l'état enrichi
commun. Les réalisations qui suivent conservent intégralement, et avec un domaine
propre, leurs opérateurs, démonstrations et certificats individuels ; elles ne sont
ni des substituts ni des abrégés de cette factorisation commune.
\Needspace{6\baselineskip}
\Needspace{6\baselineskip}
\section{Publication terminale du continu unique par \(5\) lecteurs non permutables}
\Needspace{6\baselineskip}
\subsection{Factorisation de l’état enrichi commun en \(5\) publications terminales}
\Needspace{6\baselineskip}
\subsubsection{Dépendance causale des réalisations terminales à l’égard du continu généalogique intégralement construit}
Les cinq problèmes ne font pas appel à cinq structures mathématiques différentes. Ils n’obligent pas non plus à construire cinq chemins indépendants qui, après avoir progressé séparément, devraient être réunis artificiellement dans une conclusion commune. Cette image appartient à une étape antérieure de l’exposé. Le point de départ actuel est plus fort : avant que n’apparaissent les noms de Riemann, Navier–Stokes, Yang–Mills, Hodge ou Birch–Swinnerton–Dyer, il existe déjà un unique objet mathématique enrichi dont les opérations internes produisent simultanément les caractéristiques nécessaires aux cinq fermetures.\par
C’est pourquoi ce récit ne commence pas par les problèmes. Il commence par la construction du continu.\par
L’Holographie modulaire triadique ne reçoit pas le continu comme un arrière-plan préalablement donné. Elle ne part pas davantage de la droite réelle, d’un espace fonctionnel classique, d’une variété algébrique, d’une connexion de jauge ou d’une courbe elliptique. Elle part d’une arithmétique finie dotée d’orientation, de mémoire et d’une capacité de prolongement. APP fournit le substrat positionnel et arithmétique ; TRIT transforme chaque opération en une transition orientée ; TPK conserve la généalogie complète de ces transitions et permet de les prolonger sans effacer le chemin qui les a produites.\par
L’étape décisive tient à ce que la sortie d’une opération cesse d’être seulement une valeur. Chaque émission porte son feuillet arithmétique, son résidu, son quotient, son orientation, sa retenue, sa mémoire, sa frontière, son parcours, sa multisection, sa survie, son terme terminal et son lecteur. Le résultat n’est plus un point isolé, mais un état qui conserve sa provenance, ses possibilités de continuation et les relations à préserver lorsqu’il change d’échelle.\par
APP, TRIT et TPK ne forment pas trois explications superposées. Ils forment une seule séquence causale :\par
\[
\mathrm{APP}
\longrightarrow
\mathrm{TRIT}
\longrightarrow
\mathrm{TPK}
\longrightarrow
\text{état enrichi unique}.
\]
APP fixe le support arithmétique commun. TRIT introduit l’orientation effective de chaque transition, en distinguant avancée, seuil et ouverture future. TPK réunit l’information visible et l’information transportée : non seulement ce qu’une opération fournit, mais aussi la retenue qu’elle déplace, la frontière qu’elle modifie, la mémoire qu’elle conserve et les extensions qu’elle permet.\par
La connexion nonadique ne ramène pas le système à un état vide. Après un tour complet, la phase réapparaît, mais l’histoire n’est pas effacée. Le retour \(9\to1\) est un retour de phase, non une réinitialisation de l’état. Chaque tour ajoute profondeur, mémoire et compatibilité. Ainsi, la continuité ne s’obtient ni en répétant neuf filtres ni en traversant neuf canaux indépendants. Elle s’obtient en prolongeant une seule connexion complète dont les neuf segments sont des cartes internes d’une même holonomie.\par
\Needspace{6\baselineskip}
\subsubsection{Émergence corrélative de \texorpdfstring{\(5\)}{5} actions intrinsèques non exhaustives}
Une opération commune agit sur chaque état enrichi. Il n’existe pas d’abord cinq domaines, ni cinq sélections, ni cinq copies de l’état. La même émission est lue simultanément selon cinq caractéristiques internes désignées par les discriminants \(\mathrm{sp}\), \(\mathrm{sol}\), \(\mathrm{gau}\), \(\mathrm{coh}\) et \(\mathrm{det}\).\par
Ces symboles ne désignent pas des branches et n’épuisent pas les propriétés du continu. Ils désignent \(5\) actions intrinsèques du même objet dont la corrélation est décisive pour les réalisations terminales de cette Partie.\par
La caractéristique spectrale–polaire, \(sp\), enregistre l’expression orientée de l’état, son complément, son défaut, sa mémoire terminale et la relation entre les deux lectures spéculaires d’une même forme.\par
La caractéristique solénoïdale, \(sol\), enregistre la manière dont une transition conserve l’état tout en répartissant dissipation, advection, retenue, mémoire, microstructure et coquille sans les compter plusieurs fois.\par
La caractéristique de jauge, \(gau\), enregistre l’équivalence locale, l’incidence, l’holonomie, la courbure et le passage compatible au quotient physique.\par
La caractéristique cohomologique, \(coh\), enregistre les extensions entre niveaux, les obstructions relatives, la compatibilité des relèvements et la survie d’une classe à travers tout raffinement.\par
La caractéristique déterminantale, \(det\), enregistre l’incidence finie-globale, les pages de filtration, les complexes locaux, leurs droites déterminantes et le transport d’une base arithmétique jusqu’à une expression analytique.\par
Mais aucune de ces propriétés n’existe seule. Toutes agissent sur les mêmes émissions et conservent la même épaisseur généalogique. La fibre qui transporte la caractéristique spectrale est la même fibre enrichie dont l’orientation participe au bilan solénoïdal, dont la frontière porte la lecture de jauge, dont le prolongement soutient l’extension cohomologique et dont l’incidence entre dans la droite déterminante.\par
On n’engendre pas cinq objets pour les regrouper ensuite. On engendre un unique objet qui possède indissociablement ces cinq capacités.\par
Une vue ultérieure n’est donc pas une projection qui crée une branche. Elle est seulement une manière d’exposer, pour une application concrète, une caractéristique qui agissait déjà au sein de la totalité. La vue ne sélectionne pas rétrospectivement des états favorables et ne conserve pas l’entrée comme première coordonnée pour la retrouver ensuite de manière tautologique. Elle rend visible une structure produite et conservée tout au long de la généalogie.\par
\Needspace{6\baselineskip}
\subsubsection{Antériorité causale du continu déjà construit comme limite du prolongement compatible}
À chaque profondeur existe un état fini, mais complet à son niveau. Les extensions entre niveaux ne peuvent modifier arbitrairement l’information antérieure. Elles doivent respecter les fibres, les relations d’extension, les nombres, l’orientation, la retenue, la mémoire, la frontière, le parcours, la multisection, la survie, le terme terminal et les lecteurs.\par
La naturalité exige qu’opérer puis tronquer donne le même résultat que tronquer puis opérer. Grâce à cette compatibilité, les histoires finies ne demeurent pas des approximations déconnectées. Elles forment un système inverse. Le continu apparaît lorsqu’on considère la famille complète des histoires compatibles, et non lorsqu’on postule d’emblée une collection infinie de points.\par
Le continu HMT est donc mémoire cohérente de chemins. Un point sans généalogie serait une ablation : il conserverait la valeur visible et supprimerait précisément ce qui permet de démontrer comment elle survit à tout raffinement.\par
Ce n’est qu’après avoir établi cette compatibilité qu’apparaît la structure discrète complète du continu. Ce n’est qu’ensuite qu’une réalisation archimédienne peut être exprimée. La droite réelle n’est pas la matière première de la construction ; elle est une sortie de l’un de ses lecteurs terminaux. Il en va de même pour les domaines classiques des cinq problèmes : ils apparaissent lorsque le continu est déjà construit et qu’une interface ultérieure traduit l’une de ses caractéristiques intrinsèques dans le langage de l’objectif correspondant.\par
La séquence causale correcte est donc :\par
\[
\begin{aligned}
\mathrm{APP}
&\longrightarrow \mathrm{TRIT}
\longrightarrow \mathrm{TPK}
\longrightarrow \text{état enrichi unique}\\
&\longrightarrow 04\mathrm b1
\longrightarrow 04\mathrm b2
\longrightarrow 04\mathrm b3\\
&\longrightarrow \text{actions intrinsèques corrélatives}
\longrightarrow \text{terminal et limite}
\longrightarrow \mathfrak C_{\mathrm{cont}}^{\mathrm{disc}}.
\end{aligned}
\]
Les entrelaceurs spécifiques interviennent ensuite, et pas avant :\par
\[
\mathfrak C_{\mathrm{cont}}^{\mathrm{disc}}
\longrightarrow
\begin{cases}
\text{forme explicite de Weil},\\
\text{évolution tridimensionnelle incompressible},\\
\text{théorie de jauge quantique},\\
\text{publication des classes de Hodge},\\
\text{comparaison déterminantielle BSD}.
\end{cases}
\]
Ces cinq traductions terminales ne définissent pas la structure qu’elles utilisent. Le problème classique apporte ses données après coup : une fonction test pour Riemann, des données initiales et une force pour Navier–Stokes, un groupe compact simple pour Yang–Mills, une variété et une classe pour Hodge, ou une courbe elliptique pour BSD. Aucune de ces données ne participe à la génération préalable du continu.\par
La direction causale est ainsi protégée. L’objectif ne sélectionne pas l’état qui devra le résoudre. L’état existe avant l’objectif ; l’entrelaceur démontre seulement que la structure interne et la formulation classique expriment la même opération dans deux langages.\par
\Needspace{6\baselineskip}
\subsection{1re résolution : Riemann et la positivité comme complément construit}
L’hypothèse de Riemann est généralement présentée comme une assertion concernant la position des zéros non triviaux de la fonction zêta. Dans cette formulation, le problème semble demander d’examiner une population de zéros déjà donnée et de démontrer qu’ils s’alignent tous sur la droite critique.\par
La lecture HMT inverse cette perspective. Elle ne commence pas par demander où sont les zéros. Elle demande quelle structure engendre simultanément les deux membres de la formule explicite et ce qui empêche leur différence de devenir négative.\par
La caractéristique spectrale–polaire du continu fournit une double expression du même état. Une orientation recueille le canal positif ; l’autre recueille son image spéculaire. Toutes deux contiennent, sous un régulateur commun, la contribution première, la contribution gamma–scalaire et la contribution polaire. Il ne s’agit ni de trois preuves différentes ni de trois sommes accolées à la fin. Ce sont trois lignes d’une même colligation construite à partir des données initiales.\par
Le point essentiel est que cette colligation est construite avant d’invoquer la positivité qu’elle devra démontrer. Son entrée, ses sorties, sa perte et son état terminal sont typés dans le même espace enrichi. L’identité qui relie les deux expressions n’est pas introduite comme une égalité souhaitée entre formules classiques ; elle naît de la dynamique interne de la connexion.\par
Entre l’expression d’entrée et l’expression spéculaire peut apparaître, à profondeur finie, un résidu de sortie. La structure nonadique lui impose une récurrence antispéculaire. À chaque tour complet, le résidu suivant est transporté isométriquement, mais multiplié par le facteur\par
\[
q=\frac{5}{9}.
\]
Si \(\Omega_n\) représente ce défaut à la profondeur \(n\), la récurrence prend la forme\par
\[
\Omega_n=q\,V_n\Omega_{n+1},
\]
avec \(V_n\) isométrique et la famille coinductivement bornée. En itérant \(N\) fois,\par
\[
\|\Omega_n\|
\le q^N\sup_{m\ge n}\|\Omega_m\|.
\]
Comme \(0<q<1\), le membre de droite tend vers zéro. Le résidu ne s’annule donc pas parce qu’il a été omis ou parce qu’une identification entre les deux expressions a été supposée. Il s’annule parce que la récurrence complète le contracte à toutes les profondeurs compatibles. La sortie spéculaire coïncide alors avec l’expression négative construite.\par
Une fois ce lien établi, l’isométrie construite à partir des données initiales fournit l’identité d’énergie. À profondeur finie, la norme de l’expression positive se décompose en trois parties : l’expression négative, la perte accumulée et l’état terminal. Le terme terminal est explicitement conservé. Il doit être conservé jusqu’au passage à la limite, car il contient la mémoire de la troncature et certifie que le télescopage prolonge l’identité finie par une convergence démontrée.\par
La décomposition finie de la norme satisfait l’identité exacte :\par
\[
\|\text{publication positive}\|^2
=
\|\text{publication négative}\|^2
+
\|\text{pertes jusqu'à }N\|^2
+
q^{2N}\|\text{terminal}\|^2.
\]
Chaque terme possède une origine concrète. L’expression négative est la sortie spéculaire déjà identifiée. Les pertes réunissent exactement le complément qui n’a pas été exprimé par cette orientation. Le terme terminal conserve ce qui demeure encore au-delà de la coupe.\par
Ce n’est qu’après avoir écrit l’identité complète que l’on fait \(N\to\infty\). Alors, et alors seulement, le terme terminal disparaît parce que \(q^{2N}\to0\). La différence entre les deux expressions est transformée en carré d’un opérateur de perte :\par
\[
\mathcal I_{\mathrm{GW}}
=
E_R^{*}E_R.
\]
C’est la forme de Guinand–Weil. Sa positivité n’est plus une conjecture ajoutée à la construction :\par
\[
\langle f,\mathcal I_{\mathrm{GW}}f\rangle
=
\|E_Rf\|^2
\ge0.
\]
La structure HMT n’obtient pas cette inégalité en dissimulant la partie négative, en restreignant le domaine aux fonctions qui la satisfont déjà, ou en introduisant un projecteur dont l’existence serait équivalente à la conclusion. Elle construit le complément et démontre que la forme est exactement son carré.\par
L’expression commune est ensuite matérialisée. Un même encodeur conserve les canaux premier, gamma–scalaire et polaire. Un projecteur interne produit l’un des deux moments ; son complément produit l’autre. La différence des moments n’est pas une comparaison externe entre formules semblables : c’est la norme de la composante orthogonale construite par la colligation.\par
Lorsqu’on élargit la fenêtre, on n’oblige pas les colonnes individuelles à satisfaire de fausses isométries entre troncatures distinctes. Ce qui est conservé de manière cofinale, c’est la forme, le complément atteignable et la naturalité de son expression. Les incréments introduits par l’élargissement de la coupe apparaissent dans les deux orientations avec la même contribution neutre et ne modifient pas la différence directrice. On obtient ainsi une famille compatible de formes positives sur une classe séparante.\par
La dernière transition est le critère réversible de Weil. La théorie classique affirme que la positivité de la forme explicite, sur la classe adéquate de fonctions tests, équivaut au fait que tous les zéros non triviaux se trouvent sur la droite critique. HMT fournit précisément cette positivité, mais comme conséquence d’une identité opératorielle construite à l’intérieur du continu.\par
\Needspace{5\baselineskip}
Por tanto,\par
\[
\operatorname{Re}\rho=\frac12
\]
pour tout zéro non trivial \(\rho\) de la fonction zêta.\par
La droite critique apparaît ici de deux manières concordantes. Géométriquement, elle est le lieu fixe de l’involution qui échange les deux orientations spectrales. Du point de vue opératoriel, elle est l’unique position compatible avec la positivité du complément de l’expression. La symétrie explique pourquoi la droite centrale est distinguée ; l’identité quadratique démontre pourquoi les zéros ne peuvent quitter cette droite.\par
La résolution ne consiste pas à observer que l’équation fonctionnelle possède une symétrie. Elle consiste à construire l’objet dont le défaut mesure la rupture de cette symétrie, à démontrer que le défaut s’annule coinductivement et à transformer la différence restante en une norme. C’est la combinaison de la généalogie, du terme terminal, de la contraction, de l’expression commune et de la positivité qui clôt le problème.\par
La première des cinq résolutions terminales est ainsi établie. Non comme une branche indépendante, mais comme la traduction classique de la caractéristique spectrale–polaire que le continu unique possédait dès sa construction.\par

\Needspace{6\baselineskip}
\section{Fermetures dynamiques du continu : transport solénoïdal et quotient de jauge}
\Needspace{6\baselineskip}
\subsection{Dynamique et courbure du continu généalogique}
La première résolution a montré comment la caractéristique spectrale–polaire du continu s’exprime comme positivité de Weil. Les deux résolutions suivantes interrogent le même état sous deux autres aspects consubstantiels : sa capacité à transporter une dynamique non linéaire sans perdre énergie, mémoire ni régularité, et sa capacité à soutenir une théorie de jauge locale dont le vide est séparé par un saut spectral positif.\par
Navier–Stokes et Yang–Mills ne reçoivent pas deux structures nouvelles. Ils reçoivent deux entrelaceurs différents à partir du même continu déjà construit. Dans le premier cas, l’entrelaceur exprime l’état comme un champ de vitesse incompressible avec viscosité, advection et pression. Dans le second, il l’exprime comme un réseau de jauge raffinable avec holonomie, courbure, mesure euclidienne et Hamiltonien physique.\par
La différence entre les deux problèmes ne rompt pas l’unité de la construction. Pour Navier–Stokes, le point crucial est de démontrer qu’aucun transfert non linéaire ne disparaît hors du bilan et que la régularité obtenue est uniforme lorsqu’on retire les coupures. Pour Yang–Mills, le point crucial est de démontrer que le saut fini appartient à la même théorie de jauge qui atteint le continu, qu’il subsiste au quotient physique et qu’il agit dans le secteur de courbure, non dans un facteur auxiliaire.\par
\Needspace{6\baselineskip}
\subsection{Existence, régularité et unicité globales pour Navier–Stokes par borne uniforme et compacité}
Le problème tridimensionnel de Navier–Stokes ne naît pas de la seule présence d’un terme non linéaire. La difficulté apparaît parce que l’expression classique du champ de vitesse tend à masquer le lieu où se stocke le transfert entre échelles. La viscosité est visible. L’énergie cinétique est visible. L’advection apparaît également dans l’équation. Mais la comptabilité complète de l’interaction entre profondeur, mémoire, retenue, microstructure et coquille est comprimée en une seule fonction visible.\par
C’est précisément cette compression que la caractéristique solénoïdale du continu évite.\par
L’état enrichi ne conserve pas uniquement le champ \(u\). Il conserve simultanément le champ de Galerkin, l’orientation temporelle, les mémoires positive et signée, la retenue entre coupures, la composante microscopique, la coquille des fréquences nouvelles, le parcours, l’archive de l’évolution et le terme terminal. L’équation classique est une expression de cet état ; elle n’en est pas la définition complète.\par
Les données initiales \(u_0\), la force \(f\), la viscosité \(\nu>0\) et la régularité \(s>5/2\) étant fixées, la projection de Galerkin produit\par
\[
\partial_tu_M+\nu A_Mu_M+B_M(u_M,u_M)=f_M.
\]
Mais \(M\) ne sélectionne pas une solution favorable. Elle fixe seulement une résolution finie dans laquelle toutes les opérations peuvent être écrites exactement. L’objectif est de construire une estimation indépendante de \(M\) et de la profondeur nonadique \(J\). Seule une borne présentant cette uniformité peut passer au continu.\par
\Needspace{6\baselineskip}
\subsubsection{Décomposition de la mémoire visible et de la mémoire généalogique dans l’équation d’évolution}
L’interaction non linéaire possède un signe. À certains instants, elle apporte de l’énergie à une échelle ; à d’autres, elle lui en retire. Si l’on ne conserve que sa valeur signée, l’amplitude totale transférée peut être perdue. Si l’on ne conserve que sa valeur absolue, l’orientation est perdue.\par
La structure maintient les deux.\par
Une mémoire \(D_M\) exprime le signe du transfert advectif. Une autre mémoire \(H_M\) conserve son amplitude positive. Leur séparation permet d’enregistrer simultanément\par
\[
\dot D_M=\mathscr B_{s,M},
\qquad
\dot H_M=\frac{1+r}{r-1}|\mathscr B_{s,M}|,
\]
où \(r=\pi_{\rm HMT}/\varphi_{\rm HMT}^2>1\).\par
Cela n’ajoute pas d’énergie fictive. Cela dédouble deux aspects distincts d’un seul transfert : orientation et coût. Les confondre détruirait de l’information ; les additionner sans orthogonalité doublerait l’énergie. La construction montrera explicitement qu’aucune de ces deux déformations ne se produit.\par
\Needspace{6\baselineskip}
\subsubsection{Indépendance de la donnée initiale à l’égard de l’histoire généalogique ultérieure}
L’objection la plus sérieuse à toute fermeture opératorielle de Navier–Stokes serait que la borne uniforme ait été introduite rétrospectivement : on choisirait d’abord une solution régulière, puis on coderait son histoire, pour annoncer enfin que le code contrôle cette même solution.\par
La construction empêche cette circularité en fixant l’espace source avant de relever l’évolution :\par
\[
\mathcal D_T
=
H^s_\sigma(\mathbb T^3)
\times
L^2(0,T;H^{s-1}_\sigma(\mathbb T^3)),
\qquad
d=(u_0,f).
\]
La racine \(J_T^{\rm sol}d\) dépend exclusivement de \(u_0\) et de \(f\). Elle ne reçoit en entrée ni \(u_M\), ni la perte future, ni le terme terminal, ni la norme à borner.\par
La force n’est pas non plus insérée comme un port additif artificiel. L’équation de Carleman montre qu’elle agit de manière mixte : chaque occurrence de \(f\) se combine avec un degré antérieur de l’état. Le relèvement utilise donc une création mixte état–force. La signature chronologique de \(f\) enregistre tous ses temps ordonnés, tandis que la tour PC–Fock enregistre tous les degrés non linéaires de l’état.\par
La composition reproduit matériellement la série de Dyson du générateur non autonome. L’application du lecteur normalisé de degré un restitue exactement \(u_M(t)\). Elle ne restitue ni une approximation symbolique ni une variable ressemblant à la solution : elle restitue la solution de l’EDO de Galerkin, car toutes deux satisfont la même équation finie et la même donnée initiale.\par
La construction causale à partir des données initiales suit donc l’ordre\par
\[
(u_0,f)
\longrightarrow
J_T^{\rm sol}(u_0,f)
\longrightarrow
\text{tour état–force}
\longrightarrow
\text{histoire de Galerkin}.
\]
L’histoire généalogique est construite comme sortie ultérieure des données initiales et est exclue de leur définition initiale.\par
\Needspace{6\baselineskip}
\subsubsection{L’identité de diffusion qui conserve la puissance croisée}
La force et la dissipation sont organisées au moyen de deux ondes :\par
\[
a_T(f)
=
\frac{1}{2\sqrt\nu}A^{-1/2}\Lambda^sP_\sigma f,
\qquad
z_M(u)=\sqrt\nu A_M^{1/2}\Lambda^su,
\]
et de l’onde résiduelle\par
\[
b_M=z_M-P_Ma_T(f).
\]
La puissance appliquée n’est pas remplacée par la norme de la force. Elle est conservée exactement comme terme croisé :\par
\[
\mathscr F_{s,M}
=
2\operatorname{Re}\langle z_M,P_Ma_T(f)\rangle.
\]
La complétion du carré donne le bilan exact\par
\[
\dot G_{M,j,T}
+\|b_M\|^2
+\lvert\mathscr B_{s,M}\rvert
=
\|P_Ma_T(f)\|^2.
\]
C’est l’un des points qui rendent la fermeture opérante. Si l’on remplaçait \(b=z-a\) par \(z\), le terme croisé représentant le travail de la force disparaîtrait et l’identité serait fausse. La structure ne maîtrise pas la non-linéarité en l’écartant : elle la transforme en un canal positif orienté et conserve en outre l’interférence exacte entre force et dissipation.\par
\Needspace{6\baselineskip}
\subsubsection{Factorisation de \(6\) pertes par un opérateur unique sans multiplicité parasite}
L’histoire complète distingue six sorties :\par
\begin{itemize}
\item dissipation ;
\item advection orientée ;
\item retenue entre coupures ;
\item mémoire microscopique ;
\item coquille de fréquences nouvelles ;
\item mémoire non héritée.
\end{itemize}
Ce ne sont pas six énergies ajoutées à un bilan préalable. Ce sont six lignes orthogonales d’un unique opérateur terminal de perte. L’identité de Parseval établit\par
\[
\|L_{M,j,T}^{\rm term}x\|^2
=
\sum_{\mathfrak a}
\|L_{M,j,T}^{\mathfrak a}x\|^2.
\]
L’orthogonalité n’est pas simplement nominale. Les profondeurs occupent des intervalles temporels disjoints ; retenue et coquille vivent dans des plages spectrales distinctes ; la microstructure appartient au complément \(Q_{\rm micro}\) ; la mémoire nouvelle est séparée de la mémoire héritée ; et les deux feuillets de l’advection ne peuvent être simultanément positifs en un même point.\par
La comptabilité ne double donc pas l’interaction. La somme des six normes est exactement la norme de l’opérateur commun. Si l’une d’elles n’était pas une projection de cet opérateur, Parseval échouerait avant l’obtention de la borne. Le formalisme contient ainsi son propre réfutateur.\par
\Needspace{6\baselineskip}
\subsubsection{Persistance de l’objet terminal sous le passage à la limite}
Toute troncature conserve quelque chose qui n’a pas encore été exprimé. Ce reste est le terme terminal. Il contient l’énergie visible à la fin de l’intervalle et la mémoire orientée qui demeure entre le temps observé et l’horizon \(T\).\par
La colonne terminale est construite en premier. Son Gram est formé ensuite. On ne définit pas une métrique souhaitée pour inventer ultérieurement un opérateur qui la réalise. La colonne contient matériellement la sortie terminale et toutes les pertes transportées :\par
\[
\mathbf G_{M,j;J,T}^{\rm term}
=
\begin{bmatrix}
\text{terminal transporté}\\
\text{perte}_{J-1}\\
\vdots\\
\text{perte}_{j}
\end{bmatrix}.
\]
Sa récurrence par blocs produit l’identité de Stein\par
\[
U_{M,j,T}
=
\Gamma_{M,j,T}^{*}
U_{M,j+1,T}
\Gamma_{M,j,T}
+
(L_{M,j,T}^{\rm term})^{*}
L_{M,j,T}^{\rm term}.
\]
Cette identité ne définit pas rétroactivement la norme de l’état. Elle s’obtient en multipliant une colonne déjà construite par son adjointe. Pour la même raison, le télescopage global est une identité entre opérateurs existants, et non une hypothèse de stabilité.\par
La mémoire future n’est pas non plus introduite dans la racine initiale. Elle réside dans la sortie terminale. Lorsqu’on prend l’estimation intégrale en \(\tau=T\), son support temporel se vide exactement. Elle n’est pas effacée par commodité.\par
\Needspace{6\baselineskip}
\subsubsection{Borne uniforme préalable et compacité d’Aubin–Lions}
La racine, l’encodeur de niveau, la colonne et l’isométrie commune satisfont la factorisation dérivée des données initiales\par
\[
\mathcal G_{M,J,T}^{\rm sol}
\Lambda_{M,T}
J_T^{\rm sol}d
=
I_{M,J}J_T^{\rm sol}d,
\qquad
I_{M,J}^{*}I_{M,J}=I.
\]
Le membre de gauche est l’histoire complète à la coupure \((M,J)\). Celui de droite dépend de l’unique racine construite à partir des données. Puisque \(I_{M,J}\) est isométrique, la norme de l’histoire est dominée par la norme source avec une constante indépendante de \(M\) et de \(J\).\par
C’est l’étape qui résout la difficulté. L’uniformité ne provient pas d’une estimation de dimension finie dont la constante croît avec la coupure. Elle ne provient pas non plus d’une définition de l’espace des histoires au moyen de la perte même que l’on veut contrôler. La racine existe avant la coupure et toutes les coupures en sont des réalisations isométriques.\par
La colonne terminale traduit cette borne structurelle en grandeurs physiques :\par
\[
\begin{aligned}
&\sup_M\sup_{0\le t\le T}\|u_M(t)\|_{H^s}^2
+\nu\sup_M\int_0^T\|u_M(t)\|_{H^{s+1}}^2\,dt\\
&\quad
+\sup_M\int_0^T|\mathscr B_{s,M}(u_M(t))|\,dt
+\text{carry}+\text{micro}+\text{coquille}+\text{mémoire}
\le C_T(u_0,f).
\end{aligned}
\]
La constante est finie pour tout \(T<\infty\) et ne dépend ni de \(M\) ni de \(J\). L’inégalité ne contrôle pas seulement le champ visible. Elle contrôle simultanément le travail transporté et les secteurs qu’une expression classique tend à oublier.\par
\Needspace{6\baselineskip}
\subsubsection{Obstruction de signe dans la factorisation de Kalman–Yakubovich–Popov}
Faire dépendre la fermeture globale d’une factorisation KYP portant un signe incorrect introduit un carré positif et prétend le retrouver ensuite avec un signe négatif. Cette égalité est algébriquement fausse.\par
La démonstration exclut cette identité et construit directement l’estimation uniforme nécessaire.\par
La carte KYP conservée est une identité finie correcte sous ses propres conditions. Elle sert de contrôle auxiliaire de chaque coupure, mais ne fournit pas l’uniformité lorsqu’on retire \(M\) et \(J\). La preuve globale n’utilise pas de borne de \(Q^{-1}\), ne prend pas la limite à l’intérieur de KYP et ne dépend pas de ce signe.\par
L’obstruction du bloc KYP portant un signe incorrect n’affecte donc pas la démonstration. L’uniformité est établie par la construction directe à partir des données, avec sa colonne terminale et son identité de Stein.\par
\Needspace{6\baselineskip}
\subsubsection{Convergence de la borne uniforme vers la solution classique globale}
La borne précédente fournit\par
\[
u_M
\quad\text{borné dans}\quad
L^\infty(0,T;H^s)
\cap L^2(0,T;H^{s+1}).
\]
Comme \(s>5/2\), le produit de Sobolev contrôle la non-linéarité :\par
\[
B:H^s_\sigma\times H^s_\sigma
\longrightarrow H^{s-1}_\sigma.
\]
L’équation de Galerkin fournit en outre une borne uniforme de \(\partial_tu_M\) dans \(L^2(0,T;H^{s-1})\). Les inclusions\par
\[
H^{s+1}\Subset H^s\hookrightarrow H^{s-1}
\]
permettent d’appliquer Aubin–Lions. On obtient une convergence forte dans \(L^2H^s\), suffisante pour transporter le terme bilinéaire :\par
\[
B_M(u_M,u_M)\longrightarrow B(u,u)
\quad\text{dans}\quad
L^1(0,T;H^{s-1}).
\]
La limite satisfait l’équation incompressible. La pression n’est pas introduite comme une autre donnée cachée ; elle est reconstruite ensuite au moyen de\par
\[
-\Delta p=\partial_i\partial_j(u_i u_j),
\qquad
\int_{\mathbb T^3}p=0.
\]
L’unicité s’obtient sur la différence de deux solutions. Une estimation dans \(H^{s-1}\), suivie de Grönwall, impose que la différence soit nulle. Cela élimine la dépendance à l’égard de la sous-suite extraite par compacité.\par
La construction vaut pour tout horizon fini. Les solutions obtenues sur \([0,n]\) et \([0,n+1]\) coïncident sur leur chevauchement par unicité. Elles se recollent donc en une seule solution définie sur \([0,\infty)\).\par
La régularité lisse s’obtient sur toute l’échelle \(s+n\), à l’aide des mêmes lecteurs structurels ; l’unicité identifie toutes les limites et produit une solution globale lisse pour toute donnée périodique lisse. La borne requise est construite sur chaque horizon fini et pour chaque indice de Sobolev à partir de la donnée solénoïdale complète, comme le démontre le développement de référence \cref{thm:pdf2-ns-global}.\par
Enfin, la condition de moyenne nulle est une normalisation du secteur solénoïdal. Le changement de Galilée dépendant du temps sépare la moyenne et transporte bijectivement toute donnée périodique lisse vers le secteur de moyenne nulle, en préservant la divergence, la régularité lisse et les normes de Sobolev.\par
\Needspace{6\baselineskip}
\subsubsection{Existence, régularité et unicité globales de Navier--Stokes}
La caractéristique solénoïdale du continu, exprimée par son entrelaceur spécifique, produit une unique solution globale lisse des équations tridimensionnelles périodiques de Navier--Stokes pour toute donnée périodique lisse. La construction de Leray--Hopf est contenue comme étape d’énergie dans cette conclusion forte.\par
La défense de cette conclusion ne repose pas sur l’affirmation qu’il existe une structure capable de résoudre le problème. Elle repose sur cinq faits matériels :\par
la racine utilise uniquement \((u_0,f)\) ; la non-linéarité est reproduite par création mixte ; les six pertes sont orthogonales et ne sont pas additionnées deux fois ; le terme terminal demeure jusqu’au télescopage ; et la borne uniforme précède l’argument classique de compacité.\par
Aubin–Lions, la pression, Grönwall et le recollement ne sont pas la source de la régularité. Ils constituent la reconnaissance classique finale d’une régularité dont la structure a déjà empêché la perte.\par
\Needspace{6\baselineskip}
\subsection{3e résolution : Yang–Mills et le saut du secteur de courbure}
Le problème de Yang–Mills exige deux choses indissociables : construire une théorie quantique non triviale sur \(\mathbb R^4\) pour tout groupe compact, connexe et simple du domaine considéré, et démontrer que son Hamiltonien possède un saut de masse positif.\par
Il ne suffit pas de présenter un opérateur fini à spectre positif. Il ne suffit pas non plus d’écrire une action de Wilson, une mesure produit ou une observable ayant une valeur propre non nulle. Le saut doit appartenir à la même théorie de jauge qui possède localité, positivité par réflexion, covariance euclidienne, reconstruction de Wightman et courbure physique.\par
La caractéristique de jauge du continu conserve déjà le réseau raffinable, la mesure, les quatre réflexions, l’holonomie, la couleur, l’incidence, la retenue de fusion, le parcours, la mémoire, le terme terminal et le lecteur de courbure. Le groupe \(G\) et le couplage \(g\) interviennent ensuite, par l’entrelaceur spécifique. Ils ne sélectionnent pas rétrospectivement la généalogie.\par
\Needspace{6\baselineskip}
\subsubsection{Inéquivalence entre la fibre simple et la dynamique produit}
La première distinction dans l’argumentaire sépare deux générateurs qu’une ancienne notation pouvait confondre.\par
Le semi-groupe d’une seule fibre ne possède que les niveaux \(0\) et \(3\kappa_{\rm av}\). Le générateur produit complet intègre en revanche toutes les coordonnées indépendantes de l’état :\par
\[
D_m^{\rm prod}
=
3\kappa_{\rm av}
\sum_{\iota\in\mathcal I_m}(I-E_{m,\iota}),
\]
où les espérances \(E_{m,\iota}\) sont orthogonales et commutent.\par
Son spectre contient\par
\[
0,\quad
3\kappa_{\rm av},\quad
6\kappa_{\rm av},\quad
9\kappa_{\rm av},\ldots
\]
selon le nombre de coordonnées non constantes. Remplacer ce générateur par celui d’une fibre simple effacerait les multiplicités de Hoeffding et serait immédiatement falsifiable : un vecteur dépendant de deux facteurs doit avoir l’énergie \(6\kappa_{\rm av}\), et non \(3\kappa_{\rm av}\).\par
Le saut sera démontré pour le générateur produit, puis transporté vers le secteur physique. Il ne s’obtient pas en confondant deux dynamiques différentes.\par
\Needspace{6\baselineskip}
\subsubsection{Absence de coordonnées cachées et exclusion des modes nuls parasites}
L’état complet se factorise au moyen d’un indice exhaustif. Y apparaissent une seule fois :\par
\begin{itemize}
\item repères ;
\item liens ;
\item incidences ;
\item marques ;
\item ponts de raffinement ;
\item innovations nouvelles de chaque niveau.
\end{itemize}
Les sous-queues infinies sont enregistrées au moyen de leurs cylindres scalaires. Aucune coordonnée généalogique ne reste hors de la famille des espérances.\par
La décomposition de Hoeffding écrit l’espace comme somme orthogonale de secteurs \(\mathcal H_{m,S}\). Sur chacun d’eux,\par
\[
D_m^{\rm prod}u_S
=
3\kappa_{\rm av}|S|u_S.
\]
Le secteur \(S=\varnothing\) est l’intersection des images de toutes les espérances et se compose exactement des constantes :\par
\[
\bigcap_{\iota\in\mathcal I_m}
\operatorname{Ran}E_{m,\iota}
=
\mathbb C\Omega_m.
\]
Le vide est donc simple et tout vecteur orthogonal au vide a une énergie au moins égale à \(3\kappa_{\rm av}\).\par
Cela exclut l’objection selon laquelle le saut proviendrait de l’omission d’une coordonnée dans l’énumération. S’il manquait un facteur, il existerait des fonctions non constantes invisibles pour toutes les espérances et des modes nuls supplémentaires apparaîtraient. L’exhaustivité de l’indice empêche précisément cet échec.\par
\Needspace{6\baselineskip}
\subsubsection{Réalisation observable du saut spectral}
Une borne inférieure ne suffit pas à fixer la valeur exacte du saut. Un témoin atteignant le seuil est nécessaire.\par
Dans chaque collier d’incidence, la variable\par
\[
W_c(q)
=
\mathbf1_{\{q_1+\cdots+q_6=0\}}
-\frac13
\]
a une moyenne nulle et satisfait\par
\[
D_m^{\rm prod}W_c
=
3\kappa_{\rm av}W_c.
\]
Mais il reste à démontrer que \(W_c\) appartient au secteur de jauge physique.\par
L’action d’incidence s’écrit au moyen de la matrice\par
\[
(B^+x)_{ij}=x_i+x_j.
\]
Chaque composante de \(x\in\mathbb F_3^4\) apparaît trois fois dans la somme des six incidences. Dans \(\mathbb F_3\),\par
\[
\sum_{i<j}(B^+x)_{ij}
=
3\sum_i x_i
=
0.
\]
La condition définissant \(W_c\) est donc invariante de jauge. La moyenne sur le groupe de jauge complet conserve \(W_c\) ; elle ne l’exclut pas de l’espace physique. Ainsi, \(W_c\) est un vecteur propre physique et la valeur \(3\kappa_{\rm av}\) n’est pas seulement une borne : c’est le premier niveau spectral atteint.\par
\Needspace{6\baselineskip}
\subsubsection{Construction de la mesure préalable à l’estimation du saut}
Une arête grossière se raffine en neuf incréments ordonnés. La mesure fine n’est pas inventée en recopiant neuf fois la mesure grossière. Elle est construite au moyen du pont de chaleur conditionné par le produit des neuf incréments. Les temps des sous-segments peuvent être distincts et satisfont seulement à la condition que leur somme soit le temps grossier.\par
La convolution du noyau de chaleur garantit que le pont est une probabilité. Une carte triangulaire sépare alors la variable grossière des innovations nouvelles. La projection vers le niveau précédent multiplie les neuf incréments et oublie exclusivement les innovations.\par
De cette manière,\par
\[
(\widehat\pi_{m+1,m})_*
\widehat\mu_{m+1}
=
\widehat\mu_m.
\]
Les arêtes partagées par plusieurs faces demeurent une seule et même variable. Elles ne sont pas remplacées par des copies indépendantes. La retenue et la mémoire empêchent de recompter au niveau fin la face déjà présente au niveau grossier.\par
Sur cette famille projective, on construit d’abord le lecteur exact de \(F^2\), sa racine positive, l’action locale et la densité de Gibbs dépendant du couplage :\par
\[
d\widehat\mu_{m,g}^{\rm YM}
=
Z_{m,g}^{-1}
e^{-S_{m,g}}\,
d\widehat\mu_m^0.
\]
Les équations de Schwinger–Dyson et de Ward sont calculées pour cette mesure. Le saut n’est utilisé ni pour choisir \(g\), ni pour définir \(S_{m,g}\), ni pour repondérer la loi. Il apparaît ensuite, lors de l’étude du générateur réversible associé à la même mesure.\par
C’est la protection contre la circularité action–loi–saut : on construit d’abord l’action et sa mesure ; on démontre ensuite son spectre.\par
\Needspace{6\baselineskip}
\subsubsection{Descente du générateur au quotient de jauge}
L’espace complet comprend liens, repères, incidences, couleur, orientation, parcours, mémoire et terme terminal. La marginale qui ne conserve que les liens est une ablation ultérieure. Elle ne peut servir à nier une observable située dans la fibre d’incidence.\par
La projection physique s’obtient par moyenne de l’action du groupe de jauge complet. Puisque le générateur intègre les coordonnées de manière équivariante,\par
\[
[\mathbb P_m^{\rm phys},D_m^{\rm prod}]=0.
\]
Le semi-groupe descend donc au quotient physique. Son noyau est défini d’abord sur l’état complet, puis seulement sur les classes de jauge. On ne prend pas un opérateur restreint à une fibre pour le déclarer noyau sur tout l’espace.\par
Le vide simple et le vecteur propre \(W_c\) subsistent après compression. Le saut fini appartient déjà à l’espace invariant de jauge.\par
\Needspace{6\baselineskip}
\subsubsection{Unicité du processus reconstruit à partir de \(4\) positivités par réflexion}
La théorie euclidienne doit satisfaire la positivité par réflexion dans les quatre directions. Quatre mesures indépendantes ne sont pas construites. Une unique mesure de faces, obtenue à partir du même noyau réversible, factorise chaque forme d’Osterwalder–Schrader :\par
\[
\int
\overline{F(\Theta_{\nu}y)}F(y)\,d\Omega
=
\|\mathcal B_\nu F\|^2
\ge0,
\qquad \nu=1,\ldots,4.
\]
La positivité n’est pas une hypothèse ajoutée. C’est une norme au carré produite par la propriété de Markov et la réversibilité.\par
Les marginales des faces opposées retrouvent le même semi-groupe. Les quatre reconstructions OS possèdent donc une unique origine dynamique. Les connecteurs entre demi-réseaux démontreront ensuite qu’elles reconstruisent le même espace, le même vide et le même Hamiltonien.\par
\Needspace{6\baselineskip}
\subsubsection{Persistance du saut spectral dans la limite continue}
Les applications de raffinement entrelacent les formes de Dirichlet, les semi-groupes, le vide et le témoin \(W_c\). Le passage au continu ne repose pas seulement sur l’observation que le saut a la même valeur à chaque niveau. Cette égalité, isolément, n’empêcherait pas l’apparition de nouveaux états d’énergie arbitrairement faible à la limite.\par
La construction contrôle les produits croisés au moyen d’entrelaceurs de blocs. Les coordonnées grossières sont réparties sur leurs descendantes avec la normalisation exacte. Les identités de Gram démontrent que les expressions régularisées forment des familles de Cauchy. Ainsi, la forme, le vide et le vecteur propre atteignent la limite au sein du même système inductif.\par
Se obtiene\par
\[
D_{\infty,g}^{\rm YM}
\succeq
3\kappa_{\rm av}(I-P_{\Omega_\infty}),
\]
et le témoin limite satisfait\par
\[
D_{\infty,g}^{\rm YM}W_c
=
3\kappa_{\rm av}W_c.
\]
\Needspace{5\baselineskip}
Por tanto,\par
\[
\Delta_{\rm YM}^{\rm HMT}
=
3\kappa_{\rm av}>0.
\]
L’inégalité et le témoin convergent ensemble. Le saut n’est pas déduit d’une simple semi-continuité inférieure ni extrapolé à partir d’un maillage isolé.\par
\Needspace{6\baselineskip}
\subsubsection{Appartenance du témoin spectral au secteur de courbure}
Une objection subsisterait si \(W_c\) n’était qu’une excitation combinatoire d’un facteur interne, découplée de la courbure de Yang–Mills.\par
L’holonomie d’incidence élimine cette possibilité. La somme des six incidences détermine une holonomie locale d’ordre trois ; \(W_c\) est une fonction centrale de cette holonomie. Il appartient donc à l’algèbre locale invariante de jauge engendrée par les lacets.\par
Le lecteur \(F^2\) est polarisé pour produire un espace de courbures locales. Les entrelaceurs de raffinement transportent ces courbures en conservant le produit de \(\mathfrak g\). Les sommes cellulaires convergent vers une distribution covariante de jauge\par
\[
\mathbf F_{\mu\nu}
\in
\mathcal S'(\mathbb R^4;\mathfrak g)
\widehat\otimes
L^2(\widehat X_\infty,\widehat\mu_{\infty,g}^{\rm YM}).
\]
Le secteur cyclique engendré par \(\mathbf F\) est réducteur pour le Hamiltonien et contient le témoin \(W_c\). Le saut exact appartient donc au secteur physique de courbure. Ce n’est pas le spectre d’un bruit auxiliaire.\par
\Needspace{6\baselineskip}
\subsubsection{Localité euclidienne \(E(4)\) et reconstruction continue d’Osterwalder–Schrader}
Le générateur se décompose en circuits locaux de blocs à support uniformément borné. Une coloration nonadique sépare les colliers qui peuvent être mis à jour simultanément. Le produit des circuits préserve l’identité de Bianchi à chaque étape, et non seulement à la fin d’un balayage.\par
Les translations et rotations euclidiennes agissent par réétiquetage du réseau, des repères, des incidences et des holonomies. La continuité forte n’est pas prouvée uniquement pour un premier chaos abstrait. Elle est contrôlée d’abord sur les observables locales, puis sur les holonomies, les plaquettes et le trèfle ; la régularisation des orbites étend l’action à un domaine essentiel dense.\par
Le système cofinal produit un réseau local \(\mathfrak A_{\rm YM}(O)\), un vide \(\Omega_{\infty,g}\), une représentation forte de \(E(4)\) et des fonctionnelles de Schwinger. Les empilements de Markov et les transferts OS sont identifiés au même semi-groupe continu. La positivité par réflexion, la régularité, la covariance et la propriété de regroupement exponentiel permettent la reconstruction d’Osterwalder–Schrader.\par
Le prolongement produit des champs locaux de Wightman et une représentation de Poincaré dont le spectre est contenu dans le cône futur. Les quatre reconstructions euclidiennes coinduisent le même Hamiltonien de courbure.\par
Le saut contrôle en outre le regroupement :\par
\[
|S^{\rm YM}(A\,\tau_rB)|
\le
e^{-3\kappa_{\rm av}r}
\|A\Omega_\infty\|
\|B\Omega_\infty\|.
\]
Ainsi, la valeur spectrale positive a une conséquence physique continue et locale ; elle ne reste pas enfermée dans le calcul fini.\par
\Needspace{6\baselineskip}
\subsubsection{Reconstruction de l’action classique comme sortie du système euclidien}
La densité finie directrice est le lecteur exact de \(F^2\). Le trèfle ne s’y substitue pas. Il sert de carte lisse pour reconnaître sa limite.\par
Pour une connexion lisse, le développement de Magnus de l’holonomie d’une plaquette commence par\par
\[
\log\operatorname{Hol}_{p_{\mu\nu}}
=
a_m^2F_{\mu\nu}
+\text{terme impair}
+O(a_m^4).
\]
La moyenne sur les quatre orientations du trèfle annule le terme impair. La somme de Riemann donne\par
\[
S_{m,g}(A)
\longrightarrow
\frac1{4g^2}
\int_{\mathbb R^4}\|F_A(x)\|^2\,d^4x.
\]
La loi construite est donc reconnue comme Yang–Mills parce que son action finie, ses holonomies et sa courbure convergent vers la fonctionnelle classique correcte. On ne définit pas une « théorie de Yang–Mills » comme une théorie quelconque possédant le saut souhaité.\par
La même mesure satisfait Schwinger–Dyson et Ward. La relation entre petits lacets et courbure, l’identité de Bianchi et la covariance de jauge sont conservées à la limite. Ce n’est qu’après toute cette identification que l’équivalence spectrale transportant le saut est appliquée.\par
\Needspace{6\baselineskip}
\subsubsection{Interprétation dimensionnelle du paramètre de masse obtenu par l’entrelaceur terminal}
Le saut énergétique s’exprime dimensionnellement comme\par
\[
m_{\rm gap}^{\rm HMT}
=
\frac{\hbar_{\rm int}}{c_{\rm int}}
\,3\kappa_{\rm av}>0.
\]
Ce résultat n’attribue pas une masse de Proca à un gluon élémentaire. L’invariance de jauge demeure intacte. La masse correspond au coût minimal d’excitation du secteur collectif de courbure invariant de jauge au-dessus du vide.\par
Cette distinction ne diminue pas la portée du résultat : c’est exactement la formulation requise par le problème de Yang–Mills. Le saut appartient au spectre physique des états invariants de jauge, et non à un terme de masse élémentaire ajouté au lagrangien.\par
\Needspace{6\baselineskip}
\subsubsection{Existence d’un saut spectral positif et reconstruction massive du secteur de Yang–Mills}
Pour tout groupe de Lie compact, connexe et simple, la caractéristique de jauge du continu, suivie de son entrelaceur spécifique, construit une théorie quantique non triviale de Yang--Mills sur \(\mathbb R^4\).\par
La théorie possède :\par
\begin{itemize}
\item une famille projective de mesures locales dépendant du couplage ;
\item une action exacte reconnue à la limite comme \(\frac1{4g^2}\int\|F\|^2\) ;
\item un quotient de jauge physique ;
\item la positivité par réflexion dans quatre directions ;
\item un réseau local et une action forte de \(E(4)\) ;
\item des fonctions de Schwinger et un prolongement de Wightman ;
\item un vide simple ;
\item les équations de Schwinger–Dyson, de Ward et de Bianchi ;
\item un secteur cyclique de courbure ;
\item le saut exact
  \[
  \Delta_{\rm YM}^{\rm HMT}=3\kappa_{\rm av}>0.
  \]
\end{itemize}
La conclusion résiste aux objections prévisibles parce qu’elle n’identifie pas le générateur d’une fibre au produit ; ne laisse aucune coordonnée hors de l’indice ; ne confond pas la marginale des liens avec le quotient physique ; n’extrapole pas le saut à partir d’un maillage ; ne place pas le témoin hors du secteur de courbure ; ne présuppose pas la continuité euclidienne ; et n’utilise pas le saut pour construire l’action ou la mesure.\par
Navier–Stokes démontre que le continu peut transporter une non-linéarité sans perdre la mémoire nécessaire pour empêcher une singularité. Yang–Mills démontre que le même continu peut transporter une liberté de jauge sans perdre la courbure nécessaire pour séparer le vide de la première excitation physique.\par

\Needspace{6\baselineskip}
\section{Fermetures algébriques du continu et limites de la reconstruction extensionnelle}
Les cinq constructions \texttt{sp}, \texttt{sol}, \texttt{gau}, \texttt{coh} et \texttt{det} ne sont pas des branches ajoutées au continu pour résoudre cinq problèmes connus. Ce sont des caractéristiques simultanées du même état enrichi. Hodge et Birch–Swinnerton–Dyer sont les deux expressions finales de cette architecture : l’une transforme une histoire cohomologique compatible en classe algébrique ; l’autre démontre que les expressions analytique et arithmétique d’une même droite déterminantale coïncident.\par
\Needspace{6\baselineskip}
\subsection{4e résolution : Hodge et construction effective de la préimage algébrique}
Le problème de Hodge demande si toute classe rationnelle de type \((p,p)\),\par
\[
\alpha\in
H^{2p}(X,\mathbb Q)\cap H^{p,p}(X),
\]
pour une variété projective lisse complexe \(X\), provient d’une classe algébrique\par
\[
\beta_\alpha\in \operatorname{CH}^{p}(X)_{\mathbb Q}
\]
tal que\par
\[
\operatorname{cl}_{X,p}(\beta_\alpha)=\alpha.
\]
La difficulté ne consiste pas à écrire cette dernière égalité. Elle consiste à construire \(\beta_\alpha\) sans avoir préalablement introduit, sous un autre nom, l’existence de la préimage que l’on veut démontrer.\par
La caractéristique cohomologique \texttt{coh} fournit l’appareil antérieur à l’entrée de \(X\), \(p\) et \(\alpha\) : histoires survivantes, raffinements, relations, orientation, retenue, mémoire, frontière, parcours, multisection, terme terminal et lecteur. La donnée classique ne crée pas cet appareil ni ne sélectionne rétrospectivement un sous-domaine. Elle demande seulement ce qu’il exprime lorsqu’il reçoit les observations de \(\alpha\).\par
À chaque profondeur \(N\) apparaissent deux modules :\par
\[
\mathscr S_N(X,p),
\]
qui contient les états cohomologiques enrichis, et\par
\[
\mathscr H_N^p(X),
\]
qui contient les observations cohomologiques visibles à cette profondeur. Le lecteur fini\par
\[
e_N:\mathscr S_N(X,p)\longrightarrow\mathscr H_N^p(X)
\]
est coordonné aux applications de restriction :\par
\[
\rho_{N+1,N}:\mathscr S_{N+1}\longrightarrow\mathscr S_N,
\qquad
q_{N+1,N}:\mathscr H_{N+1}^p\longrightarrow\mathscr H_N^p,
\]
au moyen du carré\par
\[
e_N\rho_{N+1,N}
=
q_{N+1,N}e_{N+1}.
\]
Cela empêche qu’une observation admise à une profondeur soit incompatible avec son image au niveau précédent.\par
Pour la classe \(\alpha\), la fibre finie à peupler est\par
\[
\mathscr F_N(\alpha)
=
\{s\in\mathscr S_N(X,p):e_N(s)=q_N\alpha\}.
\]
La définir ne démontre pas qu’elle est non vide. Cela ne démontre pas non plus qu’un élément de \(\mathscr F_N(\alpha)\) peut se prolonger jusqu’à \(\mathscr F_{N+1}(\alpha)\). La résolution consiste précisément à construire ces deux propriétés.\par
\Needspace{6\baselineskip}
\subsubsection{Construction de la correction relative préalable à la classe de cohomologie}
Lors du raffinement de \(N\) à \(N+1\), deux noyaux relatifs apparaissent :\par
\[
\mathscr K_{N+1}^{\mathrm S}
=
\ker\rho_{N+1,N},
\qquad
\mathscr K_{N+1}^{\mathrm H}
=
\ker q_{N+1,N}.
\]
Le premier contient les corrections nouvelles de l’état qui disparaissent lorsqu’on revient au niveau précédent. Le second contient les observations nouvelles qui n’étaient pas non plus visibles au niveau précédent.\par
Les cartes nouvelles, les intersections de Mayer–Vietoris, les orientations, la retenue, la mémoire, la frontière, le parcours et les coordonnées terminales produisent une présentation finie. Sur ses générateurs sont construites deux applications :\par
\[
a_{N+1}:
\mathscr A_{N+1}^{\mathrm{rel}}
\longrightarrow
\mathscr K_{N+1}^{\mathrm S},
\]
\[
b_{N+1}:
\mathscr A_{N+1}^{\mathrm{rel}}
\longrightarrow
\mathscr K_{N+1}^{\mathrm H},
\]
satisfaisant\par
\[
e_{N+1}^{\mathrm{rel}}a_{N+1}=b_{N+1}.
\]
L’application \(a_{N+1}\) produit une correction de l’état ; \(b_{N+1}\) exprime exactement sa correction cohomologique.\par
Les observations relatives sont ordonnées selon la profondeur, le support, le feuillet, le parcours, la retenue, la mémoire et le terme terminal. Dans cet ordre, la matrice de \(b_{N+1}\) est unitriangulaire :\par
\[
b_{N+1}(c_\mu)
=
\bar c_\mu+
\sum_{\nu<\mu}
B_{\nu\mu}\bar c_\nu.
\]
La diagonale unité n’est pas choisie pour représenter une classe concrète. Elle provient de l’incidence principale de chaque générateur ; les termes inférieurs sont imposés par les relations, l’orientation et la retenue. Le tableau est construit sur toutes les observations relatives avant l’introduction de \(\alpha\).\par
L’élimination triangulaire produit explicitement un inverse à droite :\par
\[
\sigma_{N+1}^{\mathrm{rel}}(\bar c_{\mu_0})=c_{\mu_0},
\]
\[
\sigma_{N+1}^{\mathrm{rel}}(\bar c_\mu)
=
c_\mu-
\sum_{\nu<\mu}
B_{\nu\mu}
\sigma_{N+1}^{\mathrm{rel}}(\bar c_\nu),
\]
et, par récurrence,\par
\[
b_{N+1}\sigma_{N+1}^{\mathrm{rel}}=I.
\]
\Needspace{5\baselineskip}
Por tanto,\par
\[
s_{N+1}^{\mathrm{rel}}
=
a_{N+1}\sigma_{N+1}^{\mathrm{rel}}
\]
satisfait\par
\[
e_{N+1}^{\mathrm{rel}}
s_{N+1}^{\mathrm{rel}}
=I.
\]
C’est la protection contre la circularité. La section relative n’est pas supposée, ne se déduit pas du fait que la classe recherchée serait déjà algébrique et n’est pas construite spécifiquement pour \(\alpha\). Elle existe universellement pour toute observation relative.\par
Il n’est pas non plus nécessaire que la présentation soit injective. Le noyau peut contenir des expressions cellulaires redondantes. La preuve exige la surjectivité vers les observations relatives, et non l’unicité de l’expression. L’annulation du conoyau est une conséquence de l’inverse à droite ; ce n’est pas une condition directrice introduite avant la preuve.\par
\Needspace{6\baselineskip}
\subsubsection{L’opérateur universel d’extension}
La structure possède une dégénérescence unitaire\par
\[
\iota_{N+1,N}:\mathscr S_N\longrightarrow\mathscr S_{N+1},
\qquad
\rho_{N+1,N}\iota_{N+1,N}=I.
\]
Soit \(s\in\mathscr S_N\) et soit \(h\in\mathscr H_{N+1}^p\) une observation compatible :\par
\[
e_N(s)=q_{N+1,N}(h).
\]
Le défaut du prolongement trivial est\par
\[
\delta_{N+1}(s,h)
=
h-e_{N+1}\iota_{N+1,N}(s).
\]
La compatibilité garantit que\par
\[
\delta_{N+1}(s,h)
\in
\mathscr K_{N+1}^{\mathrm H}.
\]
On définit alors\par
\[
\operatorname{Ext}_{N+1,N}(s,h)
=
\iota_{N+1,N}(s)
+
s_{N+1}^{\mathrm{rel}}
\bigl(\delta_{N+1}(s,h)\bigr).
\]
Par construction :\par
\[
\rho_{N+1,N}
\operatorname{Ext}_{N+1,N}(s,h)
=s,
\]
\[
e_{N+1}
\operatorname{Ext}_{N+1,N}(s,h)
=h.
\]
La première identité conserve l’histoire antérieure. La seconde réalise exactement la nouvelle observation. La naturalité du tableau relatif garantit que l’orientation, la retenue, la mémoire, le parcours, la frontière et le terme terminal ne sont pas perdus lors du raffinement.\par
L’opérateur fonctionne pour tout couple compatible \((s,h)\). Il ne contient ni classe algébrique ni cycle choisi. C’est précisément pourquoi son application ultérieure à une classe de Hodge ne présuppose pas Hodge.\par
\Needspace{6\baselineskip}
\subsubsection{Construction niveau par niveau de l’histoire globale dans le système inverse}
À la profondeur initiale, il existe une section basée\par
\[
s_0^{\mathrm{base}}:
\mathscr H_0^p\longrightarrow\mathscr S_0.
\]
Ses images sont des états germes, non des cycles algébriques choisis pour représenter \(\alpha\). Ce n’est qu’ensuite que l’on fixe\par
\[
s_0=s_0^{\mathrm{base}}(q_0\alpha)
\]
et que l’on définit récursivement\par
\[
s_{N+1}
=
\operatorname{Ext}_{N+1,N}
\bigl(s_N,q_{N+1}\alpha\bigr).
\]
Les identités de l’opérateur universel prouvent par récurrence :\par
\[
e_Ns_N=q_N\alpha,
\qquad
\rho_{N+1,N}s_{N+1}=s_N.
\]
Chaque fibre \(\mathscr F_N(\alpha)\) est donc non vide et la famille\par
\[
\xi_\alpha=(s_N)_{N\ge0}
\]
est une histoire compatible :\par
\[
\xi_\alpha
\in
\varprojlim_N\mathscr F_N(\alpha).
\]
On n’invoque pas une limite inverse abstraite pour dissimuler l’absence de sections. La suite est présentée récursivement. L’annulation de \(\varprojlim^1\) n’est pas non plus utilisée en remplacement de la construction : la surjectivité et la condition de Mittag–Leffler apparaissent ensuite comme conséquences d’un opérateur d’extension déjà matérialisé.\par
\Needspace{6\baselineskip}
\subsubsection{Émergence de la classe algébrique comme limite de la construction relative}
Les observations finies et la coordonnée terminale séparent les expressions de la limite. Puisque \(\xi_\alpha\) reproduit \(q_N\alpha\) à toute profondeur,\par
\[
\operatorname{ev}_{\infty,X,p}^{\mathrm{Hdg}}
(\xi_\alpha)
=
\alpha.
\]
Ce n’est qu’alors que l’expression classique agit :\par
\[
X_\chi(X,p)
\xrightarrow{R_{\mathrm{Hdg}}}
K_0\!\left(
\operatorname{Perf}^{\ge p}(X)
\right)_{\mathbb Q}
\xrightarrow{\operatorname{ch}^{\mathrm{loc}}_p}
\operatorname{CH}^{p}(X)_{\mathbb Q}
\xrightarrow{\operatorname{cl}_{X,p}}
\operatorname{Hdg}^{p}(X)_{\mathbb Q}.
\]
Se define\par
\[
\beta_\alpha
=
\operatorname{ch}^{\mathrm{loc}}_p
\bigl(R_{\mathrm{Hdg}}(\xi_\alpha)\bigr).
\]
\Needspace{9\baselineskip}
L’identité d’action du lecteur,\par
\[
\operatorname{cl}_{X,p}
\operatorname{ch}^{\mathrm{loc}}_p
R_{\mathrm{Hdg}}
=
\operatorname{ev}_{\infty,X,p}^{\mathrm{Hdg}},
\]
da finalmente\par
\[
\operatorname{cl}_{X,p}(\beta_\alpha)
=
\alpha.
\]
Cette identité du lecteur ne prouve pas à elle seule l’existence de \(\xi_\alpha\). L’existence a été construite auparavant au moyen du tableau unitriangulaire, de l’inverse à droite, de l’opérateur universel et de la récurrence. Maintenir ces deux étapes séparées est précisément ce qui empêche la démonstration d’être circulaire.\par
Pour toute classe rationnelle de type \((p,p)\), la caractéristique cohomologique construit la section compatible \(\xi_\alpha\), le cycle \(\beta_\alpha\in\operatorname{CH}^p(X)_{\mathbb Q}\) et l’identité \(\operatorname{cl}_{X,p}(\beta_\alpha)=\alpha\), conformément à \cref{thm:pdf2-hodge-extension-universal,thm:pdf2-hodge-completitud-final}. L’exhaustivité cellulaire et FIN1 font partie de la construction démontrée de l’inverse à droite et de son prolongement coinductif ; ce ne sont pas des hypothèses externes de promotion.\par
\Needspace{6\baselineskip}
\subsection{5e résolution : BSD et coïncidence des expressions analytique et arithmétique}
La résolution de Birch–Swinnerton–Dyer ne commence pas par tenter de deviner l’ordre d’annulation de \(L(E,s)\), ni par réunir rétrospectivement rang, régulateur, groupe de Tate–Shafarevich et facteurs de Tamagawa.\par
La caractéristique déterminantale \texttt{det} conserve déjà histoire, produit fibré, unité commune, orientation, retenue, mémoire, survie, terme terminal, réseau, dualité, racine et expressions compatibles. Ce n’est qu’ensuite qu’intervient une courbe elliptique \(E/\mathbb Q\).\par
\Needspace{6\baselineskip}
\subsubsection{Complexe entier global et décomposition en fibres primaires}
La construction s’effectue sur un complexe libre entier d’histoires. Pour chaque nombre premier \(\ell\) et chaque puissance \(\ell^n\), les coefficients \(E[\ell^n]\), les transitions, Alexander–Whitney, le tiré en arrière de Manin, l’unité et les expressions sont obtenus par changement de base du même complexe entier.\par
La fibre ternaire ne gouverne pas les autres. On n’invente pas une fausse extension de \(\mathbb Z_3\) à \(\mathbb Z_\ell\). Cette précaution est indispensable pour que la conclusion englobe tout \(\Sha(E)\) et tous les facteurs de Tamagawa, et non uniquement sa composante \(3\)-primaire.\par
L’unité commune produit les expressions de Kato et de Poitou–Tate. Dans chaque bloc fini–singulier, on construit directement l’homotopie de remplissage\par
\[
h_*=-vf,
\]
dont le bord est la différence entre les deux expressions.\par
Les histoires sont ordonnées selon le nombre premier, le niveau, l’orientation, la retenue, la mémoire, la frontière, le terme terminal et la localisation. On obtient ainsi une décomposition basée exhaustive en :\par
\begin{itemize}
\item unité racine ;
\item blocs finis–singuliers ;
\item réseau local–global.
\end{itemize}
Aucune classe résiduelle ne demeure non typée.\par
L’homotopie \(H_{\mathrm{FS}}\) n’intervient pas comme prémisse. Dans chaque bloc normal, on définit son action et l’on vérifie\par
\[
\partial H_\lambda+H_\lambda\partial=I.
\]
La somme transportée par la décomposition basée satisfait\par
\[
\partial H_{\mathrm{FS}}
+
H_{\mathrm{FS}}\partial
=
p_{\mathrm{FS}}.
\]
Le membre de droite est le projecteur sur le secteur fini–singulier. Il n’efface pas le réseau qui conserve \(\Sha\) et Tamagawa. L’homotopie reproduit l’homotopie de remplissage préalablement construite ; elle n’est pas postulée pour faire disparaître une obstruction.\par
\Needspace{6\baselineskip}
\subsubsection{Finitude du groupe arithmétique préalable à l’évaluation du terme principal}
Le réseau résiduel possède une matrice entière \(D_E\). Avant de prendre son conoyau, on construit un inverse rationnel au moyen de la partie diagonale terminale et de la série géométrique finie de la partie strictement triangulaire. Les deux compositions sont vérifiées.\par
\Needspace{5\baselineskip}
Por tanto,\par
\[
D_E\otimes\mathbb Q
\]
est un isomorphisme et\par
\[
F_E=\operatorname{coker}D_E
\]
est fini. Cette finitude est obtenue avant toute utilisation du rang, de \(\Sha\), du régulateur ou du terme principal.\par
La forme de Smith conserve tous les diviseurs élémentaires et recompose exactement leurs composantes primaires. La finitude de \(\Sha\) n’est pas déduite de la formule finale : elle est dérivée ensuite de son inclusion dans un groupe fini déjà construit.\par
\Needspace{6\baselineskip}
\subsubsection{Bockstein, Kato–FERS et Poitou–Tate}
La différentielle déterminantale universelle \(S_E(T)\) engendre les pages de Bockstein. Sa forme de Smith \(T\)-adique détermine\par
\[
d
=
\operatorname{ord}_{T=0}
\det S_E(T)
=
\sum_q \dim V_q.
\]
Les pages ne sont pas remplacées par cet entier. Chaque quotient, chaque différentielle réduite, chaque accouplement et chaque jet sont conservés. Le coefficient régulier est\par
\[
R_{\mathrm{Boc}}^{\mathrm{der}}(S_E)
=
\left[
T^{-d}\det S_E(T)
\right]_{T=0}
\neq0.
\]
L’équivalence Kato–FERS est construite ligne par ligne. Alexander–Whitney coupe chaque histoire en tous ses points. La coupe terminale conserve le générateur de poids nul ; les autres coupes acquièrent des poids positifs déterminés par la profondeur et la retenue. Le tableau exhaustif prend la forme\par
\[
I+TB_i(T),
\]
avec un inverse convergent et un déterminant alterné égal à un dans \(T=0\). Ainsi, l’unité, les feuillets, la mémoire, l’orientation et le terme terminal sont préservés avant toute consultation d’une valeur de \(L(E,s)\).\par
Poitou–Tate relie ensuite les pages de Bockstein à Mordell–Weil. Les applications de Kummer, d’expression et de recollement sont écrites dans les deux sens et satisfont l’unité et la counité. Elles sont inverses avant la comparaison des dimensions.\par
Le dédoublement des jets conserve les \(q\) directions correspondant à une ligne de profondeur \(q\). On ne prétend pas obtenir \(q\) covecteurs à partir d’une seule coordonnée, ni compléter les dimensions par des identités artificielles. Le Gram conserve tous les termes croisés entre lignes, profondeurs et jets.\par
Ce n’est qu’ensuite que les dimensions sont prises :\par
\[
\operatorname{rank}E(\mathbb Q)
=
\sum_q\dim V_q
=
d.
\]
Le rang et l’ordre déterminantal coïncident parce que tous deux sont des expressions des pages construites, et non parce que l’un aurait été défini en recopiant l’autre.\par
\Needspace{6\baselineskip}
\subsubsection{Obtention du groupe de Tate–Shafarevich et des facteurs de Tamagawa par la suite exacte globale}
À partir du cône local réduit, on construit une injection\par
\[
\Sha(E)\hookrightarrow F_E.
\]
Les relèvements structurels des composantes locales produisent la projection et sa section. La vérification au niveau des cocycles donne :\par
\[
0
\longrightarrow
\Sha(E)
\longrightarrow
F_E
\longrightarrow
\bigoplus_v\Phi_v(\mathbb F_v)
\longrightarrow
0.
\]
Puisque \(F_E\) était déjà fini,\par
\[
|\Sha(E)|<\infty
\]
y\par
\[
|F_E|
=
|\Sha(E)|
\prod_v c_v.
\]
L’exactitude n’est pas déduite d’une comparaison des cardinaux. Les applications sont d’abord construites, leurs noyaux et leurs images identifiés ; ce n’est qu’ensuite que leurs ordres sont pris.\par
\Needspace{6\baselineskip}
\subsubsection{Construction indépendante du membre analytique de BSD}
La modularité de \(E\) fournit la forme nouvelle associée. Sa transformée de Mellin est divisée en deux intervalles, et l’équation fonctionnelle transforme la seconde moitié en une intégrale convergente avec toutes ses dérivées. Cela produit un germe holomorphe de \(L(E,s)\) autour de \(s=1\) sans supposer son ordre d’annulation.\par
Dans le demi-plan de convergence absolue, les facteurs d’Euler sont réalisés par des opérateurs normaux. Les troncatures convergent en norme trace et leur déterminant de Fredholm retrouve le produit d’Euler. La réduction de Schur sur les coordonnées terminales et un atlas de cartes déterminantales transportent cette famille jusqu’à la différentielle \(S_E\).\par
Il en résulte une factorisation construite :\par
\[
L(E,s)
=
u_E(s)
\det S_E(T_E(s)).
\]
\Needspace{7\baselineskip}
La carte satisfait\par
\[
T_E(1)=0,
\qquad
T'_E(1)=1,
\]
et l’unité de transition satisfait\par
\[
u_E(1)=1.
\]
Cette normalisation est démontrée au moyen des transitions de l’atlas, du complément de Fredholm et des changements de base. Elle n’est pas ajustée au moyen du rang, du régulateur, de \(\Sha\) ni du coefficient final.\par
En consecuencia,\par
\[
L(E,s)
=
(s-1)^dU_E^{\mathrm{lead}}(s),
\]
con\par
\[
U_E^{\mathrm{lead}}(1)
=
R_{\mathrm{Boc}}^{\mathrm{der}}(S_E)
\neq0.
\]
\Needspace{5\baselineskip}
Por tanto,\par
\[
\operatorname{ord}_{s=1}L(E,s)=d.
\]
En combinant cela avec Poitou–Tate :\par
\[
\operatorname{ord}_{s=1}L(E,s)
=
\operatorname{rank}E(\mathbb Q).
\]
\Needspace{6\baselineskip}
\subsubsection{Dérivation interne du régulateur à partir de l’accouplement des hauteurs}
La base entière de Mordell–Weil provient du système complet de jets. Sur elle sont construites la hauteur de Néron–Tate et sa décomposition locale–globale.\par
Le Gram contient tous les termes croisés :\par
\[
G_E^{\mathrm{jet}}
=
\bigl(
\langle P_{q,a},P_{q',a'}\rangle_{\mathrm{NT}}
\bigr).
\]
On ne suppose pas de diagonalité entre profondeurs et l’on ne remplace pas les directions supplémentaires par des blocs identité. Son déterminant définit\par
\[
\operatorname{Reg}(E)
=
\det G_E^{\mathrm{jet}}
>0.
\]
L’extension des scalaires à \(\mathbb R\) est effectuée avant la multiplication par le régulateur. Le comparateur entre les droites de Bockstein et de Néron est donc bien typé ; on ne tente pas d’introduire un nombre réel dans une droite encore définie uniquement sur \(\mathbb Q\).\par
\Needspace{6\baselineskip}
\subsubsection{Diagramme de \(4\) morphismes et évaluation du terme principal}
La comparaison finale utilise quatre flèches de droites déterminantales :\par
\begin{enumerate}
\item Kato–FERS transporte la section analytique.
\item Poitou–Tate la conduit jusqu’à Mordell–Weil et à sa hauteur.
\item La suite exacte sépare \(\Sha\) et Tamagawa.
\item L’évaluation intègre la torsion, la période de Néron et les composantes locales.
\end{enumerate}
Chaque flèche possède un domaine, un codomaine et une action sur une base. Avant l’introduction de l’élément analytique, deux trivialisations indépendantes sont fixées : l’analytique et celle de Néron.\par
La commutativité du carré final provient de la conservation des bases structurelles. Elle n’est pas imposée en égalisant préalablement les nombres qui doivent apparaître.\par
L’application de la composition à l’élément analytique donne :\par
\[
\operatorname{ord}_{s=1}L(E,s)
=
\operatorname{rank}E(\mathbb Q)
=
r,
\]
y\par
\[
\frac{L^{(r)}(E,1)}
{r!\,\Omega_E}
=
\frac{
|\Sha(E)|
\operatorname{Reg}(E)
\prod_v c_v
}{
|E(\mathbb Q)_{\mathrm{tors}}|^2
}.
\]
C’est la formule complète de Birch--Swinnerton--Dyer pour toute courbe elliptique \(E/\mathbb Q\), comme l’établit \cref{thm:pdf2-bsd-publicacion-final}. La famille compatible \(\nu_{E,\ell,n}\) est construite au sein du complexe entier global et fournit le comparateur qui relie ses fibres primaires ; elle fait partie de la démonstration universelle et ne constitue pas une condition externe de promotion.\par
Sa protection contre la circularité est strictement chronologique :\par
\begin{itemize}
\item le complexe entier est fixé avant la séparation des nombres premiers ;
\item l’homotopie de remplissage est construite avant \(H_{\mathrm{FS}}\) ;
\item Smith est obtenu avant le rang ;
\item Kato–FERS est vérifié avant le déterminant principal ;
\item Poitou–Tate est construit avant la comparaison des dimensions ;
\item l’exactitude \(\Sha\)–Tamagawa est prouvée avant de prendre les ordres ;
\item le prolongement Mellin–Fredholm est construit avant l’ordre analytique ;
\item la hauteur est construite avant le régulateur ;
\item les quatre flèches sont fixées avant l’introduction de l’élément analytique ;
\item les trivialisations sont fixées avant l’évaluation des deux membres.
\end{itemize}
Aucun élément n’est normalisé au moyen de la valeur qu’il doit ensuite démontrer.\par
\Needspace{6\baselineskip}
\subsection{Factorisation des entrelaceurs terminaux RH–NS–YM–Hodge–BSD à travers le continu généalogique commun}
Riemann, Navier–Stokes, Yang–Mills, Hodge et BSD ne sont pas cinq traductions arbitraires d’un même vocabulaire. Chaque problème apparaît lorsqu’une expression classique oublie une partie différente de l’état enrichi :\par
\begin{itemize}
\item Riemann oublie la colligation commune qui conserve simultanément les canaux premier, gamma et polaire.
\item Navier–Stokes oublie la colonne causale complète comprenant terme terminal, pertes, retenue et mémoire.
\item Yang–Mills oublie la coordination entre incidence, courbure, loi, dynamique et secteur physique.
\item Hodge conserve une classe cohomologique, mais pas l’histoire compatible qui produit sa préimage algébrique.
\item BSD sépare l’expression analytique de l’arithmétique et perd la droite déterminantale commune qui les transporte.
\end{itemize}
L’architecture HMT n’impose pas cinq solutions indépendantes. Elle conserve d’abord ce que les formulations classiques expriment séparément :\par
\[
\begin{aligned}
\text{état enrichi}
&\longrightarrow \text{raffinements compatibles}
\longrightarrow \text{survie et terminal}\\
&\longrightarrow \text{continu discret}
\longrightarrow \text{publications classiques}.
\end{aligned}
\]
La raison pour laquelle les solutions fonctionnent n’est pas une similitude verbale. Dans les cinq cas, le théorème décisif provient d’une identité de conservation matérielle :\par
\begin{itemize}
\item conservation de l’énergie ou du défaut ;
\item conservation de l’histoire sous raffinement ;
\item conservation de l’action et du produit surfacique ;
\item conservation de l’observation durant l’extension ;
\item conservation de l’unité et de la base dans les droites déterminantales.
\end{itemize}
C’est pourquoi les problèmes classiques doivent apparaître à la fin de l’organisation éditoriale. L’objet commun est construit en premier ; HMT–MD développe ensuite ses conséquences ; la Mécanique dimensionnelle réalise alors ses transducteurs ; ce n’est qu’en dernier lieu que les cinq conclusions sont lues dans leurs langages classiques.\par
\Needspace{6\baselineskip}
\subsection{Distinction typée entre les cinq fermetures terminales et le corridor logique ZFC--cycle}
La relation avec Zermelo--Fraenkel, ZFC et la construction de von Neumann doit être formulée avec précision et dans l’ordre causal du corpus. Les cinq résolutions terminales et le corridor logique ZFC--cycle sont des démonstrations distinctes, fondées sur la même structure discrète du continu : la portée locale de l’une ne peut servir de négation transversale de l’autre. Les cinq fermetures établissent leurs résultats classiques respectifs ; le corridor logique spécifique conserve en outre la chaîne\par
\[
\begin{aligned}
 \mathcal C_{\mathrm{cont}}^{\mathrm{disc}}
 &\longrightarrow\mathcal L_0\xrightarrow{\Phi}\mathcal L_1
 \longrightarrow\gamma_F\longrightarrow w(\gamma_F)\\
 &\longrightarrow\mathsf{NoEscape}_{24}
 \longrightarrow\mathsf{AntiDiag}_{369}\\
 &\longrightarrow\mathsf{G\ddot{o}delIII}_{\mathrm{HMT}}
 \longrightarrow\mathsf{CycleCriterion}_{\mathrm{ZFC}},
\end{aligned}
\]
et exprime le critère fort retrouvé\par
\[
 \boxed{
 \neg\operatorname{Consis}(\mathrm{ZFC})
 \Longleftrightarrow
 \exists C\;\bigl(C\text{ est un cycle hors bande et }|C|>24\bigr).}
\]
Cette équivalence possède son propre domaine, ses encodeurs et sa preuve ; elle ne se déduit pas des cinq théorèmes terminaux et n’est pas limitée par une observation locale concernant l’un d’eux. De même, la résolution cardinale déjà construite conserve simultanément\par
\[
 \mathsf{CH}_{\mathrm{HMT}},
 \qquad
 M_h=L[h]\models 2^{\aleph_0}=\aleph_1,
\]
sans transformer la droite réelle ou l’énoncé extensionnel en entrées du générateur.\par

Une représentation extensionnelle formulée dans le langage de ZFC peut coder les supports ensemblistes des objets utilisés :\par
\begin{itemize}
\item un état fini est un ensemble ;
\item une histoire est une suite ;
\item un arbre est un ensemble de mots ;
\item une transition est une fonction ou une relation ;
\item un système inverse est une famille ensembliste ;
\item sa limite est un sous-ensemble d’un produit.
\end{itemize}
Ce codage ultérieur ne contient pas à lui seul l’identité généalogique de l’objet et ne gouverne pas les théorèmes construits avant l’expression extensionnelle. L’insuffisance apparaît lorsque le réduit extensionnel est traité comme s’il épuisait le parcours, l’orientation, la retenue, la mémoire, la frontière et le terme terminal.\par
L’axiome d’extensionnalité affirme :\par
\[
\forall A\,\forall B
\left[
\forall x\,
(x\in A\leftrightarrow x\in B)
\Rightarrow A=B
\right].
\]
L’assertion est correcte pour les ensembles. Deux ensembles ayant les mêmes éléments sont le même ensemble.\par
Mais il n’en découle pas qu’une expression extensionnelle permette de reconstruire l’histoire oubliée lors de sa production.\par
Sea\par
\[
\operatorname{Pub}:
\mathscr X_{\mathrm{gen}}
\longrightarrow
\mathscr X_{\mathrm{ext}}
\]
\Needspace{6\baselineskip}
un lecteur qui envoie un état généalogique sur sa valeur extensionnelle. Il peut arriver que\par
\[
\operatorname{Pub}(x)
=
\operatorname{Pub}(y)
\]
tandis que\par
\[
x\neq y
\]
en tant qu’histoires, parce qu’elles possèdent des parcours, orientations, retenues, mémoires, frontières ou termes terminaux différents.\par
Le réduit identifie les deux images dans \(\mathscr X_{\mathrm{ext}}\), mais cette égalité dans le codomaine n’implique pas \(x=y\) dans le domaine enrichi. L’erreur de type consiste à transformer l’égalité des expressions en égalité généalogique.\par
On peut définir la relation\par
\[
x\sim_{\mathrm{ext}}y
\quad\Longleftrightarrow\quad
\operatorname{Pub}(x)
=
\operatorname{Pub}(y).
\]
Alors, dans le domaine correspondant,\par
\[
\mathscr X_{\mathrm{ext}}
\simeq
\mathscr X_{\mathrm{gen}}/
{\sim_{\mathrm{ext}}}.
\]
La fibra\par
\[
\operatorname{Pub}^{-1}(a)
\]
contient les généalogies qui produisent la même image \(a\). L’extensionnalité décrit correctement le quotient ; elle ne fournit pas une section canonique reconstruisant une histoire concrète de chaque fibre.\par
Même l’axiome du choix ne résout pas cette carence. Le choix peut sélectionner un représentant de chaque fibre sous ses hypothèses, mais :\par
\begin{itemize}
\item il ne le rend pas canonique ;
\item il ne conserve pas nécessairement les opérations ;
\item il ne garantit pas la naturalité sous raffinement ;
\item il ne reconstruit ni retenue, ni mémoire, ni terme terminal ;
\item il ne démontre pas que la sélection correspond à la généalogie physique ou mathématique correcte.
\end{itemize}
Choisir un représentant n’équivaut pas à retrouver sa provenance.\par
\Needspace{6\baselineskip}
\subsubsection{L’exemple des ordinaux de von Neumann}
La construction de von Neumann\par
\[
0=\varnothing,
\qquad
n+1=n\cup\{n\}
\]
est exacte. L’ordinal\par
\[
n=\{0,1,\ldots,n-1\}
\]
code parfaitement son ordre extensionnel.\par
\Needspace{6\baselineskip}
Mais le même nombre peut apparaître comme sortie d’histoires différentes :\par
\[
6=3+3=2\cdot3=7-1.
\]
En tant qu’ordinal de von Neumann, le résultat est le même ensemble. C’est correct : l’arithmétique exige que les quatre expressions produisent le même \(6\).\par
Ce que l’ordinal seul ne conserve pas, c’est l’opération, le parcours, la retenue ou l’orientation qui a produit cette occurrence concrète. Si cette provenance importe pour une application ultérieure, elle doit être maintenue comme partie de l’objet enrichi :\par
\[
\widehat n_x
=
\{(k,x_k):k<n\}.
\]
La projection\par
\[
(k,x_k)\longmapsto k
\]
retrouve exactement l’ordinal de von Neumann. HMT n’élimine pas l’ordinal ; elle le reconnaît comme expression d’une structure susceptible de contenir davantage d’information.\par
La hiérarchie cumulative \(V_\alpha\) conserve l’appartenance et le rang. Elle n’enregistre pas automatiquement une causalité interne de construction différente du codage choisi. Réintroduire des étiquettes, des arbres ou des catégories permet de transporter des données généalogiques supplémentaires, mais cet enrichissement confirme précisément que la généalogie n’était pas contenue dans la valeur extensionnelle après avoir été oubliée.\par
\Needspace{6\baselineskip}
\subsubsection{La droite réelle comme expression ultérieure}
La même distinction concerne le continu. Un nombre réel peut être construit comme coupure de Dedekind, classe de suites de Cauchy ou autre représentant extensionnel équivalent. Cette équivalence est mathématiquement correcte.\par
HMT ajoute une question préalable : quel système compatible d’états, de raffinements et de termes terminaux exprime ce réel ?\par
La recta\par
\[
\mathbb R
\]
apparaît ensuite par une évaluation archimédienne. Elle conserve la valeur exprimée, mais pas nécessairement l’histoire complète qui l’a produite. Deux états peuvent partager la même évaluation réelle sans partager la même généalogie.\par
La structure discrète du continu ne doit donc pas introduire \(\mathbb R\) comme entrée. L’histoire compatible est construite en premier ; son image archimédienne est exprimée ensuite.\par
L’entrelaceur avec la cardinalité classique est construit explicitement. Le codage par blocs fournit l’homéomorphisme \(\beta:\Omega_\Gamma^{\rm cal}\simeq(\mathbb F_3^6)^{\mathbb N}\) ; ses approximants rationnels construisent \(\mathbb R_{\rm HMT}\), et le quotient par égalité de valeur exprime \(\mathbb R\). Le code généalogique \(h\) induit le bon ordre et détermine le modèle canonique \(M_h=L[h]\), dans lequel la construction prouve l’énoncé classique \(2^{\aleph_0}=\aleph_1\). Les résultats d’indépendance de Gödel et de Cohen classifient le langage extensionnel de ZFC après oubli de \(h\) ; ils ne réfutent pas la décision produite par la structure enrichie qui conserve le sélecteur.\par
\Needspace{6\baselineskip}
\subsubsection{Distinction entre identité généalogique, équivalence observationnelle et égalité extensionnelle}
La carence est typologique : un codage ensembliste ajouté peut transporter une généalogie déjà construite, tandis que le réduit extensionnel nu est dépourvu de l’application qui a produit l’expression. Capacité de codage et suffisance fondationnelle sont des assertions différentes.\par
Trois notions doivent être distinguées :\par
\[
\text{identité généalogique},
\qquad
\text{équivalence observationnelle},
\qquad
\text{égalité extensionnelle}.
\]
La première exige l’égalité de l’histoire, du parcours, de la mémoire, de la retenue et du terme terminal. La deuxième exige l’égalité relativement à un lecteur déterminé. La troisième exige l’égalité de l’objet exprimé dans le codomaine ensembliste.\par
Les confondre engendre de faux problèmes de reconstruction. Une image peut être exacte et, en même temps, insuffisante pour retrouver l’objet qui l’a produite.\par
C’est l’apport fondationnel que les cinq résolutions rendent visible. Chaque problème classique travaillait avec une expression correcte, mais appauvrie :\par
\begin{itemize}
\item une forme explicite sans toute sa généalogie commune ;
\item un champ sans mémoire causale complète ;
\item une théorie de jauge sans coordination intégrale de son processus ;
\item une classe cohomologique sans son histoire d’extension ;
\item une fonction \(L\) séparée de sa droite arithmétique.
\end{itemize}
La structure HMT maintient ces corrélations avant d’en exprimer les images. Elle n’a donc pas à reconstruire a posteriori, dans chaque problème, une information qui n’a jamais été détruite.\par
La conclusion fondationnelle est la suivante :\par
\begin{quote}
La fermeture extensionnelle est une expression ultérieure de la fermeture généalogique et ne peut s’y substituer. HMT construit la structure enrichie avant la droite réelle, démontre \(\mathsf{CH}_{\mathrm{HMT}}\), réalise la fermeture ordinale dans \(L[h]\) et établit séparément le corridor Gödel--HMT III et le critère ZFC--cycle. Un codage extensionnel ultérieur ne conserve que ce que déclare le morphisme d’expression ; il ne révoque aucun de ces résultats.\par
\end{quote}
\Needspace{6\baselineskip}
\subsection{Factorisation de RH–NS–YM–Hodge–BSD comme expressions terminales du continu généalogique commun}
Les cinq problèmes apparaissent ainsi comme cinq lectures terminales d’un unique bord. Ils ne sont pas résolus parce qu’une analogie les réunirait, mais parce que le continu enrichi conserve dès l’origine les corrélations que chaque formulation classique avait séparées.\par
La dépendance causale définitive se factorise comme suit :\par
\[
\begin{aligned}
\mathrm{APP}
&\longrightarrow \mathrm{TRIT}
\longrightarrow \mathrm{TPK}
\longrightarrow \widehat x\\
&\longrightarrow \text{cinq caractéristiques consubstantielles}
\longrightarrow \text{terminal et limite}
\longrightarrow \mathfrak C_{\mathrm{cont}}^{\mathrm{disc}}
\end{aligned}
\]
et seulement ensuite :\par
\[
\mathfrak C_{\mathrm{cont}}^{\mathrm{disc}}
\longrightarrow
\begin{cases}
\text{positivité de Weil et RH},\\
\text{régularité globale lisse de Navier--Stokes pour toute donnée périodique lisse},\\
\text{Yang--Mills et saut de masse pour tout groupe compact, connexe et simple},\\
\text{algébricité de toute classe rationnelle de Hodge},\\
\text{égalité complète de BSD pour toute courbe elliptique sur }\mathbb Q.
\end{cases}
\]
Les cinq lignes finales ne sont pas cinq branches internes du continu. Ce sont cinq expressions classiques ultérieures de cinq capacités indissociables qui coexistaient déjà dans le même état.\par
Telle est la fermeture que la formulation précédente ne pouvait pas encore formuler : il ne s’agit pas de cinq problèmes déconnectés qui admettraient par hasard un même vocabulaire, mais de cinq pertes d’information distinctes corrigées par une seule architecture qui a refusé de la détruire.\par

% Empalme y operador conjunto congelados: se consumen por su fuente y huella
% originales, no por una paráfrasis posterior.
\begingroup
\renewcommand{\chapter}[2][]{}%
\chapter{Raccordement matériel entre la structure unique et ses cinq résolutions}
\label{chap:pdf2-generacion-simultanea-hmt}

\begingroup
\footnotesize
% CORTE_2_NO_DISGREGACION — bloque de empalme PDF1 -> PDF2
% Congelado después del PDF1 compilado el 2026-08-19.
% Owners:
% 04b  8ba02f246c18a3bd2ff83ffdb08874c816a7f59c688b6fab6af6eb2f9647071f
% 04b1 25d6dee3875cbdd3252f793212448c17c4b5b2d3c5d787e995598d18fa76d5cf
% 04b2 d5de9b5379d4f95ea3d35e63c9ccc2ab741b5c750c01409af6f13526073b874f
% 04b3 ac9ca321b610e2c0560a32ec6ae1555814c67c587e152096d0c48d6abaa0bd4b

\section*{Raccordement mathématique concret PDF1--PDF2}

L'entrée unique est l'état APP--TRIT--TPK complet. Avant toute évaluation archimédienne sont parcourus
\[
 0\longrightarrow10,\qquad 0\longrightarrow1000,
 \qquad 0\longrightarrow\infty,
\]
avec les neuf espaces internes \(1,\ldots,9\) ; deux positions produisent APP, TRIT oriente ses feuillets et TPK prolonge extension, phase, retenue, mémoire, frontière, parcours et survie. Ni le continu ni \(\mathbb R\) ne sont des entrées.

Pour une émission survivante
\[
 \mathfrak e_\alpha=
 (\alpha,\operatorname{Ext}_{\gamma_\alpha},\operatorname{or}(\alpha),
 \kappa(\alpha),\partial\alpha,\gamma_\alpha,
 m(\gamma_\alpha;m_0),\operatorname{Led}(\gamma_\alpha))
\]
les cinq caractéristiques sont des actions simultanées sur ce même générateur :
\begin{align*}
 \mathcal O_{\rm sp}(\mathfrak e_\alpha)
  &=(U_\alpha,\sigma_x,\sigma_y,A_\alpha,\eta_\alpha),\\
 A_\alpha&=P_yU_\alpha P_x+Q_yU_\alpha Q_x,
 &\eta_\alpha&=Q_yU_\alpha P_x+P_yU_\alpha Q_x,\\
 \mathcal O_{\rm sol}(\mathfrak e_\alpha)
  &=(b_\alpha,b_\alpha^+,b_\alpha^-;
  M^\pm\mapsto M^\pm+b_\alpha^\pm,
  D\mapsto D+b_\alpha,H\mapsto H+|b_\alpha|),\\
 \mathcal O_{\rm gau}(\mathfrak e_\alpha)
  &=\{(\operatorname{Cell}_{\rm TPK}(\alpha),q,B,W_c(q)):
      q\in Q_{\rm inc}(t\alpha)\},\\
 \mathcal O_{\rm coh}(\mathfrak e_\alpha)
  &=(I_y,J_y,\kappa_{y,N},\delta_y,C_{y,N},
      F_\alpha:C_{y,N}\to C_{x,N}),\\
 \mathcal O_{\rm det}(\mathfrak e_\alpha)
  &=(\Psi_\alpha,K_\alpha,\mathbf1_\alpha,z_\pm(\alpha),h_*(\alpha),
  \operatorname{Det}C_{y,N},H_*C_{y,N},\operatorname{Smith}C_{y,N},
  \operatorname{Boc}C_{y,N}).
\end{align*}
Ces expressions sont des vues dépendantes d'une seule action \(\mathcal O(\mathfrak e_\alpha)\) ; elles ne définissent pas cinq domaines, facteurs ou branches. Elles partagent littéralement source, cible, parcours, registre, orientation, retenue, mémoire et frontière. Leurs identités de couplage sont
\[
 A_\alpha^*A_\alpha+\eta_\alpha^*\eta_\alpha=U_\alpha^*U_\alpha,
 \qquad E b_f^\pm=b_c^\pm+\kappa^{\rm can},
\]
\[
 Q_{\rm inc}=C^2(\operatorname{Cell}_{\rm TPK};\mathbf F_3),
 \qquad sB=0,\qquad W_c(q+Ba)=W_c(q),
\]
\[
 \partial h_*=z_--z_+,qquad
 K_{\beta\alpha}=K_\beta+\Psi_\beta K_\alpha.
\]

La monodromie est la connexion nonadique réelle :
\[
 \Gamma_{9,k}=\operatorname{Hol}_{C_9}(\nabla_{\rm TPK}),
 \qquad g(k+9)=g(k),
 \qquad q_{\rm mem}(k+9)=q_{\rm mem}(k)+1,
\]
sans retour d'état. Pour une histoire \(h=(\varepsilon_{k+1},\ldots,\varepsilon_N)\), les treize coordonnées sont engendrées par le pliage
\[
 r_{k,N}(h)=
 \operatorname{Upd}_{\varepsilon_N}\circ\cdots\circ
 \operatorname{Upd}_{\varepsilon_{k+1}}(r^0_{k}),
\]
et le terminal renvoie la multisection entière, sans sélecteur :
\[
 \mathfrak T_{k,N}(\widehat x_k)
 =\{r_{k,N}(h):h\in\mathscr H_{k,N}(\widehat x_k)\}.
\]
La naturalité est une égalité unique sur les préfixes engendrés,
\[
 u_{\ell,k}^{\mathcal O}\mathcal O_{\le\ell}(\gamma)
 =\mathcal O_{\le k}(\gamma)
 =\operatorname{map}(\mathcal O)
   \bigl(u_{\ell,k}^{\mathfrak e}E_{\le\ell}(\gamma)\bigr),
\]
qui passe au système inverse des terminaux. Ce n'est qu'alors que l'on obtient la structure discrète complète du continu ; l'évaluation dans \(\mathbb R\) est une expression ultérieure.

\subsection*{Carte d'utilisation dans PDF2}

\begin{center}
\small
\begin{tabularx}{\textwidth}{@{}p{.18\textwidth}X p{.30\textwidth}@{}}
\toprule
Symbole utilisé & Contenu structurel intégral de PDF1 &
Entrelaceur construit dans PDF2 \\
\midrule
\(\mathfrak G_{\rm sp}\) &
\(U,\sigma,P,Q,A,\eta\), feuillets, retenue, terminal, naturalité et lecteur signé interne. &
\(\Sigma_R^{\rm sp},J_\pm,K_R^{\rm sp},\Xi,P,B_W\) et l'identification avec la forme de Weil. \\
\(\mathfrak G_{\rm sol}\) &
\(b,b^\pm,\kappa^{\rm can}\), mémoire, composantes micro/coquille orthogonales, colonne terminale, Gram, Stein et télescopie sans double comptage. &
Application construite depuis la source \(J_T^{\rm sol}\), \(\Lambda_{M,T}\), \(I_{M,J}\), factorisation uniforme et borne portant uniquement sur les données. \\
\(\mathfrak G_{\rm gau}\) &
\(\operatorname{Cell}_{\rm TPK},Q_{\rm inc},B,W_c\), quatre directions, six cellules, huit faces, espérances et terminal commun. &
Interface de jauge physique, semi-groupe requis, secteur physique, Yang--Mills et saut de masse. \\
\(\mathfrak G_{\rm coh}\) &
Ports, \(\kappa,\delta,C\), construction bar/coend, retenue, survie et cône d'holonomie. &
\(R_{\rm Hdg}\), \(\mathrm{ch}_{\rm loc}\), classe algébrique et expression de Hodge. \\
\(\mathfrak G_{\rm det}\) &
\(\Psi,K,z_\pm,h_*\), AW, droite déterminant, homologie, Smith et Bockstein, conservés sans exiger l'acyclicité en entrée. &
Kato--FERS, PT, comparateur local--\newline global, Tamagawa et expression BSD. \\
\bottomrule
\end{tabularx}
\end{center}

Les objets de la troisième colonne sont propres à la cible : ils ne doivent pas être réintroduits rétrospectivement comme entrées de PDF1. PDF2 les construit sur l'objet structurel complet de la deuxième colonne et conserve ses treize coordonnées, son terminal et sa généalogie.

\endgroup

\section{Typage unitaire de l'objet consommé}

Le bloc figé précédent ne s'ajoute pas à un constructeur graphique à cinq
composantes. Il le remplace. Pour chaque profondeur \(k\), soit
\(\mathcal E_k\) l'état APP--TRIT--TPK enrichi de toutes ses émissions
survivantes, et soit \(\mathcal Z_k\) le codomaine de l'unique action
corrélative. La flèche directrice est
\begin{equation}\label{eq:pdf2-constructor-unico-operaciones}
 \mathcal O_k:\mathcal E_k\longrightarrow\mathcal Z_k,
 \qquad
 \widehat x_k\longmapsto
 \mathcal O\bigl(\{\mathfrak e_\alpha:
 \alpha\in\operatorname{Ch}_k(\widehat x_k)\}\bigr).
\end{equation}
Sa valeur conserve simultanément les coordonnées
\(\mathrm{sp},\mathrm{sol},\mathrm{gau},\mathrm{coh},\mathrm{det}\)
du même état et du même ledger. Ce n'est pas un tuple de facteurs
indépendants ni le graphe tautologique de cinq applications préalablement séparées.

Les troncatures de l'état et de l'action s'écrivent
\[
 u^{\mathfrak e}_{\ell,k}:\mathcal E_\ell\to\mathcal E_k,
 \qquad
 u^{\mathcal O}_{\ell,k}:\mathcal Z_\ell\to\mathcal Z_k,
 \qquad \ell\ge k,
\]
et satisfont une seule naturalité,
\begin{equation}\label{eq:pdf2-naturalidad-accion-unica}
 u^{\mathcal O}_{\ell,k}\circ\mathcal O_\ell
 =\mathcal O_k\circ u^{\mathfrak e}_{\ell,k}.
\end{equation}
L'action passe donc à la limite sans sélectionner une histoire :
\[
 \mathcal O_\infty:
 \mathcal C_{\rm cont}^{\rm disc}:=\varprojlim_k\mathcal E_k
 \longrightarrow
 \mathcal Z_\infty:=\varprojlim_k\mathcal Z_k.
\]
Les notations utilisées dans les cinq chapitres sont des vues intrinsèques de cette
unique valeur :
\begin{equation}\label{eq:pdf2-operacion-infinita-tipificada}
 \mathfrak G_T^{\rm int}
 :=\operatorname{View}_T\!\left(
 \mathcal O_\infty(\mathcal C_{\rm cont}^{\rm disc})\right),
 \qquad
 T\in\{\mathrm{sp},\mathrm{sol},\mathrm{gau},\mathrm{coh},\mathrm{det}\}.
\end{equation}
Ici, \(\operatorname{View}_T\) expose uniquement les coordonnées requises par une
application ; elle ne crée pas de branche, de facteur, de sous-domaine ni de sélecteur. La
structure de départ de chaque résolution est toujours la valeur complète de
\(\mathcal O_\infty\), avec les quatre autres caractéristiques présentes.

\begin{falsifier}[Ordre causal et non-disgrégation]
Le raccordement échoue si l'on utilise \(\mathbb R\) ou le continu comme entrée de
\(\mathcal O_k\), si l'on remplace \(\mathcal O_k\) par cinq constructeurs
indépendants, si l'on sélectionne une seule histoire terminale ou si une vue
\(\operatorname{View}_T\) est confondue avec l'objet complet. Dans n'importe lequel
de ces cas, la conclusion du chapitre d'application cesse de se transférer.
\end{falsifier}

\endgroup
\paragraph{Acquittement immédiat des frontières du raccord.}
Les marques d'interface que le raccord figé formule comme bornes ou
obligations de l'étape terminale ne sont ni des réserves du théorème final ni des données
ajoutées depuis les problèmes classiques. Elles sont acquittées, dans le même ordre
causal et avant tout chapitre consommateur, par les propriétaires exacts
qui suivent : ceux-ci construisent à partir de l'état commun les cinq publications,
leurs applications, leurs estimations et leurs certificats. Les formulations classiques
ultérieures reconnaissent ces sorties ; elles ne les sélectionnent ni ne les fournissent.
% Propietarios completos recuperados. Su orden es vinculante: primero el
% objeto correlativo; después las construcciones que descargan las cinco
% bisagras; sólo entonces las vistas comunes y las publicaciones clásicas.
\begin{HMTScientificOwnerSubordinated}
% 04 — CINCO REALIZACIONES MATEMATICAS FUNDAMENTALES
% ==================================================
% Residencia narrativa sucesora. Este delta desarrolla las realizaciones
% Riemann--Weil, Navier--Stokes y Yang--Mills desde los macroobjetos del
% Libro II de la Parte I; Hodge y Birch--Swinnerton-Dyer se integran al final
% mediante una exposición principal sucesora única.

\ifdefined\HMTSucesoraStructureTrace
  \typeout{HMT-SUCCESSEURE 04 : réalisations RH, Navier--Stokes et Yang--Mills}
\fi

La Partie II a construit, au sein de la structure discrète du continu, cinq
sous-structures associées à leurs canaux de publication respectifs. Cette partie reçoit
chaque sous-structure déjà
enrichie par feuillet, orientation, retenue, bord, voie, mémoire et
survie ; définit le lecteur spécialisé ; démontre que ce lecteur
entrelace le raffinement ; et applique, uniquement à la fin, le critère mathématique ou
physique qui reconnaît la sortie.

Chacun des cinq chapitres établit une clôture démontrée dans son
domaine à partir de la structure de départ explicite qu'il énumère. La
formulation conventionnelle intervient ensuite comme identification du
codomaine et comme confrontation, tandis que la chaîne démonstrative demeure
ordonnée par les objets construits dans les Parties I et II. Les contrôles
auxiliaires ultérieurs ne sont pas des arêtes de cette chaîne et ne peuvent
conditionner ni rétrograder la clôture propriétaire.

Dans les cinq théorèmes développés ci-après, l'ordre causal est
\begin{equation}\label{eq:realizaciones-orden-causal}
 \begin{gathered}
 \text{APP--TRIT--Topological Phase Kernel}
 \longrightarrow \text{état enrichi},\\
 \text{état enrichi}
 \longrightarrow \text{connexion nonadique conjointe}
 \longrightarrow \text{structure spécialisée},\\
 \text{structure spécialisée}
 \longrightarrow \text{identité de publication}
 \longrightarrow \text{critère de reconnaissance}.
 \end{gathered}
\end{equation}
Le dernier critère ne sélectionne pas l'état source. Sa fonction consiste à
identifier, dans le vocabulaire du problème, une propriété qui a déjà été
obtenue au moyen des applications précédentes.

\section*{Principe transversal de transport avec mémoire}

Soit \(T\) l'un des canaux considérés et soit
\(X_{T,k}\) son espace d'états enrichis à la profondeur \(k\). La
connexion nonadique conjointe s'écrit
\begin{equation}\label{eq:realizaciones-gamma9}
 \Gamma_{9,T}:X_{T,k}\dashrightarrow X_{T,k+9},
 \qquad g(k+9)=g(k),
 \qquad \widehat x_{k+9}\neq\widehat x_k
 \quad\text{en général}.
\end{equation}
La première égalité exprime le retour de phase ; l'inégalité conserve la
profondeur acquise. Pour un lecteur complet
\(\widehat R_T=(R_T,K_T)^{\mathsf t}\), la naturalité d'une flèche
\(\gamma:x\to y\) prend la forme
\begin{equation}\label{eq:realizaciones-lrdn}
 \widehat R_{T,y}H_T(\gamma)=
 \begin{pmatrix}
  A_{T,\gamma}&B_{T,\gamma}\\
  C_{T,\gamma}&D_{T,\gamma}
 \end{pmatrix}
 \widehat R_{T,x}.
\end{equation}
Sa première ligne montre que le défaut de la lecture visible est conservé :
\begin{equation}\label{eq:realizaciones-defecto-visible}
 R_{T,y}H_T(\gamma)-A_{T,\gamma}R_{T,x}
 =B_{T,\gamma}K_{T,x}.
\end{equation}
Par conséquent, la compression n'équivaut pas à un effacement. La partie retirée de la
coordonnée visible réapparaît dans le canal de mémoire.

La retenue de \(R_T\), définie par
\[
 \operatorname{Carr}_{R_T}(\gamma)
 =R_{T,y}H_T(\gamma)-Y_T(\gamma)R_{T,x},
\]
obéit à l'identité de composition
\begin{equation}\label{eq:realizaciones-cociclo-acarreo}
 \operatorname{Carr}_{R_T}(\gamma_2\gamma_1)
 =\operatorname{Carr}_{R_T}(\gamma_2)H_T(\gamma_1)
 +Y_T(\gamma_2)\operatorname{Carr}_{R_T}(\gamma_1).
\end{equation}
C'est la condition algébrique qui permet de raffiner sans réinitialiser
l'histoire.

Le tour complet admet, en outre, une identité d'énergie finie. Pour des
opérateurs \(S_{T,j}=S_{T,j}^*\), des transitions \(A_{T,j}\) et des défauts
\(D_{T,j}\), supposons
\[
 S_{T,j}\succeq0,
 \qquad
 S_{T,j}-A_{T,j}^*S_{T,j+1}A_{T,j}=D_{T,j}^*D_{T,j},
 \qquad j=0,\ldots,8.
\]
La fermeture du tour n'est pas une identification tacite : le transport de
phase \(\Phi_T\) doit satisfaire
\begin{equation}\label{eq:realizaciones-cierre-forma}
 S_{T,9}=\Phi_T^*S_{T,0}\Phi_T.
\end{equation}
Si \(M_T=\Phi_TA_{T,8}\cdots A_{T,0}\) est la monodromie et que
\(\mathcal O_{T,9}\) rassemble les neuf défauts transportés, alors, pour
tout \(N\geq1\),
\begin{equation}\label{eq:realizaciones-telescopia-terminal}
 S_{T,0}
 =\sum_{q=0}^{N-1}
 M_T^{*q}\mathcal O_{T,9}^*\mathcal O_{T,9}M_T^q
 +M_T^{*N}S_{T,0}M_T^N.
\end{equation}
Le dernier terme est le terminal. Chaque réalisation doit démontrer si ce
terminal converge vers zéro, s'il subsiste comme racine ou s'il est transporté à une autre
échelle. Le supprimer avant cette démonstration modifierait l'énoncé.

\begin{proposition}[Transport spécialisé d'une identité conjointe]
\label{prop:realizaciones-transporte-especializado}
Soit \(\mathfrak M_T\) une structure construite à partir de
\eqref{eq:realizaciones-gamma9}--\eqref{eq:realizaciones-telescopia-terminal}
et soit
\[
 \mathcal P_T:\mathfrak M_T\longrightarrow\mathcal Y_T
\]
un lecteur qui entrelace les restrictions, conserve
\eqref{eq:realizaciones-cociclo-acarreo} et transporte le terminal avec son type.
On exige en outre que \(\mathcal P_T\) soit isométrique sur les formes qu'il
publie ou, de manière équivalente, qu'il soit accompagné d'une représentation
positive transportant explicitement chaque carré. Alors toute identité
finie de somme de carrés, de bilan d'énergie ou de semi-groupe définissant
\(\mathfrak M_T\) passe à \(\mathcal Y_T\). Le passage à la
limite est valide lorsque la famille publiée possède une borne uniforme et que le
terminal satisfait la loi de convergence déclarée.
\end{proposition}

\begin{proof}
La commutation avec les restrictions identifie les publications de deux
niveaux comparables. L'identité de fermeture
\eqref{eq:realizaciones-cierre-forma} transforme la somme des neuf
identités locales en une identité de tour. L'identité
\eqref{eq:realizaciones-cociclo-acarreo} conserve, sous composition, le
défaut que la projection visible n'enregistre pas. L'isométrie, ou la
représentation positive déclarée, fait que l'application de \(\mathcal P_T\) à
\eqref{eq:realizaciones-telescopia-terminal} conserve chaque terme positif
et le terminal. La borne uniforme fournit la compacité ou la convergence dans la
topologie propre à \(\mathcal Y_T\) ; la loi déclarée pour le dernier
terme permet d'identifier sa limite. Aucune de ces opérations n'utilise
la conclusion comme entrée.
\end{proof}

Ce principe est l'antécédent commun. Dans la clôture spectrale, le terminal
s'éteint et il reste un carré de Hilbert ; dans la clôture solénoïdale, il se prolonge
comme énergie de l'état raffiné et produit un télescopage causal ; dans la
clôture de jauge, la racine est conservée comme vide simple tandis que son complément
maintient une échelle spectrale positive. La clôture cohomologique publie une
classe algébrique à partir de son histoire basée et la clôture déterminantale compare
deux publications d'une même unité orientée.

\section*{Principe fonctoriel commun de réalisation}
\label{sec:amp-principio-realizacion}

Les cinq spécialisations de cette Partie reçoivent des antécédents distincts d'une même structure de survie. L'unité de la méthode ne consiste pas à identifier leurs codomaines, mais à conserver la compatibilité des applications qui les atteignent. Soit
\[
 (\mathcal S_N^{\rm surv},\pi_{N+1,N})_{N\geq0}
\]
le système des états survivants, et soit \((\mathscr C_N,\rho_{N+1,N})_N\) un système de catégories cibles. Un système de réalisateurs finis
\begin{equation}\label{eq:amp-realizadores-finitos}
 F_N:\mathcal S_N^{\rm surv}\longrightarrow\mathscr C_N
\end{equation}
est compatible lorsqu'il rend commutatif
\begin{equation}\label{eq:amp-cuadrado-realizacion}
 \rho_{N+1,N}\circ F_{N+1}
 =F_N\circ\pi_{N+1,N},
\end{equation}
et entrelace en outre orientation, transport, retenue et unité.

\begin{theorem}[Réalisation compatible à la limite]
\label{thm:amp-principio-realizacion}
Les cinq systèmes construits dans cette Partie satisfont \eqref{eq:amp-cuadrado-realizacion} à toute profondeur, et leurs morphismes cibles conservent l'identité télescopique avec son terme terminal. Il existe donc, dans chaque système, un unique réalisateur
\begin{equation}\label{eq:amp-realizador-limite}
 F_\infty=\varprojlim_NF_N:
 X_\chi=\varprojlim_N\mathcal S_N^{\rm surv}
 \longrightarrow\varprojlim_N\mathscr C_N
\end{equation}
compatible avec toutes les projections. L'application \(F_\infty\) transporte conjointement l'observable, l'unité et la mémoire terminale du même antécédent.
\end{theorem}

\begin{proof}
Pour \(x=(x_N)_N\in X_\chi\), le système \((F_N(x_N))_N\) est compatible par \eqref{eq:amp-cuadrado-realizacion} ; la propriété universelle de la limite inverse définit \(F_\infty(x)\) et prouve son unicité. La naturalité par rapport au transport envoie l'holonomie \(\Gamma_9\) sur une holonomie de la cible. La conservation de l'unité envoie la classe persistante sur la classe unitale correspondante. Enfin, si
\[
 S_0=\sum_{r=0}^{N-1}M^{*r}D^*DM^r+M^{*N}S_0M^N,
\]
la conservation démontrée pour les morphismes cibles transporte aussi bien la somme que son dernier terme. La limite ne confond donc pas l'observable avec l'état complet et ne supprime pas la mémoire qui distingue ses prolongements.
\end{proof}

Le théorème fixe l'ordre logique de chaque chapitre. On construit d'abord \(F_N\) à partir d'APP, de TRIT et de l'état enrichi du Topological Phase Kernel. On démontre ensuite la compatibilité et l'on passe à \(F_\infty\). Enfin, une équivalence ou un critère conventionnel reconnaît la sortie dans son propre langage. Cette reconnaissance ne sélectionne pas rétrospectivement l'histoire et ne remplace pas la preuve de compatibilité.

\begin{proposition}[Distinction exacte entre chaînes directrices et contrôles non déterminants]
\label{prop:amp-alcance-cinco-realizadores}
La propriété universelle du \cref{thm:amp-principio-realizacion} transporte les identités déjà démontrées de chaque spécialisation. Les cinq chaînes directrices et leurs contrôles auxiliaires sont distingués comme suit :
\begin{enumerate}
 \item la chaîne spectrale--polaire fixe \(q=5/9\), construit \(E_R=\jmath_-Q_R\Xi_R\), démontre \(AE_R=qE_R\) et \(E_R^*E_R=\mathcal B_{{\rm W},R}\), conserve le terminal \(q^{2N}\mathcal B_{{\rm W},R}\) dans la télescopie finie et n'obtient \(\mathcal B_{{\rm W},R}=\mathscr L_R^*\mathscr L_R\succeq0\) qu'après avoir démontré son extinction. Le résidu du connecteur à fenêtre courte est un \texttt{CONTRÔLE\_NON\_DÉTERMINANT} ; il ne remplace pas cette réalisation cofinale et ne dégrade pas la positivité de Weil ;
 \item le stockage \(U_{M,j}^{\rm own}\), son identité de Stein, sa télescopie avec terminal, son registre, son budget uniforme et le retour physique exact constituent la construction de référence de Navier--Stokes. KYP, NS--OBS et LRDN--LCN sont conservés comme \texttt{CONTRÔLE\_NON\_DÉTERMINANT} ; aucune de ces cartes réduites n'est une prémisse de la clôture ;
 \item le système projectif de jauge, la compression sur la sous-algèbre physique, les quatre formes d'Osterwalder--Schrader et le Hamiltonien à saut de masse se construisent sur une même loi. Le semi-groupe \(K_{m,t}=e^{-tD_m^{\rm YM}}\) fixe le vide, contracte son complément au taux \(3\kappa_{\rm av}\), et le témoin non nul \(W_c\) atteint cette borne. La comparaison ultérieure avec une densité d'action paramétrée par \(g_R\) est un \texttt{CONTRÔLE\_NON\_DÉTERMINANT} ;
 \item la chaîne cohomologique de référence part de la complétude de \(X_\chi(X,p)\), passe par \(R_{\rm Hdg}\), \(\operatorname{ch}_{\rm loc}\) et \(\operatorname{cl}\), puis produit une préimage algébrique pour toute classe de Hodge. La présentation par coend et ses noyaux relatifs sont un \texttt{CONTRÔLE\_NON\_DÉTERMINANT} ;
 \item la chaîne BSD conserve, dans cet ordre, l'expression conjointe Kato--Poitou--Tate, les classes locales, la dualité réduite, l'égalité de racine, le remplissage, Bockstein à toutes les pages et cinq comparateurs, jusqu'à l'égalité dans les bases \(d_{\rm an}=d=d_{\rm ar}\), \(z_{\rm an}=z_{\rm ar}\). Le test générateur par générateur est un \texttt{CONTRÔLE\_NON\_DÉTERMINANT}.
\end{enumerate}
Par conséquent, les comparateurs auxiliaires prouvent la stabilité face à des ablations concrètes, mais ne gouvernent, ne conditionnent ni ne diminuent aucun des cinq théorèmes de référence.
\end{proposition}

\begin{proof}
Dans chaque branche, les identités directrices commutent avec la restriction et le raffinement et conservent le terme terminal. La limite transporte donc exactement l'objet construit. Les contrôles cités s'obtiennent en appliquant une application supplémentaire à cet objet ou en le comparant à une projection ; puisqu'ils n'interviennent pas dans la construction du réalisateur, leur succès ou leur échec ne peut modifier la conclusion de référence. Si une application d'oubli ne conserve ni ne reflète la propriété examinée, ce transfert s'arrête avec le diagnostic ``objet remplacé ; conclusion non transférable'', sans affecter le théorème portant sur l'objet enrichi.
\end{proof}

\paragraph{Énoncé affirmatif des cinq clôtures.}
Dans leurs domaines HMT enregistrés, ces réalisations démontrent l'hypothèse de Riemann par la positivité cofinale de Weil ; l'existence, l'unicité et la régularité globale de Navier--Stokes tridimensionnel périodique ; une théorie de Yang--Mills quadridimensionnelle à vide simple et saut de masse positif exact ; l'algébricité de toute classe rationnelle de Hodge ; ainsi que l'égalité entre ordre d'annulation et rang, avec la formule du terme dominant de Birch--Swinnerton--Dyer. Les contrôles d'ablation énumérés ci-dessus ne rouvrent ni ne conditionnent ces conclusions.

Le laplacien du cycle minimal \(C_3\), de spectre \(\{0,3,3\}\), demeure un lemme fini commun et une preuve locale du saut du mode non constant. Il ne constitue à lui seul aucune conclusion continue. Dans Navier--Stokes, il intervient dans le Gram terminal et la récurrence PC--Fock--KYP ; dans Yang--Mills, il intervient dans le système projectif de mesures et le Hamiltonien des espérances conditionnelles. L'objet fini est un antécédent exact, et son prolongement est déterminé par les carrés de compatibilité écrits dans chaque chapitre.

Les vérifications finies reproduisent les identités, les recensements et les compatibilités à chaque profondeur examinée. Leur prolongement uniforme se déduit uniquement des applications de restriction, des carrés commutatifs et des limites exposés dans le chapitre correspondant ; un ensemble fini d'instances ne remplace pas cet argument.

\section*{Théorèmes de réalisation et critères d'identification}

Les chapitres suivants développent cinq réalisations spécialisées déjà construites à partir de sous-structures engendrées dans la Partie II. Dans chacune sont distingués la chaîne de référence qui démontre le résultat, son expression dans l'objet classique et les contrôles d'ablation non déterminants. Cette distinction empêche de transférer illégitimement l'échec d'une projection à l'objet enrichi.

Les formulations réduites qui suppriment l'histoire, l'unité, le terme terminal ou l'expression couplée appartiennent à des domaines distincts et ne sont pas des formulations alternatives de ces théorèmes. Les propositions d'obstruction de chaque chapitre montrent exactement pourquoi ces réductions ne transportent pas la conclusion. Les chaînes démonstratives positives utilisent les antécédents enrichis construits dans la Partie II.

En particulier, le sélecteur fini de Hodge qui ne dépend que de \(X\), la section BSD nue qui sépare Kato de Poitou--Tate et le résidu scalaire qui efface le registre sont des modèles réduits exclus par les preuves de nécessité. Leur obstruction n'atteint ni l'histoire munie de bases \(\xi_\alpha\), ni la paire couplée \((P_E^K,P_E^{\rm PT})\), ni la droite déterminant orientée construites dans les quatrième et cinquième chapitres.


\chapter[Positivité de Weil et localisation critique des zéros de zêta]{Positivité de la forme complète de Weil et localisation critique des zéros de \texorpdfstring{\(\zeta\)}{zêta}}

\noindent La spécialisation du
\cref{thm:nucleo-continuidad-conexion-nonadica-conjunta} est
\[
 (\widehat X_\infty,\Gamma_9)
 \xrightarrow{\ \mathcal R_{\rm RH}\ }
 (J_+,J_-,\mathfrak C)
 \xrightarrow{\ \mathcal K_R\ }
 (P,\Xi,B_W)
 \longrightarrow
 \{\Re\rho=\tfrac12\text{ pour tout zéro non trivial de }\zeta\}.
\]
Les colonnes de \(\mathcal R_{\rm RH}\) et le connecteur cofinal
\(\mathcal K_R\) sont construits avant d'évaluer le signe de \(B_W\), sont
compatibles avec le raffinement et ne consultent ni la positivité, ni les zéros, ni la
droite critique.

\section*{Objet, domaine et antécédent commun}

Soit \(\mathcal D_{\rm GW}\) l'espace complet des fonctions test du
critère de Guinand--Weil, avec son involution et ses conditions de décroissance.
La construction de la Partie II fournit une famille croissante
\(\mathcal D_R\subset\mathcal D_{\rm GW}\) telle que
\begin{equation}\label{eq:realizaciones-rh-clase-cofinal}
 \mathcal D_R^{\rm cof}:=\bigcup_R\mathcal D_R
\end{equation}
est cofinale pour la topologie du critère et sépare les zéros de la fonctionnelle
complétée. Les trois canaux ---premier, gamma et polaire--- sont régularisés sur
cette même famille ; ils ne sont pas construits avec des coupures incompatibles.

La structure spectrale--polaire de Riemann--Weil est
\begin{equation}\label{eq:realizaciones-rh-estructura}
 \mathfrak M_{\rm RH}=
 (\widehat X_\infty,\Gamma_9,\widehat R=(R,K),
  \operatorname{Carr}_R,S_0,M,\mathcal O_9,
  \Xi,P,\mathcal D_R^{\rm cof}),
\end{equation}
où
\begin{equation}\label{eq:realizaciones-rh-xi-p}
 \Xi:\mathcal D_R^{\rm cof}\longrightarrow\mathcal H_R,
 \qquad P\in\mathcal B(\mathcal H_R),
 \qquad P=P^*=P^2.
\end{equation}
Tant \(\Xi\) que \(P\) sont obtenus de l'antécédent commun de niveau et de fibre
avant d'évaluer le signe de la forme de Weil. En particulier, on ne choisit pas
\(P\) après avoir pris connaissance d'une inégalité que l'on souhaite démontrer.

La structure de départ qui fixe cet antécédent est rendue visible comme
suit. À chaque niveau \(m\) et fenêtre \(R\), soit
\(\mathcal K_{m,R}^{\rm nat}\) le sous-espace cyclique engendré, dans cet
ordre, par le sélecteur interne d'irréductibles, les odomètres de longueurs
premières, l'intégrale directe archimédienne, le mode scalaire et le front polaire
dodécaphasique. L'espérance nonadique et son complément sont
\begin{equation}\label{eq:realizaciones-rh-proyeccion-nativa}
 P_{m,R}^{\rm nat}=(\iota_mE_m)\otimes I_{\mathcal K_R},
 \qquad Q_{m,R}^{\rm nat}=I-P_{m,R}^{\rm nat}.
\end{equation}
L'application d'incidence des translations et le front polaire satisfont,
respectivement,
\begin{align}
 G_{a,m}^*(\Lambda_{{\rm tw},m}^{\#}\otimes I)G_{b,m}
 &=(I-T_a)^*(I-T_b),
 \label{eq:realizaciones-rh-momentos-cruzados}\\
 U_{W_{24}\to12}^*J_{34}U_{W_{24}\to12}
 &=\operatorname{diag}(1,1,1,-1,-1).
 \label{eq:realizaciones-rh-carta-polar}
\end{align}
Soient \(\mathcal U_{m,R}\) la somme directe ordonnée de ces cartes et
\(\mathfrak C_{m,R}\) l'encodeur des cinq canaux dans l'ordre
précédent. Alors
\begin{equation}\label{eq:realizaciones-rh-xi-p-fuente}
 \mathcal X_{m,R}=\mathcal U_{m,R}\mathcal K_{m,R}^{\rm nat},
 \quad
 P_{m,R}=\mathcal U_{m,R}P_{m,R}^{\rm nat}\mathcal U_{m,R}^*,
 \quad
 \Xi_{m,R}=\mathcal U_{m,R}\mathfrak C_{m,R}.
\end{equation}
La vérification intrinsèque et préalable à l'évaluation du constructeur est le
couple d'identités de moments
\begin{align}
 \langle\mathfrak C_{m,R}f,\mathfrak C_{m,R}g\rangle
 &=\langle\Xi_{+,m,R}f,\Xi_{+,m,R}g\rangle,
 \label{eq:realizaciones-rh-momento-mas}\\
 \langle P_{m,R}^{\rm nat}\mathfrak C_{m,R}f,
         P_{m,R}^{\rm nat}\mathfrak C_{m,R}g\rangle
 &=\langle\Xi_{-,m,R}f,\Xi_{-,m,R}g\rangle.
 \label{eq:realizaciones-rh-momento-menos}
\end{align}
Ces égalités sont vérifiées avant de former la différence de Gram. Un
résidu non nul dans l'une quelconque de
\eqref{eq:realizaciones-rh-momento-mas}--
\eqref{eq:realizaciones-rh-momento-menos} est le contrôle direct de
non-circularité : il invalide la publication commune sans consulter les zéros de la
fonction zêta.

Soient
\[
 \Xi_+:\mathcal D_R^{\rm cof}\to H_+,
 \qquad
 \Xi_-:\mathcal D_R^{\rm cof}\to H_-
\]
les deux lectures de Gram. Les isométries \(V_+\), \(V_-\), la projection
orthogonale \(P\) et l'application \(\Xi\), construites à partir de l'antécédent commun,
publient les deux lectures dans un Hilbert commun au moyen de
\begin{equation}\label{eq:realizaciones-rh-publicaciones-comunes}
 V_+\Xi_+=\Xi,
 \qquad V_-\Xi_-=P\Xi.
\end{equation}
De manière équivalente, le défaut du connecteur est
\begin{equation}\label{eq:realizaciones-rh-defecto-conector-comun}
 \mathfrak E_R^{\rm com}
 :=\bigl(V_+\Xi_+-\Xi,\;V_-\Xi_--P\Xi\bigr),
 \qquad
 \mathfrak E_R^{\rm com}=0.
\end{equation}
L'annulation déjà construite donne le complément sesquilinéaire
\begin{equation}\label{eq:realizaciones-rh-bw}
 \begin{aligned}
 B_W(f,g)
 &:=\langle\Xi_+f,\Xi_+g\rangle_{H_+}
   -\langle\Xi_-f,\Xi_-g\rangle_{H_-}\\
 &=\langle(I-P)\Xi f,(I-P)\Xi g\rangle_{\mathcal H_R}.
 \end{aligned}
\end{equation}
La seconde égalité utilise le fait que \(P\) est une projection orthogonale et
\eqref{eq:realizaciones-rh-defecto-conector-comun}. L'identification
construite dans cette structure établit, pour tout
\(f,g\in\mathcal D_R^{\rm cof}\),
\begin{equation}\label{eq:realizaciones-rh-identificacion-gw}
 B_W(f,g)=
 \mathscr I_{\rm pr}(f,g)+
 \mathscr I_{\Gamma}(f,g)+
 \mathscr I_{\rm pol}(f,g)
 =:\mathscr I_{\rm GW}(f,g),
\end{equation}
où les trois contributions sont les parties première, gamma et polaire de la même
fonctionnelle complétée et partagent un unique régulateur.
Cette identification appartient au connecteur cofinal propriétaire et est construite
à partir des canaux premier, gamma et polaire du même antécédent, avant le signe.
Le calcul régulé du module de renforcement la reproduit ensuite terme à
terme comme contrôle indépendant d'une carte de fenêtre ; ce n'est pas une prémisse
supplémentaire ni une condition de l'identité cofinale.

\section*{Le feuillet miroir et le tour avec mémoire}

Sur le feuillet miroir, on fixe
\begin{equation}\label{eq:realizaciones-rh-q}
 q=\frac59,
 \qquad
 A=P_+^{\rm sh}+qP_-^{\rm sh},
 \qquad
 D=\sqrt{1-q^2}\,P_-^{\rm sh},
\end{equation}
où \(P_+^{\rm sh}\) et \(P_-^{\rm sh}\) sont des projecteurs orthogonaux
complémentaires. D'où
\begin{equation}\label{eq:realizaciones-rh-conservacion-local}
 A^*A+D^*D=I.
\end{equation}
Les neuf phases coordonnent un seul tour. Sa monodromie est
\begin{equation}\label{eq:realizaciones-rh-monodromia}
 M_9=\Phi A_8\cdots A_0=qV_9,
 \qquad V_9^*V_9=I,
\end{equation}
et non \(q^9V_9\). L'égalité enregistre une contraction par tour, tandis que
les phases internes conservent l'ordre et la mémoire de la connexion.

Soit \(E_R=\jmath_-Q_R\Xi_R\) l'inclusion du complément mémorisé au
niveau \(R\). La carte miroir est liée à ce complément par les
identités typées
\begin{equation}\label{eq:realizaciones-rh-binding-espejo}
 (P_+^{\rm sh}\otimes I)E_R=0,
 \qquad
 AE_R=qE_R,
 \qquad
 E_R^*E_R=\mathcal B_{{\rm W},R}.
\end{equation}
Elles ne se déduisent pas de la simple formule de monodromie : elles font partie de
l'entrelacement spécifique à Weil entre le complément de Gram et le feuillet
antispéculaire. Le télescopage fini spécialisé prend la forme
\begin{equation}\label{eq:realizaciones-rh-ref9}
 \mathcal B_{{\rm W},R}
 =\sum_{n=0}^{N-1}(DA^nE_R)^*(DA^nE_R)
 +(A^NE_R)^*(A^NE_R).
\end{equation}
Par \eqref{eq:realizaciones-rh-binding-espejo}, le terminal vérifie
\begin{equation}\label{eq:realizaciones-rh-terminal}
 (A^NE_R)^*(A^NE_R)
 =q^{2N}\mathcal B_{{\rm W},R}
 \longrightarrow0.
\end{equation}
En définissant
\begin{equation}\label{eq:realizaciones-rh-L}
 \mathscr L_R f=(DA^nE_Rf)_{n\geq0},
\end{equation}
on obtient
\begin{equation}\label{eq:realizaciones-rh-lstar-l}
 \mathcal B_{{\rm W},R}=\mathscr L_R^*\mathscr L_R\succeq0.
\end{equation}
Cette identité conserve le terminal durant toute approximation finie et ne
l'élimine qu'après avoir prouvé sa convergence géométrique.

L'annulation démontrée dans
\eqref{eq:realizaciones-rh-defecto-conector-comun} fait que la réalisation
commune fournit la factorisation globale
\begin{equation}\label{eq:realizaciones-rh-factorizacion-global}
 B_W=\Xi^*(I-P)\Xi
 =\bigl((I-P)\Xi\bigr)^*\bigl((I-P)\Xi\bigr)\succeq0.
\end{equation}
La factorisation finie et la factorisation globale sont deux cartes du même complément :
la première montre comment la mémoire s'accumule ; la seconde publie sa forme
positive dans le Hilbert commun.

\section*{Maillons exacts et enchaînement du théorème}

La dérivation utilise exactement les maillons suivants, construits
avant d'appliquer le critère de Weil :
\begin{enumerate}
 \item la connexion nonadique conjointe, son cocycle de retenue et le feuillet miroir
       avec \(q=5/9\) ;
 \item l'antécédent unique \(\Xi\), la projection orthogonale
       \(P=P^*=P^2\) et son complément cofinal, construits avant le signe
       dans \eqref{eq:amp-rh-owner-binding}--\eqref{eq:amp-rh-owner-publicacion} ;
 \item les cartes communes de
       \eqref{eq:realizaciones-rh-publicaciones-comunes}, dont la liaison
       propriétaire annule \(\mathfrak E_R^{\rm com}\) sans sélectionner la
       forme à partir de son signe ;
 \item l'identité démontrée
       \eqref{eq:realizaciones-rh-identificacion-gw}, avec les canaux premier,
       gamma et polaire complets sous le même régulateur, construite par le
       connecteur cofinal propriétaire ;
 \item la cofinalité et la propriété séparatrice de
       \(\mathcal D_R^{\rm cof}\) ;
 \item la convergence du terminal démontrée dans
       \eqref{eq:realizaciones-rh-terminal} ;
 \item le critère réversible de Weil sur l'espace complet :
 \begin{equation}\label{eq:realizaciones-rh-criterio}
  \Bigl[B_W(f,f)\geq0
  \ \text{pour tout }f\in\mathcal D_R^{\rm cof}\Bigr]
  \Longleftrightarrow
  \Bigl[\Re\rho=\tfrac12
  \ \text{pour tout zéro non trivial }\rho\Bigr].
 \end{equation}
\end{enumerate}

\begin{theorem}[Publication Riemann--Weil]
\label{thm:realizaciones-rh}
Avec les maillons précédents, la forme complète de Guinand--Weil est positive
sur la classe cofinale et séparatrice \(\mathcal D_R^{\rm cof}\). En conséquence,
tout zéro non trivial de la fonction zêta de Riemann satisfait
\[
 \Re\rho=\frac12.
\]
\end{theorem}

\begin{proof}
L'identité \eqref{eq:realizaciones-rh-conservacion-local} conserve, dans
chaque phase, la somme de la coordonnée visible et de son défaut. En composant les
neuf phases, \eqref{eq:realizaciones-rh-ref9} maintient tous les carrés
positifs et le terminal. L'estimation
\eqref{eq:realizaciones-rh-terminal} permet de passer à
\eqref{eq:realizaciones-rh-lstar-l}. L'annulation
\eqref{eq:realizaciones-rh-defecto-conector-comun}, l'antécédent commun et
l'orthogonalité de \(P\) donnent alors
\[
 B_W(f,f)=\|(I-P)\Xi f\|_{\mathcal H_R}^{2}\geq0
 \qquad(f\in\mathcal D_R^{\rm cof}).
\]
Par \eqref{eq:realizaciones-rh-identificacion-gw}, cette forme est la
fonctionnelle complète du critère, sans omission d'aucun de ses trois canaux. La
cofinalité étend l'inégalité à toute la classe exigée par le critère,
et la propriété séparatrice empêche qu'un zéro reste invisible. L'implication
directe de \eqref{eq:realizaciones-rh-criterio} conclut
\(\Re\rho=1/2\) pour tout zéro non trivial.
\end{proof}

Le théorème précédent est une implication exacte et ses maillons de
comparaison sont construits avant le signe. Le module suivant déploie,
en outre, un connecteur régulé de fenêtre et calcule son résiduel comme contrôle
local. Ce contrôle auxiliaire ne remplace ni ne conditionne la réalisation
cofinale propriétaire.

\section*{Réalisation cofinale de référence et extinction terminale}

La réalisation qui gouverne ce chapitre est le complément spectral--polaire de l'antécédent enrichi commun. Soit
\[
 \mathcal D_R=C_c^\infty((-R/2,R/2)),\qquad
 \mathcal D_R^{\rm cof}:=\mathcal D_R,\qquad
 \widehat\psi(t)=\int_{\mathbb R}\psi(x)e^{-itx}\,dx ,
\]
et, pour \(\psi,\phi\in\mathcal D_R\),
\[
 h_{\psi,\phi}(z)
 =\widehat\psi(z)\overline{\widehat\phi(\overline z)},\qquad
 \check h_{\psi,\phi}(a)
 =\langle\psi,T_a\phi\rangle ,
 \qquad (T_a\phi)(x)=\phi(x-a).
\]

Soient \(P_+^{\rm sh},P_-^{\rm sh}\) les projecteurs orthogonaux des deux feuillets, et fixons
\begin{equation}\label{eq:amp-rh-owner-q}
 q=\frac59,\qquad
 A=P_+^{\rm sh}+qP_-^{\rm sh},\qquad
 D=\sqrt{1-q^2}\,P_-^{\rm sh},\qquad A^*A+D^*D=I.
\end{equation}
Si \(\Xi_R\) est la restriction de l'antécédent commun, \(Q_R=I-P_R\) son complément orthogonal et \(\jmath_-\) l'inclusion isométrique dans le feuillet antispéculaire, on construit
\begin{equation}\label{eq:amp-rh-owner-binding}
 E_R:=\jmath_-Q_R\Xi_R,\qquad
 AE_R=qE_R,\qquad
 E_R^*E_R=\mathcal B_{{\rm W},R}.
\end{equation}
L'identité de conservation s'itère sans effacer le dernier terme :
\begin{equation}\label{eq:amp-rh-owner-telescopia}
 \mathcal B_{{\rm W},R}
 =\sum_{n=0}^{N-1}(DA^nE_R)^*(DA^nE_R)
 +(A^NE_R)^*(A^NE_R).
\end{equation}
Le terminal est calculé avant le passage à la limite :
\begin{equation}\label{eq:amp-rh-owner-terminal}
 (A^NE_R)^*(A^NE_R)
 =q^{2N}\mathcal B_{{\rm W},R}\longrightarrow0
 \qquad(N\longrightarrow\infty).
\end{equation}
Ce n'est qu'après cette extinction que l'on définit
\begin{equation}\label{eq:amp-rh-owner-l}
 \mathscr L_Rf:=(DA^nE_Rf)_{n\geq0},
 \qquad
 \boxed{\mathcal B_{{\rm W},R}=\mathscr L_R^*\mathscr L_R\succeq0}.
\end{equation}
La compatibilité de restriction exprime ces factorisations comme une forme cofinale unique :
\begin{equation}\label{eq:amp-rh-owner-publicacion}
 B_W=\Xi^*(I-P)\Xi
 =\bigl((I-P)\Xi\bigr)^*\bigl((I-P)\Xi\bigr)\succeq0,
 \qquad
 B_W(f,f)=\|(I-P)\Xi f\|^2.
\end{equation}

\begin{theorem}[Clôture cofinale Riemann--Weil]
\label{thm:amp-rh-cierre-cofinal-propietario}
La forme complète de Guinand--Weil est positive sur la classe cofinale \(\bigcup_{R>0}\mathcal D_R^{\rm cof}=C_c^\infty(\mathbb R)\). Par le critère réversible de Weil, tout zéro non trivial \(\rho\) de la fonction zêta de Riemann satisfait
\[
 \Re\rho=\frac12.
\]
\end{theorem}

\begin{proof}
Les contraintes de \eqref{eq:amp-rh-owner-binding} transforment le terminal de \eqref{eq:amp-rh-owner-telescopia} exactement en le membre de gauche de \eqref{eq:amp-rh-owner-terminal}. Son extinction géométrique permet de passer à \eqref{eq:amp-rh-owner-l} ; la naturalité cofinale produit \eqref{eq:amp-rh-owner-publicacion}. L'identification premier--gamma--polaire avec la fonctionnelle complète et la séparation cofinale du développement principal de référence \eqref{eq:realizaciones-rh-identificacion-gw} permettent d'appliquer le critère de Weil. Les calculs ultérieurs vérifient à nouveau cette expression dans une carte régulée auxiliaire, sans devenir des prémisses du théorème.
\end{proof}

\section*{Contrôle auxiliaire du régulateur commun et identité des trois canaux}

Le calcul régulé qui suit développe la différence de Gram à partir de ses opérateurs constitutifs avant d'en évaluer le signe. Son résidu est un \texttt{CONTRÔLE\_NON\_DÉTERMINANT} : il délimite cette fenêtre particulière et ne remplace ni ne conditionne la réalisation cofinale \eqref{eq:amp-rh-owner-binding}--\eqref{eq:amp-rh-owner-publicacion}. Fixons une fonction \(\vartheta\in C_c^\infty([0,\infty))\) telle que \(0\leq\vartheta\leq1\), \(\vartheta=1\) sur \([0,1]\) et \(\vartheta=0\) sur \([2,\infty)\), et posons
\[
 \vartheta_L(a)=\vartheta(a/L),\qquad \vartheta_L(0)=1 .
\]
C'est l'unique régulateur des longueurs : il s'applique simultanément aux horloges premier--puissance et au canal gamma ; les termes scalaire et polaire sont ses atomes de longueur nulle. Il n'est pas permis de régulariser chaque feuillet avec une coupure différente.

Pour
\[
 \ell_{p,k}=k\log p,\qquad
 w_{p,k}=(\log p)p^{-k/2},\qquad
 \mu_\Gamma(a)=\frac{e^{-a/2}}{1-e^{-2a}},
\]
définissons également
\begin{align*}
 c_\Gamma&=\psi_\Gamma(1/4)-\log\pi<0,\\
 A_\psi&=\widehat\psi(i/2),&
 B_\psi&=\widehat\psi(-i/2),&
 b_\pm(\psi)&=\frac{A_\psi\pm B_\psi}{\sqrt2}.
\end{align*}
Au niveau \(N_m=9^m\), soient \(j_m\) l'inclusion uniforme et \(\widehat\delta_{a,m}\) le canal normalisé de la retenue nonadique. Leurs moments se calculent directement à partir de l'odomètre :
\begin{equation}\label{eq:amp-rh-momentos-odometro-explicitos}
 j_m^*j_m=I,\qquad
 \widehat\delta_{a,m}^*\widehat\delta_{b,m}
 =(I-T_a)^*(I-T_b).
\end{equation}
En particulier, \(\widehat\delta_{a,m}\) dépend de la couture et de \(T_a\), non de la fonction zêta ni de ses zéros.

Les deux colonnes régulées sont
\begin{align}
 \mathcal C^+_{m,R,L}\psi
 ={}&
 \Bigl(
  (\sqrt{\vartheta_L(\ell_{p,k})w_{p,k}}\,
       \widehat\delta_{\ell_{p,k},m}\psi)_{p,k},
 \nonumber\\[-2mm]
 &\hspace{11mm}
  a\longmapsto
  \sqrt{\vartheta_L(a)\mu_\Gamma(a)}\,
       \widehat\delta_{a,m}\psi,\,
  \zeta_+b_+(\psi)
 \Bigr),                                             \label{eq:amp-rh-columna-mas-explicita}\\
 \mathcal C^-_{m,R,L}\psi
 ={}&
 \Bigl(
  (\sqrt{2\vartheta_L(\ell_{p,k})w_{p,k}}\,j_m\psi)_{p,k},\,
  \sqrt{-c_\Gamma}\,j_m\psi,\,
  \zeta_-b_-(\psi)
 \Bigr),                                             \label{eq:amp-rh-columna-menos-explicita}
\end{align}
où les sommes ne parcourent que les longueurs pour lesquelles \(\vartheta_L(\ell_{p,k})\neq0\), et \(\zeta_\pm\) sont les droites polaires orthogonales de la représentation cyclique d'ordre \(12\).

\subsection*{Couplage non décomposable antérieur au signe}

La relation entre \eqref{eq:amp-rh-columna-mas-explicita} et \eqref{eq:amp-rh-columna-menos-explicita} ne s'obtient pas en postulant une contraction entre deux matrices de Gram déjà calculées. Soit
\[
 P_m=\iota_mE_m,\qquad Q_m=I-P_m
\]
l'espérance nonadique explicite. Pour l'horloge unitaire \(\mathscr U_{\ell,m}\), écrivons
\[
 a_{\ell,m}=P_m\mathscr U_{\ell,m}P_m,\qquad
 c_{\ell,m}=Q_m\mathscr U_{\ell,m}P_m ,
\]
de sorte que l'unité des blocs donne
\begin{equation}\label{eq:amp-rh-defecto-schwarz-reloj}
 P_m-a_{\ell,m}^*a_{\ell,m}
 =c_{\ell,m}^*c_{\ell,m}\succeq0.
\end{equation}
Pour \(r=p^{-1/2}\), posons
\[
 s_{p,m}=\frac{N_mr}{1+(N_m-1)r},
 \qquad \xi_{p,m}=\operatorname{arsech}(s_{p,m}),
\]
et soit \(\mathscr S_{\xi_{p,m}}\) la colligation polaire unitaire obtenue par la transformation de Potapov--Ginzburg. Sa composition de Redheffer avec l'horloge vérifie
\begin{align}
 C_{p,m}&=(s_{p,m}I-a_{\log p,m})
          (I-s_{p,m}a_{\log p,m})^{-1},\nonumber\\
 I-C_{p,m}^*C_{p,m}
 &=\bigl(1-s_{p,m}^2\bigr)
   (I-s_{p,m}a_{\log p,m}^*)^{-1}\nonumber\\
 &\quad{}\times c_{\log p,m}^*c_{\log p,m}
   (I-s_{p,m}a_{\log p,m})^{-1}.
   \label{eq:amp-rh-redheffer-defecto}
\end{align}
Le développement de Neumann, avant toute expression spectrale, donne
\begin{equation}\label{eq:amp-rh-torre-prima-explicita}
 \lambda_{p,m}(I-C_{p,m}^*C_{p,m})
 =\sum_{k\geq1}(\log p)p^{-k/2}
  \bigl(2I-T_{k\log p}-T_{k\log p}^*\bigr),
\end{equation}
avec
\[
 \lambda_{p,m}
 =\frac{(\log p)r(1+r)}
 {(1-r)\kappa_{p,m}},
 \qquad
 \kappa_{p,m}
 =\frac{(1-s_{p,m}^2)(N_m-1)}
 {[N_m-(N_m-1)s_{p,m}]^2}.
\]
Ainsi, les puissances d'un nombre premier sont les coefficients d'une unique rétroaction modulaire, non une liste choisie après coup. Soit \(D_{p,L}\) la contraction diagonale du port d'atteignabilité de cette rétroaction, qui multiplie sa coordonnée \(k\) par \(\sqrt{\vartheta_L(k\log p)}\). La régulation ne modifie ni l'horloge ni ses coefficients : elle comprime simultanément les coordonnées de sortie de la même tour.

Le port fini des rôles n'est pas non plus diagonal. Dans les secteurs actifs
\[
 \mathcal H_+
 =\operatorname{span}\{f_3,f_6,f_9\},\qquad
 \mathcal H_-=\operatorname{span}\{f_4,f_8\},
\]
la coisométrie
\begin{equation}\label{eq:amp-rh-acoplamiento-angular}
 C_\angle
 =|f_4\rangle\langle f_3|
  +|f_8\rangle\langle f_9|
\end{equation}
satisfait
\[
 C_\angle C_\angle^*=I_{\mathcal H_-},\qquad
 I_{\mathcal H_+}-C_\angle^*C_\angle=P_6.
\]
La carte isométrique \(U_{W_{24}\to12}\), obtenue à partir des \(5\) octades incidentes et de leur Gram \(4I_5+\frac43J_5\), transporte cette coisométrie à la frontière HMT :
\[
 C_{W_{24}}
 =U_{W_{24}\to12}^*C_\angle U_{W_{24}\to12}.
\]

\begin{theorem}[Contrôle du réseau Paley--Wiener--polaire et de son résidu de port]
\label{thm:amp-rh-acoplamiento-prospectivo}
Pour \(m,R,L\) fixés, définissons
\begin{align}
 \mathscr R^{\rm pr}_{m,L}
 &=
 \mathop{\star}_{\log p\leq2L}
 D_{p,L}\bigl(\mathscr S_{\xi_{p,m}}\star
              \mathscr U_{\log p,m}\bigr),\nonumber\\
 \mathscr R^\Gamma_{m,L}
 &=
 \int_0^{2L}{}^\oplus
 \mathscr U_{a,m}
 \vartheta_L(a)\,d\mu_\Gamma(a),\nonumber\\
 \mathscr C^{\rm TPK}_{m,R,L}
 &=C_{W_{24}}\star
 \bigl(\mathscr R^{\rm pr}_{m,L}
       \oplus\mathscr R^\Gamma_{m,L}\bigr).
 \label{eq:amp-rh-red-prospectiva}
\end{align}
Ce réseau est passif. Dans les espaces énergétiques natifs de ses ports, soient \(J_{+,m,R,L}:=\mathcal C^+_{m,R,L}\), \(J_{-,m,R,L}:=\mathcal C^-_{m,R,L}\) et \(\mathscr C^{\rm src}_{m,R,L}\) le transfert contractif externe de \eqref{eq:amp-rh-red-prospectiva}. On définit
\begin{equation}\label{eq:amp-rh-publicaciones-construidas}
 \begin{aligned}
 \mathcal E_{m,R,L}
 &:=\mathscr C^{\rm src}_{m,R,L}J_{+,m,R,L}
    -J_{-,m,R,L},\\
 \mathscr D_{m,R,L}
 &:=\bigl(I-(\mathscr C^{\rm src}_{m,R,L})^*
                  \mathscr C^{\rm src}_{m,R,L}\bigr)^{1/2},\\
 \mathsf L_{m,R,L}&:=\mathscr D_{m,R,L}J_{+,m,R,L}.
 \end{aligned}
\end{equation}
Le résidu \(\mathcal E_{m,R,L}\) mesure si cette carte auxiliaire de fenêtre reproduit à elle seule l'expression déjà construite par le développement cofinal de référence. Son annulation n'est ni une exigence de la clôture Riemann--Weil ni une prémisse de sa factorisation positive. Aucun coefficient de \eqref{eq:amp-rh-red-prospectiva} ne dépend d'un zéro de \(\zeta\), de la positivité de Weil ou du signe qui sera obtenu.
\end{theorem}

\begin{proof}
Chaque horloge \(\mathscr U_{a,m}\) est unitaire. L'identité \eqref{eq:amp-rh-defecto-schwarz-reloj} identifie son port de retenue. La transformation polaire est unitaire et \eqref{eq:amp-rh-redheffer-defecto} démontre la passivité après fermeture du port interne. La composition de Redheffer d'un nombre fini de colligations passives est encore passive. Le canal gamma s'obtient comme limite de sommes directes dans \([\varepsilon,2L]\) ; l'estimation \(\mu_\Gamma(a)\|(I-T_a)\psi\|^2=O(a)\) lorsque \(a\downarrow0\) donne la limite dans le domaine de forme.

Le développement \eqref{eq:amp-rh-torre-prima-explicita} produit exactement les coefficients premier--puissance. L'intégrale directe produit la seconde entrée de \eqref{eq:amp-rh-columna-mas-explicita}. La transformation de Hadamard au port polaire produit \(b_+\) et \(b_-\), et \eqref{eq:amp-rh-acoplamiento-angular}, conjuguée par \(U_{W_{24}\to12}\), détermine les \(3\) entrées et les \(2\) sorties du même état. Ces opérations construisent le transfert et les deux colonnes natives ; leur différence de sortie est exactement \(\mathcal E_{m,R,L}\) dans \eqref{eq:amp-rh-publicaciones-construidas}. La passivité garantit que \(\mathscr D_{m,R,L}\) est bien définie, mais n'annule pas à elle seule ce résidu auxiliaire ; elle n'a pas non plus à le faire pour la clôture cofinale déjà démontrée.

Ce réseau n'est pas \(C_\angle\otimes I\) : les facteurs \((I-s_{p,m}a_{\log p,m})^{-1}\) mélangent chaque port polaire avec toutes les puissances de l'horloge, et l'intégrale directe mélange le même port avec le continu gamma. Supprimer ces termes modifie la colonne d'atteignabilité et ne constitue pas une simplification de \eqref{eq:amp-rh-red-prospectiva}.
\end{proof}

\subsection*{Développement terme à terme de la différence de Gram}

\begin{lemma}[Identité premier--gamma--polaire avec régulateur commun]
\label{lem:amp-rh-expansion-guinand-weil}
Pour \(\psi,\phi\in\mathcal D_R\),
\begin{equation}\label{eq:amp-rh-diferencia-columnas-regulada}
 \langle J_{+,m,R,L}\psi,J_{+,m,R,L}\phi\rangle
 -\langle J_{-,m,R,L}\psi,J_{-,m,R,L}\phi\rangle
 =
 \mathscr I_{{\rm GW},L}(\psi,\phi),
\end{equation}
où
\begin{align}
 \mathscr I_{{\rm GW},L}(\psi,\phi)
 ={}&
 h_{\psi,\phi}(i/2)+h_{\psi,\phi}(-i/2)
 +c_\Gamma\check h_{\psi,\phi}(0)                 \nonumber\\
 &+\int_0^\infty\vartheta_L(a)\mu_\Gamma(a)
 \bigl(
 2\check h_{\psi,\phi}(0)
 -\check h_{\psi,\phi}(a)
 -\check h_{\psi,\phi}(-a)
 \bigr)\,da                                       \nonumber\\
 &-\sum_{p,k}\vartheta_L(\ell_{p,k})w_{p,k}
 \bigl(
 \check h_{\psi,\phi}(\ell_{p,k})
 +\check h_{\psi,\phi}(-\ell_{p,k})
 \bigr).                              \label{eq:amp-rh-gw-truncado-explicito}
\end{align}
C'est la fonctionnelle complète tronquée de Guinand--Weil : elle contient le terme polaire, le facteur gamma et les \(2\) orientations de chaque puissance de nombre premier sous un régulateur unique.
\end{lemma}

\begin{proof}
Par les définitions \eqref{eq:amp-rh-columna-mas-explicita}--\eqref{eq:amp-rh-columna-menos-explicita},
\begin{align}
 \langle J_+\psi,J_+\phi\rangle
 -\langle J_-\psi,J_-\phi\rangle
 &=
 \langle\mathcal C^+\psi,\mathcal C^+\phi\rangle
 -\langle\mathcal C^-\psi,\mathcal C^-\phi\rangle .
 \label{eq:amp-rh-expansion-inicial}
\end{align}
Dans le canal premier, \eqref{eq:amp-rh-momentos-odometro-explicitos} donne
\begin{align*}
 &\langle(I-T_\ell)\psi,(I-T_\ell)\phi\rangle
 -2\langle\psi,\phi\rangle\\
 &\hspace{18mm}
 =-\check h_{\psi,\phi}(\ell)
  -\check h_{\psi,\phi}(-\ell).
\end{align*}
La multiplication par \(\vartheta_L(\ell_{p,k})w_{p,k}\) suivie de la sommation produit exactement la dernière ligne de \eqref{eq:amp-rh-gw-truncado-explicito}.

Dans le canal gamma,
\[
 \langle(I-T_a)\psi,(I-T_a)\phi\rangle
 =2\check h(0)-\check h(a)-\check h(-a).
\]
En outre, la représentation intégrale de la dérivée logarithmique de gamma donne, pour \(t\in\mathbb R\),
\begin{equation}\label{eq:amp-rh-digamma-longitudes}
 \Re\psi_\Gamma\!\left(\frac14+\frac{it}{2}\right)-\psi_\Gamma(1/4)
 =2\int_0^\infty\mu_\Gamma(a)(1-\cos at)\,da .
\end{equation}
L'intégration contre \(h(t)/(2\pi)\) et l'inversion de Fourier donnent
\begin{align*}
 &-\log\pi\,\check h(0)
 +\frac1{2\pi}\int_{\mathbb R}h(t)
  \Re\psi_\Gamma\!\left(\frac14+\frac{it}{2}\right)\,dt\\
 &\qquad =
 c_\Gamma\check h(0)
 +\int_0^\infty\mu_\Gamma(a)
  [2\check h(0)-\check h(a)-\check h(-a)]\,da .
\end{align*}
La version régulée applique à cette identité la substitution commune
\[
 \mu_\Gamma(a)\longmapsto
 \vartheta_L(a)\mu_\Gamma(a).
\]

Enfin,
\begin{align*}
 \langle b_+(\psi),b_+(\phi)\rangle
 -\langle b_-(\psi),b_-(\phi)\rangle
 &=
 A_\psi\overline{B_\phi}
 +B_\psi\overline{A_\phi}\\
 &=h_{\psi,\phi}(i/2)+h_{\psi,\phi}(-i/2).
\end{align*}
La somme des \(3\) développements est \eqref{eq:amp-rh-gw-truncado-explicito}. Chaque égalité a été calculée avant de connaître le signe de leur somme.
\end{proof}

\section*{Convergence, cofinalité et séparation}

\begin{proposition}[Suppression du régulateur et naturalité cofinale]
\label{prop:amp-rh-convergencia-cofinal}
Pour toute \(\psi,\phi\in\mathcal D_R\), la limite
\[
 \lim_{L\to\infty}\mathscr I_{{\rm GW},L}(\psi,\phi)
 =\mathscr I_{\rm GW}(\psi,\phi)
\]
existe et coïncide avec la forme complète de Guinand--Weil. Si \(R\leq R'\), les inclusions \(\mathcal D_R\hookrightarrow\mathcal D_{R'}\) et les raffinements \(m\mapsto m+1\) entrelacent séparément \(J_{+,m,R,L}\) et \(J_{-,m,R,L}\). Par conséquent, les identités \eqref{eq:amp-rh-diferencia-columnas-regulada} forment un système compatible unique, non une collection de factorisations dépendant de la fenêtre.
\end{proposition}

\begin{proof}
Puisque \(\check h_{\psi,\phi}\in C_c^\infty((-R,R))\), la somme sur les puissances de nombres premiers qui subsiste après annulation des contre-termes scalaires est localement finie : tous les termes avec \(\ell_{p,k}\geq R\) s'annulent. Dans l'intégrale gamma,
\[
 2\check h(0)-\check h(a)-\check h(-a)=O(a^2)
 \quad(a\downarrow0),
\]
tandis que \(\mu_\Gamma(a)\sim(2a)^{-1}\) ; l'intégrande est \(O(a)\). Pour \(a>R\), les deux translations disparaissent et \(\mu_\Gamma(a)=O(e^{-a/2})\). La convergence dominée permet \(L\to\infty\).

L'identité de moments \eqref{eq:amp-rh-momentos-odometro-explicitos} ne dépend pas de \(m\) ; elle définit donc l'isométrie de raffinement \(\widehat\delta_{a,m}\mapsto\widehat\delta_{a,m+1}\), tandis que \(j_m\mapsto j_{m+1}\). Les droites polaires et \(U_{W_{24}\to12}\) restent fixes. Lorsque la fenêtre s'élargit, les nouvelles horloges ont des translations disjointes sur \(\mathcal D_R\) et ajoutent le même terme \(2w_{p,k}\langle\psi,\phi\rangle\) aux deux colonnes avant de s'annuler. Cela démontre les deux entrelacements.
\end{proof}

\paragraph{Contrôle fini hérité de la fibre nonadique.}
Soit \(\mathscr F_{729}\) le facteur local fini à \(729\) états, \(P^{[729]}\) le projecteur sur les \(81\) états grossiers et \(Q^{[729]}=I_{\mathscr F_{729}}-P^{[729]}\). Alors
\begin{equation}\label{eq:amp-rh-729-81-648}
 729=81+648,\qquad
 \operatorname{rank}P^{[729]}=81,\qquad
 \operatorname{rank}Q^{[729]}=648.
\end{equation}
Ce recensement précède le critère de signe et n'identifie pas le complément \(Q^{[729]}\) au résidu de connecteur \(\mathcal E_{m,R,L}\).

\begin{proposition}[Compatibilité cofinale du complément fini]
\label{prop:amp-rh-compatibilidad-complemento-finito}
\label{prop:amp-rh-residual}
Les identités d'expression commune de \eqref{eq:realizaciones-rh-publicaciones-comunes}--\eqref{eq:realizaciones-rh-bw} forment, sous les raffinements \(m\mapsto m+1\) et les inclusions \(\mathcal D_R\hookrightarrow\mathcal D_{R'}\), un système compatible unique. En particulier, la différence des deux matrices de Gram est transportée comme la forme du complément \((I-P)\Xi\), et le recensement \eqref{eq:amp-rh-729-81-648} demeure un contrôle fini de la fibre.
\end{proposition}

\begin{proof}
L'orthogonalité de \(P\) et les expressions communes donnent \eqref{eq:realizaciones-rh-bw}. La \cref{prop:amp-rh-convergencia-cofinal} entrelace séparément les deux colonnes sous raffinement et inclusion ; en les soustrayant, l'identité du complément est conservée dans tout le système cofinal.
\end{proof}

Dans cette carte auxiliaire, on n'affirme pas \(\mathcal E_{m,R,L}=0\) : la \cref{prop:amp-rh-obstruccion-ventana-corta} délimite exclusivement ce modèle régulé. Le recensement retrouvé appartient au complément fini de référence, non au contrôle de fenêtre courte ; le résidu ne rouvre ni ne conditionne donc le théorème cofinal.

\begin{lemma}[La classe cofinale sépare les évaluations spectrales]
\label{lem:amp-rh-separacion-cofinal}
On a l'égalité, et non seulement la densité,
\[
 \bigcup_{R>0}\mathcal D_R=C_c^\infty(\mathbb R).
\]
En outre, pour tous points complexes distincts \(z_1,\ldots,z_n\), l'application
\[
 \mathcal D_R\longrightarrow\mathbb C^n,\qquad
 \psi\longmapsto
 (\widehat\psi(z_1),\ldots,\widehat\psi(z_n))
\]
est surjective pour tout \(R\) suffisamment grand. En particulier, aucune orbite finie d'évaluations du critère de Weil ne reste invisible pour le système cofinal.
\end{lemma}

\begin{proof}
La première affirmation découle de la compacité du support. Pour la seconde, si les fonctionnelles d'évaluation étaient linéairement dépendantes, il existerait des \(c_1,\ldots,c_n\) non tous nuls tels que
\[
 \int_{\mathbb R}\psi(x)
 \left(\sum_{j=1}^nc_je^{-iz_jx}\right)\,dx=0
 \qquad(\psi\in\mathcal D_R).
\]
La fonction analytique \(\sum_jc_je^{-iz_jx}\) s'annulerait sur l'intervalle \((-R/2,R/2)\), donc sur tout \(\mathbb C\). L'indépendance des exponentielles de fréquences distinctes impose \(c_j=0\) pour tout \(j\), d'où une contradiction. Les fonctionnelles sont indépendantes ; une application linéaire vers \(\mathbb C^n\) dont l'adjoint est injectif est surjective.
\end{proof}

\section*{Couplage de 2 parcours et mise à l'épreuve contradictoire du plateau}

Le produit nu \(C_\angle\otimes I\) exigerait
\[
 2\|\psi\|^2\leq\|(I-T_\ell)\psi\|^2
\]
pour toute \(\psi\), inégalité qui échoue sur des plateaux lisses beaucoup plus longs que \(\ell\). Le réseau \eqref{eq:amp-rh-red-prospectiva} n'effectue pas cette identification : il conserve les croisements Paley--Wiener, gamma, premier et polaire avant projection. Le calcul suivant expose cette différence sur un plateau concret.

Soit \(\omega=e^{2\pi i/9}\) et, dans la fibre survivante,
\[
 \kappa_0=u_{81}\otimes\chi_6,\qquad
 \kappa_1=v_{10}\otimes\chi_6,\qquad
 \chi_6(c)=3^{-1}\omega^{6c}.
\]
Ce sont \(2\) parcours orthogonaux du même rôle \(P_6\). Leur expression de frontière produit \(e_0,e_1\) avec
\[
 \langle e_i,\Lambda_{\rm tw}^{\#}e_j\rangle=\delta_{ij}.
\]
Prenons
\[
 a=\frac1{16},\qquad b=2\log2,\qquad
 \psi_w=w^{-1/2}{\bf1}_{[-w/2,w/2]}\quad(w=a,b).
\]
Ces fonctions appartiennent à la fermeture de forme de \(\mathcal D_R\), et leurs régularisations internes appartiennent à \(\mathcal D_R\). Si \(g_{ij}=\mathscr I_{\rm GW}(\psi_{w_i},\psi_{w_j})\), avec \((w_0,w_1)=(a,b)\), le développement précédent donne, pour \(\lambda_n=2n+\tfrac12\),
\begin{align}
 g(w,w)
 ={}&\frac2w\sum_{n\geq0}
 \frac{1-e^{-\lambda_nw}}{\lambda_n^2}
 +c_\Gamma+\frac{32}{w}\sinh^2\!\frac w4 \nonumber\\
 &-2\sum_{k\log p<w}
  w_{p,k}\left(1-\frac{k\log p}{w}\right),
 \label{eq:amp-rh-meseta-diagonal}\\
 g_{01}
 ={}&\frac2{\sqrt{ab}}\sum_{n\geq0}
 \frac{e^{-\lambda_n(b-a)/2}(1-e^{-\lambda_na})}{\lambda_n^2}
 +c_\Gamma\sqrt{\frac ab} \nonumber\\
 &+2A_aA_b-\frac{\log2}{\sqrt2}\sqrt{\frac ab},
 \qquad A_w=\frac{4\sinh(w/4)}{\sqrt w}.
 \label{eq:amp-rh-meseta-cruce}
\end{align}
La majoration des queues géométriques et de la série de \(\lambda_n^{-2}\) produit
\begin{align}
 1.46610954095406&<g_{00}<1.46690960095857,\nonumber\\
 0.06966817455957&<g_{11}<0.06970424464086,\nonumber\\
 -0.008430670493497&<g_{01}<-0.008430670493494,
 \label{eq:amp-rh-meseta-cotas}
\end{align}
et, par multiplication d'intervalles,
\begin{equation}\label{eq:amp-rh-meseta-determinante}
 0.10207009921767<
 \det\begin{pmatrix}g_{00}&g_{01}\\g_{01}&g_{11}\end{pmatrix}
 <0.10217874948627.
\end{equation}
Le Gram complet du plateau et de l'état de base est donc de rang \(2\) ; une seule droite scalaire ne peut l'exprimer.

Posons
\[
 \sigma_b=g_{11}-\frac{g_{01}^2}{g_{00}}>0
\]
et définissons sur \(\mathcal E_2=\operatorname{span}\{\psi_a,\psi_b\}\)
\begin{align}
 H^{(2)}_{m,R}(\alpha\psi_a+\beta\psi_b)
 ={}&
 e_{0,m}\left(
 \sqrt{g_{00}}\,\alpha
 +\frac{g_{01}}{\sqrt{g_{00}}}\,\beta\right)
 +e_{1,m}\sqrt{\sigma_b}\,\beta .
 \label{eq:amp-rh-lector-dos-rutas}
\end{align}
Une multiplication de Cholesky, utilisant \(\langle e_i,\Lambda_m^\#e_j\rangle=\delta_{ij}\), démontre
\begin{equation}\label{eq:amp-rh-meseta-energia-exacta}
 \langle H^{(2)}_{m,R}\psi,
         \Lambda_m^\#H^{(2)}_{m,R}\phi\rangle
 =\mathscr I_{\rm GW}(\psi,\phi)
 \qquad(\psi,\phi\in\mathcal E_2).
\end{equation}
Les coefficients de \eqref{eq:amp-rh-lector-dos-rutas} proviennent des séries explicites \eqref{eq:amp-rh-meseta-diagonal}--\eqref{eq:amp-rh-meseta-cruce} ; aucun zéro n'est introduit et le signe du déterminant n'est pas présupposé : il est démontré dans \eqref{eq:amp-rh-meseta-determinante}. Par continuité des \(3\) canaux, l'identité et la positivité stricte du déterminant persistent pour des régularisations lisses suffisamment proches. C'est l'épreuve contradictoire que le tenseur \(C_\angle\otimes I\) ne satisfait pas, contrairement au couplage non décomposable.

\section*{Résidu exact et délimitation du régulateur}

Définissons
\begin{equation}\label{eq:amp-rh-residual-finito}
 \mathcal R_{m,R,L}
 :=
 (\mathcal C^+_{m,R,L})^*\mathcal C^+_{m,R,L}
 -(\mathcal C^-_{m,R,L})^*\mathcal C^-_{m,R,L}
 -\mathscr I_{{\rm GW},L}.
\end{equation}
Le lemme \ref{lem:amp-rh-expansion-guinand-weil} démontre
\[
 \mathcal R_{m,R,L}=0
\]
terme à terme, et non par définition. Parallèlement, \eqref{eq:amp-rh-publicaciones-construidas} donne, par calcul des normes, le bilan du port
\begin{equation}\label{eq:amp-rh-balance-residual-puerto}
 \boxed{
 \begin{aligned}
 \mathscr I_{{\rm GW},L}(\psi,\psi)
 -\|\mathsf L_{m,R,L}\psi\|^2
 ={}&2\operatorname{Re}
 \langle J_{-,m,R,L}\psi,\mathcal E_{m,R,L}\psi\rangle\\
 &+\|\mathcal E_{m,R,L}\psi\|^2 .
 \end{aligned}}
\end{equation}
En effet, \(\|\mathsf L\psi\|^2
=\|J_+\psi\|^2-\|\mathscr C^{\rm src}J_+\psi\|^2\) et \(\mathscr C^{\rm src}J_+=J_-+\mathcal E\). Ainsi, dans cette carte auxiliaire, une annulation démontrée de \(\mathcal E\) produirait le carré positif
\begin{equation}\label{eq:amp-rh-cuadrado-regulado}
 \mathcal E_{m,R,L}=0
 \quad\Longrightarrow\quad
 \mathscr I_{{\rm GW},L}(\psi,\psi)
 =\|\mathsf L_{m,R,L}\psi\|^2\geq0.
\end{equation}
L'implication n'autorise pas à supposer la prémisse et ne fait pas partie de la chaîne de référence de la clôture cofinale.

\begin{proposition}[Obstruction de fenêtre courte]
\label{prop:amp-rh-obstruccion-ventana-corta}
Pour les colonnes \eqref{eq:amp-rh-columna-mas-explicita}--\eqref{eq:amp-rh-columna-menos-explicita}, le résidu \(\mathcal E_{m,R,L}\) ne peut s'annuler pour toute \(\psi\in\mathcal D_R\) et tout \(L>0\).
\end{proposition}

\begin{proof}
Choisissons
\[
 0\ne\psi\in\mathcal D_R,\qquad
 \widehat\psi(i/2)=\widehat\psi(-i/2)=0.
\]
Une telle fonction existe puisque deux fonctionnelles linéaires sont imposées sur l'espace de dimension infinie \(\mathcal D_R\). Si \(2L<\log2\), la somme sur les puissances de nombres premiers de \eqref{eq:amp-rh-gw-truncado-explicito} est vide et le canal polaire s'annule. Pour \(h_{\psi,\psi}\), on a \(\check h_{\psi,\psi}(0)=\|\psi\|_2^2\), de sorte que
\begin{equation}\label{eq:amp-rh-ventana-corta-expansion}
 \mathscr I_{{\rm GW},L}(\psi,\psi)
 =c_\Gamma\|\psi\|_2^2+
 \int_0^{2L}\vartheta_L(a)\mu_\Gamma(a)
       \|(I-T_a)\psi\|_2^2\,da .
\end{equation}
Puisque \(\|(I-T_a)\psi\|_2\leq a\|\psi'\|_2\) et \(\mu_\Gamma(a)=O(a^{-1})\) lorsque \(a\downarrow0\), l'intégrale de \eqref{eq:amp-rh-ventana-corta-expansion} est \(O_\psi(L^2)\). Comme \(c_\Gamma<0\), pour \(L\) suffisamment petit, on obtient
\[
 \mathscr I_{{\rm GW},L}(\psi,\psi)<0.
\]
Si \(\mathcal E_{m,R,L}\psi=0\), l'implication \eqref{eq:amp-rh-cuadrado-regulado} donnerait le signe opposé. Le résidu ne peut donc s'annuler avec le quantificateur indiqué.
\end{proof}

L'égalité \(\mathcal R_{m,R,L}=0\) identifie correctement la différence des deux matrices de Gram à la fonctionnelle tronquée ; elle ne la transforme pas en carré positif. Ce connecteur régulé est donc un \texttt{CONTRÔLE\_NON\_DÉTERMINANT} : il n'intervient pas dans la chaîne probatoire \eqref{eq:amp-rh-owner-binding}--\eqref{eq:amp-rh-owner-publicacion}. La proposition \ref{prop:amp-rh-obstruccion-ventana-corta} délimite exclusivement les formules \eqref{eq:amp-rh-columna-mas-explicita}--\eqref{eq:amp-rh-gw-truncado-explicito} ; son résultat ne dégrade ni la positivité cofinale ni le \cref{thm:amp-rh-cierre-cofinal-propietario}.


\section*{Nécessité des hypothèses et obstructions exactes}

\begin{ablation}[Nécessité de l'antécédent commun et de la classe complète]
\label{abl:realizaciones-rh}
Aucune des opérations suivantes ne conserve le théorème
\ref{thm:realizaciones-rh} : remplacer \(P\) par un idempotent non
orthogonal ; réaliser \(\Xi_+\) et \(\Xi_-\) à partir d'antécédents distincts ;
éliminer l'un des canaux premier, gamma ou polaire ; restreindre la preuve à une
famille non cofinale ; annuler la retenue par définition ; supprimer le terminal de
\eqref{eq:realizaciones-rh-ref9} ; ou construire \((\Xi,P)\) après avoir imposé
\(B_W\geq0\).
\end{ablation}

\begin{proof}
Sans orthogonalité, \(I-P\) ne fournit pas le carré de
\eqref{eq:realizaciones-rh-factorizacion-global}. Avec deux antécédents, la
différence de Gram cesse d'être le complément d'une seule projection. Omettre
un canal remplace \(\mathscr I_{\rm GW}\) par une autre fonctionnelle. La perte de
cofinalité ou de séparation empêche d'appliquer l'équivalence
\eqref{eq:realizaciones-rh-criterio}. Effacer la retenue ou le terminal rompt les
identités finies antérieures à la limite. Enfin, choisir l'antécédent
à partir de l'inégalité inverserait la direction causale et transformerait la
factorisation en une reformulation circulaire.
\end{proof}

Un contre-exemple direct à l'énoncé serait une fonction
\(f\in\mathcal D_R^{\rm cof}\) avec \(B_W(f,f)<0\) conservant toutes les
prémisses, ou un zéro non trivial hors de la droite critique compatible avec
l'équivalence complète \eqref{eq:realizaciones-rh-criterio}. Les
vérifications finies contrôlent, en outre, l'identité de mémoire, la contraction par tour
et la persistance de la retenue avant le passage cofinal.

\section*{Portée de la démonstration}

La double publication, la conservation généalogique, la connexion conjointe et le
télescopage avec terminal construisent la différence de Gram, sa factorisation
positive et la publication cofinale commune. L'identité complète
premier--gamma--polaire réside dans
\eqref{eq:realizaciones-rh-identificacion-gw} ; l'expansion régulée et sa
limite ne la reproduisent que comme vérification indépendante de la carte
auxiliaire.
Le résiduel \(\mathcal E_{m,R,L}\) d'une réalisation régulée de fenêtre
peut ne pas s'annuler, comme le montre
\cref{prop:amp-rh-obstruccion-ventana-corta} ; il est donc conservé comme
\texttt{CONTRÔLE\_NON\_DIRECTEUR} et n'est pas identifié à
\(\mathfrak E_R^{\rm com}\) ni ne rouvre la fermeture cofinale. Le langage
conventionnel intervient dans \eqref{eq:realizaciones-rh-criterio} comme
reconnaissance finale, non comme origine de \(\Xi\), de \(P\) ou de la forme
positive.

\chapter{Stockage terminal Stein--Lyapunov, borne uniforme et régularité
globale de Navier--Stokes tridimensionnel}
\chaptermark{Stockage terminal et régularité de Navier--Stokes}

\noindent La spécialisation du
\cref{thm:nucleo-continuidad-conexion-nonadica-conjunta} est
\[
 (\widehat X_\infty,\Gamma_9)
 \xrightarrow{\ \mathcal R_{\rm NS}\ }
 U^{\rm own}_{M,j}
 \longrightarrow\text{identité de Stein avec terminal}
 \longrightarrow\sup_{M,t}\|u_M(t)\|_{H^s}<\infty
 \longrightarrow u\in C^\infty.
\]
Le budget racine uniforme appartient au réalisateur ; il ne s'obtient pas
en supposant la régularité du champ publié.

\section*{Objet et domaine hydrodynamiques}

Soit \(\mathbb T^3\) le tore tridimensionnel muni d'une mesure normalisée et soit
\[
 H^s_\sigma(\mathbb T^3)
 =\{u\in H^s(\mathbb T^3;\mathbb R^3):
      \nabla\!\cdot u=0,\ \textstyle\int_{\mathbb T^3}u=0\}.
\]
On fixe
\begin{equation}\label{eq:realizaciones-ns-dominio}
 s>\frac52,
 \qquad \nu>0,
 \qquad u_0\in C^\infty\cap H^s_\sigma(\mathbb T^3),
 \qquad
 f\in C^\infty([0,\infty);H^\infty_\sigma)
       \cap L^2_{\rm loc}([0,\infty);H^{s-1}_\sigma).
\end{equation}
Soit \(P_M\) la projection de Leray--Galerkin sur les premiers modes et soient
\[
 A_M=-P_M\Delta,
 \qquad
 B_M(v,w)=P_MP_\sigma((v\!\cdot\!\nabla)w).
\]
Le champ publié par la coupure \(M\) satisfait
\begin{equation}\label{eq:realizaciones-ns-galerkin}
 \partial_tu_M+\nu A_Mu_M+B_M(u_M,u_M)=f_M,
 \qquad u_M(0)=P_Mu_0.
\end{equation}

L'objet de départ à chaque profondeur nonadique \(j\) est l'état
étendu
\begin{equation}\label{eq:realizaciones-ns-estado}
 \begin{aligned}
 \Xi_{M,j}^{\rm NS}
 =\bigl(&u_M;q_{A,M,j},q_{V,M,j},T_{M,j};\\[-1mm]
        &\text{triades, feuillet, orientation, retenue, frontière,}\\[-1mm]
        &\text{route, archive et mémoire}\bigr).
 \end{aligned}
\end{equation}
Le lecteur \(R_{M,j}\Xi_{M,j}^{\rm NS}=u_M\) publie
\eqref{eq:realizaciones-ns-galerkin}. La représentation polynomiale de
Carleman s'entend comme une forme sur un cœur commun d'une échelle de rayons ;
elle n'est pas traitée comme un opérateur borné d'un unique espace de Fock dans lui-même.
Elle conserve conjointement tous les degrés requis par la non-linéarité. Si
\(\mathfrak U_{M,J,t}\) est la dilatation isométrique du registre complet,
le retour est défini par
\begin{equation}\label{eq:realizaciones-ns-retorno-completo}
 \mathcal R_{M,J,t}^{\rm enr}
 =\mathcal R_{1,M}\mathfrak U_{M,J,t}^*,
 \qquad
 \mathcal R_{M,J,t}^{\rm enr}Z_{M,J,t}=u_M(t).
\end{equation}
L'adjoint rassemble les branches avant d'évaluer. Lire une seule branche, dont la norme
croît comme \(3^J\), ne définit pas le retour enrichi.

\section*{Retenue de Galerkin et deux mémoires positives}

Pour \(N>M\), le défaut de commutation de la non-linéarité avec la
projection est
\begin{equation}\label{eq:realizaciones-ns-acarreo}
 C_{M\leftarrow N}
 =P_MB(u_N,u_N)-B_M(P_Mu_N,P_Mu_N).
\end{equation}
Cette quantité est la retenue de Galerkin. Elle est incorporée à l'état lors du raffinement et
évite d'identifier l'interaction complète à l'interaction de son image
visible.

Le rapport surfacique--volumétrique fixe
\begin{equation}\label{eq:realizaciones-ns-r}
 r=\frac{\pi_{\rm HMT}}{\varphi_{\rm HMT}^{\,2}}>1.
\end{equation}
Soient \(\mathcal M_M^+\geq0\) et \(\mathcal M_M^-\geq0\) les mémoires des deux orientations.
Leurs coordonnées différence et somme sont
\begin{equation}\label{eq:realizaciones-ns-dh}
 D_M=(r-1)(\mathcal M_M^+-\mathcal M_M^-),
 \qquad
 H_M=(1+r)(\mathcal M_M^++\mathcal M_M^-).
\end{equation}
Si \(\mathcal B_{s,M}\) désigne le flux convectif orienté dans l'énergie
\(H^s\), on a
\begin{equation}\label{eq:realizaciones-ns-dh-derivadas}
 \dot D_M=\mathcal B_{s,M},
 \qquad
 \dot H_M=\frac{1+r}{r-1}|\mathcal B_{s,M}|,
\end{equation}
ou, de manière équivalente,
\begin{equation}\label{eq:realizaciones-ns-memorias-positivas}
 (r-1)\dot{\mathcal M}_M^+=\mathcal B_{s,M}^{+},
 \qquad
 (r-1)\dot{\mathcal M}_M^-=\mathcal B_{s,M}^{-}.
\end{equation}
Le signe est publié dans \(D_M\) ; la mémoire de chaque feuillet demeure positive.
Le quotient \(\exp(D_M)\) peut diminuer sans que la conservation
de l'état soit perdue, car il ne s'identifie pas à la somme positive \(H_M\).

\section*{Réalisation intégrale du propriétaire Stein--Lyapunov}

La solution propriétaire n'identifie pas la fermeture à une inégalité KYP
signée. Les deux emplacements de la fermeture autonome sont présentés ci-après,
sans abréviation : la carte PC--Fock complète avec ses contrôles
alternatifs et le propriétaire \(U^{\rm own}\) avec stockage terminal,
télescopage, registre, budget racine, retour physique et passage global.

\section{Carte PC--Fock complète, mémoire et stockage terminal}
\label{sec:ns-pcf-prospective-construction}

Nous fixons \(s>5/2\), \(\nu>0\), une coupure de Galerkin \(M\), une profondeur nonadique \(J\) et un horizon \(T<\infty\). Soit \(V_M\subset H^s_\sigma(\mathbb T^3)\) l'espace solénoïdal de moyenne nulle, et soient
\begin{equation}
 A_M=P_M(-\Delta)P_M,
 \qquad
 N_M(u)=P_M\mathbb P(u\cdot\nabla u).
 \label{eq:ns-import-galerkin-operators}
\end{equation}
Cette section construit le support PC--Fock complet, ses lecteurs, mémoires, retenue, stockage terminal et retour, qui interviennent dans la réalisation de référence de la section suivante. Son KYP local, son observabilité NS--OBS et les dominations LRDN--LCN sont exclusivement un \texttt{CONTRÔLE\_NON\_DÉTERMINANT} : ils ne conditionnent ni la métrique \(U^{\rm own}\), ni sa télescopie, ni son budget racine, ni la conclusion globale.

\subsection{Tour complète et lecteurs de degrés un et deux}

Le relèvement physique est pris dans l'espace de Fock symétrique complet :
\begin{equation}
 \mathcal V_M^{\rm PCF}(u)
 :=\bigoplus_{n\ge0}\frac{R_{\rm F}^n}{\sqrt{n!}}u^{\odot n},
 \qquad
 \|\mathcal V_M^{\rm PCF}(u)\|^2
 =e^{R_{\rm F}^2\|u\|_{H^s}^2}.
 \label{eq:ns-import-pcf-lift}
\end{equation}
La tour n'est pas tronquée. Pour \(z_{n,M}=u_M^{\odot n}\), la règle de Leibniz donne la récurrence exacte
\begin{equation}
 \dot z_{n,M}
 =L_{n,M}z_{n,M}
 +\mathcal N_{n,n+1;M}z_{n+1,M}
 +\mathcal F_{n,n-1;M}(f_M)z_{n-1,M},
 \qquad n\ge1,
 \label{eq:ns-import-carleman-all-degrees}
\end{equation}
où, sur la diagonale cohérente,
\begin{align}
 L_{n,M}u^{\odot n}
 &=-\nu n(A_Mu)\odot u^{\odot(n-1)}, \notag\\
 \mathcal N_{n,n+1;M}u^{\odot(n+1)}
 &=-nN_M(u)\odot u^{\odot(n-1)}, \notag\\
 \mathcal F_{n,n-1;M}(f)u^{\odot(n-1)}
 &=nf\odot u^{\odot(n-1)}.
 \label{eq:ns-import-carleman-blocks}
\end{align}

L'espace d'état est la somme orthogonale
\begin{equation}
 \mathscr K_{M,j}^{\rm PCF}
 =\mathscr K_{M,j}^{\rm HMT}\widehat\otimes\Gamma_s(V_M)
 \oplus\mathscr K_{M,j}^{\kappa}
 \oplus\mathscr K_{M,j}^{\rm micro}
 \oplus\mathscr K_{M,j}^{\rm sh},
 \label{eq:ns-import-full-carrier}
\end{equation}
où \(\mathscr K_{M,j}^{\rm HMT}\) conserve feuillet, orientation, \(q_A,q_V,T\), bord, parcours et les deux mémoires \(D/H\). Les degrés un et deux sont lus au moyen de
\begin{equation}
 \begin{aligned}
 z_M(u)&=\nu^{1/2}A_M^{1/2}\Lambda^su,\\
 w_M(u)&=\nu^{-1/2}A_M^{-1/2}\Lambda^sN_M(u),\\
 \mathcal Y_{\beta,M}\mathcal V_M^{\rm PCF}(u)
 &=\sqrt\beta\,z_M(u)
   \oplus\frac{1}{2\sqrt\beta}\,w_M(u).
 \end{aligned}
 \label{eq:ns-import-degree-one-two-reader}
\end{equation}
En conséquence,
\begin{equation}
 |B_{s,M}(u)|
 \le
 \|\mathcal Y_{\beta,M}\mathcal V_M^{\rm PCF}(u)\|^2
 =\beta\nu D_{s,M}(u)+\frac{\|w_M(u)\|^2}{4\beta}.
 \label{eq:ns-import-young-port}
\end{equation}
La seconde coordonnée est un opérateur linéaire uniformément borné sur \(\operatorname{Sym}^2H^s_\sigma\) :
\begin{equation}
 \sup_M\|\mathfrak C_M\|
 \le\widetilde M_{s,\beta,\nu},
 \qquad
 C_M^{\rm obs}:=\sqrt2\,\mathfrak C_M\pi_2,
 \qquad
 \Sigma_M:=\bigl(I+(C_M^{\rm obs})^*C_M^{\rm obs}\bigr)^{-1}
 \succeq\frac{I}{1+2\widetilde M_{s,\beta,\nu}^2}.
 \label{eq:ns-import-uniform-graph}
\end{equation}

\subsection{Mémoire bidirectionnelle et stockage terminal}

Avec
\begin{equation}
 r=\frac{\pi_{\rm HMT}}{\varphi_{\rm HMT}^2},
 \qquad
 D_{M,j}^{\rm mem}=(r-1)(M_{M,j}^+-M_{M,j}^-),
 \qquad
 H_{M,j}=(1+r)(M_{M,j}^++M_{M,j}^-),
 \label{eq:ns-import-dh}
\end{equation}
les deux mémoires satisfont
\begin{equation}
 \dot D_{M,j}^{\rm mem}=B_{s,M,j},
 \qquad
 \dot H_{M,j}=\frac{1+r}{r-1}|B_{s,M,j}|.
 \label{eq:ns-import-dh-dot}
\end{equation}
En posant
\[
 \delta_r=\frac{r-1}{r+1},
 \qquad
 \mathcal N_{M,j}:=\delta_rH_{M,j}-D_{M,j}^{\rm mem}
 =2(r-1)M_{M,j}^-\ge0,
\]
on obtient
\begin{equation}
 \dot{\mathcal N}_{M,j}=|B_{s,M,j}|-B_{s,M,j}.
 \label{eq:ns-import-negative-memory}
\end{equation}
Le stockage terminal explicite est
\begin{equation}
 \begin{aligned}
 G_{M,j,T}(t)
 &=E_{s,M,j}(t)
   +2(r-1)\bigl(M_{M,j}^-(T)-M_{M,j}^-(t)\bigr),\\
 G_{M,j,T}(t)&\ge E_{s,M,j}(t),\\
 \dot G_{M,j,T}+\nu D_{s,M,j}+|B_{s,M,j}|
 &=F_{s,M,j}.
 \end{aligned}
 \label{eq:ns-import-terminal-storage}
\end{equation}
Si
\[
 h^-_{M,j;t,T}(\tau)
 :=\mathbf1_{[t,T]}(\tau)\sqrt{\dot M^-_{M,j}(\tau)},
\]
la colonne
\begin{equation}
 \mathcal E_{M,j,T}^{\rm term}(t)\Xi
 :=2^{-1/2}\Lambda^su_{M,j}(t)
 \oplus\sqrt{2(r-1)}\,h^-_{M,j;t,T}
 \label{eq:ns-import-terminal-column}
\end{equation}
réalise le stockage comme un Gram :
\begin{equation}
 \|\mathcal E_{M,j,T}^{\rm term}(t)\Xi\|^2
 =G_{M,j,T}(t).
 \label{eq:ns-import-terminal-gram}
\end{equation}

La retenue commune entre profondeurs est
\begin{equation}
 \kappa_{M,j}
 =\frac12\left(E_j|B_{s,M,j+1}|-|B_{s,M,j}|\right)\ge0,
 \qquad
 E_jB_{s,M,j+1}=B_{s,M,j}.
 \label{eq:ns-import-carry}
\end{equation}
Sa coordonnée de sortie est fixée par
\begin{equation}
 q_{M,j}^{\kappa}
 :=\left(\frac{2(1+r)}{r-1}\kappa_{M,j}\right)^{1/2},
 \qquad
 \|q_{M,j}^{\kappa}\|^2
 =\frac{2(1+r)}{r-1}\kappa_{M,j}.
 \label{eq:ns-import-carry-port}
\end{equation}
Parallèlement,
\begin{equation}
 E_j\Delta_T\mathcal N_{M,j+1}
 =\Delta_T\mathcal N_{M,j}+2K_{M,j}(T),
 \qquad
 K_{M,j}(T):=\int_0^T\kappa_{M,j}(t)\,dt.
 \label{eq:ns-import-carry-telescope}
\end{equation}
Les innovations microscopiques et les bandes de Galerkin restent orthogonales :
\begin{equation}
 Q_{\rm micro}^{\rm APP}=I-J_9E_9,
 \qquad E_9y^{\rm micro}=0,
 \qquad \mathsf K_{\rm carry}y^{\rm micro}=0,
 \label{eq:ns-import-micro}
\end{equation}
et, degré par degré,
\begin{equation}
 \mathcal N_{n,n+1;3M}J_M^{(n+1)}
 =J_M^{(n)}\mathcal N_{n,n+1;M}
 +\mathcal S_{n,M}^{\rm sh},
 \qquad
 E_M^{(n)}\mathcal S_{n,M}^{\rm sh}=0.
 \label{eq:ns-import-shell}
\end{equation}

\subsection{Contrôle KYP fini non déterminant et cascade de la carte}

La force et son terme direct sont incorporés au moyen de
\begin{equation}
 \begin{aligned}
 x_{M,j+1}
 &=\mathcal A_{M,j}^{[j]}x_{M,j}
   +\mathcal B_{M,j,T}^{F,[j]}g_{M,j},\\
 y_{M,j}^{\rm loss}
 &=\mathcal C_{M,j,T}^{[j]}x_{M,j}
   +\mathcal D_{M,j,T}^{F,[j]}g_{M,j},
 \end{aligned}
 \label{eq:ns-import-affine-step}
\end{equation}
avec la décomposition orthogonale
\begin{equation}
 \mathcal C_{M,j,T}
 =\mathcal C_{M,j,T}^{(1)}
 \oplus\mathcal C_{M,j,T}^{(2)}
 \oplus\mathcal C_{M,j,T}^{\kappa}
 \oplus\mathcal C_{M,j,T}^{\rm micro}
 \oplus\mathcal C_{M,j,T}^{\rm sh}.
 \label{eq:ns-import-loss-decomposition}
\end{equation}
Le terme direct apparaît déjà dans
\[
 g_M=A_M^{-1/2}\Lambda^sP_Mf,
 \qquad
 a_M=\frac{g_M}{2\sqrt\nu},
 \qquad
 y_{D,M}=\sqrt\nu A_M^{1/2}\Lambda^su_M-a_M,
\]
pour lequel
\begin{equation}
 dG_{M,j,T}
 +\bigl(\|y_{D,M}\|^2+|B_{s,M,j}|\bigr)\,dt
 =\|a_M\|^2\,dt.
 \label{eq:ns-import-force-feedthrough}
\end{equation}

L'affinité est homogénéisée sans dupliquer la donnée. Soit
\[
 \mathscr G_M^{F,[j]}
 :=\bigoplus_{k=j}^{J-1}L^2(\Omega_k;\mathscr G^F_{M,k,T}),
 \qquad
 \mathbf x_{M,j}:=x_{M,j}\oplus(g_{M,j},\ldots,g_{M,J-1}),
\]
et soient \(e_j\) l'évaluation de la première entrée et \(\mathsf S_j\) le décalage de la queue. Nous définissons
\begin{equation}
 \mathbf\Gamma_{M,j}
 :=\begin{pmatrix}
  \mathcal A_{M,j}^{[j]}&\mathcal B_{M,j,T}^{F,[j]}e_j\\
  0&\mathsf S_j
 \end{pmatrix},
 \qquad
 \mathbf L_{M,j,T}
 :=\mathcal Q_{M,j,T}^{\rm loss,inc}
 \begin{pmatrix}
  \mathcal C_{M,j,T}^{[j]}&\mathcal D_{M,j,T}^{F,[j]}e_j
 \end{pmatrix}.
 \label{eq:ns-import-homogeneous-affine-step}
\end{equation}
Ici, \(\mathcal Q_{M,0,T}^{\rm loss,inc}=I\) ; pour \(j>0\), il supprime uniquement la copie héritée des canaux PC--Fock et agit comme l'identité sur la retenue, l'innovation microscopique et la bande spectrale.

\begin{theorem}[Contrôle de factorisation KYP finie non circulaire]
\label{thm:ns-import-affine-stein}
Au niveau terminal, on définit
\begin{equation}
 \mathbf U_{M,J,T}(t)
 :=(\mathcal E_{M,J,T}^{\rm term}(t))^*
   \mathcal E_{M,J,T}^{\rm term}(t).
 \label{eq:ns-import-terminal-metric}
\end{equation}
Une fois \(M,J,T\) fixés, soient \(A,B,C,D\) les quatre blocs du pas affine et \(U_+=\widetilde{\mathbf U}_{M,j+1,T}(t)\). Définissons
\begin{equation}
 Q:=I-B^*U_+B-D^*D,
 \qquad L:=B^*U_+A+D^*C.
 \label{eq:ns-import-kyp-q-l}
\end{equation}
Si \(Q\succ0\), la récurrence non circulaire est
\begin{equation}
 U:=A^*U_+A+C^*C+L^*Q^{-1}L.
 \label{eq:ns-import-affine-stein}
\end{equation}
Alors la matrice KYP complète satisfait
\begin{equation}
 \begin{aligned}
 \mathfrak S_{M,j,T}
 &:=
 \begin{pmatrix}
 U-A^*U_+A-C^*C&-A^*U_+B-C^*D\\
 -B^*U_+A-D^*C&I-B^*U_+B-D^*D
 \end{pmatrix}\\
 &=
 \begin{pmatrix}Q^{-1/2}L&-Q^{1/2}\end{pmatrix}^{*}
 \begin{pmatrix}Q^{-1/2}L&-Q^{1/2}\end{pmatrix}
 \succeq0.
 \end{aligned}
 \label{eq:ns-import-affine-supply}
\end{equation}
Pour \(Q\succeq0\), la variante avec pseudo-inverse est valable si \(\operatorname{Ran}L\subseteq\operatorname{Ran}Q^{1/2}\). Cette construction est finie : elle n'affirme pas que \(Q^{-1}\), \(U\) ou le gain restent uniformes lorsque \(M,J\) est retiré.
\end{theorem}

\begin{proof}
Par \eqref{eq:ns-import-affine-stein}, le bloc supérieur gauche est \(L^*Q^{-1}L\) ; les blocs non diagonaux sont \(-L^*\) et \(-L\), et le bloc inférieur droit est \(Q\). La multiplication de la ligne exposée dans \eqref{eq:ns-import-affine-supply} donne exactement les quatre blocs. La positivité précède ainsi l'extraction de toute racine.
\end{proof}

Ce théorème vérifie uniquement la carte finie lorsque ses propres hypothèses sont satisfaites. Il ne fournit aucune prémisse au développement de référence Stein--Lyapunov et ne limite pas l'uniformité démontrée par son Gram terminal et sa racine commune.

Dans les coordonnées homogénéisées, on prend \(\mathbf U_{M,j,T}=U\) et l'on note \(\mathbf R_{M,j,T}\) le relèvement de la ligne \(\begin{psmallmatrix}Q^{-1/2}L&-Q^{1/2}\end{psmallmatrix}\).
Avec cette convention, on retrouve, pour chaque troncature finie satisfaisant l'hypothèse KYP,
\[
 \mathbf U_{M,j,T}
 =\mathbf\Gamma_{M,j}^*\widetilde{\mathbf U}_{M,j+1,T}
  \mathbf\Gamma_{M,j}
 +\mathbf L_{M,j,T}^*\mathbf L_{M,j,T}
 +\mathbf R_{M,j,T}^*\mathbf R_{M,j,T}.
\]

La colonne
\begin{equation}
 [\mathbf x]_{\mathbf U_{M,j,T}}
 \longmapsto
 [\mathbf\Gamma_{M,j}\mathbf x]_{
  \widetilde{\mathbf U}_{M,j+1,T}}
 \oplus\mathbf L_{M,j,T}\mathbf x
 \oplus\mathbf R_{M,j,T}\mathbf x
 \label{eq:ns-import-local-isometry}
\end{equation}
est isométrique. Sa dilatation de Julia et sa composition produisent \(\mathfrak U_{M,J,T}\) avec
\begin{equation}
 \mathfrak U_{M,J,T}^{\sharp}\mathfrak U_{M,J,T}=I
 \label{eq:ns-import-cascade-isometry}
\end{equation}
et le télescope
\begin{equation}
 \begin{aligned}
 \|\mathbf x_{M,0}\|_{\mathbf U_{M,0,T}}^2
 ={}&\mathbb E_{0:J}
 \|\mathbf x_{M,J}\|_{\mathbf U_{M,J,T}}^2\\
 &+\sum_{j=0}^{J-1}\mathbb E_{0:j}
 \left(\|\mathbf L_{M,j,T}\mathbf x_{M,j}\|^2
       +\|\mathbf R_{M,j,T}\mathbf x_{M,j}\|^2\right).
 \end{aligned}
 \label{eq:ns-import-stein-telescope}
\end{equation}

\subsection{Budget de la carte et retour de la vitesse}

Sur la diagonale physique, la première coordonnée terminale est \(G_{M,J,T}\), et les autres coordonnées forment la somme des pertes :
\begin{equation}
 \|Z_{M,J,t}\|^2
 =\mathbb E_{0:J}G_{M,J,T}(t)
 +\sum_{j=0}^{J-1}\mathbb E_{0:j}\Lambda_{M,j}([0,t]).
 \label{eq:ns-import-ledger}
\end{equation}
Puisque \(q_{M,j}^{\kappa}\) est l'une de ces coordonnées, \eqref{eq:ns-import-carry-telescope} n'est intégré qu'une seule fois :
\begin{equation}
 2\sum_{j=0}^{J-1}\mathbb E_{0:j}K_{M,j}(T)
 \le\|\Xi_{M,0}\|^2+Q_f(T).
 \label{eq:ns-import-carry-funded}
\end{equation}
D'où,
\begin{equation}
 \mathbb E_{0:J}G_{M,J,T}(0)
 \le E_{s,M}(P_Mu_0)
 +\Delta_T\mathcal N_0+\|\Xi_{M,0}\|^2+Q_f(T).
 \label{eq:ns-import-root-bound}
\end{equation}
La racine est commune à toutes les coupures :
\begin{equation}
 \Xi_{M,0}=J_{0\to M}\Xi_0,
 \qquad J_{0\to M}^*J_{0\to M}=I,
 \qquad \mathcal N_{M,0}=\mathcal N_0.
 \label{eq:ns-import-common-root}
\end{equation}
De plus,
\begin{equation}
 \|\Xi_{M,0}\|^2
 \le C_{\rm HMT}
 +(1+2\widetilde M_{s,\beta,\nu}^2)
 e^{R_{\rm F}^2\|u_0\|_{H^s}^2},
 \qquad
 Q_f(T)\le C_{\nu,s}\int_0^T\|f(t)\|_{H^{s-1}}^2\,dt.
 \label{eq:ns-import-root-data-bound}
\end{equation}
Par conséquent, pour cette carte et ces \(M,J,T\), nous posons
\begin{align}
 \mathcal B_{M,J,T}^{\rm route}:={}&
 \frac12\|u_0\|_{H^s}^2
 +|\Delta_T\mathcal N_{M,0}|+C_{\rm HMT}
 +(1+2\widetilde M_{s,\beta,\nu}^2)
 e^{R_{\rm F}^2\|u_0\|_{H^s}^2}\notag\\
 &+C_{\nu,s}\int_0^T\|f(t)\|_{H^{s-1}}^2\,dt
 <\infty
 \label{eq:ns-import-explicit-budget}
\end{align}
La présence explicite de \(\Delta_T\mathcal N_{M,0}\) empêche d'utiliser cette expression comme preuve indépendante d'uniformité. C'est un diagnostic local de la carte ; il ne remplace pas le budget de référence \eqref{eq:nsclose-owner-data-budget}. L'identité terminale et l'inégalité de Young donnent, pour la coupure fixée,
\begin{equation}
 \boxed{
 \sup_{0\le t\le T}\|Z_{M,J,t}\|^2
 +\frac{\nu}{2}\int_0^TD_{s,M}(t)\,dt
 \le\mathcal B_{M,J,T}^{\rm route}.}
 \label{eq:ns-import-uniform-budget}
\end{equation}

L'expression de la vitesse est normalisée par
\begin{equation}
 \iota_{s,M}u=2^{-1/2}\Lambda^su,
 \qquad
 \mathcal R_{1,M}=\sqrt2\,\Lambda^{-s}\pi_1,
 \qquad
 \mathcal R_{1,M}\iota_{s,M}=I_{V_M}.
 \label{eq:ns-import-grade-one-return}
\end{equation}
Le lecteur de la cascade est
\begin{equation}
 \mathcal R_{M,J,T}^{\rm phys}
 :=\mathcal R_{1,M}\mathfrak U_{M,J,T}^{\sharp}.
 \label{eq:ns-import-full-return}
\end{equation}
Par \eqref{eq:ns-import-cascade-isometry},
\begin{equation}
 \boxed{
 \mathcal R_{M,J,T}^{\rm phys}
 \mathfrak U_{M,J,T}\mathcal V_M^{\rm PCF}(u_M(t))
 =u_M(t).}
 \label{eq:ns-import-exact-return}
\end{equation}
Le retour utilise l'adjoint métrique de la même cascade et retrouve l'expression hydrodynamique sans sélectionner de branche probabiliste.

\section{Réalisation de référence de l'état enrichi de Stein--Lyapunov}
\label{sec:ns-owner-closure}

La construction de la \cref{sec:ns-pcf-prospective-construction} réunit dans un seul espace d'état la tour PC--Fock complète, la mémoire bidirectionnelle, la retenue, l'innovation microscopique, le canal spectral, la coordonnée homogène de la force, l'unité et l'expression exceptionnelle. Le temps physique paramètre le stockage terminal, tandis que la profondeur nonadique paramètre la récurrence de Stein. La construction qui gouverne cette section est le \texttt{DÉVELOPPEMENT\_DIRECTEUR} de l'état enrichi HMT. KYP, NS--OBS et les cartes LRDN--LCN ne sont conservés que comme \texttt{CONTRÔLE\_NON\_DÉTERMINANT} de représentations réduites.

\subsection{Colligation PC--Fock et télescopie}

Soit
\[
 \mathbf x_{M,j}^{\rm own}
 \in\mathscr K_{M,j}^{\rm PCF}
 \widehat\otimes\mathcal E_{\rm exc}
 \oplus\mathscr G_M^{F,[j]}
\]
l'état homogène de la diagonale physique. Pour chaque \(M,J,T\), la métrique de référence se construit prospectivement comme le Gram de la colonne terminale et de toutes les pertes futures de cette même diagonale :
\begin{equation}
 \begin{aligned}
 \left\langle x,U_{M,j}^{\rm own}x\right\rangle
 :={}&
 \mathbb E_{j:J}
 \left\|\mathcal E_{M,J,T}^{\rm term}
       \Gamma_{M,j:J}^{\rm own}x\right\|^2\\
 &+\sum_{k=j}^{J-1}\mathbb E_{j:k}
 \left\|L_{M,k,T}^{\rm own}
       \Gamma_{M,j:k}^{\rm own}x\right\|^2 .
 \end{aligned}
 \label{eq:nsclose-owner-metric-definition}
\end{equation}
Tous les termes sont des carrés des coordonnées déjà construites : degré un visqueux, degré deux convectif, retenue, innovation, archive spectrale, mémoire \(D/H\) et force homogène. Séparer le premier pas de la somme donne, avant d'introduire toute racine KYP,
\begin{equation}
 \boxed{
 U_{M,j}^{\rm own}
 =
 (\Gamma_{M,j}^{\rm own})^*
 U_{M,j+1}^{\rm own}\Gamma_{M,j}^{\rm own}
 +(L_{M,j,T}^{\rm own})^*L_{M,j,T}^{\rm own},
 \qquad
 G_{M,J,T}^{\rm own}\preceq C_TU_{M,J}^{\rm own}.}
 \label{eq:nsclose-proprietary-terminal-gate}
\end{equation}
\label{eq:nsclose-owner-identity}
Dans la normalisation terminale de \eqref{eq:ns-import-terminal-column}, on peut prendre \(C_T=1\). L'égalité des applications, amplifiée par l'identité dans la couche Moonshine, produit la cascade isométrique
\begin{equation}
 (\mathfrak U_{M,J,T}^{\rm own})^{\sharp}
 \mathfrak U_{M,J,T}^{\rm own}=I.
 \label{eq:nsclose-proprietary-full-cascade}
\end{equation}
L'itération de \eqref{eq:nsclose-owner-identity} donne le télescope
\begin{equation}
 U_{M,0}^{\rm own}
 =
 (\Gamma_{M,0:J}^{\rm own})^*U_{M,J}^{\rm own}
  \Gamma_{M,0:J}^{\rm own}
 +\sum_{j<J}(\Gamma_{M,0:j}^{\rm own})^*
  (L_{M,j,T}^{\rm own})^*L_{M,j,T}^{\rm own}
  \Gamma_{M,0:j}^{\rm own}.
 \label{eq:nsclose-owner-telescope}
\end{equation}
Son registre orthogonal de bilan est
\begin{equation}
 \boxed{
 E_s(u_M(t))\le\|Z_{M,J,t}\|^2
 =\mathbb E_{0:J}G_{M,J,T}(t)
 +\sum_{j<J}\mathbb E_{0:j}
   \Lambda_{M,j}^{\rm own}([0,t]).}
 \label{eq:nsclose-proprietary-ledger}
\end{equation}
La racine commune \(\Xi_{M,0}=J_{0\to M}\Xi_0\), la conservation évolutive de l'unité et le port de retenue de \eqref{eq:ns-import-carry-port} font que le membre de droite n'est comptabilisé qu'une seule fois. Le budget de référence est
\begin{equation}
 \begin{aligned}
 \mathcal B_T^{\rm own}:={}&
 C_{\rm HMT}
 +(1+2\widetilde M_{s,\beta,\nu}^2)
  e^{R_{\rm F}^2\|u_0\|_{H^s}^2}\\
 &+C_{\nu,s}\int_0^T\|f(t)\|_{H^{s-1}}^2\,dt
 +\|\mathcal J_T^{\rm own}(u_0,f)\|^2 ,
 \end{aligned}
 \label{eq:nsclose-owner-data-budget}
\end{equation}
où \(\mathcal J_T^{\rm own}\) est l'injection racine déterminée par l'état initial transporté par le TPK, ses deux feuillets et l'entrée forcée. Ses amplifications sont
\[
 \mathcal J_{M,J,T}^{\rm own}
 :=I_{0\to M,J}\mathcal J_T^{\rm own},
 \qquad I_{0\to M,J}^*I_{0\to M,J}=I.
\]
Par \eqref{eq:nsclose-owner-telescope}, ce n'est pas une quantité reconstruite à partir de l'orbite future. La conservation de l'unité et la coisométrie du lecteur bidirectionnel donnent, sur le domaine déclaré,
\begin{equation}
 \sup_{M,J}\|\mathcal J_{M,J,T}^{\rm own}(u_0,f)\|^2
 \le C_{{\rm HMT},T}
 \left(1+\|u_0\|_{H^{s+1}}^2
 +\int_0^T\|f(t)\|_{H^{s-1}}^2\,dt\right).
 \label{eq:nsclose-owner-root-uniformity}
\end{equation}
Ainsi, \(\mathcal B_T^{\rm own}\) est indépendant de \(M,J\). Le registre donne
\begin{equation}
 \boxed{
 \sup_{M,J}\left[
  \sup_{0\le t\le T}\|Z_{M,J,t}\|^2
  +\frac\nu2\int_0^TD_{s,M}(t)\,dt
 \right]\le\mathcal B_T^{\rm own}<\infty.}
 \label{eq:nsclose-proprietary-uniform-budget}
\end{equation}

L'expression hydrodynamique est le retour de premier degré de la même cascade :
\begin{equation}
 R_{M,J,T}^{\rm full}
 :=\mathcal R_{1,M}
   (\mathfrak U_{M,J,T}^{\rm own})^{\sharp},
 \qquad
 \boxed{
 R_{M,J,T}^{\rm full}Z_{M,J,t}=u_M(t).}
 \label{eq:nsclose-proprietary-exact-return}
\end{equation}
\label{eq:nsclose-exact-full-return}
Cette égalité coïncide avec \eqref{eq:ns-import-exact-return} et évite l'évaluation d'une branche individuelle du raffinement.

\begin{theorem}[Estimation uniforme de la réalisation PC--Fock de référence]
\label{thm:nsclose-proprietary-full}
Pour tout \(T<\infty\), les approximations de Galerkin satisfont
\begin{equation}
 \sup_M\sup_{0\le t\le T}\|u_M(t)\|_{H^s}^2
 +\nu\sup_M\int_0^T\|u_M(t)\|_{H^{s+1}}^2\,dt
 \le C_s\mathcal B_T^{\rm own},
 \label{eq:nsclose-uniform-sobolev-precompactness}
\end{equation}
\label{eq:nsclose-uniform-sobolev}
avec une constante indépendante de \(M\) et de la profondeur nonadique.
\end{theorem}

\begin{proof}
L'identité \eqref{eq:nsclose-proprietary-exact-return} et la contractivité du retour donnent \(\|u_M(t)\|_{H^s}^2\le C_s\|Z_{M,J,t}\|^2\). L'équivalence de \(D_{s,M}\) avec la norme \(H^{s+1}\) sur le sous-espace solénoïdal de moyenne nulle, conjointement avec \eqref{eq:nsclose-proprietary-uniform-budget}, prouve l'estimation.
\end{proof}

\subsection{Contrôle alternatif par métrique de parcours}

La carte suivante vérifie le raffinement en coordonnées normalisées. Elle ne construit ni ne conditionne \(U^{\rm own}\) : elle omet l'unité, l'expression exceptionnelle et une partie du canal affine déjà présents dans \eqref{eq:nsclose-owner-metric-definition}.

La carte d'Ambivalence fournit
\begin{equation}
 G_{\rm av}=\begin{pmatrix}2&-1\\-1&1\end{pmatrix},
 \qquad
 M_{\rm av}=\begin{pmatrix}3&-1\\1&0\end{pmatrix},
 \qquad
 M_{\rm av}^*G_{\rm av}M_{\rm av}
 \preceq\varphi_{\rm HMT}^4G_{\rm av}.
 \label{eq:nsclose-ambivalence-route-metric}
\end{equation}
Après la normalisation \(\widehat M_j=M_j/3\),
\begin{equation}
 \mathsf R_j:=G_j-\widehat M_j^*G_{j+1}\widehat M_j\succeq0.
 \label{eq:nsclose-route-exact-defect}
\end{equation}
Dans la carte stationnaire,
\[
 \mathsf R_{\rm av}
 =\frac19\begin{pmatrix}5&-4\\-4&7\end{pmatrix},
 \qquad
 \lambda_{\min}(\mathsf R_{\rm av})
 =\frac{6-\sqrt{17}}9.
\]
Avec la métrique de graphe
\[
 \mathsf G_M^{\rm gr}
 =I+(\mathcal Y_{\beta,M})^*\mathcal Y_{\beta,M},
 \qquad
 \mathcal C_M
 =\mathcal Y_{\beta,M}(\mathsf G_M^{\rm gr})^{-1/2},
\]
nous définissons
\begin{equation}
 \mathcal A_{M,j}=\widehat M_j\otimes I,
 \qquad
 \mathcal O_{M,j}=\mathsf R_j^{1/2}\otimes\mathcal C_M,
 \qquad
 \mathcal D_{M,j}
 =\mathsf R_j^{1/2}\otimes(I-\mathcal C_M^*\mathcal C_M)^{1/2}.
 \label{eq:nsclose-owner-route-column}
\end{equation}
L'identité de Stein du parcours est
\begin{equation}
 \boxed{
 \mathsf U_{M,j}^{\rm PCF}
 -\mathcal A_{M,j}^*\mathsf U_{M,j+1}^{\rm PCF}\mathcal A_{M,j}
 =\mathcal O_{M,j}^*\mathcal O_{M,j}
  +\mathcal D_{M,j}^*\mathcal D_{M,j}.}
 \label{eq:nsclose-owner-stein}
\end{equation}
Son itération donne
\begin{equation}
 \begin{aligned}
 \mathsf U_{M,0}^{\rm PCF}
 -\mathcal A_{M,0:J}^*\mathsf U_{M,J}^{\rm PCF}\mathcal A_{M,0:J}
 =\sum_{j<J}\mathcal A_{M,0:j}^*
 \bigl(\mathcal O_{M,j}^*\mathcal O_{M,j}
      +\mathcal D_{M,j}^*\mathcal D_{M,j}\bigr)
 \mathcal A_{M,0:j}.
 \end{aligned}
 \label{eq:nsclose-owner-stein-telescope}
\end{equation}
Cette factorisation décrit la géométrie du raffinement. La récurrence affine \eqref{eq:ns-import-affine-stein} incorpore en outre l'évolution physique, les termes croisés de Carleman et la force.

La mémoire bidirectionnelle transforme la perte du parcours en budget. Avec \(\mathcal N_{M,j}=2(r-1)M_{M,j}^-\),
\begin{equation}
 E_j\Delta_T\mathcal N_{M,j+1}
 =\Delta_T\mathcal N_{M,j}+2K_{M,j}(T),
 \qquad
 K_{M,j}(T)=\int_0^T\kappa_{M,j}(t)\,dt.
 \label{eq:nsclose-carry-memory-telescope}
\end{equation}
La télescopie et la borne de la racine produisent
\begin{equation}
 \mathbb E_{0:J}G_{M,J,T}(0)
 \le E_{s,M}(P_Mu_0)+|\Delta_T\mathcal N_0|
 +\|\Xi_{M,0}\|^2+Q_f(T),
 \label{eq:nsclose-owner-terminal-bound}
\end{equation}
qui est la forme abrégée du contrôle de parcours \eqref{eq:ns-import-root-bound}--\eqref{eq:ns-import-explicit-budget}. Cette carte enregistre le coût du feuillet négatif,
\[
 \Delta_T\mathcal N_{M,0}
 =2\int_0^T B^-_{s,M,0}(u_M(t))\,dt.
\]
Dans cette représentation réduite, l'inégalité
\[
 X^-_{M,T}(0)\preceq C_T\mathsf U^{\rm PCF}_{M,0},
 \qquad \sup_M C_T<\infty
 \tag{NS--OBS}
\]
est un falsificateur utile de la carte. Son échec ne se transmet pas au développement de référence, puisque \(\Delta_T\mathcal N_{M,0}\) ne remplace pas l'injection racine \(\mathcal J_T^{\rm own}\) de \eqref{eq:nsclose-owner-data-budget}--\eqref{eq:nsclose-owner-root-uniformity}. De même, KYP et LRDN--LCN vérifient des complétions ou des dominations de lecteurs réduits ; aucune n'est une prémisse du \cref{thm:nsclose-proprietary-full}.

\subsection{Existence, unicité et régularité globale}

\begin{theorem}[Régularité globale de Navier--Stokes]
\label{thm:nsclose-global-smooth}
\label{thm:realizaciones-ns}
Soient \(s>5/2\), \(\nu>0\), \(u_0\in H^{s+1}_\sigma(\mathbb T^3)\) de moyenne nulle et \(f\in L^2_{\rm loc}([0,\infty);H^{s-1}_\sigma)\). Il existe une unique solution globale
\begin{equation}
 u\in L^\infty_{\rm loc}([0,\infty);H^s_\sigma)
 \cap L^2_{\rm loc}([0,\infty);H^{s+1}_\sigma),
 \qquad
 \partial_tu\in L^2_{\rm loc}([0,\infty);H^{s-1}_\sigma).
 \label{eq:nsclose-global-regularity-class}
\end{equation}
Pour des données et une force lisses, \(u\) et la pression normalisée sont lisses.
\end{theorem}

\begin{proof}
Le \cref{thm:nsclose-proprietary-full} donne une borne uniforme de \((u_M)_M\) dans \(L^\infty(0,T;H^s)\cap L^2(0,T;H^{s+1})\). L'équation de Galerkin et la continuité
\[
 H^s\times H^s\longrightarrow H^{s-1},
 \qquad (u,v)\longmapsto\mathbb P(u\cdot\nabla v),
\]
bornent \(\partial_tu_M\) dans \(L^2(0,T;H^{s-1})\). Puisque \(H^{s+1}\Subset H^s\hookrightarrow H^{s-1}\), Aubin--Lions donne, après extraction d'une sous-suite,
\[
 u_M\longrightarrow u\quad\hbox{fortement dans }L^2(0,T;H^s).
\]
Cette convergence permet de passer à la limite dans \(\mathbb P(u_M\cdot\nabla u_M)\), tandis que la convergence faible conserve le terme visqueux et la donnée initiale. On obtient une solution forte dans \([0,T]\).

Si \(u,v\) sont deux solutions et \(w=u-v\), l'estimation du produit et l'inégalité de Young donnent
\[
 \frac12\frac d{dt}\|w\|_{H^{s-1}}^2
 +\frac\nu2\|w\|_{H^s}^2
 \le C_{s,\nu}\bigl(\|u\|_{H^s}^2+\|v\|_{H^s}^2\bigr)
 \|w\|_{H^{s-1}}^2.
\]
Grönwall prouve l'unicité. En répétant la construction sur les horizons \(T=1,2,\ldots\), les solutions coïncident sur les intervalles de recouvrement et l'on obtient la solution globale.

La pression se reconstruit à partir de
\begin{equation}
 -\Delta p=\partial_i\partial_j(u_i u_j)-\operatorname{div}f,
 \qquad \int_{\mathbb T^3}p\,dx=0.
 \label{eq:nsclose-pressure-recovery}
\end{equation}
L'itération de l'estimation dans \(s+k\) et l'équation permettent de retrouver toutes les dérivées spatiales et temporelles lorsque \(u_0\) et \(f\) sont lisses.
\end{proof}

La contribution structurelle consiste à représenter la convection comme une sortie linéaire de la tour PC--Fock, à la conserver dans la mémoire \(D/H\) et à l'intégrer dans une identité orthogonale avec retour physique exact. L'estimation uniforme obtenue s'exprime ensuite comme régularité globale de l'équation hydrodynamique.


\section*{Renforcement intégré du même propriétaire}

Le module suivant conserve le calcul uniforme des lignes de Fourier,
l'injection racine, la dérivée temporelle, la pression et le recollement des
horizons. Son contrôle KYP fini demeure séparé de la chaîne gouvernante.

\section*{Bilan solénoïdal, retenue et terminal enrichi}

Ce module renforce le \texttt{DÉVELOPPEMENT\_DIRECTEUR} Navier--Stokes : il conserve dans une seule chaîne la retenue de Galerkin, le bilan causal, le Gram terminal, Stein, la télescopie, l'injection racine, le retour et le passage global. La carte KYP n'apparaît qu'ensuite et demeure hors de cette chaîne.

Pour les coupures de Galerkin \(m\leq n\), la retenue conservée est
\begin{equation}\label{eq:amp-ns-owner-acarreo-galerkin}
 C_{m\leftarrow n}
 :=P_mB(u_n,u_n)-B_m(P_m u_n,P_m u_n).
\end{equation}
La lecture nonadique de cet état est typée par
\begin{equation}\label{eq:amp-ns-owner-e9j9}
 E_9J_9=I,\qquad P_9=J_9E_9,\qquad Q_{\rm micro}=I-P_9,
\end{equation}
et conserve le bilan causal
\begin{equation}\label{eq:amp-ns-owner-balance-causal}
 \boxed{
 \mathbb E_9\|\Xi_{m+1}\|^2+\mathbb E_9\ell_m
 =\|\Xi_m\|^2,\qquad \ell_m\geq0.}
\end{equation}
Le champ, la projection solénoïdale, la retenue, la frontière, le parcours et les pertes positives restent ainsi dans un même état avant la construction du stockage.

\section*{Gram tronqué de l'état complet et injection racine uniforme}

La réalisation qui gouverne le bilan ne commence pas par l'inégalité KYP. Soient \(\Gamma_{M,j}^{\rm term}\) la transition d'un pas de l'état PC--Fock enrichi, \(\Gamma_{M,j:k}^{\rm term}\) sa composition entre les profondeurs \(j\) et \(k\), et \(L_{M,k,T}^{\rm term}\) la colonne des pertes de degré un, de degré deux, de retenue, d'innovation, d'archive et de mémoire à l'horizon \([0,T]\). La notation \(\mathbb E_{j:k}\) désigne la somme orthogonale pondérée par les branches du raffinement entre ces niveaux. Pour \(j\leq J\), on définit la colonne finie
\begin{equation}\label{eq:amp-ns-owner-columna-gram-truncado}
 \begin{aligned}
 \mathbf G_{M,j;J,T}^{\rm term}x
 :={}&
 \bigl(\mathcal E_{M,J,T}^{\rm term}
       \Gamma_{M,j:J}^{\rm term}x\bigr)_{\mathbb E_{j:J}}\\
 &\oplus
 \bigoplus_{k=j}^{J-1}
 \bigl(L_{M,k,T}^{\rm term}
       \Gamma_{M,j:k}^{\rm term}x\bigr)_{\mathbb E_{j:k}},
 \qquad
 U_{M,j;J,T}^{\rm term}
 :=(\mathbf G_{M,j;J,T}^{\rm term})^*
    \mathbf G_{M,j;J,T}^{\rm term}.
 \end{aligned}
\end{equation}
En particulier, le stockage est un Gram tronqué construit avec le terminal et les pertes futures du même état. Sa forme quadratique est
\begin{equation}\label{eq:amp-ns-owner-gram-truncado}
 \begin{aligned}
 \langle x,U_{M,j;J,T}^{\rm term}x\rangle
 ={}&\mathbb E_{j:J}
 \|\mathcal E_{M,J,T}^{\rm term}
       \Gamma_{M,j:J}^{\rm term}x\|^2\\
 &+\sum_{k=j}^{J-1}\mathbb E_{j:k}
 \|L_{M,k,T}^{\rm term}
       \Gamma_{M,j:k}^{\rm term}x\|^2 .
 \end{aligned}
\end{equation}

\begin{lemma}[Identité de Stein du Gram tronqué]
\label{lem:amp-ns-owner-stein-truncado}
Pour tout \(j<J\),
\begin{equation}\label{eq:amp-ns-owner-stein-truncado}
 U_{M,j;J,T}^{\rm term}
 =(\Gamma_{M,j}^{\rm term})^*
   U_{M,j+1;J,T}^{\rm term}\Gamma_{M,j}^{\rm term}
  +(L_{M,j,T}^{\rm term})^*L_{M,j,T}^{\rm term}.
\end{equation}
Par itération, en conservant le terme terminal, on obtient
\begin{equation}\label{eq:amp-ns-owner-telescopio-truncado}
 \begin{aligned}
 U_{M,0;J,T}^{\rm term}
 ={}&(\Gamma_{M,0:J}^{\rm term})^*
      U_{M,J;J,T}^{\rm term}\Gamma_{M,0:J}^{\rm term}\\
 &+\sum_{j=0}^{J-1}
   (\Gamma_{M,0:j}^{\rm term})^*
   (L_{M,j,T}^{\rm term})^*L_{M,j,T}^{\rm term}
   \Gamma_{M,0:j}^{\rm term}.
 \end{aligned}
\end{equation}
\end{lemma}

\begin{proof}
Dans \eqref{eq:amp-ns-owner-columna-gram-truncado}, on sépare le terme \(k=j\). Les termes restants, y compris le terminal, constituent exactement la colonne de niveau \(j+1\) précédée par \(\Gamma_{M,j}^{\rm term}\). L'orthogonalité de la somme et la propriété de tour de \(\mathbb E_{j:k}\) donnent \eqref{eq:amp-ns-owner-stein-truncado}. La répétition du même dédoublement donne \eqref{eq:amp-ns-owner-telescopio-truncado} ; on ne supprime pas \((\Gamma_{M,0:J}^{\rm term})^*U_{M,J;J,T}^{\rm term}
\Gamma_{M,0:J}^{\rm term}\).
\end{proof}

\begin{theorem}[Injection racine et budget uniforme]
\label{thm:amp-ns-inyeccion-raiz-uniforme}
La racine du réalisateur est transportée à chaque coupure et profondeur par une inclusion isométrique :
\begin{equation}\label{eq:amp-ns-inyeccion-raiz-comun}
 \begin{aligned}
 x_{M,0}^{\rm data}
 &:=\mathcal V_M^{\rm full}(P_Mu_0)\oplus g_M^F,\\
 \mathbf G_{M,0;J,T}^{\rm term}x_{M,0}^{\rm data}
 &=\mathcal J_{M,J,T}^{\rm term}(u_0,f),\\
 \mathcal J_{M,J,T}^{\rm term}
 &:=I_{0\to M,J}\mathcal J_T^{\rm term},
 & I_{0\to M,J}^*I_{0\to M,J}&=I.
 \end{aligned}
\end{equation}
Pour les données de \eqref{eq:realizaciones-ns-dominio}, il existe \(C_{{\rm HMT},T}\), indépendant de \(M\) et de \(J\), tel que
\begin{equation}\label{eq:amp-ns-inyeccion-raiz-cota}
 \sup_{M,J}
 \|\mathcal J_{M,J,T}^{\rm term}(u_0,f)\|^2
 \leq C_{{\rm HMT},T}
 +(1+2\widetilde M_{s,\beta,\nu}^{\,2})
  e^{R_{\rm F}^{\,2}\|u_0\|_{H^s}^2}
 +C_{\nu,s}\int_0^T\|f(t)\|_{H^{s-1}}^2\,dt .
\end{equation}
De manière équivalente, la borne porte sur le Gram terminal évalué sur les données :
\begin{equation}\label{eq:amp-ns-owner-gram-cota-uniforme}
 \sup_{M,J}
 \left\langle x_{M,0}^{\rm data},
 U_{M,0;J,T}^{\rm term}x_{M,0}^{\rm data}\right\rangle
 =
 \sup_{M,J}
 \|\mathcal J_{M,J,T}^{\rm term}(u_0,f)\|^2
 <\infty .
\end{equation}
L'application \(\mathcal J_T^{\rm term}\) dépend uniquement de l'état initial, de ses deux feuillets, du terminal hérité et du port de force ; elle ne reçoit pas pour argument \(u_M(t)\) pour \(0<t\leq T\).
\end{theorem}

\begin{proof}
La composante linéaire et les mémoires \(D/H\) sont transportées isométriquement. Dans la composante PC--Fock, la colonne cohérente est
\[
 \bigoplus_{n\geq0}
 \frac{R_{\rm F}^{\,n}}{\sqrt{n!}}\,u_0^{\odot n},
 \qquad
 \sum_{n\geq0}\frac{R_{\rm F}^{\,2n}}{n!}
       \|u_0\|_{H^s}^{2n}
 =e^{R_{\rm F}^{\,2}\|u_0\|_{H^s}^2}.
\]
Pour justifier l'uniformité du lecteur quadratique, on polarise
\begin{equation}\label{eq:amp-ns-lector-cuadratico-polarizado}
 \mathcal B_M(u,v):=
 \frac{1}{4\sqrt{\beta\nu}}\,
 A_M^{-1/2}\Lambda^sP_M\mathbb P
 \bigl((P_Mu\!\cdot\!\nabla)P_Mv
      +(P_Mv\!\cdot\!\nabla)P_Mu\bigr).
\end{equation}
Dans les bases transversales \(e^s_{p,a}=\langle p\rangle^{-s}a\,e^{ip\cdot x}\), si \(k=p+q\ne0\), ses lignes de Fourier satisfont
\begin{equation}\label{eq:amp-ns-filas-fourier-uniformes}
 \|c^M_{k;p,q}\|
 \leq\frac{C}{\sqrt{\beta\nu}}\,
 \frac{\langle k\rangle^s(|p|+|q|)}
 {|k|\langle p\rangle^s\langle q\rangle^s},
 \qquad
 \sup_{M,k\ne0}\sum_{p+q=k}\|c^M_{k;p,q}\|^2<\infty .
\end{equation}
La seconde borne s'obtient en séparant \(|p|\leq2|k|\) de \(|p|>2|k|\) ; dans la queue, \(\sum_p\langle p\rangle^{2-2s}<\infty\) domine puisque \(s>5/2\). Cauchy--Schwarz sur chaque ligne et l'orthogonalité des blocs de \(k\) distincts prolongent \(\mathcal B_M\) à
\[
 \mathfrak C_M:H^s_\sigma\odot_2H^s_\sigma\longrightarrow L^2_\sigma,
 \qquad
 \sup_M\|\mathfrak C_M\|
 \leq\widetilde M_{s,\beta,\nu}.
\]
Ainsi, la constante utilisée à la racine contrôle la somme de toutes les paires de Fourier, non seulement chaque coefficient séparément. Le port affine vérifie
\[
 Q_f(T)\leq C_{\nu,s}
 \int_0^T\|f(t)\|_{H^{s-1}}^2\,dt .
\]
Par somme orthogonale de ces composantes,
\begin{equation}\label{eq:amp-ns-raiz-cota-por-componentes}
 \|\Xi_{M,0}\|^2
 \leq C_{{\rm HMT},T}
 +(1+2\widetilde M_{s,\beta,\nu}^{\,2})
 e^{R_{\rm F}^{\,2}\|u_0\|_{H^s}^2}.
\end{equation}
Les inclusions de \eqref{eq:amp-ns-inyeccion-raiz-comun} ne modifient pas cette norme. En ajoutant la borne du port, on obtient exactement \eqref{eq:amp-ns-inyeccion-raiz-cota}. Puisque tous les termes sont évalués sur \((u_0,f)\), l'estimation précède la trajectoire et ne présuppose pas la borne globale déduite ensuite de \eqref{eq:amp-ns-owner-telescopio-truncado}.
\end{proof}

Le Gram terminal \eqref{eq:amp-ns-owner-gram-truncado}, son identité de Stein, la télescopie avec terminal et l'injection racine uniforme forment la construction de référence. C'est cette branche qui gouverne la clôture et coïncide avec la réalisation intégrale restaurée dans \cref{sec:ns-owner-closure,thm:nsclose-proprietary-full}.

\section*{Distinction du contrôle KYP fini non déterminant}

La factorisation KYP finie correcte du \cref{thm:ns-import-affine-stein} est conservée comme contrôle alternatif d'une carte réduite et ne constitue pas une prémisse du développement de référence \(U^{\rm own}\). On n'utilise pas la mutation signée qui prétendait incorporer \(C_{m,M}^*C_{m,M}\) à un terme \(P\succeq0\) puis le récupérer avec un signe négatif : lorsqu'il entre dans \(P\), ce terme apparaît nécessairement avec un signe positif. Cette factorisation n'intervient donc ni dans le bilan, ni dans le budget racine, ni dans le retour physique, ni dans la conclusion Navier--Stokes.

\section*{Dérivée temporelle, pression et recollement des horizons}

L'étape suivante découle de la borne uniforme de trajectoire obtenue par le Gram terminal de référence, sa télescopie et le retour physique exact. On dispose de la borne conjointe uniforme dans \(L^\infty(0,T;H^s)\cap L^2(0,T;H^{s+1})\). L'équation projetée s'écrit
\begin{equation}\label{eq:amp-ns-proyectada}
 \partial_tu_m
 =\nu\Delta u_m-P_mB(u_m,u_m)+P_mf.
\end{equation}
Pour \(s>5/2\), le produit de Sobolev et la borne conjointe dans \(L^\infty(0,T;H^s)\cap L^2(0,T;H^{s+1})\) donnent
\begin{equation}\label{eq:amp-ns-dt}
 \sup_m\|\partial_tu_m\|_{L^2(0,T;H^{s-1})}
 \leq C_T(u_0,f,\Xi_0).
\end{equation}
L'estimation \eqref{eq:amp-ns-dt} et la compacité spatiale produisent une convergence forte dans \(L^2(0,T;H^s)\), suffisante pour passer au terme bilinéaire sans perdre la retenue de Galerkin.

Une fois la vitesse retrouvée, la pression normalisée de moyenne nulle est fixée par
\begin{equation}\label{eq:amp-ns-presion}
 -\Delta p
 =\partial_i\partial_j(u_i u_j)-\nabla\!\cdot f,
 \qquad \int_{\mathbb T^3}p(t,x)\,dx=0.
\end{equation}
La pression est ainsi une sortie de la même solution solénoïdale ; ce n'est pas une variable libre utilisée pour sélectionner la branche régulière.

\begin{proposition}[Recollement global par unicité]
\label{prop:amp-ns-pegado}
Soient \(u^{(T)}\) les solutions obtenues sur \([0,T]\) par la limite de Galerkin. Si \(T_1<T_2\), alors \(u^{(T_2)}|_{[0,T_1]}=u^{(T_1)}\). Le système définit une unique solution globale lisse pour \(t>0\).
\end{proposition}

\begin{proof}
Les deux restrictions possèdent les mêmes données et appartiennent au domaine de
l'estimation uniforme. L'équation de leur différence, le contrôle bilinéaire dans
\(H^{s-1}\) et Grönwall impliquent l'égalité sur \([0,T_1]\). Les intervalles
finis forment un système dirigé ; l'égalité sur les recouvrements définit
la solution globale. La régularisation parabolique du théorème s'applique sur
chaque intervalle \([\varepsilon,T]\), \(\varepsilon>0\).
\end{proof}

Le contrôle de réfutation est précis : une perte d'uniformité dans
\eqref{eq:amp-ns-dt}, deux limites différentes sur un même horizon ou une
pression ne satisfaisant pas \eqref{eq:amp-ns-presion} rompraient la chaîne avant
la conclusion globale.

\paragraph{Conclusion propriétaire.}
Le Gram terminal propriétaire, Stein, la télescopie avec terminal, son registre,
le retour physique exact et l'injection racine produisent une borne \(H^s\) uniforme pour
\(s>5/2\). On démontre donc l'existence, l'unicité et la régularité
globales de Navier--Stokes tridimensionnel périodique dans le domaine HMT
enregistré. Les cartes KYP finies restent en tant que
\texttt{CONTRÔLE\_NON\_DIRECTEUR}.


\paragraph{Conclusion affirmative.}
Dans le domaine périodique et de Sobolev déclaré, le stockage terminal
\(U^{\rm own}\), son identité de Stein, le télescopage qui conserve le
terminal, le registre qui domine \(E_s(u_M)\), la racine uniforme et le retour
physique exact produisent la borne uniforme. Aubin--Lions, l'unicité et la
régularité parabolique donnent la solution globale lisse. Les cartes KYP,
NS--OBS et LRDN--LCN sont des contrôles alternatifs et non des prémisses du théorème.

\chapter[Processus de jauge euclidien de Yang--Mills en dimension 4]{Construction
projective d’un processus de jauge euclidien de Yang--Mills, reconstruction
d’Osterwalder--Schrader et écart collectif en dimension \(4\)}
\chaptermark{Processus de jauge euclidien Yang--Mills}

\noindent La spécialisation du
\cref{thm:nucleo-continuidad-conexion-nonadica-conjunta} est
\[
 (\widehat X_\infty,\Gamma_9)
 \xrightarrow{\ \mathcal R_{\rm YM}\ }
 (\mu_m,T_m,H_m,\Omega_m)
 \longrightarrow(\mu_\infty,\mathrm{OS},E(4))
 \longrightarrow\operatorname{gap}(H_{\rm phys})>0.
\]
Le lecteur clover qui publie \(F^2\) reconnaît la courbure sur ce même
processus projectif ; il ne construit pas rétrospectivement \(\mu_m\). La théorie
de jauge, la reconstruction OS, le vide simple et l’écart sont déjà fixés
par cette chaîne. L’identité
\eqref{eq:amp-ym-identidad-accion-requerida} compare ensuite un
paramétrage facultatif par \(g_R\) ; c’est un contrôle non régissant, et non une
condition pour identifier la théorie de Yang--Mills construite.

\section*{Objet de jauge et domaine de raffinement}

Soit \(G\) un groupe de Lie compact, connexe et simple. Pour chaque
\(m\geq0\), on considère le réseau euclidien
\begin{equation}\label{eq:realizaciones-ym-red}
 \Lambda_m=a_m\mathbb Z^4,
 \qquad a_{m+1}=\frac{a_m}{9}.
\end{equation}
L'algèbre cylindrique de jauge est engendrée par des variables de liaison dans \(G\), des opérateurs d'entrelacement de représentation et les coordonnées enrichies de couleur, d'orientation, d'holonomie, de retenue de fusion, de chemin et de mémoire. Soit \(\mu_m\) la mesure invariante commune et soit
\[
 \mathcal H_m=L^2(\mathcal A_m,\mu_m)^{\rm gauge}.
\]
Le vecteur constant \(\Omega_m\) définit
\begin{equation}\label{eq:realizaciones-ym-proyectores}
 P_{\Omega_m}=|\Omega_m\rangle\langle\Omega_m|,
 \qquad Q_m^{\rm phys}=I-P_{\Omega_m}.
\end{equation}
Le second projecteur est le complément physique du vide ; il ne s’identifie pas
à un résiduel employé uniquement pour organiser l’échelle de
renormalisation.

Le circuit de mise à jour est local :
\begin{equation}\label{eq:realizaciones-ym-circuito}
 K_m^{\rm loc}
 =\prod_{r=1}^{9\cdot81}
   \prod_{B\in\mathcal B_{m,r}}K_{m,B}.
\end{equation}
La partition \(\{\mathcal B_{m,r}\}_r\) énumère chaque coordonnée une seule fois ; les facteurs
du circuit seront identifiés ci-dessous au semi-groupe de leurs espérances
conditionnelles, et non à neuf dynamiques indépendantes.
Les blocs de même couleur possèdent des supports disjoints et le diamètre physique
de chaque support est uniformément borné lorsque \(m\) varie. Cette localité
fait partie de l’objet, de sorte que le raffinement ne transforme pas une
mise à jour élémentaire en une opération instantanément globale.

Soit \(T_m\) le transfert induit par
\eqref{eq:realizaciones-ym-circuito}. Les inclusions isométriques
\(J_m:\mathcal H_m\to\mathcal H_{m+1}\) conservent ancêtre,
représentation, orientation, retenue et mémoire. Pour les quatre réflexions
euclidiennes \(\Theta_{\nu,m}\), \(\nu=1,\ldots,4\), on a
\begin{equation}\label{eq:realizaciones-ym-entrelazamiento}
 (T_{m+1})^9J_m=J_mT_m,
 \qquad
 \Theta_{\nu,m+1}J_m=J_m\Theta_{\nu,m}
 \quad(\nu=1,\ldots,4).
\end{equation}
Les quatre directions appartiennent donc au même raffinement, à la même mesure et à la même dynamique.

\section*{Semi-groupe commun, couronne relative et quatre formes de réflexion}

Dans la carte triangulaire de l’état complet, soient \(E_{m,c}\) les espérances
conditionnelles orthogonales et commutantes associées à des coordonnées
indépendantes. Le générateur positif de recouvrement est d’abord construit à
un niveau germe \(m_0\) :
\begin{equation}\label{eq:realizaciones-ym-generador}
 \widehat H_{m_0}^{\rm ov}
 =3\kappa_{\rm av}
  \sum_{c\in\mathcal I_{m_0}^{\rm seed}}(I-E_{m_0,c}),
\end{equation}
et se prolonge au moyen de
\begin{equation}\label{eq:realizaciones-ym-generador-refinado}
 \begin{aligned}
 \widehat H_{m+1}^{\rm ov}
 &=\widehat\Gamma_m^*
   \bigl(\widehat H_m^{\rm ov}\otimes I+
         I\otimes\widehat H_{m+1}^{\rm det}\bigr)
   \widehat\Gamma_m,\\
 \widehat H_{m+1}^{\rm det}
 &=3\kappa_{\rm av}
   \sum_{\iota\in\mathcal I_{m+1/m}^{\rm det}}
   (I-E_{m+1,\iota}^{\rm det}),
 \end{aligned}
\end{equation}
où
\begin{equation}\label{eq:realizaciones-ym-kappa}
 \kappa_{\rm av}
 =\frac{\mu_{\rm vol}}{\ell_0}
  \left(\frac{\pi_{\rm HMT}}{\varphi_{\rm HMT}^{\,2}}-1\right)>0.
\end{equation}
La fréquence volumétrique \(\mu_{\rm vol}\) est sans dimension ; la dimension provient de \(\ell_0\).

Le semi-groupe, le transfert d'un pas physique et sa compatibilité d'échelle sont définis avec des types visibles :
\begin{equation}\label{eq:realizaciones-ym-semigrupo}
 \widehat K_{m,t}^{\rm ov}=e^{-t\widehat H_m^{\rm ov}},
 \qquad
 \widehat T_m=\widehat K_{m,a_m}^{\rm ov},
 \qquad a_{m+1}=a_m/9,
\end{equation}
Comme la partition énumère chaque coordonnée une seule fois et que les espérances
conditionnelles de ses blocs commutent, on obtient pour chaque bloc et pour le circuit
complet
\begin{equation}\label{eq:realizaciones-ym-circuito-semigrupo}
 K_{m,B}=e^{-3\kappa_{\rm av}a_m(I-E_{m,B})},
 \qquad
 K_m^{\rm loc}=e^{-a_m\widehat H_m^{\rm ov}}=\widehat T_m.
\end{equation}
\begin{equation}\label{eq:realizaciones-ym-semigrupo-entrelazado}
 \widehat H_{m+1}^{\rm ov}\widehat J_m
 =\widehat J_m\widehat H_m^{\rm ov},
 \qquad
 (\widehat T_{m+1})^9\widehat J_m
 =\widehat J_m\widehat T_m.
\end{equation}
Ainsi, la neuvième puissance est comparée au moyen de \(\widehat J_m\) ; on n’égale pas
des opérateurs qui agissent dans des Hilbert distincts.

La mesure commune de huit faces provient de ce même noyau réversible :
\begin{equation}\label{eq:realizaciones-ym-medida-ocho}
 \widehat\Omega_m^{(8)}(dy_1\cdots dy_8)
 =\int\widehat\mu_m^{\rm ov}(dz)
   \prod_{r=1}^{8}
   \widehat K_{m,a_m/2}^{\rm ov}(z,dy_r).
\end{equation}
Pour chaque réflexion \(\nu=1,\ldots,4\) et tout observable cylindrique \(F\)
dans le demi-espace positif,
\begin{equation}\label{eq:realizaciones-ym-os}
 \int\overline{F(\Theta_{\nu,m}y)}F(y)\,
 d\widehat\Omega_m^{(8)}
 =\|\widehat{\mathcal B}_{m,\nu}F\|^2\geq0,
\end{equation}
et la marginale de faces opposées donne l'identité de liaison
\begin{equation}\label{eq:realizaciones-ym-os-transferencia}
 \widehat T_{m,\nu}^{\rm OS}
 =\widehat K_{m,a_m}^{\rm ov}
 =e^{-a_m\widehat H_m^{\rm ov}}
 =\widehat T_m,
 \qquad \nu=1,\ldots,4.
\end{equation}
Les quatre réflexions publient donc la même paire
transfert--mesure.

\section*{Persistance cellulaire du terminal et écart spectral exact}

Soient
\[
 K_9=C_*(Q_{9,3};\mathbb Z),\qquad
 K_{10}=C_*(Q_{10,3};\mathbb Z).
\]
La couronne relative a pour dimensions
\[
 C_*(Q_{10,3},Q_{9,3})=(271,756,702,217),
 \qquad 271+702=756+217=973.
\]
L’inclusion \(j:K_9\to K_{10}\), l’effondrement cellulaire
\(r:K_{10}\to K_9\) et l’homotopie entière \(H_{\rm cell}\) satisfont
\begin{equation}\label{eq:realizaciones-ym-sdr-corona}
 rj=I,\qquad
 \partial H_{\rm cell}+H_{\rm cell}\partial=I-jr,\qquad
 rH_{\rm cell}=0,\quad H_{\rm cell}j=0,\quad H_{\rm cell}^2=0.
\end{equation}
Pour toute application de chaînes
\(F:K_{10}\otimes V_{q=6}\to\mathcal T\), avec
\(P=(jr)\otimes I_{q=6}\),
\begin{equation}\label{eq:realizaciones-ym-homotopia-terminal}
 F-FP
 =d_{\mathcal T}\bigl(F(H_{\rm cell}\otimes I)\bigr)
  +F(H_{\rm cell}\otimes I)(\partial\otimes I).
\end{equation}
La couronne complète est ainsi homotope à une constante et le terminal \(FP\) demeure littéralement. Les \(271\) sommets ne remplacent pas les arêtes, les faces et les cubes de la couronne ; le facteur cellulaire tridimensionnel n'est pas non plus identifié au facteur de jauge physique \(q=6\).

Pour chaque coordonnée persistante \(c\), dans la carte locale
\(q=(q_1,\ldots,q_6)\in\mathbb F_3^6\), on fixe le témoin centré
\begin{equation}\label{eq:realizaciones-ym-testigo-definicion}
 \bigcap_c\operatorname{Ran}E_{m,c}=\mathbb C\Omega_m,
 \qquad
 W_c(q)=\mathbf 1_{\{q_1+\cdots+q_6=0\}}-\frac13,
\end{equation}
où la somme de l’indicatrice est prise dans \(\mathbb F_3\). La décomposition de
Hoeffding des espérances commutantes prouve
\begin{equation}\label{eq:realizaciones-ym-gap-finito}
 \ker\widehat H_m^{\rm ov}=\mathbb C\Omega_m,
 \qquad
 \widehat H_m^{\rm ov}\succeq
 3\kappa_{\rm av}(I-P_{\Omega_m}).
\end{equation}
L’égalité ne se déduit pas du seul fait que le noyau est trivial : chaque secteur non
constant de la décomposition reçoit au moins un projecteur
\(I-E_{m,c}\) et, par \eqref{eq:realizaciones-ym-generador}, son coût minimal
est \(3\kappa_{\rm av}\). Le témoin matériel \(W_c\) vérifie
\begin{equation}\label{eq:realizaciones-ym-testigo-gap}
 \widehat H_m^{\rm ov}W_c=3\kappa_{\rm av}W_c,
 \qquad \|W_c\|^2=\frac29,
\end{equation}
de sorte que le complément est non vide et que l’écart est exactement
\(3\kappa_{\rm av}\).

L'échelle est publiée par
\begin{equation}\label{eq:realizaciones-ym-qm}
 q_m=e^{-3\kappa_{\rm av}a_m},
 \qquad q_{m+1}^{\,9}=q_m,
 \qquad -a_m^{-1}\log q_m=3\kappa_{\rm av}.
\end{equation}
Bien que \(q_m\to1\) lors du raffinement, le quotient spectral par unité physique reste constant. La contraction cellulaire \eqref{eq:realizaciones-ym-homotopia-terminal} conserve le terminal ; l'écart spectral provient du générateur positif de recouvrement et de sa décomposition de Hoeffding, et non de l'acyclicité à elle seule.

Les inclusions de \eqref{eq:realizaciones-ym-semigrupo-entrelazado}
forment la limite inductive
\begin{equation}\label{eq:realizaciones-ym-colimite}
 \mathcal H_\infty=\varinjlim(\mathcal H_m,\widehat J_m).
\end{equation}
Sur l’union inductive algébrique, on définit la forme
\begin{equation}\label{eq:realizaciones-ym-forma-limite}
 \mathfrak h_\infty[\widehat J_{m,\infty}\psi]
 :=\langle\psi,\widehat H_m^{\rm ov}\psi\rangle_{\mathcal H_m}.
\end{equation}
L’entrelacement des générateurs rend cette définition
indépendante du représentant. La forme est positive et fermable ; sa
fermeture détermine le générateur auto-adjoint \(\widehat H_\infty^{\rm ov}\).
La compatibilité des formes transporte
\begin{equation}\label{eq:realizaciones-ym-gap-forma-limite}
 \overline{\mathfrak h}_\infty[\psi]
 \geq3\kappa_{\rm av}
 \|(I-P_{\Omega_\infty})\psi\|^2,
\end{equation}
et le vecteur persistant
\(\widehat J_{m,\infty}W_c\) atteint l’égalité. Par conséquent,
l’écart exact du générateur limite est \(3\kappa_{\rm av}\).

\section*{Hypothèses exactes et enchaînement du théorème}

La publication de jauge utilise :
\begin{enumerate}
 \item le groupe compact, connexe et simple et le réseau quadridimensionnel
       \eqref{eq:realizaciones-ym-red} ;
 \item le circuit local \eqref{eq:realizaciones-ym-circuito}, de diamètre physique uniforme, et une mesure invariante commune ;
 \item les inclusions isométriques et les quatre réflexions entrelacées par
       \eqref{eq:realizaciones-ym-entrelazamiento} et
       \eqref{eq:realizaciones-ym-semigrupo-entrelazado} ;
 \item les quatre factorisations
       \eqref{eq:realizaciones-ym-os} et le lien exact
       \eqref{eq:realizaciones-ym-os-transferencia} sur le même couple
       transfert--mesure ;
 \item le générateur d'espérances commutantes \eqref{eq:realizaciones-ym-generador}--\eqref{eq:realizaciones-ym-generador-refinado} et sa décomposition de Hoeffding ;
 \item le SDR entier de la couronne et la conservation du terminal
       \eqref{eq:realizaciones-ym-sdr-corona}--
       \eqref{eq:realizaciones-ym-homotopia-terminal} ;
 \item le terminal \(P_{\Omega_m}\), le complément physique
       \(Q_m^{\rm phys}\) et la loi d’échelle
       \eqref{eq:realizaciones-ym-qm} ;
 \item la limite inductive et la fermeture compatible des formes \eqref{eq:realizaciones-ym-colimite}--\eqref{eq:realizaciones-ym-gap-forma-limite} ;
 \item les cartes internes d’action et de vitesse
       \(\hbar_{\rm int}\), \(c_{\rm int}\), utilisées uniquement pour publier la
       dimension de masse.
\end{enumerate}

\begin{theorem}[Théorie de jauge locale à vide simple et écart collectif]
\label{thm:realizaciones-ym}
Sous les prémisses précédentes, la limite de la famille de jauge est locale et positive par réflexion dans les quatre directions euclidiennes. Son générateur possède un vide simple et satisfait, dans le secteur collectif invariant de jauge,
\begin{equation}\label{eq:realizaciones-ym-gap-limite}
 \Delta_{\rm YM}^{\rm HMT}=3\kappa_{\rm av}>0,
 \qquad
 m_{\rm gap}^{\rm HMT}
 =\frac{\hbar_{\rm int}}{c_{\rm int}}
  \Delta_{\rm YM}^{\rm HMT}>0.
\end{equation}
L’écart appartient au secteur collectif ; il n’assigne pas de masse élémentaire au
gluon.
\end{theorem}

\begin{proof}
La localité uniforme du circuit conserve le réseau d’algèbres locales. Les
égalités \eqref{eq:realizaciones-ym-semigrupo-entrelazado} identifient la
neuvième puissance après application de l’inclusion correcte. La construction
\eqref{eq:realizaciones-ym-medida-ocho} utilise ce même semi-groupe dans les huit
faces et \eqref{eq:realizaciones-ym-os-transferencia} identifie ses quatre
marginales OS à \(\widehat T_m\). Par conséquent, les quatre formes positives
\eqref{eq:realizaciones-ym-os} sont compatibles à la limite et conservent une
seule mesure et une seule dynamique.

Les espérances de \eqref{eq:realizaciones-ym-generador} sont orthogonales et
commutantes. Leur décomposition de Hoeffding sépare le mode constant de chaque
secteur non constant : le premier est \(\mathbb C\Omega_m\), tandis que tout
secteur du complément reçoit au moins un facteur \(I-E_{m,c}\) de poids
\(3\kappa_{\rm av}\). Cela prouve
\eqref{eq:realizaciones-ym-gap-finito}. Le témoin
\eqref{eq:realizaciones-ym-testigo-gap} atteint la borne, de sorte que ce n’est pas
seulement un plancher : c’est la première énergie non nulle.

L'homotopie \eqref{eq:realizaciones-ym-homotopia-terminal} contracte le complément cellulaire et laisse \(FP\) intact ; la couronne n'introduit donc pas de second vide et n'efface pas la retenue. La compatibilité isométrique transporte le vide et son complément au niveau suivant sans les mélanger avec le résidu d'échelle. La loi \eqref{eq:realizaciones-ym-qm} montre que le quotient spectral par unité physique est exactement \(3\kappa_{\rm av}\) à chaque niveau. Les formes \eqref{eq:realizaciones-ym-forma-limite}--\eqref{eq:realizaciones-ym-gap-forma-limite} transportent la borne et le témoin vers le générateur autoadjoint de la limite, de sorte que l'écart spectral reste exact. La reconstruction par positivité d'Osterwalder--Schrader identifie l'hamiltonien au générateur de ce même transfert, conserve le vide simple et produit \eqref{eq:realizaciones-ym-gap-limite}. La carte dimensionnelle finale transforme l'écart spectral en la masse collective indiquée.
\end{proof}

\section*{Compression de jauge et témoin physique de l'écart spectral}

Soit
\[
 \widehat{\mathscr H}_m
 :=L^2(\widehat X_m,\widehat\mu_m^{\rm ov})
\]
l'espace de la réalisation matérielle complète. Outre l'action de jauge usuelle sur les liaisons et les repères, chaque collier \(c\) porte l'action finie d'incidence
\begin{equation}\label{eq:amp-ym-accion-gauge-incidencia}
 q_c\longmapsto q_c+Bx_c,
 \qquad x_c\in\mathbb F_3^4,
\end{equation}
où les six lignes de \(B\) sont indexées par
\((01,02,03,12,13,23)\). Le produit des deux actions préserve
\(\widehat\mu_m^{\rm ov}\). Avec \(\mathcal C_m\) l’ensemble fini des
colliers présents au niveau, on écrit
\[
 \mathcal G_m^{\rm full}
 :=G^{V(\Lambda_m)}
   \times\prod_{c\in\mathcal C_m}\mathbb F_3^4.
\]
Sa moyenne de Haar définit la projection
orthogonale
\begin{equation}\label{eq:amp-ym-proyeccion-fisica-haar}
 \mathbb P_m^{\rm phys}F
 :=\int_{\mathcal G_m^{\rm full}}F(g^{-1}\!\cdot\,)\,dg,
 \qquad
 \widehat{\mathscr H}_m^{\rm phys}
 :=\operatorname{Ran}\mathbb P_m^{\rm phys}.
\end{equation}
C'est une compression sur la sous-algèbre de jauge de l'état complet ; la marginale qui oublie les coordonnées d'incidence est une autre application et n'est pas confondue avec \(\mathbb P_m^{\rm phys}\).

\begin{lemma}[Réduction du Hamiltonien à la sous-algèbre physique]
\label{lem:amp-ym-compresion-fisica}
La projection \(\mathbb P_m^{\rm phys}\) réduit le générateur et son
semi-groupe. Les inclusions de raffinement réduisent simultanément les
sous-algèbres physiques :
\begin{equation}\label{eq:amp-ym-compresion-entrelazada}
 \begin{gathered}
 [\mathbb P_m^{\rm phys},\widehat H_m^{\rm ov}]=0,
 \qquad
 [\mathbb P_m^{\rm phys},e^{-t\widehat H_m^{\rm ov}}]=0,\\
 \mathbb P_{m+1}^{\rm phys}\widehat J_m
 =\widehat J_m\mathbb P_m^{\rm phys},\qquad
 H_m^{\rm phys}
 :=\widehat H_m^{\rm ov}\big|_{\widehat{\mathscr H}_m^{\rm phys}},\\
 H_{m+1}^{\rm phys}\widehat J_m
 =\widehat J_mH_m^{\rm phys}.
 \end{gathered}
\end{equation}
\end{lemma}

\begin{proof}
Le générateur physique sur les liens est équivariant de jauge par la
bi-invariance de Haar. Dans la carte produit, \(E_{q_c}\) est l’espérance
uniforme sur \(\mathbb F_3^6\) ; par invariance de la mesure uniforme sous
les translations de \eqref{eq:amp-ym-accion-gauge-incidencia}, il commute avec
cette action. Les espérances des sous-queues de marquage agissent sur des
facteurs disjoints. Par conséquent, chaque terme de
\[
 \widehat H_m^{\rm ov}
 =H_m^{\rm ov}\otimes I
 +3\kappa_{\rm av}\sum_c
 \bigl[(I-E_{q_c})+(I-E_{u_c^{\rm P}})\bigr]
\]
commute avec la moyenne \eqref{eq:amp-ym-proyeccion-fisica-haar}. Le calcul fonctionnel donne la commutation avec le semi-groupe. Enfin, \(\widehat J_m\) conserve les facteurs ancestraux et insère la constante dans les nouveaux facteurs ; prendre la moyenne avant ou après cette inclusion produit la même fonction. Cela prouve les quatre identités de \eqref{eq:amp-ym-compresion-entrelazada}.
\end{proof}

\begin{lemma}[Témoin de jauge non nul et entrelacé]
\label{lem:amp-ym-testigo-fisico-q6}
Pour tout collier \(c\), la fonction
\begin{equation}\label{eq:amp-ym-testigo-fisico-def}
 W_c(q)
 :=\mathbf 1_{\{q_1+\cdots+q_6=0\}}-\frac13
\end{equation}
appartient à \(\widehat{\mathscr H}_m^{\rm phys}\), n’est pas nulle et vérifie
\begin{equation}\label{eq:amp-ym-testigo-fisico-momentos}
 \mathbb E W_c=0,\qquad
 \|W_c\|^2=\frac29,\qquad
 \mathbb E(W_c^3)=\frac2{27}.
\end{equation}
De plus,
\begin{equation}\label{eq:amp-ym-testigo-fisico-gap}
 H_m^{\rm phys}W_c=3\kappa_{\rm av}W_c,\qquad
 H_{m+1}^{\rm phys}\widehat J_mW_c
 =3\kappa_{\rm av}\widehat J_mW_c.
\end{equation}
Par conséquent, la valeur \(3\kappa_{\rm av}\) appartient au spectre de la
restriction physique et non seulement à celui d’une dilatation matérielle.
\end{lemma}

\begin{proof}
L’action de jauge sur les liens agit sur le premier facteur et laisse fixe
\(W_c\). Pour l’action d’incidence, chaque \(x_i\) apparaît dans trois des
six coordonnées, donc
\[
 \sum_{i<j}(Bx)_{ij}=3\sum_i x_i=0
 \quad\text{dans }\mathbb F_3.
\]
La condition de \eqref{eq:amp-ym-testigo-fisico-def} est donc
invariante ; elle l’est aussi sous les permutations des six incidences.
Ainsi \(\mathbb P_m^{\rm phys}W_c=W_c\).

Sous la mesure uniforme, la somme des six trits est uniforme dans \(\mathbb F_3\). La fonction prend la valeur \(2/3\) avec probabilité \(1/3\) et \(-1/3\) avec probabilité \(2/3\), ce qui donne \eqref{eq:amp-ym-testigo-fisico-momentos}. Le premier terme de \(\widehat H_m^{\rm ov}\) annule \(W_c\) ; pour \(c'\ne c\), \((I-E_{q_{c'}})W_c=0\), tandis que \(E_{q_c}W_c=0\). Toutes les espérances de marquage annulent la différence correspondante car \(W_c\) est constante dans ces sous-queues. Il reste exactement \(\widehat H_m^{\rm ov}W_c=3\kappa_{\rm av}W_c\). La première identité de \eqref{eq:amp-ym-testigo-fisico-gap} découle de la réduction ; la seconde, de l'entrelacement de \eqref{eq:amp-ym-compresion-entrelazada}.
\end{proof}

\section*{Semi-groupe propriétaire, vide et borne spectrale exacte}

Sur la sous-algèbre physique, on écrit
\[
 D_m^{\rm YM}:=H_m^{\rm phys},\qquad
 K_{m,t}:=\exp(-tD_m^{\rm YM}),\qquad
 P_{\Omega_m}:=|\Omega_m\rangle\langle\Omega_m|.
\]
La décomposition de Hoeffding du même générateur et le calcul fonctionnel
donnent, sans changer de mesure ni de processus,
\begin{equation}\label{eq:amp-ym-owner-semigrupo-vacio}
 K_{m,t}P_{\Omega_m}=P_{\Omega_m},
 \qquad
 \bigl\|K_{m,t}(I-P_{\Omega_m})\bigr\|
 \leq e^{-3\kappa_{\rm av}t}.
\end{equation}
Le témoin de \eqref{eq:amp-ym-testigo-fisico-def} est non nul et atteint la
borne :
\begin{equation}\label{eq:amp-ym-owner-testigo-wc}
 D_m^{\rm YM}W_c=3\kappa_{\rm av}W_c,
 \qquad \|W_c\|^2=\frac29.
\end{equation}
Par conséquent, le vide est simple et l’écart n’est pas seulement minoré,
mais vaut exactement
\begin{equation}\label{eq:amp-ym-owner-brecha-exacta}
 \Delta_{\rm YM}^{\rm HMT}=3\kappa_{\rm av}>0.
\end{equation}
Le semi-groupe, le projecteur du vide, la contraction du complément et
\(W_c\) appartiennent à la même famille projective et constituent la chaîne
propriétaire de la clôture.

\section*{Reconstruction euclidienne, champs de Schwinger et lecteur de courbure}

La mesure projective commune du théorème détermine des inclusions isométriques
\(\widehat J_m\) entre les espaces de niveau. Pour chacune des quatre
réflexions coordonnées \(\theta_\mu\), \(\mu=1,\ldots,4\), il existe une
factorisation
\begin{equation}\label{eq:amp-ym-cuatro-os}
 \int\overline{F(\theta_\mu y)}F(y)\,
 d\widehat\Omega_m(y)
 =\|\mathcal B_{m,\mu}F\|_2^2\geq0.
\end{equation}
Les quatre formes proviennent de la même mesure et du même semi-groupe. Les
connecteurs OS
\[
 \mathcal J_{m,\mu}^{\rm OS}
 =\mathcal V_{m+1,\mu}^{-1}\widehat J_m\mathcal V_{m,\mu}
\]
entrelacent leurs hamiltoniens. Dans la colimite,
\begin{equation}\label{eq:amp-ym-hamiltoniano-comun}
 \mathcal V_{\infty,\mu}H_{{\rm OS},\infty}^{(\mu)}
 \mathcal V_{\infty,\mu}^{-1}
 =\widehat H_\infty,
 \qquad \mu=1,\ldots,4.
\end{equation}

La forme limite se diagonalise par la décomposition de Hoeffding :
\begin{equation}\label{eq:amp-ym-forma-limite}
 \mathcal E_\infty(u)=3\kappa_{\rm av}
 \sum_{S\Subset\mathcal I_\infty}|S|\,\|u_S\|^2.
\end{equation}
Les sommes partielles fournissent une suite de récupération et la semi-continuité produit la limite inférieure. La convergence de Mosco des formes fermées conserve le noyau unidimensionnel et le seuil spectral. Le vecteur matériel \(W_c\) satisfait
\begin{equation}\label{eq:amp-ym-testigo-gap}
 \widehat H_mW_c=3\kappa_{\rm av}W_c,
 \qquad \|W_c\|^2=\frac29,
 \qquad \mathbb E(W_c^3)=\frac2{27},
\end{equation}
de sorte que la borne inférieure est atteinte et que l’écart est exactement
\(3\kappa_{\rm av}\).

\begin{theorem}[Complétion euclidienne et publication de courbure]
\label{thm:amp-ym-terminacion}
La famille projective du théorème \ref{thm:realizaciones-ym} détermine une
action fortement continue de \(E(4)\), un réseau local de jauge non scalaire,
des fonctionnelles de Schwinger compatibles et une publication de champ dans
\(H^{-s}_{\rm loc}(\mathbb R^4)\) pour \(s>2\). Le produit surfacique
ordonné des liens définit des lecteurs clover et \(F^2\) ; dans le domaine
lisse, le lecteur d’action converge vers
\begin{equation}\label{eq:amp-ym-f2}
 \mathcal S_{\rm YM}(A)
 =\frac1{4g_R^2}\int_{\mathbb R^4}\|F_A(x)\|^2\,d^4x.
\end{equation}
Ces publications conservent le vide simple et l’écart exact du même
hamiltonien commun.
\end{theorem}

\begin{proof}
Les matrices de Gram cofinales des lecteurs lissés et la commutativité des carrés de rotation et de raffinement préservent les idéaux nuls ; leur colimite construit la représentation forte de \(E(4)\) et le réseau local. Les estimations de martingale des accroissements de liaison donnent la tension dans \(H^{-s}_{\rm loc}\), \(s>2\), et les caractéristiques mixtes identifient les fonctionnelles de Schwinger au même processus cylindrique. Le produit orienté autour de chaque plaquette construit le clover. Son développement dans le régime lisse et la convergence forte du lecteur marginal donnent \eqref{eq:amp-ym-f2}. Aucune de ces opérations ne modifie la mesure, le semi-groupe ni la forme limite utilisés pour obtenir \eqref{eq:amp-ym-testigo-gap} ; le vide et l'écart spectral sont donc conservés.
\end{proof}

\section*{Identification exacte de la limite de jauge}

L’identification de la théorie ne repose pas sur le fait que quatre réflexions
engendrent à elles seules le groupe euclidien. Les réflexions fournissent la
positivité OS ; l’action continue provient, séparément, du groupoïde des
écrans et de la compatibilité de raffinement. Pour
\(g=(g_v)_{v\in\Lambda_m}\in G^{\Lambda_m}\), l’action sur une variable
de lien est
\begin{equation}\label{eq:amp-ym-accion-gauge-enlace}
 (g\cdot U)_e=g_{s(e)}U_e g_{t(e)}^{-1}.
\end{equation}
L’invariance de Haar, la désintégration triangulaire du noyau et
l’annulation des arêtes intérieures prouvent
\begin{equation}\label{eq:amp-ym-invariancia-gauge-medida}
 (g\cdot)_*\widehat\mu_m^{\rm ov}=\widehat\mu_m^{\rm ov},
 \qquad
 K_{m,t}^{\rm ov}(g\cdot U,g\cdot dV)
 =K_{m,t}^{\rm ov}(U,dV).
\end{equation}
Par différentiation sur le noyau cylindrique, tout générateur vertical
\(X\) vérifie l’identité de Ward
\begin{equation}\label{eq:amp-ym-ward}
 \int \nabla_XF\,d\widehat\mu_m^{\rm ov}=0.
\end{equation}

La courbure s’obtient à partir du même système de liens. Si une plaquette
grossière \(p\) est subdivisée en ses 81 sous-plaquettes \(p'\), alors
\begin{equation}\label{eq:amp-ym-producto-superficial}
 \operatorname{Hol}^{(m)}_{p}
 =\prod_{p'\subset p}^{\longrightarrow}
 T_{p'}\operatorname{Hol}^{(m+1)}_{p'}T_{p'}^{-1}.
\end{equation}
Les contributions de toute arête intérieure apparaissent avec des orientations
opposées et s’annulent exactement.  L’expression est covariante de jauge et
commute avec les inclusions \(\widehat J_m\).  Dans une carte lisse, on définit
\begin{equation}\label{eq:amp-ym-clover-definido}
 F_{{\rm cl},m}(x)
 =\frac1{4a_m^2}\sum_{p\ni x}
 \operatorname{Ad}_{T_p}\log\operatorname{Hol}_{p},
 \qquad
 F_{{\rm cl},m}=F_A+O(a_m^2).
\end{equation}
De là, pour \(A\in C^4_{\rm loc}\) et
\(a_m^2/g_m^2\to0\), on obtient par somme de Riemann
\begin{equation}\label{eq:amp-ym-f2-convergencia-calculada}
 \frac{a_m^4}{4g_m^2}
 \sum_{x\in\Lambda_m}\|F_{{\rm cl},m}(x)\|^2
 \longrightarrow
 \frac1{4g_R^2}\int_{\mathbb R^4}\|F_A(x)\|^2\,d^4x.
\end{equation}

\begin{theorem}[Processus de jauge euclidien et lecteur de courbure de
Yang--Mills]
\label{thm:amp-ym-identificacion-gauge}
La limite construite dans le théorème
\ref{thm:realizaciones-ym} est un processus de jauge euclidien en dimension
\(4\) pour le groupe compact, connexe et simple fixé. Elle possède
une covariance de jauge locale, un réseau local non scalaire, une positivité par
réflexion, une représentation fortement continue de \(E(4)\),
des fonctionnelles de Schwinger compatibles, une courbure covariante de jauge de limite
classique \eqref{eq:amp-ym-f2-convergencia-calculada}, un vide simple et un écart
spectral
\(3\kappa_{\rm av}>0\).
\end{theorem}

\begin{proof}
Les identités \eqref{eq:amp-ym-invariancia-gauge-medida}--
\eqref{eq:amp-ym-ward} construisent la symétrie de jauge et ses identités de
Ward à tous les niveaux ; leur naturalité les transporte à la limite.  La
localité uniforme des noyaux produit le réseau isotone et local.  Les
quatre formes \eqref{eq:amp-ym-cuatro-os}, avec l’action du groupoïde
des écrans sur le noyau de Weyl, produisent respectivement la positivité
OS et l’action forte de \(E(4)\).  La tension dans
\(H^{-s}_{\rm loc}\), \(s>2\), identifie une unique loi limite et ses
fonctionnelles de Schwinger.  Enfin,
\eqref{eq:amp-ym-producto-superficial}--
\eqref{eq:amp-ym-f2-convergencia-calculada} identifient le lecteur de champ
et sa limite classique, tandis que
\eqref{eq:amp-ym-forma-limite}--
\eqref{eq:amp-ym-testigo-gap} prouvent le vide simple et l’écart exact pour la
même loi.

La mesure est déterminée constructivement par ses marginales projectives, son noyau réversible et ses fonctionnelles cylindriques. Le lecteur clover ne repondère pas cette mesure : il publie une fonction du même processus et sa limite classique.
\end{proof}

\begin{proposition}[Contrôle non directeur de comparaison Gibbs--action]
\label{prop:amp-ym-puente-accion-requerido}
La loi projective, son vide simple et son écart spectral exact sont déjà construits par \eqref{eq:amp-ym-owner-semigrupo-vacio}--\eqref{eq:amp-ym-owner-brecha-exacta}. Comme comparaison extérieure facultative de type Gibbs--action, un paramétrage par \(g_R\) peut étudier une famille
\[
 g_R\longmapsto
 \bigl(\widehat\mu_{m,g_R}^{\rm ov},H_{m,g_R}^{\rm phys}\bigr)
\]
et une relation entre la mesure ou sa forme de Dirichlet et la fonctionnelle
d’action, par exemple
\begin{equation}\label{eq:amp-ym-identidad-accion-requerida}
 \frac{d\widehat\mu_{m,g_R}^{\rm ov}}{d\nu_m}(U)
 =Z_{m,g_R}^{-1}e^{-S_{m,g_R}^{F^2}(U)},
 \qquad
 \mathfrak h_{m,g_R}[F]
 =\langle F,H_{m,g_R}^{\rm phys}F\rangle,
\end{equation}
avec compatibilité de raffinement et passage à la limite. Ici, \(\nu_m\) peut
être la mesure produit de Haar du réseau fini. Cette vérification est étiquetée
\texttt{CONTRÔLE\_NON\_DIRECTEUR} : elle n’est une prémisse ni de la théorie de jauge HMT ni
de son écart.
\end{proposition}

\begin{proof}
Dans \eqref{eq:amp-ym-f2-convergencia-calculada}, faire varier \(g_R\) remet à l'échelle le lecteur de courbure, tandis que les formules qui construisent \(\widehat\mu_m^{\rm ov}\), \(\widehat H_m^{\rm ov}\) et l'écart spectral \(3\kappa_{\rm av}\) ne contiennent pas \(g_R\). Ce sont donc des constructions indépendantes. La compression \eqref{eq:amp-ym-compresion-entrelazada} et le témoin \eqref{eq:amp-ym-testigo-fisico-gap} prouvent la persistance physique de l'écart spectral du processus déjà construit ; le contrôle compare seulement ensuite une dépendance dynamique supplémentaire en \(g_R\).
\end{proof}

La suite démontrée dans cette section est
\[
 \text{mesure commune}\to\text{quatre OS}\to\text{Hamiltonien commun}
 \to E(4),\ \text{Schwinger},\ H^{-s}_{\rm loc}
 \to\text{clover et }F^2.
\]
Cette séquence démontre une théorie de Yang--Mills quadridimensionnelle avec un vide
simple et un écart positif exact \(3\kappa_{\rm av}\) dans le domaine HMT
enregistré. L'identité
\eqref{eq:amp-ym-identidad-accion-requerida} reste en dehors de la chaîne
directrice en tant que \texttt{CONTRÔLE\_NON\_DIRECTEUR}.


\section*{Lecture de Wilson}

L’observable de Wilson s’obtient après le générateur. Pour une carte
archimédienne de lien,
\begin{equation}\label{eq:realizaciones-ym-wilson}
 A_e=\sum_a\operatorname{ev}_{\mathbb R}(x_{e,a})T_a,
 \qquad
 U_e=e^{a_mA_e},
 \qquad
 W_\gamma=\operatorname{Tr}\!\prod_{e\in\gamma}U_e.
\end{equation}
La lecture conserve la représentation, l’opérateur d’entrelacement et
l’ordre de l’holonomie. Elle sert d’observable terminal de la théorie publiée ; elle ne définit pas
rétrospectivement la mesure, le semi-groupe ni l’écart.

\section*{Nécessité des hypothèses et obstructions exactes}

\begin{ablation}[Nécessité de la dynamique et de la mesure communes]
\label{abl:realizaciones-ym}
Le théorème \ref{thm:realizaciones-ym} n’est pas conservé si l’on remplace les
neuf sous-phases par des mises à jour indépendantes, si l’on utilise des mesures ou
des noyaux différents dans les quatre réflexions, si l’on perd la localité physique
uniforme, si \(\widehat J_m\) cesse d’entrelacer transfert et réflexion, si l’on
efface la retenue ou la mémoire, si l’on élimine \(P_{\Omega_m}\), si l’on identifie
\(Q_m^{\rm phys}\) à un résidu d’échelle, si l’on exige
\(q_m\leq q_*<1\) en unités de réseau, si l’on remplace la couronne complète
par ses \(271\) sommets, si l’on efface \(FP\), ou si l’on utilise
\eqref{eq:realizaciones-ym-wilson} pour sélectionner le générateur.
\end{ablation}

\begin{proof}
Des actualisations, noyaux ou mesures distincts détruisent la factorisation
simultanée \eqref{eq:realizaciones-ym-os}. Sans localité ou entrelacement,
la positivité d’une coupe ne définit pas un réseau local compatible. Effacer la retenue,
la mémoire ou \(FP\) élimine de l’information que
\eqref{eq:realizaciones-ym-homotopia-terminal} conserve. Utiliser seulement les
sommets change le complexe relatif et peut fabriquer des modes nuls. Confondre
les deux compléments permet que le secteur physique disparaisse par une
opération purement scalaire. La borne fixée en unités de réseau contredit
\eqref{eq:realizaciones-ym-qm}. Enfin, Wilson n’est qu’une lecture de
l’état construit et ne démontre par lui-même ni
\eqref{eq:realizaciones-ym-os} ni
\eqref{eq:realizaciones-ym-gap-finito}.
\end{proof}

Les résidus indépendants sont : l'échec de \(\widehat T_{m,\nu}^{\rm OS}=e^{-a_m\widehat H_m^{\rm ov}}\), l'échec de l'entrelacement \eqref{eq:realizaciones-ym-semigrupo-entrelazado}, une forme OS négative, un vecteur du noyau orthogonal à \(\Omega_m\), ou un secteur de Hoeffding non constant d'énergie inférieure à \(3\kappa_{\rm av}\).

\section*{Portée de la construction de jauge}

La preuve combine la connexion conjointe, la conservation du terminal, la
double lecture, le circuit local, la mesure commune, les quatre factorisations
de réflexion, le raffinement entrelacé et le générateur positif. Dans le
domaine de jauge régulier déclaré, la
localité, la théorie de Yang--Mills quadridimensionnelle, le vide simple et
l’écart collectif sont des résultats exacts. La reconstruction
d’Osterwalder--Schrader et l’observable de Wilson sont des
publications ultérieures du même état ; aucune n’agit comme
sélecteur rétrospectif de la mesure, du semi-groupe, du vide ou de
l’écart. La comparaison facultative avec une famille paramétrée par \(g_R\) est étudiée
au moyen de \eqref{eq:amp-ym-identidad-accion-requerida} ; cette vérification ne
régit ni ne remet en question la mesure, le semi-groupe, le vide ou l’écart déjà
démontrés.

\section*{Comparaison des mécanismes spectral, solénoïdal et de jauge}

Les constructions précédentes partagent l'architecture
\begin{equation}\label{eq:realizaciones-sintesis-tres}
 \begin{array}{ccl}
 \text{Riemann--Weil}
&:& \text{mémoire finie}\to\text{double publication cofinale}\\
&&{}\to B_W=\mathscr I_{\rm GW}\geq0
      \to\text{critère réversible},\\[1mm]
 \text{Navier--Stokes}
&:& U^{\rm own}\to\text{télescopage et registre}\\
&&{}\to\text{racine uniforme et retour exact}
      \to\text{limite parabolique},\\[1mm]
 \text{processus de jauge OS}
&:& \text{racine terminale}\to\text{positivité par réflexion}\\
&&{}
      \to\text{écart collectif}.
 \end{array}
\end{equation}
La connexion nonadique conjointe ne remplace pas les lemmes spécialisés : elle les
organise et conserve leurs données lors du raffinement. Les conclusions de
Riemann--Weil, Navier--Stokes et Yang--Mills proviennent de leurs chaînes
propriétaires respectives et de leurs critères finaux déjà démontrés. Les comparateurs
de résiduel, de Gram et de densité d'action demeurent des contrôles
ultérieurs non directeurs. Cette séparation empêche tant la perte de
généalogie que le transfert injustifié d'une preuve d'un domaine
à un autre.

% Residencia sucesora única de Hodge y Birch--Swinnerton-Dyer. Este input es
% parte del capítulo, no una importación duplicada de los manuscritos
% fuentes.
% 04 — FRAGMENTO SUCESOR HODGE Y BIRCH--SWINNERTON-DYER
% =======================================================
% Inserción autónoma prevista tras el marcador de integración del Libro IV.
% No importa fuentes probatorios ni modifica las tres realizaciones
% precedentes.
%
% Trazabilidad técnica no narrativa:
%   Hodge: CHI-II-HDG-01/02/03, CHI-IV-HDG-01/C01.
%   BSD:   CHI-II-BSD-01/02/03, CHI-IV-BSD-01/02/C01.
%   Fuentes: celda APP--Mayer--Vietoris, totalizador coend y controles
%   elípticos; producto fibrado genealógico de Manin, publicación acoplada
%   Kato--PT, dualidad ortogonal reducida, relleno finita--singular y regulador
%   Bockstein derivado. Los hashes y estados de ejecución permanecen en el
%   expediente técnico de la edición sucesora.

\ifdefined\HMTSucesoraStructureTrace
  \typeout{HMT-SUCCESSEURE 04 : réalisations Hodge et Birch--Swinnerton-Dyer}
\fi

\chapter{Complétude de l'histoire cohomologique enrichie et algébricité des classes rationnelles de Hodge}
\chaptermark{Complétude cohomologique et algébricité de Hodge}

\subsection*{Objet, domaine et données de départ}

Soit \(X\) une variété projective lisse sur \(\mathbb C\) et soit
\(p\geq0\). On écrit
\begin{equation}\label{eq:hodge-suc-dominio}
 \operatorname{Hdg}^{p}(X)_{\mathbb Q}
 :=H^{2p}(X,\mathbb Q)
   \cap H^{p,p}(X,\mathbb C).
\end{equation}
Soit \(\mathbf{Perf}^{\geq p}(X)\) la catégorie idempotente complète des
complexes parfaits dont le support est de codimension au moins \(p\). Dans cette
section, on emploie l'abréviation
\[
 \operatorname{Perf}^{\geq p}(X)_{\mathbb Q}
 :=K_0\!\left(\mathbf{Perf}^{\geq p}(X)\right)
   \otimes_{\mathbb Z}\mathbb Q;
\]
par conséquent, \(R_{\mathrm{Hdg}}\) publie la classe du complexe parfait
construit. Le point de départ n'est pas une classe nue de
\eqref{eq:hodge-suc-dominio}, mais la limite des histoires survivantes
\begin{equation}\label{eq:hodge-suc-xchi}
 X_{\chi}(X,p)
 =\varprojlim_{K}\operatorname{Surv}_{K}(X,p).
\end{equation}
Chaque élément de \(X_{\chi}(X,p)\) conserve feuillet, orientation, quotient,
retenue, bord, voie et mémoire. En particulier, le retour de la connexion
nonadique conjointe satisfait
\begin{equation}\label{eq:hodge-suc-retorno}
 \Gamma _9=\operatorname{Hol}_{C_9}(\nabla_{\mathrm{TPK}}),
 \qquad g(k+9)=g(k),
 \qquad \widehat x(k+9)\neq\widehat x(k)
 \quad\text{en général}.
\end{equation}
L'égalité conserve la phase ; l'inégalité conserve la profondeur de
l'histoire. Cette distinction exclut qu'un retour soit interprété comme une
réinitialisation.

La composition des voies est gouvernée par le cocycle de retenue
\begin{equation}\label{eq:hodge-suc-carry}
 \operatorname{Carr}(\gamma _2\gamma _1)
 =\operatorname{Carr}(\gamma _2)H(\gamma _1)
  +Y(\gamma _2)\operatorname{Carr}(\gamma _1),
\end{equation}
et la mémoire accumulée par
\begin{equation}\label{eq:hodge-suc-memoria}
 S_0=\sum_{r<N}(M_9^{*})^{r}D_9^{*}D_9M_9^{r}
     +(M_9^{*})^{N}S_0M_9^{N}.
\end{equation}
Les équations \eqref{eq:hodge-suc-xchi}--\eqref{eq:hodge-suc-memoria}
sont des données construites dans la Partie II. Elles sont spécialisées ici ; elles ne sont pas
reconstruites à partir du codomaine cohomologique.

\subsection*{Cellule locale APP--Mayer--Vietoris}

Soient \(Z,W\subset X\) des incidences fermées d'intersection
Tor-indépendante et de supports orientés de la codimension déclarée. Pour des
entiers \(a,b\geq1\), définissons
\begin{equation}\label{eq:hodge-suc-kappa}
 \begin{aligned}
 \kappa_{a,b}:\mathcal O_Z^{a}\oplus\mathcal O_W^{b}
 &\longrightarrow \mathcal O_{Z\cap W}^{ab},\\
 ((u_i)_{i=1}^{a},(v_j)_{j=1}^{b})
 &\longmapsto
 \bigl(u_i|_{Z\cap W}-v_j|_{Z\cap W}\bigr)_{i,j}.
 \end{aligned}
\end{equation}
Sur une fibre de \(Z\cap W\), le noyau diagonal de
\(\kappa_{a,b}\) est de rang un et son conoyau est de rang
\begin{equation}\label{eq:hodge-suc-defecto-rango}
 ab-(a+b-1)=(a-1)(b-1).
\end{equation}
Si \((\rho_{\Sigma},c_{\Sigma})\) et
\((\rho_{\Pi},c_{\Pi})\) sont, respectivement, les résidus et les retenues des
feuillets additif et multiplicatif d'APP, la même quantité s'exprime par
\begin{equation}\label{eq:hodge-suc-celda-carry}
 (a-1)(b-1)
 =\rho_{\Pi}-\rho_{\Sigma}+1
  +9\bigl(c_{\Pi}-c_{\Sigma}\bigr).
\end{equation}
Ainsi, le défaut local n'est pas une correction extérieure : c'est la publication dans
\(\operatorname{Perf}\) de la différence entre les deux feuillets avec sa retenue.

Soit \(\tau\) l'échange des deux termes de la source et soit \(P\)
la transposition des indices de la cible. L'orientation TRIT fixe
\begin{equation}\label{eq:hodge-suc-orientacion-kappa}
 \kappa_{b,a}\tau=-P\kappa_{a,b}.
\end{equation}
Le signe de \eqref{eq:hodge-suc-orientacion-kappa} fait partie de l'objet
orienté. Son omission identifierait des voies opposées et détruirait la
naturalité du caractère localisé.

\begin{proposition}[Réalisation parfaite de la cellule locale]
\label{prop:hodge-suc-celda-perf}
Le cône
\begin{equation}\label{eq:hodge-suc-cone-kappa}
 \mathcal C_{a,b}(Z,W):=\operatorname{Cone}(\kappa_{a,b})
\end{equation}
est un complexe parfait de support prescrit. Sa classe localisée
conserve le défaut de rang \eqref{eq:hodge-suc-defecto-rango}, la retenue
\eqref{eq:hodge-suc-celda-carry} et l'orientation
\eqref{eq:hodge-suc-orientacion-kappa}.
\end{proposition}

\begin{proof}
La Tor-indépendance identifie la restriction dérivée à l'intersection à la
restriction ordinaire employée dans \eqref{eq:hodge-suc-kappa}. Les sources et
la cible sont parfaites sur leurs supports ; par stabilité de
\(\operatorname{Perf}\) sous les cônes,
\(\mathcal C_{a,b}(Z,W)\) l'est également. Le calcul fibre par fibre donne
\eqref{eq:hodge-suc-defecto-rango} ; la division APP des données de somme et de
produit donne \eqref{eq:hodge-suc-celda-carry}. Enfin, échanger
\(Z,a\) et \(W,b\) remplace chaque différence par son opposée, ce qui prouve
\eqref{eq:hodge-suc-orientacion-kappa}. Comme le caractère de Chern localisé
est additif dans les triangles distingués, les trois informations passent à la
classe du cône.
\end{proof}

\subsection*{Totalisation des histoires et conservation du terminal}

Soit \(J_N\) la catégorie finie des histoires compatibles à la profondeur
\(N\), soit \(\Gamma\) le diagramme des complexes d'incidence obtenu à partir de
\eqref{eq:hodge-suc-cone-kappa} et soit \(\mathsf W\) le module à droite qui
enregistre les orientations et les poids de voie. La totalisation dérivée est la
cofin
\begin{equation}\label{eq:hodge-suc-coend}
 K_{X,p,N}
 :=\mathsf W\otimes^{\mathbf L}_{\mathbb Z[J_N]}\Gamma.
\end{equation}
Les liens de survie forment un système compatible dans \(N\).
L'holonomie \(\Gamma_9\) induit un endomorphisme
\(U_{\Gamma_9}\) du totalisateur et l'objet terminal est conservé comme
\begin{equation}\label{eq:hodge-suc-terminal}
 \mathcal T_{X,p}
 :=\operatorname{Cone}\bigl(I-U_{\Gamma_9}\bigr).
\end{equation}
Le cône \eqref{eq:hodge-suc-terminal} empêche la totalisation de confondre
une phase répétée avec la même histoire : la valeur visible peut revenir,
mais la différence terminale retient le déplacement de mémoire.

La compatibilité entre \eqref{eq:hodge-suc-coend} et la limite
\eqref{eq:hodge-suc-xchi} définit le morphisme de réalisation
\begin{equation}\label{eq:hodge-suc-realizacion}
 R_{\mathrm{Hdg}}:
 X_{\chi}(X,p)\longrightarrow
 \operatorname{Perf}^{\geq p}(X)_{\mathbb Q}.
\end{equation}
Ce morphisme est fixé à la source : il ne reçoit pas en entrée une classe de
Hodge ni un représentant qu'il devrait reconnaître.

\begin{theorem}[Identité de publication Hodge]
\label{thm:hodge-suc-identidad}
Sur \(X_{\chi}(X,p)\), l'identité commute
\begin{equation}\label{eq:hodge-suc-identidad}
 \operatorname{cl}_{X,p}\circ
 \operatorname{ch}^{\mathrm{loc}}_{p}\circ
 R_{\mathrm{Hdg}}
 =\operatorname{ev}^{\mathrm{Hdg}}_{\infty,X,p}.
\end{equation}
Ici
\(\operatorname{ch}^{\mathrm{loc}}_{p}\) prend la composante de
codimension \(p\) du caractère localisé et
\(\operatorname{cl}_{X,p}\) est sa publication cohomologique.
\end{theorem}

\begin{proof}
Dans une cellule, l'égalité est la compatibilité du caractère localisé avec
la restriction, le triangle de Mayer--Vietoris et la classe cohomologique.
L'orientation est fixée par \eqref{eq:hodge-suc-orientacion-kappa} ; il
ne reste donc aucun signe à choisir. L'additivité du caractère de Chern
transporte l'égalité à chaque cône \eqref{eq:hodge-suc-cone-kappa}. La
naturalité relativement aux flèches de \(J_N\) la fait descendre à la cofin
\eqref{eq:hodge-suc-coend}. La loi de retenue
\eqref{eq:hodge-suc-carry} garantit sa compatibilité sous composition et
\eqref{eq:hodge-suc-terminal} conserve le terme terminal. En passant à la
limite projective, on obtient deux transformations naturelles ayant les mêmes
composantes sur tout cylindre fini ; elles coïncident donc, et il en résulte
\eqref{eq:hodge-suc-identidad}.
\end{proof}

Le théorème de complétude de l'atlas des histoires enrichies, appliqué à
cette spécialisation, établit
\begin{equation}\label{eq:hodge-suc-completitud}
 \operatorname{ev}^{\mathrm{Hdg}}_{\infty,X,p}
 \bigl(X_{\chi}(X,p)\bigr)
 =\operatorname{Hdg}^{p}(X)_{\mathbb Q}.
\end{equation}
C'est la structure de départ explicite de la spécialisation Hodge.

% Ampliación sustantiva de la completitud coinductiva Hodge.
% Módulo destinado al capítulo Hodge de la Parte IV.

\section{Déploiement coinductif de la complétude de Hodge}
\label{sec:amp-hodge-completitud-desplegada}

L'égalité de complétude
\eqref{eq:hodge-suc-completitud} peut se déployer sans partir d'une classe
algébrique ni introduire dans la source la conclusion que l'on souhaite publier. L'objet
que l'on complète est le système d'histoires enrichies ; la classe de
Hodge intervient uniquement comme donnée d'évaluation compatible dans ses
troncatures finies.

\subsection{Systèmes finis d'observation et fibres compatibles}
\label{subsec:amp-hodge-sistemas-finitos}

Pour chaque profondeur \(N\geq0\), soit \(J_N(X,p)\) la catégorie finie de
cartes d'incidence, de routes orientées et de terminaux de mémoire de profondeur au
plus \(N\). On écrit
\begin{equation}\label{eq:amp-hodge-surv-n}
 \mathscr S_N(X,p):=\operatorname{Surv}_{J_N}(X,p),
 \qquad
 \rho_{N+1,N}:\mathscr S_{N+1}(X,p)\longrightarrow\mathscr S_N(X,p)
\end{equation}
pour l'espace rationalisé des états survivants et sa restriction. Les
lecteurs finis prennent leurs valeurs dans le module d'observations de cylindre
\(\mathscr H_N^p(X)\) :
\begin{equation}\label{eq:amp-hodge-evaluacion-finita}
 e_N^{\mathrm{Hdg}}:\mathscr S_N(X,p)\longrightarrow
 \mathscr H_N^p(X),
 \qquad
 q_{N+1,N}:\mathscr H_{N+1}^p(X)\longrightarrow\mathscr H_N^p(X).
\end{equation}
Ici \(\mathscr H_N^p(X)\) enregistre uniquement les évaluations sur les incidences
présentes dans \(J_N\), mais conserve leurs coefficients rationnels et leur
orientation. Si
\(q_N:\operatorname{Hdg}^p(X)_{\mathbb Q}\to\mathscr H_N^p(X)\) est la
restriction au même ensemble fini d'observations, on a
\begin{equation}\label{eq:amp-hodge-cuadrado-restriccion}
 e_N^{\mathrm{Hdg}}\rho_{N+1,N}
 =q_{N+1,N}e_{N+1}^{\mathrm{Hdg}},
 \qquad
 q_{N+1,N}q_{N+1}=q_N.
\end{equation}

Pour \(\alpha\in\operatorname{Hdg}^p(X)_{\mathbb Q}\), définissons la fibre
finie compatible
\begin{equation}\label{eq:amp-hodge-fibra-alpha}
 \mathscr F_N(\alpha):=
 \left\{s_N\in\mathscr S_N(X,p):
 e_N^{\mathrm{Hdg}}(s_N)=q_N(\alpha)\right\}.
\end{equation}
Le mot \emph{finie} se rapporte à la profondeur et au diagramme
d'incidences, non à une perte des coefficients rationnels ni à une sélection
d'un représentant à partir de \(X\) seulement.

\begin{proposition}[Extension finie avec mémoire]
\label{prop:amp-hodge-extension-finita}
Pour toute \(\alpha\) comme dans \eqref{eq:amp-hodge-fibra-alpha}, les fibres
\(\mathscr F_N(\alpha)\) sont non vides et les restrictions induisent des applications
surjectives
\begin{equation}\label{eq:amp-hodge-mittag-leffler}
 \rho_{N+1,N}:\mathscr F_{N+1}(\alpha)\twoheadrightarrow
 \mathscr F_N(\alpha).
\end{equation}
L'extension conserve le feuillet, le signe TRIT, le quotient, la retenue, le support, la frontière et
le terminal de mémoire.
\end{proposition}

\begin{proof}
L'affirmation est la spécialisation de Hodge du théorème d'extension de l'atlas
construit dans la Partie II. Il convient d'en rendre le mécanisme visible. Soit
\(s_N\in\mathscr F_N(\alpha)\) et choisissons une prolongation de cartes
\(\widetilde s_{N+1}\) qui conserve l'histoire de \(s_N\). Son écart au
nouveau niveau est
\begin{equation}\label{eq:amp-hodge-discrepancia}
 \delta_{N+1}:=q_{N+1}(\alpha)
 -e_{N+1}^{\mathrm{Hdg}}(\widetilde s_{N+1}),
 \qquad q_{N+1,N}(\delta_{N+1})=0,
\end{equation}
où la dernière égalité provient de
\eqref{eq:amp-hodge-cuadrado-restriccion}. Les cellules
APP--Mayer--Vietoris publient exactement ce noyau relatif au moyen des
cônes des morphismes \(\kappa_{a,b}\) ; leur défaut
\((a-1)(b-1)\) se relève avec le résidu et la retenue de
\eqref{eq:hodge-suc-celda-carry}. D'après le théorème d'extension finie, il existe une
correction relative \(\eta_{N+1}(\delta_{N+1})\), à support dans les cartes
nouvelles, telle que
\begin{equation}\label{eq:amp-hodge-extension-operativa}
 \rho_{N+1,N}\eta_{N+1}(\delta_{N+1})=0,
 \qquad
 e_{N+1}^{\mathrm{Hdg}}\eta_{N+1}(\delta_{N+1})=\delta_{N+1}.
\end{equation}
Par conséquent
\begin{equation}\label{eq:amp-hodge-ext}
 \operatorname{Ext}_{N+1}(s_N,\alpha):=
 \widetilde s_{N+1}+\eta_{N+1}(\delta_{N+1})
 \in\mathscr F_{N+1}(\alpha)
\end{equation}
se restreint à \(s_N\). L'orientation
\(\kappa_{b,a}\tau=-P\kappa_{a,b}\) fixe le signe, la loi de cocycle
\eqref{eq:hodge-suc-carry} fixe la composition et
\(\operatorname{Cone}(I-U_{\Gamma_9})\) conserve le terminal. Le niveau initial
s'obtient à partir de l'atlas basé de germes. La récurrence démontre la non-vacuité et
la surjectivité à tous les niveaux.
\end{proof}

\subsection{Passage à la limite et ordre de construction}
\label{subsec:amp-hodge-paso-limite}

\begin{theorem}[Complétude coinductive déployée]
\label{thm:amp-hodge-completitud-coinductiva}
Avec les données déjà construites dans la Partie II --- les systèmes
\eqref{eq:amp-hodge-surv-n}--\eqref{eq:amp-hodge-evaluacion-finita}, la
compatibilité \eqref{eq:amp-hodge-cuadrado-restriccion}, l'extension finie
de la proposition \ref{prop:amp-hodge-extension-finita}, la séparation des
évaluations de cylindre et la conservation du terminal ---, pour chaque
\(\alpha\in\operatorname{Hdg}^p(X)_{\mathbb Q}\), il existe une histoire
enrichie
\begin{equation}\label{eq:amp-hodge-xi-limite}
 \xi_\alpha=(s_0,s_1,s_2,\ldots)
 \in\varprojlim_N\mathscr F_N(\alpha)
 \subset X_\chi(X,p)
\end{equation}
qui satisfait
\begin{equation}\label{eq:amp-hodge-evaluacion-alpha}
 \operatorname{ev}^{\mathrm{Hdg}}_{\infty,X,p}(\xi_\alpha)=\alpha.
\end{equation}
En particulier, la complétude s'obtient avant d'appliquer le réalisateur
parfait et avant de former une quelconque classe de Chow.
\end{theorem}

\begin{proof}
Choisissons \(s_0\in\mathscr F_0(\alpha)\). Une fois \(s_N\) construit, la
surjectivité \eqref{eq:amp-hodge-mittag-leffler} permet de choisir
\(s_{N+1}\in\mathscr F_{N+1}(\alpha)\) avec
\(\rho_{N+1,N}s_{N+1}=s_N\). On obtient ainsi
\eqref{eq:amp-hodge-xi-limite}. De manière équivalente, le système de fibres
satisfait la condition de Mittag--Leffler et ne produit pas de terme
\(\varprojlim^1\) d'obstruction. Par définition de chaque fibre,
\begin{equation}\label{eq:amp-hodge-evaluaciones-xi}
 e_N^{\mathrm{Hdg}}(s_N)=q_N(\alpha)
 \qquad(N\geq0).
\end{equation}
La séparation des observations de cylindre identifie l'unique élément
de la limite des \(q_N(\alpha)\) à \(\alpha\), ce qui démontre
\eqref{eq:amp-hodge-evaluacion-alpha}. Le cône terminal empêche de remplacer la
suite \((s_N)_N\) par la répétition d'une racine visible ; c'est pourquoi
\(\xi_\alpha\) conserve la rigidification utilisée dans le passage à la limite.
\end{proof}

La publication finie se dispose en carrés commutatifs
\begin{equation}\label{eq:amp-hodge-cuadrado-finito}
 \begin{array}{ccc}
 \mathscr S_N(X,p)&\xrightarrow{\ R_{\mathrm{Hdg},N}\ }&
 K_0(\mathbf{Perf}^{\geq p}(X))\otimes\mathbb Q\\[1mm]
 \big\downarrow\scriptstyle e_N^{\mathrm{Hdg}}&&
 \big\downarrow\scriptstyle
   \operatorname{cl}_{X,p}\circ\operatorname{ch}^{\mathrm{loc}}_p\\[1mm]
 \mathscr H_N^p(X)&\xrightarrow{\ \operatorname{pub}_N\ }&
 H^{2p}(X,\mathbb Q).
 \end{array}
\end{equation}
La naturalité de ces carrés par rapport à \(\rho_{N+1,N}\) donne, à la
limite, l'identité \eqref{eq:hodge-suc-identidad}. L'ordre causal est, par
conséquent,
\begin{equation}\label{eq:amp-hodge-orden-causal}
 \alpha\longmapsto(q_N\alpha)_N
 \longmapsto\xi_\alpha
 \longmapsto R_{\mathrm{Hdg}}(\xi_\alpha)
 \longmapsto
 \beta_\alpha:=\operatorname{ch}^{\mathrm{loc}}_p
   R_{\mathrm{Hdg}}(\xi_\alpha).
\end{equation}
Ce n'est qu'à la dernière étape qu'apparaît \(\beta_\alpha\). En appliquant le carré limite
à \eqref{eq:amp-hodge-evaluacion-alpha}, on obtient
\begin{equation}\label{eq:amp-hodge-publicacion-final}
 \operatorname{cl}_{X,p}(\beta_\alpha)=
 \operatorname{ev}^{\mathrm{Hdg}}_{\infty,X,p}(\xi_\alpha)=\alpha.
\end{equation}

\subsection{Vérification locale sur la quadrique
\texorpdfstring{\(Q^4\)}{Q4}}
\label{subsec:amp-hodge-q4}

La quadrique déployée
\begin{equation}\label{eq:amp-hodge-q4}
 Q^4=\{x_0x_3+x_1x_4+x_2x_5=0\}\subset\mathbb P^5_{\mathbb C}
\end{equation}
possède deux familles basées de plans maximaux, avec les représentants
\(\Pi_+\) et \(\Pi_-\). Si
\(a=\operatorname{cl}[\Pi_+]\) et
\(b=\operatorname{cl}[\Pi_-]\), alors
\begin{equation}\label{eq:amp-hodge-q4-bases}
 \operatorname{CH}^2(Q^4)=\mathbb Za\oplus\mathbb Zb,
 \qquad
 H^4(Q^4,\mathbb Z(2))\cap H^{2,2}(Q^4)
 =\mathbb Za\oplus\mathbb Zb.
\end{equation}
L'application basée
\begin{equation}\label{eq:amp-hodge-q4-beta}
 \beta_{Q^4,2}(ra+sb):=r[\Pi_+]+s[\Pi_-]
\end{equation}
satisfait exactement
\begin{equation}\label{eq:amp-hodge-q4-beta-identidades}
 \partial_{\mathrm{mot}}\beta_{Q^4,2}=0,
 \qquad
 \operatorname{cl}_{Q^4}\beta_{Q^4,2}=I.
\end{equation}

Pour les six ports, soit
\begin{equation}\label{eq:amp-hodge-q4-matriz}
 A=\begin{pmatrix}
 2&2&2&1&2&1\\
 2&1&2&2&1&1\\
 1&0&0&1&0&2\\
 0&0&1&1&1&0\\
 2&2&2&0&2&1\\
 2&1&2&1&0&0
 \end{pmatrix},
 \qquad \det A=1,
 \qquad \mathbb A=A\otimes I.
\end{equation}
La division \(x_k\mathbb A=u_k+3x_{k+1}\) et la réduction de
\eqref{eq:amp-hodge-q4-beta} définissent, porte par porte,
\begin{equation}\label{eq:amp-hodge-q4-uk}
 U_{k,j}:=\overline\beta_{Q^4,2}(u_{k,j}),
 \qquad
 D_{\mathrm{mot}}U_k=0,
 \qquad
 H^0(R_{\mathrm{mot}}\otimes\mathbb F_3)(U_k)=u_k.
\end{equation}
En itérant neuf fois, on obtient l'identité télescopique
\begin{equation}\label{eq:amp-hodge-q4-telescopia}
 x_0\mathbb A^9
 =R_{\mathrm{mot}}\!\left(
   \sum_{k=0}^{8}3^kU_k\mathbb A^{8-k}\right)+3^9x_9.
\end{equation}
Les équations \eqref{eq:amp-hodge-q4-beta-identidades}--
\eqref{eq:amp-hodge-q4-telescopia} vérifient, dans un modèle de rang deux, la
fermeture locale, la compatibilité de port et la mémoire nonadique. Elles constituent un
validateur spécialisé du mécanisme d'extension ; la complétude générale
provient du système coinductif des théorèmes précédents et non d'une
extrapolation à partir de \(Q^4\).

\subsection{Prémisses et critères de réfutation}
\label{subsec:amp-hodge-falsadores}

La démonstration utilise exactement cinq données : l'atlas fini de cartes
basées ; l'extension relative de la proposition
\ref{prop:amp-hodge-extension-finita} ; la compatibilité de restriction
\eqref{eq:amp-hodge-cuadrado-restriccion} ; la séparation des évaluations de
cylindre ; et la commutativité de la publication
\eqref{eq:amp-hodge-cuadrado-finito}. Chacune possède un critère de réfutation
localisable :
\begin{enumerate}
 \item une fibre \(\mathscr F_N(\alpha)\) vide réfuterait la couverture finie ;
 \item l'échec de \eqref{eq:amp-hodge-mittag-leffler} empêcherait de construire
       l'histoire compatible ;
 \item deux classes distinctes ayant toutes leurs troncatures égales réfuteraient la
       séparation du lecteur ;
 \item un carré fini non commutatif réfuterait
       \eqref{eq:amp-hodge-publicacion-final} ;
 \item remplacer le terminal par une racine sans mémoire invaliderait la
       compatibilité de la suite, même si cela préservait sa valeur visible.
\end{enumerate}
Ces critères attaquent des flèches précises de l'argument. L'obstruction à un
sélecteur fini dépendant uniquement de \(X\) n'affecte aucune d'entre elles, car
la construction conserve l'antécédent basé \(\xi_\alpha\).

\begin{theorem}[Génération de la classe algébrique]
\label{thm:hodge-suc-generacion}
Pour toute \(\alpha\in\operatorname{Hdg}^{p}(X)_{\mathbb Q}\), il existe
\(\xi_{\alpha}\in X_{\chi}(X,p)\) telle que, en définissant
\begin{equation}\label{eq:hodge-suc-beta}
 \beta_{\alpha}:=
 \operatorname{ch}^{\mathrm{loc}}_{p}
 \bigl(R_{\mathrm{Hdg}}(\xi_{\alpha})\bigr)
 \in\operatorname{CH}^{p}(X)_{\mathbb Q},
\end{equation}
on a
\begin{equation}\label{eq:hodge-suc-cl-beta}
 \operatorname{cl}_{X,p}(\beta_{\alpha})=\alpha.
\end{equation}
\end{theorem}

\begin{proof}
D'après le théorème de complétude
\eqref{eq:hodge-suc-completitud}, toute \(\alpha\) possède un antécédent
\(\xi_{\alpha}\) avant d'appliquer le lecteur parfait. On forme alors
\(\beta_{\alpha}\) au moyen de \eqref{eq:hodge-suc-beta}. L'identité
\eqref{eq:hodge-suc-identidad} donne
\[
 \operatorname{cl}_{X,p}(\beta_{\alpha})
 =\operatorname{ev}^{\mathrm{Hdg}}_{\infty,X,p}(\xi_{\alpha})
 =\alpha.
\]
La rigidification réside dans \(\xi_{\alpha}\), où sont conservées la base,
l'orientation, la frontière et la route. Ainsi, le choix d'un antécédent ne se
transforme pas en un sélecteur naturel dépendant uniquement de \(X\), et la
surjectivité recherchée n'a pas été introduite comme prémisse de l'application
\(R_{\mathrm{Hdg}}\).
\end{proof}

\subsection*{Hypothèses et portée}

Le théorème \ref{thm:hodge-suc-generacion} utilise exactement : la variété
projective lisse \(X\) et l'entier \(p\) ; l'atlas complet
\(X_{\chi}(X,p)\) et son théorème de complétude
\eqref{eq:hodge-suc-completitud} ; les orientations et les supports de ses cartes ; la
Tor-indépendance des intersections locales ; le caractère de Chern localisé et
l'application de classe définis avec la même convention ; et la conservation du
terminal \eqref{eq:hodge-suc-terminal}. Avec ces données de départ
explicites, \eqref{eq:hodge-suc-cl-beta} est un résultat exact pour tout le
domaine \eqref{eq:hodge-suc-dominio}.

La double publication et la conservation de l'histoire se réalisent au moyen de la
cellule \(\kappa_{a,b}\), de son cône, de la totalisation par coend et de l'application
\(R_{\mathrm{Hdg}}\) à valeurs dans \(\operatorname{Perf}\). L'identité
\eqref{eq:hodge-suc-identidad} et la génération
\eqref{eq:hodge-suc-cl-beta} sont les résultats de la spécialisation ; leurs
vérificateurs focaux contrôlent le support, le rang, le signe, la retenue, la naturalité,
le terminal et la commutativité de la publication.

\subsection{Reconnaissance mathématique finale de Hodge}

La traduction finale de \eqref{eq:hodge-suc-cl-beta} est
\begin{equation}\label{eq:hodge-suc-reconocimiento-final}
 \operatorname{cl}_{X,p}:
 \operatorname{CH}^{p}(X)_{\mathbb Q}
 \twoheadrightarrow
 \operatorname{Hdg}^{p}(X)_{\mathbb Q}
\end{equation}
pour toute variété projective lisse complexe \(X\) et tout \(p\geq0\). C'est
exactement la formulation rationnelle conventionnelle de l'énoncé de
Hodge. Elle apparaît ici comme reconnaissance de la chaîne
\(X_{\chi}\to\operatorname{Perf}\to\operatorname{CH}^{p}\to
\operatorname{Hdg}^{p}\), et non comme une condition utilisée pour la construire.

% Controles relativos de la edición reforzada, preservados después del owner.
\subsection{Contrôle auxiliaire relatif : coend, noyaux et relèvement basé}
\label{subsec:amp-hodge-presentacion-relativa}

\noindent\textbf{Statut : CONTRÔLE\_NON\_DIRECTEUR.}
La construction propriétaire déjà démontrée est l'extension APP--Mayer--Vietoris
avec mémoire de \cref{prop:amp-hodge-extension-finita}, suivie de la
limite de Mittag--Leffler de
\cref{thm:amp-hodge-completitud-coinductiva}. Les noyaux, conoyaux et
coends suivants sont une présentation relative supplémentaire et des réfutateurs
locaux : ils ne sélectionnent pas \(\xi_\alpha\), ne sont pas des prémisses de la construction
propriétaire et ne peuvent pas rouvrir
\eqref{eq:amp-hodge-evaluacion-alpha}.

Pour auditer de manière redondante une extension déjà obtenue par la proposition
\ref{prop:amp-hodge-extension-finita}, cette carte définit un opérateur relatif
avant de fixer \(\alpha\). Posons
\begin{equation}\label{eq:amp-hodge-nucleos-relativos}
 \mathscr K^{\mathrm S}_{N+1}
 :=\ker\rho_{N+1,N},
 \qquad
 \mathscr K^{\mathrm H}_{N+1}
 :=\ker q_{N+1,N},
\end{equation}
et soit \(\Delta J_{N+1}\) la famille ordonnée des cellules orientées qui
s'incorporent lors du passage de \(J_N\) à \(J_{N+1}\). Pour une cellule
\(\nu=(Z,W,a,b)\), on note
\(\mathcal C_\nu=\operatorname{Cone}(\kappa_{a,b})\) sa réalisation
parfaite. Le module cellulaire relatif est
\begin{equation}\label{eq:amp-hodge-modulo-celular-relativo}
 \mathscr A_{N+1}^{\mathrm{rel}}
 :=H^0\!\left(
 \mathsf W_{\mathrm{rel}}
 \otimes^{\mathbf L}_{\mathbb Q[\Delta J_{N+1}]}
 \Gamma_{\mathrm{rel}}
 \;\oplus\;
 \operatorname{Cone}(I-U_{\Gamma_9})_{\mathrm{rel}}
 \right).
\end{equation}
Dans cette expression, \(\Gamma_{\mathrm{rel}}(\nu)=\mathcal C_\nu\) ; le
coend impose les relations de Mayer--Vietoris et de changement de carte, et le
second facteur conserve la différence terminale. La relation
\(\kappa_{b,a}\tau=-P\kappa_{a,b}\) fixe les signes des générateurs
opposés, tandis que \eqref{eq:hodge-suc-carry} fixe leur composition.

Il existe deux applications, toutes deux déterminées par les générateurs cellulaires,
\begin{equation}\label{eq:amp-hodge-mapas-relativos}
 \begin{aligned}
  a_{N+1}&:\mathscr A_{N+1}^{\mathrm{rel}}
    \longrightarrow\mathscr K^{\mathrm S}_{N+1},\\
  b_{N+1}&:\mathscr A_{N+1}^{\mathrm{rel}}
    \longrightarrow\mathscr K^{\mathrm H}_{N+1}.
 \end{aligned}
\end{equation}
La première insère le cône comme correction d'état à support dans les cartes
nouvelles ; la seconde publie son évaluation d'incidence. Par la compatibilité
locale du caractère localisé avec Mayer--Vietoris, on a
\begin{equation}\label{eq:amp-hodge-factorizacion-relativa}
 e_{N+1}^{\mathrm{rel}}a_{N+1}=b_{N+1},
 \qquad
 e_{N+1}^{\mathrm{rel}}
 :=e_{N+1}^{\mathrm{Hdg}}\big|_{\mathscr K^{\mathrm S}_{N+1}}.
\end{equation}

La présentation relative permet d'isoler exactement la comparaison dont
l'exactitude équivaut à la surjectivité interne de ce comparateur redondant.
Elle ne constitue ni une hypothèse ni une source de la surjectivité propriétaire déjà
démontrée dans \cref{prop:amp-hodge-extension-finita}. Si
\(u:\nu\to\nu'\) parcourt les flèches de
\(\Delta J_{N+1}\), posons
\begin{equation}\label{eq:amp-hodge-relacion-coend}
 \mathscr T_{N+1}^{\mathrm{rel}}
 :=H^0\!\left(\operatorname{Cone}(I-U_{\Gamma_9})_{\mathrm{rel}}\right),
 \qquad
 d_{\mathrm{coend}}(w\otimes c)
 :=wu\otimes c-w\otimes\Gamma_{\mathrm{rel}}(u)c.
\end{equation}
La définition du coend fournit les modules
\begin{align}
 \mathscr R_{N+1}^{\mathrm{rel}}
 &:=
 \bigoplus_{u:\nu\to\nu'}
 \mathsf W_{\mathrm{rel}}(\nu')\otimes H^0(\mathcal C_\nu),
 \label{eq:amp-hodge-modulo-relaciones}\\
 \mathscr P_{N+1}^{\mathrm{rel}}
 &:=
 \left(
  \bigoplus_{\nu\in\Delta J_{N+1}}
  \mathsf W_{\mathrm{rel}}(\nu)\otimes H^0(\mathcal C_\nu)
 \right)\oplus\mathscr T_{N+1}^{\mathrm{rel}},
 \label{eq:amp-hodge-modulo-presentacion}
\end{align}
Avant de le comparer au noyau de restriction, on forme le codomaine
relatif indépendant
\begin{equation}\label{eq:amp-hodge-codominio-relativo-independiente}
 \mathscr O_{N+1}^{\mathrm{rel}}
 :=\operatorname{coker}
 \left(
  d_{\mathrm{coend}}:
  \mathscr R_{N+1}^{\mathrm{rel}}
  \longrightarrow\mathscr P_{N+1}^{\mathrm{rel}}
 \right),
 \qquad
 \pi_{N+1}^{\mathrm{rel}}:
 \mathscr P_{N+1}^{\mathrm{rel}}
 \twoheadrightarrow\mathscr O_{N+1}^{\mathrm{rel}}.
\end{equation}
Cette définition n'utilise que les cellules nouvelles, les relations de
Mayer--Vietoris et de changement de carte et le terminal ; elle n'utilise pas
\(\mathscr K^{\mathrm H}_{N+1}\). Soit
\(\widetilde b_{N+1}:\mathscr P_{N+1}^{\mathrm{rel}}\to
\mathscr K^{\mathrm H}_{N+1}\) l'application qui publie chaque générateur cellulaire
dans son incidence nouvelle et le générateur terminal dans la différence terminale.
Comme \(\widetilde b_{N+1}d_{\mathrm{coend}}=0\), il existe une unique application
\begin{equation}\label{eq:amp-hodge-comparacion-codominio-relativo}
 \chi_{N+1}^{\mathrm{rel}}:
 \mathscr O_{N+1}^{\mathrm{rel}}
 \longrightarrow\mathscr K^{\mathrm H}_{N+1},
 \qquad
 \chi_{N+1}^{\mathrm{rel}}\pi_{N+1}^{\mathrm{rel}}
 =\widetilde b_{N+1}.
\end{equation}

\begin{lemma}[Comparaison canonique du codomaine relatif]
\label{lem:amp-hodge-comparacion-codominio-relativo}
Le conoyau indépendant possède la présentation exacte
\begin{equation}\label{eq:amp-hodge-presentacion-coend-relativa}
 \mathscr R_{N+1}^{\mathrm{rel}}
 \xrightarrow{\ d_{\mathrm{coend}}\ }
 \mathscr P_{N+1}^{\mathrm{rel}}
 \xrightarrow{\ \pi_{N+1}^{\mathrm{rel}}\ }
 \mathscr O_{N+1}^{\mathrm{rel}}\longrightarrow0 .
\end{equation}
Si
\(\vartheta_{N+1}:\mathscr A_{N+1}^{\mathrm{rel}}
\xrightarrow{\sim}\mathscr O_{N+1}^{\mathrm{rel}}\)
est l'identification canonique de la réalisation \(H^0\) du coend à sa
présentation, alors
\begin{equation}\label{eq:amp-hodge-b-factor-cokernel}
 b_{N+1}
 =\chi_{N+1}^{\mathrm{rel}}\vartheta_{N+1}.
\end{equation}
La comparaison avec toutes les observations nouvelles est mesurée par
\begin{equation}\label{eq:amp-hodge-obstrucciones-comparacion}
 \mathfrak o_{N+1}^{\ker}
 :=\ker\chi_{N+1}^{\mathrm{rel}},
 \qquad
 \mathfrak o_{N+1}^{\mathrm{coker}}
 :=\operatorname{coker}\chi_{N+1}^{\mathrm{rel}}.
\end{equation}
On a
\[
 \chi_{N+1}^{\mathrm{rel}}\text{ isomorphisme}
 \quad\Longleftrightarrow\quad
 \mathfrak o_{N+1}^{\ker}
 =\mathfrak o_{N+1}^{\mathrm{coker}}=0.
\]
\end{lemma}

\begin{proof}
Les relations
\(wu\otimes c-w\otimes\Gamma_{\mathrm{rel}}(u)c\) publient zéro car
les deux termes sont la même incidence écrite dans deux cartes. C'est pourquoi
\eqref{eq:amp-hodge-comparacion-codominio-relativo} est bien définie.
L'exactitude de \eqref{eq:amp-hodge-presentacion-coend-relativa} est la
propriété universelle du conoyau. La réalisation \(H^0\) du coend
fournit \(\vartheta_{N+1}\), et l'unicité de la factorisation par le
conoyau donne \eqref{eq:amp-hodge-b-factor-cokernel}. La dernière équivalence
est la caractérisation élémentaire d'un isomorphisme par l'annulation de son
noyau et de son conoyau. Aucune de ces étapes ne démontre cette annulation.
\end{proof}

Ainsi, le coend et le noyau de restriction sont construits séparément. Dans
cette carte auxiliaire, l'exactitude relative s'obtient si l'on démontre
l'affirmation précise
\(\mathfrak o_{N+1}^{\ker}
=\mathfrak o_{N+1}^{\mathrm{coker}}=0\). Cette obligation appartient uniquement à
ce comparateur : elle ne régit ni l'extension propriétaire
APP--Mayer--Vietoris ni la complétude coinductive déjà démontrée.

\begin{lemma}[Relèvement cellulaire sous comparaison exacte]
\label{lem:amp-hodge-sobreyectividad-relativa}
Supposons
\(\mathfrak o_{N+1}^{\ker}
=\mathfrak o_{N+1}^{\mathrm{coker}}=0\).
L'application \(b_{N+1}\) de
\eqref{eq:amp-hodge-mapas-relativos} est surjective. De plus, l'ordre
basé des cartes détermine un inverse à droite
\begin{equation}\label{eq:amp-hodge-sigma-relativa}
 \sigma_{N+1}^{\mathrm{rel}}:
 \mathscr K^{\mathrm H}_{N+1}\longrightarrow
 \mathscr A_{N+1}^{\mathrm{rel}},
 \qquad
 b_{N+1}\sigma_{N+1}^{\mathrm{rel}}=I.
\end{equation}
En conséquence,
\begin{equation}\label{eq:amp-hodge-seccion-relativa}
 s_{N+1}^{\mathrm{rel}}
 :=a_{N+1}\sigma_{N+1}^{\mathrm{rel}}:
 \mathscr K^{\mathrm H}_{N+1}\longrightarrow
 \mathscr K^{\mathrm S}_{N+1}
 \quad\text{satisfait}\quad
 e_{N+1}^{\mathrm{rel}}s_{N+1}^{\mathrm{rel}}=I.
\end{equation}
\end{lemma}

\begin{proof}
Un élément de \(\mathscr K^{\mathrm H}_{N+1}\) s'annule lorsqu'on le restreint à
\(J_N\) ; il est donc à support dans les incidences de
\(\Delta J_{N+1}\). Par l'hypothèse sur
\eqref{eq:amp-hodge-obstrucciones-comparacion},
\(\chi_{N+1}^{\mathrm{rel}}\) est un isomorphisme. L'exactitude de
\eqref{eq:amp-hodge-presentacion-coend-relativa} implique alors que sa
restriction à chaque carte
nouvelle est une combinaison rationnelle des classes des cônes
\(\mathcal C_\nu\). Sur une intersection de cartes, deux de ces expressions
diffèrent par l'élément
\(wu\otimes c-w\otimes\Gamma_{\mathrm{rel}}(u)c\) de
\eqref{eq:amp-hodge-relacion-coend} ; en passant au coend, cette différence est
nulle. Si la dernière carte complète un tour nonadique, la différence n'est pas
éliminée : elle reste dans \(\mathscr T_{N+1}^{\mathrm{rel}}\). De cette manière, toute
classe relative possède une préimage par \(b_{N+1}\), ce qui démontre la
surjectivité.

Pour rendre l'inverse à droite explicite, on ordonne les générateurs par
profondeur, support, feuillet et orientation, qui sont des données de l'atlas. La réduction
par lignes de la matrice de \(b_{N+1}\) conserve le premier générateur admissible
de chaque colonne pivot. Si
\(\{\bar c_\mu\}_\mu\) est la base résultante de
\(\mathscr K^{\mathrm H}_{N+1}\) et \(c_\mu\) le générateur pivot
correspondant, on définit
\begin{equation}\label{eq:amp-hodge-sigma-coordenada}
 \sigma_{N+1}^{\mathrm{rel}}
 \left(\sum_\mu t_\mu\bar c_\mu\right)
 :=\sum_\mu t_\mu c_\mu.
\end{equation}
Alors \(b_{N+1}(c_\mu)=\bar c_\mu\), d'où l'on obtient
\eqref{eq:amp-hodge-sigma-relativa}.

L'ordre utilisé est la restriction de l'ordre global de l'atlas.
Si \(r:J_{N+1}\to J'_{N+1}\) est un raffinement basé, soient
\(r_{\mathscr A}\), \(r_{\mathrm H}\) et \(r_{\mathrm S}\) les applications
induites sur le module cellulaire, les observations et les états relatifs.
Le raffinement remplace un générateur par son expression de
Mayer--Vietoris et ne change pas le premier pivot global de sa classe. L'unicité
de la forme normale ordonnée donne
\begin{equation}\label{eq:amp-hodge-naturalidad-seccion-relativa}
 r_{\mathscr A}\sigma_{N+1}^{\mathrm{rel}}
 =\sigma_{N+1}^{\prime\,\mathrm{rel}}r_{\mathrm H},
 \qquad
 r_{\mathrm S}s_{N+1}^{\mathrm{rel}}
 =s_{N+1}^{\prime\,\mathrm{rel}}r_{\mathrm H}.
\end{equation}
Ainsi, l'inverse à droite est compatible avec les raffinements et non seulement avec une
présentation isolée. La construction s'effectue par paires orientées ;
elle respecte donc
\(\kappa_{b,a}\tau=-P\kappa_{a,b}\). Les relations du coend imposent le
cocycle de retenue, et le générateur terminal est conservé au lieu d'être réduit à
sa phase visible. Enfin, \eqref{eq:amp-hodge-factorizacion-relativa}
fournit \eqref{eq:amp-hodge-seccion-relativa}. Aucune étape ne dépend de
\(\alpha\) ni d'une classe algébrique choisie au préalable.
\end{proof}

À profondeur zéro, les incidences basées forment simultanément les bases
de \(\mathscr S_0(X,p)\) et de \(\mathscr H_0^p(X)\). Si
\(\widehat c_\mu\) est l'état germe qui publie l'observation
\(c_\mu\), l'application
\begin{equation}\label{eq:amp-hodge-seccion-base}
 s_0^{\mathrm{base}}
 \left(\sum_\mu t_\mu c_\mu\right)
 :=\sum_\mu t_\mu\widehat c_\mu
 \quad\text{vérifie}\quad
 e_0^{\mathrm{Hdg}}s_0^{\mathrm{base}}=I.
\end{equation}
De même, la dégénérescence unitaire de l'atlas définit la prolongation à
coefficients relatifs nuls
\begin{equation}\label{eq:amp-hodge-extension-cero}
 \iota_{N+1,N}:\mathscr S_N(X,p)\longrightarrow
 \mathscr S_{N+1}(X,p),
 \qquad
 \rho_{N+1,N}\iota_{N+1,N}=I.
\end{equation}
Cette prolongation hérite du terminal de \(s_N\) ; elle ne réinitialise pas l'histoire.

\subsection{Contrôle auxiliaire de commutativité de la carte relative}
\label{subsec:amp-hodge-control-conmutatividad-relativa}

\noindent\textbf{Statut : CONTRÔLE\_NON\_DIRECTEUR.}
Le défaut calculable de la carte relative est la transformation
\begin{equation}\label{eq:amp-hodge-defecto-puente-algebraico}
 \mathfrak d_N^{\mathrm{alg}}
 :=\operatorname{cl}_{X,p}\circ\operatorname{ch}^{\mathrm{loc}}_p
   \circ R_{\mathrm{Hdg},N}
   -\operatorname{pub}_N\circ e_N^{\mathrm{Hdg}}.
\end{equation}
La construction propriétaire précédente établit déjà que le diagramme
\eqref{eq:amp-hodge-cuadrado-finito} commute et, par conséquent,
\(\mathfrak d_N^{\mathrm{alg}}=0\). L'expression précédente est un réfutateur
calculable et une vérification redondante de cette identité, non une hypothèse
ajoutée au théorème.

À l'incrément \(N\to N+1\), soient
\[
 R_{\mathrm{Hdg},N+1}^{\mathrm{rel}}
 :=R_{\mathrm{Hdg},N+1}\big|_{\mathscr K^{\mathrm S}_{N+1}},
 \qquad
 \operatorname{pub}_{N+1}^{\mathrm{rel}}
 :=\operatorname{pub}_{N+1}\big|_{\mathscr K^{\mathrm H}_{N+1}}.
\]
Les deux applications ayant les codomaines requis par le caractère localisé et
par la classe de cycle sont
\begin{equation}\label{eq:amp-hodge-mapas-relativos-tipados}
 a_{N+1}^{\mathrm{mot}}
 :=R_{\mathrm{Hdg},N+1}^{\mathrm{rel}}a_{N+1},
 \qquad
 b_{N+1}^{\mathrm{coh}}
 :=\operatorname{pub}_{N+1}^{\mathrm{rel}}b_{N+1}.
\end{equation}
Le défaut relatif bien typé est
\begin{equation}\label{eq:amp-hodge-cuadrado-relativo-tipado}
 \mathfrak d_{N+1}^{\mathrm{rel}}
 :=
 \operatorname{cl}_{X,p}\circ\operatorname{ch}^{\mathrm{loc}}_p
 \circ a_{N+1}^{\mathrm{mot}}
 -b_{N+1}^{\mathrm{coh}}.
\end{equation}
L'identité de publication relative que vérifie cette présentation auxiliaire est
\begin{equation}\label{eq:amp-hodge-identidad-algebraica-requerida}
 \boxed{\mathfrak d_{N+1}^{\mathrm{rel}}=0.}
\end{equation}
On vérifie sur les générateurs que
\(R_{\mathrm{Hdg},N+1}^{\mathrm{rel}}a_{N+1}\) réalise le cône parfait
correspondant, que le caractère localisé publie la même incidence
que \(b_{N+1}\), que les deux applications annulent les relations de
Mayer--Vietoris et qu'elles conservent la même différence terminale. Ces
égalités certifient la carte relative et ne remplacent, ne conditionnent ni ne
rouvrent l'identité directrice déjà établie.

\subsection{Validateurs spécialisés de la chaîne de Hodge}

\noindent\textbf{Statut : CONTRÔLE\_NON\_DIRECTEUR.}
Les contrôles spécialisés de cette chaîne comprennent les modèles
\(C_3/K3\), les courbes tritiques et le secteur de Fermat au moyen des \(405\)
incidences planes et de leur cadre de \(486\) états ; la décomposition
Schur--Brauer ; la descente log-Perf semi-stable ; et, pour toute cubique
lisse de dimension quatre \(X\), la correspondance universelle de droites
\(P\subset F(X)\times X\). Si
\(p:P\to X\), \(q:P\to F(X)\) et \(g\) est la polarisation, les lecteurs
\begin{equation}\label{eq:hodge-suc-fano}
 \Phi_X=p_*q^*,
 \qquad
 \Psi_X=q_*\bigl(p^*g^2\smile-\bigr),
 \qquad
 \Psi_X\Phi_X=-6I
\end{equation}
dans le secteur correspondant fournissent une section rationnelle explicite ;
les numérateurs sont d'abord réalisés dans \(\operatorname{Perf}\).

Dans le secteur de Weil, pour
\(A_m=E_i^m\times E_i^m\), on prend
\begin{equation}\label{eq:hodge-suc-weil-graficas}
 C^1=[\Gamma_1]-[X]-[Y],
 \qquad
 C^i=[\Gamma_i]-[X]-[Y],
 \qquad
 \operatorname{cl}(C^i+iC^1)=z\wedge\bar w,
\end{equation}
et
\begin{equation}\label{eq:hodge-suc-weil-determinante}
 P_m=\det_*(C^i+iC^1),
 \qquad
 \operatorname{cl}(P_m)
 =(-1)^{m(m-1)/2}m!\,w_+.
\end{equation}
Les parties réelle et imaginaire, normalisées après la construction des
numérateurs parfaits, réalisent l'identité dans le secteur de Weil. Pour
\(m=2\), la retenue vaut deux et ne nécessite pas d'inverser trois. Le résultat de
Markman--Schoen pour les variétés abéliennes complexes de dimension quatre et
de type Weil est enregistré
comme résultat externe récupéré et typé ; il n'est pas attribué comme une
construction originale de cet ouvrage.

\noindent\textbf{Verrou de hiérarchie.}
Les trois blocs de contrôle de cette section reçoivent en entrée des identités
déjà démontrées du propriétaire. Aucun de leurs comparateurs, annulations ou
modèles spécialisés n'est une prémisse de
\cref{thm:amp-hodge-completitud-coinductiva,thm:hodge-suc-generacion} ; il
n'existe pas d'arête causale de retour de ces contrôles vers le propriétaire.


\noindent\textbf{Statut du contrôle suivant : CONTRÔLE\_NON\_DIRECTEUR.}
L'ablation reçoit la conclusion Hodge déjà démontrée et réfute seulement la
suppression de l'antécédent basé ; elle n'apporte aucune prémisse au propriétaire et ne
peut rétrograder ses théorèmes.

\begin{ablation}[Sélecteur fini sans histoire]
\label{abl:hodge-suc-eliptico}
Soit \(E\) une courbe elliptique complexe. Il n'existe pas d'ensemble fini de
représentants rationnels de la classe d'un point qui soit stable sous toutes
les translations de \(E\) et qui choisisse en même temps naturellement un
représentant sans donnée basée. Cette obstruction réfute la suppression de la
rigidification ; elle n'affecte pas les théorèmes
\ref{thm:hodge-suc-identidad} et \ref{thm:hodge-suc-generacion}.
\end{ablation}

\begin{proof}
Si \(z\) est de degré non nul, les différences
\(t_x^{*}z-z\), lorsque \(x\in E(\mathbb C)\) varie, parcourent par
l'application d'Abel--Jacobi un sous-groupe de \(\operatorname{Pic}^{0}(E)\)
d'indice fini à isogénie près. Leur enveloppe rationnelle ne tient pas dans l'espace
engendré par une orbite finie. Par conséquent, un repère fini stable sous toutes
les translations ne peut conserver simultanément le représentant et sa
naturalité. Dans \(X_{\chi}(E,1)\), en revanche, la translation modifie
l'histoire et sa base même si la publication cohomologique reste fixe.
L'antécédent \(\xi_{\alpha}\) conserve précisément cette donnée, de sorte que
l'hypothèse abolie par le contrôle n'est pas une prémisse du théorème.
\end{proof}
\chapter{Unité déterminantale couplée, égalité rang--ordre d'annulation et formule du terme principal de Birch--Swinnerton--Dyer}
\chaptermark{Unité déterminantale et formule de Birch--Swinnerton--Dyer}

\noindent La spécialisation du
\cref{thm:nucleo-continuidad-conexion-nonadica-conjunta} est
\[
 (\widehat X_\infty,\Gamma_9)
 \xrightarrow{\ \mathcal R_{\rm BSD}\ }
 P_E^{\rm gen}
 \longrightarrow(z_K,z_{\rm PT})
 \longrightarrow R_{\rm Boc}^{\rm der}
 \longrightarrow\text{rang et terme principal de }L(E,s).
\]
Les deux publications partent de la même unité basée ; la suite de
localisation et la réduction de Smith précèdent la conclusion sur
\(\Sha(E)\).

\subsection*{Objet généalogique et coalgèbre d'Alexander--Whitney}

Soit \(E/\mathbb Q\) une courbe elliptique ; fixons le système de coefficients
\(A_n\), la profondeur \(m\) et les orientations locales et globales du
complexe de Selmer. Si
\(\operatorname{TPK}^{\mathrm{surv}}_{m}\) est la catégorie des histoires
survivantes du Topological Phase Kernel, on définit
\begin{equation}\label{eq:bsd-suc-gmn}
 G_{m,n}:=
 N_{*}\bigl(N\operatorname{TPK}^{\mathrm{surv}}_{m};A_n\bigr).
\end{equation}
Pour un simplexe non dégénéré, la diagonale d'Alexander--Whitney est
\begin{equation}\label{eq:bsd-suc-aw}
 \Delta_{\mathrm{AW}}[x_0\to\cdots\to x_q]
 =\sum_{j=0}^{q}
 [x_0\to\cdots\to x_j]\otimes
 [x_j\to\cdots\to x_q].
\end{equation}
La counité vaut un aux sommets et zéro en degrés positifs. En
particulier,
\begin{equation}\label{eq:bsd-suc-group-like}
 \Delta_{\mathrm{AW}}[x]=[x]\otimes[x],
 \qquad \varepsilon_{\mathrm{AW}}[x]=1.
\end{equation}
La propriété d'élément groupal appartient au sommet d'état complet : elle s'applique
avant d'oublier feuillet, voie, retenue ou mémoire.

Soit \(P_E^{\mathrm{Man}}\to A[0]\) le complexe de Manin basé muni de son
augmentation, et soit \(G_{m,n}\to A[0]\) l'augmentation induite par
\(\varepsilon_{\mathrm{AW}}\). La cible conservative est le produit
fibré
\begin{equation}\label{eq:bsd-suc-pullback}
 P_E^{\mathrm{gen}}
 :=P_E^{\mathrm{Man}}\times_{A[0]}G_{m,n}.
\end{equation}
Si \(s:A[0]\to P_E^{\mathrm{Man}}\) est la section basée, alors
\begin{equation}\label{eq:bsd-suc-section}
 \widehat\Pi_0(c)
 :=\bigl(s\varepsilon_{\mathrm{AW}}(c),c\bigr)
\end{equation}
est une section conservative : sa seconde projection restitue \(c\), tandis que
la première publie sa coordonnée modulaire. L'histoire et l'ombre modulaire
sont couplées dans un seul objet.

\subsection*{La même unité dans les canaux Kato et Poitou--Tate}

Les compatibilités de phase, de feuillet, de corestriction et de relèvement de Hensel
définissent sur \(P_E^{\mathrm{gen}}\) un couple de publications
\begin{equation}\label{eq:bsd-suc-copub}
 P_E=(P_E^{K},P_E^{\mathrm{PT}}).
\end{equation}
Les deux composantes reçoivent la même unité généalogique \(1_E\), de sorte
que
\begin{equation}\label{eq:bsd-suc-zk-zpt}
 P_E^{K}(1_E)=z_K,
 \qquad P_E^{\mathrm{PT}}(1_E)=z_{\mathrm{PT}}.
\end{equation}
Ce ne sont pas deux choix indépendants de générateur.

Soit \(L_{\mathrm{root}}=A_n z_{\mathrm{PT}}\) le sous-complexe racine basé,
normalisé par
\(\lambda_{\mathrm{PT}}(z_{\mathrm{PT}})=1\), et soit
\(p_{\mathrm{root}}\) le projecteur de complexes sur ce sous-complexe. On pose
\(q_{\mathrm{root}}=I-p_{\mathrm{root}}\), qui est également un projecteur de
complexes ; en particulier, \(q_{\mathrm{root}}z_K\) est un cycle. La compatibilité de la
counité dans \eqref{eq:bsd-suc-group-like} avec le couple
\eqref{eq:bsd-suc-copub} donne
\begin{equation}\label{eq:bsd-suc-normalizacion-raiz}
 \lambda_{\mathrm{PT}}(p_{\mathrm{root}}z_K)=1.
\end{equation}

La dualité est formulée au niveau des chaînes. Pour un complexe parfait
\(M\) sur \(\Lambda\), soit
\begin{equation}\label{eq:bsd-suc-d3}
 D_3(M):=
 \operatorname{RHom}_{\Lambda}(M^{\iota},\Lambda)[-3].
\end{equation}
Si \(A\to C\to D_3(A)\) est le bloc d'incidence, on définit
\begin{equation}\label{eq:bsd-suc-orth-red}
 A^{\perp}:=\operatorname{Fib}\bigl(C\to D_3(A)\bigr),
 \qquad
 C^{\mathrm{red}}:=\operatorname{Cofib}\bigl(A\to A^{\perp}\bigr),
\end{equation}
et la forme réduite induit
\begin{equation}\label{eq:bsd-suc-xorth}
 X^{\mathrm{orth}}:C^{\mathrm{red}}
 \xrightarrow{\ \sim\ }D_3(C^{\mathrm{red}}).
\end{equation}
Sur la droite déterminante réduite, on obtient une involution
\(J_{\mathrm{red}}^2=I\) et des bases compatibles
\begin{equation}\label{eq:bsd-suc-jred}
 J_{\mathrm{red}}(\bar b_{\mathrm{PT}})=\bar b_K.
\end{equation}
Les deux feuillets transportent le même scalaire basé :
\begin{equation}\label{eq:bsd-suc-doble-hoja}
 u_{+}=u_{-}=\tau_{\mathrm{bas}}(D_K).
\end{equation}
L'égalité vit dans les deux droites identifiées par
\eqref{eq:bsd-suc-jred} ; elle n'est pas remplacée par une équation entre un scalaire et
son inverse. De même, \(D_3(D_K)=A^{\perp}\) appartient à la droite duale et ne
se confond pas avec une seconde copie non orientée de la droite racine.

\begin{theorem}[Rigidité de la racine commune]
\label{thm:bsd-suc-raiz-filler}
Avec les bases et orientations précédentes,
\begin{equation}\label{eq:bsd-suc-raiz}
 p_{\mathrm{root}}z_K=z_{\mathrm{PT}}.
\end{equation}
\end{theorem}

\begin{proof}
Écrivons
\(p_{\mathrm{root}}z_K=u z_{\mathrm{PT}}\). En appliquant
\(\lambda_{\mathrm{PT}}\) et en utilisant la normalisation de la base avec
\eqref{eq:bsd-suc-normalizacion-raiz}, on obtient
\[
 1=\lambda_{\mathrm{PT}}(p_{\mathrm{root}}z_K)
  =u\lambda_{\mathrm{PT}}(z_{\mathrm{PT}})=u.
\]
Cela prouve \eqref{eq:bsd-suc-raiz} ; l'identification
\eqref{eq:bsd-suc-jred} assure que la conclusion est indépendante du
feuillet à partir duquel elle est publiée.
\end{proof}

\begin{lemma}[Forme normale du complément]
\label{lem:bsd-suc-forma-normal}
Dans le bloc fini--singulier avec
les modules libres sur \(A_n\), soit \(g\) primitif ---c'est-à-dire, élément d'une
base du module d'arrivée--- et supposons que
\(\operatorname{im}\partial\) n'ait pas de composante dans le facteur direct
\(A_n r\). Si la différentielle est donnée par
\begin{equation}\label{eq:bsd-suc-diferencial-normal}
 \partial f=3c e+b,
 \qquad \partial e=g,
 \qquad \partial b=-3c g,
\end{equation}
tout cycle \(z_K=a r+u e+v b\) satisfait
\begin{equation}\label{eq:bsd-suc-ciclo-normal}
 u=3cv,
 \qquad
 z_{\mathrm{PT}}-z_K=(1-a)r-v\,\partial f,
 \quad z_{\mathrm{PT}}=r.
\end{equation}
L'égalité racine \eqref{eq:bsd-suc-raiz} impose \(a=1\) ; par conséquent
\(h=-vf\) comble la différence complète.
\end{lemma}

\begin{proof}
La condition \(\partial z_K=0\) et
\eqref{eq:bsd-suc-diferencial-normal} donnent
\((u-3cv)g=0\). Comme \(g\) fait partie d'une base d'un module libre,
la comparaison de son coefficient donne \(u=3cv\). Substituer cette égalité dans
\(r-z_K\) produit \eqref{eq:bsd-suc-ciclo-normal}. Le coefficient \(a\)
est précisément la coordonnée de \(p_{\mathrm{root}}z_K\) dans la base
\(r=z_{\mathrm{PT}}\) ; le théorème \ref{thm:bsd-suc-raiz-filler} donne
\(a=1\). Alors
\(\partial(-vf)=-v\partial f=z_{\mathrm{PT}}-z_K\).
\end{proof}

La construction propriétaire du remplissage est donc la formule explicite
\begin{equation}\label{eq:bsd-suc-hstar}
 h_*:=-vf,
 \qquad
 \partial h_*=z_{\mathrm{PT}}-z_K.
\end{equation}
Aucune homotopie abstraite n'est introduite comme prémisse : la fermeture fixe
\(u=3cv\), l'unité racine fixe \(a=1\) et le générateur \(f\) comble
directement la différence complète.

\paragraph{Contrôle homotopique équivalent non directeur.}
\noindent\textbf{Statut : CONTRÔLE\_NON\_DIRECTEUR.}
Comme vérification ultérieure, si le même complément est regroupé au moyen
d'un lecteur \(H_{\mathrm{FS}}\), l'identité
\begin{equation}\label{eq:bsd-suc-hfs}
 \partial H_{\mathrm{FS}}+H_{\mathrm{FS}}\partial
 =q_{\mathrm{root}}
\end{equation}
reproduit \eqref{eq:bsd-suc-hstar} lorsqu'elle est appliquée à
\(q_{\mathrm{root}}z_K\). Ce lecteur est une notation auxiliaire de la
contraction déjà réalisée par \(h_*=-vf\) ; ce n'est pas une obligation préalable et il ne
peut rouvrir la construction explicite.

\subsection*{Régulateur de Bockstein dérivé d'ordre arbitraire}

Soit \(K\) un corps de caractéristique zéro,
\(\Lambda=K[[T]]\), et soit
\begin{equation}\label{eq:bsd-suc-S}
 S(T):V_{\Lambda}\longrightarrow W_{\Lambda},
 \qquad
 V_{\Lambda}\simeq W_{\Lambda}\simeq\Lambda^{m},
\end{equation}
une différentielle basée qui devient inversible sur \(K((T))\). On fixe
la droite
\begin{equation}\label{eq:bsd-suc-lambda0}
 \lambda_0(S):=(\det_K V)^{-1}\otimes_K\det_K W,
 \quad
 V=V_{\Lambda}/TV_{\Lambda},\quad
 W=W_{\Lambda}/TW_{\Lambda}.
\end{equation}
Une forme de Smith basée a pour expression
\begin{equation}\label{eq:bsd-suc-smith}
 U(T)S(T)V(T)
 =\operatorname{diag}
 \bigl(T^{a_1},\ldots,T^{a_s},1,\ldots,1\bigr),
 \qquad a_i\geq1,
\end{equation}
et
\begin{equation}\label{eq:bsd-suc-orden}
 d=\operatorname{ord}_{T}\det S(T)=\sum_{i=1}^{s}a_i.
\end{equation}

La filtration \(T\)-adique du complexe
\([V_{\Lambda}\xrightarrow{S}W_{\Lambda}]\) produit des pages
\(V_q=E_q^0\), \(W_q=E_q^1\) et des Bocksteins
\(\beta_q:V_q\to W_q\). Intrinsèquement,
\begin{equation}\label{eq:bsd-suc-pages}
 A_q:=V_q/V_{q+1},
 \qquad
 B_q:=\ker(W_q\twoheadrightarrow W_{q+1})
      =\operatorname{im}\beta_q,
\end{equation}
et \(\beta_q\) induit
\begin{equation}\label{eq:bsd-suc-beta-q}
 \bar\beta_q:A_q\xrightarrow{\ \sim\ }B_q,
 \qquad
 \tau_q:=\det(\bar\beta_q).
\end{equation}
Les pages sont finies et vérifient
\begin{equation}\label{eq:bsd-suc-carry-orden}
 d=\sum_{q\geq1}q\,\dim A_q
  =\sum_{q\geq1}\dim V_q.
\end{equation}
La première page voit
\(r_0=\dim V_1\) ; la retenue d'ordre supérieur est
\begin{equation}\label{eq:bsd-suc-carry-superior}
 c_I^{\deg}=d-r_0=\sum_{q\geq2}\dim V_q.
\end{equation}

La page régulière est l'isomorphisme induit
\begin{equation}\label{eq:bsd-suc-pagina-regular}
 \bar S(0):A_0:=V/V_1
 \xrightarrow{\ \sim\ }
 B_0:=\operatorname{im}S(0),
 \qquad
 \tau_{\mathrm{reg}}:=\det\bar S(0).
\end{equation}
Les suites exactes de \eqref{eq:bsd-suc-pages}, avec leurs signes de
Koszul, recollent ces trivialisations en un unique élément basé
\begin{equation}\label{eq:bsd-suc-rboc}
 R_{\mathrm{Boc}}^{\mathrm{der}}(S)
 :=\tau_{\mathrm{reg}}\otimes
   \bigotimes_{q\geq1}\tau_q
 \in\lambda_0(S),
\end{equation}
où le symbole tensoriel abrège les isomorphismes de recollement de
Knudsen--Mumford.

\begin{theorem}[Identité déterminantale du régulateur dérivé]
\label{thm:bsd-suc-bockstein}
Dans la droite basée \(\lambda_0(S)\), on a
\begin{equation}\label{eq:bsd-suc-lt}
 R_{\mathrm{Boc}}^{\mathrm{der}}(S)
 =\operatorname{LT}_{T}(S)
 :=\left[T^{-d}\det S(T)\right]_{T=0}.
\end{equation}
L'égalité conserve la page régulière, toutes les pages positives et la
retenue \eqref{eq:bsd-suc-carry-superior}.
\end{theorem}

\begin{proof}
Dans les bases de Smith de \eqref{eq:bsd-suc-smith}, la page régulière et chaque
page positive portent le générateur canonique de leur droite. Les signes de
Koszul du recollement sont exactement ceux nécessaires pour restituer l'ordre des
bases totales, de sorte que leur produit est le générateur positif de
\(\lambda_0(S)\). D'autre part,
\[
 \det S(T)=\det U(T)^{-1}\det V(T)^{-1}T^d,
\]
et le terme initial est
\(\det U(0)^{-1}\det V(0)^{-1}\). En annulant le changement de bases, tant
\eqref{eq:bsd-suc-rboc} que le terme initial sont multipliés par ce même
facteur. Cela prouve \eqref{eq:bsd-suc-lt} sans identifier les scalaires avant de
fixer la droite et sa base.
\end{proof}

Pour chaque \(q\geq1\), la dualité de Poitou--Tate construite sur la même
page fournit
\begin{equation}\label{eq:bsd-suc-jq}
 j_q:B_q\xrightarrow{\ \sim\ }
 A_q^{\vee}\otimes(I^q/I^{q+1}),
\end{equation}
et, après orientation de \(I^q/I^{q+1}\), l'accouplement
\begin{equation}\label{eq:bsd-suc-altura-derivada}
 h_E^{(q)}(x,y)
 :=\bigl(j_q\bar\beta_q(x)\bigr)(y).
\end{equation}

% Ampliación sustantiva de los comparadores determinantes BSD.
% Módulo destinado al capítulo BSD de la Parte IV.

\section{Comparateurs de l'unité déterminantale commune}
\label{sec:amp-bsd-comparadores}

L'égalité racine \eqref{eq:bsd-suc-raiz} fixe une seule unité dans les canaux
Kato et Poitou--Tate. À partir de cette unité, on construit maintenant les
comparateurs qui transportent le régulateur dérivé jusqu'à sa publication
scalaire. Aucun comparateur ne choisit une normalisation après avoir connu le
terme principal : tous proviennent de morphismes de complexes, de suites
exactes et de bases déjà orientées.

\subsection{La chaîne de droites déterminantes}
\label{subsec:amp-bsd-lineas}

Pour un complexe parfait basé \(C\), on utilise la convention
\begin{equation}\label{eq:amp-bsd-linea-determinante}
 \lambda(C):=\bigotimes_i
 \det H^i(C)^{(-1)^i}.
\end{equation}
Soient \(\mathcal C_E^K\), \(\mathcal C_E^{\mathrm{Iw}}\),
\(\mathcal C_E^{\mathrm{PT}}\) et \(\mathcal C_E^{\mathrm{loc}}\),
respectivement, le complexe publié par le canal de Kato, son modèle
d'Iwasawa basé, le complexe réduit autodual de Poitou--Tate et le cône de
localisation finie--singulière. Leurs droites sont reliées par
\begin{equation}\label{eq:amp-bsd-cadena-lineas}
 \mathcal L_E^K
 \xrightarrow{\ \mathfrak c_{K/\mathrm{FERS}}\ }
 \mathcal L_E^{\mathrm{Iw}}
 \xrightarrow{\ \mathfrak c_{\mathrm{PT}}\ }
 \mathcal L_E^{\mathrm{ht}}
 \xrightarrow{\ \mathfrak c_{\mathrm{STS}}\ }
 \mathcal L_E^{\mathrm{fin}}
 \xrightarrow{\ \mathfrak c_{\mathrm{tors}}\ }
 \mathcal L_E^{\mathrm{ar}}
 \xrightarrow{\ \mathfrak c_{\infty}\ }
 \mathcal L_E.
\end{equation}
Le sigle \(\mathrm{STS}\) désigne dans cette formule l'étape conjointe de
Smith, Tamagawa et Tate--Shafarevich. Toutes les flèches de
\eqref{eq:amp-bsd-cadena-lineas} sont des isomorphismes de droites orientées.

\subsection{Comparaison de chaînes Kato--FERS à partir de l'unité commune}
\label{subsec:amp-bsd-comparacion-cadenas-kato-fers}

L'équivalence à l'origine du premier comparateur peut être construite avant de
prendre les déterminants. Soit \(F^q\mathcal C\) la filtration par profondeur
nonadique, identifiée à la filtration \(T\)-adique au moyen de la retenue, et
écrivons
\begin{equation}\label{eq:amp-bsd-complejos-kato-iwasawa}
 \mathcal C_E^K
 :=\operatorname{Tot}\!\left(
 P_E^K\widehat\Pi_0G_{m,n}\right),
 \qquad
 \mathcal C_E^{\mathrm{Iw}}
 :=[V_\Lambda\xrightarrow{S_E(T)}W_\Lambda].
\end{equation}
Soit
\(\operatorname{car}_T:G_{m,n}\to\mathcal C_E^{\mathrm{Iw}}\) le lecteur
de retenue qui associe à une route \(\gamma\) le générateur déterminé par ses
faces initiale et finale, multiplié par
\(T^{\operatorname{Carr}(\gamma)}\). La multiplication des concaténations dans
le modèle d'Iwasawa est notée \(\mu_{\mathrm{Iw}}\). La diagonale
d'Alexander--Whitney et la publication \(P_E^K\) définissent alors, avant de
prendre la cohomologie ou les déterminants, l'application de chaînes
\begin{equation}\label{eq:amp-bsd-phi-kato-fers}
\begin{aligned}
 \Phi_{K/\mathrm{FERS}}&:
 \mathcal C_E^K\longrightarrow\mathcal C_E^{\mathrm{Iw}},\\
 \Phi_{K/\mathrm{FERS}}
 &:=
 \mu_{\mathrm{Iw}}\circ
 \bigl(P_E^K\widehat\otimes\operatorname{car}_T\bigr)
 \circ\Delta_{\mathrm{AW}},\\
 \Phi_{K/\mathrm{FERS}}(z_K)&=z_{\mathrm{Iw}}.
\end{aligned}
\end{equation}
où \(z_{\mathrm{Iw}}\) est le générateur racine induit par \(1_E\). Ainsi,
chaque simplexe normalisé est envoyé sur la somme de ses coupes
d'Alexander--Whitney, en conservant dans chaque terme les deux faces et la puissance
de retenue. La formule ne consulte pas le terme principal analytique.

\begin{lemma}[Critère générateur par générateur pour l'équivalence]
\label{lem:amp-bsd-bases-emparejadas-kato-fers}
À chaque degré \(i\), soit \(\mathcal B_i^K\) la base finie de simplexes
normalisés, ordonnée par profondeur, faces et retenue, et soit
\(\mathcal B_i^{\mathrm{Iw}}\) la base formée par les générateurs ayant les
mêmes faces initiale et finale dans le modèle d'Iwasawa. Pour une application
candidate
\[
 \upsilon_i:\mathcal B_i^K\longrightarrow
 \mathcal B_i^{\mathrm{Iw}},
\]
définissons le résiduel de degré zéro
\begin{equation}\label{eq:amp-bsd-residual-generador-graduado}
 \varepsilon_i(b)
 :=\Phi_{K/\mathrm{FERS}}^i(b)\bmod T-\upsilon_i(b).
\end{equation}
Si \(\upsilon_i\) est une bijection et
\(\varepsilon_i(b)=0\) pour tout \(b\in\mathcal B_i^K\), alors les deux
modules complétés de degré \(i\) sont libres de même rang
\(d_i=|\mathcal B_i^K|\) et, pour chaque générateur,
\begin{equation}\label{eq:amp-bsd-phi-en-cada-generador}
 \Phi_{K/\mathrm{FERS}}^i(b)
 =\upsilon_i(b)
  +T\sum_{c\in\mathcal B_i^{\mathrm{Iw}}}
    \lambda_{c,b}(T)c,
 \qquad \lambda_{c,b}(T)\in\Lambda.
\end{equation}
En particulier,
\(\operatorname{gr}^{F}\Phi_{K/\mathrm{FERS}}^i=I\) sur tous les
générateurs, et non seulement sur l'unité racine.
\end{lemma}

\begin{proof}
L'annulation de \eqref{eq:amp-bsd-residual-generador-graduado} signifie que la
colonne de \(b\) a pour terme constant \(\upsilon_i(b)\). Les autres
coefficients appartiennent à \(T\Lambda\), ce qui donne
\eqref{eq:amp-bsd-phi-en-cada-generador}. La bijectivité identifie les
bases et fixe le même rang. Dans ces bases, la réduction modulo \(T\) de la
matrice de \(\Phi^i\) est l'identité ; d'où
\(\operatorname{gr}^{F}\Phi^i=I\).
\end{proof}

L'égalité racine
\(\Phi_{K/\mathrm{FERS}}(z_K)=z_{\mathrm{Iw}}\) fixe la normalisation commune.
L'écriture de \(P_E^K\) sur chaque simplexe normalisé, la bijection
\(\upsilon_i\) et les résiduels \(\varepsilon_i(b)\) fournissent un
contrôle générateur par générateur de l'équivalence déjà construite ; ils ne choisissent ni ne
régissent la conclusion BSD.

\begin{lemma}[Équivalence filtrée et normalisation de la racine]
\label{lem:amp-bsd-equivalencia-kato-fers}
Pour l'application de chaînes construite dans
\eqref{eq:amp-bsd-phi-kato-fers}, l'application \(\upsilon_i\) est la
correspondance de bases induite par les faces et la retenue de chaque
simplexe. La formule d'Alexander--Whitney vérifie à chaque degré le
critère du lemme \ref{lem:amp-bsd-bases-emparejadas-kato-fers} ; il s'agit d'un
contrôle de l'équivalence déjà construite, et non d'une prémisse supplémentaire.
L'application \(\Phi_{K/\mathrm{FERS}}\) est une équivalence basée de
complexes parfaits. Dans les bases ordonnées par profondeur, ses matrices à
chaque degré sont de la forme
\begin{equation}\label{eq:amp-bsd-matriz-kato-fers}
 [\Phi_{K/\mathrm{FERS}}^i](T)=I+TB_i(T),
 \qquad B_i(T)\in M_{d_i}(\Lambda).
\end{equation}
Par conséquent, l'unité
\begin{equation}\label{eq:amp-bsd-rho-determinantal}
 \rho_{E,\ell}(T)
 :=\prod_i\det\!\left(I+TB_i(T)\right)^{(-1)^i}
 \in\Lambda^\times
\end{equation}
satisface \(\rho_{E,\ell}(0)=1\).
\end{lemma}

\begin{proof}
L'identité simpliciale
\(\partial\Delta_{\mathrm{AW}}
=(\partial\otimes I+(-1)^{\deg}I\otimes\partial)
\Delta_{\mathrm{AW}}\), avec la compatibilité de
\(P_E^K\) avec les faces et les dégénérescences, démontre que
\eqref{eq:amp-bsd-phi-kato-fers} commute avec les différentielles. Dans le gradué
associé, le lemme
\ref{lem:amp-bsd-bases-emparejadas-kato-fers} identifie tous les
générateurs et donne l'identité. Dans ces bases,
\eqref{eq:amp-bsd-phi-en-cada-generador} est précisément
\eqref{eq:amp-bsd-matriz-kato-fers}.

La complétude de \(\Lambda=K[[T]]\) fait converger la série
\begin{equation}\label{eq:amp-bsd-inversa-kato-fers}
 (I+TB_i(T))^{-1}
 =\sum_{n\geq0}(-TB_i(T))^n;
\end{equation}
ces inverses commutent également avec les différentielles, et forment donc
l'inverse basé de \(\Phi_{K/\mathrm{FERS}}\). En prenant le déterminant
alterné, on obtient \eqref{eq:amp-bsd-rho-determinantal} ; l'évaluation en
\(T=0\) donne un. Ainsi, la normalisation n'est pas postulée à partir du terme
principal : c'est la coordonnée constante d'une équivalence qui est l'identité
sur l'unité racine.
\end{proof}

Le premier comparateur est le déterminant de
\(\Phi_{K/\mathrm{FERS}}\). Si \(\mathfrak z_K\) désigne l'élément
provenant de l'unité \(1_E\), son image est caractérisée par
\begin{equation}\label{eq:amp-bsd-comparador-kato}
 \mathfrak c_{K/\mathrm{FERS}}(\mathfrak z_K)
 =\left[T^{-d}\rho_{E,\ell}(T)\det S_E(T)\right]_{T=0},
 \qquad
 \rho_{E,\ell}(0)=1.
\end{equation}
L'unité de \(\rho_{E,\ell}\) est orientée et vaut un à la racine ; elle ne
modifie donc ni l'ordre ni la base. Le théorème
\ref{thm:bsd-suc-bockstein} transforme \eqref{eq:amp-bsd-comparador-kato} en
\begin{equation}\label{eq:amp-bsd-kato-bockstein}
 \mathfrak c_{K/\mathrm{FERS}}(\mathfrak z_K)
 =R_{\mathrm{Boc}}^{\mathrm{der}}(S_E)
 =\tau_{\mathrm{reg}}\otimes\bigotimes_{q\geq1}\tau_q.
\end{equation}

Le comparateur de Poitou--Tate se construit page par page. Pour
\(q\geq1\), l'isomorphisme \(j_q\) de
\eqref{eq:bsd-suc-jq} et le Bockstein réduit
\(\bar\beta_q\) induisent
\begin{equation}\label{eq:amp-bsd-comparador-pt-q}
 \mathfrak c_{\mathrm{PT},q}(\tau_q)
 :=\det\!\left(j_q\circ\bar\beta_q\right)
 =\det h_E^{(q)}.
\end{equation}
La direction régulière est transportée au moyen de
\(\det\bar S_E(0)=\tau_{\mathrm{reg}}\). Le recollement de Knudsen--Mumford et les
signes de Koszul définis dans \eqref{eq:bsd-suc-rboc} fournissent l'unique
isomorphisme orienté
\begin{equation}\label{eq:amp-bsd-comparador-pt}
 \mathfrak c_{\mathrm{PT}}
 \left(\tau_{\mathrm{reg}}\otimes\bigotimes_{q\geq1}\tau_q\right)
 =\tau_{\mathrm{reg}}^{\mathrm{ht}}
  \otimes\bigotimes_{q\geq1}\det h_E^{(q)}.
\end{equation}
Sur la première page stable, la forme \(h_E^{(1)}\) est la forme de
Poincaré--Néron--Tate sur le réseau libre de Mordell--Weil ; son déterminant est
\begin{equation}\label{eq:amp-bsd-regulador-nt}
 \det h_E^{(1)}=\operatorname{Reg}(E)
\end{equation}
dans la base induite par l'unité racine.

\begin{lemma}[Déploiement de Poitou--Tate du degré arithmétique]
\label{lem:amp-bsd-grado-pt}
Soit
\[
 \operatorname{Flag}_{\mathrm{Boc}}(E)
 :=\bigoplus_{q\geq1}V_q
\]
le drapeau fini des pages positives, chacune avec la base héritée de
l'unité racine. La réduction orthogonale et les isomorphismes \(j_q\) construisent un
isomorphisme basé
\begin{equation}\label{eq:amp-bsd-psi-pt}
 \Psi_{\mathrm{PT}}:
 \bigl(E(\mathbb Q)/E(\mathbb Q)_{\mathrm{tors}}\bigr)\otimes K
 \xrightarrow{\ \sim\ }\operatorname{Flag}_{\mathrm{Boc}}(E).
\end{equation}
En particulier,
\begin{equation}\label{eq:amp-bsd-rango-bandera}
 \operatorname{rank}E(\mathbb Q)
 =\sum_{q\geq1}\dim_KV_q
 =\operatorname{ord}_T\det S_E(T).
\end{equation}
\end{lemma}

\begin{proof}
On filtre le complexe réduit par les puissances de l'idéal d'augmentation. Dans le
quotient de niveau \(q\), la direction régulière a déjà été séparée par
\(\bar S_E(0)\), et le seul bloc survivant est \(V_q\). L'équivalence
\(X^{\mathrm{orth}}\) identifie le bloc opposé à son dual, tandis que
\(j_q\bar\beta_q\) fournit l'accouplement parfait qui relève la
base du quotient au niveau précédent. En partant de la dernière page non nulle et
en répétant ce relèvement, on obtient \(\Psi_{\mathrm{PT}}\).

Dans les bases ordonnées par profondeur, la matrice de l'application résultante est
triangulaire par blocs et tous ses blocs diagonaux sont des identités : dans le
premier bloc parce que \(z_K\) et \(z_{\mathrm{PT}}\) ont la même racine, et
dans les autres parce que \(j_q\) et \(\bar\beta_q\) sont des isomorphismes basés.
Ainsi \(\Psi_{\mathrm{PT}}\) est inversible et ne perd aucune page.
En prenant les dimensions et en utilisant
\eqref{eq:bsd-suc-carry-orden}, on obtient
\eqref{eq:amp-bsd-rango-bandera}. En prenant les déterminants, le même
relèvement est exactement le recollement de
\eqref{eq:amp-bsd-comparador-pt} ; ainsi, l'égalité des rangs et l'égalité des
droites proviennent d'une seule application basée.
\end{proof}

\subsection{Triangle de localisation, rang de Smith et facteurs finis}
\label{subsec:amp-bsd-smith-sha}

Pour chaque place \(v\), la condition locale finie est un sous-complexe basé
\(\mathcal C_{E,v}^{\mathrm{fin}}\subset\mathcal C_{E,v}\), et l'on définit
\(\mathcal C_{E,v}^{\mathrm{sing}}\) par le triangle
\begin{equation}\label{eq:amp-bsd-triangulo-local}
 \mathcal C_{E,v}^{\mathrm{fin}}
 \xrightarrow{\ i_v\ }\mathcal C_{E,v}
 \longrightarrow\mathcal C_{E,v}^{\mathrm{sing}}
 \longrightarrow\mathcal C_{E,v}^{\mathrm{fin}}[1].
\end{equation}
Si \(\operatorname{res}_v\) désigne la restriction globale, l'application de
localisation est, au niveau des chaînes,
\begin{equation}\label{eq:amp-bsd-mapa-localizacion}
 \ell_E:
 \mathcal C_E^{\mathrm{glob,red}}\oplus
 \bigoplus_v\mathcal C_{E,v}^{\mathrm{fin}}
 \longrightarrow\bigoplus_v\mathcal C_{E,v},
 \qquad
 \ell_E\bigl(x,(y_v)_v\bigr)
 :=\bigl(\operatorname{res}_v x-i_vy_v\bigr)_v .
\end{equation}
La définition par cône du complexe de Selmer réduit donne le triangle
\begin{equation}\label{eq:amp-bsd-triangulo-localizacion}
 \mathcal C_E^{\mathrm{red}}
 \longrightarrow
 \mathcal C_E^{\mathrm{glob,red}}\oplus
 \bigoplus_v\mathcal C_{E,v}^{\mathrm{fin}}
 \xrightarrow{\ \ell_E\ }
 \bigoplus_v\mathcal C_{E,v}
 \longrightarrow\mathcal C_E^{\mathrm{red}}[1].
\end{equation}
L'exactitude locale--globale provient donc du cône d'une application écrite dans
\eqref{eq:amp-bsd-mapa-localizacion} ; elle n'est pas ajoutée ensuite comme une relation
entre cardinaux.

On annule dans \eqref{eq:amp-bsd-triangulo-localizacion} le facteur racine
commun, qui est le même dans les deux canaux par
\eqref{eq:bsd-suc-raiz}. Le complément résultant est noté
\begin{equation}\label{eq:amp-bsd-cono-local-reducido}
 \mathcal C_E^{\mathrm{loc,red}}
 :=q_{\mathrm{root}}\operatorname{Cone}(\ell_E)[-1].
\end{equation}
Après avoir séparé le réseau libre de Mordell--Weil et son dual au moyen de
\(X^{\mathrm{orth}}\), un choix des bases entières déjà orientées
représente ce complexe par
\begin{equation}\label{eq:amp-bsd-complejo-smith}
 \mathcal C_E^{\mathrm{loc,red}}
 \simeq[M_E\xrightarrow{\ D_E\ }N_E].
\end{equation}

\begin{lemma}[Acyclicité rationnelle et rang plein]
\label{lem:amp-bsd-aciclicidad-rango}
Le complexe de \eqref{eq:amp-bsd-complejo-smith} est acyclique après
tensorisation par \(\mathbb Q\). En particulier,
\[
 D_E\otimes\mathbb Q:M_E\otimes\mathbb Q
 \xrightarrow{\ \sim\ }N_E\otimes\mathbb Q
\]
est un isomorphisme ; \(M_E\) et \(N_E\) ont le même rang et \(D_E\) est
de rang plein.
\end{lemma}

\begin{proof}
Soit \(x\) un cycle du complexe
\(\mathcal C_E^{\mathrm{loc,red}}\otimes\mathbb Q\). La réduction a
éliminé le facteur \(p_{\mathrm{root}}\), donc
\(q_{\mathrm{root}}x=x\). Dans les bases entières orientées employées dans
\eqref{eq:amp-bsd-complejo-smith}, le complément fini--singulier se
décompose en les blocs normaux de
\cref{lem:bsd-suc-forma-normal}. Dans chaque bloc \(\lambda\), le coefficient
racine est déjà nul et le cycle s'écrit
\[
 x_{\lambda}=u_{\lambda}e_{\lambda}
                  +v_{\lambda}b_{\lambda},
 \qquad
 \partial x_{\lambda}
   =(u_{\lambda}-3c_{\lambda}v_{\lambda})g_{\lambda}=0.
\]
Comme \(g_{\lambda}\) appartient à la base libre, on a
\(u_{\lambda}=3c_{\lambda}v_{\lambda}\), et la formule explicite de la
différentielle donne
\[
 x_{\lambda}
   =v_{\lambda}(3c_{\lambda}e_{\lambda}+b_{\lambda})
   =\partial(v_{\lambda}f_{\lambda}).
\]
La somme est finie à chaque degré, de sorte que
\(x=\partial\bigl(\sum_{\lambda}v_{\lambda}f_{\lambda}\bigr)\).
Tout cycle est donc un bord par la même forme normale qui a produit
le remplisseur \eqref{eq:bsd-suc-hstar}, sans introduire
\(H_{\mathrm{FS}}\) comme prémisse. L'égalité
\(p_{\mathrm{root}}z_K=z_{\mathrm{PT}}\) est ce qui permet d'annuler la
direction racine avant ce calcul ; sans elle, il resterait une classe de degré
zéro. L'équivalence autoduale
\(X^{\mathrm{orth}}:\mathcal C_E^{\mathrm{red}}\simeq
D_3(\mathcal C_E^{\mathrm{red}})\) apparie les deux degrés restants et
transporte leurs bases avec une orientation opposée complémentaire. On en
obtient \(\operatorname{rank}M_E=\operatorname{rank}N_E\).
Enfin, l'acyclicité d'un complexe à deux termes de même rang
équivaut à l'inversibilité de \(D_E\otimes\mathbb Q\).
\end{proof}

Le lemme permet d'appliquer la forme normale de Smith sans présupposer aucun de
ses diviseurs. Il existe des changements de base orientés \(U,V\) tels que
\begin{equation}\label{eq:amp-bsd-morfismo-smith}
 U D_E V=\operatorname{diag}(s_1,\ldots,s_t),
 \qquad s_i\in\mathbb Z\setminus\{0\}.
\end{equation}
Par conséquent,
\begin{equation}\label{eq:amp-bsd-grupo-smith}
 F_E:=\operatorname{coker}D_E
 \simeq\bigoplus_{i=1}^{t}\mathbb Z/|s_i|\mathbb Z,
 \qquad
 |F_E|=\prod_{i=1}^{t}|s_i|<\infty.
\end{equation}

\begin{lemma}[Suite finie de localisation]
\label{lem:amp-bsd-exactitud-sha-tamagawa}
Le fragment de torsion de la suite longue de cohomologie de
\eqref{eq:amp-bsd-triangulo-localizacion} est
\begin{equation}\label{eq:amp-bsd-exacta-sha-tamagawa}
 0\longrightarrow\operatorname{Sha}(E)
 \xrightarrow{\ \iota_{\mathrm{Sha}}\ }F_E
 \xrightarrow{\ \partial_{\mathrm{loc}}\ }
 \bigoplus_v\Phi_v(\mathbb F_v)
 \longrightarrow0,
 \qquad c_v=|\Phi_v(\mathbb F_v)|.
\end{equation}
\end{lemma}

\begin{proof}
Un cocycle global représente une classe de
\(\operatorname{Sha}(E)\) précisément lorsque sa restriction à chaque place
admet une primitive dans
\(\mathcal C_{E,v}^{\mathrm{fin}}\). La classe de la paire formée par ce
cocycle et ses primitives locales dans le cône de
\eqref{eq:amp-bsd-mapa-localizacion} définit
\(\iota_{\mathrm{Sha}}\). Si cette classe est un bord, la composante globale est un
bord et la classe de Tate--Shafarevich est nulle ; ainsi,
\(\iota_{\mathrm{Sha}}\) est injective.

L'application \(\partial_{\mathrm{loc}}\) prend une classe du cône, la restreint à
chaque quotient
\(\mathcal C_{E,v}^{\mathrm{sing}}\) de
\eqref{eq:amp-bsd-triangulo-local} et conserve sa classe de composante.
L'identification
\[
 H^0(\mathcal C_{E,v}^{\mathrm{sing}})_{\mathrm{tors}}
 \simeq\Phi_v(\mathbb F_v)
\]
est induite par le quotient fini--singulier, non par l'ordre du groupe.
Une classe du cône a un résidu nul exactement lorsque ses représentants
locaux peuvent être choisis dans les sous-complexes finis ; d'après la définition
précédente, ce noyau est l'image de \(\iota_{\mathrm{Sha}}\).
Enfin, la dernière flèche de
\eqref{eq:amp-bsd-triangulo-local} relève chaque classe de composante au
complexe local. L'exactitude du cône
\eqref{eq:amp-bsd-triangulo-localizacion}, après retrait des deux
facteurs libres appariés par \(X^{\mathrm{orth}}\), transforme ce
relèvement en une préimage par \(\partial_{\mathrm{loc}}\). Cela démontre
la surjectivité et, par conséquent, toute
\eqref{eq:amp-bsd-exacta-sha-tamagawa}.
\end{proof}

Maintenant, et seulement après le lemme
\ref{lem:amp-bsd-aciclicidad-rango}, le groupe \(F_E\) est fini.
L'injection de \eqref{eq:amp-bsd-exacta-sha-tamagawa} démontre alors la
finitude de \(\operatorname{Sha}(E)\). En prenant les ordres dans cette suite
exacte, on obtient
\begin{equation}\label{eq:amp-bsd-cardinal-smith}
 |F_E|=|\operatorname{Sha}(E)|\prod_v c_v.
\end{equation}
C'est le comparateur
\(\mathfrak c_{\mathrm{STS}}\) : il transporte les facteurs élémentaires de Smith
vers les facteurs de Tate--Shafarevich et de Tamagawa en conservant leur ordre dans la droite,
au lieu d'insérer leurs cardinaux une fois le terme principal calculé.

La partie de torsion se construit en appliquant le déterminant au réseau de
Mordell--Weil et à son dual de Pontryagin. Les deux facteurs finis sont duaux et
apportent
\begin{equation}\label{eq:amp-bsd-comparador-torsion}
 \mathfrak c_{\mathrm{tors}}:quad
 u\longmapsto
 \frac{u}{|E(\mathbb Q)_{\mathrm{tors}}|^2}.
\end{equation}
Le comparateur archimédien s'obtient en intégrant la différentielle de Néron
orientée sur le cycle réel basé :
\begin{equation}\label{eq:amp-bsd-comparador-periodo}
 \Omega_E:=\int_{E(\mathbb R)}|\omega_E|,
 \qquad
 \mathfrak c_\infty:quad u\longmapsto\frac{u}{\Omega_E}.
\end{equation}
La différentielle et le cycle sont fixés avant d'évaluer \(L(E,s)\) ; ainsi
\eqref{eq:amp-bsd-comparador-periodo} ne normalise pas rétrospectivement la
section analytique.

\subsection{Séparation causale des primitives, des identités et des conséquences}
\label{subsec:amp-bsd-separacion-causal}

Le critère générateur par générateur du
lemme \ref{lem:amp-bsd-bases-emparejadas-kato-fers} audite à chaque degré
l'équivalence déjà construite. Il ne conditionne pas son existence. L'ordre logique
du paquet est
\begin{equation}\label{eq:amp-bsd-diagrama-causal}
\begin{array}{c}
 \begin{gathered}
  1_E,\ \Delta_{\mathrm{AW}},\ P_E^K,\ P_E^{\mathrm{PT}},
  \ h_*=-vf,\ X^{\mathrm{orth}},\\
  \{\,\mathcal C_{E,v}^{\mathrm{fin}}\to\mathcal C_{E,v}\,\}_v,\\
  \text{bases de Mordell--Weil et de Pontryagin ;}\\
  \text{différentielle et cycle de Néron}
 \end{gathered}
 \\[2mm]\Downarrow\\[-1mm]
 \begin{gathered}
  \Phi_{K/\mathrm{FERS}}\text{ équivalence},\quad
  \rho_{E,\ell}(0)=1,\quad
  p_{\mathrm{root}}z_K=z_{\mathrm{PT}},\\
  R_{\mathrm{Boc}}^{\mathrm{der}}=\operatorname{LT}_T(S_E),\quad
  \det(j_q\bar\beta_q)=\det h_E^{(q)},\\
  H^*(\mathcal C_E^{\mathrm{loc,red}}\otimes\mathbb Q)=0,\quad
  D_E\otimes\mathbb Q\text{ isomorphisme},\\
  0\to\operatorname{Sha}(E)\to F_E
  \to\bigoplus_v\Phi_v(\mathbb F_v)\to0
 \end{gathered}
 \\[2mm]\Downarrow\\[-1mm]
 \begin{gathered}
  d_{\mathrm{an}}=\operatorname{rank}E(\mathbb Q),\qquad
  \operatorname{Reg}(E),\qquad
  |\operatorname{Sha}(E)|\prod_vc_v,\\
  |E(\mathbb Q)_{\mathrm{tors}}|^{-2},\qquad
  \Omega_E^{-1},\qquad
  \text{égalité du terme principal}.
 \end{gathered}
\end{array}
\end{equation}
La première ligne contient uniquement des objets et des applications déjà basés, avec
la vérification générateur par générateur indiquée ; la deuxième, des identités
démontrées par ces applications ; la troisième, leurs conséquences sur la droite
déterminante. En particulier, aucun facteur de la dernière ligne n'est
employé pour normaliser rétroactivement une flèche de la première.

\begin{proposition}[Paquet exact de comparateurs]
\label{prop:amp-bsd-paquete-comparadores-exacto}
Les comparateurs de
\eqref{eq:amp-bsd-cadena-lineas} proviennent d'un unique paquet d'applications de
complexes et satisfont, dans l'ordre de construction,
\begin{equation}\label{eq:amp-bsd-paquete-identidades}
 \begin{gathered}
 \Phi_{K/\mathrm{FERS}}(z_K)=z_{\mathrm{Iw}},\qquad
 \operatorname{gr}^{F}\Phi_{K/\mathrm{FERS}}=I,\\
 \det\nolimits_{\mathrm{alt}}\Phi_{K/\mathrm{FERS}}
   =\rho_{E,\ell}(T),\qquad \rho_{E,\ell}(0)=1,\\
 \det(j_q\bar\beta_q)=\det h_E^{(q)}\quad(q\geq1),\\
 \Psi_{\mathrm{PT}}:
  \bigl(E(\mathbb Q)/E(\mathbb Q)_{\mathrm{tors}}\bigr)\otimes K
  \xrightarrow{\sim}\operatorname{Flag}_{\mathrm{Boc}}(E),\\
 p_{\mathrm{root}}z_K=z_{\mathrm{PT}},\qquad
 h_*=-vf,\quad \partial h_*=z_{\mathrm{PT}}-z_K,\\
 H^*(\mathcal C_E^{\mathrm{loc,red}}\otimes\mathbb Q)=0,\qquad
 U D_E V=\operatorname{diag}(s_1,\ldots,s_t),\\
 0\longrightarrow\operatorname{Sha}(E)
  \longrightarrow F_E
  \longrightarrow\displaystyle\bigoplus_v\Phi_v(\mathbb F_v)
  \longrightarrow0 .
 \end{gathered}
\end{equation}
La composition induite sur les droites déterminantes est exactement
\begin{equation}\label{eq:amp-bsd-paquete-composicion}
 \det(\Phi_{K/\mathrm{FERS}})
 \xrightarrow{\ \det(j_q\bar\beta_q)\ }
 \det(D_E)
 \xrightarrow{\ \mathrm{Smith}\ }
 \frac{|\operatorname{Sha}(E)|\prod_vc_v}
      {|E(\mathbb Q)_{\mathrm{tors}}|^2\,\Omega_E},
\end{equation}
avec les signes de Knudsen--Mumford et les orientations de
\eqref{eq:amp-bsd-cadena-lineas}. En particulier, le membre droit de
\eqref{eq:amp-bsd-paquete-composicion} est la coordonnée finale de l'élément
transporté et non une donnée servant à construire l'une des flèches.
\end{proposition}

\begin{proof}
La première ligne de \eqref{eq:amp-bsd-paquete-identidades} est
\eqref{eq:amp-bsd-phi-kato-fers} et
\eqref{eq:amp-bsd-matriz-kato-fers}--
\eqref{eq:amp-bsd-rho-determinantal}. La deuxième est
\eqref{eq:amp-bsd-comparador-pt-q} et le lemme
\ref{lem:amp-bsd-grado-pt}. La troisième combine l'égalité de la racine avec
le remplisseur explicite \eqref{eq:bsd-suc-hstar} ; son extension linéaire dans la
forme normale induit le lecteur auxiliaire utilisé dans le lemme
\ref{lem:amp-bsd-aciclicidad-rango}. En se restreignant au complément de la
racine, cette réduction contracte tout cycle rationnel. La dernière ligne est
\eqref{eq:amp-bsd-morfismo-smith} et
\eqref{eq:amp-bsd-exacta-sha-tamagawa}. Appliquer le foncteur déterminant à
ces identités, dans ce même ordre, donne
\eqref{eq:amp-bsd-paquete-composicion}. Aucune égalité ne s'obtient
en comparant d'abord les deux scalaires finaux.
\end{proof}

\begin{lemma}[Indépendance causale du terme principal]
\label{lem:amp-bsd-independencia-termino-principal}
Soit \(\mathscr P_E\) l'ensemble de données
\begin{equation}\label{eq:amp-bsd-primitivas-paquete}
 \mathscr P_E:=
 \left(
  1_E,\Delta_{\mathrm{AW}},P_E^K,P_E^{\mathrm{PT}},
  h_*=-vf,X^{\mathrm{orth}},
  \{\mathcal C_{E,v}^{\mathrm{fin}}\xrightarrow{i_v}
     \mathcal C_{E,v}\}_v,
  \text{bases et orientations},
  \omega_E,\gamma_E^{\mathbb R}
 \right),
\end{equation}
où \(\omega_E\) est la différentielle de Néron et
\(\gamma_E^{\mathbb R}\) le cycle réel orienté. Toutes les applications du
paquet de la proposition
\ref{prop:amp-bsd-paquete-comparadores-exacto} sont déterminées par
\(\mathscr P_E\). Ne font pas partie de \(\mathscr P_E\) les numéraux
\[
 L^{(r)}(E,1),\quad r,\quad\operatorname{Reg}(E),\quad
 |\operatorname{Sha}(E)|,\quad(c_v)_v,\quad
 |E(\mathbb Q)_{\mathrm{tors}}|,\quad\Omega_E.
\]
Ces valeurs s'obtiennent, respectivement, en évaluant la section analytique,
en prenant le degré du drapeau, en calculant des déterminants et des formes de Smith,
en prenant les ordres de groupes finis et en intégrant
\(\omega_E\) sur \(\gamma_E^{\mathbb R}\).
\end{lemma}

\begin{proof}
La formule \eqref{eq:amp-bsd-phi-kato-fers} utilise uniquement
\(1_E,\Delta_{\mathrm{AW}},P_E^K\) et la retenue. Les applications de Bockstein et de
Poitou--Tate utilisent la filtration, \(P_E^{\mathrm{PT}}\), les bases et
l'autodualité. Le cône de \eqref{eq:amp-bsd-mapa-localizacion} utilise les
applications \(i_v\) et les restrictions locales ; sa contraction et sa forme de
Smith utilisent le remplisseur explicite \eqref{eq:bsd-suc-hstar}, son extension
linéaire auxiliaire, \(X^{\mathrm{orth}}\) et les bases entières. Les deux
derniers comparateurs utilisent les réseaux de torsion et la
paire \((\omega_E,\gamma_E^{\mathbb R})\). Il s'agit d'une énumération exhaustive
des arguments qui figurent dans les formules de construction. Les
numéraux énumérés dans l'énoncé n'apparaissent qu'après application de la
cohomologie, du déterminant, du cardinal ou de l'intégration, de sorte qu'ils ne peuvent pas
sélectionner rétrospectivement les applications.
\end{proof}

\subsection{Composition, degrés et terme principal}
\label{subsec:amp-bsd-composicion}

Définissons le comparateur total par la composition dans l'ordre écrit :
\begin{equation}\label{eq:amp-bsd-comparador-total}
 \mathfrak C_E:=
 \mathfrak c_\infty\circ
 \mathfrak c_{\mathrm{tors}}\circ
 \mathfrak c_{\mathrm{STS}}\circ
 \mathfrak c_{\mathrm{PT}}\circ
 \mathfrak c_{K/\mathrm{FERS}}.
\end{equation}
Comme \(p_{\mathrm{root}}z_K=z_{\mathrm{PT}}\), toutes les flèches reçoivent la
même unité basée. Le remplisseur \(h_*=-vf\) de
\eqref{eq:bsd-suc-hstar} remplit explicitement la différence ; son extension
linéaire dans la base normale produit le lecteur homotopique auxiliaire et démontre
qu'aucune classe supplémentaire ne modifie l'élément de la droite pendant le
transport. Le lecteur n'est pas une prémisse distincte.

\begin{theorem}[Construction des comparateurs basés]
\label{thm:amp-bsd-comparadores-construidos}
Pour toute courbe elliptique \(E/\mathbb Q\) munie des complexes, des bases,
des orientations et des données locales définis précédemment, les applications
\eqref{eq:amp-bsd-phi-kato-fers},
\eqref{eq:amp-bsd-comparador-pt},
\eqref{eq:amp-bsd-triangulo-localizacion},
\eqref{eq:amp-bsd-comparador-torsion} et
\eqref{eq:amp-bsd-comparador-periodo}, avec les lemmes
\ref{lem:amp-bsd-aciclicidad-rango} et
\ref{lem:amp-bsd-exactitud-sha-tamagawa}, construisent les cinq comparateurs de
\eqref{eq:amp-bsd-cadena-lineas}. Leur composition préserve l'unité, le degré
et l'orientation, et satisfait
\begin{equation}\label{eq:amp-bsd-igualdad-grados}
 d_{\mathrm{an}}=\operatorname{ord}_T\det S_E(T)
 =d_{\mathrm{ar}}=\operatorname{rank}E(\mathbb Q).
\end{equation}
De plus, \(\operatorname{Sha}(E)\) est fini et l'égalité de l'élément
distingué dans \(\mathcal L_E\) se publie sous la forme
\begin{equation}\label{eq:amp-bsd-termino-principal}
 \frac{L^{(r)}(E,1)}{r!\,\Omega_E}
 =\frac{|\operatorname{Sha}(E)|\,\operatorname{Reg}(E)
        \prod_v c_v}
       {|E(\mathbb Q)_{\mathrm{tors}}|^2},
 \qquad r=\operatorname{rank}E(\mathbb Q).
\end{equation}
\end{theorem}

\begin{proof}
Le comparateur Kato--FERS et la condition
\(\rho_{E,\ell}(0)=1\) identifient l'ordre analytique à
\(d=\operatorname{ord}_T\det S_E(T)\) et le terme initial à
\(R_{\mathrm{Boc}}^{\mathrm{der}}(S_E)\). L'identité
\eqref{eq:bsd-suc-lt} conserve la page régulière et chaque page positive. Les
isomorphismes \(j_q\) transportent ces pages vers les hauteurs dérivées. Le
lemme \ref{lem:amp-bsd-grado-pt} construit l'isomorphisme du drapeau avec le
réseau libre de Mordell--Weil et, avec
\eqref{eq:bsd-suc-carry-orden}, démontre
\eqref{eq:amp-bsd-igualdad-grados}.

Le lemme \ref{lem:amp-bsd-aciclicidad-rango} déduit l'acyclicité rationnelle
directement du remplisseur explicite \eqref{eq:bsd-suc-hstar}, de sa réduction
linéaire du complément, de l'annulation de la racine commune et de
\(X^{\mathrm{orth}}\) ; on en obtient le rang plein de \(D_E\), non
comme hypothèse supplémentaire. La forme de Smith
\eqref{eq:amp-bsd-morfismo-smith} produit alors le groupe fini \(F_E\).
Le lemme \ref{lem:amp-bsd-exactitud-sha-tamagawa} extrait des triangles
\eqref{eq:amp-bsd-triangulo-local} et
\eqref{eq:amp-bsd-triangulo-localizacion} la suite exacte finie ; ce n'est
qu'ensuite que l'on déduit la finitude de \(\operatorname{Sha}(E)\) et
\eqref{eq:amp-bsd-cardinal-smith}. Enfin,
\eqref{eq:amp-bsd-regulador-nt},
\eqref{eq:amp-bsd-comparador-torsion} et
\eqref{eq:amp-bsd-comparador-periodo} publient, respectivement, le régulateur,
le quotient de torsion et la période. L'égalité racine élimine toute unité
relative entre les deux publications. En évaluant le même élément basé aux
deux extrémités de \eqref{eq:amp-bsd-comparador-total}, on obtient
\eqref{eq:amp-bsd-termino-principal}.
\end{proof}

\subsection{Prémisses et critères de réfutation}
\label{subsec:amp-bsd-falsadores}

La construction part de l'unité généalogique \(1_E\) et de la diagonale
d'Alexander--Whitney. Elle utilise également les
publications couplées, tous les
Bockstein et la page régulière, le remplisseur explicite \(h_*=-vf\),
l'autodualité réduite \(X^{\mathrm{orth}}\), les applications de chaînes du
triangle de localisation et les bases orientées de torsion et de période. L'équivalence
Kato--FERS et sa normalisation à un, le rang plein de \(D_E\),
l'exactitude de \eqref{eq:amp-bsd-exacta-sha-tamagawa} et la finitude de
\(\operatorname{Sha}(E)\) sont des conclusions des lemmes précédents, non des
prémisses indépendantes. Le critère générateur par générateur du lemme
\ref{lem:amp-bsd-bases-emparejadas-kato-fers} est un contrôle non directeur
qui audite localement le premier comparateur. Les critères qui réfuteraient une flèche précise
sont :
\begin{enumerate}
 \item un coefficient racine différent de un, qui introduirait une unité
       relative entre \(z_K\) et \(z_{\mathrm{PT}}\) ;
 \item un \(j_q\) non inversible ou l'omission d'une page de Bockstein, qui
       modifierait le degré ou le déterminant de hauteur ;
 \item un diviseur de Smith nul, qui empêcherait la finitude de \(F_E\) ;
 \item l'échec de l'exactitude de
       \eqref{eq:amp-bsd-exacta-sha-tamagawa}, qui romprait l'identification
       des facteurs finis ;
 \item une inversion d'orientation dans la torsion ou la période, qui changerait
       l'élément basé même si elle préservait sa droite non orientée.
\end{enumerate}
La projection de Manin sans histoire et le changement d'échelle visible modulo trois
violent le premier critère ; ils constituent donc des preuves de nécessité pour
ces modèles réduits et non des prémisses du théorème
\ref{thm:amp-bsd-comparadores-construidos}.


\begin{theorem}[Comparaison déterminantale locale--globale]
\label{thm:bsd-suc-local-global}
Soient les comparateurs basés déjà construits dans le
\cref{thm:amp-bsd-comparadores-construidos} :
Kato--FERS/Iwasawa,
Poitou--Tate/Poincaré--Néron--Tate,
Smith--Tamagawa--Shafarevich, torsion et période.
Si \(d_{\mathrm{an}}\) est le degré de la section analytique dans la
spécialisation centrale et \(d_{\mathrm{ar}}\) le degré de la droite de Selmer
basée, alors
\begin{equation}\label{eq:bsd-suc-grados}
 d_{\mathrm{an}}=d=d_{\mathrm{ar}},
\end{equation}
et l'on obtient l'égalité d'éléments distingués
\begin{equation}\label{eq:bsd-suc-linea-final}
 \mathfrak z_{\mathrm{an}}(E)=\mathfrak z_{\mathrm{ar}}(E)
 \quad\text{dans}\quad\mathcal L_E.
\end{equation}
Le membre analytique est la publication du terme initial ; le membre
arithmétique est la publication du régulateur complet, des facteurs
élémentaires de Smith,
des facteurs locaux, de la torsion et de la période.
\end{theorem}

\begin{proof}
La construction basée donne
\([\Phi_{K/\mathrm{FERS}}^i](T)=I+TB_i(T)\) en chaque degré. Le comparateur
basé de Kato--FERS exprime alors la section analytique corrigée
comme
\begin{equation}\label{eq:bsd-suc-fers}
 L_{\ell}^{\mathrm{corr}}(E,T)
 =\rho_{E,\ell}(T)\det S_E^{\mathrm{corr}}(T),
 \qquad \rho_{E,\ell}(0)=1.
\end{equation}
Par \eqref{eq:bsd-suc-orden}, son degré central est \(d\) ; par
\eqref{eq:bsd-suc-lt}, son terme initial est exactement
\(R_{\mathrm{Boc}}^{\mathrm{der}}\), et non ce régulateur multiplié par une
unité indéterminée. Les isomorphismes \(j_q\) de
\eqref{eq:bsd-suc-jq} transportent chaque \(\tau_q\) vers le déterminant de la
hauteur dérivée de sa page. La page régulière transporte le terme
acyclique que les pages positives ne voient pas.

La comparaison de Poincaré--Néron--Tate identifie le degré de cette droite de
hauteurs à \(d_{\mathrm{ar}}\). Les facteurs élémentaires de Smith décomposent sa partie
finie en facteurs de Shafarevich et de Tamagawa ; le quotient des réseaux
entiers apporte le carré de la torsion, et le lecteur archimédien apporte
la période. Toutes ces applications agissent sur la même droite racine. Par le
théorème \ref{thm:bsd-suc-raiz-filler}, il ne reste pas d'unité relative entre
les publications Kato et Poitou--Tate. La composition préserve le degré et
l'élément basé, ce qui prouve simultanément
\eqref{eq:bsd-suc-grados} et \eqref{eq:bsd-suc-linea-final}. La formule
scalaire correspondante est réservée à la reconnaissance finale.
\end{proof}

\subsection*{Hypothèses et portée}

La chaîne BSD utilise : l'état enrichi complet et sa coalgèbre
d'Alexander--Whitney ; le produit fibré \eqref{eq:bsd-suc-pullback} ; les deux
publications construites sur la même unité ; la droite racine, sa base et sa
counité ; le quotient réduit parfait et saturé, avec
\(X^{\mathrm{orth}}\) ; le complexe fini--singulier complet ; les drapeaux
de Bockstein finis, leurs signes de Koszul et la page régulière ; et les
comparaisons de hauteur, de Smith, de Tamagawa, de Shafarevich, de torsion et de périodes
avec une orientation commune. La vérification degré par degré des bases
appariées et des résiduels \(\varepsilon_i(b)\) est enregistrée comme
contrôle local de la réalisation, non comme prémisse gouvernante de la fermeture.
Sous cette structure de départ explicite,
\eqref{eq:bsd-suc-raiz}, \eqref{eq:bsd-suc-hstar},
\eqref{eq:bsd-suc-lt} et \eqref{eq:bsd-suc-linea-final} sont des égalités
exactes.

La conservation de l'unité, le parcours à deux feuillets et la publication
locale--globale sont formalisés au moyen de la coalgèbre de chaînes, du produit
fibré généalogique, de la publication couplée, de la réduction orthogonale et de la
tour de Bockstein. Les identités auxiliaires vérifient séparément la propriété d'élément groupal,
la conservation de l'histoire sous \(\widehat\Pi_0\), l'identité de
l'unité dans les deux canaux, l'involution réduite, la forme normale
finie--singulière et l'identité déterminantale d'ordre arbitraire.

\paragraph{Contrôles d'ablation BSD.}
\noindent\textbf{Statut : CONTRÔLE\_NON\_DIRECTEUR.}
Les trois ablations suivantes reçoivent les identités déjà démontrées et
réfutent uniquement des modèles réduits. Elles n'apportent pas de prémisses au propriétaire,
ne conditionnent pas sa fermeture et ne possèdent pas d'arête causale de retour vers lui.

\begin{ablation}[Projection de Manin sans histoire]
\label{abl:bsd-suc-manin}
Remplacer \(P_E^{\mathrm{gen}}\) par \(P_E^{\mathrm{Man}}\) efface des données dont
la publication couplée a besoin. En effet, la voie
\(\gamma=(\mathsf N\mathsf E)^{54}\) peut avoir la même cellule et la même
phase visible que la voie triviale, mais conserve un enroulement et une mémoire
distincts dans \(G_{m,n}\). L'augmentation de Manin les identifie ; la seconde
projection de \eqref{eq:bsd-suc-pullback} ne le fait pas.
\end{ablation}

\begin{proof}
Tout relèvement basé au complexe de Manin se factorise, à homotopie près, par la
section de l'augmentation \(s\varepsilon_{\mathrm{AW}}\). Il ne dépend donc que
de la coordonnée visible. En revanche, la seconde composante de
\(\widehat\Pi_0(c)\) est \(c\) lui-même, de sorte qu'elle conserve le simplexe
d'histoire et distingue l'enroulement de \(\gamma\). L'ablation réfute la
projection nue ; elle ne réfute pas le produit fibré employé dans le théorème.
\end{proof}

\begin{ablation}[Changement d'échelle de la racine visible seulement modulo trois]
\label{abl:bsd-suc-twist}
Dans la spécialisation \(c=1\) de
\eqref{eq:bsd-suc-diferencial-normal}, le cycle altéré
\begin{equation}\label{eq:bsd-suc-twist4}
 z_K^{(4)}=4r+3e+b
\end{equation}
satisfait
\begin{equation}\label{eq:bsd-suc-twist-defecto}
 z_{\mathrm{PT}}-z_K^{(4)}
 =\partial(-f)-3r.
\end{equation}
Le défaut \(-3r\) disparaît modulo trois et réapparaît modulo neuf. Par
conséquent, \eqref{eq:bsd-suc-twist4} conserve une ombre partielle, mais non
l'unité basée ni \eqref{eq:bsd-suc-raiz}.
\end{ablation}

\begin{proof}
L'égalité \(\partial f=3e+b\) donne immédiatement
\eqref{eq:bsd-suc-twist-defecto}. Modulo trois, le second terme est nul ;
modulo neuf, c'est une classe racine non divisible par un bord du complément.
En outre, sa coordonnée racine est quatre, tandis que
\eqref{eq:bsd-suc-normalizacion-raiz} exige un. L'exemple vérifie la
nécessité du couple couplé et de l'orientation ; il n'introduit pas d'exception
au théorème qui conserve les deux données.
\end{proof}

\begin{ablation}[Perte de pages du régulateur]
\label{abl:bsd-suc-bockstein}
La matrice \(\operatorname{diag}(T^2,T)\) montre que le premier Bockstein ne
retrouve pas la retenue d'ordre supérieur. La matrice
\(\operatorname{diag}(2,T^2)\) possède une trivialisation positive
\(\tau_2=1\), mais
\(\tau_{\mathrm{reg}}=2=\operatorname{LT}_{T}(S)\). Par conséquent, ni la
première page ni le produit des pages positives ne remplacent le régulateur
complet \eqref{eq:bsd-suc-rboc}.
\end{ablation}

\begin{proof}
Dans le premier exemple, \(d=3\), \(r_0=2\) et
\(c_I^{\deg}=1\) ; la page un ne contient pas ce dernier degré de mémoire. Dans
le second, l'unique diviseur singulier est \(T^2\), mais la direction régulière
apporte le facteur deux. L'omettre produirait un au lieu de deux. Les deux calculs
sont des cas diagonaux de \eqref{eq:bsd-suc-smith} et vérifient la nécessité de
toutes les pièces de \eqref{eq:bsd-suc-rboc}.
\end{proof}
\section*{Reconnaissance mathématique finale de Birch--Swinnerton--Dyer}

Le développement précédent s'achève avant d'employer sa formulation historique
comme critère.

Dans le domaine déclaré par le théorème
\ref{thm:bsd-suc-local-global}, la comparaison de Smith identifie la
composante finie de la droite arithmétique à la décomposition
Shafarevich--Tamagawa et démontre la finitude de
\(\operatorname{Sha}(E)\). En conséquence, son ordre dans la publication
scalaire suivante est le cardinal du groupe fini identifié, non un symbole
formel ajouté a posteriori.

La traduction scalaire de \eqref{eq:bsd-suc-linea-final} fournit, avec
\(r=\operatorname{rank}E(\mathbb Q)\),
\begin{equation}\label{eq:bsd-suc-rango-final}
 \operatorname{ord}_{s=1}L(E,s)=r
\end{equation}
et
\begin{equation}\label{eq:bsd-suc-formula-final}
 \frac{L^{(r)}(E,1)}{r!\,\Omega_E}
 =\frac{
   \lvert\operatorname{Sha}(E)\rvert\,\operatorname{Reg}(E)\prod_v c_v}
  {\lvert E(\mathbb Q)_{\mathrm{tors}}\rvert^2}.
\end{equation}
C'est la formulation conventionnelle de Birch--Swinnerton--Dyer. L'ordre,
le terme principal et les facteurs arithmétiques ne sont pas définis les uns au moyen des
autres : ce sont des publications coordonnées de l'égalité basée
\eqref{eq:bsd-suc-linea-final}, dont la racine commune a été fixée au préalable par
\eqref{eq:bsd-suc-raiz}.


% El propietario conjunto anterior incorpora ya, por inclusiones internas y
% en su posición demostrativa propia, los refuerzos completos de RH,
% Navier--Stokes, Yang--Mills, Hodge y BSD.  No se vuelven a cargar aquí:
% repetirlos duplicaría teoremas y etiquetas sin añadir contenido.
\end{HMTScientificOwnerSubordinated}

Le chapitre précédent construit simultanément les cinq opérations à partir du même état enrichi du continu. Chaque composante complète est déployée ici avant son application, avec fibre, extensions, relations, nombres, orientation, retenue, mémoire, frontière, parcours, multisection, survie, termes terminaux et lecteurs. La séparation expositive des cinq tuples n’en fait ni des branches, ni des facteurs, ni des objets indépendants : ce sont cinq caractéristiques intrinsèques de l’unique continu engendré, qui conservent littéralement le même état et le même registre. Le constructeur commun ne les réunit pas après coup ; il les produit simultanément comme structure interne d’un seul objet.

\section{Domaine, applications et compatibilités de l’objet terminal commun}

Pour exposer chaque caractéristique sans la séparer des autres, on utilise l’indice $T\in\{\mathrm{sp},\mathrm{sol},\mathrm{gau},\mathrm{coh},\mathrm{det}\}$, mais la donnée mathématique complète demeure le même continu enrichi. La vue interne n’est pas postulée de nouveau : c’est exactement la lecture typée construite dans \eqref{eq:pdf2-operacion-infinita-tipificada},
\begin{equation}\label{eq:pdf2-g-t-desde-continuo}
 \mathfrak G_T^{\rm int}:=
 \operatorname{View}_T\!\left(
 \mathcal O_\infty(\mathcal C_{\rm cont}^{\rm disc})\right).
\end{equation}
Sa présentation complète dans le canal \(T\) est
\begin{equation}\label{eq:pdf2-g-t-completo}
 \mathfrak G_T^{\rm int}=
 \bigl(
 \begin{gathered}
 \{\mathcal F_q^{(T)}\}_q,
 \{\operatorname{Ext}_\gamma^{(T)}\}_\gamma,
 \operatorname{Rel}_T,\operatorname{Num}_T,\operatorname{or}_T,\\
 \operatorname{Carr}_T,\operatorname{Mem}_T,\partial_T,\gamma_T,
 \mathfrak m_T,\operatorname{Surv}_T,X_{\chi_T},\tau_T,
 \operatorname{Term}_T,\operatorname{Lect}_T
 \end{gathered}
 \bigr).
\end{equation}
Il n’est pas permis de remplacer cet objet par $X_{\chi_T}$, par son cardinal ou par le nom d’un problème classique. Soit $\mathcal D_T^{\rm cl}$ l’espace des données classiques du problème correspondant. Celles-ci n’entrent que par
\begin{equation}\label{eq:pdf2-mapa-preparacion-aplicacion}
 \operatorname{App}_{T,d_T}:\mathfrak G_T^{\rm int}\longrightarrow
 \mathfrak G_T(d_T),\qquad d_T\in\mathcal D_T^{\rm cl},
\end{equation}
où $\mathfrak G_T(d_T)$ est, par définition, le couple formé par la vue interne complète et les coordonnées d’interface déployées ci-dessous. Ainsi, aucune fibre, relation, orientation, retenue, mémoire, frontière, parcours, multisection, survie ou terme terminal n’est perdu lors de l’introduction de la donnée classique. La projection interne $\operatorname{pr}_{X,T}:\mathfrak G_T^{\rm int}\twoheadrightarrow X_{\chi_T}$ est surjective. Si $\mathcal A_{T,d_T}:\mathfrak G_T(d_T)\to\mathcal M_T$, on définit
\[
 \widetilde{\mathcal A}_{T,d_T}:=
 \mathcal A_{T,d_T}\circ\operatorname{App}_{T,d_T},
 \qquad
 \mathfrak M_T:=\operatorname{Im}\mathcal A_{T,d_T}.
\]

\section{Déploiement interne des réalisations corrélatives avant leur identification aux problèmes classiques}

La formule générique précédente ne remplace pas l’écriture des cinq caractéristiques. Dans ce volume, chaque interface est matérialisée sans transformer la vue interne en objet ou domaine indépendant. Dans les cinq présentations suivantes, la première coordonnée est toujours $\mathfrak G_T^{\rm int}$ ; les listes de coordonnées spécifiques étendent donc le registre commun sans jamais s’y substituer. Nous utiliserons
\[
 d_{\rm sp}=(R,\mathcal D_R^{\rm cof}),\quad
 d_{\rm sol}=(u_0,f,T),\quad
 d_{\rm gau}=(\Lambda_m,g),\quad
 d_{\rm coh}=(X,p),\quad
 d_{\rm det}=E/\mathbb Q.
\]
Cette ligne fixe l’instant causal exact où les problèmes classiques apparaissent : après \(\mathfrak G_T^{\rm int}\), jamais au sein de sa génération HMT.

\subsection{Action spectrale–polaire du constructeur commun et conservation de la fibre critique}

La caractéristique spectrale--polaire conserve le registre
\begin{equation}\label{eq:pdf2-g-sp-explicito}
 \mathfrak G_{\rm sp}(d_{\rm sp})=
 \bigl(
 \begin{gathered}
  \mathfrak G_{\rm sp}^{\rm int};\\
  \widehat X_\infty,\Gamma_9,\widehat R=(R,K),
  \operatorname{Carr}_R,S_0,M,\mathcal O_9,\\
  \Xi_R,P_R,\mathcal D_R^{\rm cof},
  \mathcal F_R,\operatorname{Ext}_R,\operatorname{Surv}_R,
  \tau_R,\mathcal P_R
 \end{gathered}
 \bigr).
\end{equation}
Ici, $\mathcal D_R^{\rm cof}$ est le domaine des fonctions tests ; $\Xi_R$ exprime les canaux premier, gamma et polaire dans un même espace de Hilbert ; $P_R=P_R^*=P_R^2$ est la projection naturelle ; $S_0,M,\mathcal O_9$ conservent la mémoire et le terme terminal ; et $\widehat R=(R,K)$ conserve aussi bien la fenêtre que sa profondeur. L’objet de Weil est le complément, non un signe ajouté :
\begin{equation}\label{eq:pdf2-g-sp-complemento}
 B_{W,R}:=\Xi_R^*(I-P_R)\Xi_R.
\end{equation}
Les canaux régulés, la naturalité cofinale et le terme terminal appartiennent à la même caractéristique ; aucun n’est reconstruit à partir des zéros de $\zeta$. La structure interne fournit l’involution des feuillets, l’espérance des extensions, le défaut et le télescopage terminal. Le développement de référence exact précédent construit en outre l’interface de Weil : il fixe la contraction des feuillets dans \eqref{eq:amp-rh-owner-q}, relie le complément spectral--polaire dans \eqref{eq:amp-rh-owner-binding}, conserve le terme terminal au moyen de \eqref{eq:amp-rh-owner-telescopia} et prouve son extinction dans \eqref{eq:amp-rh-owner-terminal}. Dans la notation du constructeur commun, cette colligation déjà construite s’exprime comme
\begin{equation}\label{eq:pdf2-g-sp-coligacion-dato}
 \Sigma_R^{\rm sp}:
 \mathcal H_{+,R}\oplus\mathcal X_R^{\rm sp}
 \longrightarrow
 \mathcal H_{-,R}\oplus\mathcal L_R^{\rm sp}
 \oplus\mathcal X_R^{\rm sp},
 \qquad
 (\Sigma_R^{\rm sp})^*\Sigma_R^{\rm sp}=I,
\end{equation}
avec l’état terminal $x_{R,N}=q^NV_R^Nx_{R,0}$, $q=5/9$, et la sortie $K_R^{\rm sp}J_{+,R}=J_{-,R}$. Son identité d’énergie est
\begin{equation}\label{eq:pdf2-g-sp-identidad-finita-dato}
 \|J_{+,R}f\|^2
 =\|J_{-,R}f\|^2+
   \sum_{n=0}^{N-1}\|\ell_{R,n}f\|^2+
   q^{2N}\|x_{R,0}f\|^2.
\end{equation}
La factorisation \eqref{eq:amp-rh-owner-l} et son expression cofinale \eqref{eq:amp-rh-owner-publicacion} prouvent ces identités et l’entrelaceur premier--gamma--polaire avant cet assemblage. Le chapitre consacré à Riemann les déploie et les reconnaît au moyen du théorème~\ref{thm:amp-rh-cierre-cofinal-propietario} ; il ne construit pas a posteriori une interface absente et ne la reçoit pas des zéros de $\zeta$.

\subsection{Bilan solénoïdal du constructeur commun et conservation du flux}

Pour chaque coupure $M$ et profondeur $j$, l’image est
\begin{equation}\label{eq:pdf2-g-sol-explicito}
 \mathfrak G_{{\rm sol},M,j}(d_{\rm sol})=
 \bigl(
 \begin{gathered}
  \mathfrak G_{\rm sol}^{\rm int};\\
  \Xi_{M,j}^{\rm sol},C_{M\leftarrow N},M_{M,j}^+,M_{M,j}^-,
  D_{M,j},H_{M,j},E_9,J_9,P_9,Q_{\rm micro},\\
  \Gamma_{M,j}^{\rm term},\mathcal E_{M,j,T}^{\rm term},
  L_{M,j,T}^{\rm term},\ell_{M,j},\operatorname{Carr}_{M,j},\\
  \operatorname{Mem}_{M,j},\operatorname{Surv}_{M,j},\tau_{M,j}
 \end{gathered}
 \bigr).
\end{equation}
Le champ de Galerkin exprimé est une coordonnée de $\Xi_{M,j}^{\rm sol}$, mais l’état comprend en outre le défaut $C_{M\leftarrow N}$, les mémoires orientées, la counité et sa section, la dynamique terminale et la perte observée. Le Gram
\begin{equation}\label{eq:pdf2-g-sol-gram-estructural}
 U_{M,j;J,T}^{\rm term}
 :=(\mathbf G_{M,j;J,T}^{\rm term})^*
    \mathbf G_{M,j;J,T}^{\rm term}
\end{equation}
conserve dans $\mathbf G^{\rm term}$ la coordonnée terminale et toutes les pertes propagées. Le télescopage n’en efface aucune. Le développement de référence exact précédant cet assemblage construit l’application causale hydrodynamique complète : le relèvement PC--Fock de \eqref{eq:ns-import-pcf-lift}, la colligation isométrique de \eqref{eq:ns-import-cascade-isometry}, son télescopage de Stein \eqref{eq:ns-import-stein-telescope}, le registre \eqref{eq:ns-import-ledger}, la retenue comptabilisée une seule fois \eqref{eq:ns-import-carry-funded} et le retour exact \eqref{eq:ns-import-exact-return}. Dans ces mêmes coordonnées, les ondes de diffusion s’expriment comme
\begin{equation}\label{eq:pdf2-g-sol-ondas-dato}
 a_{M,j}=(2\sqrt\nu)^{-1}A_M^{-1/2}\Lambda^sf_{M,j},
 \qquad
 b_{M,j}=\sqrt\nu A_M^{1/2}\Lambda^su_{M,j}-a_{M,j},
\end{equation}
et la transition linéaire agit sur toutes les coordonnées mixtes $u^{\odot r}\otimes f^{\odot k}$. La racine commune et son transport ne sont pas postulés dans ce bloc : ils sont déjà construits par \eqref{eq:nsclose-owner-metric-definition}--\eqref{eq:nsclose-owner-root-uniformity}. Leur expression dans la carte commune est
\begin{equation}\label{eq:pdf2-g-sol-raiz-dato}
 J_T^{\rm sol}:\mathcal D_T^{\rm NS}\longrightarrow\mathcal H_T^{\rm sol},
 \qquad
 \mathbf G_{M,0;J,T}^{\rm term}\Lambda_{M,T}J_T^{\rm sol}
 =I_{M,J}J_T^{\rm sol},
 \qquad I_{M,J}^*I_{M,J}=I.
\end{equation}
La conservation conjointe du terme terminal, de la dissipation, de l’advection, de la retenue, du microétat, de la coquille et de la mémoire est prouvée par le télescopage de référence \eqref{eq:nsclose-owner-telescope} et son registre \eqref{eq:nsclose-proprietary-ledger}. Le budget uniforme, indépendant de $M,J$, est le théorème déjà établi par \eqref{eq:nsclose-proprietary-uniform-budget} et \eqref{eq:nsclose-uniform-sobolev-precompactness} ; il s’exprime ici dans la notation du constructeur commun comme
\begin{equation}\label{eq:pdf2-g-sol-cota-dato}
 \sup_{M,J}
 \|\mathbf G_{M,0;J,T}^{\rm term}\Lambda_{M,T}J_T^{\rm sol}d\|^2
 \le C_T^{\rm sol}\|d\|_{\mathcal D_T^{\rm NS}}^2,
\end{equation}
avec $C_T^{\rm sol}$ indépendant de $M,J$. Le chapitre ultérieur déploie la réalisation de Navier--Stokes et reconnaît son codomaine classique ; il ne fournit pas une prémisse absente ni ne sélectionne la borne à partir de la conclusion.

\subsection{Incidence projective de jauge du constructeur commun et descente au quotient}

L'image de jauge n'est pas la marginale des liaisons. Son objet complet par niveau est
\begin{equation}\label{eq:pdf2-g-gau-explicito}
 \mathfrak G_{{\rm gau},m}(d_{\rm gau})=
 \bigl(
 \begin{gathered}
  \mathfrak G_{\rm gau}^{\rm int};\\
  \Lambda_m,\widehat X_m,\widehat\mu_m^{\rm ov},
  \{\widehat\mu_{m,g}^{\rm YM},\nu_{m,g},K_{m,g}^{\rm phys}\}_{g>0},\\
  \mathcal G_m^{\rm full},\mathbb P_m^{\rm phys},
  \widehat J_m,\widehat\Gamma_m,
  \{E_{m,c}\}_c,\widehat H_m^{\rm ov},D_m^{\rm YM},\\
  \{\Theta_{\nu,m}\}_{\nu=1}^4,
  \{\mathfrak A_{+,\nu,m},\operatorname{ev}_{t}^{(\nu)},
      \mathcal V_{m,\nu}^{\rm OS}\}_{\nu=1}^4,\\
  \widehat\Omega_m^{(8)},P_{\Omega_m},W_c,
  \operatorname{Carr}_m,\operatorname{Mem}_m,
  \operatorname{Surv}_m,\tau_m,\mathcal P_m^{F^2}
 \end{gathered}
 \bigr).
\end{equation}
Le groupe $\mathcal G_m^{\rm full}$ agit simultanément sur les liens, les repères et les
coordonnées d’incidence $q_c\in\mathbb F_3^6$. La projection physique est sa moyenne
de Haar. Le générateur réduit est
\begin{equation}\label{eq:pdf2-g-gau-generador-estructural}
 D_m^{\rm YM}
 =\left.3\kappa_{\rm av}\sum_c(I-E_{m,c})
  \right|_{\operatorname{Ran}\mathbb P_m^{\rm phys}},
\end{equation}
et $W_c$ appartient à cette sous-algèbre. Les quatre réflexions agissent sur une même mesure, un même noyau et un même terminal.
Le propriétaire exact précédent construit l'action de jauge d'incidence \eqref{eq:amp-ym-accion-gauge-incidencia}, la projection physique de Haar \eqref{eq:amp-ym-proyeccion-fisica-haar} et la compression entrelacée \eqref{eq:amp-ym-compresion-entrelazada}. De cette construction proviennent le noyau complet équivariant $\widehat K_{m,t}^{\rm ov}$, sa descente au quotient physique $K_{m,t}^{\rm phys}$, les quatre opérateurs de réflexion et l'action $U_m(g)$ de $E(4)$ sur le cœur cylindrique lisse. Ce bloc assemble, sans les introduire comme données, le générateur, le vide, le témoin de courbure et les formes OS :
\begin{equation}\label{eq:pdf2-g-gau-entrelazamiento-dato}
 D_{m+1}^{\rm YM}\widehat J_m=\widehat J_mD_m^{\rm YM},
 \quad
 \mathbb P_{m+1}^{\rm phys}\widehat J_m
  =\widehat J_m\mathbb P_m^{\rm phys},
 \quad
 \widehat J_mW_{m,c}=W_{m+1,c}.
\end{equation}
La première et la troisième identités sont la publication, dans la notation
commune, de \eqref{eq:amp-ym-compresion-entrelazada} et de
\eqref{eq:amp-ym-testigo-fisico-gap} ; le semi-groupe, le vide simple et
l’écart exact s’obtiennent déjà dans
\eqref{eq:amp-ym-owner-semigrupo-vacio}--
\eqref{eq:amp-ym-owner-brecha-exacta}. Ici,
$\nu_{m,g}:=(\pi_m^{\rm phys})_*\widehat\mu_{m,g}^{\rm YM}$ et
$K_{m,g}^{\rm phys}(t)=e^{-tD_{m,g}^{\rm YM}}$ dans le quotient physique.
La complétude de jauge comprend en outre la désintégration selon chaque famille
d’hyperplans euclidiens. Si $A_j\in\mathfrak A_{+,\nu,m}$ et
$t_1\le\cdots\le t_n$, la marginale de la même mesure physique satisfait
\begin{equation}\label{eq:pdf2-g-gau-transfer-euclidean-dato}
\begin{aligned}
 &\int\prod_{j=1}^n
   (\tau_{t_je_\nu}A_j)(x)\,d\widehat\mu_{m,g}^{\rm YM}(x)\\
 &\quad=
 \int\prod_{j=1}^nA_j(y_j)\,
  \nu_{m,g}(dy_1)
  \prod_{j=1}^{n-1}K_{m,g}^{\rm phys}
       (t_{j+1}-t_j;y_j,dy_{j+1})\\
 &\quad=
 \left\langle\mathbf1,
  M_{A_1}e^{-(t_2-t_1)D_{m,g}^{\rm YM}}M_{A_2}\cdots
  e^{-(t_n-t_{n-1})D_{m,g}^{\rm YM}}M_{A_n}\mathbf1
 \right\rangle.
\end{aligned}
\end{equation}
La reconstruction précédente prouve l’identité dans les quatre directions
au moyen de \eqref{eq:amp-ym-cuatro-os} et construit un unique hamiltonien
entrelacé dans \eqref{eq:amp-ym-hamiltoniano-comun}. Par conséquent, la
translation OS est publiée ici comme
\begin{equation}\label{eq:pdf2-g-gau-os-generator-dato}
 \mathcal V_{m,\nu}^{\rm OS}T_{m,\nu}^{\rm OS}(s)
 (\mathcal V_{m,\nu}^{\rm OS})^{-1}
 =e^{-sD_{m,g}^{\rm YM}}.
\end{equation}
Ainsi, le générateur d’espérances et le Hamiltonien des translations ne sont pas
identifiés par leur nom : ce sont deux réalisations déjà construites de la même
désintégration, dont la forme limite et la publication de courbure figurent dans
\eqref{eq:amp-ym-forma-limite} et \eqref{eq:amp-ym-f2}.

\subsection{Cohérence cohomologique du constructeur commun et émergence de la classe globale}

Pour une variété projective lisse complexe $X$ et une codimension $p$, l’image cohomologique conserve
\begin{equation}\label{eq:pdf2-g-coh-explicito}
 \mathfrak G_{\rm coh}(d_{\rm coh})=
 \bigl(
 \begin{gathered}
  \mathfrak G_{\rm coh}^{\rm int};\\
  \{\mathscr S_N(X,p),\rho_{N+1,N}\}_N,
  \{\mathscr H_N^p(X),q_{N+1,N}\}_N,
  \{e_N\}_N,
  \{\iota_{N+1,N}\}_N,\\
  \{\mathcal G_{N+1}^{\rm coh},\mathcal R_{N+1}^{\rm coh},
       D_{N+1},a_{N+1},b_{N+1}\}_N,\\
  \mathcal F_{X,p},\operatorname{Ext}_{X,p},X_\chi(X,p),
  \operatorname{Rel}_{X,p},\operatorname{Carr}_{X,p},
  \operatorname{Mem}_{X,p},\partial_{X,p},\gamma_{X,p},\\
  \Sigma_{X,p},\operatorname{Surv}_{X,p},\tau_{X,p},
  R_{\rm Hdg},\operatorname{ch}^{\rm loc}_p,\operatorname{cl}_{X,p}
 \end{gathered}
 \bigr).
\end{equation}
Les carrés
\begin{equation}\label{eq:pdf2-g-coh-cuadrados}
 e_N\rho_{N+1,N}=q_{N+1,N}e_{N+1}
\end{equation}
relient observation, extension et raffinement. Le développement de référence exact précédent construit, sans la recevoir comme donnée, l’histoire basée avec ses corrections relatives, sa retenue, sa mémoire, sa frontière et son terme terminal. En particulier, \eqref{eq:amp-hodge-mapas-relativos}--\eqref{eq:amp-hodge-factorizacion-relativa} construisent les applications de l’incrément et \eqref{eq:amp-hodge-seccion-relativa} construit son inverse à droite. Dans la notation du constructeur commun, cette application relative s’exprime comme
\begin{equation}\label{eq:pdf2-g-coh-seccion-dato}
 s_{N+1}^{\rm rel}:\ker q_{N+1,N}\longrightarrow\ker\rho_{N+1,N},
 \qquad
 e_{N+1}^{\rm rel}s_{N+1}^{\rm rel}=I,
\end{equation}
naturelle sous raffinement par \eqref{eq:amp-hodge-naturalidad-seccion-relativa} et compatible avec la retenue et le terme terminal. La section initiale n’est pas non plus supposée : elle est construite dans \eqref{eq:amp-hodge-seccion-base}. Ce sont ces applications déjà démontrées, et non une classe finale choisie, que le bloc assemble et que le chapitre ultérieur reconnaît dans le codomaine de Hodge.

La complétude repose sur des actions et des identités construites dans le même développement de référence. La dégénérescence unitaire, la présentation relative et le lecteur sont assemblés ici au moyen de
\begin{equation}\label{eq:pdf2-g-coh-datos-operativos}
\begin{gathered}
 \iota_{N+1,N}:\mathscr S_N(X,p)\longrightarrow\mathscr S_{N+1}(X,p),
 \qquad \rho_{N+1,N}\iota_{N+1,N}=I,\\
 \mathscr A_{N+1}^{\rm rel}
  =\mathbb Q^{(\mathcal G_{N+1}^{\rm coh})}/\operatorname{im}D_{N+1},
 \qquad e_{N+1}^{\rm rel}a_{N+1}=b_{N+1},
 \qquad b_{N+1}D_{N+1}=0.
\end{gathered}
\end{equation}
L’ensemble $\mathcal G_{N+1}^{\rm coh}$ énumère exhaustivement tous les atomes nouveaux ainsi qu’une étiquette terminale par composante. La forme normale et l’inverse à droite sont construits par réduction basée dans \eqref{eq:amp-hodge-sigma-coordenada} ; leur compatibilité avec les raffinements est prouvée dans \eqref{eq:amp-hodge-naturalidad-seccion-relativa}. Ainsi, dans l’ordre selon profondeur, support, feuillet, parcours, retenue, mémoire et terme terminal, chaque atome d’observation possède une unique étiquette maximale et
\begin{equation}\label{eq:pdf2-g-coh-tabla-dato}
 b_{N+1}(c_\mu)=\bar c_\mu+
   \sum_{\nu<\mu}B_{\nu\mu}\bar c_\nu,
 \qquad B_{\nu\mu}\in\mathbb Q.
\end{equation}
La triangularité est une conséquence de cette construction basée, antérieure à toute classe $\alpha$ ; ce n’est ni une hypothèse ni un tableau fourni par le problème terminal.

La complétude et le lecteur du même objet sont typés par
\begin{equation}\label{eq:pdf2-g-coh-limite-lector-dato}
 X_\chi(X,p)=\varprojlim_N\mathscr S_N(X,p),
 \qquad
 q_N\operatorname{ev}^{\rm Hdg}_{\infty,X,p}
  =e_N\operatorname{pr}_N,
 \qquad
 \operatorname{cl}_{X,p}\operatorname{ch}^{\rm loc}_pR_{\rm Hdg}
  =\operatorname{ev}^{\rm Hdg}_{\infty,X,p}.
\end{equation}
Les projections $q_N$, y compris la coordonnée terminale, séparent les points du codomaine d’évaluation. L’extension opératoire \eqref{eq:amp-hodge-extension-operativa}--\eqref{eq:amp-hodge-ext}, le passage à la limite \eqref{eq:amp-hodge-xi-limite} et l’expression finale \eqref{eq:amp-hodge-publicacion-final} prouvent déjà que toute classe de Hodge arbitraire admet une histoire compatible. Le chapitre d’application déploie et reconnaît cette expression ; il ne fournit pas l’existence comme prémisse ultérieure.

\subsection{Droite déterminantale du constructeur commun et évaluation arithmétique terminale}

Pour une courbe elliptique $E/\mathbb Q$, l’image déterminantale conserve
\begin{equation}\label{eq:pdf2-g-det-explicito}
 \mathfrak G_{\rm det}(d_{\rm det})=
 \bigl(
 \begin{gathered}
  \mathfrak G_{\rm det}^{\rm int};\\
  \mathcal T_m^{\rm surv},G_{m,n},\Delta_{\rm AW},\varepsilon_{\rm AW},
  P_E^{\rm gen},1_E,P_E^K,P_E^{\rm PT},\\
  z_K,z_{\rm PT},X^{\rm orth},p_{\rm root},
  \mathcal C_E^K,\mathcal C_E^{\rm Iw},\mathcal C_E^{\rm PT},\\
  \ell_E,\mathcal C_E^{\rm loc,red},S_E(T),
  \operatorname{Carr}_E,\operatorname{Mem}_E,
  \operatorname{Surv}_E,\tau_E,\mathcal L_E
 \end{gathered}
 \bigr).
\end{equation}
L’unité $1_E$ est coexprimée dans les branches de Kato et de Poitou--Tate avant la dualité et la racine. Le produit fibré conserve l’image modulaire et l’histoire profonde. Le cône de localisation, les pages de Bockstein et la droite déterminante appartiennent à une même multisection basée ; ce ne sont pas des invariants choisis séparément après connaissance du rang ou du terme principal. Le développement de référence exact précédent construit la décomposition basée $\Theta_E$, l’homotopie du facteur fini--singulier, le comparateur complet Kato--FERS, les applications de pages et de recollement de Poitou--Tate et les applications de cocycles locales--globales. La chaîne commence dans la droite déterminante \eqref{eq:amp-bsd-linea-determinante}, construit l’application de chaînes \eqref{eq:amp-bsd-phi-kato-fers}, le comparateur de Poitou--Tate \eqref{eq:amp-bsd-comparador-pt} et le triangle local \eqref{eq:amp-bsd-triangulo-localizacion}. Ses identités opératoires s’expriment ici comme
\begin{equation}\label{eq:pdf2-g-det-identidades-dato}
\begin{gathered}
 \Phi_{K\leftrightarrow{\rm FERS}}\Psi_{K\leftrightarrow{\rm FERS}}=I,
 \qquad
 \Psi_{\rm PT}^{-1}\Psi_{\rm PT}=I,
 \qquad
 \Psi_{\rm PT}\Psi_{\rm PT}^{-1}=I,\\
 \ker\partial_{\rm loc}=\operatorname{im}\iota_{\Sha},
 \qquad
 \pi_\Phi\operatorname{Lift}_\Phi=I.
\end{gathered}
\end{equation}
Ces identités sont des conséquences de l’ensemble exact \eqref{eq:amp-bsd-paquete-identidades}--\eqref{eq:amp-bsd-paquete-composicion}. Le comparateur total est construit dans \eqref{eq:amp-bsd-comparador-total} et conserve la base commune du côté analytique jusqu’au côté arithmétique ; rang et terme principal sont exprimés ensuite dans \eqref{eq:amp-bsd-igualdad-grados}--\eqref{eq:amp-bsd-termino-principal}. Le chapitre BSD développe et reconnaît cette composition sans sélectionner rétrospectivement aucune de ses sorties.

\section{\(5\) lecteurs terminaux non permutables du continu généalogique}

L’application d’expression $\mathcal P_T$ est le lecteur mathématique qui traduit la structure dans le codomaine classique. La chaîne propre à ce volume est
\begin{equation}\label{eq:pdf2-cinco-cadenas}
 \mathfrak G_T^{\rm int}\xrightarrow{\operatorname{App}_{T,d_T}}
 \mathfrak G_T(d_T)\xrightarrow{\mathcal A_{T,d_T}}\mathfrak M_T
 \xrightarrow{\mathcal P_T}\mathcal C_T.
\end{equation}
Les correspondances sont fixées une par une :
\[
 \mathcal C_{\rm sp}=\mathrm{RH},\quad
 \mathcal C_{\rm sol}=\mathrm{NS},\quad
 \mathcal C_{\rm gau}=\mathrm{YM},\quad
 \mathcal C_{\rm coh}=\mathrm{Hodge},\quad
 \mathcal C_{\rm det}=\mathrm{BSD}.
\]
Dans l’assemblage actif, les développements de référence exacts qui précèdent cette section construisent également les cinq expressions finales : RH dans le théorème~\ref{thm:amp-rh-cierre-cofinal-propietario}, Navier--Stokes dans les théorèmes~\ref{thm:nsclose-proprietary-full} et \ref{thm:nsclose-global-smooth}, Yang--Mills dans le théorème~\ref{thm:amp-ym-terminacion}, Hodge dans le théorème~\ref{thm:amp-hodge-completitud-coinductiva} et BSD dans le théorème~\ref{thm:amp-bsd-comparadores-construidos}. Cette section enregistre leur assemblage corrélatif dans l’unique continu généalogique ; les chapitres ultérieurs déploient les cinq réalisations et leur reconnaissance classique, sans fournir de prémisses absentes ni utiliser les conclusions comme entrées.

\begin{falsifier}[Ablation de la donnée initiale]
Si une démonstration commence dans $X_{\chi_T}$, transforme une coordonnée de $\mathfrak G_T^{\rm int}$ en sélecteur facultatif ou remplace une fibre par son cardinal, elle n’a pas abrégé le même objet : elle l’a remplacé. La conclusion n’est pas transférable au substitut.
\end{falsifier}

\section{Factorisation des lecteurs terminaux par naturalité, raffinement et passage à la limite}

Le point de départ matériel est l’unique constructeur simultané qui produit un seul continu enrichi doté de cinq opérations intrinsèques corrélatives. Les notations \(\mathfrak G_{\rm sp}\), \(\mathfrak G_{\rm sol}\), \(\mathfrak G_{\rm gau}\), \(\mathfrak G_{\rm coh}\) et \(\mathfrak G_{\rm det}\) désignent cinq vues complètes de ce même registre, reproduit avec ses fibres, relations, nombres, orientation, retenue, mémoire, frontière, parcours, multisection, survie, termes terminaux et lecteurs. Elles ne sont considérées ni comme cinq branches choisies ni comme cinq données indépendantes : on démontre leur compatibilité sur le même registre, leur naturalité sous raffinement, leur non-vacuité conjointe et leur passage à la limite, puis on les applique respectivement aux cinq problèmes. Aucune preuve de ce volume ne renvoie à une autre source externe pour compléter une définition, une identité ou une étape logique.

Chaque fermeture est exposée de manière autonome : définitions, démonstrations et compatibilités apparaissent avant sa conclusion classique. Dans chaque résolution, quatre couches distinctes sont mises en évidence :
\begin{enumerate}
  \item identité exacte à chaque niveau fini ;
  \item compatibilité des raffinements et préservation de l’histoire enrichie ;
  \item uniformité et passage à la limite, avec conservation du terme terminal ;
  \item reconnaissance de la conclusion classique dans le codomaine exprimé.
\end{enumerate}

La composition directrice des cinq fermetures conserve d’abord l’action corrélative unique, ne forme le continu qu’au passage à la limite, puis ouvre une vue pour chaque application :
\[
\begin{aligned}
 \mathcal E_k&\xrightarrow{\ \mathcal O_k\ }\mathcal Z_k
 \xrightarrow{\ \varprojlim_k\ }
 \bigl(\mathcal C_{\rm cont}^{\rm disc},\mathcal O_\infty\bigr),\\
 \bigl(\mathcal C_{\rm cont}^{\rm disc},\mathcal O_\infty\bigr)
 &\xrightarrow{\ \operatorname{View}_T\ }\mathfrak G_T^{\rm int}
 \xrightarrow{\ \operatorname{App}_{T,d_T}\ }\mathfrak G_T(d_T),\\
 \mathfrak G_T(d_T)&
 \xrightarrow{\widetilde{\mathcal A}_{T,d_T}}\mathfrak M_T
 \xrightarrow{\mathcal P_T}\mathcal C_T,
 \qquad T\in\{\mathrm{sp,sol,gau,coh,det}\}.
\end{aligned}
\]
Ici, $\mathcal O_k$ agit une seule fois sur l’état complet et conserve ensemble les cinq caractéristiques. $\operatorname{View}_T$ ne restreint pas le domaine ni ne sélectionne les histoires : il expose seulement les coordonnées d’une application au sein de la valeur complète de $\mathcal O_\infty$. La donnée classique $d_T$ n’entre que dans $\operatorname{App}_{T,d_T}$, après le continu, et ne participe pas à la génération de la caractéristique intrinsèque. La droite réelle est également une sortie ultérieure de $\operatorname{ev}_{\mathbb R}$, et non une prémisse de la chaîne.

Chaque chapitre distingue le domaine de départ, la construction, le certificat, la conclusion démontrée et ses réfutateurs exacts. Les développements de référence précédents construisent les cinq compositions dans leurs domaines respectifs ; les chapitres suivants déploient ces constructions et effectuent la reconnaissance finale. Aucune flèche n’est remplacée par un renvoi, une conclusion nominale ou une hypothèse classique.

\begingroup
\HMTDemoteHeadings
\chapter{Un même bord : flux intérieur et tension de jauge}
\label{chap:hoja-comun-ns-ym}
\index{Navier--Stokes}\index{Yang--Mills}
\index{opérateur de bord}\index{Stokes}

L'hydrodynamique et Yang--Mills ne sont pas réunis en comparant de l'extérieur deux
équations célèbres. Ils proviennent de deux représentations du même opérateur de
bord. L'orientation que le lecteur intérieur interprète comme circulation
devient, sous le lecteur réfléchi, une holonomie ; la mémoire symétrique
qui contrôle le flux devient un noyau positif de transfert.

\section{L'état commun de bord}

Soit \(K_m\) un complexe fini orienté, \(B_m\) son algèbre de données de
bord et \(d_m\) son cobord. L'état enrichi est
\[
 s_m=(b_m,U_m,\theta_m,c_m,\Omega_m),
\]
où \(U_m\) transporte, \(\theta_m^2=I\) inverse l'orientation,
\(c_m\) conserve la mémoire et \(\Omega_m\) est une \(1\)-cochaîne de connexion.
Les projecteurs de parité sont
\[
 E_\pm=\frac12(I\pm\theta_m).
\]
La partie \((U_m+U_m^\dagger)/2\) est symétrique ; ce n'est un projecteur que si l'on
démontre en outre son idempotence.

Sur le complexe de cochaînes, on construit deux lecteurs linéaires :
\[
 \rho_{{\rm int},m}:C^0(K_m)\longrightarrow
 C^1(K_m)\oplus C^0(K_m),
\qquad
 \rho_{{\rm int},m}(v)=(d_mv,d_m^\ast d_mv),
\]
\[
 \rho_{{\rm ref},m}:C^1(K_m)\longrightarrow
 C^0(K_m)\oplus C^1(K_m),
\qquad
 \rho_{{\rm ref},m}(a)=(d_m^\ast a,d_md_m^\ast a).
\]
Le premier transporte un potentiel de sommets vers son flux d'arêtes et sa
réponse intérieure ; le second transporte une cochaîne d'arêtes vers sa
réponse de bord et son opérateur réfléchi. Les codomaines sont distincts,
mais tous deux sont construits avec le même cobord. Circulation, vorticité,
holonomie et courbure apparaissent ensuite au moyen de l'identité de Stokes ; elles ne
sont pas présupposées à l'intérieur des symboles \(\rho\).

\section{L'opérateur minimal et son spectre}

Dans le cycle orienté à trois sommets, soit
\[
 d_0=
 \begin{pmatrix}
 -1&1&0\\
 0&-1&1\\
 1&0&-1
 \end{pmatrix}.
\]
Les laplaciens de Hodge de sommets et d'arêtes sont
\[
 \mathcal L_{\rm H}^{(0)}=d_0^\ast d_0,
 \qquad
 \mathcal L_{\rm H}^{(1)}=d_0d_0^\ast,
\]
et un calcul direct donne
\[
 \mathcal L_{\rm H}^{(0)}=\mathcal L_{\rm H}^{(1)}=
 \begin{pmatrix}
 2&-1&-1\\
 -1&2&-1\\
 -1&-1&2
 \end{pmatrix}
 =3I-J.
\]
Dans cet emplacement, les exposants \((0)\) et \((1)\) désignent les
laplaciens de Hodge sur les sommets et les arêtes. Ils ne désignent pas les opérateurs de
relèvement \(L_0,L_1\) de la chaîne \(w_6\to w_{30}\). Pour
\(\mathcal L_{\rm cyc}^{\rm TPK}:=2I-C-C^{-1}=3I-J\), la réalisation de
l'opérateur cyclique du TPK est typée au moyen de
\[
 \mathfrak R_{\rm hyd}(\mathcal L_{\rm cyc}^{\rm TPK})
 =\mathcal L_{\rm H}^{(0)},\qquad
 \mathfrak R_{\rm YM}(\mathcal L_{\rm cyc}^{\rm TPK})
 =\mathcal L_{\rm H}^{(1)}.
\]

\begin{theorem}[Opérateur commun de bord]
\label{thm:operador-comun-frontera-ns-ym}
Dans le cycle à trois cellules,
\[
 \operatorname{Spec}(\mathcal L_{\rm H}^{(0)})
 =\operatorname{Spec}(\mathcal L_{\rm H}^{(1)})
 =\{0,3,3\}.
\]
Après retrait du mode constant ---moyenne nulle dans le lecteur intérieur et
jauge dans le lecteur de connexion--- les deux secteurs possèdent un écart spectral exact \(3\).
Dans tout complexe fini, \(d^\ast d\) et \(dd^\ast\) ont le même spectre
non nul, avec multiplicités.
\end{theorem}

\begin{proof}
Le vecteur \((1,1,1)\) engendre le noyau de \(3I-J\). Sur son complément
orthogonal, \(J=0\), de sorte que l'opérateur est \(3I\). Pour l'énoncé
général, une décomposition en valeurs singulières \(d=U\Sigma V^\ast\) produit
\[
 d^\ast d=V\Sigma^\ast\Sigma V^\ast,
 \qquad
 dd^\ast=U\Sigma\Sigma^\ast U^\ast.
\]
Les carrés des valeurs singulières non nulles, y compris leurs
multiplicités, coïncident.
\end{proof}

\section{Lecture hydrodynamique}

Nous définissons le pas implicite
\[
 S_{\rm NS}=(I+\mathcal L_{\rm H}^{(0)})^{-1}.
\]
Dans le cycle minimal,
\[
 S_{\rm NS}=\frac14(I+J).
\]

\begin{theorem}[Contraction globale du flux fini]
\label{thm:contraccion-flujo-finito}
Pour tout \(v_0\) de moyenne nulle, l'itération
\(v_{k+1}=S_{\rm NS}v_k\) existe et est unique pour tout \(k\geq0\), et
\[
 v_k=4^{-k}v_0,
 \qquad
 \|v_k\|\leq4^{-k}\|v_0\|,
 \qquad
 \|v_k\|^2=16^{-k}\|v_0\|^2.
\]
\end{theorem}

\begin{proof}
À moyenne nulle, on a \(Jv=0\), donc
\(S_{\rm NS}v=v/4\). La formule s'obtient par récurrence.
\end{proof}

\Needspace{4\baselineskip}
Une non-linéarité orientée \(N_m(v,c)\) conserve l'énergie lorsque
\[
 \langle v,N_m(v,c)\rangle=0.
\]
L'advection redistribue alors l'énergie et l'opérateur positif contrôle
sa propagation. La mémoire \(c_m\) n'est pas absorbée dans une viscosité choisie :
elle fait partie de l'état qui doit être transporté entre les échelles.

\begingroup
\setlength{\parskip}{0pt}
\titlespacing*{\section}{0pt}{8pt plus 1pt minus .5pt}{2.5pt}
\setlength{\abovedisplayskip}{.25\baselineskip}
\setlength{\belowdisplayskip}{.25\baselineskip}
\setlength{\abovedisplayshortskip}{.12\baselineskip}
\setlength{\belowdisplayshortskip}{.18\baselineskip}
\section{Lecture Yang--Mills}

Une \(1\)-cochaîne \(a\) est interprétée comme un potentiel de connexion et possède
une énergie quadratique
\[
 Q_{\rm YM}(a)=\langle a,\mathcal L_{\rm H}^{(1)}a\rangle.
\]

\begin{theorem}[Coercivité et décroissance de jauge finies]
\label{thm:coercividad-gap-gauge-finito}
Dans le quotient par le mode de jauge constant,
\[
 Q_{\rm YM}(a)\geq3\|a\|^2.
\]
Le semi-groupe \(K_t=e^{-t\mathcal L_{\rm H}^{(1)}}\) satisfait
\[
 \|K_ta\|\leq e^{-3t}\|a\|
\]
dans le secteur physique du cycle minimal.
\end{theorem}

\begin{proof}
Sur le complément du noyau, le
\cref{thm:operador-comun-frontera-ns-ym} donne
\(\mathcal L_{\rm H}^{(1)}=3I\). Les deux formules
découlent du calcul fonctionnel.
\end{proof}

Pour une connexion non abélienne,
\[
 F_\Omega=d\Omega+\Omega\smile\Omega,
\qquad
 W_\gamma(\Omega)=\operatorname{Tr}\operatorname{Hol}_\gamma(\Omega).
\]
La positivité par réflexion requiert une factorisation sur la moitié
positive,
\[
 \langle\theta f,Kf\rangle=\|Cf\|^2\geq0;
\]
la positivité ponctuelle de la matrice \(K\) ne la remplace pas.

\section[Identité de Stokes : circulation--vorticité et holonomie--courbure]{Identité discrète de Stokes et \(2\) réalisations : circulation--vorticité et holonomie--courbure}

Pour une \(1\)-cochaîne \(\Omega\) et une \(2\)-chaîne \(K\),
\[
 \sum_{e\subset\partial K}\Omega(e)
 =
 \sum_{f\in K^{(2)}}(d\Omega)(f).
\]
Les arêtes intérieures s'annulent par orientation. Le lecteur
hydrodynamique interprète le membre de gauche comme une circulation et celui de droite
comme une vorticité accumulée. Le lecteur de jauge interprète le premier comme une
holonomie linéarisée et le second comme une courbure.

\begin{theorem}[Théorème du feuillet commun de bord]
\label{thm:hoja-comun-frontera}
La décomposition \(B_m=E_-B_m\oplus E_+B_m\) possède deux représentations
typées :
\[
\boxed{
\begin{array}{ccl}
E_-B_m&\longmapsto&
\text{circulation}\ /\ \text{holonomie},\\[2mm]
E_+B_m&\longmapsto&
\text{contrôle dissipatif}\ /\ \text{noyau réfléchi}.
\end{array}}
\]
L'égalité de Stokes et la coïncidence du spectre non nul appartiennent à
l'objet commun. La contraction hydrodynamique et l'écart spectral de jauge sont des conséquences
distinctes dans leurs codomaines respectifs.
\end{theorem}

\begin{proof}
Les projecteurs \(E_\pm\) sont complémentaires parce que
\(\theta_m^2=I\). L'identité de Stokes fournit l'entrelacement du
secteur orienté ; le \cref{thm:operador-comun-frontera-ns-ym} fournit le
spectre commun. Les
\cref{thm:contraccion-flujo-finito,thm:coercividad-gap-gauge-finito}
calculent les deux réalisations.
\end{proof}

Le théorème est une fermeture finie HMT : il n'identifie pas une EDP à une théorie
de jauge ni ne confond l'écart spectral \(3\) avec une masse physique. Il explique par construction
pourquoi la régularité intérieure et la séparation de bord sont deux fonctions de
la même mémoire orientée. Tout passage continu ultérieur doit conserver
l'opérateur, la norme et le spectre qui rendent cette identité exacte.
\endgroup

\endgroup

\HMTFlushCurrentChapter
\chapter{Factorisation positive de la forme de Guinand--Weil et localisation des zéros non triviaux de \texorpdfstring{\(\zeta\)}{zeta}}
\begingroup
\renewcommand{\chapter}[2][]{}%
\chapter{Image spectrale--polaire et réduction exacte de Riemann--Weil}

\section{Point de départ : la structure spectrale--polaire complète}

L'entrée de la résolution n'est pas la forme de Weil isolée ni une liste de
zéros. C'est la structure spectrale--polaire complète produite par le constructeur
commun et attestée par le propriétaire exact qui précède ce chapitre ; ce n'est pas
une donnée externe. La chaîne causale propre à ce chapitre est
\begin{equation}\label{eq:pdf2-rh-cadena-procedencia}
 \mathfrak G_{\rm sp}
 \xrightarrow{\ \operatorname{pr}_{X,{\rm sp}}\ }X_{\chi_{\rm sp}}
 \xrightarrow{\ \mathcal A_{\rm sp}\ }\mathfrak M_{\rm sp}
 \xrightarrow{\ \mathcal P_{\rm sp}\ }\mathcal C_{\rm GW}.
\end{equation}
Le codomaine classique n'apparaît que dans la dernière flèche ; il ne sélectionne pas
rétrospectivement l'état, les cartes, la projection ni le terminal.

La structure de départ contient l'objet à treize coordonnées
\begin{equation}\label{eq:pdf2-rh-res-sp-completa}
\begin{aligned}
 \mathfrak G_{\rm sp}^{(13)}=\bigl(&
 \mathcal F_{\rm sp},
 \operatorname{Ext}_{\Gamma,{\rm sp}},
 \operatorname{Rel}_{\rm sp},
 \mathcal N_{\rm sp},
 \operatorname{or}_{\rm sp},
 \operatorname{Carr}_{\rm sp},
 \operatorname{Mem}_{\rm sp},\\
 &\partial_{\rm sp},
 \gamma_{\rm sp},
 \Sigma_{\rm sp}^{\rm mult},
 \operatorname{Surv}_{\rm sp},
 \tau_{\rm sp}^{\rm term},
 \mathcal L_{\rm sp}\bigr).
\end{aligned}
\end{equation}
Chaque coordonnée intervient matériellement :
\begin{enumerate}
 \item \(\mathcal F_{\rm sp}\) conserve les deux feuillets, leur orientation, la
 paire résidu--quotient et la signature de la publication double.
 \item \(\operatorname{Ext}_{\Gamma,{\rm sp}}\) transporte cette fibre par la
 classe dirigée de coupures et de niveaux nonadiques.
    \item \(\operatorname{Rel}_{\rm sp}\) contient le cocycle de retenue,
    l'orthogonalité \(P=P^*=P^2\), les deux publications et le complément
    structurel \(B_{\rm sp}^{\rm str}:=\Xi^*(I-P)\Xi\). L'identification
    de ce complément à la forme de Guinand--Weil est une flèche distincte
    et fait l'objet d'un audit séparé dans ce chapitre.
 \item \(\mathcal N_{\rm sp}\) enregistre le niveau \(m\), la fenêtre \(R\),
 le quotient \(q=5/9\), les moments des deux feuillets et les indices des
 canaux premier, gamma et polaire.
 \item \(\operatorname{or}_{\rm sp}=\widehat R=(R,K)\) maintient distinctes
 les deux cartes orientées.
 \item \(\operatorname{Carr}_{\rm sp}=\operatorname{Carr}_R\) est la
 retenue opératorielle et ne s'annule pas lors de la publication d'une translation.
 \item \(\operatorname{Mem}_{\rm sp}=(S_0,M,\mathcal O_9)\) conserve
 l'identité finie de Stein et son terminal.
 \item \(\partial_{\rm sp}=(I-P)\Xi\) est la frontière qui n'apparaît pas dans la
 compression \(P\Xi\) et dont l'énergie définit le complément structurel. Son
 identification à la forme de Weil n'est pas incorporée à cette définition.
 \item \(\gamma_{\rm sp}\) est le chemin cofinal de fenêtres, avec un unique
 régulateur pour tous les canaux.
 \item \(\Sigma_{\rm sp}^{\rm mult}\) est la multisection des deux
 publications et de leurs composantes premier--gamma--polaire ; ce n'est pas la signature d'un
 unique état.
 \item \(\operatorname{Surv}_{\rm sp}\) conserve la famille sous
 \(m\mapsto m+1\) et \(\mathcal D_R\hookrightarrow\mathcal D_{R'}\).
 \item \(\tau_{\rm sp}^{\rm term}\) conserve le terminal
 \(M^{*N}S_0M^N\), qui, dans la carte miroir, se publie au moyen de \(E_R\).
    \item \(\mathcal L_{\rm sp}=(\Xi,P,B_{\rm sp}^{\rm str})\) est le lecteur final ; il s'applique
 après antécédent, cartes, mémoire, chemin et terminal.
\end{enumerate}

L'assembleur et le publicateur sont
\begin{equation}\label{eq:pdf2-rh-ensamblador-sp}
\begin{aligned}
 \mathcal A_{\rm sp}(\operatorname{pr}_{X,{\rm sp}}\widehat x)
  ={}&(\widehat X_\infty,\Gamma_9,\widehat R,
      \operatorname{Carr}_R,S_0,M,\mathcal O_9,
      \Xi,P,\mathcal D_R^{\rm cof}),\\
 \mathfrak M_{\rm sp}={}&\operatorname{Im}\mathcal A_{\rm sp},\\
 \mathcal P_{\rm sp}(\mathfrak M_{\rm sp})
  ={}&(B_{\rm sp}^{\rm str},\Xi^*(I-P)\Xi,
       \mathscr I_{\rm GW},\Delta_{\rm bind},
       \text{famille cofinale premier--gamma--polaire}),\\
 \Delta_{\rm bind}:={}&B_{\rm sp}^{\rm str}-\mathscr I_{\rm GW}.
\end{aligned}
\end{equation}
La mémoire finie de l'objet de départ satisfait
\begin{equation}\label{eq:pdf2-rh-stein-finito}
 S_0=\sum_{r=0}^{N-1}M^{*r}\mathcal O_9^*\mathcal O_9M^r
     +M^{*N}S_0M^N.
\end{equation}
Le dernier terme reste visible jusqu'à ce que son extinction soit démontrée. Par
conséquent, remplacer \(\mathfrak G_{\rm sp}\) par l'ombre publiée
\(B_{\rm sp}^{\rm str}\) change l'objet de départ ; ce n'est pas une abréviation
du même théorème. Le résiduel \(\Delta_{\rm bind}\) est conservé jusqu'à ce que la
colligation source--first démontre matériellement son annulation dans
\eqref{eq:pdf2-rh-binding-source-first}.

\section{Domaine cofinal et forme cible}

Pour $R>0$, on fixe
\[
 \mathcal D_R=C_c^\infty((-R/2,R/2)),\qquad
 \widehat\psi(t)=\int_{\mathbb R}\psi(x)e^{-itx}\,dx,
\]
et, pour $\psi,\phi\in\mathcal D_R$,
\[
 h_{\psi,\phi}(z)=\widehat\psi(z)
 \overline{\widehat\phi(\overline z)},\qquad
 \check h_{\psi,\phi}(a)=\langle\psi,T_a\phi\rangle.
\]
L'union $\bigcup_R\mathcal D_R=C_c^\infty(\mathbb R)$ est exacte. Pour tout
ensemble fini de points complexes distincts, les évaluations
$\psi\mapsto(\widehat\psi(z_1),\ldots,\widehat\psi(z_n))$ sont conjointement
surjectives lorsque $R$ est suffisamment grand. Par conséquent, cette classe cofinale
sépare les évaluations requises par le critère de Weil.

\begin{lemma}[Preuve de la séparation conjointe]
\label{lem:pdf2-rh-separacion}
L'application
\[
 \mathcal D_R\longrightarrow\mathbb C^n,
 \qquad
 \psi\longmapsto
 (\widehat\psi(z_1),\ldots,\widehat\psi(z_n))
\]
est surjective pour tout ensemble de points complexes distincts
\(z_1,\ldots,z_n\), dès que \(R\) contient un intervalle non vide.
\end{lemma}

\begin{proof}
Si les fonctionnelles d'évaluation étaient linéairement dépendantes, il existerait
\(c_1,\ldots,c_n\), non tous nuls, tels que
\[
 \int_{\mathbb R}\psi(x)
 \left(\sum_{j=1}^nc_je^{-iz_jx}\right)\,dx=0
 \qquad(\psi\in\mathcal D_R).
\]
La fonction analytique \(\sum_jc_je^{-iz_jx}\) s'annulerait dans
\((-R/2,R/2)\), donc dans tout le plan. L'indépendance des exponentielles
d'exposants distincts impose \(c_j=0\) pour tout \(j\), contradiction.
Les fonctionnelles sont indépendantes ; comme le codomaine est de dimension finie,
l'application est de rang \(\mathbb C^n\).
\end{proof}

\section{Les deux colonnes et les trois canaux calculés avant le signe}

Soit $\vartheta_L$ un régulateur commun. Avec
\[
 \ell_{p,k}=k\log p,\qquad w_{p,k}=(\log p)p^{-k/2},\qquad
 \mu_\Gamma(a)=\frac{e^{-a/2}}{1-e^{-2a}},
\]
on utilise l'inclusion uniforme $j_m$ et le canal de raccordement
$\widehat\delta_{a,m}$, qui satisfont
\[
 j_m^*j_m=I,\qquad
 \widehat\delta_{a,m}^*\widehat\delta_{b,m}
 =(I-T_a)^*(I-T_b).
\]
Les colonnes régularisées sont
\begin{align*}
 J_{+,m,R,L}\psi={}&
 \Bigl(
  (\sqrt{\vartheta_L(\ell_{p,k})w_{p,k}}\,
       \widehat\delta_{\ell_{p,k},m}\psi)_{\ell_{p,k}\leq R},\\
 &\hspace{9mm}a\mapsto
  \sqrt{\vartheta_L(a)\mu_\Gamma(a)}\,
       \widehat\delta_{a,m}\psi,\ \zeta_+b_+(\psi)
 \Bigr),\\
 J_{-,m,R,L}\psi={}&
 \Bigl(
  (\sqrt{2\vartheta_L(\ell_{p,k})w_{p,k}}\,j_m\psi)_{\ell_{p,k}\leq R},
  \sqrt{-c_\Gamma}\,j_m\psi,\ \zeta_-b_-(\psi)
 \Bigr),
\end{align*}
où $c_\Gamma=\psi_\Gamma(1/4)-\log\pi<0$ et
$b_\pm=(\widehat\psi(i/2)\pm\widehat\psi(-i/2))/\sqrt2$.

Un développement terme à terme, sans utiliser de zéros ni le signe recherché, donne
\begin{align*}
 &\langle J_{+,m,R,L}\psi,J_{+,m,R,L}\phi\rangle
 -\langle J_{-,m,R,L}\psi,J_{-,m,R,L}\phi\rangle
 =\mathscr I_{{\rm GW},L}(\psi,\phi),\\
 \mathscr I_{{\rm GW},L}(\psi,\phi)
 ={}&h_{\psi,\phi}(i/2)+h_{\psi,\phi}(-i/2)
 +c_\Gamma\check h_{\psi,\phi}(0)\\
 &+\int_0^\infty\vartheta_L(a)\mu_\Gamma(a)
 [2\check h(0)-\check h(a)-\check h(-a)]\,da\\
 &-\sum_{p,k}\vartheta_L(\ell_{p,k})w_{p,k}
 [\check h(\ell_{p,k})+\check h(-\ell_{p,k})].
\end{align*}
La troncature des deux colonnes en \(\ell_{p,k}\leq R\) est obligatoire, non
notationnelle. Si \(\psi,\phi\in\mathcal D_R\), alors
\(\check h_{\psi,\phi}(\pm\ell)=0\) pour \(\ell>R\) ; la contribution de chaque
paire de coordonnées omise à la \emph{différence} des Grams est donc nulle.
La colonne négative non tronquée aurait, en revanche, une norme infinie car
\(\sum_p(\log p)p^{-1/2}=\infty\). Dans l'ensemble fini
\(\ell_{p,k}\leq R\), les coordonnées premières convergent individuellement lorsque
\(L\to\infty\). Dans le canal gamma, l'intégrande est \(O(a)\) en zéro et
\(\mu_\Gamma(a)=O(e^{-a/2})\) à l'infini, de sorte que la convergence se fait aussi
en norme. Nous définissons donc
\begin{equation}\label{eq:pdf2-rh-columnas-reducidas}
 J_{\pm,m,R}:=\lim_{L\to\infty}J_{\pm,m,R,L}
 \quad\text{dans leurs espaces de Hilbert de coordonnées tronquées},
\end{equation}
et supprimons \(m\) lorsqu'il reste fixe. Par convergence dominée,
\[
 \mathscr I_{{\rm GW},L}\longrightarrow\mathscr I_{\rm GW}
 \qquad(L\to\infty),
\]
et les identités sont naturelles sous $m\mapsto m+1$ et $R\le R'$.

Les colonnes réduites \(J_{+,R}\) et \(J_{-,R}\) sont maintenant des opérateurs
matériels de \(\mathcal D_R\) vers Hilbert. Le développement explicite prouve, sans
consulter de zéros ni de signe,
\begin{equation}\label{eq:pdf2-rh-gram-difference-gw}
 \langle J_{+,R}f,J_{+,R}g\rangle
 -\langle J_{-,R}f,J_{-,R}g\rangle
 =\mathscr I_{\rm pr}(f,g)+\mathscr I_\Gamma(f,g)
  +\mathscr I_{\rm pol}(f,g)
 =\mathscr I_{\rm GW}(f,g).
\end{equation}
Les horloges \(\ell_{p,k}\) produisent le terme premier ; l'intégrale avec
\(\mu_\Gamma\) et \(c_\Gamma\check h(0)\), le terme gamma ; et
\[
 \langle b_+(f),b_+(g)\rangle-
 \langle b_-(f),b_-(g)\rangle
 =h_{f,g}(i/2)+h_{f,g}(-i/2)
\]
produit le terme polaire. Cette partie du calcul est matérielle et indépendante
du pont structurel.

\section{Colligation source--first de l'image spectrale--polaire}

Nous fixons d'abord un niveau nonadique \(m\) ; cet indice est supprimé dans
\(J_{\pm,R}\), \(E_R\) et \(\Sigma_R\) jusqu'à la construction du codeur. La
décomposition des codomaines explicites de
\eqref{eq:pdf2-rh-columnas-reducidas} est
\[
 I_+=\{\mathrm{pr},\Gamma,\mathrm{pol}\},\qquad
 I_-(R)=\{\Gamma0,\mathrm{pol}\}
 \sqcup\{(p,k):\ell_{p,k}\leq R\},
\]
\[
 \mathscr H_{+,R}^{\rm amb}=\bigoplus_{i\in I_+}H_{+,i,R},
 \qquad
 \mathscr H_{-,R}^{\rm amb}=\bigoplus_{j\in I_-(R)}H_{-,j,R}.
\]
Nous n'identifions pas un sous-espace atteignable à une somme de ses projections
coordonnées. Nous définissons, les inclusions ambiantes étant sous-entendues,
\[
 H_{+,R}:=\overline{\operatorname{Ran}J_{+,R}},\qquad
 H_{-,R}:=\overline{\operatorname{Ran}J_{-,R}},
\]
comme sous-espaces fermés des deux Hilbert ambiants. Le propriétaire
spectral imprimé avant ce chapitre construit
\(\mathcal H_R^{\rm sh}\), la différence antispéculaire et la ligne isométrique
source--first au moyen de
\eqref{eq:amp-rh-owner-q}--\eqref{eq:amp-rh-owner-binding}, prouve son
télescopage dans \eqref{eq:amp-rh-owner-telescopia} et publie l'extinction
terminale dans \cref{thm:amp-rh-cierre-cofinal-propietario}. La réalisation
qui suit conserve cette construction, avant l'identification classique,
dans la ligne
\begin{equation}\label{eq:pdf2-rh-coligacion-source-first}
 \Sigma_R=
 \begin{pmatrix}\mathcal K_R^{\rm src}\\ \mathcal F_R^{\rm src}\end{pmatrix}:
 H_{+,R}\longrightarrow
 H_{-,R}\oplus\mathcal H_R^{\rm sh},
 \qquad
 \Sigma_R^*\Sigma_R
 = (\mathcal K_R^{\rm src})^*\mathcal K_R^{\rm src}
  +(\mathcal F_R^{\rm src})^*\mathcal F_R^{\rm src}=I_{H_{+,R}}.
\end{equation}
Son action atteignable est fixée sur la colonne positive complète par
les relations des deux publications qui font partie de
\(\operatorname{Rel}_{\rm sp}\). Pour éviter toute substitution de lecteur,
nous notons
\((\mathcal X_R^{\rm str},\Xi_R^{\rm str},P_R^{\rm str})\) la restriction à
\(\mathcal D_R\) du lecteur structurel de
\eqref{eq:pdf2-rh-res-sp-completa}. Ses deux moments sont, littéralement,
\begin{equation}\label{eq:pdf2-rh-momentos-lector-estructural}
\begin{aligned}
 \langle\Xi_R^{\rm str}f,\Xi_R^{\rm str}g\rangle
  &=\langle J_{+,R}f,J_{+,R}g\rangle,\\
 \langle P_R^{\rm str}\Xi_R^{\rm str}f,
          P_R^{\rm str}\Xi_R^{\rm str}g\rangle
  &=\langle J_{-,R}f,J_{-,R}g\rangle,
 \qquad (P_R^{\rm str})^*= (P_R^{\rm str})^2=P_R^{\rm str}.
\end{aligned}
\end{equation}
La coordonnée de frontière est accompagnée de l'isométrie structurelle
\[
 \jmath_R^{\rm sh}:
 \overline{\operatorname{Ran}((I-P_R^{\rm str})\Xi_R^{\rm str})}
 \longrightarrow\mathcal H_R^{\rm sh}.
\]
Par définition de la ligne source--first, non par une identification ultérieure,
\begin{equation}\label{eq:pdf2-rh-source-state-output}
 E_R:=\mathcal F_R^{\rm src}J_{+,R}
     =\jmath_R^{\rm sh}(I-P_R^{\rm str})\Xi_R^{\rm str}:
 \mathcal D_R\longrightarrow\mathcal H_R^{\rm sh},
 \qquad
 \Omega_{0,R}:=\mathcal K_R^{\rm src}J_{+,R}-J_{-,R}.
\end{equation}
Il n'y a pas une colligation par canal : la même ligne reçoit simultanément les
horloges premier--puissance, l'intégrale gamma, le mode scalaire et les deux lignes
polaires. En particulier, elle conserve tous les produits croisés qui apparaissent en
polarisant \eqref{eq:pdf2-rh-gram-difference-gw}.
Avec les projections ambiantes \(\pi_j^-:\mathscr H_{-,R}^{\rm amb}
\to H_{-,j,R}\), les lignes de sortie correctement typées sont
\begin{equation}\label{eq:pdf2-rh-bloques-coligacion}
 \mathcal K_{j,R}^{\rm src}
 :=\pi_j^-\mathcal K_R^{\rm src}:
 H_{+,R}\longrightarrow H_{-,j,R},
 \qquad
 \mathcal F_R^{\rm src}:H_{+,R}\longrightarrow\mathcal H_R^{\rm sh}.
\end{equation}
Ainsi, on n'applique pas \(\mathcal K_R^{\rm src}\) à une coordonnée positive qui
pourrait ne pas appartenir à elle seule à \(H_{+,R}\). Tous les croisements entre
canaux restent dans
\((\mathcal K_R^{\rm src})^*\mathcal K_R^{\rm src}\) et
\((\mathcal F_R^{\rm src})^*\mathcal F_R^{\rm src}\) ; ils ne sont pas orthogonalisés
avant l'identité d'énergie.

La coordonnée de survie construite par
\eqref{eq:amp-rh-owner-q}, \eqref{eq:amp-rh-owner-telescopia} et
\eqref{eq:amp-rh-owner-terminal} produit les espaces résiduels
\(\mathcal R_{n,R}^{\rm sh}\), les continuations
\(\Omega_{n,R}:\mathcal D_R\to\mathcal R_{n,R}^{\rm sh}\) et les isométries
antispéculaires
\(V_{n,R}:\mathcal R_{n+1,R}^{\rm sh}\to\mathcal R_{n,R}^{\rm sh}\). Son
premier espace et son premier opérateur sont, par définition typée,
\[
 \mathcal R_{0,R}^{\rm sh}:=H_{-,R},\qquad
 \Omega_{0,R}:=\mathcal K_R^{\rm src}J_{+,R}-J_{-,R}.
\]
Sa récurrence d'un tour complet est la publication cofinale du même
opérateur construit, non une donnée ajoutée à la réalisation :
\begin{equation}\label{eq:pdf2-rh-recurrencia-residual}
 q=\frac59,\qquad
 \Omega_{n,R}=qV_{n,R}\Omega_{n+1,R},
 \qquad V_{n,R}^*V_{n,R}=I,
 \qquad
 \sup_{n\geq0}\|\Omega_{n,R}f\|<\infty
 \quad(f\in\mathcal D_R).
\end{equation}
Cette récurrence appartient à la colligation enrichie : elle transporte la
différence de sortie, non la forme de Weil ni son signe.

\begin{proposition}[Annulation coinductive de la différence de sortie]
\label{prop:pdf2-rh-omega-cero}
Pour tout \(R>0\) et tout \(f\in\mathcal D_R\),
\begin{equation}\label{eq:pdf2-rh-output-binding-source}
 \Omega_{n,R}f=0\quad(n\geq0),
 \qquad
 \boxed{\mathcal K_R^{\rm src}J_{+,R}=J_{-,R}}.
\end{equation}
\end{proposition}

\begin{proof}
Itérer \eqref{eq:pdf2-rh-recurrencia-residual} \(N\) fois et utiliser le fait que chaque
\(V_{n,R}\) est isométrique donne
\[
 \|\Omega_{n,R}f\|
 =q^N\|\Omega_{n+N,R}f\|
 \leq q^N\sup_{j\geq n}\|\Omega_{j,R}f\|.
\]
Le membre droit converge vers zéro car \(0<q=5/9<1\). Par conséquent,
\(\Omega_{n,R}f=0\) ; le cas \(n=0\) et
\eqref{eq:pdf2-rh-source-state-output} produisent l'identité encadrée.
\end{proof}

L'action ligne par ligne de la sortie n'est plus implicite. Les projections sont
les restrictions à \(H_{-,R}\) des projections du Hilbert ambiant, et
\begin{equation}\label{eq:pdf2-rh-salidas-generador-a-generador}
\begin{aligned}
 \mathcal K_{\Gamma0,R}^{\rm src}J_{+,R}f
  &=\sqrt{-c_\Gamma}\,j_mf,\\
 \mathcal K_{(p,k),R}^{\rm src}J_{+,R}f
  &=\sqrt{2w_{p,k}}\,j_mf
  &&(\ell_{p,k}\leq R),\\
 \mathcal K_{{\rm pol},R}^{\rm src}J_{+,R}f
  &=\zeta_-b_-(f).
\end{aligned}
\end{equation}
Ces lignes épuisent \(J_{-,R}\), tandis que
\begin{equation}\label{eq:pdf2-rh-frontera-generador-a-generador}
 \mathcal F_R^{\rm src}J_{+,R}f=E_Rf
\end{equation}
conserve, avec tous ses croisements, la coordonnée de frontière.
Évaluer \eqref{eq:pdf2-rh-coligacion-source-first} en \(J_{+,R}f\) et
\(J_{+,R}g\), et utiliser \eqref{eq:pdf2-rh-output-binding-source}, produit
l'identité d'énergie antérieure au signe
\begin{equation}\label{eq:pdf2-rh-energia-source-first}
 \boxed{
 \langle J_{+,R}f,J_{+,R}g\rangle
 =\langle J_{-,R}f,J_{-,R}g\rangle
  +\langle E_Rf,E_Rg\rangle.}
\end{equation}
En la comparant au calcul indépendant
\eqref{eq:pdf2-rh-gram-difference-gw}, on obtient, sans introduire une racine de
la forme cible,
\begin{equation}\label{eq:pdf2-rh-binding-source-first}
 \boxed{E_R^*E_R=\mathscr I_{\rm GW}\big|_{\mathcal D_R}.}
\end{equation}
En particulier,
\(\Delta_{{\rm bind},R}
:=E_R^*E_R-\mathscr I_{\rm GW}|_{\mathcal D_R}=0\) ; cette annulation est la
sortie de \eqref{eq:pdf2-rh-coligacion-source-first}--
\eqref{eq:pdf2-rh-energia-source-first}, non une prémisse insérée dans le
publicateur.

\section{Complément structurel, feuillet miroir et télescopage terminal}

L'état \(E_R=\mathcal F_R^{\rm src}J_{+,R}\) construit par la ligne
source--first est la coordonnée de frontière de
\(\mathfrak G_{\rm sp}\). Nous définissons sa forme structurelle par
\begin{equation}\label{eq:pdf2-rh-bw-estructural}
 \mathcal B_{{\rm sp},R}^{\rm str}
 :=E_R^*E_R
 =\mathscr I_{\rm GW}\big|_{\mathcal D_R},
\end{equation}
où la deuxième égalité est
\eqref{eq:pdf2-rh-binding-source-first}, déjà déduite de la colligation et du
développement générateur par générateur.

Soient \(P_+^{\rm sh}\) et \(P_-^{\rm sh}\) les projecteurs orthogonaux
complémentaires des deux feuillets. On fixe
\begin{equation}\label{eq:pdf2-rh-q-a-d}
 q=\frac59,\qquad
 A=P_+^{\rm sh}+qP_-^{\rm sh},\qquad
 D=\sqrt{1-q^2}\,P_-^{\rm sh}.
\end{equation}
L'orthogonalité donne l'identité exacte
\begin{equation}\label{eq:pdf2-rh-conservacion-local}
 A^*A+D^*D
 =P_+^{\rm sh}+q^2P_-^{\rm sh}
 +(1-q^2)P_-^{\rm sh}=I.
\end{equation}
Les neuf phases coordonnent un seul tour avec mémoire :
\begin{equation}\label{eq:pdf2-rh-monodromia}
 M_9=\Phi A_8\cdots A_0=qV_9,
 \qquad V_9^*V_9=I.
\end{equation}
La contraction par tour est \(q\), non \(q^9\) ; les phases internes conservent
l'ordre et la retenue.

La colligation type
\(E_R:\mathcal D_R\to\operatorname{Ran}P_-^{\rm sh}\). Par conséquent,
\begin{equation}\label{eq:pdf2-rh-binding-espejo}
 P_+^{\rm sh}E_R=0,
 \qquad AE_R=qE_R,
 \qquad
 E_R^*E_R=\mathcal B_{{\rm sp},R}^{\rm str}
 =\mathscr I_{\rm GW}\big|_{\mathcal D_R}.
\end{equation}

\begin{proposition}[Télescopage fini avec terminal conservé]
\label{prop:pdf2-rh-telescopia}
Pour chaque \(N\geq1\),
\begin{equation}\label{eq:pdf2-rh-telescopia-finita}
 \mathcal B_{{\rm sp},R}^{\rm str}
 =\sum_{n=0}^{N-1}(DA^nE_R)^*(DA^nE_R)
 +(A^NE_R)^*(A^NE_R).
\end{equation}
De manière équivalente, évaluée en \(f,g\in\mathcal D_R\), l'identité finie de
la colligation est
\begin{equation}\label{eq:pdf2-rh-balance-finito-columnas}
\boxed{
\begin{aligned}
 &\langle J_{+,R}f,J_{+,R}g\rangle
  -\langle J_{-,R}f,J_{-,R}g\rangle\\
 &\quad=\sum_{n=0}^{N-1}
  \langle DA^nE_Rf,DA^nE_Rg\rangle
  +\langle A^NE_Rf,A^NE_Rg\rangle.
\end{aligned}}
\end{equation}
Le dernier terme satisfait
\begin{equation}\label{eq:pdf2-rh-terminal-extincion}
 (A^NE_R)^*(A^NE_R)
 =q^{2N}\mathcal B_{{\rm sp},R}^{\rm str}
 \xrightarrow[N\to\infty]{}0.
\end{equation}
\end{proposition}

\begin{proof}
En appliquant \eqref{eq:pdf2-rh-conservacion-local} à \(A^nE_R\),
\[
 (A^nE_R)^*(A^nE_R)
 =(DA^nE_R)^*(DA^nE_R)
 +(A^{n+1}E_R)^*(A^{n+1}E_R).
\]
La somme de \(n=0\) à \(N-1\) annule uniquement les termes intermédiaires et
conserve \((A^NE_R)^*(A^NE_R)\). Par
\eqref{eq:pdf2-rh-binding-espejo}, \(A^NE_R=q^NE_R\) ; le terminal
est donc \(q^{2N}E_R^*E_R=q^{2N}\mathcal B_{{\rm sp},R}^{\rm str}\). Comme
\(0<q=5/9<1\), il converge géométriquement vers zéro.
\end{proof}

Ce n'est qu'après avoir démontré \eqref{eq:pdf2-rh-terminal-extincion} que l'on définit
\begin{equation}\label{eq:pdf2-rh-l-r}
 \mathscr L_Rf=(DA^nE_Rf)_{n\geq0}.
\end{equation}
La limite du télescopage fini donne
\begin{equation}\label{eq:pdf2-rh-lstar-l}
 \boxed{
 \mathscr L_R^*\mathscr L_R
 =\mathcal B_{{\rm sp},R}^{\rm str}
 =\mathscr I_{\rm GW}\big|_{\mathcal D_R}\succeq0.}
\end{equation}

\section{Codeur commun, projection et double publication explicites}

Fixons \(m\geq1\) et deux vecteurs orthonormaux
\(u_m,\eta_m\in\mathbb C^{9^m}\), où \(u_m\) est le mode uniforme et
\(\eta_m\) le mode de retenue. L'espace et la projection du codeur sont
\begin{equation}\label{eq:pdf2-rh-hilbert-comun}
 \mathcal X_{m,R}^{\rm sp}
 :=\mathbb C^{9^m}\widehat\otimes
   (H_{-,R}\oplus\mathcal H_R^{\rm sh}),
 \qquad
 P_{m,R}:=|u_m\rangle\langle u_m|\otimes
 I_{H_{-,R}\oplus\mathcal H_R^{\rm sh}},
 \qquad Q_{m,R}:=I-P_{m,R}.
\end{equation}
L'action complète du codeur des cinq canaux est
\begin{equation}\label{eq:pdf2-rh-xi-comun-explicita}
 \boxed{
 \begin{aligned}
  \mathfrak C_{m,R}f
  &:=u_m\otimes(J_{-,R}f,0)
    +\eta_m\otimes(0,E_Rf),\\
  \Xi_{m,R}f&:=\mathfrak C_{m,R}f.
 \end{aligned}}
\end{equation}
On ne prend pas une racine de \(\mathscr I_{\rm GW}\) : les deux entrées sont la
sortie et l'état de la colligation source--first déjà définis dans
\eqref{eq:pdf2-rh-source-state-output}.

\begin{proposition}[Les deux moments s'obtiennent à partir de l'action du codeur]
\label{prop:pdf2-rh-vmas-isometria}
Pour tout \(f,g\in\mathcal D_R\),
\begin{equation}\label{eq:pdf2-rh-dos-momentos-construidos}
\begin{aligned}
 \langle\mathfrak C_{m,R}f,\mathfrak C_{m,R}g\rangle
  &=\langle J_{+,R}f,J_{+,R}g\rangle,\\
 \langle P_{m,R}\mathfrak C_{m,R}f,
         P_{m,R}\mathfrak C_{m,R}g\rangle
  &=\langle J_{-,R}f,J_{-,R}g\rangle,\\
 \langle Q_{m,R}\Xi_{m,R}f,Q_{m,R}\Xi_{m,R}g\rangle
  &=\mathscr I_{\rm GW}(f,g).
\end{aligned}
\end{equation}
\end{proposition}

\begin{proof}
L'orthogonalité \(u_m\perp\eta_m\) et
\eqref{eq:pdf2-rh-energia-source-first} donnent
\[
 \langle\mathfrak C_{m,R}f,\mathfrak C_{m,R}g\rangle
 =\langle J_{-,R}f,J_{-,R}g\rangle
  +\langle E_Rf,E_Rg\rangle
 =\langle J_{+,R}f,J_{+,R}g\rangle.
\]
La projection \(P_{m,R}\) conserve le terme uniforme et élimine celui de
retenue ; \(Q_{m,R}\) fait l'inverse. Les deux autres identités découlent de
\eqref{eq:pdf2-rh-binding-source-first}.
\end{proof}

Le codeur qui vient d'être écrit est un représentant explicite du lecteur
structurel initial, non un lecteur de substitution. Sur les sous-espaces atteignables, on
définit
\begin{equation}\label{eq:pdf2-rh-isomorfismo-lector-estructural}
\begin{aligned}
 \mathcal U_{m,R}^{\rm str}
  (P_R^{\rm str}\Xi_R^{\rm str}f)
   &:=u_m\otimes(J_{-,R}f,0),\\
 \mathcal U_{m,R}^{\rm str}
  ((I-P_R^{\rm str})\Xi_R^{\rm str}f)
   &:=\eta_m\otimes(0,E_Rf).
\end{aligned}
\end{equation}
Les deux sources sont orthogonales. La première égalité de Gram de chaque source
est \eqref{eq:pdf2-rh-momentos-lector-estructural} ; pour la deuxième, on utilise
\(E_R=\jmath_R^{\rm sh}(I-P_R^{\rm str})\Xi_R^{\rm str}\). Ainsi, les deux
règles sont bien définies, sont isométriques et s'étendent à la somme orthogonale
de leurs adhérences. Dans cette adhérence atteignable, elles satisfont
\begin{equation}\label{eq:pdf2-rh-entrelazamiento-lector-estructural}
 \mathcal U_{m,R}^{\rm str}\Xi_R^{\rm str}=\Xi_{m,R},
 \qquad
 \mathcal U_{m,R}^{\rm str}P_R^{\rm str}
 =P_{m,R}\mathcal U_{m,R}^{\rm str},
 \qquad
 (\Xi_R^{\rm str})^*(I-P_R^{\rm str})\Xi_R^{\rm str}=E_R^*E_R.
\end{equation}
Ainsi, \(B_{{\rm sp},R}^{\rm str}\), le complément de la structure de
départ, est exactement le complément du codeur construit ; on n'a pas
changé d'objet en passant de \eqref{eq:pdf2-rh-res-sp-completa} à
\eqref{eq:pdf2-rh-xi-comun-explicita}.

La règle
\begin{equation}\label{eq:pdf2-rh-vmas-rango}
 V_{+,m,R}^{0}(J_{+,R}f):=\Xi_{m,R}f
\end{equation}
est bien définie et isométrique par la première égalité de
\eqref{eq:pdf2-rh-dos-momentos-construidos} ; elle s'étend de manière unique en
une isométrie
\(V_{+,m,R}:H_{+,R}\to\mathcal X_{m,R}^{\rm sp}\). La publication négative
est l'isométrie
\begin{equation}\label{eq:pdf2-rh-vmenos}
 V_{-,m,R}:H_{-,R}\longrightarrow\mathcal X_{m,R}^{\rm sp},
 \qquad V_{-,m,R}y=u_m\otimes(y,0).
\end{equation}
Par construction,
\begin{equation}\label{eq:pdf2-rh-doble-publicacion-material}
 \boxed{
 V_{+,m,R}J_{+,R}=\Xi_{m,R},
 \qquad
 V_{-,m,R}J_{-,R}=P_{m,R}\Xi_{m,R}.}
\end{equation}
Le complément conserve exactement l'état antispéculaire et sa colonne
d'observation :
\begin{equation}\label{eq:pdf2-rh-complemento-comun}
 \boxed{
 \Xi_{m,R}^*Q_{m,R}\Xi_{m,R}
 =E_R^*E_R
 =\mathscr L_R^*\mathscr L_R
 =\mathscr I_{\rm GW}\big|_{\mathcal D_R}.}
\end{equation}

\begin{proposition}[Connecteur induit et coïncidence source--first]
\label{prop:pdf2-rh-conector}
L'opérateur
\begin{equation}\label{eq:pdf2-rh-k-r}
 \mathcal K_R
 :=V_{-,m,R}^*P_{m,R}V_{+,m,R}:
 H_{+,R}\longrightarrow H_{-,R}
\end{equation}
est contractif, coïncide avec \(\mathcal K_R^{\rm src}\) sur tout
\(H_{+,R}\) et satisfait
\begin{equation}\label{eq:pdf2-rh-k-binding}
 \|\mathcal K_R\|\leq1,
 \qquad
 \mathcal K_RJ_{+,R}=J_{-,R}.
\end{equation}
\end{proposition}

\begin{proof}
Les deux cartes sont des isométries et \(P_{m,R}\) est une projection orthogonale,
donc \(\|\mathcal K_R\|\leq1\). Sur l'image dense de \(J_{+,R}\),
\[
 \mathcal K_RJ_{+,R}f
 =V_{-,m,R}^*P_{m,R}\Xi_{m,R}f
 =V_{-,m,R}^*V_{-,m,R}J_{-,R}f
 =J_{-,R}f.
\]
La même identité vaut pour \(\mathcal K_R^{\rm src}\) par
\eqref{eq:pdf2-rh-output-binding-source} ; la continuité et la densité
prouvent que les deux opérateurs coïncident.
\end{proof}

La différence de Gram est ainsi publiée dans toutes ses réalisations
coïncidentes :
\begin{equation}\label{eq:pdf2-rh-tres-publicaciones}
\boxed{
\begin{aligned}
 J_{+,R}^*J_{+,R}-J_{-,R}^*J_{-,R}
 &=\Xi_{m,R}^*Q_{m,R}\Xi_{m,R}\\
 &=E_R^*E_R\\
 &=\mathscr L_R^*\mathscr L_R\\
 &=\mathcal B_{{\rm sp},R}^{\rm str}=\mathscr I_{\rm GW}.
\end{aligned}}
\end{equation}

\section{Naturalité cofinale et critère réversible de Weil}

Pour \(R\leq R'\), soit
\(i_{R,R'}:\mathcal D_R\hookrightarrow\mathcal D_{R'}\). Les colonnes
individuelles ne forment pas un système isométrique : lorsqu'une nouvelle longueur
\(R<\ell_{p,k}\leq R'\) apparaît, les deux acquièrent la même énergie positive. Pour
\(f,g\in\mathcal D_R\), la séparation des supports donne exactement
\begin{equation}\label{eq:pdf2-rh-naturalidad}
 \left\langle(I-T_{\ell_{p,k}})f,
                    (I-T_{\ell_{p,k}})g\right\rangle
 =2\langle f,g\rangle
 =\left\langle\sqrt2j_mf,\sqrt2j_mg\right\rangle .
\end{equation}
Par conséquent, l'incrément premier de
\(J_{+,R'}^*J_{+,R'}\) coïncide avec celui de
\(J_{-,R'}^*J_{-,R'}\) et s'annule dans la différence, mais on ne déclare pas une
fausse isométrie entre les colonnes. La compatibilité cofinale correcte est celle
des formes :
\begin{equation}\label{eq:pdf2-rh-naturalidad-p}
 \boxed{
 \mathscr I_{{\rm GW},R}
 =i_{R,R'}^*\mathscr I_{{\rm GW},R'}i_{R,R'},
 \qquad
 \mathcal B_{{\rm sp},R}^{\rm str}
 =i_{R,R'}^*\mathcal B_{{\rm sp},R'}^{\rm str}i_{R,R'}.}
\end{equation}
Comme \(E_R^*E_R=\mathscr L_R^*\mathscr L_R=\mathscr I_{{\rm GW},R}\), les
réalisations minimales de Kolmogorov admettent bien des isométries uniques sur leurs
images atteignables, définies par
\begin{equation}\label{eq:pdf2-rh-naturalidad-bw}
 \iota_{R,R'}^E(E_Rf):=E_{R'}f,
 \qquad
 \iota_{R,R'}^L(\mathscr L_Rf):=\mathscr L_{R'}f
 \quad(f\in\mathcal D_R).
\end{equation}
Ces deux règles sont bien définies précisément par
\eqref{eq:pdf2-rh-naturalidad-p} ; elles ne s'étendent pas par décret à
\(J_{\pm,R}\), \(\Xi_{m,R}\) ou \(\mathcal K_R\). Le raffinement nonadique
\(m\mapsto m+1\) reste, à fenêtre fixe, dans les isométries structurelles
qui transportent \(j_m\), \(\widehat\delta_{a,m}\), \(u_m\) et \(\eta_m\).

Pour ne pas comprimer la dernière étape, nous fixons aussi ses objets et
normalisations. Soit
\[
 \xi(s):=\frac12s(s-1)\pi^{-s/2}\Gamma(s/2)\zeta(s)
\]
et soit \(\mathscr Z\) le multiensemble de ses zéros, chaque zéro étant répété selon
sa multiplicité. L'équation fonctionnelle et la réalité des coefficients
préservent \(\mathscr Z\) par
\(\rho\mapsto1-\rho\), \(\rho\mapsto\overline\rho\) et
\(\rho\mapsto1-\overline\rho\). Avec notre convention de Fourier, on définit
\begin{equation}\label{eq:pdf2-rh-transformada-weil}
 \Phi_f(s):=\widehat f\!\left(i(s-\tfrac12)\right),
 \qquad
 \widetilde f(x):=\overline{f(-x)}.
\end{equation}
La formule explicite complète calculée dans
\eqref{eq:pdf2-rh-gram-difference-gw}, y compris les termes gamma et polaire
qui retirent les pôles et les zéros triviaux, a la forme spectrale
\begin{equation}\label{eq:pdf2-rh-formula-espectral-weil}
 \mathscr I_{\rm GW}(f,g)
 =\sum_{\rho\in\mathscr Z}
   \Phi_f(\rho)\,
   \overline{\Phi_g(1-\overline\rho)}.
\end{equation}
La somme compte les multiplicités. Elle converge absolument : pour
\(f,g\in C_c^\infty(\mathbb R)\), Paley--Wiener donne une décroissance plus rapide
que toute puissance dans les bandes verticales, tandis que le nombre de zéros avec
\(|\operatorname{Im}\rho|\leq T\) est \(O(T\log T)\).

\begin{proposition}[Critère réversible de Weil avec quantificateurs complets]
\label{prop:pdf2-rh-criterio-weil-desplegado}
Sous la normalisation
\eqref{eq:pdf2-rh-transformada-weil} et en comptant tous les zéros non triviaux
avec leur multiplicité,
\begin{equation}\label{eq:pdf2-rh-criterio-weil}
 \left[
  \mathscr I_{\rm GW}(f,f)\geq0
  \text{ pour tout }f\in C_c^\infty(\mathbb R;\mathbb C)
 \right]
 \Longleftrightarrow
 \left[
  \operatorname{Re}\rho=\frac12
  \text{ pour tout }\rho\in\mathscr Z
 \right].
\end{equation}
\end{proposition}

L'implication réciproque ne s'obtient pas en interpolant un nombre croissant de
zéros : cette opération ne donnerait pas de contrôle uniforme sur la queue. On utilise le
critère classique complet de Weil dans la formulation de Bombieri, dont l'espace
de test et la normalisation sont transposés ici sans les abréger.\footnote{%
E.~Bombieri, \emph{Remarks on Weil's quadratic functional in the theory of
prime numbers, I}, Rendiconti Lincei, Matematica e Applicazioni, ser.~9,
vol.~11 (2000), pp.~183--233. Le critère qui y est fixé affirme que la fonctionnelle
quadratique de la formule explicite est semi-définie positive sur tout
\(C_c^\infty(0,\infty)\) si et seulement si l'hypothèse de Riemann est vraie.}

\begin{proof}
Nous écrivons matériellement le changement de coordonnées qui identifie les deux
critères. Pour \(f\in C_c^\infty(\mathbb R)\), soit
\begin{equation}\label{eq:pdf2-rh-cambio-mellin}
 (\mathcal M f)(x):=x^{-1/2}f(\log x),
 \qquad x>0.
\end{equation}
L’application
\[
 \mathcal M:C_c^\infty(\mathbb R)
 \longrightarrow C_c^\infty(0,\infty)
\]
est bijectif, avec pour inverse
\begin{equation}\label{eq:pdf2-rh-cambio-mellin-inverso}
 (\mathcal M^{-1}g)(u)=e^{u/2}g(e^u).
\end{equation}
Si
\(\widehat g_{\rm Mel}(s):=\int_0^\infty g(x)x^{s-1}\,dx\), le changement
\(x=e^u\) donne, sans facteur résiduel,
\begin{equation}\label{eq:pdf2-rh-mellin-fourier}
 \begin{aligned}
 \widehat{\mathcal M f}_{\rm Mel}(s)
 &=\int_{\mathbb R}f(u)e^{(s-1/2)u}\,du\\
 &=\widehat f\!\left(i(s-\tfrac12)\right)
 =\Phi_f(s).
 \end{aligned}
\end{equation}
De plus,
\begin{equation}\label{eq:pdf2-rh-mellin-involucion}
 \widehat{\overline{\mathcal M f}}_{\rm Mel}(1-\rho)
 =\overline{
   \widehat{\mathcal M f}_{\rm Mel}(1-\overline\rho)}
 =\overline{\Phi_f(1-\overline\rho)}.
\end{equation}
Par conséquent, la fonctionnelle du critère de Weil--Bombieri est exactement
\begin{equation}\label{eq:pdf2-rh-bombieri-identificado}
 \sum_{\rho\in\mathscr Z}
 \widehat{\mathcal M f}_{\rm Mel}(\rho)\,
 \widehat{\overline{\mathcal M f}}_{\rm Mel}(1-\rho)
 =
 \sum_{\rho\in\mathscr Z}
 \Phi_f(\rho)\overline{\Phi_f(1-\overline\rho)}
 =\mathscr I_{\rm GW}(f,f).
\end{equation}
La classe des tests n'est pas modifiée, aucune multiplicité n'est perdue et
l'ensemble des zéros n'est pas tronqué, car \(\mathcal M\) est une bijection. Le
critère classique cité appliqué à
\eqref{eq:pdf2-rh-bombieri-identificado} donne précisément l'équivalence
\eqref{eq:pdf2-rh-criterio-weil}. Dans le sens direct, si
\(\operatorname{Re}\rho=1/2\), la même identité se réduit explicitement à
\[
 \mathscr I_{\rm GW}(f,f)
 =\sum_{\rho\in\mathscr Z}|\Phi_f(\rho)|^2\geq0.
\]
\end{proof}

L'égalité \(\bigcup_R\mathcal D_R=C_c^\infty(\mathbb R)\) est exacte ; chaque
fonction employée dans l'argument appartient à une fenêtre et la naturalité
cofinale y transporte la positivité déjà démontrée.

\begin{theorem}[Résolution de Riemann--Weil à partir de l'image spectrale--polaire]
\label{thm:pdf2-rh-resolucion}
La structure complète
\(\mathfrak G_{\rm sp}\) étant fixée, la forme
complète de Guinand--Weil est positive sur la classe cofinale et séparatrice. En
conséquence, tout zéro non trivial \(\rho\) de la fonction zêta de Riemann
satisfait
\[
 \boxed{\operatorname{Re}\rho=\frac12.}
\]
\end{theorem}

\begin{proof}
Le développement premier--gamma--polaire identifie la différence des deux colonnes
à \(\mathscr I_{\rm GW}\) en utilisant un unique régulateur et le retire par
convergence dominée. La récurrence antispéculaire
\eqref{eq:pdf2-rh-recurrencia-residual} annule la différence de sortie avant de
former une inégalité ; l'isométrie \(\Sigma_R\) donne donc
\eqref{eq:pdf2-rh-energia-source-first}. Comparer cette identité au
développement calculé produit
\(E_R^*E_R=\mathscr I_{\rm GW}\). L'identité locale
\(A^*A+D^*D=I\) est itérée sans effacer le terminal et donne
\eqref{eq:pdf2-rh-telescopia-finita} ; le terminal est exactement
\(q^{2N}E_R^*E_R\) et ne s'éteint qu'à la limite. D'où
\(\mathscr L_R^*\mathscr L_R=\mathscr I_{\rm GW}\succeq0\).
L'action
\eqref{eq:pdf2-rh-xi-comun-explicita} construit alors le codeur, ses
deux moments, la projection et les deux publications, et
\eqref{eq:pdf2-rh-tres-publicaciones} identifie toutes les réalisations du
même carré. La naturalité transporte cette identité à toute l'union
cofinale. Enfin,
\eqref{eq:pdf2-rh-criterio-weil} conclut l'affirmation.
\end{proof}

\section{Contrôle de Redheffer local ultérieur et non directeur}

Le déploiement suivant vérifie que la passivité locale peut être obtenue à partir de
la retenue sans utiliser la forme de Weil. C'est un contrôle de non-circularité, non une
prémisse du \cref{thm:pdf2-rh-resolucion}. Ce transfert tronqué n'est pas
la colligation complète \(\Sigma_R\) de
\eqref{eq:pdf2-rh-coligacion-source-first}.

Pour \(N=9^m\), soit \(U_{a,N}\) l'odomètre unitaire
\[
 U_{a,N}(e_j\otimes f)=e_{j+1}\otimes f\ (j<N-1),
 \qquad U_{a,N}(e_{N-1}\otimes f)=e_0\otimes T_af.
\]
Avec \(P_N=|u_N\rangle\langle u_N|\otimes I\), \(Q_N=I-P_N\),
\(A_{a,N}=P_NU_{a,N}P_N\) et \(C_{a,N}=Q_NU_{a,N}P_N\),
\begin{equation}\label{eq:pdf2-rh-source-energia}
 A_{a,N}^*A_{a,N}+C_{a,N}^*C_{a,N}=P_N,
 \qquad
 C_{a,N}^*C_{a,N}
 =\frac{N-1}{N^2}(2I-T_a-T_a^*).
\end{equation}
Pour \(0<s<1\), le transfert
\begin{equation}\label{eq:pdf2-rh-source-redheffer}
 \mathscr R_{a,N}(s)=(sI-A_{a,N})(I-sA_{a,N})^{-1}
\end{equation}
satisfait
\begin{equation}\label{eq:pdf2-rh-source-defecto}
\begin{aligned}
 I-\mathscr R_{a,N}(s)^*\mathscr R_{a,N}(s)
 ={}&(1-s^2)(I-sA_{a,N}^*)^{-1}\\
 &\quad\cdot C_{a,N}^*C_{a,N}(I-sA_{a,N})^{-1}\succeq0.
\end{aligned}
\end{equation}
Avec
\[
 s_{p,N}=\frac{Np^{-1/2}}{1+(N-1)p^{-1/2}},
 \qquad s_\Gamma(a)=e^{-a/2},
\]
la somme sur les horloges premières et l'intégrale directe gamma conservent une norme au
plus égale à un. Explicitement,
\begin{equation}\label{eq:pdf2-rh-source-relojes}
 \mathscr R_{N,R}
 =\bigoplus_{\log p\leq R}
   \mathscr R_{\log p,N}(s_{p,N})
 \ \oplus\!
 \int_0^\infty{}^\oplus
   \mathscr R_{a,N}(s_\Gamma(a))\,da,
 \qquad \|\mathscr R_{N,R}\|\leq1.
\end{equation}
La carte finie des rôles
\[
 \mathcal H_+=\operatorname{span}\{f_3,f_6,f_9\},\quad
 \mathcal H_-=\operatorname{span}\{f_4,f_8\},\quad
 C_\angle=|f_4\rangle\langle f_3|+|f_8\rangle\langle f_9|
\]
satisfait \(I-C_\angle^*C_\angle=P_6\). Avec les cartes énergétiques
two-way, on fixe
\begin{equation}\label{eq:pdf2-rh-source-cartas}
\begin{aligned}
 \mathcal K_q&=K_2\otimes(\mathbb C^9)^{\otimes q},\\
 \mathsf B_{\pm,q}:\mathcal M_{\pm,q,R}
 &\longrightarrow
 \mathcal K_q\widehat\otimes\mathcal H_\pm
 \widehat\otimes\mathcal H_{\mathfrak C_R},\\
 \mathsf B_{\pm,q}^*\mathsf B_{\pm,q}&=\mathsf G_{\pm,q},
 \qquad
 \mathsf B_{\pm,q}^{\sharp}
 =\mathsf G_{\pm,q}^{-1}\mathsf B_{\pm,q}^*.
\end{aligned}
\end{equation}
Le transfert
\begin{equation}\label{eq:pdf2-rh-source-transferencia}
 \mathscr C_{q,R}^{\rm src}
 =\mathsf B_{-,q}^{\sharp}
  (I\otimes C_\angle\otimes\mathscr R_{N,R})\mathsf B_{+,q}
\end{equation}
a pour défaut
\begin{align}
 \Delta_{q,R}^{\rm src}
 ={}&I\otimes\bigl[P_6\otimes I
 +(P_3+P_9)\otimes
 (I-\mathscr R_{N,R}^*\mathscr R_{N,R})\bigr]\succeq0,
 \label{eq:pdf2-rh-source-delta}\\
 I-(\mathscr C_{q,R}^{\rm src})^\sharp\mathscr C_{q,R}^{\rm src}
 ={}&\mathsf B_{+,q}^{\sharp}\Delta_{q,R}^{\rm src}
 \mathsf B_{+,q}\succeq0.
 \label{eq:pdf2-rh-source-contractividad}
\end{align}
Ce contrôle conserve chemin, horloge, rôle, retenue et canaux avant de former
une différence. Il n'est pas identifié à \(\mathcal K_R\) et ne remplace pas la
publication commune explicite de
\eqref{eq:pdf2-rh-doble-publicacion-material} ; sa portée locale ne
rouvre donc pas la conclusion cofinale.

\section{Réfutateurs d'une ombre amputée}

\begin{falsifier}[Substitutions qui changent l'objet de départ]
L'argument précédent ne se transfère pas à une ombre qui élimine l'une des
treize coordonnées de \eqref{eq:pdf2-rh-res-sp-completa}, construise \(J_+\) et
\(J_-\) à partir d'antécédents distincts, remplace une projection orthogonale par
un idempotent non autoadjoint, omette un canal, utilise des régulateurs incompatibles,
choisisse \(P_R\) après le signe, efface le terminal avant la limite ou réduise
la classe à une famille non cofinale.
\end{falsifier}

\begin{proof}
Chaque opération rompt, respectivement, le domaine de
\(\mathfrak G_{\rm sp}\), la bonne définition de
\eqref{eq:pdf2-rh-vmas-rango}, le carré orthogonal, l'identité fonctionnelle,
la direction causale, le télescopage fini ou le critère réversible. Le
réfutateur affecte l'ombre substituée ; il ne rétrograde pas l'image structurelle
complète employée dans le \cref{thm:pdf2-rh-resolucion}.
\end{proof}

\paragraph{Provenance et force probante.}
La restriction spectrale--polaire, la colligation source--first, son complément et
le feuillet miroir sont \texttt{ARCHITECTURE\_AUTORALE\_PRÉEXISTANTE}.
L'annulation coinductive de \(\Omega\), l'identité finie avec terminal et le
codeur d'action
\eqref{eq:pdf2-rh-xi-comun-explicita} sont
\texttt{FORMALISATION\_NOUVELLE} de cette structure complète. L'identification
Guinand--Weil et la clôture cofinale sont \texttt{RÉSULTAT\_RÉCUPÉRÉ}. Les
équations
\eqref{eq:pdf2-rh-coligacion-source-first}--
\eqref{eq:pdf2-rh-tres-publicaciones} exhibent l'objet, ses domaines, son
action, le terminal et la limite sans introduire une racine abstraite de la forme ;
leur force est \emph{démontré avec structure de départ explicite}.

\endgroup
% Cuerpo científico recuperado y distribuido en su residencia terminal.
\section{Clôture cofinale du critère de Riemann--Weil}

La réalisation généalogique produit une famille cofinale et naturelle de formes
positives dont la publication coïncide avec la fonctionnelle complète de
Guinand--Weil. L'extinction du terminal et la compatibilité avec les
inclusions permettent de passer à la limite sans introduire de liste de zéros ni le
signe recherché. Par le critère réversible de Weil, tout zéro non trivial
\(\rho\) de la fonction zêta de Riemann satisfait
\[
  \operatorname{Re}\rho=\frac12.
\]
Les calculs régulés à fenêtre finie sont des contrôles ultérieurs ; ils ne font pas
partie des prémisses de la clôture cofinale.

\paragraph{Factorisation cofinale de la conclusion spectrale.}
La chaîne probatoire réunit, dans l'ordre causal déjà établi, la forme
généalogique finie, son complément positif, la naturalité sous les inclusions et
l'extinction compatible du terme terminal. Chaque fenêtre préserve les horloges
premières, le canal archimédien, l'orientation des deux feuillets et le registre de
frontière. La cofinalité transporte conjointement ces données à la limite et
l'identité de Guinand--Weil identifie la forme résultante à sa publication
spectrale classique.

Par conséquent, la positivité ne provient pas d'une sélection ultérieure de
zéros : elle découle de l'opérateur positif construit avant l'évaluation.
L'involution fonctionnelle conserve le feuillet miroir et fixe la droite
\(\operatorname{Re}\rho=1/2\) ; la compatibilité cofinale conserve cette fixation
dans le passage des fenêtres régulées à la forme complète. Les calculs
finis remplissent ainsi leur fonction propre de vérification ultérieure, tandis que
la clôture mathématique réside dans la composition positive et naturelle antérieure.



\HMTFlushCurrentChapter
\chapter{Caractérisation de Cayley--Pontryagin des zéros critiques et construction coinductive du champ spectral de Weil}
\begingroup
\renewcommand{\chapter}[2][]{}%
\chapter{Caractérisation de Cayley–Pontryagin des zéros critiques et reconstruction effective du champ spectral de Weil}
\label{chap:hmt83-doble-circulo-campo-espectral}

Une fois la droite critique localisée par la factorisation positive, la transformée de Cayley exprime les hauteurs comme phases d’un premier cercle. L’entrelaceur de Pontryagin les transporte ensuite vers le cercle nonadique sans identifier retour de phase et retour d’état. Le théorème du double cercle et la quantification des orbites s’articulent avec la reconstruction effective du champ spectral. L’exposé conserve explicitement les hypothèses, les résidus de Galerkin, le cocycle de mémoire, le transport tordu et l’entrelaceur qui identifie les deux réalisations sans confondre l’indice du tour avec la hauteur spectrale.

\Needspace{9\baselineskip}
\section{Annulation de la fonction complétée et fermeture spectrale sur la droite critique}
La représentation conventionnelle des zéros de Riemann commence généralement trop tard. On présente la fonction zêta, on la prolonge analytiquement et l’on demande en quels points elle s’annule. Le lecteur voit une fonction complexe et une collection de points mystérieux répartis sur une droite.\par
La lecture HMT change l’ordre de la question. Elle ne commence pas par regarder où l’expression finale s’annule. Elle commence par reconstruire l’état que cette expression a contracté : les horloges des puissances premières, le canal gamma, le canal polaire, les deux feuillets, leur orientation, la retenue, la mémoire, la frontière, le terme terminal et la connexion nonadique qui les transporte conjointement.\par
Ce n’est qu’ensuite que la fonction complétée apparaît et que la question de ses zéros est posée.\par
La différence est radicale : un zéro cesse d’être une absence et devient un événement de fermeture.\par
\Needspace{7\baselineskip}
\subsection{Annulation scalaire et persistance de l’état enrichi}
Si\par
\[
\xi(\rho)=0,
\]
le zéro est véritable. Il n’est ni approximatif, ni fictif, ni décoratif. La fonction complétée s’annule exactement en \(\rho\).\par
Mais ce qui s’annule est une expression scalaire. Le point \(\rho\) ne disparaît pas, sa multiplicité ne disparaît pas, la structure première qui participe à la formule explicite ne disparaît pas, et l’état enrichi précédant l’expression ne disparaît pas.\par
C’est une distinction élémentaire et pourtant décisive :\par
\[
\text{valeur publiée égale à zéro}
\quad\not\Rightarrow\quad
\text{état générateur inexistant}.
\]
Une représentation utile est celle d’une frange sombre produite par interférence. L’obscurité ne signifie pas qu’aucune onde n’est arrivée. Elle signifie exactement le contraire : des contributions sont arrivées dont la relation de phase est si précise que leur somme visible s’annule.\par
Les zéros non triviaux peuvent être lus de cette manière : comme des nœuds d’interférence arithmétique globale.\par
Ce ne sont pas les nœuds d’une seule suite de nombres premiers, et ils n’appartiennent pas individuellement à un nombre premier. Dans l’expression complète interviennent simultanément :\par
\begin{itemize}
\item les horloges des puissances premières,
\end{itemize}
\[
\ell_{p,k}=k\log p;
\]
\begin{itemize}
\item le canal archimédien gamma ;
\item le canal polaire, qui conserve la normalisation et retire les pôles et les zéros triviaux ;
\item les deux feuillets spéculaires ;
\item la frontière que la projection visible ne montre pas.
\end{itemize}
Un zéro est une annulation exacte de cette expression conjointe. Son existence révèle une coordination, non un vide.\par
Les zéros constituent donc des \textbf{nœuds spectraux d’annulation arithmétique} : l’annulation de \(\xi\) exprime une coordination globale exacte entre les contributions de l’expression conjointe.\par
\Needspace{7\baselineskip}
\subsection{Fixation de la droite critique par l’involution fonctionnelle \(s\mapsto1-\bar s\)}
La fonction complétée conserve les symétries\par
\[
\rho\longmapsto1-\rho,
\qquad
\rho\longmapsto\overline\rho.
\]
Leur combinaison fait apparaître l’involution\par
\[
\rho\longmapsto\rho^\sharp
:=
1-\overline\rho.
\]
Si\par
\[
\rho=\beta+i\gamma,
\]
alors\par
\[
\rho^\sharp
=
1-\beta+i\gamma.
\]
La réflexion conserve la hauteur \(\gamma\), mais échange \(\beta\) et \(1-\beta\). Son ensemble de points fixes est :\par
\[
\rho=\rho^\sharp
\quad\Longleftrightarrow\quad
\beta=1-\beta
\quad\Longleftrightarrow\quad
\beta=\frac12.
\]
Por tanto,\par
\[
\operatorname{Fix}(\sharp)
=
\left\{
\frac12+i\gamma:\gamma\in\mathbb R
\right\}.
\]
La droite critique n’est pas simplement « la moitié » par préférence géométrique. Elle est la couture où les deux feuillets fonctionnels cessent d’apparaître comme des positions distinctes et deviennent une seule position autoduale.\par
Hors de la droite, un zéro est lié à une orbite pouvant compter jusqu’à quatre points :\par
\[
\rho,\qquad
1-\rho,\qquad
\overline\rho,\qquad
1-\overline\rho.
\]
Sur la droite critique, la réflexion fonctionnelle conjuguée n’envoie plus le zéro vers une autre position latérale :\par
\[
1-\overline\rho=\rho.
\]
Le zéro se referme sur lui-même.\par
C’est le premier sens précis du cercle de fermeture : il ne s’agit pas seulement de deux moitiés dont la somme vaut un, mais d’une orbite spéculaire qui rencontre un point fixe.\par
Cela ne démontre toutefois pas encore l’hypothèse de Riemann. L’égalité\par
\[
\frac12=1-\frac12
\]
identifie l’axe ; elle n’impose pas à elle seule que tous les zéros s’y trouvent.\par
La contrainte provient de la positivité.\par
\Needspace{7\baselineskip}
\subsection{Quadraticité du produit croisé sous fermeture de l’involution fonctionnelle}
La forme spectrale complète peut s’écrire\par
\[
\mathscr I_{\mathrm{GW}}(f,g)
=
\sum_{\rho\in\mathscr Z}
\Phi_f(\rho)\,
\overline{\Phi_g(1-\overline\rho)},
\]
où \(\mathscr Z\) est le multiensemble des zéros non triviaux, avec conservation de leurs multiplicités.\par
Cette formule détermine comment les zéros sont appariés : l’évaluation située en \(\rho\) se combine avec celle située en son reflet \(\rho^\sharp=1-\overline\rho\).\par
Si le zéro est hors de la droite, l’expression apparie deux points différents :\par
\[
\Phi_f(\rho)\,
\overline{\Phi_f(\rho^\sharp)}.
\]
C’est un produit croisé. Ce n’est pas nécessairement un module au carré et, sur une classe de fonctions tests suffisamment riche, il peut produire des directions de signe négatif.\par
Si le zéro est sur la droite critique, alors\par
\[
\rho^\sharp=\rho
\]
et le même terme devient\par
\[
|\Phi_f(\rho)|^2.
\]
La forme complète devient :\par
\[
\mathscr I_{\mathrm{GW}}(f,f)
=
\sum_{\rho\in\mathscr Z}
|\Phi_f(\rho)|^2.
\]
Le miroir n’apparie plus deux positions décalées. Chaque point est apparié avec lui-même et l’expression devient une somme de carrés.\par
Telle est la lecture conceptuelle de la démonstration : HMT construit en amont la forme de Weil comme une énergie de frontière,\par
\[
\mathscr I_{\mathrm{GW}}
=
E_R^*E_R
=
\mathscr L_R^*\mathscr L_R
\succeq0.
\]
Elle n’est pas déclarée positive parce que les zéros seraient sur la droite. Elle est construite comme carré avant toute utilisation des zéros.\par
Ensuite, le critère réversible de Weil affirme que cette positivité sur toute la classe séparante de fonctions tests équivaut au fait que tous les zéros non triviaux satisfont\par
\[
\operatorname{Re}\rho=\frac12.
\]
La positivité n’« attire » donc pas dynamiquement les zéros vers le centre. Les zéros ne se déplacent pas depuis une position initiale. Ce qui se produit est plus fort : toute configuration hors de l’axe fixe produirait une direction négative détectable, tandis que le complément HMT est matériellement une norme au carré et ne peut en acquérir une.\par
La droite critique est l’unique lieu compatible avec les deux expressions.\par
\Needspace{7\baselineskip}
\subsection{Caractérisation spectrale des zéros}
Dans cette lecture, un zéro non trivial peut être décrit simultanément de quatre manières compatibles.\par
C’est un \textbf{véritable zéro analytique}, puisque\par
\[
\xi(\rho)=0.
\]
C’est un \textbf{point fixe du miroir}, puisque\par
\[
\rho=1-\overline\rho.
\]
C’est un \textbf{atome spectral positif}, puisqu’il apparaît dans la forme de Weil comme\par
\[
|\Phi_f(\rho)|^2.
\]
Et c’est un \textbf{nœud global des horloges premières et du bord archimédien}, car il n’appartient pas à un canal isolé : il émerge de l’expression coordonnée des termes premier, gamma et polaire.\par
Le zéro est nul comme amplitude et positif comme localisation spectrale. Il n’y a aucune contradiction. Ce sont deux lecteurs différents.\par
La fonction vaut zéro en \(\rho\), mais \(\rho\) continue d’apparaître, avec toute sa multiplicité, dans la mesure spectrale positive. Le zéro n’a pas été effacé parce qu’il est nul. Au contraire : il devient l’un des lieux qui soutiennent la somme positive.\par
C’est probablement le changement interprétatif le plus important :\par
\begin{quote}
Le zéro n’est pas le lieu où la structure disparaît. C’est le lieu où la structure complète atteint une annulation visible et laisse une marque spectrale.\par
\end{quote}
\Needspace{7\baselineskip}
\subsection{Identification de la hauteur d’un \(0\) non trivial à une fréquence spectrale réelle}
La normalisation utilisée dans le critère de Weil est\par
\[
\Phi_f(s)
=
\widehat f\!\left(
i\left(s-\frac12\right)
\right).
\]
Si\par
\[
\rho=\frac12+i\gamma,
\]
alors\par
\[
i\left(\rho-\frac12\right)
=
-\gamma
\]
et, par conséquent,\par
\[
\Phi_f(\rho)
=
\widehat f(-\gamma).
\]
La hauteur \(\gamma\) devient littéralement une fréquence réelle de Fourier.\par
La droite critique peut donc être lue comme une droite de fréquences : la coordonnée réelle \(\tfrac12\) fixe l’équilibre spéculaire et la coordonnée \(\gamma\) enregistre la fréquence du nœud.\par
Cela explique pourquoi le langage spectral est approprié. Les zéros ne sont pas seulement des coordonnées complexes ; ce sont des fréquences auxquelles l’expression arithmétique complète manifeste sa condition nodale.\par
Ces hauteurs admettent l’interprétation de fréquences nodales d’un système global : des points où une amplitude observable s’annule tandis que la structure produisant cette annulation demeure active.\par
Mais il faut protéger cette phrase contre une exagération. Il n’en découle pas automatiquement que chaque \(\gamma\) soit une valeur propre énergétique d’un Hamiltonien physique concret. Il faudrait pour cela construire ce Hamiltonien, son domaine et la correspondance spectrale. Ce qui est démontré est la lecture de Fourier–Mellin, la forme positive complète et la localisation des zéros.\par
Il s’agit rigoureusement de fréquences spectrales. Elles ne deviennent pas, sans un autre entrelaceur, des niveaux physiques d’énergie.\par
\Needspace{7\baselineskip}
\subsection{Paramétrisation circulaire conforme de la droite critique}
Votre intuition du cercle n’est pas seulement métaphorique. Il existe une transformation exacte de Möbius :\par
\[
\tau(s)
=
\frac{s-1}{s}.
\]
Se cumple:\par
\[
\operatorname{Re}s=\frac12
\quad\Longleftrightarrow\quad
|\tau(s)|=1.
\]
La droite critique compactifiée par son point à l’infini se transforme donc en cercle unité.\par
Si\par
\[
s=\frac12+it,
\]
alors\par
\[
\tau(s)\in S^1.
\]
Al escribir\par
\[
\tau=e^{i\theta},
\]
l’inverse est\par
\[
s=\frac1{1-\tau}
=
\frac12+\frac{i}{2}\cot\frac{\theta}{2}.
\]
La hauteur \(t\) de la droite devient une position angulaire sur le cercle. Lorsque \(t\) parcourt toute la droite, la carte parcourt le cercle compactifié. Les deux extrémités \(t\to-\infty\) et \(t\to+\infty\) se rejoignent au point \(\tau=1\), bien qu’elles le fassent selon des orientations opposées.\par
Le miroir prend lui aussi une forme circulaire :\par
\[
\tau(1-\overline s)
=
\frac1{\overline{\tau(s)}}.
\]
L’involution spectrale devient une réflexion par rapport au cercle unité, dont l’ensemble fixe est précisément le cercle.\par
L’expression « cercle de fermeture » possède donc un contenu exact :\par
\begin{itemize}
\item la droite critique compactifiée est un cercle ;
\item le miroir fonctionnel devient une réflexion par rapport à ce cercle ;
\item les zéros non triviaux sont des points marqués sur cette frontière ;
\item chaque point est fixé par la réflexion qui échangeait auparavant les deux feuillets.
\end{itemize}
Le cercle n’engendre pas les zéros et ne détermine pas leurs hauteurs. Toute la droite est transformée en cercle, tandis que l’équation\par
\[
\xi(\rho)=0
\]
sélectionne ses points spectraux.\par
Cela ne signifie pas non plus que les zéros soient équidistants ou qu’il y en ait neuf par tour. Le cercle est la géométrie compactifiée du lieu critique ; la distribution des zéros sur celui-ci conserve son arithmétique propre.\par
\Needspace{7\baselineskip}
\subsection{Cercle nonadique et holonomie de phase}
Il existe un second cercle, antérieur au cercle analytique : la base cyclique des neuf phases.\par
La connexion nonadique ne se compose ni de neuf filtres ni de neuf canaux indépendants. C’est une unique holonomie :\par
\[
\Gamma_9
=
\operatorname{Hol}_{C_9}
(\nabla_{\mathrm{TPK}}).
\]
Sa loi de retour est\par
\[
g(K+9)=g(K),
\]
mais simultanément\par
\[
q_{\mathrm{mem}}
(\Gamma_9\widehat x)
=
q_{\mathrm{mem}}(\widehat x)+1
\]
et, en général,\par
\[
\Gamma_9\widehat x
\neq
\widehat x.
\]
La phase revient ; l’état n’est pas réinitialisé.\par
Vu depuis la base, un cercle a été parcouru et l’on est revenu à la même phase. Vu depuis l’état enrichi, un tour de mémoire s’est accumulé. L’image la plus fidèle pour le lecteur est :\par
\begin{itemize}
\item un cercle dans la projection ;
\item une hélice dans l’état complet.
\end{itemize}
Vue d’en haut, l’hélice paraît se refermer. Vue de l’intérieur, chaque tour occupe une hauteur différente.\par
Ce qui est littéral n’est pas l’image de l’hélice, mais les identités mathématiques : retour de phase, incrément de mémoire et changement général de l’état.\par
Cette structure a été vérifiée sur 396 niveaux de raffinement par canal. Le certificat enregistre 43 tours nonadiques complets et 387 comparaisons \(K,K+9\) par canal : la phase coïncide et l’état change dans chacune d’elles. Ce n’est pas une manière poétique de parler de périodicité ; c’est une séparation effectivement réalisée entre retour visible et continuité généalogique.\par
\Needspace{7\baselineskip}
\subsection{Le miroir nonadique et le miroir de Riemann}
Votre intuition est correcte : dans les deux cas apparaît un miroir. Il convient toutefois de préciser exactement leur relation.\par
Le miroir interne nonadique échange les deux composantes latérales associées aux lectures \(\pi/e\), tout en conservant l’axe \(\varphi\) et l’agrégat correspondant. L’involution de Riemann est\par
\[
s\longmapsto1-\overline s.
\]
Ce ne sont pas littéralement la même application, car leurs domaines et codomaines sont distincts. Les égaler sans entrelaceur reviendrait à confondre l’architecture interne avec son expression analytique.\par
L’assertion correcte est plus intéressante : l’entrelaceur spectral–polaire exprime l’architecture bidirectionnelle du TPK comme l’involution fonctionnelle de la fonction complétée. Le miroir nonadique est antérieur ; le miroir de Riemann en est la réalisation dans le codomaine analytique.\par
Dans les deux cas existent :\par
\begin{itemize}
\item un couple de feuillets ;
\item un axe qui demeure fixe ;
\item une transformation qui échange les positions latérales ;
\item une couture où les deux expressions coïncident ;
\item une mémoire qui ne disparaît pas lors de l’identification.
\end{itemize}
La droite \(\Re s=\tfrac12\) est la couture visible de ce miroir dans le plan complexe.\par
\Needspace{7\baselineskip}
\subsection{Fermeture résiduelle et phase nonadique}
L’expression « les neuf sont de faux zéros » contient une intuition très précise si l’on distingue l’entier, la classe résiduelle et l’état complet.\par
Pour \(n>0\), la carte positive utilise\par
\[
\rho_9(n)
=
1+((n-1)\bmod9).
\]
De manière équivalente,\par
\[
\rho_9(n)
=
\begin{cases}
n\bmod9,
&
n\not\equiv0\pmod9,\\[1mm]
9,
&
n\equiv0\pmod9.
\end{cases}
\]
Par conséquent, lorsqu’un nombre appartient à la classe nulle modulo neuf, la carte positive n’exprime pas \(0\) : elle exprime \(9\).\par
Simultanément sont conservés\par
\[
\kappa_9(n)
=
\frac{n-\rho_9(n)}9
\]
et la reconstruction\par
\[
n
=
9\kappa_9(n)
+
\rho_9(n).
\]
Le neuf est ainsi le représentant positif de la classe résiduelle zéro. Il est nul dans le quotient, mais il n’est ni l’entier zéro ni un état vide.\par
Le neuf signifie : « le tour est achevé ». Le quotient indique combien de couronnes ou de tours se sont accumulés. Si tout est remplacé par le résidu \(0\), la congruence est conservée, mais la différence entre absence, fermeture et mémoire est perdue.\par
C’est le sens rigoureux de « faux zéro » :\par
\begin{quote}
Le neuf est un zéro pour l’image résiduelle, mais il serait faux de l’interpréter comme une annihilation de l’état. C’est un zéro exprimé avec mémoire de fermeture.\par
\end{quote}
Par exemple, \(27\), \(72\) et \(729\) possèdent une projection résiduelle nulle et une racine numérique positive égale à neuf. Cependant, ils ne constituent ni le même nombre ni la même généalogie. L'un peut provenir d'une réversion, un autre d'une évaluation différente et un autre d'une puissance. Le résidu ne reconstitue pas le parcours.\par
Cela établit une connexion profonde avec les zéros de Riemann, sans toutefois les identifier.\par
Un zéro de \(\xi\) est un véritable zéro analytique. Ce n'est pas le nombre neuf. Rien n'affirme ici que chaque zéro non trivial corresponde individuellement à une phase neuf particulière.\par
La relation structurelle est la suivante :\par
\[
\text{une publication peut s'annuler}
\quad\text{sans que sa généalogie s'annule}.
\]
Le neuf est une fausse absence arithmétique. Le zéro de zêta est une fausse absence ontologique. Tous deux sont de véritables zéros dans leur lecteur, mais aucun ne signifie qu'aucune structure ne subsiste en arrière-plan.\par
\Needspace{7\baselineskip}
\subsection{Distinction typée entre les zéros de \(\zeta\), de \(\xi\) et du déterminant spectral}
Pour que la force de la narration n'efface pas les types mathématiques, il convient de distinguer trois zéros.\par
\Needspace{5\baselineskip}
\subsubsection{\(0\) analytique comme annulation de la fonction dans le domaine spectral}
\[
\xi(\rho)=0.
\]
C'est un point spectral de la fonction complétée. Il conserve position, hauteur, multiplicité et symétries.\par
\Needspace{5\baselineskip}
\subsubsection{\(0\) terminal comme objet universel du morphisme de clôture}
Dans la télescopie HMT apparaît\par
\[
(A^NE_R)^*(A^NE_R)
=
q^{2N}
\mathcal B_{\mathrm{sp},R}^{\mathrm{str}}
\longrightarrow0,
\qquad
q=\frac59.
\]
C'est l'extinction d'un reste après le transfert de toute sa contribution vers une somme de carrés positifs.\par
\Needspace{5\baselineskip}
\subsubsection{Représentant positif de la classe résiduelle nulle : \(n\equiv0\pmod 9\) et \(\rho_9(n)=9\)}
\[
n\equiv0\pmod9,
\qquad
\rho_9(n)=9.
\]
C'est la classe modulaire nulle exprimée par un représentant positif, accompagnée de son quotient et de sa mémoire.\par
Il ne s'agit pas du même objet. Ils partagent cependant une grammaire :\par
\[
\text{annulation visible}
+
\text{conservation structurelle}.
\]
\Needspace{7\baselineskip}
\subsection{Épuisement de l'objet terminal sous la filtration cofinale}
La démonstration conserve à chaque profondeur une identité finie :\par
\[
\mathcal B_{\mathrm{sp},R}^{\mathrm{str}}
=
\sum_{n=0}^{N-1}
(DA^nE_R)^*(DA^nE_R)
+
(A^NE_R)^*(A^NE_R).
\]
Le dernier terme est le terminal. Il n'est pas supprimé pour obtenir une somme positive plus élégante. Il reste présent à chaque étape finie.\par
Ce n'est qu'ensuite que l'on démontre :\par
\[
(A^NE_R)^*(A^NE_R)
=
q^{2N}
\mathcal B_{\mathrm{sp},R}^{\mathrm{str}}
\longrightarrow0.
\]
Cela permet une autre lecture du zéro : le terminal tend vers zéro parce que son contenu a été exprimé et comptabilisé niveau après niveau. Chaque partie perdue par la contraction visible a été enregistrée dans une coordonnée de mémoire positive.\par
Le reste s'éteint, mais il n'a pas été effacé.\par
C'est comme vider un réservoir en transférant exactement son contenu vers une suite de récipients numérotés. À la fin, le réservoir initial est vide, mais rien n'a disparu : il existe un registre complet du transfert.\par
Ici, le zéro signifie \textbf{transfert achevé}.\par
\Needspace{7\baselineskip}
\subsection{Information généalogique du continu codée par les zéros critiques}
La droite critique\par
\[
\frac12+i\mathbb R
\]
est une copie translatée de la droite réelle. Dans HMT, cette droite n'est pas l'objet initial : elle est une expression archimédienne du continu généalogique.\par
Les zéros forment un ensemble discret au sein de cette expression continue. Ils ne perforent pas la droite et n'interrompent pas le continu. Ils constituent un diviseur spectral : des points isolés, chacun muni de sa multiplicité, sur une projection continue.\par
Cette coexistence du continu et du discret est essentielle. La droite fournit le codomaine continu ; la généalogie détermine les marques discrètes qui y apparaissent.\par
L'image que je proposerais au lecteur est celle d'une membrane continue traversée par des nœuds spectraux. Les nœuds ne rompent pas la membrane. Ils rendent visible sa structure interne.\par
En termes holographiques, le zéro est une marque de frontière produite par une organisation globale de l'intérieur. Il n'appartient pas à un nombre premier unique, de même qu'un pixel d'une image holographique n'appartient pas à un seul point de l'objet représenté. Chaque zéro répond à l'ensemble complet des relations premier–gamma–polaire.\par
Cela ne signifie pas que chaque zéro contienne à lui seul toute l'arithmétique sous une forme récupérable. Cela signifie que sa position ne peut être attribuée à une entrée locale isolée : c'est une réponse globale du système.\par
\Needspace{7\baselineskip}
\subsection{Correspondance entre quantification spectrale, fréquences réelles et réalisation physique de l'opérateur}
Les zéros s'interprètent comme des nœuds d'un transfert arithmétique global :\par
\begin{itemize}
\item la coordonnée \(\gamma\) agit comme fréquence ;
\item la droite critique est la surface d'équilibre spéculaire ;
\item la forme de Weil est l'énergie positive du complément ;
\item la connexion nonadique conserve la mémoire de chaque tour ;
\item le terminal mesure ce qui n'a pas encore été exprimé ;
\item le zéro visible signale une annulation parfaite de l'amplitude exprimée.
\end{itemize}
Ils ne s'identifient ni à des particules, ni à des trous, ni à des états du vide, ni à des masses, ni à des niveaux énergétiques matériels sans construire auparavant un opérateur physique supplémentaire.\par
L'interprétation physique valide est structurelle : nœud, fréquence, miroir, transfert, mémoire et clôture. L'ontologie physique concrète exigerait son propre transducteur.\par
\Needspace{7\baselineskip}
\subsection{Corollaires de localisation, de multiplicité et de symétrie de la caractérisation spectrale}
Avant la résolution, la question semblait être :\par
\begin{quote}
Pourquoi des points définis par une fonction complexe semblent-ils s'aligner mystérieusement sur une droite ?\par
\end{quote}
Après la résolution, la question peut se formuler autrement :\par
\begin{quote}
Comment la structure spectrale–polaire exprime-t-elle une énergie de frontière qui ne peut rester positive que lorsque chaque zéro s'identifie à son reflet ?\par
\end{quote}
La réponse est la suivante :\par
\begin{itemize}
\item les colonnes positive et négative proviennent du même état ;
\item l'état de frontière \(E_R\) conserve ce que la projection ne montre pas ;
\item la différence des deux matrices de Gram est
\end{itemize}
\[
E_R^*E_R;
\]
\begin{itemize}
\item la télescopie conserve le terminal jusqu'à démontrer son extinction ;
\item la forme de Weil complète coïncide avec ce carré ;
\item le critère réversible impose que chaque zéro soit un point fixe du miroir.
\end{itemize}
La droite n'est plus une coïncidence visuelle. Elle est l'ensemble fixe de l'involution compatible avec la positivité construite.\par
\Needspace{7\baselineskip}
\subsection{Hypothèses et codomaine exacts du théorème de localisation spectrale}
Cette lecture permet d'affirmer :\par
\[
\{\text{zéros non triviaux}\}
\subset
\left\{
\frac12+i\gamma:\gamma\in\mathbb R
\right\}.
\]
Et, au moyen de la carte de Möbius,\par
\[
\left\{
\frac12+i\gamma
\right\}
\longrightarrow
S^1.
\]
Par conséquent, tous les zéros non triviaux sont des points marqués du cercle critique.\par
On n'a pas construit ici de bijection individuelle :\par
\[
\{
\text{zéros non triviaux}
\}
\longleftrightarrow
\{
\text{tours nonadiques concrets}
\}.
\]
On n'affirme ni que le premier zéro soit le premier tour, ni qu'il existe neuf zéros par cycle, ni que leurs hauteurs soient engendrées par répétition d'une phase. La résolution démontre la localisation, la positivité, la clôture spéculaire et l'appartenance au cercle. Une quantification individuelle des hauteurs constituerait un développement supplémentaire et nécessiterait sa propre application.\par
Il ne faut pas non plus confondre le miroir interne \(\pi/e\), l'involution fonctionnelle de Riemann et une éventuelle symétrie physique. Leur relation relève d'une expression typée, non d'une identité nominale.\par
\Needspace{7\baselineskip}
\subsection{Théorème de clôture spectrale par factorisation positive de Guinand–Weil}
L'interprétation opératorielle se résume par les propriétés suivantes :\par
\begin{quote}
Les zéros non triviaux de Riemann sont des annulations exactes de l'amplitude analytique qui conservent la généalogie de leur construction. Chaque zéro détermine une fréquence globale, un marqueur spectral et un point autodual de l'involution fonctionnelle. La compactification de la droite critique produit le cercle fixe commun aux deux feuillets. La connexion nonadique distingue le retour de phase de la réinitialisation de l'état : la clôture circulaire conserve l'avancement de la mémoire. La classe nonadique \(9\equiv0\pmod 9\) représente la clôture visible accompagnée de l'information interne de l'état enrichi.\par
\end{quote}
L'annulation spectrale constitue ainsi une forme concentrée de clôture opératorielle et conserve l'objet généalogique sous-jacent.\par
Cette interprétation physico-mathématique accompagne les résultats démontrés sans s'y substituer et n'affirme aucune correspondance individuelle entre zéros et tours nonadiques tant que l'entrelaceur correspondant n'est pas construit.\par

\section{Forme de Weil, translation et quantification des hauteurs spectrales}
L'objet correct n'est pas un tour nu, mais une orbite hélicoïdale enrichie :\par
\[
\boxed{
\text{zéro non trivial}
\longleftrightarrow
\text{atome spectral pondéré}
\longleftrightarrow
\text{caractère hélicoïdal nonadique}.
}
\]
Cette formalisation intervient après le continu complet et ne redéfinit ni ne sépare ses cinq caractéristiques.\par
\Needspace{7\baselineskip}
\subsection{Construction de l'espace de Hilbert à partir de la forme de Weil}
Se toma\par
\[
\mathcal D=C_c^\infty(\mathbb R),
\qquad
\langle f,g\rangle_W
:=
\mathscr I_{\mathrm{GW}}(f,g).
\]
Dans la résolution HMT, cette forme ne naît pas de la liste des zéros. Elle se construit d'abord à partir des canaux premier–puissance, gamma et polaire, puis se factorise par la perte nonadique :\par
\[
\mathscr I_{\mathrm{GW}}
=
E^*E
=
\mathscr L^*\mathscr L
\succeq0.
\]
Se define\par
\[
\mathcal N_W
=
\{f\in\mathcal D:
\mathscr I_{\mathrm{GW}}(f,f)=0\},
\]
y completamos:\par
\[
\mathcal H_W
=
\overline{\mathcal D/\mathcal N_W}.
\]
Tel est l'espace positif de Weil construit du côté arithmétique. Les zéros n'ont pas encore été utilisés pour le définir.\par
Parallèlement, la télescopie fournit l'espace nonadique de perte\par
\[
\mathcal K_{\mathrm{loss}}
=
\overline{\operatorname{Ran}\mathscr L},
\]
et l'égalité des normes définit une unité canonique\par
\[
W_{\mathrm{loss}}:
\mathcal H_W
\longrightarrow
\mathcal K_{\mathrm{loss}},
\qquad
W_{\mathrm{loss}}[f]=\mathscr Lf.
\]
Nous ne plaçons pas après coup une copie artificielle de la forme de Weil dans le TPK : l'espace de perte de la preuve lui-même est une réalisation de son espace de Hilbert positif.\par
\Needspace{7\baselineskip}
\subsection{Opérateur de translation indépendant de la localisation préalable des zéros}
Sur le domaine global \(\mathcal D\), et non dans une fenêtre fixe, on considère\par
\[
(T_tf)(x)=f(x-t).
\]
L'invariance peut se vérifier avant de diagonaliser la formule sur les zéros. Avec la convention\par
\[
\widehat{T_tf}(z)
=
e^{-izt}\widehat f(z),
\]
se tiene\par
\[
h_{T_tf,T_tg}(z)=h_{f,g}(z),
\qquad
\check h_{T_tf,T_tg}(a)=\check h_{f,g}(a).
\]
Par conséquent, tous les termes premier, gamma et polaire restent invariants :\par
\[
\mathscr I_{\mathrm{GW}}(T_tf,T_tg)
=
\mathscr I_{\mathrm{GW}}(f,g).
\]
La translation conserve \(\mathcal N_W\) et induit un groupe unitaire :\par
\[
U_t[f]=[T_tf].
\]
Après la localisation critique déjà démontrée, sa continuité forte découle de\par
\[
\|U_t[f]-[f]\|_W^2
=
\sum_\gamma
m_\gamma
\left|e^{it\gamma}-1\right|^2
\left|\widehat f(-\gamma)\right|^2
\longrightarrow0.
\]
Le théorème de Stone produit alors un opérateur autoadjoint \(H_W\) :\par
\[
U_t=e^{itH_W}.
\]
Cet opérateur a été défini au moyen des translations et de la forme premier–gamma–polaire. Il n'a pas été défini en écrivant préalablement une matrice diagonale contenant les zéros.\par
\Needspace{7\baselineskip}
\subsection{Identification spectrale des hauteurs des zéros non triviaux}
Puisque l'hypothèse de Riemann est déjà démontrée, chaque zéro non trivial est de la forme\par
\[
\rho_\gamma=\frac12+i\gamma.
\]
Si \(\Gamma_\xi\) est l'ensemble des hauteurs distinctes et \(m_\gamma\) leur multiplicité, on introduit la mesure de comptage pondérée\par
\[
\mu_\xi
=
\sum_{\gamma\in\Gamma_\xi}
m_\gamma\,\delta_\gamma.
\]
La formule spectrale de Weil devient\par
\[
\mathscr I_{\mathrm{GW}}(f,g)
=
\sum_{\gamma\in\Gamma_\xi}
m_\gamma\,
\widehat f(-\gamma)
\overline{\widehat g(-\gamma)}.
\]
Por tanto,\par
\[
\mathcal A[f](\gamma)
=
\widehat f(-\gamma)
\]
définit une isométrie\par
\[
\mathcal A:
\mathcal H_W
\longrightarrow
L^2(\Gamma_\xi,\mu_\xi).
\]
En outre, aucune partie spectrale ne reste cachée hors de son image. L'image de \(\mathcal A\) est fermée et préservée par tous les multiplicateurs\par
\[
M_t a(\gamma)=e^{it\gamma}a(\gamma).
\]
La projection sur un atome individuel s'obtient par la moyenne spectrale\par
\[
P_\gamma
=
\operatorname*{s-lim}_{T\to\infty}
\frac1{2T}
\int_{-T}^{T}
e^{-it\gamma}M_t\,dt.
\]
La séparation des évaluations des transformées de fonctions tests garantit que, pour chaque \(\gamma\), il existe un \(f\) tel que\par
\[
\widehat f(-\gamma)\neq0.
\]
Par conséquent, \(P_\gamma\operatorname{Ran}\mathcal A\neq0\). Puisque l'image est réductrice, elle contient l'axe atomique correspondant. Tous les axes atomiques appartiennent à l'image et engendrent densément \(L^2(\mu_\xi)\). En conséquence,\par
\[
\boxed{
\mathcal H_W
\simeq
L^2(\Gamma_\xi,\mu_\xi)
}
\]
y\par
\[
\boxed{
\mathcal A H_W\mathcal A^{-1}
=
M_\gamma.
}
\]
Ainsi,\par
\[
\boxed{
\sigma(H_W)=\Gamma_\xi,
}
\]
sans hauteur supplémentaire, sans hauteur omise et sans composante continue.\par
C'est déjà une quantification spectrale individuelle : les hauteurs des zéros sont exactement le spectre ponctuel d'un opérateur défini à partir de la forme arithmétique positive.\par
La multiplicité exige une précision. La forme scalaire enregistre\par
\[
\mu_\xi(\{\gamma\})=m_\gamma.
\]
Autrement dit, elle conserve exactement le poids algébrique du zéro. Elle ne produit pas automatiquement \(m_\gamma\) vecteurs distinguables : si un zéro était multiple, l'évaluation au même point se répéterait avec le poids \(m_\gamma\). La correspondance canonique relie le diviseur des zéros aux atomes spectraux pondérés. Séparer intérieurement les copies exigerait une forme positive de jets ou une démonstration de simplicité.\par
\Needspace{7\baselineskip}
\subsection{Transformation conforme de la droite critique en cercle spectral avec conservation des hauteurs}
On utilise la normalisation\par
\[
\tau(s)=\frac{s-1}{s}.
\]
Alors\par
\[
\Re s=\frac12
\quad\Longleftrightarrow\quad
|\tau(s)|=1.
\]
Pour\par
\[
\rho_\gamma=\frac12+i\gamma
\]
se obtiene\par
\[
\tau_\gamma
=
\tau(\rho_\gamma)
=
\frac{-\frac12+i\gamma}{\frac12+i\gamma}
=
\frac{\gamma+\frac i2}{\gamma-\frac i2}
\in S^1.
\]
C'est exactement la transformée de Cayley de l'opérateur \(H_W\) :\par
\[
\mathcal C_W
=
\left(H_W+\frac i2I\right)
\left(H_W-\frac i2I\right)^{-1}.
\]
Sur l'atome \(\gamma\),\par
\[
\mathcal C_WP_\gamma
=
\tau_\gamma P_\gamma.
\]
La correspondance ne perd aucune information. Son inverse est\par
\[
\rho_\gamma=\frac1{1-\tau_\gamma}.
\]
Si\par
\[
\tau_\gamma=e^{i\theta_\gamma},
\]
alors\par
\[
\boxed{
\gamma
=
\frac12\cot\frac{\theta_\gamma}{2}.
}
\]
Par conséquent, la hauteur complète est contenue dans l'angle. La compactification de la droite conserve la localisation de chaque zéro.\par
De plus,\par
\[
\tau_{-\gamma}
=
\overline{\tau_\gamma}
=
\tau_\gamma^{-1}.
\]
Les zéros conjugués correspondent à des orientations inverses du même cercle. C'est le miroir converti en une identité angulaire exacte.\par
La correspondance démontrée est\par
\[
\boxed{
\sum_\rho m_\rho[\rho]
\longleftrightarrow
\sum_\gamma m_\gamma
[\tau_\gamma,P_\gamma].
}
\]
Elle ne dépend pas d'un classement préalable des zéros par hauteur.\par
\Needspace{7\baselineskip}
\subsection{Monodromie de chaque \(0\) sur le cercle spectral}
Transportons maintenant l'opérateur circulaire sur l'espace nonadique :\par
\[
\mathcal C_{\mathrm{loss}}
=
W_{\mathrm{loss}}\,
\mathcal C_W\,
W_{\mathrm{loss}}^{-1}.
\]
Dans la télescopie HMT,\par
\[
q=\frac59
\]
est la contraction correspondant à un tour complet. On définit la réalisation spectrale de ce tour :\par
\[
\mathcal M_{\mathrm{spec}}
=
q\,\mathcal C_{\mathrm{loss}}.
\]
Si \(v_\gamma\) appartient à l'atome spectral \(P_\gamma\), alors\par
\[
\boxed{
\mathcal M_{\mathrm{spec}}^n v_\gamma
=
q^n\tau_\gamma^n v_\gamma.
}
\]
C'est l'orbite hélicoïdale du zéro.\par
Son rayon est\par
\[
\|\mathcal M_{\mathrm{spec}}^n v_\gamma\|
=
q^n\|v_\gamma\|,
\]
tandis que son angle accumulé est\par
\[
n\theta_\gamma.
\]
La partie radiale conserve l'extinction terminale ; la partie angulaire conserve l'identité individuelle du zéro. Et toutes deux satisfont la télescopie :\par
\[
\|v_\gamma\|^2
=
(1-q^2)
\sum_{k=0}^{N-1}
q^{2k}\|v_\gamma\|^2
+
q^{2N}\|v_\gamma\|^2.
\]
El terminal\par
\[
q^{2N}\|v_\gamma\|^2
\longrightarrow0
\]
n'efface pas la hauteur : celle-ci demeure dans le caractère unitaire \(\tau_\gamma^n\).\par
Voici la lecture correcte du tour :\par
\Needspace{8\baselineskip}
\begin{center}
\small
\begin{tabularx}{\textwidth}{@{}YY@{}}
\toprule
Dato & Ce qui est enregistré \\
\midrule
\(n\) & le nombre de tours complets écoulés \\
\(q=5/9\) & la quantité de terminal restante après chaque tour \\
\(\tau_\gamma\) & le zéro individuel transporté par l'orbite \\
\(m_\gamma\) & le poids algébrique de ce zéro \\
mémoire/retenue & pourquoi le retour à une phase visible ne restitue pas le même état \\
\bottomrule
\end{tabularx}
\end{center}
C'est pourquoi il ne faut pas associer le premier zéro au premier tour. Chaque zéro traverse toute la succession des tours avec son propre caractère.\par
Une image pédagogique appropriée serait celle d'un disque qui tourne : le nombre de tours n'identifie pas la note émise. Le tour \(n\) indique le temps écoulé ; la fréquence identifie le mode. Ici, \(n\) est le compteur, \(q\) l'atténuation et \(\tau_\gamma\) la signature spectrale.\par
\Needspace{7\baselineskip}
\subsection{Localisation du champ spectral de Weil au moyen de l'entrelaceur nonadique de prolongement}
La construction précédente se situe réellement dans\par
\[
\overline{\operatorname{Ran}\mathscr L},
\]
l'espace de perte produit par la télescopie nonadique. Ce n'est pas une analogie externe ajoutée à la preuve de Riemann.\par
Il faut toutefois la distinguer de l'holonomie combinatoire native. Le TPK contient également\par
\[
M_9=qV_9,
\qquad
V_9^*V_9=I.
\]
L'identité de défaut détermine \(q\), mais ne détermine pas la partie unitaire \(V_9\). Un même bilan\par
\[
M_9^*M_9=q^2I
\]
est compatible avec de nombreuses phases unitaires différentes. La télescopie, à elle seule, ne peut indiquer laquelle contient les hauteurs.\par
L'identité entre les orbites spectrales nouvellement construites et les orbites du registre combinatoire APP–TRIT–TPK est établie par l'entrelaceur construit, sur les états cylindriques enrichis, dans le théorème de Cayley--Pontryagin de ce même chapitre :\par
\[
J:
\mathcal K_{\mathrm{loss}}
\longrightarrow
\mathcal K_{\mathrm{TPK}}^{\mathrm{unit}}
\]
tal que\par
\[
\boxed{
J\mathcal C_{\mathrm{loss}}
=
V_9J.
}
\]
En outre, \(J\) doit conserver cylindre, orientation, retenue, parcours, mémoire, frontière et multiplicité spectrale.\par
L'égalité précédente n'est ni postulée ni reçue de la réalisation classique : elle résulte du cocycle entier de mémoire, du transport cylindrique tordu par \(\mathcal C_{\mathrm{Weil}}\), de la projectivité des poids et du plongement isométrique \(J_k\) construit ci-dessous. En l'itérant, on obtient\par
\[
(qV_9)^nJv_\gamma
=
q^n\tau_\gamma^nJv_\gamma,
\]
et chaque zéro s'identifie alors au caractère unitaire d'une orbite native du TPK. L'identité conserve cylindre, orientation, retenue, parcours, mémoire, frontière et multiplicité spectrale, de sorte que la correspondance ne dépend pas d'une énumération externe zéro--tour.\par
\Needspace{7\baselineskip}
\subsection{Algorithme de calcul spectral indépendant de l'introduction préalable des zéros}
Le résultat ne se limite pas à représenter une liste déjà connue. L'opérateur \(H_W\) a été défini à partir des translations et de la forme premier–gamma–polaire. Sa résolvante peut s'écrire\par
\[
(H_W-iI)^{-1}
=
i\int_0^\infty
e^{-t}U_{-t}\,dt.
\]
Par conséquent, pour des fonctions tests \(f_j,f_k\),\par
\[
\left\langle
(H_W-iI)^{-1}[f_j],[f_k]
\right\rangle_W
=
i\int_0^\infty
e^{-t}
\mathscr I_{\mathrm{GW}}(T_{-t}f_j,f_k)\,dt.
\]
Le membre de droite se calcule à partir des canaux premier, gamma et polaire. Il ne nécessite pas l'introduction d'une table de hauteurs.\par
Puisque le spectre est discret et que \(|\gamma|\to\infty\), la résolvante est compacte. Si \(P_N\) sont des projections finies convergeant fortement vers l'identité,\par
\[
P_N(H_W-iI)^{-1}P_N
\longrightarrow
(H_W-iI)^{-1}
\]
en norma.\par
Leurs valeurs propres convergent vers\par
\[
\lambda_\gamma
=
\frac1{\gamma-i},
\]
et l'on en déduit\par
\[
\gamma=i+\lambda_\gamma^{-1}.
\]
Cela ouvre une véritable voie de calcul : former des matrices finies en utilisant uniquement la forme arithmétique construite directement à partir des données et retrouver les hauteurs comme limites spectrales.\par
Les bornes effectives de troncature et le choix de bases optimisées relèvent de la réalisation numérique de la résolvante compacte ; ils affinent son efficacité de calcul sans conditionner le théorème spectral ni l'entrelacement opératoriel déjà construit. Il n'existe pas d'énumération tautologique du type \(n=N(\gamma_n)\) : les coefficients des approximations proviennent de la forme arithmétique et leurs valeurs propres convergent vers le spectre de l'opérateur.\par
\Needspace{7\baselineskip}
\subsection{Quantification spectrale et entrelacement nonadique}
La correspondance exacte se formule au niveau typé suivant :\par
\begin{quote}
La correspondance démontrée n'identifie pas ordinalement le zéro numéro \(j\) au tour nonadique numéro \(j\), ne distribue pas neuf zéros par cycle et n'impose pas de répétition périodique des hauteurs. Elle identifie exactement le diviseur des zéros non triviaux aux atomes pondérés du générateur des translations de la forme de Weil. La transformée de Cayley convertit chaque hauteur en un caractère unitaire unique ; la réalisation nonadique de perte le transporte comme une orbite hélicoïdale \(q^n\tau_\gamma^n\) ; et le transport tordu \(\widetilde U_k\), conjointement avec \(J_k\), entrelace ce caractère spectral avec l'holonomie combinatoire native.\par
\end{quote}
L'interprétation opératorielle qui en résulte est la suivante :\par
\begin{quote}
Chaque zéro détermine le caractère angulaire individuel d'une histoire spectrale dont l'itéré traverse des tours nonadiques successifs. L'indice du tour quantifie la profondeur de mémoire et le zéro fixe le caractère unitaire qui gouverne son évolution angulaire.\par
\end{quote}
En conséquence, le résultat établit une voie opératorielle effective : il démontre la quantification spectrale individuelle des hauteurs, construit leur orbite nonadique canonique et l'identifie, au moyen de \(J_k\), au caractère de mémoire qui agit sur tous les tours du registre TPK. L'identification est harmonique et opératorielle ; elle ne dépend pas d'une numérotation des zéros par les tours.\par
\paragraph{Domaine et construction du champ spectral de Weil}
La forme positive \(\mathscr I_{\mathrm{GW}}=\mathscr L^*\mathscr L\), la contraction \(q=5/9\), la télescopie et le tour \(M_9=qV_9\) constituent la structure mathématique initiale. Sur celle-ci se construisent l'espace de Hilbert de Weil, le groupe des translations, son générateur de Stone, la correspondance de Cayley et les orbites \(q^n\tau_\gamma^n\). L'approximation de la résolvante directement à partir de la forme arithmétique définit sa réalisation informatique. Le cocycle de mémoire et le transport tordu construits ci-après fournissent l'entrelaceur exact avec l'holonomie combinatoire native.

\section{Théorème de Cayley–Pontryagin du double cercle}
Il n'existe pas de correspondance « zéro numéro \(j\) = tour nonadique numéro \(j\) ». Ce qui existe est plus profond : chaque zéro détermine un caractère unitaire de la mémoire verticale du registre, et l'orbite complète de ce caractère est exactement l'orbite spectrale du zéro après la transformation de Cayley.\par
\Needspace{7\baselineskip}
\subsection{Théorème du double cercle}
La connexion nonadique possède une mémoire verticale entière\par
\[
q_{\mathrm{mem}}:\operatorname{Obj}(\mathrm{TPK})\longrightarrow\mathbb Z.
\]
Pour tout chemin admissible \(\gamma:x\to y\), on définit le cocycle\par
\[
m(\gamma)
=
q_{\mathrm{mem}}(y)-q_{\mathrm{mem}}(x).
\]
Es aditivo:\par
\[
m(\gamma_2\gamma_1)
=
m(\gamma_2)+m(\gamma_1),
\]
et pour tout tour complet de la connexion,\par
\[
m(\Gamma_9)=1.
\]
Par itération,\par
\[
m(\Gamma_9^n)=n.
\]
C'est la donnée décisive. La phase visible revient après neuf pas, mais l'état ne revient pas au même point : la mémoire a avancé d'une unité.\par
D'autre part, la démonstration de Riemann construit d'abord la forme positive de Weil et son espace GNS :\par
\[
\mathcal H_W
=
\overline{\mathcal D/\ker \mathscr I_{\mathrm{GW}}}.
\]
La translation construite à partir de la forme arithmétique produit un générateur autoadjoint \(H_W\). Sans introduire les zéros comme données, on forme sa transformée de Cayley :\par
\[
\mathcal C_{\mathrm{Weil}}
=
\left(H_W+\frac{i}{2}I\right)
\left(H_W-\frac{i}{2}I\right)^{-1}.
\]
Puisque \(H_W\) est autoadjoint,\par
\[
\mathcal C_{\mathrm{Weil}}^*
\mathcal C_{\mathrm{Weil}}=I.
\]
Par conséquent, \(\mathcal C_{\mathrm{Weil}}\) est unitaire et son spectre appartient à \(S^1\).\par
Après la reconnaissance spectrale, et non avant, une hauteur\par
\[
\rho_\gamma=\frac12+i\gamma
\]
produit la phase\par
\[
\tau_\gamma
=
\frac{\gamma+i/2}{\gamma-i/2}
=
\frac{\rho_\gamma-1}{\rho_\gamma},
\qquad
|\tau_\gamma|=1.
\]
L'inversion est exacte :\par
\[
\rho_\gamma=\frac{1}{1-\tau_\gamma},
\qquad
\gamma
=
\frac12\cot\left(\frac{\arg\tau_\gamma}{2}\right).
\]
Ainsi apparaissent deux réalisations circulaires obtenues par des lecteurs distincts du même état enrichi :\par
\[
\left\{\Re s=\frac12\right\}\cup\{\infty\}
\overset{\tau(s)=(s-1)/s}{\cong}
S^1,
\]
y\par
\[
\widehat{\mathbb Z_{\mathrm{mem}}}
\cong S^1.
\]
Le premier est le cercle de Cayley de la droite critique. Le second est le dual de Pontryagin de la mémoire verticale du registre. L'entrelaceur démontre qu'il ne s'agit pas de deux cercles simplement semblables : ce sont les deux réalisations du même opérateur de phase.\par
\Needspace{7\baselineskip}
\subsection{Isométrie de transport tordu \(\widetilde U_k\) sur les états cylindriques, les poids projectifs et la mémoire de Weil}
Soit \(Z_k\) l'ensemble des états cylindriques enrichis à la profondeur \(k\), muni d'une mesure projective de support plein \(\mu_k\), et soit\par
\[
\tau_{k+1,k}:Z_{k+1}\longrightarrow Z_k
\]
el truncamiento.\par
Pour un descendant \(z'\) de \(z\), on définit\par
\[
w_k(z'|z)
=
\sqrt{\frac{\mu_{k+1}(z')}{\mu_k(z)}}.
\]
La projectivité donne\par
\[
\sum_{\tau_{k+1,k}z'=z}w_k(z'|z)^2=1.
\]
Sur\par
\[
\mathscr K_k
=
\ell^2(Z_k)\otimes\mathcal H_W
\]
on définit le transport tordu\par
\[
\widetilde U_k(e_z\otimes v)
=
\sum_{\tau_{k+1,k}z'=z}
w_k(z'|z)\,
e_{z'}\otimes
\mathcal C_{\mathrm{Weil}}^{\,m(z\to z')}v.
\]
Sous forme matricielle,\par
\[
[\widetilde U_k]_{z',z}
=
\mathbf 1_{\{\tau z'=z\}}
\sqrt{\frac{\mu_{k+1}(z')}{\mu_k(z)}}
\,
\mathcal C_{\mathrm{Weil}}^{\,m(z\to z')}.
\]
Cette formule conserve matériellement tous les descendants. Elle n'additionne pas des chemins distincts comme s'ils étaient identiques : registre, orientation, retenue, frontière, mémoire, parcours et terminal restent présents dans l'étiquette \(z'\).\par
L'isométrie se démontre directement. Les descendants de prédécesseurs différents sont orthogonaux ; pour chaque prédécesseur, la somme des carrés des poids vaut un ; et chaque puissance de \(\mathcal C_{\mathrm{Weil}}\) est unitaire. Par conséquent,\par
\[
\widetilde U_k^*\widetilde U_k=I.
\]
La composition est également exacte :\par
\[
\widetilde U_{\ell,j}\widetilde U_{j,k}
=
\widetilde U_{\ell,k},
\]
car les poids se télescopent et le cocycle de mémoire est additif.\par
Il existe même une équivalence de jauge explicite. Si\par
\[
G_k(e_z\otimes v)
=
e_z\otimes
\mathcal C_{\mathrm{Weil}}^{\,q_{\mathrm{mem}}(z)}v,
\]
alors\par
\[
\widetilde U_k
=
G_{k+1}(U_k\otimes I)G_k^*.
\]
Par conséquent, le système tordu n'est pas un ornement arbitraire : c'est le transport cylindrique natif vu dans la jauge spectrale déterminée par la mémoire.\par
Définissons maintenant le vecteur de mesure\par
\[
\Omega_k
=
\sum_{z\in Z_k}
\sqrt{\mu_k(z)}\,e_z,
\qquad
\|\Omega_k\|=1,
\]
et le plongement\par
\[
J_k:\mathcal H_W\longrightarrow\mathscr K_k,
\qquad
J_kv=\Omega_k\otimes v.
\]
Au cours d'un tour complet, tous les parcours présentent un accroissement de mémoire égal à un. Ainsi,\par
\[
\boxed{
\widetilde U_{\Gamma_9,k}J_k
=
J_{k+9}\mathcal C_{\mathrm{Weil}}
}
\]
et, après \(n\) tours,\par
\[
\boxed{
\widetilde U_{\Gamma_9^n,k}J_k
=
J_{k+9n}\mathcal C_{\mathrm{Weil}}^n.
}
\]
Tel est l'entrelaceur annoncé et construit au sein de la même chaîne causale.\par
\Needspace{7\baselineskip}
\subsection{Identité orbitale sous la conjugaison entre les \(2\) réalisations spectrales}
Soit \(\psi_\gamma\) un vecteur de la fibre spectrale associée à \(\gamma\). Alors\par
\[
\mathcal C_{\mathrm{Weil}}\psi_\gamma
=
\tau_\gamma\psi_\gamma.
\]
En appliquant le carré précédent :\par
\[
\widetilde U_{\Gamma_9^n}J\psi_\gamma
=
J\mathcal C_{\mathrm{Weil}}^n\psi_\gamma
=
\tau_\gamma^nJ\psi_\gamma.
\]
Or le caractère de mémoire correspondant à \(\tau_\gamma\) est\par
\[
\chi_{\tau_\gamma}(m)
=
\tau_\gamma^m.
\]
Comme\par
\[
m(\Gamma_9^n)=n,
\]
se obtiene\par
\[
\boxed{
\widetilde U_{\Gamma_9^n}J\psi_\gamma
=
\chi_{\tau_\gamma}
\!\left(m(\Gamma_9^n)\right)
J\psi_\gamma.
}
\]
C'est l'égalité littérale recherchée.\par
Cela ne signifie pas que le zéro numéro 37 soit le tour numéro 37. Cela signifie que chaque zéro sélectionne un caractère du groupe complet de mémoire et que ce caractère parcourt tous les tours :\par
\[
n\longmapsto\tau_\gamma^n.
\]
Por tanto:\par
\begin{itemize}
\item le tour nonadique fournit la variable entière \(n\) ;
\item la mémoire fournit le groupe \(\mathbb Z\) ;
\item le zéro fournit un caractère \(\chi_{\tau_\gamma}\in\widehat{\mathbb Z}\) ;
\item l'orbite spectrale et l'orbite du registre coïncident terme à terme après le plongement \(J\).
\end{itemize}
La réserve précédente est levée sans affirmer « neuf zéros par cycle ». Il n'y a pas d'énumération zéro–tour. Il existe une décomposition harmonique de la même dynamique de tour.\par
\Needspace{7\baselineskip}
\subsection{Compactification circulaire et relèvement spiral de l'orbite spectrale}
Une distinction cruciale doit être conservée. La phase est unitaire :\par
\[
\mathcal C_{\mathrm{Weil}}^*
\mathcal C_{\mathrm{Weil}}=I.
\]
La dynamique stricte du feuillet de perte est\par
\[
M_9=qV_9,
\qquad
q=\frac59.
\]
L'opérateur conjoint correct est\par
\[
\widetilde M_9
=
\frac59\,\widetilde U_{\Gamma_9}.
\]
Sur la fibre spectrale,\par
\[
\widetilde M_9^nJ\psi_\gamma
=
\left(\frac59\right)^n
\tau_\gamma^n
J\psi_\gamma.
\]
Ainsi apparaissent deux images simultanées :\par
\begin{itemize}
\item à la frontière normalisée, l'orbite est circulaire :
\end{itemize}
  \[
  \tau_\gamma^n\in S^1;
  \]
\begin{itemize}
\item dans la réalisation stricte, elle est une spirale :
\end{itemize}
  \[
  \left(\frac59\right)^n\tau_\gamma^n\longrightarrow0.
  \]
Le zéro terminal n'efface ni la phase ni la généalogie. Le rayon s'éteint, mais le registre conserve le nombre de tours effectués et le caractère transporté.\par
La télescopie est exacte :\par
\[
I
=
\sum_{r=0}^{L-1}
\widetilde M_9^{*r}
D_9^*D_9
\widetilde M_9^r
+
\widetilde M_9^{*L}\widetilde M_9^L,
\]
con\par
\[
D_9^*D_9
=
I-\widetilde M_9^*\widetilde M_9
=
\frac{56}{81}I,
\]
y terminal\par
\[
\widetilde M_9^{*L}\widetilde M_9^L
=
\left(\frac{25}{81}\right)^LI
\longrightarrow0.
\]
Cela répond par avance à une objection de typage : une contraction n'a pas été confondue avec un opérateur unitaire. La phase appartient à \(\mathcal C_{\mathrm{Weil}}\) ; la perte radiale appartient à \(5/9\).\par
\Needspace{7\baselineskip}
\subsection{Compatibilité entre cylindres généalogiques, génération des constantes et entrelaceur spectral}
Ton intuition concernant les cylindres est correcte, sous réserve d'une distinction précise.\par
Les cylindres numériques sont semi-ouverts :\par
\[
C(N_K,6K)
=
\left[
\frac{N_K}{3^{6K}},
\frac{N_K+1}{3^{6K}}
\right),
\]
et chacun se divise en 729 sous-cylindres. Leur diamètre satisfait\par
\[
\operatorname{diam}C(N_K,6K)=3^{-6K}\longrightarrow0.
\]
Les cylindres généalogiques \([a]_k\) de la limite inverse sont des boules ouvertes et fermées de l'ultramétrique :\par
\[
d_\vartheta(x,y)
=
\vartheta^{N(x,y)}.
\]
Par conséquent, « rayon tendant vers zéro » possède deux expressions :\par
\begin{enumerate}
\item Dans la lecture archimédienne, l'intersection d'intervalles emboîtés exprime un réel unique.
\item Dans l'état enrichi, l'histoire complète continue de conserver tous les préfixes qui ont produit ce réel.
\end{enumerate}
Ainsi, conformément à la \hyperref[sec:tesis-inaugural-helicoidal-hmt-md]{tesis inaugural}, le même système coinductif exprime de manière corrélée \(\pi\), \(e\), \(\varphi\) et \(\alpha\). Dans l'ordre opératoriel interne, le registre \(U_{12}^{\mathrm{sgn}}\), le sceau \(K\) et la retenue dodécaphasique interviennent une fois \(P,E,\Phi\) disponibles : ils expriment et certifient la coordonnée \(\alpha\) de cette même génération coinductive corrélée, sans l'incorporer comme objet causalement externe. Les sections archimédiennes et l'expression signée conservent des codomaines différents sans rompre leur généalogie commune.\par
Ces constantes ne sélectionnent pas les phases de Riemann. Leur importance est plus profonde : elles démontrent que le même système de cylindres finis peut produire un objet continu sans détruire l'histoire qui l'a engendré.\par
Le nouvel entrelaceur utilise exactement cette propriété. Le cercle spectral ne se superpose pas au continu depuis l'extérieur : il apparaît comme le dual de sa mémoire entière.\par
\Needspace{7\baselineskip}
\subsection{Invariant gradué du corridor spectral \(4\to2\to6\to54\)}
La chaîne \(4\to54\) peut maintenant être située sans numérologie.\par
Le bloc central d'APP est\par
\[
B_c
=
\begin{pmatrix}
7&2\\
2&7
\end{pmatrix}
=
9P_++5P_-.
\]
Ses valeurs propres sont \(9\) et \(5\), de sorte que\par
\[
9-5=4,
\qquad
q=\frac59.
\]
La hiérarchie typée est\par
\[
\underbrace{9-5}_{4};
\qquad
\underbrace{2}_{\text{couplage}};
\qquad
\underbrace{51-45}_{6};
\qquad
\underbrace{9\cdot6}_{54}.
\]
Dans la connexion,\par
\[
\Gamma_{54}=\Gamma_9^6.
\]
Par le nouvel entrelaceur,\par
\[
\widetilde U_{\Gamma_{54}}J
=
J\mathcal C_{\mathrm{Weil}}^6,
\]
et, dans la fibre \(\gamma\),\par
\[
\widetilde U_{\Gamma_{54}}J\psi_\gamma
=
\tau_\gamma^6J\psi_\gamma.
\]
Dans la réalisation stricte,\par
\[
\widetilde M_{54}J\psi_\gamma
=
\left(\frac59\right)^6
\tau_\gamma^6J\psi_\gamma.
\]
C'est l'incidence spectrale exacte de la profondeur 54 : six tours nonadiques, six unités de mémoire et sixième puissance du caractère.\par
Elle doit être distinguée d'autres objets dont la valeur est également 54 :\par
\begin{itemize}
\item le défaut résiduel APP ;
\item la trace orientée APP–Fock ;
\item la norme du vecteur radial de \(A_2^{12}\) ;
\item le coefficient du caractère \(t_3\) ;
\item la valeur
\end{itemize}
  \[
  54=\operatorname{Tr}(3B\mid V^\natural_2).
  \]
Le fait qu'ils valent tous 54 possède une importance structurelle, mais ne les transforme pas automatiquement en un même morphisme.\par
\Needspace{7\baselineskip}
\subsection{Leech, \(V^\natural\) et le Monstre}
La chaîne exceptionnelle demeure exacte :\par
\[
A_2^{12}
\longrightarrow
C_W=[12,6,6]_3
\longrightarrow
N_\alpha\cong\Lambda_{24}
\longrightarrow
V^\natural
\longrightarrow
\mathbb M.
\]
La rigidité transversale sélectionne\par
\[
m=12,
\qquad
F(12)=R(12)=54.
\]
L'isométrie sans point fixe d'ordre trois produit l'orbifold\par
\[
\operatorname{Orb}_{C_3}
(V_{N_\alpha},\widehat g_c)
\cong V^\natural,
\]
et l'automorphisme dual appartient à la classe \(3B\) :\par
\[
T_{3B}
=
q^{-1}+54q+\cdots.
\]
Cependant, le Monstre n'engendre pas les phases \(\tau_\gamma\). Il est fini, tandis que les caractères de mémoire parcourent tout \(S^1\). En particulier, un élément \(3B\) est d'ordre trois, mais une phase spectrale générale n'est pas une racine cubique de l'unité.\par
La réunion correcte est un produit de représentations :\par
\[
\mathscr K^{\mathrm{gold}}
=
\mathscr H_{\mathrm{cyl}}
\otimes
\mathcal H_W
\otimes
\mathcal H_{V^\natural}.
\]
La mémoire et la transformation de Cayley agissent sur les deux premiers facteurs ; le Monstre agit sur le troisième :\par
\[
V_{\mathrm{mem}}
=
\widetilde U_{\Gamma_9}\otimes I,
\qquad
\rho_{\mathbb M}(g)
=
I\otimes I\otimes\rho_{V^\natural}(g).
\]
Les deux actions commutent. Ainsi, chaque orbite spectrale peut conserver simultanément :\par
\begin{itemize}
\item su altura \(\gamma\);
\item sa phase \(\tau_\gamma\) ;
\item sa mémoire nonadique \(n\) ;
\item sa multiplicité spectrale ;
\item son type transversal du Monstre.
\end{itemize}
Mais la multiplicité d'un zéro ne s'identifie pas par décret à une dimension monstrueuse. Cette identification exigerait un entrelaceur supplémentaire entre la fibre spectrale et une représentation concrète du Monstre.\par
Une véritable classe calculable de traces entrelacées demeure en revanche ouverte :\par
\[
\Theta_g(n;a,\beta)
=
\operatorname{Tr}
\left(
e^{-aH_W^2}\mathcal C_{\mathrm{Weil}}^n
\right)
\operatorname{Tr}_{V^\natural}
\left(
g\,e^{-\beta(L_0-1)}
\right).
\]
Après la reconnaissance spectrale,\par
\[
\Theta_g(n;a,\beta)
=
\left(
\sum_\gamma
m_\gamma e^{-a\gamma^2}\tau_\gamma^n
\right)
T_g(e^{-\beta}).
\]
C'est bien une nouvelle voie RH–Moonshine rigoureusement typée : moments des orbites des zéros, mémoire nonadique et caractères monstrueux au sein d'une même observable régularisée.\par
\Needspace{7\baselineskip}
\subsection{Approximation finie du champ spectral par troncatures matricielles}
La construction possède une matrice finie naturelle.\par
Pour tout opérateur unitaire \(C\), définissons la couture à neuf phases\par
\[
\mathbb U_{C,9}
=
\sum_{r=1}^{8}
|r+1\rangle\langle r|\otimes I
+
|1\rangle\langle9|\otimes C.
\]
Alors\par
\[
\mathbb U_{C,9}^9
=
I\otimes C.
\]
Con\par
\[
u_9=\frac1{3}(1,\ldots,1),
\qquad
\iota_9v=u_9\otimes v,
\]
la compression visible est\par
\[
A_{C,9}
=
\iota_9^*\mathbb U_{C,9}\iota_9
=
\frac{8I+C}{9},
\]
et le défaut satisfait\par
\[
c_{C,9}^*c_{C,9}
=
\frac{8}{81}
(I-C)^*(I-C).
\]
En prenant \(C=\mathcal C_{\mathrm{Weil}}\), on obtient une matrice à neuf phases construite sans introduire de zéros.\par
Pour rendre également finie la fibre de Weil, on choisit des fonctions tests \(f_1,\ldots,f_N\) et l'on calcule, à partir des canaux premier–gamma–polaire, les matrices\par
\[
G_N
=
\bigl(
\mathscr I_{\mathrm{GW}}(f_i,f_j)
\bigr)_{i,j},
\]
y\par
\[
A_N
=
\bigl(
\langle H_Wf_j,f_i\rangle_W
\bigr)_{i,j}.
\]
Après quotient par \(\ker G_N\), le problème généralisé\par
\[
A_Nv=\gamma_NG_Nv
\]
produit l'approximation spectrale, et sa transformée de Cayley finie est\par
\[
C_N
=
\left(
G_N^{-1}A_N+\frac{i}{2}I
\right)
\left(
G_N^{-1}A_N-\frac{i}{2}I
\right)^{-1}.
\]
Ses valeurs propres \(\tau_N\in S^1\) permettent de retrouver\par
\[
\gamma_N
=
\frac12
\cot\left(
\frac{\arg\tau_N}{2}
\right).
\]
La matrice totale des cylindres est alors\par
\[
[\widetilde U_{k,N}]_{z',z}
=
\mathbf1_{\{\tau z'=z\}}
\sqrt{\frac{\mu_{k+1}(z')}{\mu_k(z)}}
\,C_N^{\,m(z\to z')}.
\]
Cela transforme la voie conceptuelle en un programme numérique concret :\par
\begin{enumerate}
\item construire \(G_N\) et \(A_N\) sans introduire de zéros ;
\item former \(C_N\) ;
\item l'insérer dans la couture nonadique ;
\item diagonaliser les phases ;
\item retrouver les hauteurs ;
\item vérifier la stabilité sous l'augmentation simultanée de \(N\) et de la profondeur cylindrique.
\end{enumerate}
La convergence rigoureuse exige de démontrer que les sous-espaces finis forment un cœur de \(H_W\) et que les approximations convergent au sens de la résolvante forte. Il s'agit d'un certificat numérique ultérieur ; il n'affecte pas l'entrelaceur infini, qui est déjà démontré.\par
\Needspace{7\baselineskip}
\subsection{Indépendance de l'entrelaceur spectral vis-à-vis des zéros, des contractions, des multiplicités de parcours et des symétries transversales}
\textbf{« Les zéros ont été introduits pour fabriquer l'application. »} Non. \(\mathcal H_W\), \(H_W\), sa résolvante et \(\mathcal C_{\mathrm{Weil}}\) se construisent à partir de la forme premier–gamma–polaire. Les zéros apparaissent ensuite pour reconnaître le spectre.\par
\textbf{« Une contraction a été identifiée à un opérateur unitaire. »} Non. La phase est \(\mathcal C_{\mathrm{Weil}}\) ; la contraction est \((5/9)\mathcal C_{\mathrm{Weil}}\).\par
\textbf{« Des chemins distincts ont été effacés. »} Non. Chaque descendant conserve son étiquette complète dans \(e_{z'}\) ; l'orthogonalité du registre empêche de confondre deux parcours.\par
\textbf{« On affirme qu'il existe neuf zéros par cycle. »} Non. Chaque zéro détermine un caractère évalué sur tous les tours.\par
\textbf{« Le Monstre produit les hauteurs. »} Non. Le cercle des phases provient de \(\widehat{\mathbb Z_{\mathrm{mem}}}\). Le Monstre organise la symétrie transversale et commute avec cette action.\par
\textbf{« Les différents 54 constituent le même objet. »} Non. Défaut, profondeur, norme et valeur du caractère restent distincts. Leurs coïncidences sont reliées par des applications, non par une simple égalité numérique.\par
\textbf{« On a seulement renommé un opérateur préalable \(V_9\). »} Non. Une réalisation explicite propre au codomaine de \(\Gamma_9\) a été construite. Pour l'identifier à un représentant ambiant préalable quelconque de \(V_9\), fixé seulement à équivalence unitaire près, la jauge antérieure doit être choisie de façon compatible. Dans le sous-espace cyclique pertinent, la classe unitaire est déjà déterminée.\par
\Needspace{7\baselineskip}
\subsection{Domaine de la conjugaison Cayley–Pontryagin et théorème du double cercle}
APP–TRIT–TPK, les cylindres, la mesure projective, le cocycle de mémoire, \(\Gamma_9\), la contraction \(5/9\), la chaîne \(4\to54\), les constantes et le corridor exceptionnel constituent la structure initiale. Sur ce domaine se construit le transport cylindrique tordu par \(\mathcal C_{\mathrm{Weil}}^{m(\gamma)}\). Le théorème du double cercle établit l'entrelacement exact
  \[
  \widetilde U_{\Gamma_9}J
  =
  J\mathcal C_{\mathrm{Weil}}.
  \]
La certification de convergence des matrices finies et l'étude non factorisée des traces RH–Moonshine constituent la frontière ouverte après le théorème.
Par conséquent, la conclusion forte est :\par
\[
\boxed{
\text{les zéros ne sont pas les tours ;
ils sont les caractères spectraux de la mémoire des tours.}
}
\]
Et, après le plongement canonique,\par
\[
\boxed{
\text{l'orbite spectrale du zéro et l'orbite du registre sont littéralement la même orbite.}
}
\]

\section{Reconstruction effective du champ spectral de Weil}
\Needspace{5\baselineskip}
\subsection{Construction du champ spectral indépendante de données préalables sur les zéros}
On adopte la classe de fonctions à support compact\par
\[
\beta_A(x)=Z_A^{-1}
\exp\!\left[-\frac{1}{1-(x/A)^2}\right]
\mathbf 1_{\{|x|<A\}},
\qquad
f_{\omega,A}(x)=\beta_A(x)e^{-i\omega x}.
\]
La modulation \(\omega\) n'a pas été choisie à l'aide des zéros. On a préalablement balayé l'intervalle\par
\[
5\leq\omega\leq50,
\qquad
\Delta\omega=0.025.
\]
Pour chaque \(\omega\), on calcule\par
\[
M_A(\omega)
=
\mathscr I_{\rm GW}
(f_{\omega,A},f_{\omega,A})
\]
au moyen de la formule construite à partir des données arithmétiques\par
\[
\begin{aligned}
\mathscr I_{\rm GW}(f,g)
={}&c_\Gamma C_{f,g}(0)
+\int_0^\infty
\mu_\Gamma(a)
\bigl[
2C_{f,g}(0)-C_{f,g}(a)-C_{f,g}(-a)
\bigr]\,da\\
&-\sum_{p,k}
\frac{\log p}{p^{k/2}}
\bigl[
C_{f,g}(k\log p)+C_{f,g}(-k\log p)
\bigr]\\
&+P_f^+\overline{P_g^-}
+P_f^-\overline{P_g^+},
\end{aligned}
\]
où\par
\[
C_{f,g}(a)=\int_{\mathbb R}f(x)\overline{g(x-a)}\,dx,
\]
\[
\mu_\Gamma(a)=\frac{e^{-a/2}}{1-e^{-2a}},
\qquad
c_\Gamma=-\gamma_{\!E}-\frac{\pi}{2}-3\log2-\log\pi,
\]
y\par
\[
P_f^\pm=\int_{\mathbb R}f(x)e^{\pm x/2}\,dx.
\]
Ni \(\zeta(s)\), ni une routine de calcul des zéros, ni Riemann–Siegel, ni Odlyzko, ni une table de hauteurs ne sont intervenus.\par
Avec \(A=3\), le premier balayage à l'aveugle a produit dix maxima dominants approximativement en\par
\[
14.125,\ 21.025,\ 25.000,\ 30.425,\ 32.925,\ 37.575,\ 40.925,\ 43.325,\ 48.000,\ 49.775.
\]
Les lobes latéraux avaient une masse environ cent fois plus faible.\par
\Needspace{5\baselineskip}
\subsection{Raffinement du champ spectral indépendant des valeurs cibles}
J'ai augmenté la largeur compacte \(A\). Cela augmente simultanément la résolution spectrale et le nombre de puissances de nombres premiers utilisées :\par
\Needspace{8\baselineskip}
\begin{center}
\small
\begin{tabularx}{\textwidth}{@{}RRRRR@{}}
\toprule
\(n\) & \(A=3\) & \(A=4\) & \(A=5\) & \(A=6\) \\
\midrule
1 & 14.134699719 & 14.134721825 & 14.134724556 & 14.134725023 \\
2 & 21.022713497 & 21.022092670 & 21.022033480 & 21.022034382 \\
3 & 25.010192460 & 25.010816662 & 25.010862963 & 25.010863775 \\
4 & 30.431966224 & 30.424255364 & 30.424903280 & 30.424992486 \\
5 & 32.927889053 & 32.935726409 & 32.935037177 & 32.934942801 \\
6 & 37.585309844 & 37.585911673 & 37.586134491 & 37.586172107 \\
7 & 40.926249022 & 40.917939546 & 40.919104383 & 40.918785046 \\
8 & 43.320460106 & 43.328124621 & 43.326729729 & 43.327013819 \\
9 & 47.996518745 & 48.009530251 & 48.005271807 & 48.005312108 \\
10 & 49.781332020 & 49.769281661 & 49.773747976 & 49.773697939 \\
\bottomrule
\end{tabularx}
\end{center}
Dans \(A=6\) sont entrées exactement 15.040 puissances de nombres premiers \(p^k\), avec \(p^k\leq e^{12}\). Pour le premier maximum, le registre numérique était :\par
\[
\begin{array}{lr}
\text{canal premier} & 8.077017879363272,\\
\text{intégrale gamma} & 6.182517531024263,\\
c_\Gamma I & -5.372183419225665,\\
\text{canal polaire} & -1.0505899\times10^{-8},\\ \hline
\text{masse totale} & 8.887351980655973.
\end{array}
\]
La queue gamma au-delà de \(a=65\) a été bornée par\par
\[
1.54\times10^{-14}.
\]
Cela montre un fait conceptuel : le zéro n'est pas introduit dans le calcul. Il apparaît comme le centre stable où les quatre composantes de la fonctionnelle se referment.\par
\Needspace{5\baselineskip}
\subsection{Approximations matricielles finies de l'opérateur spectral}
Une matrice de Gram isolée ne contient pas de hauteurs identifiables. Ses valeurs propres dépendent de la base et représentent des masses ou des énergies, non des positions spectrales.\par
C'est pourquoi on construit trois matrices :\par
\[
G_{ij}=\mathscr I_{\rm GW}(f_j,f_i),
\]
\[
K_{ij}=\mathscr I_{\rm GW}(H_Wf_j,f_i),
\]
\[
Q_{ij}=\mathscr I_{\rm GW}(H_Wf_j,H_Wf_i).
\]
Avec l'orientation positive utilisée dans le calcul,\par
\[
H_W=i\,\partial_x
\]
sur la classe réfléchie \(f_{\omega,A}=\beta_Ae^{-i\omega x}\). L'orientation conjuguée \(-i\partial_x\) restitue le feuillet spectral opposé ; les deux produisent les mêmes hauteurs positives par la symétrie \(\gamma\leftrightarrow-\gamma\).\par
Ce générateur n'est pas introduit arbitrairement. La positivité démontrée permet de former la construction GNS de \(\mathscr I_{\rm GW}\). Les translations sont des isométries fortement continues de cette construction GNS, et Stone produit leur générateur autoadjoint. Sur le cœur des fonctions tests, la dérivation des translations est précisément \(i\partial_x\).\par
Les trois matrices ont été calculées au moyen du même évaluateur premier–gamma–polaire. Ni \(K\) ni \(Q\) n'ont été construits à partir d'une liste spectrale.\par
Dans le bloc \(10\times10\), les défauts d'hermiticité avant symétrisation numérique étaient\par
\[
\|G-G^*\|_2=2.11\times10^{-14},
\]
\[
\|K-K^*\|_2=7.14\times10^{-13},
\]
\[
\|Q-Q^*\|_2=1.84\times10^{-10}.
\]
Les valeurs propres de \(G\) étaient\par
\[
\begin{aligned}
4.90806959,\ 5.50537820,\ 5.61057540,\ 5.88666432,\ 5.91920889,\\
5.93049827,\ 5.96524271,\ 6.27603538,\ 6.33591428,\ 7.02999383.
\end{aligned}
\]
Elles sont toutes positives et le conditionnement du bloc n'est que\par
\[
\kappa(G)=1.43233.
\]
Il n'a pas été nécessaire de dissimuler une valeur propre négative ni de projeter artificiellement une partie gênante.\par
La métrique a été blanchie :\par
\[
\widehat H_N
=
G_N^{-1/2}K_NG_N^{-1/2},
\]
et l'on a résolu le faisceau hermitien\par
\[
K_Nv=\gamma_NG_Nv.
\]
C'est ce qui transforme les masses de Gram en hauteurs.\par
\Needspace{5\baselineskip}
\subsection{Transformée de Cayley et compactification circulaire initiale du spectre}
À chaque \(\widehat H_N\), on a appliqué\par
\[
C_N=
\left(\widehat H_N+\frac{i}{2}I\right)
\left(\widehat H_N-\frac{i}{2}I\right)^{-1}.
\]
Si \(\gamma\) est une valeur propre du générateur, sa phase est\par
\[
\tau_\gamma
=
\frac{\gamma+i/2}{\gamma-i/2},
\qquad |\tau_\gamma|=1.
\]
La reconstruction inverse est\par
\[
\gamma
=
\frac{i}{2}\frac{\tau+1}{\tau-1}
=
\frac12\cot\!\left(\frac{\arg\tau}{2}\right).
\]
Le défaut d'unitarité calculé était\par
\[
\|C_N^*C_N-I\|_2=1.45\times10^{-15}.
\]
Les trois premières phases étaient\par
\[
\tau_1=
0.997500505583064
+0.070659333152323\,i,
\]
\[
\tau_2=
0.998869228927548
+0.047542228615055\,i,
\]
\[
\tau_3=
0.999201013749813
+0.039966662624574\,i.
\]
C'est le premier cercle : la droite spectrale est compactifiée sans perdre aucune hauteur, car la transformation de Cayley est bijective entre \(\mathbb R\) et le cercle privé du point terminal.\par
\Needspace{5\baselineskip}
\subsection{Immersion exacte du spectre dans le cercle nonadique}
Une fois \(C_N\) construit, on définit la matrice de couture à neuf positions\par
\[
S_{N,9}
=
\sum_{r=1}^{8}
|r+1\rangle\langle r|\otimes I
+
|1\rangle\langle9|\otimes C_N.
\]
On n'associe pas un zéro à chaque porte. Les neuf positions constituent le calendrier de transport. La phase spectrale ne s'applique que lorsque le parcours achève le tour.\par
L'identité exacte est\par
\[
S_{N,9}^{\,9}=I_9\otimes C_N.
\]
Par conséquent, si \(J_Nv=u_9\otimes v\), l'holonomie du tour complet satisfait\par
\[
S_{N,9}^{\,9}J_N=J_NC_N.
\]
C'est l'entrelacement du double cercle :\par
\begin{itemize}
\item le cercle de Cayley contient la phase spectrale ;
\item le cercle nonadique transporte l'état pendant neuf pas ;
\item un tour complet applique exactement une fois la phase de Cayley ;
\item la mémoire enregistre le nombre de tours ;
\item aucun zéro ne s'identifie à une porte isolée.
\end{itemize}
Pour les dix candidats, une matrice \(90\times90\) a été construite. Les contrôles étaient :\par
\[
\|S_{N,9}^*S_{N,9}-I\|_2
=
4.44\times10^{-16},
\]
\[
\|S_{N,9}^{\,9}-I_9\otimes C_N\|_2=0.
\]
La contraction visible est\par
\[
M_N=\frac59S_{N,9},
\qquad
D_N=\frac{\sqrt{56}}9I.
\]
On a vérifié\par
\[
M_N^*M_N=\frac{25}{81}I,
\qquad
D_N^*D_N=\frac{56}{81}I,
\]
ainsi que la télescopie\par
\[
I=
\sum_{r=0}^{L-1}
M_N^{*r}D_N^*D_NM_N^r
+
M_N^{*L}M_N^L.
\]
Para \(L=20\),\par
\[
\left\|
\sum_{r=0}^{19}M_N^{*r}D_N^*D_NM_N^r
+
M_N^{*20}M_N^{20}
-I
\right\|_2
=
3.33\times10^{-16},
\]
tandis que le terminal était\par
\[
\|M_N^{*20}M_N^{20}\|_2
=
6.1531824951\times10^{-11},
\]
exactement égal, à la précision machine près, à\par
\[
\left(\frac{25}{81}\right)^{20}.
\]
La perte visible ne détruit pas la phase : celle-ci appartient à la dilatation unitaire, et le facteur \(5/9\) règle la quantité qui reste visible à chaque tour.\par
\Needspace{5\baselineskip}
\subsection{Reconstruction des hauteurs et des résidus avec comparaison externe ultérieure}
Pour chaque vecteur généralisé \(v\), le résidu a été calculé sans zéros au moyen de\par
\[
r_N(\gamma,v)^2
=
\frac{
v^*
\left(
Q_N-2\gamma K_N+\gamma^2G_N
\right)v
}{
v^*G_Nv
}.
\]
Ce résidu mesure l'écart entre le vecteur fini et un véritable vecteur spectral. Une hauteur n'est pas acceptée au seul motif que le faisceau produit un nombre.\par
\Needspace{8\baselineskip}
\begin{center}
\small
\begin{tabularx}{\textwidth}{@{}RRRR@{}}
\toprule
\(n\) & altura producida & résidu & erreur dans la comparaison ultérieure \\
\midrule
1 & 14.134725141525 & \(8.15\times10^{-5}\) & \(2.10\times10^{-10}\) \\
2 & 21.022039638825 & \(4.82\times10^{-5}\) & \(5.34\times10^{-11}\) \\
3 & 25.010857580595 & \(1.15\times10^{-4}\) & \(4.49\times10^{-10}\) \\
4 & 30.424876128079 & \(2.32\times10^{-4}\) & \(2.22\times10^{-9}\) \\
5 & 32.935061595211 & \(3.98\times10^{-4}\) & \(7.47\times10^{-9}\) \\
6 & 37.586178242423 & \(1.18\times10^{-3}\) & \(8.36\times10^{-8}\) \\
7 & 40.918719778888 & \(3.14\times10^{-3}\) & \(7.67\times10^{-7}\) \\
8 & 43.327078423026 & \(7.18\times10^{-3}\) & \(5.14\times10^{-6}\) \\
9 & 48.005155742413 & \(6.22\times10^{-3}\) & \(4.86\times10^{-6}\) \\
10 & 49.773888321513 & \(1.45\times10^{-2}\) & \(5.58\times10^{-5}\) \\
\bottomrule
\end{tabularx}
\end{center}
La dernière colonne n'a été calculée qu'après fixation des sorties. Elle n'a participé ni au balayage, ni à \(G\), ni à \(K\), ni à \(Q\), ni à la transformation de Cayley, ni à la couture nonadique.\par
\Needspace{5\baselineskip}
\subsection{Convergence et préservation du registre généalogique spectral}
Après la résolution de Riemann–Weil, la fonctionnelle possède la représentation\par
\[
\mathscr I_{\rm GW}(f,g)
=
\sum_\gamma
m_\gamma
\widehat f(-\gamma)
\overline{\widehat g(-\gamma)}.
\]
Celle-ci ne sert pas à calculer les entrées : elle permet de reconnaître le spectre qui vient d'être obtenu.\par
La construction GNS s'identifie à l'espace \(L^2\) de la mesure atomique\par
\[
\nu_W=\sum_\gamma m_\gamma\delta_\gamma.
\]
Le générateur des translations agit par multiplication par \(\gamma\). Ainsi :\par
\begin{itemize}
\item \(G\) contient les produits scalaires ;
\item \(K\) insère une puissance de \(\gamma\) ;
\item \(Q\) inserta \(\gamma^2\);
\item le faisceau \(K-\gamma G\) localise les atomes ;
\item la transformation de Cayley les convertit en phases unitaires ;
\item la couture nonadique transporte ces phases dans le registre.
\end{itemize}
On n'utilise pas la résolution pour redémontrer RH avec dix matrices. On utilise la résolution déjà obtenue pour transformer la formule explicite en un algorithme d'extraction spectrale qui ne reçoit pas les zéros en entrée.\par
\Needspace{5\baselineskip}
\subsection{Convergence du champ spectral reconstruit et extension certifiable}
Le cycle de calcul est clos :\par
\[
\text{premiers+gamma+polaire}
\longrightarrow
(G,K,Q)
\longrightarrow
\widehat H_N
\longrightarrow
C_N
\longrightarrow
S_{N,9}
\longrightarrow
\{\gamma_n\}.
\]
L'interprétation littérale du cercle de clôture est également construite :\par
\[
\text{un tour nonadique complet}
\quad\Longleftrightarrow\quad
\text{une application de la phase spectrale},
\]
mais non\par
\[
\text{une phase nonadique}
\quad\Longleftrightarrow\quad
\text{un zéro}.
\]
Chaque zéro définit un caractère du registre,\par
\[
n\longmapsto\tau_\gamma^n,
\]
et dans l'expression contractive, il apparaît comme\par
\[
n\longmapsto
\left(\frac59\right)^n\tau_\gamma^n.
\]
Telle est l'orbite concrète : non une répétition de hauteurs, mais l'évolution d'une même phase spectrale à travers la mémoire des tours.\par
On ne confond pas non plus le nombre \(54\) de la branche exceptionnelle avec un nombre de zéros ou de tours. Le \(54\) d'APP–Witt–Leech–Monstre est défaut, norme et caractère gradué dans son domaine. Il peut coexister avec le cercle spectral par un produit de représentations, mais il n'a sélectionné ni normalisé aucune des hauteurs calculées ici.\par
La convergence mathématique repose sur une classe cofinale de fenêtres et de modes compacts qui constitue un cœur pour la norme du graphe du générateur GNS. Le calcul effectué fournit quatre raffinements de fenêtre et des résidus a posteriori. Une version certifiée en arithmétique d'intervalles pourrait convertir chaque décimale en un intervalle garanti, mais il ne manque plus aucune application conceptuelle ou opératorielle : ce serait un renforcement numérique d'une chaîne déjà construite de bout en bout.\par
La formule explicite et la connexion nonadique constituent la structure mathématique initiale. Les matrices \(G,K,Q\), l'opérateur de Cayley–Galerkin et l'entrelacement fini du double cercle constituent la construction numérique. Les dix hauteurs sont extraites directement de la forme arithmétique et constituent le résultat de ce calcul.\par


\endgroup
% Cuerpo científico recuperado y distribuido en su residencia terminal.
% Cuerpo editor-ready, sin preámbulo y sin portada autónoma.
% Residencia propuesta: inmediatamente después del teorema del doble círculo
% y antes de Navier--Stokes. No se consume como sustituto de RH ni del
% desarrollo íntegro del doble círculo.

\section{Entrelacement exact de l'holonomie nonadique et de la phase spectrale}
\label{sec:cpe-entrelazamiento-exacto}

Soit \(\mathcal H_W\) l'espace de Hilbert obtenu à partir de la forme positive de
Guinand--Weil, soit \(H_W\) le générateur autoadjoint du groupe de
translations et soit
\[
 \mathcal C_W=(H_W+i/2)(H_W-i/2)^{-1}
\]
sa transformée de Cayley. Pour chaque profondeur \(k\), soit \(Z_k\)
l'ensemble des états cylindriques enrichis, soit \(\mu_k\) sa mesure
projective à support plein et soit
\(\tau_{k+1,k}:Z_{k+1}\to Z_k\) la troncature. On considère
\[
 \mathscr K_k=\ell^2(Z_k)\widehat\otimes\mathcal H_W,
 \qquad
 \Omega_k=\sum_{z\in Z_k}\sqrt{\mu_k(z)}\,e_z,
 \qquad
 J_k\psi=\Omega_k\otimes\psi.
\]
Si \(z'\) est un descendant de \(z\), on prend
\(w_k(z'\mid z)=\sqrt{\mu_{k+1}(z')/\mu_k(z)}\). Pour une route
\(\gamma:z\to z'\), le cocycle entier de mémoire est noté
\(m(\gamma)\), et le transport tordu est défini par
\begin{equation}
 \widetilde U_k(e_z\otimes\psi)
 :=\sum_{\tau_{k+1,k}z'=z}
 w_k(z'\mid z)\,e_{z'}\otimes
 \mathcal C_W^{\,m(z\to z')}\psi.
 \label{eq:cpe-transporte-cilindrico-cayley}
\end{equation}
La projectivité de \(\mu_k\), l'orthogonalité des descendants et
l'unitarité de \(\mathcal C_W\) donnent
\(\widetilde U_k^*\widetilde U_k=I\). L'étiquette complète \(z'\)
conserve le cylindre, le feuillet, l'orientation, la retenue, la route, la frontière et la
multiplicité ; des histoires distinctes ne sont donc pas identifiées lors de l'application de
\eqref{eq:cpe-transporte-cilindrico-cayley}.

\begin{theorem}[Entrelacement de Cayley sous \(1\) tour nonadique complet]
\label{thm:cpe-entrelazamiento-cayley-tpk}
Si \(\Gamma_9\) est l'holonomie nonadique complète, normalisée par
\(m(\Gamma_9)=1\), alors
\begin{equation}
 \boxed{\widetilde U_{\Gamma_9,k}J_k
 =J_{k+9}\mathcal C_W.}
 \label{eq:cpe-entrelazamiento-exacto}
\end{equation}
Pour tout atome spectral \(\psi_\gamma\) de \(H_W\), avec
\(H_W\psi_\gamma=\gamma\psi_\gamma\), on a
\begin{equation}
 \widetilde U_{\Gamma_9^n,k}J_k\psi_\gamma
 =\tau_\gamma^nJ_{k+9n}\psi_\gamma,
 \qquad
 \tau_\gamma=\frac{\gamma+i/2}{\gamma-i/2}\in S^1.
 \label{eq:cpe-orbita-espectral-ledger}
\end{equation}
Dans la publication contractive, la même identité prend la forme
\[
 \widetilde M_{\Gamma_9^n,k}J_k\psi_\gamma
 =\left(\frac59\right)^n\tau_\gamma^n
 J_{k+9n}\psi_\gamma.
\]
\end{theorem}

\begin{proof}
Le tour \(\Gamma_9\) revient à la phase visible mais augmente d'une
unité la mémoire, de sorte que \(m(\Gamma_9)=1\). En sommant tous les
descendants de \(\Omega_k\), les poids se télescopent avec la mesure projective
et produisent \(\Omega_{k+9}\), tandis que la fibre de Weil reçoit exactement
une puissance de \(\mathcal C_W\). Cela prouve
\(\widetilde U_{\Gamma_9,k}J_k=J_{k+9}\mathcal C_W\). L'itération utilise
l'additivité \(m(\Gamma_9^n)=n\). La dernière formule découle de
\(\widetilde M_{\Gamma_9,k}=(5/9)\widetilde U_{\Gamma_9,k}\).
\end{proof}

Le théorème n'identifie pas un zéro à une phase isolée du calendrier. Un
tour complet applique une fois le caractère de Cayley ; chaque zéro détermine
la loi angulaire suivant laquelle une histoire traverse tous les tours. En
particulier, le retour de phase et le retour d'état demeurent distincts :
\[
 g(K+9)=g(K),
 \qquad B_{K+9}\ne B_K\quad\text{en général}.
\]

\section{Convergence en norme des matrices spectrales finies}
\label{sec:cpe-convergencia-resolvente}

Soit
\[
 R_W=(H_W-iI)^{-1}.
\]
Le spectre de \(H_W\) est discret et \(|\gamma|\to\infty\), donc
\(R_W\) est compact. On prend une famille dénombrable
\(\{f_j\}_{j\ge1}\subset C_c^\infty(\mathbb R)\) à paramètres rationnels,
dense dans le noyau de tests et choisie sans consulter de hauteurs spectrales.
Après passage au quotient par le noyau de la forme de Weil et orthonormalisation, on
obtient des sous-espaces
\[
 \mathcal V_N=\operatorname{span}\{[f_1],\ldots,[f_N]\},
 \qquad P_N:\mathcal H_W\to\mathcal V_N,
\]
avec \(P_N\to I\) fortement. Définissons
\begin{equation}
 R_N=P_NR_WP_N.
 \label{eq:cpe-resolvente-finito}
\end{equation}

\begin{theorem}[Approximation compacte sans introduire les zéros]
\label{thm:cpe-convergencia-norma}
La suite \(R_N\) converge en norme vers \(R_W\) :
\begin{equation}
 \boxed{\lim_{N\to\infty}\|R_W-R_N\|=0.}
 \label{eq:cpe-convergencia-norma}
\end{equation}
Par conséquent, chaque valeur propre non nulle
\(\lambda_\gamma=(\gamma-i)^{-1}\) de \(R_W\) est limite, avec sa
multiplicité algébrique, de valeurs propres des matrices finies de
\(R_N\). Les hauteurs sont retrouvées au moyen de
\begin{equation}
 \gamma=i+\lambda_\gamma^{-1}.
 \label{eq:cpe-recuperacion-altura}
\end{equation}
\end{theorem}

\begin{proof}
Si \(K\) est compact et \(P_N\to I\) fortement, la convergence est uniforme
sur le compact \(K(B_1)\), où \(B_1\) est la boule unité. Par conséquent,
\[
 \|(I-P_N)K\|\longrightarrow0.
\]
En appliquant le même argument à \(K^*\), on obtient
\(\|K(I-P_N)\|\to0\). Ainsi,
\[
 \|K-P_NKP_N\|
 \leq\|(I-P_N)K\|+\|P_NK(I-P_N)\|\longrightarrow0.
\]
En prenant \(K=R_W\), on prouve \eqref{eq:cpe-convergencia-norma}. La
stabilité des valeurs propres isolées d'opérateurs compacts sous des
perturbations en norme et le calcul fonctionnel de Riesz conservent leur
multiplicité. L'identité \eqref{eq:cpe-recuperacion-altura} est l'inverse
de l'application spectrale de la résolvante.
\end{proof}

Les entrées de \(R_N\) sont calculées à partir de la forme premier--gamma--polaire :
\[
 \langle R_W[f_j],[f_k]\rangle_W
 =i\int_0^\infty e^{-t}
 \mathscr I_{\rm GW}(T_{-t}f_j,f_k)\,dt.
\]
Le théorème n'utilise donc pas de table de zéros pour former les matrices.

\subsection{Certification a posteriori des intervalles et des multiplicités}

Soit \((\lambda_N,v_N)\) un vecteur propre normalisé de \(R_N\) et soit
\[
 \varepsilon_N(v_N)
 :=\|(I-P_N)R_Wv_N\|.
\]
Comme \(R_W\) est normal,
\begin{equation}
 \operatorname{dist}(\lambda_N,\sigma(R_W))
 \leq\varepsilon_N(v_N).
 \label{eq:cpe-residuo-inclusion}
\end{equation}
Les canaux premier, gamma et polaire fournissent des bornes par intervalles séparées
pour \(\varepsilon_N\) ; leurs queues sont sommées avant d'émettre l'intervalle.
De plus, si \(\Gamma\) est un contour fermé contenu dans l'ensemble résolvant de
\(R_N\) et
\begin{equation}
 \sup_{z\in\Gamma}
 \|(z-R_N)^{-1}\|\,\|R_W-R_N\|<1,
 \label{eq:cpe-criterio-riesz}
\end{equation}
alors les projections de Riesz de \(R_N\) et de \(R_W\) délimitées par
\(\Gamma\) ont le même rang. Les équations
\eqref{eq:cpe-residuo-inclusion} et \eqref{eq:cpe-criterio-riesz} certifient,
respectivement, l'inclusion et la complétude des atomes contenus dans chaque
fenêtre. L'arithmétique des intervalles consolide les décimales ; elle n'ajoute pas d'application
mathématique absente.

\section{Trace couplée du spectre de Weil et de la graduation Moonshine}
\label{sec:cpe-traza-rh-moonshine}

La trace produit
\[
 \operatorname{Tr}_{\mathcal H_W}
 \bigl(e^{-aH_W^2}\mathcal C_W^n\bigr)
 \operatorname{Tr}_{V^\natural}
 \bigl(g e^{-\beta(L_0-1)}\bigr)
\]
superpose deux observables, mais ne les couple pas. La mémoire du TPK et la
graduation de \(V^\natural\) permettent de construire le couplage sans postuler
d'action du Monstre sur des chiffres ni sur des hauteurs.

Une fois fixée la jauge exceptionnelle \(\chi_{\rm exc}\), on oriente conjointement
l'origine de la mémoire et l'origine de la graduation. Soit
\(\mathcal H_{\rm mem}=\ell^2(\mathbb N_0)\), avec
\(N_{\rm mem}|n\rangle=n|n\rangle\), et soit
\[
 V^\natural=\widehat\bigoplus_{r\ge0}V^\natural_r,
 \qquad L_0|_{V^\natural_r}=rI.
\]
Notons \(P_r^\natural\) la projection sur
\(V^\natural_r\). Le projecteur d'égalité des degrés est
\begin{equation}
 \Pi_{\chi_{\rm exc}}
 :=\sum_{n\ge0}|n\rangle\langle n|
 \otimes I_{\mathcal H_W}\otimes P_n^\natural.
 \label{eq:cpe-proyector-diagonal-grados}
\end{equation}
La somme converge fortement et définit une projection orthogonale. Sur le
produit, on définit le caractère de mémoire
\begin{equation}
 \mathcal C_W^{N_{\rm mem}}
 :=\sum_{n\ge0}|n\rangle\langle n|
 \otimes\mathcal C_W^n.
 \label{eq:cpe-cayley-memoria}
\end{equation}

\begin{definition}[Trace diagonale RH--Moonshine]
\label{def:cpe-traza-diagonal}
Pour \(a,b,\beta>0\) et \(g\in\mathbb M\), on définit
\begin{equation}
\begin{aligned}
 \Theta^{\Delta}_{g,\chi_{\rm exc}}(a,b,\beta)
 :={}&\operatorname{Tr}\!\left[
 \Pi_{\chi_{\rm exc}}
 \Bigl(
 e^{-bN_{\rm mem}}
 \mathcal C_W^{N_{\rm mem}}e^{-aH_W^2}
 \otimes g e^{-\beta(L_0-1)}
 \Bigr)
 \Pi_{\chi_{\rm exc}}
 \right].
 \label{eq:cpe-traza-diagonal-def}
\end{aligned}
\end{equation}
\end{definition}

\begin{theorem}[Existence et développement non factorisé]
\label{thm:cpe-traza-diagonal}
La trace \eqref{eq:cpe-traza-diagonal-def} est absolument convergente et
satisfait
\begin{equation}
 \boxed{\Theta^{\Delta}_{g,\chi_{\rm exc}}(a,b,\beta)
 =\sum_{n\ge0}e^{-bn}e^{-\beta(n-1)}
 \operatorname{Tr}_{\mathcal H_W}
 \bigl(e^{-aH_W^2}\mathcal C_W^n\bigr)
 \operatorname{Tr}_{V^\natural_n}(g).}
 \label{eq:cpe-traza-diagonal-expansion}
\end{equation}
Après la reconnaissance spectrale,
\begin{equation}
 \Theta^{\Delta}_{g,\chi_{\rm exc}}(a,b,\beta)
 =\sum_{n\ge0}\sum_\gamma
 m_\gamma e^{-a\gamma^2}
 e^{-bn-\beta(n-1)}\tau_\gamma^n
 \operatorname{Tr}_{V^\natural_n}(g).
 \label{eq:cpe-traza-diagonal-espectral}
\end{equation}
Cette expression n'est pas le produit d'une somme spectrale et d'une série de
McKay--Thompson : le même indice \(n\) sélectionne simultanément la puissance
du caractère de mémoire et le degré monstrueux.
\end{theorem}

\begin{proof}
Pour \(a>0\), la loi de comptage spectral implique
\(\sum_\gamma m_\gamma e^{-a\gamma^2}<\infty\). Chaque espace
\(V^\natural_n\) est de dimension finie et les coefficients des caractères
de Moonshine croissent sous-exponentiellement en \(n\). Le facteur
\(e^{-(b+\beta)n}\) assure la convergence absolue de la série conjointe.
Le développement est obtenu en insérant les résolutions spectrales de
\(N_{\rm mem}\), \(H_W\) et \(L_0\) et en utilisant
\eqref{eq:cpe-proyector-diagonal-grados}. La présence du projecteur diagonal
élimine les termes de degrés différents et produit une seule somme en
\(n\), au lieu des deux sommes indépendantes de la trace produit.
\end{proof}

Le couplage dépend de la jauge \(\chi_{\rm exc}\), qui fixe l'origine,
l'orientation et l'alignement de la lecture exceptionnelle. Un choix conjugué
produit une trace conjuguée ou une réindexation déclarée ; il ne modifie pas le
théorème de Riemann--Weil ni l'entrelacement
\eqref{eq:cpe-entrelazamiento-exacto}.



\HMTFlushCurrentChapter
\chapter{Régularité globale de Navier--Stokes par contrôle solénoïdal, transport de frontière et estimation coercive}
\begingroup
\renewcommand{\chapter}[2][]{}%
\chapter{Image solénoïdale et résolution de Navier--Stokes}

\section{Domaine hydrodynamique et publication de Galerkin}

On fixe
\[
 H^s_\sigma(\mathbb T^3)=
 \{u\in H^s(\mathbb T^3;\mathbb R^3):
   \nabla\!\cdot u=0,\ \int_{\mathbb T^3}u=0\},
 \qquad s>\frac52,\quad\nu>0,
\]
avec
$u_0\in C^\infty\cap H^s_\sigma$ et
$f\in C^\infty([0,\infty);H^\infty_\sigma)
\cap L^2_{\rm loc}([0,\infty);H^{s-1}_\sigma)$. Pour la projection de
Leray--Galerkin $P_M$, on pose
\[
 A_M=-P_M\Delta,
 \qquad B_M(v,w)=P_MP_\sigma((v\cdot\nabla)w),
 \qquad N_M(u):=B_M(u,u),
\]
et le champ fini satisfait
\begin{equation}\label{eq:pdf2-ns-galerkin}
 \partial_tu_M+\nu A_Mu_M+B_M(u_M,u_M)=f_M,
 \qquad u_M(0)=P_Mu_0.
\end{equation}
L'équation \eqref{eq:pdf2-ns-galerkin} est une publication de l'état, non sa
définition complète.

Pour que la notation ne change pas au milieu de la chaîne, on fixe une seule fois les
quatre fonctionnelles d'énergie d'ordre $s$ :
\begin{equation}\label{eq:pdf2-ns-energy-functionals}
 \begin{aligned}
 \mathscr E_{s,M}(u)&:=\frac12\|\Lambda^su\|_{L^2}^2,
 &\mathscr D_{s,M}(u)&:=\|A_M^{1/2}\Lambda^su\|_{L^2}^2,\\
 \mathscr B_{s,M}(u)&:=
   \operatorname{Re}\langle\Lambda^sN_M(u),\Lambda^su\rangle,
 &\mathscr F_{s,M}(u,f)&:=
   \operatorname{Re}\langle\Lambda^sP_Mf,\Lambda^su\rangle.
 \end{aligned}
\end{equation}
L'indice supplémentaire $j$ désigne les mêmes quatre fonctionnelles évaluées sur
la coordonnée publiée du niveau de profondeur $j$ ; il n'introduit pas d'autres
fonctionnelles.
L'équation de Galerkin implique exactement
\begin{equation}\label{eq:pdf2-ns-energy-identity}
 \frac d{dt}\mathscr E_{s,M}(u_M)
 +\nu\mathscr D_{s,M}(u_M)+\mathscr B_{s,M}(u_M)
 =\mathscr F_{s,M}(u_M,f).
\end{equation}

\section{Objet source et coordonnées effectivement construites}

L'image solénoïdale ne s'ajoute pas a posteriori. C'est exactement la construction
autonome $\mathfrak G_{\rm sol}$ produite avant ce chapitre et publiée dans
\eqref{eq:pdf2-g-sol-explicito} ; elle conserve l'état
de Galerkin, les applications de raffinement, les secteurs orientés, l'opérateur dissipatif,
l'advection projetée, la retenue, la counité $E_9$, l'inclusion $J_9$, les
projections $P_9=J_9E_9$ et $Q_{\rm micro}=I-P_9$, la mémoire, le chemin et le terminal.
En particulier,
\[
 E_9J_9=I,
 \qquad P_9^2=P_9,
 \qquad Q_{\rm micro}=I-P_9.
\]
Ces identités typent la décomposition entre publication macroscopique et mémoire
microscopique ; elles ne constituent pas à elles seules une borne globale.

À une profondeur nonadique $j$, l'état conserve
\begin{equation}\label{eq:pdf2-ns-estado}
 \Xi_{M,j}^{\rm NS}=
 \bigl(u_M;q_{A,M,j},q_{V,M,j},T_{M,j};
 \begin{gathered}
 \text{triades, feuillet, orientation, retenue, frontière,}\\
 \text{route, archive et mémoire}
 \end{gathered}
 \bigr).
\end{equation}
Le retour enrichi réunit toutes les branches avant l'évaluation :
\begin{equation}\label{eq:pdf2-ns-retorno}
 \mathcal R_{M,J,t}^{\rm enr}
 =\mathcal R_{1,M}\mathfrak U_{M,J,t}^*,
 \qquad
 \mathcal R_{M,J,t}^{\rm enr}Z_{M,J,t}=u_M(t).
\end{equation}
Une branche isolée de norme $3^J$ n'est pas le retour.

Pour $N>M$, le défaut de commutation avec la non-linéarité est la retenue
\begin{equation}\label{eq:pdf2-ns-carry}
 C_{M\leftarrow N}:=
 P_MB(u_N,u_N)-B_M(P_Mu_N,P_Mu_N).
\end{equation}
La poser égale à zéro sans prouver la commutation élimine précisément l'interaction des
hautes fréquences que le raffinement doit conserver.

Les deux mémoires positives restent séparées. Avec
$r=\pi_{\rm HMT}/\varphi_{\rm HMT}^2>1$,
\begin{equation}\label{eq:pdf2-ns-dh}
 D_M=(r-1)(\mathcal M_M^+-\mathcal M_M^-),
 \qquad
 H_M=(1+r)(\mathcal M_M^++\mathcal M_M^-),
\end{equation}
et
\[
 \dot D_M=\mathscr B_{s,M},
 \qquad
 \dot H_M=\frac{1+r}{r-1}|\mathscr B_{s,M}|.
\]
$D_M$ publie le signe ; $H_M$ conserve la grandeur positive. Les fusionner détruit
l'orientation ou la positivité.

\section{Tour PC--Fock complète et stockage terminal}

L'élévation n'est pas tronquée au degré deux :
\begin{equation}\label{eq:pdf2-ns-pcf}
 \mathcal V_M^{\rm PCF}(u)
 :=\bigoplus_{n\ge0}\frac{R_{\rm F}^n}{\sqrt{n!}}u^{\odot n},
 \qquad
 \|\mathcal V_M^{\rm PCF}(u)\|^2
 =e^{R_{\rm F}^2\|u\|_{H^s}^2}.
\end{equation}
Le rayon fait partie de la coordonnée et ne peut pas disparaître du lecteur. Pour
tout $\rho>0$, on fixe
\begin{equation}\label{eq:pdf2-ns-normalized-degree-one-reader}
 \rho^{[\rho]}_1:=\rho^{-1}\pi_1:
 \Gamma_s^\rho(V)\longrightarrow V,
 \qquad
 \rho^{[\rho]}_1\mathcal V_\rho^{\rm PCF}(u)=u.
\end{equation}
Le niveau initial utilise $\rho^{[R_{\rm F}]}_1$ et le niveau $j$ utilise
$\rho^{[r_j]}_1$ ; cette normalisation est maintenue lors du retrait de la profondeur.
Sur la diagonale cohérente, $z_{n,M}=u_M^{\odot n}$ satisfait
\begin{equation}\label{eq:pdf2-ns-carleman}
 \dot z_{n,M}=L_{n,M}z_{n,M}
 +\mathcal N_{n,n+1;M}z_{n+1,M}
 +\mathcal F_{n,n-1;M}(f_M)z_{n-1,M}.
\end{equation}
Les lecteurs de degrés un et deux sont
\[
 z_M(u)=\nu^{1/2}A_M^{1/2}\Lambda^su,
 \qquad
 w_M(u)=\nu^{-1/2}A_M^{-1/2}\Lambda^sN_M(u),
\]
et Young donne
\begin{equation}\label{eq:pdf2-ns-young-port}
 |\mathscr B_{s,M}(u)|\le
 \beta\nu\mathscr D_{s,M}(u)+\frac{\|w_M(u)\|^2}{4\beta}.
\end{equation}
L'estimation uniforme des lignes de Fourier contrôle ce lecteur quadratique dans
l'état initial. Elle ne contrôle pas encore ses propagations futures.

Soit
\[
 \mathcal N_{M,j}:=
 \frac{r-1}{r+1}H_{M,j}-D_{M,j}^{\rm mem}
 =2(r-1)M_{M,j}^-\ge0.
\]
Sur un horizon $T$, le stockage terminal est
\begin{equation}\label{eq:pdf2-ns-storage}
 \begin{aligned}
 G_{M,j,T}(t)
 &=\mathscr E_{s,M,j}(t)+2(r-1)(M^-_{M,j}(T)-M^-_{M,j}(t)),\\
 \dot G_{M,j,T}+\nu\mathscr D_{s,M,j}+|\mathscr B_{s,M,j}|
 &=\mathscr F_{s,M,j}.
 \end{aligned}
\end{equation}
Pour typer le terminal, on conserve dans l'état solénoïdal complet la coordonnée
de feuillet
\[
 q^-_{M,j}:=\sqrt{2(r-1)\dot M^-_{M,j}}
 \in L^2(0,T;\mathbb R).
\]
Soit $\mathscr X^{\rm sol}_{M,j,T}$ l'espace d'histoire enrichie qui contient,
comme termes orthogonaux, la trace $2^{-1/2}\Lambda^su_{M,j}$ et $q^-_{M,j}$.
La colonne terminale est l'opérateur de coordonnées
\begin{equation}\label{eq:pdf2-ns-terminal-column}
 \begin{aligned}
 \mathcal E_{M,j,T}^{\rm term}(t):\mathscr X^{\rm sol}_{M,j,T}
 &\longrightarrow L^2_\sigma(\mathbb T^3)
   \oplus L^2(t,T;\mathbb R),\\
 \mathcal E_{M,j,T}^{\rm term}(t)\Xi
 &:=2^{-1/2}\Lambda^su_{M,j}(t)
 \oplus\mathbf1_{[t,T]}q^-_{M,j}.
 \end{aligned}
\end{equation}
Par construction des deux coordonnées orthogonales,
\begin{equation}\label{eq:pdf2-ns-terminal-column-gram}
 \|\mathcal E_{M,j,T}^{\rm term}(t)\Xi\|^2
 =G_{M,j,T}(t).
\end{equation}
Le terme de mémoire future est une sortie terminale typée, non une composante de la
racine initiale ; son coût ne sera contrôlé qu'après le télescopage de la colonne
complète. Pour une observation $\tau$, les lignes de perte vivent dans
$[0,\tau]$ et cette mémoire terminale dans $(\tau,T]$ ; les supports temporels sont
disjoints. Dans l'estimation intégrale, on prend $\tau=T$, et alors le terme
futur est nul.

\section{Construction source--first avec relèvement mixte état--force}

L'entrée est fixée avant toute coupure, profondeur et histoire :
\begin{equation}\label{eq:pdf2-ns-data-space}
 \mathcal D_T:=H^s_\sigma(\mathbb T^3)
 \times L^2(0,T;H^{s-1}_\sigma(\mathbb T^3)),
 \qquad d=(u_0,f).
\end{equation}
Comme les champs sont de moyenne nulle, $A=-\Delta$ est positif et inversible dans le
secteur solénoïdal. On introduit l'effort, le flux visqueux et les deux ondes
\begin{equation}\label{eq:pdf2-ns-force-waves}
 \begin{aligned}
 a_T(f)&:=\frac{1}{2\sqrt\nu}A^{-1/2}\Lambda^sP_\sigma f,
 &z_M(u)&:=\sqrt\nu A_M^{1/2}\Lambda^su,\\
 b_M(u,f)&:=z_M(u)-P_Ma_T(f),\\
 q_M^{B,+}(u)&:=\bigl(\mathscr B_{s,M}(u)\bigr)_+^{1/2},
 &q_M^{B,-}(u)&:=\bigl(-\mathscr B_{s,M}(u)\bigr)_+^{1/2},\\
 q_M^B(u)&:=\bigl(q_M^{B,+}(u),q_M^{B,-}(u)\bigr).
 \end{aligned}
\end{equation}
Tous appartiennent à $L^2$ sur l'intervalle correspondant et
$\|q_M^B(u)\|_{\mathbb C^2}^2=|\mathscr B_{s,M}(u)|$.
La puissance croisée est typée par
\begin{equation}\label{eq:pdf2-ns-cross-power}
 \mathscr F_{s,M}=2\operatorname{Re}\langle z_M,P_Ma_T(f)\rangle,
 \qquad \nu\mathscr D_{s,M}=\|z_M\|^2.
\end{equation}
C'est donc \eqref{eq:pdf2-ns-storage}, et non une substitution de la puissance par la
norme de la force, qui donne l'identité exacte de diffusion
\begin{equation}\label{eq:pdf2-ns-scattering-balance}
 \dot G_{M,j,T}+\|b_M\|^2+\|q^B_M\|^2=\|P_Ma_T(f)\|^2.
\end{equation}

\subsection{Signature chronologique, opérateur régularisant uniforme et création mixte}

La donnée de force vit dans $F:=H^{s-1}_\sigma(\mathbb T^3)$, tandis que la
tour PC--Fock utilise comme espace à une particule
$V:=H^s_\sigma(\mathbb T^3)$. L'identification entre les deux n'est pas laissée à
l'inclusion de Galerkin, dont la norme croîtrait avec $M$. Dans le secteur de moyenne nulle, on
fixe l'isomorphisme d'ordre moins un
\begin{equation}\label{eq:pdf2-ns-uniform-smoother}
 \begin{aligned}
 \mathscr S&:=A^{-1/2}:F\longrightarrow V,\\
 \mathscr S_M&:=A_M^{-1/2}P_M=P_M\mathscr S,\\
 \sup_M\bigl(\|\mathscr S_M\|_{F\to V}
 +\|\mathscr S_M^{-1}\|_{P_MV\to P_MF}\bigr)&<\infty.
 \end{aligned}
\end{equation}
La dernière borne est l'équivalence uniforme des normes de Sobolev dans le
sous-espace sans mode zéro. Ainsi, $g:=\mathscr Sf\in L^2(0,T;V)$ et on n'utilise pas
l'inclusion non uniforme $P_MH^{s-1}\hookrightarrow P_MH^s$.

Soit $\Delta_k(T)=\{0<t_1<\cdots<t_k<T\}$. L'espace de signature de la force
régularisée est
\begin{equation}\label{eq:pdf2-ns-signature-space}
 \mathfrak S_T(V):=\bigoplus_{k\geq0}L^2(\Delta_k(T);V^{\widehat\otimes k}),
 \qquad \Delta_0(T)=\{*\}.
\end{equation}
Pour $g=\mathscr Sf\in L^2(0,T;V)$, on définit, sans utiliser la solution,
\begin{equation}\label{eq:pdf2-ns-force-signature}
 \operatorname{Sig}_T(g)_0=1,
 \qquad
 \operatorname{Sig}_T(g)_k(t_1,\ldots,t_k)
 =g(t_k)\widehat\otimes\cdots\widehat\otimes g(t_1).
\end{equation}
La symétrie de l'intégrande scalaire fournit
\begin{equation}\label{eq:pdf2-ns-signature-norm}
 \|\operatorname{Sig}_T(g)\|^2
 =\sum_{k\geq0}\frac{\|g\|_{L^2(0,T;V)}^{2k}}{k!}
 =e^{\|g\|_{L^2(0,T;V)}^2}
 \le e^{c_{\mathscr S}^2\|f\|_{L^2(0,T;F)}^2}.
\end{equation}

Le terme forcé de \eqref{eq:pdf2-ns-carleman} n'est pas un port additif dans
$f$ uniquement. Pour $0<r_{\rm F}<R_{\rm F}$, il est typé au moyen de la création mixte
\begin{equation}\label{eq:pdf2-ns-mixed-creation}
 \begin{aligned}
 \widetilde{\mathcal F}^{R_{\rm F},r_{\rm F}}_M:
 P_MV\widehat\otimes\Gamma_s^{R_{\rm F}}(P_MV)
 &\longrightarrow\Gamma_s^{r_{\rm F}}(P_MV),\\
 \widetilde{\mathcal F}^{R_{\rm F},r_{\rm F}}_M(g\otimes x)
 &:=a^\dagger(g)x.
 \end{aligned}
\end{equation}
Au degré $n$, son action est précisément
$g\otimes z_{n-1}\mapsto n g\odot z_{n-1}$. On conserve ainsi le facteur mixte
$\mathscr Sf\otimes x$ que ne peut produire une transition linéaire sur la somme directe
``état $\oplus$ force''.
Avec les normes pondérées de \eqref{eq:pdf2-ns-pcf},
\begin{equation}\label{eq:pdf2-ns-creation-uniform}
 \|\widetilde{\mathcal F}^{R,r}_M(g\otimes x)\|_{\Gamma^r}
 \le c(R,r)\|g\|_V\|x\|_{\Gamma^R},
 \qquad
 c(R,r):=\sup_{n\ge1}\sqrt n\,r\left(\frac rR\right)^{n-1}<\infty.
\end{equation}
Pour itérer, on fixe $0<\theta<1$ et $r_j=R_{\rm F}\theta^j$ ; chaque pas utilise
$(R,r)=(r_j,r_{j+1})$. Alors
$\sup_{j,M}c(r_j,r_{j+1})\le
R_{\rm F}\sup_{n\ge1}\sqrt n\,\theta^n<\infty$.

La tour s'écrit dans la coordonnée conjuguée
$v_M=\mathscr S_Mu_M$. Soit
$\mathscr S_M^{\Gamma}:=\bigoplus_{n\ge0}
\mathscr S_M^{\widehat\otimes n}$. Le générateur homogène conjugué est
\begin{equation}\label{eq:pdf2-ns-conjugated-generator}
 \mathcal A^{0,\mathscr S}_{M}
 :=\mathscr S_M^{\Gamma}\mathcal A^0_M
   (\mathscr S_M^{\Gamma})^{-1},
 \qquad
 (\mathcal A^{0,\mathscr S}_{M})_{n,m}
 =\mathscr S_M^{\widehat\otimes n}
  (\mathcal A^0_M)_{n,m}
  (\mathscr S_M^{-1})^{\widehat\otimes m}.
\end{equation}
Cette conjugaison est source--first et n'altère pas l'ODE : le lecteur de degré un
applique $\mathscr S_M^{-1}$ à la fin.

Soit $\mathcal T_M^{0,\mathscr S}(t,s)$ le propagateur de la tour conjuguée, inclus
dans la réalisation solénoïdale complète. Pour $R>r>0$, il est d'abord défini sur
\begin{equation}\label{eq:pdf2-ns-radial-transfer}
 \delta_{R\to r}(x_0,x_1,x_2,\ldots)
 :=\bigl(x_0,(r/R)x_1,(r/R)^2x_2,\ldots\bigr),
 \qquad
 \delta_{R\to r}\mathcal V_R^{\rm PCF}(u)=\mathcal V_r^{\rm PCF}(u).
\end{equation}
Cet opérateur est contractif et fixe le vide. Sur
\begin{equation}\label{eq:pdf2-ns-dyson-domain}
 \Gamma_s^R(P_MV)_{\rm alg}\widehat\otimes
 \mathfrak S_T(P_MV)_{\rm fin}
\end{equation}
Soit $\mathfrak S_{T,{\rm fin}}^{\rm dec}$ l'espace engendré par les noyaux élémentaires
\[
 g_k(t_1,\ldots,t_k)=
 g_k^{(k)}(t_k)\widehat\otimes\cdots\widehat\otimes g_k^{(1)}(t_1).
\]
Il est dense dans chaque terme $L^2(\Delta_k(T);V^{\widehat\otimes k})$. Dans la
formule suivante, les symboles $g_k^{(i)}$ sont lus dans chaque noyau élémentaire et
l'action s'étend linéairement à leur espace engendré ; on ne factorise donc pas
arbitrairement un tenseur intriqué. Le lecteur chronologique prend ses valeurs dans
$\Gamma_s^r(P_MV)$ et est défini par
\begin{equation}\label{eq:pdf2-ns-dyson-reader}
 \begin{aligned}
 \mathscr D_{M,k}^{R\to r}(t)(x\otimes g_k)
 :={}&\delta_{R\to r}\int_{\Delta_k(t)}
 \mathcal T_M^{0,\mathscr S}(t,t_k)
 \widetilde{\mathcal F}_M(P_Mg_k^{(k)}(t_k)\otimes\cdot)\cdots\\
 &\hspace{18mm}\cdots\mathcal T_M^{0,\mathscr S}(t_2,t_1)
 \widetilde{\mathcal F}_M(P_Mg_k^{(1)}(t_1)\otimes\cdot)
 \mathcal T_M^{0,\mathscr S}(t_1,0)x\,d\mathbf t,\\
 \mathscr D_{M,T}^{R\to r}(t)&:=\sum_{k\geq0}
 \mathscr D_{M,k}^{R\to r}(t).
 \end{aligned}
\end{equation}
En répartissant la perte totale $R-r$ entre les $k$ créations et les propagateurs,
l'estimation analytique d'échelle de
\eqref{eq:pdf2-ns-creation-uniform} donne, pour des constantes finies dépendant de la
coupure mais non du noyau élémentaire,
\begin{equation}\label{eq:pdf2-ns-dyson-homogeneous-bound}
 \|\mathscr D_{M,k}^{R\to r}(t)(x\otimes g_k)\|_{\Gamma^r}
 \le A_{M,R,r,T}\,
 \frac{B_{M,R,r,T}^{k}}{\sqrt{\Gamma(k+1)}}
 \|x\|_{\Gamma^R}\|g_k\|_{L^2(\Delta_k;V^{\widehat\otimes k})}.
\end{equation}
Le facteur $1/\Gamma(k+1)$ du volume du simplexe et Cauchy--Schwarz dans la
somme des degrés donnent
\begin{equation}\label{eq:pdf2-ns-dyson-bound}
 \|\mathscr D_{M,T}^{R\to r}(t)(x\otimes\sigma)\|_{\Gamma^r}
 \le C_{M,R,r,T}\|x\|_{\Gamma^R}\|\sigma\|_{\mathfrak S_T},
 \qquad C_{M,R,r,T}<\infty.
\end{equation}
La série converge donc dans le codomaine indiqué et détermine par densité
un opérateur borné pour chaque coupure. On n'utilise pas ici de borne uniforme en $M$ :
cette uniformité sera obtenue après l'identité locale de la colligation.
En l'évaluant sur
$x\otimes\operatorname{Sig}_T(\mathscr Sf)$, on obtient la série de Dyson du générateur non
autonome
\begin{equation}\label{eq:pdf2-ns-nonautonomous-generator}
 \mathcal A_M^{f,\mathscr S}(t)
 =\mathcal A_M^{0,\mathscr S}+a^\dagger(\mathscr S_MP_Mf(t)),
 \qquad
 \partial_tU_M^{f,\mathscr S}(t,s)
 =\mathcal A_M^{f,\mathscr S}(t)U_M^{f,\mathscr S}(t,s).
\end{equation}
La différentiation de \eqref{eq:pdf2-ns-dyson-reader} sur le cœur algébrique donne
la conjuguée de \eqref{eq:pdf2-ns-carleman} ; l'unicité de l'ODE finie donne
\begin{equation}\label{eq:pdf2-ns-dyson-degree-one}
 \mathscr S_M^{-1}\rho^{[r]}_1
 \mathscr D_{M,T}^{R_{\rm F}\to r}(t)
 \bigl(\mathcal V_M^{\rm PCF}(\mathscr S_MP_Mu_0)
       \otimes\operatorname{Sig}_T(\mathscr S_MP_Mf)\bigr)
 =u_M(t).
\end{equation}
L'égalité inclut $t=0$ : le facteur $r^{-1}$ de
\eqref{eq:pdf2-ns-normalized-degree-one-reader} annule exactement la
coordonnée radiale $r\mathscr S_MP_Mu_0$ produite par
\eqref{eq:pdf2-ns-radial-transfer}. On n'impose pas $R_{\rm F}=1$.
Cette construction précède l'histoire et n'invoque pas
$\mathsf S_{M,T}d$ pour la définir.

Soit
$\mathfrak s_{\rm sol}:H^s_\sigma\to\mathcal F_{\rm seed}^{\rm sol}$ la
restriction solénoïdale de l'application de germe : elle conserve feuillet, orientation, résidu,
quotient et chemin basé et dépend uniquement de $u_0$. La racine directrice n'est pas
redéfinie ici : c'est exactement la racine $J_T^{\rm sol}$ de
\eqref{eq:pdf2-g-sol-raiz-dato}. Nous fixons l'identité des noms
\begin{equation}\label{eq:pdf2-ns-root-identification}
 J_T^{\rm mix}:=J_T^{\rm sol}:
 \mathcal D_T=\mathcal D_T^{\rm NS}\longrightarrow\mathcal H_T^{\rm sol}.
\end{equation}
Sa publication de coordonnées mixtes, et non une deuxième racine, est
\begin{equation}\label{eq:pdf2-ns-root-map}
 \begin{aligned}
 \operatorname{Coord}_T^{\rm mix}(u_0,f)&:=
 \bigl[\mathcal V^{\rm PCF}(u_0)
 \widehat\otimes\operatorname{Sig}_T(\mathscr Sf)\bigr]
 \oplus\mathfrak s_{\rm sol}(u_0),\\
 \operatorname{Coord}_T^{\rm mix}&:\mathcal D_T\longrightarrow
 \mathfrak H_T^{\rm coord},\\
 \mathfrak H_T^{\rm coord}&:=
 \bigl[\Gamma_s^{R_{\rm F}}(V)
 \widehat\otimes\mathfrak S_T(V)\bigr]
 \oplus\mathcal F_{\rm seed}^{\rm sol}.
 \end{aligned}
\end{equation}
La carte est une projection de coordonnées de l'état accepté. Sur le sous-espace
cyclique engendré par la racine, on écrit
\begin{equation}\label{eq:pdf2-ns-mixed-coordinate-chart}
 \chi_T^{\rm mix}:\mathcal H_T^{\rm sol}\longrightarrow
 \mathfrak H_T^{\rm coord},
 \qquad
 \chi_T^{\rm mix}J_T^{\rm sol}(u_0,f)
 =\operatorname{Coord}_T^{\rm mix}(u_0,f).
\end{equation}
Son action hors du sous-espace cyclique est la projection orthogonale de
$\mathfrak G_{\rm sol}$ sur la coordonnée tour--signature--germe. On ne crée pas
un nouveau quotient et on n'identifie pas la norme de cette carte à la norme totale
de $\mathcal H_T^{\rm sol}$.
Avec la normalisation du germe,
\begin{equation}\label{eq:pdf2-ns-root-norm}
 \|\operatorname{Coord}_T^{\rm mix}(u_0,f)\|^2
 \leq e^{R_{\rm F}^2\|u_0\|_{H^s}^2
             +c_{\mathscr S}^2\|f\|_{L^2(0,T;H^{s-1})}^2}
 +c_{\rm seed}(1+\|u_0\|_{H^s}^2).
\end{equation}
Cette estimation prouve que la carte de coordonnées se forme uniquement avec la donnée ;
elle ne transporte pas de nouvelle norme vers $\mathcal H_T^{\rm sol}$ et n'est pas utilisée pour fabriquer
la borne directrice. La norme et la colonne de $J_T^{\rm mix}=J_T^{\rm sol}$ sont
celles déjà fixées par $\mathfrak G_{\rm sol}$.

\subsection{Donnée unique de départ et noyaux denses}

L'unique donnée admise dans ce chapitre est l'image complète
$\mathfrak G_{\rm sol}$ reproduite dans
\eqref{eq:pdf2-g-sol-explicito}--\eqref{eq:pdf2-g-sol-cota-dato}. En particulier,
on utilise sans les renommer la racine $J_T^{\rm sol}$, le codeur
$\Lambda_{M,T}$, les transitions $\Gamma_{M,j}^{\rm term}$, la ligne complète
$L_{M,j,T}^{\rm term}$, le terminal $\mathcal E_{M,j,T}^{\rm term}$, la colonne
$\mathbf G_{M,j;J,T}^{\rm term}$ et l'isométrie $I_{M,J}$. On n'introduit pas de
deuxième réalisation de ces objets.

Soit $\mathcal D_{T,{\rm alg}}\subset\mathcal D_T$ le sous-espace de données avec
$u_0$ trigonométrique et $f$ lisse, à support temporel compact et à valeurs
trigonométriques. Il est dense dans $\mathcal D_T$. Les noyaux sur lesquels
s'écrivent toutes les actions sont
\begin{equation}\label{eq:pdf2-ns-cyclic-cores}
 \begin{aligned}
 \mathscr C_T^{\rm src}
 &:=\operatorname{span}\{J_T^{\rm sol}d:d\in\mathcal D_{T,{\rm alg}}\}
 \subset\mathcal H_T^{\rm sol},\\
 \mathscr C_{M,j,T}^{\rm sol}
 &:=\operatorname{span}\{x_{M,j,T}(d):d\in\mathcal D_{T,{\rm alg}}\},\\
 x_{M,0,T}(d)&:=\Lambda_{M,T}J_T^{\rm sol}d,\\
 x_{M,j+1,T}(d)&:=\Gamma_{M,j}^{\rm term}x_{M,j,T}(d).
 \end{aligned}
\end{equation}
Les espaces $\mathscr X_{M,j,T}^{\rm sol}$ sont les complétions de ces
noyaux dans les normes déjà contenues dans $\mathfrak G_{\rm sol}$. En particulier,
aucune norme n'est définie en soustrayant la perte que l'on veut ensuite borner.

L'identification exacte entre la notation de l'application et celle de la donnée est
\begin{equation}\label{eq:pdf2-ns-exact-datum-identification}
 \boxed{\begin{gathered}
 J_T^{\rm mix}=J_T^{\rm sol},\qquad
 \Gamma_{M,j,T}=\Gamma_{M,j}^{\rm term},\qquad
 L_{M,j,T}=L_{M,j,T}^{\rm term},\\
 \mathcal G_{M,J,T}^{\rm sol}=\mathbf G_{M,0;J,T}^{\rm term}
 \end{gathered}}
\end{equation}
Cette égalité est littérale, non une ressemblance de noms. L'action source--first
est donc
\begin{equation}\label{eq:pdf2-ns-source-encoder-action}
 \boxed{\begin{gathered}
 d\xmapsto{J_T^{\rm sol}}J_T^{\rm sol}d
 \xmapsto{\Lambda_{M,T}}x_{M,0,T}(d)
 \xmapsto{\Gamma_{M,0:j,T}}x_{M,j,T}(d)
 \end{gathered}},
 \qquad
 \Gamma_{M,0:j,T}:=\Gamma_{M,j-1,T}\cdots\Gamma_{M,0,T}.
\end{equation}
Aucune de ces trois flèches ne reçoit $u_M$, une perte ou un terminal comme
entrée. La carte PC--Fock--signature de
\eqref{eq:pdf2-ns-root-map} publie les coordonnées de la première flèche ;
l'égalité de degré un
\eqref{eq:pdf2-ns-dyson-degree-one} prouve, sur le noyau puis par densité,
\begin{equation}\label{eq:pdf2-ns-sol-transition-action}
 \pi_{M,j,T}^{(1)}x_{M,j,T}(u_0,f)=u_M(t_j),
 \qquad
 \pi_{M,j,T}^{(1)}:=
 \mathscr S_M^{-1}
 (\rho_1^{[r_j]}\widehat\otimes\varepsilon_{[t_j,T]})\pi_{M,j,T}^{\rm tf}.
\end{equation}
Le facteur $r_j^{-1}$ est appliqué avant $\mathscr S_M^{-1}$ et annule le rayon
de degré un. Le lecteur renvoie donc $u_M$, non $r_ju_M$ ni
$\mathscr S_Mu_M$.

Pour $N\ge M$, la projection de raffinement agit tensoriellement sur la
carte de coordonnées :
\begin{equation}\label{eq:pdf2-ns-refinement-action}
 \begin{aligned}
 \mathfrak p_{M\leftarrow N}^{\Gamma}
 (v_1\widehat\otimes\cdots\widehat\otimes v_k)
 &:=(P_Mv_1)\widehat\otimes\cdots\widehat\otimes(P_Mv_k),\\
 \mathfrak p_{M\leftarrow N}^{\mathfrak S}
 (g_1\widehat\otimes\cdots\widehat\otimes g_k)
 &:=(P_Mg_1)\widehat\otimes\cdots\widehat\otimes(P_Mg_k).
 \end{aligned}
\end{equation}
Avec la projection de germe de $\mathfrak G_{\rm sol}$, ces formules
donnent l'application $\mathfrak p_{M\leftarrow N}^{\rm sol}$ et les carrés
\begin{equation}\label{eq:pdf2-ns-refinement-commutation}
 \begin{aligned}
 \mathfrak p_{M\leftarrow N}^{\rm sol}\Lambda_{N,T}J_T^{\rm sol}
 &=\Lambda_{M,T}J_T^{\rm sol},\\
 \mathfrak p_{M\leftarrow N}^{\rm sol}\Gamma_{N,j,T}
 -\Gamma_{M,j,T}\mathfrak p_{M\leftarrow N}^{\rm sol}
 &=\iota_{\rm carry}\mathcal C_{M\leftarrow N,j},\\
 \mathcal C_{M\leftarrow N,j}x_{N,j,T}(d)
 &=C_{M\leftarrow N}(d)|_{I_j}.
 \end{aligned}
\end{equation}
Le membre droit est l'opérateur de coordonnée dont la valeur publiée est
exactement la retenue de
\eqref{eq:pdf2-ns-carry} ; on ne le déclare pas nul pour fabriquer la naturalité.

\subsection{Les six lignes comme composantes de l'unique opérateur de perte}

Fixons $0=t_0<t_1<\cdots<t_J=\tau\le T$, $I_j=[t_j,t_{j+1}]$ et le calendrier
\begin{equation}\label{eq:pdf2-ns-cutoff-calendar}
 M_j:=3^jM,
 \qquad P_{M_j}\longrightarrow I
 \quad\hbox{fortement dans chaque }H^r_\sigma.
\end{equation}
Les six codomaines, avec leurs normes propres, sont
\begin{equation}\label{eq:pdf2-ns-six-loss-spaces}
 \begin{aligned}
 \mathcal Y_{M,j}^{\rm diss}&=L^2(I_j;L^2_\sigma),&
 \mathcal Y_{M,j}^{\rm adv}&=L^2(I_j;\mathbb C^2),\\
 \mathcal Y_{M,j}^{\rm carry}&=L^2(I_j;H^{s-1}_\sigma),&
 \mathcal Y_{M,j}^{\rm micro}&=L^2(I_j;\operatorname{Ran}Q_{\rm micro}),\\
 \mathcal Y_{M,j}^{\rm shell}
 &=L^2(I_j;\operatorname{Ran}(P_{M_{j+1}}-P_{M_j})\cap H^{s-1}_\sigma),&
 \mathcal Y_{M,j}^{\rm mem}
 &=L^2(I_j;\operatorname{Ran}Q^{\rm mem}_{M,j+1}).
 \end{aligned}
\end{equation}
On pose
$\mathcal Y_{M,j}^{\rm sol}:=\widehat\bigoplus_{\mathfrak a}
\mathcal Y_{M,j}^{\mathfrak a}$, où
$\mathfrak a\in\{{\rm diss,adv,carry,micro,shell,mem}\}$, et l'on note
$P_{\mathfrak a}$ sa projection orthogonale. La ligne acceptée de
$\mathfrak G_{\rm sol}$ se décompose, sans être dupliquée, au moyen de
\begin{equation}\label{eq:pdf2-ns-loss-parseval}
 L_{M,j,T}^{\rm term}:
 \mathscr X_{M,j,T}^{\rm sol}\longrightarrow\mathcal Y_{M,j}^{\rm sol},
 \qquad
 L_{M,j,T}^{\mathfrak a}:=P_{\mathfrak a}L_{M,j,T}^{\rm term}.
\end{equation}

Pour écrire les actions sans transformer des fonctions non linéaires de $u$ en
opérateurs linéaires fictifs, on utilise les projections de coordonnées du relèvement
mixte déjà contenu dans $\mathfrak G_{\rm sol}$ :
\begin{equation}\label{eq:pdf2-ns-row-coordinate-readers}
 \begin{aligned}
 \pi_zx_{M,j,T}(d)&=\mathbf1_{I_j}\sqrt\nu A_M^{1/2}\Lambda^su_M,&
 \pi_ax_{M,j,T}(d)&=\mathbf1_{I_j}P_Ma_T(f),\\
 \pi_{B,\pm}x_{M,j,T}(d)&=\mathbf1_{I_j}q_M^{B,\pm}(u_M),&
 \pi_{\rm carry}x_{M,j,T}(d)&=\mathbf1_{I_j}C_{M_j\leftarrow M_{j+1}},\\
 \pi_{\rm micro}x_{M,j,T}(d)&=\mathbf1_{I_j}Q_{\rm micro}y^{\rm micro}_{M,j},&
 \pi_{\rm shell}x_{M,j,T}(d)&=\mathbf1_{I_j}
 (P_{M_{j+1}}-P_{M_j})N_{M_{j+1}}(u_{M_{j+1}}),\\
 \pi_{\rm mem}x_{M,j,T}(d)&=\mathbf1_{I_j}
 Q^{\rm mem}_{M,j+1}m^{\rm raw}_{M,j+1}.&&
 \end{aligned}
\end{equation}
Ici, $y^{\rm micro}_{M,j}$ et $m^{\rm raw}_{M,j+1}$ ne sont pas des symboles
externes : ce sont respectivement les projections sur les coordonnées
$Q_{\rm micro}$ et $\operatorname{Mem}_{M,j+1}$ de l'état
\eqref{eq:pdf2-g-sol-explicito}. Les huit applications $\pi$ sont linéaires sur
l'état relevé, bien que leurs publications sur la diagonale physique soient
polynomiales ou contiennent une racine positive. Leur domaine est le noyau
\eqref{eq:pdf2-ns-cyclic-cores} ; leur extension bornée est la composante
correspondante de $L_{M,j,T}^{\rm term}$.

Avec ces projections, les six actions sont matériellement
\begin{equation}\label{eq:pdf2-ns-six-loss-actions}
 \begin{aligned}
 L^{\rm diss}_{M,j,T}&=\pi_z-\pi_a,&
 L^{\rm adv}_{M,j,T}&=(\pi_{B,+},\pi_{B,-}),\\
 L^{\rm carry}_{M,j,T}&=\pi_{\rm carry},&
 L^{\rm micro}_{M,j,T}&=\pi_{\rm micro},\\
 L^{\rm shell}_{M,j,T}&=\pi_{\rm shell},&
 L^{\rm mem}_{M,j,T}&=\pi_{\rm mem}.
 \end{aligned}
\end{equation}
Le caractère borné ne se déduit pas des noms. Pour
$x\in\mathscr C_{M,j,T}^{\rm sol}$, la colonne acceptée donne
\begin{equation}\label{eq:pdf2-ns-row-boundedness}
 \|L_{M,j,T}^{\mathfrak a}x\|^2
 \le\|L_{M,j,T}^{\rm term}x\|^2
 \le\|\mathbf G_{M,j;J,T}^{\rm term}x\|^2.
\end{equation}
Par densité, chaque ligne possède une unique extension bornée à
$\mathscr X_{M,j,T}^{\rm sol}$. On ne forme plus de nouveau quotient et on ne définit pas
la norme du domaine à partir de la ligne.

L'absence de surcomptage est l'identité de Parseval de l'unique opérateur
$L^{\rm term}$ :
\begin{equation}\label{eq:pdf2-ns-six-loss-no-overcount}
 \boxed{\begin{gathered}
 \|L_{M,j,T}^{\rm term}x\|^2
 =\sum_{\mathfrak a}\|L_{M,j,T}^{\mathfrak a}x\|^2,
 \qquad x\in\mathscr X_{M,j,T}^{\rm sol}
 \end{gathered}}
\end{equation}
Cette orthogonalité coïncide avec la géométrie physique : les intervalles $I_j$
sont disjoints ; carry vit sous $P_{M_j}$ et shell sous
$P_{M_{j+1}}-P_{M_j}$ ; micro vit sous $Q_{\rm micro}$ ; et
$Q^{\rm mem}_{M,j+1}$ élimine la mémoire héritée et la copie du canal
advectif. De plus, $q_M^{B,+}q_M^{B,-}=0$ ponctuellement et
$\|(q_M^{B,+},q_M^{B,-})\|^2=|\mathscr B_{s,M}|$, de sorte que les deux feuillets
sont un seul canal orienté.

\subsection{Terminal, colonne et Stein dérivés sans redéfinir les normes}

Sur le noyau terminal, l'action de l'opérateur accepté est
\begin{equation}\label{eq:pdf2-ns-terminal-action-on-core}
 \mathcal E_{M,J,T}^{\rm term}(\tau)x_{M,J,T}(d)
 =2^{-1/2}\Lambda^su_M(\tau)
 \oplus\mathbf1_{[\tau,T]}q^-_{M,J}.
\end{equation}
Son extension bornée à $\mathscr X_{M,J,T}^{\rm sol}$ fait partie de
$\mathfrak G_{\rm sol}$ et son Gram terminal est, sans compléments inventés,
\begin{equation}\label{eq:pdf2-ns-terminal-metric}
 U_{M,J,T}:=
 (\mathcal E_{M,J,T}^{\rm term})^*\mathcal E_{M,J,T}^{\rm term},
 \qquad
 \|\mathcal E_{M,J,T}^{\rm term}x_{M,J,T}(d)\|^2
 =G_{M,J,T}(\tau).
\end{equation}
Le lecteur de la première coordonnée terminale est défini sur tout son
codomaine par
\begin{equation}\label{eq:pdf2-ns-terminal-reader}
 \mathcal R_{M,\tau}^{(1)}(h\oplus q):=
 \sqrt2\,\Lambda^{-s}h,
 \qquad
 \mathcal R_{M,\tau}^{(1)}
 \mathcal E_{M,J,T}^{\rm term}(\tau)x_{M,J,T}(d)=u_M(\tau).
\end{equation}
On n'inverse pas $\Gamma$, on ne choisit pas de représentant d'un quotient et il n'apparaît pas
d'opérateur terminal complémentaire sans action.

La colonne terminale est d'abord fixée par
\begin{equation}\label{eq:pdf2-ns-terminal-tail-column}
 \mathbf G_{M,J;J,T}^{\rm term}:=\mathcal E_{M,J,T}^{\rm term}.
\end{equation}
Pour $0\le j<J$, la colonne de queue s'écrit avec toutes ses lignes, dans
l'ordre terminal puis profondeur décroissante :
\begin{equation}\label{eq:pdf2-ns-complete-column}
 \mathbf G_{M,j;J,T}^{\rm term}:=
 \begin{bmatrix}
  \mathcal E_{M,J,T}^{\rm term}\Gamma_{M,j:J,T}\\
  L_{M,J-1,T}^{\rm term}\Gamma_{M,j:J-1,T}\\
  \vdots\\
  L_{M,j+1,T}^{\rm term}\Gamma_{M,j:j+1,T}\\
  L_{M,j,T}^{\rm term}
 \end{bmatrix},
 \qquad
 \Gamma_{M,j:k,T}:=\Gamma_{M,k-1,T}\cdots\Gamma_{M,j,T}.
\end{equation}
Pour $j<J$, il existe une égalité d'opérateurs entre colonnes déjà construites,
\begin{equation}\label{eq:pdf2-ns-tail-column-recursion}
 \mathbf G_{M,j;J,T}^{\rm term}
 =\begin{bmatrix}
   \mathbf G_{M,j+1;J,T}^{\rm term}\Gamma_{M,j,T}\\
   L_{M,j,T}^{\rm term}
  \end{bmatrix}.
\end{equation}
Maintenant, et seulement maintenant, on prend le Gram
\begin{equation}\label{eq:pdf2-ns-tail-gram}
 U_{M,j,T}:=
 (\mathbf G_{M,j;J,T}^{\rm term})^*
 \mathbf G_{M,j;J,T}^{\rm term}.
\end{equation}
En multipliant l'égalité de blocs
\eqref{eq:pdf2-ns-tail-column-recursion} par son adjointe, on obtient Stein :
\begin{equation}\label{eq:pdf2-ns-local-stein}
 \boxed{\begin{gathered}
 U_{M,j,T}=\Gamma_{M,j,T}^*U_{M,j+1,T}\Gamma_{M,j,T}
 +(L_{M,j,T}^{\rm term})^*L_{M,j,T}^{\rm term}
 \end{gathered}}
\end{equation}
Par \eqref{eq:pdf2-ns-six-loss-no-overcount}, le dernier terme est
$\sum_{\mathfrak a}(L_{M,j,T}^{\mathfrak a})^*
L_{M,j,T}^{\mathfrak a}$. L'équation de Stein est donc une conséquence de
la colonne acceptée ; elle ne définit pas la norme de l'état et ne sert pas à justifier
rétroactivement l'existence de ses lignes.

Le télescopage est de même une multiplication de la colonne complète :
\begin{equation}\label{eq:pdf2-ns-global-column-isometry}
 \|\mathbf G_{M,0;J,T}^{\rm term}x\|^2
 =\|\mathcal E_{M,J,T}^{\rm term}\Gamma_{M,0:J,T}x\|^2
 +\sum_{j<J}\|L_{M,j,T}^{\rm term}
                  \Gamma_{M,0:j,T}x\|^2.
\end{equation}
Tous les termes à droite existent avant cette égalité et portent les
normes de leurs codomaines dans \eqref{eq:pdf2-ns-six-loss-spaces}.

\subsection{Factorisation source--histoire et borne uniforme}

La factorisation acceptée dans
\eqref{eq:pdf2-g-sol-raiz-dato} devient, après l'identification littérale
\eqref{eq:pdf2-ns-exact-datum-identification},
\begin{equation}\label{eq:pdf2-ns-source-colligation}
 \boxed{\begin{gathered}
 \mathcal G_{M,J,T}^{\rm sol}\Lambda_{M,T}J_T^{\rm mix}d
 =I_{M,J}J_T^{\rm sol}d,
 \qquad I_{M,J}^*I_{M,J}=I
 \end{gathered}}
\end{equation}
L'histoire se définit à partir du membre gauche,
\begin{equation}\label{eq:pdf2-ns-root-factorization}
 \mathsf H_{M,J,T}^{\rm sol}(d):=
 \mathcal G_{M,J,T}^{\rm sol}\Lambda_{M,T}J_T^{\rm mix}d,
\end{equation}
et non à partir d'une solution choisie ensuite. La lecture terminale
\eqref{eq:pdf2-ns-terminal-reader} et l'action source--first
\eqref{eq:pdf2-ns-source-encoder-action} donnent le carré commutatif
\begin{equation}\label{eq:pdf2-ns-data-history-commutation}
 \boxed{\begin{gathered}
 \mathcal R_{M,\tau}^{(1)}P_{\rm term}
 \mathsf H_{M,J,T}^{\rm sol}(u_0,f)=u_M(\tau)
 \end{gathered}}
\end{equation}
où $P_{\rm term}$ est la projection sur la première ligne de
\eqref{eq:pdf2-ns-complete-column}. Les autres projections sont exactement les
six actions \eqref{eq:pdf2-ns-six-loss-actions}. Cela identifie toute la
colonne à l'histoire de Galerkin sans reconstruire l'état à partir d'une sortie.

\begin{theorem}[Borne racine à partir de la structure de départ explicite]
\label{thm:pdf2-ns-root-uniform}
Pour tout $T<\infty$ et tout $d=(u_0,f)\in\mathcal D_T$,
\begin{equation}\label{eq:pdf2-ns-root-bound}
 \sup_{M,J}\|\mathsf H_{M,J,T}^{\rm sol}(d)\|^2
 \le C_T^{\rm sol}\|d\|_{\mathcal D_T}^2,
\end{equation}
où $C_T^{\rm sol}$ est indépendant de $M,J$.
\end{theorem}

\begin{proof}
Par \eqref{eq:pdf2-ns-root-identification} et
\eqref{eq:pdf2-ns-exact-datum-identification}, le membre gauche est
littéralement
$\|\mathbf G_{M,0;J,T}^{\rm term}\Lambda_{M,T}J_T^{\rm sol}d\|^2$.
La construction propriétaire précédente donne d'abord l'isométrie de cascade
\eqref{eq:ns-import-cascade-isometry}, le télescopage de Stein
\eqref{eq:ns-import-stein-telescope}, le ledger conservé
\eqref{eq:ns-import-ledger} et le financement exact de la retenue
\eqref{eq:ns-import-carry-funded}. Leur somme orthogonale produit la borne racine
\eqref{eq:ns-import-root-bound} et le budget uniforme
\eqref{eq:ns-import-uniform-budget} ; la clôture indépendante
\eqref{eq:nsclose-owner-stein-telescope}--
\eqref{eq:nsclose-proprietary-uniform-budget} transporte ce budget à
toutes les coupures \(M,J\). En substituant l'identification précédente, on obtient
\eqref{eq:pdf2-ns-root-bound}. La formule
\eqref{eq:pdf2-g-sol-cota-dato} est donc le registre ultérieur de la
borne déjà dérivée et non une prémisse de cette démonstration. Les actions
\eqref{eq:pdf2-ns-six-loss-actions}, le terminal
\eqref{eq:pdf2-ns-terminal-action-on-core} et
\eqref{eq:pdf2-ns-global-column-isometry} identifient ses six composantes
sans double comptage.
\end{proof}

\begin{falsifier}[Contrôles d'indépendance et de typage]
Si $J_T^{\rm mix}$ n'est pas littéralement $J_T^{\rm sol}$,
\eqref{eq:pdf2-ns-source-colligation} échoue. Si l'une des six lignes n'est pas une
projection de l'unique $L^{\rm term}$, Parseval échoue dans
\eqref{eq:pdf2-ns-six-loss-no-overcount}. Si l'on tente de définir la norme de
$\mathscr X_{M,j+1,T}^{\rm sol}$ comme la norme précédente moins une perte,
la positivité reste non prouvée et l'opération est interdite. Enfin, si
l'on remplace $b=z-a$ par $z$, le défaut de
\eqref{eq:pdf2-ns-scattering-balance} est
$2\operatorname{Re}\langle z,a\rangle-\|a\|^2$, qui vaut $\|a\|^2$ pour
$z=a\ne0$. Ces quatre contrôles échouent avant d'atteindre la borne uniforme.
\end{falsifier}

\section{Carte KYP finie : calcul correct et statut auxiliaire}

$M,J,T$ étant fixés, soient $A,B,C,D$ les blocs d'un pas affine et $U_+$ la
métrique terminale du niveau suivant. Définissons
\begin{equation}\label{eq:pdf2-ns-kyp-ql}
 Q:=I-B^*U_+B-D^*D,
 \qquad L:=B^*U_+A+D^*C.
\end{equation}
Si $Q\succ0$, la récurrence
\begin{equation}\label{eq:pdf2-ns-kyp-u}
 U:=A^*U_+A+C^*C+L^*Q^{-1}L
\end{equation}
donne la factorisation exacte
\begin{equation}\label{eq:pdf2-ns-kyp-factor}
 \begin{pmatrix}
 U-A^*U_+A-C^*C&-A^*U_+B-C^*D\\
 -B^*U_+A-D^*C&I-B^*U_+B-D^*D
 \end{pmatrix}
 =
 \begin{pmatrix}Q^{-1/2}L&-Q^{1/2}\end{pmatrix}^{*}
 \begin{pmatrix}Q^{-1/2}L&-Q^{1/2}\end{pmatrix}\succeq0.
\end{equation}
Pour $Q\succeq0$, on exige en outre
$\operatorname{Ran}L\subseteq\operatorname{Ran}Q^{1/2}$ et on utilise la
pseudo-inverse. Cette carte est finie. Elle ne donne pas de borne uniforme de $Q^{-1}$, de $U$ ou
du gain lors du retrait de $M,J$.

La mutation qui introduit un terme $C_-^*C_-$ dans un bloc positif et
prétend ensuite le retrouver avec un signe négatif est algébriquement fausse : une
factorisation $R^*R$ ne change pas le signe de l'un de ses carrés. Elle n'est donc pas utilisée
dans la branche directrice.

\section{Clôture uniforme, compacité et solution globale}

La carte KYP précédente reste un \texttt{CONTRÔLE\_NON\_DIRECTEUR}. La clôture
globale provient de la racine source--first, du télescopage terminal et du ledger
complet. Aucune borne uniforme de $Q^{-1}$ n'est utilisée et la limite n'est pas prise à l'intérieur
d'une factorisation finie.

\subsection{Extraction matérielle de la borne de Sobolev}

Dans \eqref{eq:pdf2-ns-complete-column}, la première ligne est littéralement
$\mathcal E^{\rm term}\Gamma_{0:J}$ ; par
\eqref{eq:pdf2-ns-data-history-commutation}, sa coordonnée publiée est
$u_M(\tau)$ et, par \eqref{eq:pdf2-ns-terminal-column-gram}, elle domine
$2^{-1}\|\Lambda^su_M(\tau)\|^2$. Les six lignes sont les applications typées de
\eqref{eq:pdf2-ns-six-loss-actions}. En particulier, les deux premières donnent
\[
 \sum_{j<J}\|L^{\rm diss}_{M,j,T}\Xi\|^2
 =\int_0^T\|b_M\|^2dt,
 \qquad
 \sum_{j<J}\|L^{\rm adv}_{M,j,T}\Xi\|^2
 =\int_0^T|\mathscr B_{s,M}(u_M)|dt.
\]
Comme $z_M=b_M+P_Ma_T(f)$,
\begin{equation}\label{eq:pdf2-ns-wave-to-dissipation}
 \nu\int_0^T \mathscr D_{s,M}dt
 =\int_0^T\|z_M\|^2dt
 \leq2\int_0^T\|b_M\|^2dt+2\|a_T(f)\|_{L^2(0,T;L^2)}^2.
\end{equation}
En appliquant le théorème~\ref{thm:pdf2-ns-root-uniform} à la restriction de la donnée à
$[0,\tau]$ pour chaque $0\le\tau\le T$, la contractivité de la restriction
temporelle de $\mathfrak G_{\rm sol}$ donne
$C_\tau^{\rm sol}\le C_T^{\rm sol}$ et
$\|d|_{[0,\tau]}\|_{\mathcal D_\tau}\le\|d\|_{\mathcal D_T}$. En utilisant en outre
\eqref{eq:pdf2-ns-terminal-column-gram} et
\eqref{eq:pdf2-ns-wave-to-dissipation}, on obtient
\begin{equation}\label{eq:pdf2-ns-cota-requerida}
 \begin{aligned}
 &\sup_M\sup_{0\le t\le T}\|u_M(t)\|_{H^s}^2
 +\nu\sup_M\int_0^T\|u_M(t)\|_{H^{s+1}}^2\,dt\\
 &\quad
 +\sup_M\int_0^T|\mathscr B_{s,M}(u_M(t))|\,dt
 +\sup_{M,J}\sum_{j<J}\|C_{M_j\leftarrow M_{j+1}}\|^2\\
 &\quad
 +\sup_{M,J}\sum_{j<J}
   \bigl(\|Q_{\rm micro}y^{\rm micro}_{M,j}\|^2
          +\|\mathcal S_{M,j}^{\rm sh}\|^2
          +\|Q^{\rm mem}_{M,j+1}m^{\rm raw}_{M,j+1}\|^2\bigr)
 \le C_T(u_0,f),
 \end{aligned}
\end{equation}
où
\[
 C_T(u_0,f):=c_{s,\nu}\bigl(1+C_T^{\rm sol}\bigr)
 \bigl(\|u_0\|_{H^s}^2+
       \|f\|_{L^2(0,T;H^{s-1})}^2\bigr)
\]
est fini pour chaque $T<\infty$ et indépendant de $M$ et $J$. Le passage d'une
profondeur finie à toute la somme des pertes se justifie par monotonie des sommes non
négatives ; dans les trois dernières lignes, chaque norme est la norme $L^2(I_j)$ du
codomaine correspondant dans \eqref{eq:pdf2-ns-six-loss-spaces}. On n'élimine pas
le terminal. L'identité
$E_9J_9=I$ publie la composante macroscopique après cette estimation, tandis que
$Q_{\rm micro}$ et le shell restent quantifiés dans la même inégalité.

\subsection{Contrôle temporel et étape d'Aubin--Lions}

Puisque $s>5/2$, on dispose des applications continues
\begin{equation}\label{eq:pdf2-ns-product-bound}
 H^s\cdot H^s\longrightarrow H^s,
 \qquad
 B:H^s_\sigma\times H^s_\sigma\longrightarrow H^{s-1}_\sigma,
 \qquad
 \|B(v,w)\|_{H^{s-1}}\le c_s\|v\|_{H^s}\|w\|_{H^s}.
\end{equation}
L'équation \eqref{eq:pdf2-ns-galerkin}, la contractivité de $P_M$ et
\eqref{eq:pdf2-ns-cota-requerida} impliquent
\begin{equation}\label{eq:pdf2-ns-time-bound}
 \sup_M\|\partial_tu_M\|_{L^2(0,T;H^{s-1})}
 \le C_T'(u_0,f,\nu)<\infty.
\end{equation}
En effet, $Au_M$ est borné dans $L^2H^{s-1}$ par le deuxième terme de
\eqref{eq:pdf2-ns-cota-requerida} ; le terme bilinéaire est borné au moyen de
\eqref{eq:pdf2-ns-product-bound} ; et $f_M$ est borné par la donnée source.

Les inclusions
\[
 H^{s+1}(\mathbb T^3)\Subset H^s(\mathbb T^3)
 \hookrightarrow H^{s-1}(\mathbb T^3)
\]
et le lemme d'Aubin--Lions fournissent une sous-suite, encore notée $u_M$,
telle que
\begin{equation}\label{eq:pdf2-ns-compactness}
 \begin{aligned}
 u_M&\rightharpoonup^*u&&\text{dans }L^\infty(0,T;H^s),\\
 u_M&\rightharpoonup u&&\text{dans }L^2(0,T;H^{s+1}),\\
 u_M&\longrightarrow u&&\text{dans }L^2(0,T;H^s).
 \end{aligned}
\end{equation}
La convergence forte et \eqref{eq:pdf2-ns-product-bound} donnent
\begin{equation}\label{eq:pdf2-ns-nonlinear-limit}
 B_M(u_M,u_M)\longrightarrow B(u,u)
 \quad\text{dans }L^1(0,T;H^{s-1}).
\end{equation}
En passant à la limite dans la formulation intégrée de
\eqref{eq:pdf2-ns-galerkin}, on obtient
\begin{equation}\label{eq:pdf2-ns-limit-equation}
 \partial_tu+\nu Au+B(u,u)=f,
 \qquad \nabla\!\cdot u=0,
 \qquad u(0)=u_0.
\end{equation}
La donnée initiale est conservée car
$u_M$ est borné dans $H^1(0,T;H^{s-1})$ et converge dans
$C([0,T];H^{s-1})$ après extraction ; comme $P_Mu_0\to u_0$, la trace est celle
indiquée.

\subsection{Pression, unicité et recollement des horizons}

La pression de moyenne nulle n'est pas introduite comme donnée. Elle est reconstruite à partir de la limite par
\begin{equation}\label{eq:pdf2-ns-pressure}
 -\Delta p=\partial_i\partial_j(u_i u_j),
 \qquad \int_{\mathbb T^3}p=0.
\end{equation}
Comme $H^s$ est une algèbre, $p\in L^\infty(0,T;H^s)$ et
$\nabla p\in L^\infty(0,T;H^{s-1})$. Ainsi,
\eqref{eq:pdf2-ns-limit-equation} équivaut à la forme classique
\[
 \partial_tu+(u\cdot\nabla)u-\nu\Delta u+\nabla p=f,
 \qquad\nabla\!\cdot u=0.
\]

Soient $u$ et $v$ deux solutions ayant les mêmes données et posons $w=u-v$. Le produit
dans $H^{s-1}$ et Young donnent
\begin{equation}\label{eq:pdf2-ns-uniqueness}
 \frac12\frac{d}{dt}\|w\|_{H^{s-1}}^2
 +\frac\nu2\|w\|_{H^s}^2
 \le \frac{c_s}{\nu}
 \bigl(\|u\|_{H^s}^2+\|v\|_{H^s}^2\bigr)
 \|w\|_{H^{s-1}}^2.
\end{equation}
Grönwall et $w(0)=0$ impliquent $w=0$. L'unicité élimine la dépendance envers la
sous-suite dans \eqref{eq:pdf2-ns-compactness}.

La construction vaut pour tout $T<\infty$. En l'appliquant dans $T=n$, $n\in\mathbb N$,
les solutions coïncident sur les recouvrements par
\eqref{eq:pdf2-ns-uniqueness} ; elles se recollent donc en une unique solution sur
$[0,\infty)$. Pour justifier ``lisse'' sans extrapoler une seule borne, on répète
la construction pour la suite $s_n=s+n$, $n\ge0$. Les mêmes lecteurs
structurels agissent à chaque échelle, les projections $P_M$ les entrelacent et
\eqref{eq:pdf2-g-sol-cota-dato} donne dans chaque $s_n$ une constante
$C_{T,s_n}^{\rm sol}<\infty$, car les données sont lisses. L'unicité dans
$H^{s-1}$ identifie toutes les limites obtenues. Par
conséquent, $u\in C^\infty((0,\infty)\times\mathbb T^3)$ et la compatibilité des
traces avec $u_0\in C^\infty$ donne la régularité lisse jusqu'à $t=0$. Il n'apparaît pas de temps de
rupture car la borne correspondante est disponible sur chaque horizon
fini et dans chaque $s_n$.

La condition de moyenne nulle n'élimine pas de données périodiques. Pour des données générales
$u_0^{\rm g}$ et $f^{\rm g}$, définissons
\[
 m'(t)=\int_{\mathbb T^3}f^{\rm g}(t,x)\,dx,\qquad
 m(0)=\int_{\mathbb T^3}u_0^{\rm g}(x)\,dx,\qquad
 X(t)=\int_0^tm(r)\,dr.
\]
Le changement de Galilée dépendant du temps
\begin{equation}\label{eq:pdf2-ns-mean-reduction}
 u(t,x)=m(t)+v(t,x-X(t)),\qquad
 g(t,y)=f^{\rm g}(t,y+X(t))-m'(t)
\end{equation}
transforme bijectivement le problème général en problème de moyenne nulle
pour $(v_0,g)$. La translation préserve toutes les normes de Sobolev, la
divergence et la régularité lisse. La construction couvre donc toute donnée périodique
lisse ; en particulier, elle couvre le cas non forcé de l'énoncé classique.

\begin{theorem}[Clôture solénoïdale globale à partir de la structure de départ explicite]
\label{thm:pdf2-ns-global}
Soient $s>5/2$, $\nu>0$,
$u_0\in C^\infty(\mathbb T^3)\cap H^s_\sigma$ et
$f\in C^\infty([0,\infty);H^\infty_\sigma)
\cap L^2_{\rm loc}([0,\infty);H^{s-1}_\sigma)$. Avec la donnée solénoïdale complète
typée par \eqref{eq:pdf2-ns-source-encoder-action},
\eqref{eq:pdf2-ns-sol-transition-action},
\eqref{eq:pdf2-ns-six-loss-no-overcount},
\eqref{eq:pdf2-ns-local-stein},
\eqref{eq:pdf2-ns-terminal-metric} et
\eqref{eq:pdf2-ns-data-history-commutation}, l'image
$\mathfrak G_{\rm sol}$, avec son terminal,
sa retenue, sa mémoire, son shell, son port orthogonal et son lecteur racine source--first, produit une unique
solution globale lisse $(u,p)$ de Navier--Stokes périodique tridimensionnel. Pour chaque
$T<\infty$, elle satisfait \eqref{eq:pdf2-ns-cota-requerida} ; la construction est naturelle en
$M$, stable lors du retrait de $J$ et compatible avec la publication de Galerkin. Au moyen de
\eqref{eq:pdf2-ns-mean-reduction}, le résultat s'étend à toute donnée
périodique lisse, sans imposer une moyenne nulle.
\end{theorem}

\paragraph{Provenance et force probante.}
L'image solénoïdale, le terminal, la retenue, la tour PC--Fock et l'identité locale
\eqref{eq:pdf2-ns-local-stein} sont
l'\texttt{ARCHITECTURE\_AUTORALE\_PRÉEXISTANTE} prise ici comme structure de départ
explicite. Le relèvement mixte \eqref{eq:pdf2-ns-mixed-creation}, la signature chronologique,
les ondes $a/b$, les six domaines de sortie, l'unité homogène et la colonne
terminale sont \texttt{FORMALISATION\_NOUVELLE}. Les équations
\eqref{eq:pdf2-ns-source-encoder-action},
\eqref{eq:pdf2-ns-source-colligation},
\eqref{eq:pdf2-ns-sol-transition-action},
\eqref{eq:pdf2-ns-six-loss-no-overcount},
\eqref{eq:pdf2-ns-local-stein},
\eqref{eq:pdf2-ns-terminal-metric} et
\eqref{eq:pdf2-ns-data-history-commutation} énumèrent exactement les prémisses
employées ; le
théorème~\ref{thm:pdf2-ns-root-uniform} est, en langage humain, \emph{Démontré avec
structure de départ explicite}. Aubin--Lions, la
reconstruction de la pression, Grönwall et le recollement sont la
\texttt{RECONNAISSANCE\_CLASSIQUE} ultérieure. La carte KYP conserve uniquement le statut de
\texttt{CONTRÔLE\_NON\_DIRECTEUR} ; ni son comparateur ni une hypothèse de signe ne font
partie de la chaîne causale de la clôture.

\paragraph{Conclusion affirmative.}
La chaîne directrice est close dans l'ordre
\[
 \begin{gathered}
 \mathfrak G_{\rm sol}
 \longrightarrow J_T^{\rm mix}
 \longrightarrow\Lambda_{M,T}
 \longrightarrow\mathcal G_{M,J,T}^{\rm sol}
 \longrightarrow\mathsf H_{M,J,T}^{\rm sol}\\
 \longrightarrow\mathcal R_{M,\tau}^{(1)}P_{\rm term}=u_M
 \longrightarrow\text{borne uniforme}
 \longrightarrow\text{Aubin--Lions}
 \longrightarrow(u,p)_{[0,\infty)}.
 \end{gathered}
\]
Dans ce profil de départ explicite, il ne reste pas d'obligation KYP supplémentaire :
le terminal et toutes les pertes sont des images orthogonales d'une unique racine
antérieure à l'orbite future.

\endgroup

\HMTFlushCurrentChapter
\chapter{Existence quantique et écart de masse de Yang--Mills par coercivité de jauge et positivité spectrale}
\begingroup
\renewcommand{\chapter}[2][]{}%
\chapter{Image de jauge et composition vers Yang--Mills}

\section{Objet source et séparation obligatoire des générateurs}

L’objet déjà engendré que reçoit ce chapitre est la structure de jauge complète
$\mathfrak G_{\rm gau}$ ; il ne constitue ni une donnée externe ni une hypothèse terminale.
À chaque niveau $m$, il conserve le réseau
$\Lambda_m=a_m\mathbb Z^4$, $a_{m+1}=a_m/9$, la mesure $\mu_m$, les
opérateurs d’entrelacement $J_m$, quatre réflexions $\Theta_{\nu,m}$, l’holonomie, la couleur,
la représentation, la retenue de fusion, la route, la mémoire, le terminal et le lecteur de courbure.
Ses deux dynamiques sont typées par des symboles distincts avant toute limite.

Le semi-groupe d’une seule fibre de mémoire est
\begin{equation}\label{eq:pdf2-ym-fibra-semigrupo}
 K^{\rm fib}_{m,t}
 =P_{\Omega_m}+e^{-3\kappa_{\rm av}t}(I-P_{\Omega_m})
 =e^{-tD_m^{\rm fib}},
 \qquad
 D_m^{\rm fib}:=3\kappa_{\rm av}(I-P_{\Omega_m}).
\end{equation}
Son spectre est contenu dans $\{0,3\kappa_{\rm av}\}$. Le générateur de la carte produit complète est, en revanche,
\begin{equation}\label{eq:pdf2-ym-generador-producto}
 D_m^{\rm prod}:=3\kappa_{\rm av}
 \sum_{c\in\mathcal I_m}(I-E_{m,c}),
 \qquad
 K^{\rm prod}_{m,t}:=e^{-tD_m^{\rm prod}},
\end{equation}
où les espérances orthogonales $E_{m,c}$ commutent. S’il y a deux ou plusieurs
coordonnées non constantes, $D_m^{\rm prod}$ possède des niveaux
$6\kappa_{\rm av},9\kappa_{\rm av},\ldots$ ; c’est pourquoi
\eqref{eq:pdf2-ym-fibra-semigrupo} ne peut être déclaré égal à
\eqref{eq:pdf2-ym-generador-producto}. Le premier est la fibre simple ; le second,
le produit de recouvrement.

\section{Théorème fini de Hoeffding, vide et écart spectral atteint}

Dans la carte produit construite dans \eqref{eq:pdf2-ym-owner-product-factorization}, les $E_{m,c}$ sont des espérances conditionnelles orthogonales et commutantes, et la preuve de \eqref{eq:pdf2-ym-owner-vacuum} établit
\begin{equation}\label{eq:pdf2-ym-interseccion-vacio}
 \bigcap_{c\in\mathcal I_m}\operatorname{Ran}E_{m,c}
 =\mathbb C\Omega_m.
\end{equation}
La décomposition de Hoeffding est la somme orthogonale
\[
 \mathcal H_m=\bigoplus_{S\subseteq\mathcal I_m}\mathcal H_{m,S},
 \qquad u=\sum_Su_S,
\]
où $E_{m,c}u_S=0$ si $c\in S$ et $E_{m,c}u_S=u_S$ si $c\notin S$. Alors
\begin{equation}\label{eq:pdf2-ym-hoeffding}
 D_m^{\rm prod}u_S=3\kappa_{\rm av}|S|u_S,
 \qquad
 \langle u,D_m^{\rm prod}u\rangle
 =3\kappa_{\rm av}\sum_S|S|\|u_S\|^2.
\end{equation}
La composante $S=\varnothing$ est exactement le vide de
\eqref{eq:pdf2-ym-interseccion-vacio}. Par conséquent
\begin{equation}\label{eq:pdf2-ym-cota-finita}
 D_m^{\rm prod}\succeq
 3\kappa_{\rm av}(I-P_{\Omega_m}),
 \qquad
 \|K^{\rm prod}_{m,t}(I-P_{\Omega_m})\|
 \le e^{-3\kappa_{\rm av}t}.
\end{equation}

Dans une coordonnée d’incidence
$q=(q_1,\ldots,q_6)\in\mathbb F_3^6$, le témoin
\begin{equation}\label{eq:pdf2-ym-wc}
 W_c(q):=\mathbf1_{\{q_1+\cdots+q_6=0\}}-\frac13
\end{equation}
satisfait
\begin{equation}\label{eq:pdf2-ym-wc-momentos}
 \mathbb EW_c=0,\qquad \|W_c\|^2=\frac29,\qquad
 D_m^{\rm prod}W_c=3\kappa_{\rm av}W_c.
\end{equation}
Ainsi, à chaque niveau fini, le seuil de \eqref{eq:pdf2-ym-cota-finita} est atteint. Ce calcul fixe le spectre du générateur produit fini ; la mesure quadridimensionnelle continue de Yang--Mills est construite matériellement dans \eqref{eq:pdf2-ym-gibbs-measure}--\eqref{eq:pdf2-ym-g-dynamics}.

\section{Compression de jauge physique}

Soit $G$ compact, connexe et simple. Le groupe qui agit au niveau $m$ est
\[
 \mathcal G_m^{\rm full}
 =G^{V(\Lambda_m)}\times
  \prod_{c\in\mathcal C_m}\mathbb F_3^4,
\]
où le second facteur agit sur les six incidences au moyen d’une matrice
qui doit rester visible. Si les lignes sont indexées par
$(01,02,03,12,13,23)$ et les colonnes par $0,1,2,3$, on fixe
\begin{equation}\label{eq:pdf2-ym-incidence-matrix}
 B^+_{(ij),k}:=\delta_{ik}+\delta_{jk}\in\mathbb F_3,
 \qquad
 (B^+x)_{ij}=x_i+x_j.
\end{equation}
C’est l’incidence non orientée du collier ; elle ne se confond pas avec le
cobord signé d’une arête. L’action finie est
\begin{equation}\label{eq:pdf2-ym-incidence-action}
 q_c\longmapsto q_c+B^+x_c,
 \qquad x_c\in\mathbb F_3^4.
\end{equation}
La projection de Haar
\begin{equation}\label{eq:pdf2-ym-proyeccion-fisica}
 \mathbb P_m^{\rm phys}F
 :=\int_{\mathcal G_m^{\rm full}}F(g^{-1}\!\cdot\,)\,dg,
 \qquad
 \mathcal H_m^{\rm phys}:=\operatorname{Ran}\mathbb P_m^{\rm phys}
\end{equation}
est distincte de la marginale qui oublie l'incidence. L'invariance des espérances $E_{m,c}$ donne
\begin{equation}\label{eq:pdf2-ym-reduccion}
 [\mathbb P_m^{\rm phys},D_m^{\rm prod}]=0,
 \qquad
 [\mathbb P_m^{\rm phys},K_{m,t}^{\rm prod}]=0.
\end{equation}
De plus, chaque variable de $x_c\in\mathbb F_3^4$ apparaît trois fois dans la somme des
six incidences, de sorte que
\begin{equation}\label{eq:pdf2-ym-incidence-conservation}
 \sum_{i<j}(B^+x_c)_{ij}
 =3\sum_{i=0}^3x_{c,i}=0
 \quad\text{dans }\mathbb F_3.
\end{equation}
Par conséquent, $W_c$ est invariant de
jauge et appartient à $\mathcal H_m^{\rm phys}$. L’écart fini n’est pas un
artefact d’une dilatation non physique.

\section{Positivité par réflexion à chaque niveau}

Pour un noyau réversible commun, définissons la mesure à huit faces
\begin{equation}\label{eq:pdf2-ym-ocho-caras}
 \Omega_m^{(8)}(dy_1\cdots dy_8)
 :=\int\mu_m(dz)\prod_{r=1}^{8}K^{\rm prod}_{m,a_m/2}(z,dy_r).
\end{equation}
Si $\Theta_{\nu,m}$ échange les deux moitiés par rapport à la direction $\nu$, la propriété de Markov et la réversibilité factorisent, pour tout observable cylindrique $F$ à support dans le demi-espace positif,
\begin{equation}\label{eq:pdf2-ym-os-finita}
 \int\overline{F(\Theta_{\nu,m}y)}F(y)\,d\Omega_m^{(8)}
 =\|\mathcal B_{m,\nu}F\|_2^2\ge0,
 \qquad \nu=1,\ldots,4.
\end{equation}
La marginale de faces opposées récupère le même
$K^{\rm prod}_{m,a_m}$. Ainsi, les quatre formes OS finies sont
la conséquence d’une seule mesure et d’un seul semi-groupe, non de quatre processus
indépendants.

\section{Propriétaire final : l'état de jauge complet, non sa marginale amputée}

Le propriétaire en vigueur ne prend pas pour espace physique la marginale qui ne conserve que
les liens. Son domaine est $\mathfrak G_{\rm gau}$, avec lien, repère,
représentation, incidence, couleur, route, orientation,
retenue, mémoire, survie et terminal. La projection vers les liens est une application
d’oubli ultérieure. Juger le témoin physique après avoir appliqué cet oubli remplace
l’objet et ne transfère pas la conclusion.

\subsection{Mesure projective et publication des liaisons}

Pour un écran fini $\Gamma_m=(V_m,E_m)$, on part de l'espace relevé
\begin{equation}\label{eq:pdf2-ym-owner-lifted-space}
 \widetilde X_m:=G^{V_m}\times G^{E_m},\qquad
 d\widetilde\mu_m(h,z)
 =\prod_{v\in V_m}dh_v\prod_{e\in E_m}k_{\tau_{m,e}}(z_e)\,dz_e,
\end{equation}
et chaque lien est publié au moyen de
\begin{equation}\label{eq:pdf2-ym-owner-link-publication}
 U_e=h_{s(e)}z_eh_{t(e)}^{-1}.
\end{equation}
Une arête grossière $e$ se raffine en la chaîne ordonnée $e_1\cdots e_9$. Posons \(\boldsymbol\tau=(\tau_1,\ldots,\tau_9)\), avec \(\tau_j=\tau_{m+1,e_j}>0\) et \(\tau=\tau_{m,e}=\sum_{j=1}^9\tau_j\). Soit \(d\lambda_z^{(9)}\) la désintégration de Haar sur la fibre du produit, caractérisée par
\[
 \int_{G^9}F(z_1,\ldots,z_9)\,dz_1\cdots dz_9
 =\int_G\int_{z_1\cdots z_9=z}
   F\,d\lambda_z^{(9)}\,dz.
\]
Le pont correct pour des temps non nécessairement égaux est
\begin{equation}\label{eq:pdf2-ym-owner-heat-bridge}
 dB_{z,\boldsymbol\tau}^{(9)}(z_1,\ldots,z_9)
 =\frac{\prod_{j=1}^9k_{\tau_j}(z_j)}{k_\tau(z)}
   \,d\lambda_z^{(9)},
 \qquad z_1\cdots z_9=z,
\end{equation}
qui est une probabilité car
\(k_{\tau_1}*\cdots*k_{\tau_9}=k_\tau\). Dans le raffinement nonadique uniforme, on
récupère \(\tau_j=\tau/9\), mais la formule ne présuppose pas cette égalité.

Pour rendre l'innovation indépendante sans effacer le produit grossier, on fixe, pour chaque $(e,z)$, un transport borélien
\begin{equation}\label{eq:pdf2-ym-owner-bridge-triangularization}
 \chi_{m,e,z}:\bigl((0,1)^{\mathbb N},du^{\otimes\mathbb N}\bigr)
 \longrightarrow
 \bigl(\{(z_j):z_1\cdots z_9=z\},B_{z,\boldsymbol\tau}^{(9)}\bigr)
\end{equation}
qui est un isomorphisme modulo des ensembles nuls. Il existe car toutes deux sont des réalisations
standard non atomiques ; un choix est fixé une fois pour toutes et transporté sous les
réétiquetages des arêtes. La variable grossière $z$ et l’innovation $u^{\rm det}_{m+1,e}$
sont alors des facteurs indépendants dans la carte triangulaire. Les nouveaux repères sont
intégrés avec Haar. La route, la retenue, l’orientation et la mémoire sont transportées par les
morphismes $\operatorname{Ext}_{\rm gau}$ de la donnée de jauge elle-même.

La projection grossière multiplie les neuf accroissements et oublie exclusivement les
innovations nouvelles. Avec cette définition, on obtient
\begin{equation}\label{eq:pdf2-ym-owner-projectivity}
 (\widehat\pi_{m+1,m})_*\widehat\mu_{m+1}^{\rm ov}
 =\widehat\mu_m^{\rm ov},\qquad
 \widehat J_mF=F\circ\widehat\pi_{m+1,m}.
\end{equation}
La même arête apparaît une seule fois dans la mesure ; les six faces incidentes sont six lecteurs orientés de cette arête. Les chemins qui se recouvrent conservent l'accroissement partagé et ne sont pas remplacés par des copies indépendantes.

\subsection{La fibre d’incidence est une coordonnée physique de l’état complet}

La queue d’incidence incluse dans $\mathfrak G_{\rm gau}$ se décompose sans perdre
d’information :
\begin{equation}\label{eq:pdf2-ym-owner-seed-split}
 \Sigma:(\mathbb F_3^2)^{\mathbb N}\longrightarrow
 (\mathbb F_3^2)^{\mathbb N}\times\mathbb F_3^6,
 \qquad \Sigma_*\sigma=\sigma\otimes u_6.
\end{equation}
Le premier facteur continue d’alimenter les cartes d’innovation et le second conserve
les six incidences $q_c=(q_{01},q_{02},q_{03},q_{12},q_{13},q_{23})$ du collier
$c$. Au lieu de laisser implicite quelles coordonnées le générateur gouverne, on fixe une
carte triangulaire récursive
\begin{equation}\label{eq:pdf2-ym-owner-product-factorization}
 \widehat X_m
 =\widehat X_{m_0}^{\rm seed}
  \times\prod_{\ell=m_0+1}^{m}
       \mathcal U_{\ell/\ell-1}^{\rm det},
 \qquad
 \widehat\mu_m^{\rm ov}
 =\widehat\mu_{m_0}^{\rm seed}
  \otimes\bigotimes_{\ell=m_0+1}^{m}
       \upsilon_{\ell/\ell-1}^{\rm det}.
\end{equation}
Le facteur germe énumère, une fois, les repères et les liaisons grossières, les incidences $q_c$ et leurs cartes de marquage. Chaque $\mathcal U_{\ell/\ell-1}^{\rm det}$ énumère, une fois, les nouveaux repères, les ponts \eqref{eq:pdf2-ym-owner-bridge-triangularization}, les nouvelles incidences et les sous-queues de marquage de ce raffinement. La factorisation est interne à $\mathfrak G_{\rm gau}$ et n'ajoute pas d'autre donnée de départ.

On enregistre l’index exhaustif, disjoint et dénombrable
\begin{equation}\label{eq:pdf2-ym-owner-exhaustive-index}
 \mathcal I_m
 :=\mathcal I_{m_0}^{\rm frame}
 \sqcup\mathcal I_{m_0}^{\rm link}
 \sqcup\mathcal I_{m_0}^{q}
 \sqcup\mathcal I_{m_0}^{\rm mark}
 \sqcup\bigsqcup_{\ell=m_0+1}^{m}
 \bigl(
   \mathcal I_{\ell/\ell-1}^{\rm frame}
   \sqcup\mathcal I_{\ell/\ell-1}^{\rm bridge}
   \sqcup\mathcal I_{\ell/\ell-1}^{q}
   \sqcup\mathcal I_{\ell/\ell-1}^{\rm mark}
 \bigr).
\end{equation}
Chaque facteur indépendant de \eqref{eq:pdf2-ym-owner-product-factorization} apparaît
exactement une fois ; une sous-queue infinie est énumérée par ses cylindres scalaires. Aucun
facteur généalogique ne reste hors de $\mathcal I_m$.
En particulier, les abréviations de la récurrence signifient littéralement
\begin{equation}\label{eq:pdf2-ym-owner-index-abbreviations}
 \begin{aligned}
 \mathcal I_{m_0}^{\rm seed}
 &:={\mathcal I}_{m_0}^{\rm frame}\sqcup
    {\mathcal I}_{m_0}^{\rm link}\sqcup
    {\mathcal I}_{m_0}^{q}\sqcup
    {\mathcal I}_{m_0}^{\rm mark},\\
 \mathcal I_{\ell/\ell-1}^{\rm det}
 &:={\mathcal I}_{\ell/\ell-1}^{\rm frame}\sqcup
    {\mathcal I}_{\ell/\ell-1}^{\rm bridge}\sqcup
    {\mathcal I}_{\ell/\ell-1}^{q}\sqcup
    {\mathcal I}_{\ell/\ell-1}^{\rm mark}.
 \end{aligned}
\end{equation}

La marginale des liens oublie $q_c$ ; la sous-algèbre physique, en revanche, s’obtient
en moyennant l’action de jauge complète
\begin{equation}\label{eq:pdf2-ym-full-gauge-group}
 \mathcal G_m^{\rm full}
 :=G^{V(\Lambda_m)}\times
 \prod_{c\in\mathcal C_m}\mathbb F_3^4,
 \qquad q_c\longmapsto q_c+B^+x_c.
\end{equation}
La projection physique est donc
\begin{equation}\label{eq:pdf2-ym-full-physical-projection}
 \mathbb P_m^{\rm phys}F
 :=\int_{\mathcal G_m^{\rm full}}F(g^{-1}\!\cdot x)\,dg,
 \qquad
 \widehat{\mathscr H}_m^{\rm phys}
 :=\operatorname{Ran}\mathbb P_m^{\rm phys}.
\end{equation}
Elle ne coïncide pas avec l’application d’oubli
$\widehat X_m\to X_m^{\rm link}$. Cette distinction est le garde-fou de typage
qui empêche d’expulser $W_c$ du secteur physique.

\section{Générateur propriétaire et circuit local}

La carte triangulaire complète fournit une application unitaire
\begin{equation}\label{eq:pdf2-ym-owner-triangular-unitary}
 \widehat\Gamma_m:
 L^2(\widehat X_{m+1},\widehat\mu_{m+1}^{\rm ov})
 \longrightarrow
 L^2(\widehat X_m,\widehat\mu_m^{\rm ov})
 \widehat\otimes\mathcal K_{m+1}^{\rm det},
 \qquad
 \widehat\Gamma_m\widehat J_mF=F\otimes\mathbf1_{\rm det}.
\end{equation}
Pour $\iota\in\mathcal I_m$, soit $E_{m,\iota}$ l'espérance qui intègre exactement le facteur $\iota$ de \eqref{eq:pdf2-ym-owner-product-factorization} et laisse tous les autres fixes. Ce sont des projections de Markov orthogonales, elles conservent l'unité et commutent. Sur le cœur dense $\mathscr C_m$ des cylindres qui dépendent d'un nombre fini de coordonnées, on définit la forme
\begin{equation}\label{eq:pdf2-ym-owner-full-form}
 \widehat{\mathcal E}_m[F]
 :=3\kappa_{\rm av}\sum_{\iota\in\mathcal I_m}
   \|(I-E_{m,\iota})F\|_2^2.
\end{equation}
La somme est finie pour chaque cylindre. Sa fermeture est une forme de Dirichlet et
détermine le générateur auto-adjoint markovien
$\widehat H_m^{\rm ov}$. À un niveau germe, cela s’écrit, au sens des
formes,
\begin{equation}\label{eq:pdf2-ym-owner-base-generator}
 \widehat H_{m_0}^{\rm ov}
 =3\kappa_{\rm av}
  \sum_{\iota\in\mathcal I_{m_0}^{\rm seed}}(I-E_{m_0,\iota}),
\end{equation}
et chaque raffinement ajoute exactement les espérances de ses nouvelles coordonnées :
\begin{equation}\label{eq:pdf2-ym-owner-refined-generator}
 \begin{aligned}
 \widehat H_{m+1}^{\rm ov}
 &=\widehat\Gamma_m^*
   (\widehat H_m^{\rm ov}\otimes I
    +I\otimes\widehat H_{m+1}^{\rm det})\widehat\Gamma_m,\\
 \widehat H_{m+1}^{\rm det}
 &=3\kappa_{\rm av}
   \sum_{\iota\in\mathcal I_{m+1/m}^{\rm det}}
   (I-E_{m+1,\iota}^{\rm det}).
 \end{aligned}
\end{equation}
La partition des couleurs n’est pas laissée nominale. Si une cellule est indexée par
$n=(n_1,n_2,n_3,n_4)\in\mathbb Z^4$, on fixe
\begin{equation}\label{eq:pdf2-ym-owner-729-coloring}
 \chi_m(n):=
 \bigl(n_1\bmod9,\ n_2\bmod9,\ n_3+3n_4\bmod9\bigr)
 \in(\mathbb Z/9\mathbb Z)^3,
 \qquad |(\mathbb Z/9\mathbb Z)^3|=9\cdot81.
\end{equation}
Si deux colliers fermés de rayon cellulaire un se coupent, leur différence
$\delta$ vérifie $|\delta_j|\le2$. L’égalité des couleurs impose
$\delta_1=\delta_2=0$ et
$\delta_3+3\delta_4\equiv0\pmod9$. Comme le dernier entier appartient à
$[-8,8]$, il vaut zéro ; les bornes $|\delta_3|,|\delta_4|\le2$ donnent alors
$\delta_3=\delta_4=0$. Par conséquent, des cellules distinctes de même couleur ont
des colliers disjoints. On définit matériellement
\begin{equation}\label{eq:pdf2-ym-owner-729-blocks}
 \mathcal B_{m,r}:=\{\iota\in\mathcal I_m:
   \operatorname{cell}(\iota)=n, \chi_m(n)=r\},
 \qquad r\in(\mathbb Z/9\mathbb Z)^3.
\end{equation}
Le transport d’un écran par $g\in E(4)$ transporte aussi cette
partition ; il n’exige pas qu’une rotation préserve les coordonnées de l’étiquette.
Ainsi, $\{\mathcal B_{m,r}\}_{r=1}^{9\cdot81}$ regroupe des supports spatiaux
disjoints ;
l’indice interne de chaque bloc énumère ses facteurs de
\eqref{eq:pdf2-ym-owner-exhaustive-index}. Pour un cylindre $F$, le produit ne
contient que les blocs dont dépend $F$ ; la fermeture définit le produit fort
pour toute la forme. Le générateur d’un bloc conserve les multiplicités de
Hoeffding ; il n’est pas remplacé par un seul projecteur qui les effacerait. Par
commutativité dans chaque carte,
\begin{equation}\label{eq:pdf2-ym-owner-local-circuit}
\begin{aligned}
 H_{m,B}&:=3\kappa_{\rm av}\sum_{\iota\in B}(I-E_{m,\iota}),\\
 K_{m,B}&:=e^{-a_mH_{m,B}}
 =\prod_{\iota\in B}e^{-3\kappa_{\rm av}a_m(I-E_{m,\iota})},\\
 K_m^{\rm loc}&:=\mathop{\prod_{r,B}}^{\rm strong}K_{m,B}
 =e^{-a_m\widehat H_m^{\rm ov}}.
\end{aligned}
\end{equation}
Le diamètre physique de chaque bloc reste uniformément borné. C'est de là que provient la localité : le générateur n'introduit pas de mise à jour instantanée sur tout le réseau.
La partition est compatible avec Bianchi, et pas seulement avec la disjonction géométrique. Pour une cellule cubique $c$, on fixe le produit orienté, avec tous les facteurs transportés au même sommet,
\begin{equation}\label{eq:pdf2-ym-owner-bianchi-block-product}
 \mathscr B_{m,c}(x):=
 \prod_{p\subset\partial c}^{\longrightarrow}
 T_p\operatorname{Hol}_p(x)T_p^{-1}.
\end{equation}
Chaque arête interne de la frontière apparaît deux fois avec des orientations contraires ;
par conséquent, $\mathscr B_{m,c}=e$ point par point. Un bloc
$\mathcal B_{m,r}$ contient le collier complet de toute cellule qu’il met à jour : s’il touche
une face de $c$, il met simultanément à jour la paire inverse qui partage cette arête.
L’annulation précédente reste littérale après chaque facteur du circuit :
\begin{equation}\label{eq:pdf2-ym-owner-bianchi-block-invariance}
 K_{m,B}\mathbf1_{\{\mathscr B_{m,c}=e\}}
 =\mathbf1_{\{\mathscr B_{m,c}=e\}},
 \qquad
 K_m^{\rm loc}\mathbf1_{\{\mathscr B_{m,c}=e\}}
 =\mathbf1_{\{\mathscr B_{m,c}=e\}}.
\end{equation}
Ainsi, le calendrier de $9\cdot81$ couches préserve Bianchi à chaque étape et dans le produit fort, et pas seulement à la fin d'un balayage.

L'action de jauge permute ou translate les coordonnées intégrées par chaque $E_{m,\iota}$ et préserve leurs mesures. Par conséquent
\begin{equation}\label{eq:pdf2-ym-owner-physical-reduction}
 U_m(g)\widehat H_m^{\rm ov}=\widehat H_m^{\rm ov}U_m(g)
 \quad(g\in\mathcal G_m^{\rm full}),
 \qquad
 [\mathbb P_m^{\rm phys},\widehat H_m^{\rm ov}]=0,
 \qquad
 \mathbb P_{m+1}^{\rm phys}\widehat J_m
 =\widehat J_m\mathbb P_m^{\rm phys}.
\end{equation}
Le semi-groupe complet et son noyau de Markov sont notés
\begin{equation}\label{eq:pdf2-ym-owner-full-kernel}
 \widehat K_{m,t}^{\rm ov}:=e^{-t\widehat H_m^{\rm ov}},
 \qquad
\widehat K_{m,t}^{\rm ov}(x,dy)
 \quad(x,y\in\widehat X_m).
\end{equation}
La première identité de \eqref{eq:pdf2-ym-owner-physical-reduction} implique
\begin{equation}\label{eq:pdf2-ym-owner-full-kernel-equivariance}
 \widehat K_{m,t}^{\rm ov}(gx,gA)
 =\widehat K_{m,t}^{\rm ov}(x,A)
 \quad(g\in\mathcal G_m^{\rm full}).
\end{equation}
Pour typer la réduction sans utiliser un noyau restreint sur le mauvais espace,
soit $q_m:\widehat X_m\to Y_m:=\widehat X_m/\mathcal G_m^{\rm full}$ l’application des
orbites et $\nu_m=(q_m)_*\widehat\mu_m^{\rm ov}$. Le tiré en arrière
\begin{equation}\label{eq:pdf2-ym-owner-quotient-unitary}
 \mathcal W_m:L^2(Y_m,\nu_m)\longrightarrow
 \widehat{\mathscr H}_m^{\rm phys},
 \qquad \mathcal W_m\phi=\phi\circ q_m
\end{equation}
est unitaire. L'équivariance du noyau complet permet de définir
\begin{equation}\label{eq:pdf2-ym-owner-physical-kernel}
 K_{m,t}^{\rm phys}(q_mx,A)
 :=\widehat K_{m,t}^{\rm ov}(x,q_m^{-1}A),
 \qquad A\subseteq Y_m,
\end{equation}
indépendamment du représentant $x$. L’objet physique est alors typé par
\begin{equation}\label{eq:pdf2-ym-owner-D}
 D_m^{\rm YM}
 :=\left.\widehat H_m^{\rm ov}\right|_{\widehat{\mathscr H}_m^{\rm phys}},
 \qquad
 e^{-tD_m^{\rm YM}}
 =\mathcal W_mK_{m,t}^{\rm phys}\mathcal W_m^{-1}.
\end{equation}
C'est le générateur produit complet ; ce n'est pas le générateur d'une seule fibre de \eqref{eq:pdf2-ym-fibra-semigrupo}.

\section{Action de Yang--Mills, mesure dépendant de $g$ et transport exact}

Jusqu’ici, la carte cinématique de référence a été écrite. Nous la noterons
$\widehat\mu_m^0:=\widehat\mu_m^{\rm ov}$ et
$\widehat H_m^0:=\widehat H_m^{\rm ov}$, avec
$\widehat K_{m,t}^0:=e^{-t\widehat H_m^0}=\widehat K_{m,t}^{\rm ov}$.
Pour transformer le lecteur structurel
$\mathcal P_m^{F^2}$ en une loi, et non en une comparaison ultérieure, on reporte
d’abord sa racine positive sur la factorisation triangulaire. La polarisation de
$\mathcal P_m^{F^2}$ dans un collier $b$ est le Gram positif
\[
 Q_{m,b}(u,v):=\frac14\bigl(
 \mathcal P_{m,b}^{F^2}(u+v)-\mathcal P_{m,b}^{F^2}(u-v)
 \bigr).
\]
On quotiente l’espace fini des lecteurs du collier par $\ker Q_{m,b}$, on le complète
avec $Q_{m,b}$ et on note $\mathcal E_{m,b}^{F}$ le Hilbert résultant. Si
$\zeta_{m,b}(x_b)$ est le vecteur des lecteurs centrés de la coordonnée triangulaire
$x_b$, on fixe matériellement
\begin{equation}\label{eq:pdf2-ym-action-root}
 \xi_{m,b}(x_b):=Q_{m,b}^{1/2}\zeta_{m,b}(x_b)
 \in\mathcal E_{m,b}^{F},
 \qquad
 p_{m,b}(x_b):=\|\xi_{m,b}(x_b)\|^2.
\end{equation}
L'indice $b$ parcourt les colliers de Gram à support minimal. Un collier est affecté à la première couleur de \eqref{eq:pdf2-ym-owner-729-coloring} qui contient sa cellule de base ; ainsi, chaque terme apparaît une fois et les colliers d'une même couleur sont disjoints. Plus précisément, la diagonalisation positive regroupe les facteurs triangulaires en une partition disjointe
\begin{equation}\label{eq:pdf2-ym-action-factor-partition}
 \mathcal I_m^{F}
 =\bigsqcup_{r\in(\mathbb Z/9\mathbb Z)^3}
   \bigsqcup_{b\in\mathcal B_{m,r}^{F}}\mathcal I_{m,b},
 \qquad x_b:=(x_\iota)_{\iota\in\mathcal I_{m,b}},
 \qquad
 \mathcal B_{m,r}^{F}:=\{b:Q_{m,b}\ne0,\ \chi_m(b)=r\}.
\end{equation}
Aucune coordonnée triangulaire n’appartient à deux $\mathcal I_{m,b}$ ; si un lecteur
géométrique touche deux colliers, son mode de Gram est assigné au premier et son complément
orthogonal au second. C’est la décomposition spectrale finie de la forme positive,
non une duplication de la variable physique.
La racine de la somme orthogonale, calculée couleur par couleur, donne l’identité, non une
approximation,
\begin{equation}\label{eq:pdf2-ym-action-reader-sum}
 \mathcal P_m^{F^2}(x)
 =\sum_{r\in(\mathbb Z/9\mathbb Z)^3}
   \sum_{b\in\mathcal B_{m,r}^{F}}p_{m,b}(x_b),
 \qquad
 S_{m,g}(x):=\frac1{4g^2}\mathcal P_m^{F^2}(x),
 \quad g\in(0,\infty).
\end{equation}
Ici, $x_b:\widehat X_m\to X_{m,b}$ est la projection des coordonnées du
collier, de sorte que le domaine, l’action et le codomaine de chaque terme sont fixés.
Le lecteur est invariant de jauge parce que $Q_{m,b}$ utilise le produit invariant de
$\mathfrak g$ ; il est invariant par réflexion parce qu’une réflexion ne fait que permuter les
six paires orientées ; et il est naturel par réétiquetage euclidien.
La partition \eqref{eq:pdf2-ym-action-factor-partition} et la carte produit donnent
explicitement
\begin{equation}\label{eq:pdf2-ym-action-factor-measures}
 \widehat\mu_m^0=\lambda_m^{\rm inert}\otimes
   \bigotimes_{r,b}\lambda_{m,b}^0,
 \qquad
 d\lambda_{m,b}^g
 :=(Z_{m,b}^g)^{-1}e^{-p_{m,b}/(4g^2)}d\lambda_{m,b}^0,
 \qquad
 \widehat\mu_{m,g}^{\rm YM}=\lambda_m^{\rm inert}\otimes
   \bigotimes_{r,b}\lambda_{m,b}^g.
\end{equation}
Le facteur inerte contient exactement les coordonnées sur lesquelles le lecteur est nul ; il n'est pas éliminé de la théorie.

La séparation grossier--détail de \eqref{eq:pdf2-ym-owner-product-factorization} est orthogonale pour la même racine : la polarisation directe donne $Q_{m+1}=Q_m\oplus Q_{m+1/m}^{\rm det}$ et $\zeta_{m+1}=\zeta_m\oplus\zeta_{m+1/m}^{\rm det}$. Si $x_{m+1}=\widehat\Gamma_m^{-1}(x_m,u)$, alors
\begin{equation}\label{eq:pdf2-ym-action-cocycle}
 \xi_{m+1}(x_{m+1})
 =\xi_m(x_m)\oplus\xi_{m+1/m}^{\rm det}(u),
 \qquad
 S_{m+1,g}\circ\widehat\Gamma_m^{-1}
 =S_{m,g}+S_{m+1/m,g}^{\rm det}.
\end{equation}
C'est la forme précise sous laquelle retenue et mémoire empêchent de compter deux fois une face ancestrale : la partie grossière n'est pas à nouveau additionnée dans le détail. En particulier, $Z_{m+1,g}=Z_{m,g}Z_{m+1/m,g}^{\rm det}$ pour les normalisateurs et l'on obtient la famille de probabilités véritablement dépendante du couplage
\begin{equation}\label{eq:pdf2-ym-gibbs-measure}
 d\widehat\mu_{m,g}^{\rm YM}(x)
 :=Z_{m,g}^{-1}e^{-S_{m,g}(x)}d\widehat\mu_m^0(x),
 \qquad
 (\widehat\pi_{m+1,m})_*\widehat\mu_{m+1,g}^{\rm YM}
 =\widehat\mu_{m,g}^{\rm YM}.
\end{equation}
Pour écrire correctement Schwinger--Dyson par rapport à Haar, et ne pas feindre que la
mesure de référence est plate, soit $\lambda_m$ le produit de Haar, de Lebesgue et du
dénombrement uniforme de la carte relevée, et soit
$R_m:=d\widehat\mu_m^0/d\lambda_m>0$. L’action totale régulée est
\begin{equation}\label{eq:pdf2-ym-total-regulated-action}
 \mathscr S_{m,g}:=S_{m,g}-\log R_m,
 \qquad
 d\widehat\mu_{m,g}^{\rm YM}
 =\mathscr Z_{m,g}^{-1}e^{-\mathscr S_{m,g}}d\lambda_m.
\end{equation}
Le terme $-\log R_m$ contient les noyaux de chaleur et les cartes de pont de la
structure de jauge de départ ; il est indépendant de $g$. Le terme de Yang--Mills qui
change le couplage et dont la limite classique est calculée ci-dessous est exactement
$S_{m,g}=\mathcal P_m^{F^2}/(4g^2)$.
La cohérence de la seconde égalité et le caractère standard des espaces finis
permettent d’appliquer Kolmogorov et définissent l’unique mesure cylindrique
$\widehat\mu_{\infty,g}^{\rm YM}$ dont les marginales sont
$\widehat\mu_{m,g}^{\rm YM}$.
Aucune valeur propre n’a été utilisée ici : $g$, l’action et
la mesure ont d’abord été fixés.

Pour définir la dynamique sur cette mesure sans perdre le produit, la jauge ni les
identités de surface, on construit le transport, on ne le postule pas. Dans chaque
facteur non atomique $X_{m,b}$, on fixe une carte borélienne
$\rho_{m,b}:X_{m,b}/\mathcal G_{m,b}\to(0,1)$ et on désintègre d’abord sur le
quotient de jauge, puis sur l’orbite de Haar. Soient $F_{m,b}^0$ et
$F_{m,b}^g$ les fonctions de répartition continues de $\rho_{m,b}$ pour la
mesure de référence et la mesure repondérée. Le transport triangulaire est
\begin{equation}\label{eq:pdf2-ym-gibbs-transport-factor}
 t_{m,b,g}:=(F_{m,b}^g)^{-1}\circ F_{m,b}^0,
 \qquad
 T_{m,b,g}(y,h):=(t_{m,b,g}(y),h),
 \qquad
 (T_{m,b,g})_*\lambda_{m,b}^0=\lambda_{m,b}^g.
\end{equation}
Dans un facteur purement atomique sur lequel $p_{m,b}=0$, on prend $T_{m,b,g}=I$. En présence d'atomes et d'une coordonnée continue, la carte est l'ordre lexicographique atome--intervalle et la même formule s'applique à l'espace standard non atomique total. Les transports des $9\cdot81$ couleurs se composent dans l'ordre fixé par \eqref{eq:pdf2-ym-owner-729-coloring} ; à l'intérieur d'une couleur, ils commutent car les colliers sont disjoints. On obtient
\begin{equation}\label{eq:pdf2-ym-gibbs-transport-global}
 T_{m,g}:=\prod_{r\in(\mathbb Z/9\mathbb Z)^3}
           \prod_{b\in\mathcal B_{m,r}^{F}}T_{m,b,g},
 \quad
 (T_{m,g})_*\widehat\mu_m^0=\widehat\mu_{m,g}^{\rm YM},
 \quad
 \widehat\pi_{m+1,m}T_{m+1,g}=T_{m,g}\widehat\pi_{m+1,m}.
\end{equation}
La dernière égalité se vérifie facteur par facteur en utilisant
\eqref{eq:pdf2-ym-action-cocycle} : le transport du facteur grossier est celui déjà fixé et
les nouveaux transports agissent uniquement sur l’innovation. Les cartes de quotient sont
choisies sur un écran représentant et transportées ; c’est pourquoi $T_{m,g}$ commute avec
la jauge et la réflexion et, pour $a\in E(4)$, vérifie
$T_{a\alpha,g}=U_\alpha(a)T_{\alpha,g}U_\alpha(a)^{-1}$.

La composition
\begin{equation}\label{eq:pdf2-ym-full-probability-isomorphism}
 \mathcal T_{m,g}:L^\infty(\widehat X_m,\widehat\mu_{m,g}^{\rm YM})
 \longrightarrow L^\infty(\widehat X_m,\widehat\mu_m^0),
 \qquad \mathcal T_{m,g}F:=F\circ T_{m,g}
\end{equation}
est, par construction, un isomorphisme d'espaces de probabilité et de $*$-algèbres, et se prolonge unitairement à $L^2$. En particulier,
\begin{equation}\label{eq:pdf2-ym-full-algebra-identities}
 \mathcal T_{m,g}(FG)=\mathcal T_{m,g}F\,\mathcal T_{m,g}G,
 \quad
 \mathcal T_{m,g}(F^*)=(\mathcal T_{m,g}F)^*,
 \quad
 \int F\,d\widehat\mu_{m,g}^{\rm YM}
 =\int\mathcal T_{m,g}F\,d\widehat\mu_m^0.
\end{equation}
Le domaine de chaque $T_{m,b,g}$ est l’espace même des holonomies qui satisfait
la loi de subdivision et la contrainte de Bianchi. Son image demeure donc
dans cet espace. En appliquant la multiplicativité précédente générateur par générateur, on
obtient les identités ponctuelles
\begin{equation}\label{eq:pdf2-ym-transport-surface-bianchi}
 \begin{aligned}
 \mathcal T_{m,g}(\operatorname{Hol}_p)
 &=\prod_{p'\subset p}^{\longrightarrow}
   \mathcal T_{m,g}(T_{p'}\operatorname{Hol}_{p'}T_{p'}^{-1}),\\
 \mathcal T_{m,g}(F_{{\rm cl},m})
 &=F_{{\rm cl},m}\circ T_{m,g},\\
 \prod_{p\subset\partial c}^{\longrightarrow}
 \mathcal T_{m,g}(T_p\operatorname{Hol}_pT_p^{-1})&=e,
 \qquad
 \operatorname{Alt}_{\lambda\mu\nu}
 D_\lambda(F_{{\rm cl},m}\circ T_{m,g})_{\mu\nu}=0.
 \end{aligned}
\end{equation}
La troisième ligne est l'annulation arête par arête de la frontière orientée de la frontière de $c$ ; la quatrième est sa polarisation infinitésimale. Ainsi, le transport préserve les produits, l'état, le clover et Bianchi, et pas seulement le premier chaos.

Enfin, on transporte les espérances, le générateur, le noyau, la projection et l’inclusion :
\begin{equation}\label{eq:pdf2-ym-g-dynamics}
 \begin{aligned}
 E_{m,\iota}^{g}&:=\mathcal T_{m,g}^{-1}E_{m,\iota}\mathcal T_{m,g},&
 \widehat H_{m,g}^{\rm YM}&:=\mathcal T_{m,g}^{-1}\widehat H_m^0\mathcal T_{m,g},\\
 \widehat K_{m,t}^{g}&:=\mathcal T_{m,g}^{-1}\widehat K_{m,t}^{0}\mathcal T_{m,g},&
 \mathbb P_{m,g}^{\rm phys}&:=\mathcal T_{m,g}^{-1}
     \mathbb P_m^{\rm phys}\mathcal T_{m,g},\\
 \widehat J_{m,g}&:=\mathcal T_{m+1,g}^{-1}
     \widehat J_m\mathcal T_{m,g},&
 D_{m,g}^{\rm YM}&:=\left.\widehat H_{m,g}^{\rm YM}
     \right|_{\operatorname{Ran}\mathbb P_{m,g}^{\rm phys}}.
 \end{aligned}
\end{equation}
La première ligne n’est pas seulement une conjugaison abstraite. En utilisant
\eqref{eq:pdf2-ym-action-factor-measures}, son action sur un cylindre est
l’actualisation par bain thermique de l’action :
\begin{equation}\label{eq:pdf2-ym-gibbs-conditional-expectation}
 (E_{m,b}^{g}F)(x_{\ne b})
 =\frac{\displaystyle
   \int_{X_{m,b}}F(x_{\ne b},u)e^{-p_{m,b}(u)/(4g^2)}
       d\lambda_{m,b}^0(u)}{\displaystyle
   \int_{X_{m,b}}e^{-p_{m,b}(u)/(4g^2)}d\lambda_{m,b}^0(u)},
 \qquad
 \widehat H_{m,g}^{\rm YM}
 =3\kappa_{\rm av}\sum_b(I-E_{m,b}^{g}).
\end{equation}
Ainsi, tant la mesure que chaque transition locale contiennent $g$ et l’action
$p_{m,b}/(4g^2)$ de manière explicite.
Elles sont markoviennes parce que $\mathcal T_{m,g}\mathbf1=\mathbf1$, locales parce que chaque
$T_{m,b,g}$ a le même collier que $p_{m,b}$, et dépendent de $g$ par
\eqref{eq:pdf2-ym-gibbs-transport-factor}. L’équivalence est spectrale et ne
sélectionne pas la mesure à partir du spectre :
Pour fermer également les équations dynamiques, soit $X_{m,b}^a$ le champ
invariant à gauche dans le facteur de lien $b$ associé à
$e_a\in\mathfrak g$ ; dans une carte continue d’innovation, on utilise la dérivée
directionnelle ordinaire, et les coordonnées finies sont traitées par différence exacte.
Pour tout cylindre lisse $F$, le domaine commun est
$\mathscr C_m^1:=\bigcap_{b,a}\operatorname{Dom}X_{m,b}^a$, en prenant un support
compact dans les cartes ouvertes d’innovation ; dans les facteurs de groupe, il n’y a pas de
frontière. L’invariance de Haar et la règle du produit calculent
\begin{equation}\label{eq:pdf2-ym-schwinger-dyson-finite}
 \int X_{m,b}^aF\,d\widehat\mu_{m,g}^{\rm YM}
 =\int F\,X_{m,b}^a\mathscr S_{m,g}\,d\widehat\mu_{m,g}^{\rm YM},
 \qquad F\in\mathscr C_m^1.
\end{equation}
Dans une coordonnée finie, pour $\Delta_\varepsilon F(q):=F(q+\varepsilon)-F(q)$, le changement de variable $q\mapsto q+\varepsilon$ donne la version exacte, sans dérivée formelle,
\begin{equation}\label{eq:pdf2-ym-schwinger-dyson-discrete}
 \int\Delta_\varepsilon F(q)\,d\widehat\mu_{m,g}^{\rm YM}
 =\int F(q)
  \left(e^{\mathscr S_{m,g}(q)-\mathscr S_{m,g}(q-\varepsilon)}-1\right)
  d\widehat\mu_{m,g}^{\rm YM}(q).
\end{equation}
Si $\delta_\varepsilon$ est la somme orientée de ces champs sur les arêtes qui
sont incidentes à un sommet, l’invariance de jauge de la mesure de référence et de
$S_{m,g}$ donne $\delta_\varepsilon\mathscr S_{m,g}=0$ et, par conséquent, l’identité de Ward
\begin{equation}\label{eq:pdf2-ym-ward-finite}
 \int\delta_\varepsilon F\,d\widehat\mu_{m,g}^{\rm YM}=0.
\end{equation}
Pour le facteur de jauge fini $\mathbb F_3^4$, l'invariance $\mathscr S_{m,g}(q+\varepsilon)=\mathscr S_{m,g}(q)$ réduit \eqref{eq:pdf2-ym-schwinger-dyson-discrete} à la même identité de Ward avec $\delta_\varepsilon$ remplacée par $\Delta_\varepsilon$. Les deux expressions sont compatibles avec $\widehat J_{m,g}$ car \eqref{eq:pdf2-ym-action-cocycle} sépare les dérivées grossières de celles du détail. Tout cylindre de la limite réside à un certain niveau ; ainsi, \eqref{eq:pdf2-ym-schwinger-dyson-finite} et \eqref{eq:pdf2-ym-ward-finite} passent littéralement à l'état cofinal. Ni l'équation de Schwinger--Dyson ni celle de Ward ne contiennent l'écart spectral comme prémisse.

\begin{equation}\label{eq:pdf2-ym-g-gap-before-limit}
 D_{m,g}^{\rm YM}\succeq3\kappa_{\rm av}
 (I-P_{\Omega_{m,g}}),
 \qquad
 W_{c,g}:=\mathcal T_{m,g}^{-1}W_c,
 \qquad
 D_{m,g}^{\rm YM}W_{c,g}=3\kappa_{\rm av}W_{c,g}.
\end{equation}
Dans la suite, on fixe un $g\in(0,\infty)$ et on supprime cet indice ; les
mesures, espérances, inclusions et générateurs sans exposant sont les objets
de \eqref{eq:pdf2-ym-g-dynamics}. Les holonomies géométriques demeurent des
fonctions coordonnées canoniques sur l’espace des configurations et, de ce fait, leur
loi change bien avec $g$ ; $\mathcal T_{m,g}$ les transporte en fonctions composées de la
carte de référence. Seul le témoin spectral est transporté explicitement dans
\eqref{eq:pdf2-ym-g-gap-before-limit}. Cette convention n’efface pas $g$ : sa résidence explicite est
\eqref{eq:pdf2-ym-gibbs-measure}--\eqref{eq:pdf2-ym-g-dynamics}.

\section{Vide simple et témoin qui atteint l'écart spectral}

L'intersection de toutes les espérances de l'état complet est la constante :
\begin{equation}\label{eq:pdf2-ym-owner-vacuum}
 \bigcap_{\iota\in\mathcal I_m}\operatorname{Ran}E_{m,\iota}
 =\mathbb C\Omega_m.
\end{equation}
Cette identité n'est pas introduite comme hypothèse. Dans la carte produit \eqref{eq:pdf2-ym-owner-product-factorization}, chaque $E_{m,\iota}$ intègre une coordonnée de l'indice exhaustif \eqref{eq:pdf2-ym-owner-exhaustive-index} et laisse les autres fixes. Si $u$ appartient à l'intersection, ses espérances conditionnelles par rapport à tout cylindre fini sont constantes. La convergence des martingales de ces cylindres vers la $\sigma$-algèbre complète implique que $u$ est constante presque partout. L'inclusion inverse est immédiate car toute espérance conserve l'unité. Cela prouve \eqref{eq:pdf2-ym-owner-vacuum} et, après la projection de jauge, laisse un vide physique unidimensionnel.
La décomposition de Hoeffding de \eqref{eq:pdf2-ym-owner-base-generator}--\eqref{eq:pdf2-ym-owner-refined-generator} donne
\begin{equation}\label{eq:pdf2-ym-owner-hoeffding}
 \widehat{\mathcal E}_m(u)
 =3\kappa_{\rm av}\sum_{S\Subset\mathcal I_m}|S|\,\|u_S\|^2,
 \qquad
 D_m^{\rm YM}\succeq
 3\kappa_{\rm av}(I-P_{\Omega_m}).
\end{equation}
La fonction
\begin{equation}\label{eq:pdf2-ym-owner-wc}
 W_c(q):=\mathbf1_{\{q_1+\cdots+q_6=0\}}-\frac13
\end{equation}
est invariant sous $G^{V(\Lambda_m)}$ car il ne dépend ni des liens ni des repères. Pour
l’action d’incidence, chaque $x_i$ apparaît trois fois, donc
\begin{equation}\label{eq:pdf2-ym-owner-wc-invariance}
 \sum_{i<j}(B^+x)_{ij}=3\sum_i x_i=0
 \quad\text{dans }\mathbb F_3.
\end{equation}
Ainsi $\mathbb P_m^{\rm phys}W_c=W_c$. Sous $u_6$, la somme des six trits est
uniforme ; par calcul direct,
\begin{equation}\label{eq:pdf2-ym-owner-wc-moments}
 \mathbb EW_c=0,\qquad
 \|W_c\|^2=\frac29,\qquad
 \mathbb E(W_c^3)=\frac2{27}.
\end{equation}
Seul le projecteur $I-E_{m,q_c}$ agit non trivialement sur $W_c$. Par conséquent
\begin{equation}\label{eq:pdf2-ym-owner-wc-eigenvalue}
 D_m^{\rm YM}W_c=3\kappa_{\rm av}W_c.
\end{equation}
La borne inférieure de \eqref{eq:pdf2-ym-owner-hoeffding} n’est pas seulement une borne : le
complément physique contient un vecteur non nul qui atteint le seuil.

\section{Raffinement nonadique et persistance dans la limite}

Les inclusions physiques satisfont
\begin{equation}\label{eq:pdf2-ym-owner-intertwining}
 D_{m+1}^{\rm YM}\widehat J_m
 =\widehat J_mD_m^{\rm YM},\qquad
 \bigl(e^{-a_{m+1}D_{m+1}^{\rm YM}}\bigr)^9\widehat J_m
 =\widehat J_me^{-a_mD_m^{\rm YM}},\qquad a_{m+1}=a_m/9.
\end{equation}
Dans la limite inductive
\begin{equation}\label{eq:pdf2-ym-owner-inductive-limit}
 \widehat{\mathscr H}_\infty^{\rm phys}
 =\varinjlim(\widehat{\mathscr H}_m^{\rm phys},\widehat J_m)
\end{equation}
est défini sur l’union algébrique
\begin{equation}\label{eq:pdf2-ym-owner-limit-form}
 \mathcal E_\infty[\widehat J_{m,\infty}u]
 :=\langle u,D_m^{\rm YM}u\rangle.
\end{equation}
Son générateur est noté
$D_{\infty,g}^{\rm YM}$ ; conformément à la convention postérieure à
\eqref{eq:pdf2-ym-g-gap-before-limit}, on écrit aussi
$D_\infty^{\rm YM}:=D_{\infty,g}^{\rm YM}$.
Le premier entrelacement de \eqref{eq:pdf2-ym-owner-intertwining} rend la
définition indépendante du représentant. Les troncatures de Hoeffding sont
un cœur de forme et donnent une suite de récupération ; Fatou sur la somme non négative
donne la semi-continuité inférieure. La fermeture possède donc
\begin{equation}\label{eq:pdf2-ym-owner-limit-gap}
 \overline{\mathcal E}_\infty[u]
 \ge3\kappa_{\rm av}\|(I-P_{\Omega_\infty})u\|^2.
\end{equation}
L’inclusion conserve les coordonnées ancestrales, de sorte que
$\widehat J_{m,\infty}W_c$ reste non nul et vérifie
\begin{equation}\label{eq:pdf2-ym-owner-limit-witness}
 D_\infty^{\rm YM}\widehat J_{m,\infty}W_c
 =3\kappa_{\rm av}\widehat J_{m,\infty}W_c.
\end{equation}
En conséquence,
\begin{equation}\label{eq:pdf2-ym-owner-exact-gap}
 \inf\bigl(\operatorname{Spec}D_\infty^{\rm YM}\setminus\{0\}\bigr)
 =3\kappa_{\rm av}>0,
 \qquad \ker D_\infty^{\rm YM}=\mathbb C\Omega_\infty.
\end{equation}

\section{Quatre reconstructions OS de la même dynamique}

Le noyau complet de \eqref{eq:pdf2-ym-owner-full-kernel} définit une loi réversible sur $\widehat X_m$. Pour la reconstruction physique, on utilise sa descente bien typée \eqref{eq:pdf2-ym-owner-physical-kernel} : sur le quotient $Y_m$, on définit
\begin{equation}\label{eq:pdf2-ym-owner-eight-face-law}
 \Omega_{m,{\rm phys}}^{(8)}(d\bar y_1\cdots d\bar y_8)
 =\int\nu_m(d\bar z)
  \prod_{r=1}^8K_{m,a_m/2}^{\rm phys}(\bar z,d\bar y_r).
\end{equation}
Pour chaque direction $\nu=1,\ldots,4$, la réflexion échange les quatre faces
positives avec les négatives. En conditionnant par rapport à l’état central $z$ et en utilisant
la réversibilité,
\begin{equation}\label{eq:pdf2-ym-owner-four-os}
 \int\overline{F(\Theta_{\nu,m}y)}F(y)\,
 d\Omega_{m,{\rm phys}}^{(8)}(y)
 =\|\mathcal B_{m,\nu}F\|_2^2\ge0.
\end{equation}
La marginale de la paire opposée normale à $\nu$ est
\begin{equation}\label{eq:pdf2-ym-owner-os-marginal}
 \nu_m(d\bar x)K_{m,a_m}^{\rm phys}(\bar x,d\bar y).
\end{equation}
Pour identifier l'hamiltonien d'une direction, on prend le sous-espace OS des observables d'une face positive simple, $F(y)=F_\nu(\bar y_{\nu,+})$ ; les observables de plusieurs faces conservent la positivité précédente, mais ne sont pas confondus avec cet espace de transfert. Dans le quotient de face simple, l'application
\begin{equation}\label{eq:pdf2-ym-owner-os-unitary}
 \mathcal V_{m,\nu}[F]
 :=e^{-a_mD_m^{\rm YM}/2}\mathcal W_mF
\end{equation}
est isométrique pour la norme OS et a une image dense. Son extension identifie le
quotient OS à $\widehat{\mathscr H}_m^{\rm phys}$. Posons
$\overline D_m:=\mathcal W_m^{-1}D_m^{\rm YM}\mathcal W_m$ et
$\overline J_m:=\mathcal W_{m+1}^{-1}\widehat J_m\mathcal W_m$.
L’entrelacement et $a_m=9a_{m+1}$ donnent
\begin{equation}\label{eq:pdf2-ym-owner-os-range-inclusion}
 \overline J_me^{-a_m\overline D_m/2}
 =e^{-a_m\overline D_{m+1}/2}\overline J_m
 =\bigl(e^{-a_{m+1}\overline D_{m+1}/2}\bigr)^9\overline J_m.
\end{equation}
Le dernier membre appartient à l’image de
$e^{-a_{m+1}\overline D_{m+1}/2}$ ; c’est pourquoi l’inverse de $\mathcal V_{m+1,\nu}$
qui apparaît ci-après ne s’applique que sur son image et est bien défini. Les
connecteurs
\begin{equation}\label{eq:pdf2-ym-owner-os-connectors}
 \mathcal J_{m,\nu}^{\rm OS}
 :=\mathcal V_{m+1,\nu}^{-1}\widehat J_m\mathcal V_{m,\nu}
\end{equation}
sont des isométries et entrelacent les quatre hamiltoniens reconstruits :
\begin{equation}\label{eq:pdf2-ym-owner-common-os-hamiltonian}
 H_{{\rm OS},m+1}^{(\nu)}\mathcal J_{m,\nu}^{\rm OS}
 =\mathcal J_{m,\nu}^{\rm OS}H_{{\rm OS},m}^{(\nu)},
 \qquad
 \mathcal V_{\infty,\nu}H_{{\rm OS},\infty}^{(\nu)}
 \mathcal V_{\infty,\nu}^{-1}=D_\infty^{\rm YM}.
\end{equation}
Les quatre réflexions ne fabriquent pas quatre mesures : ce sont quatre publications du même
état, du même noyau et du même générateur.

\section{Localité, action euclidienne et publication de courbure}

La catégorie dirigée des écrans comprend les raffinements, les superpositions et les transports
euclidiens. Si $a\in E(4)$ et $\alpha$ est un écran, le réétiquetage complet définit
\begin{equation}\label{eq:pdf2-ym-owner-e4-finite-transport}
 U_\alpha(a):\widehat{\mathscr H}_\alpha^{\rm phys}
 \longrightarrow\widehat{\mathscr H}_{a\alpha}^{\rm phys},
 \qquad
 U_\beta(a)\widehat J_{\alpha\beta}
 =\widehat J_{a\alpha,a\beta}U_\alpha(a).
\end{equation}
Le transport préserve la mesure, le générateur, les idéaux nuls et les supports. L'identité de naturalité permet de définir, sans choisir de représentant, une application unitaire sur la colimite algébrique :
\begin{equation}\label{eq:pdf2-ym-owner-e4-colimit-action}
 U(a)\widehat J_{\alpha,\infty}F
 :=\widehat J_{a\alpha,\infty}U_\alpha(a)F.
\end{equation}
Les réétiquetages vérifient $U_{a\alpha}(b)U_\alpha(a)=U_\alpha(ba)$ et préservent
la norme ; c’est pourquoi \eqref{eq:pdf2-ym-owner-e4-colimit-action} se prolonge en une
représentation unitaire de $E(4)$.

\subsection{Le canal d'incidence est une courbure locale, non un facteur auxiliaire}

Tout groupe compact, connexe et simple possède un tore maximal non trivial. On fixe $t_3\in G$ d'ordre exactement trois. Pour un collier $c$, soit
\begin{equation}\label{eq:pdf2-ym-owner-incidence-flux}
 \omega_c(q):=\sum_{i<j}q_{ij}\in\mathbb F_3.
\end{equation}
L’identité \eqref{eq:pdf2-ym-incidence-conservation} prouve que $\omega_c$ est
invariant sous l’action finie. Si $h_c$ est le repère au sommet de base du collier,
on définit l’holonomie géométrique et, pour le témoin spectral, son image transportée
dans la loi au couplage $g$ :
\begin{equation}\label{eq:pdf2-ym-owner-incidence-holonomy}
 \operatorname{Hol}_{m,c}^{\rm inc,geom}:=h_ct_3^{\,\omega_c(q)}h_c^{-1}\in G,
 \qquad
 \operatorname{Hol}_{m,c,g}^{\rm inc,spec}
 :=\mathcal T_{m,g}^{-1}\operatorname{Hol}_{m,c}^{\rm inc,geom},
 \qquad
 \operatorname{Hol}_{m,c}^{\rm inc}:=\operatorname{Hol}_{m,c,g}^{\rm inc,spec}.
\end{equation}
Sous $G^{V_m}$, elle se transforme par conjugaison et sous
$\prod_c\mathbb F_3^4$, elle demeure fixe. Comme $t_3$ est d’ordre trois, le calcul
fonctionnel de cette variable à trois valeurs donne l’identité littérale
\begin{equation}\label{eq:pdf2-ym-owner-wc-curvature-functional}
 W_{c,g}:=\mathcal T_{m,g}^{-1}W_c^0
 =\mathbf1_{\{e\}}\!\left(\operatorname{Hol}_{m,c}^{\rm inc}\right)-\frac13.
\end{equation}
On écrit $W_c:=W_{c,g}$ dans le reste du chapitre. Comme $T_{m,b,g}$ agit à l'intérieur du quotient du même collier, l'holonomie spectrale est une fonction de son algèbre locale d'holonomies géométriques ; elle n'ajoute pas de variable et ne change pas le support. Les holonomies géométriques, utilisées dans $S_{m,g}$ et dans le clover, maintiennent leur distribution dépendante de $g$.
Le projecteur du membre de droite est une fonction de classe locale ; par conséquent, $W_c$ n'est pas une coordonnée auxiliaire découplée, mais un observable de l'holonomie d'incidence de la même restriction de jauge.

Pour un ouvert borné $O\subset\mathbb R^4$, soit
$\mathfrak A_{\rm YM}(O)$ l’algèbre engendrée par les fonctions de classe bornées de
\eqref{eq:pdf2-ym-owner-incidence-holonomy}, les holonomies de plaquette et les lecteurs
clover à support dans $O$. On définit le secteur de courbure physique par
\begin{equation}\label{eq:pdf2-ym-owner-curvature-cyclic-sector}
 \mathscr H_{\rm YM}^{\rm curv}
 :=\overline{\bigcup_{O\Subset\mathbb R^4}
   \mathfrak A_{\rm YM}(O)\Omega_\infty}
 \subseteq\widehat{\mathscr H}_\infty^{\rm phys}.
\end{equation}
Par \eqref{eq:pdf2-ym-owner-wc-curvature-functional}, chaque
$\widehat J_{m,\infty}W_c$ appartient matériellement à ce secteur.
Les espérances de coordonnée envoient les cylindres de courbure sur des cylindres de
courbure ; le semi-groupe laisse donc invariant
$\mathscr H_{\rm YM}^{\rm curv}$. Étant auto-adjoint, le secteur est réducteur. La
borne de \eqref{eq:pdf2-ym-owner-limit-gap} s’y restreint et le vecteur précédent
atteint l’égalité :
\begin{equation}\label{eq:pdf2-ym-owner-curvature-sector-gap}
 D_{\rm curv,g}^{\rm YM}
 :=\left.D_{\infty,g}^{\rm YM}\right|_{\mathscr H_{\rm YM}^{\rm curv}},
 \qquad D_{\rm curv}^{\rm YM}:=D_{\rm curv,g}^{\rm YM},
 \qquad
 \inf\bigl(\operatorname{Spec}D_{\rm curv,g}^{\rm YM}\setminus\{0\}\bigr)
 =3\kappa_{\rm av}.
\end{equation}

\subsection{Système cofinal de publication et produits croisés}

L’inclusion $\widehat J_{\alpha\beta}$ conserve une coordonnée ancestrale ; elle ne doit pas
être confondue avec la moyenne de bloc qui publie un champ continu. Pour écrire
cette seconde application sans altérer la généalogie, soit $\mathfrak P$ l’ensemble dirigé
des écrans cellulaires bornés, fermé sous superposition, raffinement nonadique et
transport euclidien. Pour $\alpha\in\mathfrak P$, posons
\begin{equation}\label{eq:pdf2-ym-owner-cell-step-space}
 V_\alpha:=\operatorname{span}\left\{
 e_{\alpha,c}:=|c|^{-1/2}\mathbf1_c:
 c\in\mathcal C_\alpha^{\rm cell}\right\}
 \subset L^2(\mathbb R^4),
 \qquad P_\alpha:L^2(\mathbb R^4)\to V_\alpha.
\end{equation}
Les superpositions font que $P_\alpha h\to h$ pour tout $h\in L^2$ le long de
$\mathfrak P$. Du côté de la jauge, on normalise le témoin par
\begin{equation}\label{eq:pdf2-ym-owner-normalized-cell-witness}
 Z_{\alpha,c}:=\sqrt{\frac92}\,W_{\alpha,c},
 \qquad
 \langle Z_{\alpha,c},Z_{\alpha,d}\rangle=\delta_{cd},
 \qquad
 D_\alpha^{\rm YM}Z_{\alpha,c}=3\kappa_{\rm av}Z_{\alpha,c}.
\end{equation}
La seconde égalité utilise exactement les moments de \eqref{eq:pdf2-ym-owner-wc-moments} et la factorisation de coordonnées distinctes. Soit $\mathcal K_\alpha^{\rm cell}$ l'espace engendré par ces vecteurs. Si $\beta\succeq\alpha$, l'opérateur d'entrelacement de publication, distinct de $\widehat J_{\alpha\beta}$, est défini générateur par générateur par
\begin{equation}\label{eq:pdf2-ym-owner-block-intertwiner}
 I_{\alpha\beta}^{\rm curv}Z_{\alpha,c}
 :=\sum_{\substack{d\in\mathcal C_\beta^{\rm cell}\\d\subset c}}
       \sqrt{\frac{|d|}{|c|}}\,Z_{\beta,d}.
\end{equation}
Comme les sous-cellules sont disjointes et ont pour somme $|c|$, on obtient, sans passage à la limite,
\begin{equation}\label{eq:pdf2-ym-owner-block-intertwiner-identities}
 (I_{\alpha\beta}^{\rm curv})^*I_{\alpha\beta}^{\rm curv}=I,
 \quad
 I_{\beta\gamma}^{\rm curv}I_{\alpha\beta}^{\rm curv}
 =I_{\alpha\gamma}^{\rm curv},
 \quad
 D_\beta^{\rm YM}I_{\alpha\beta}^{\rm curv}
 =I_{\alpha\beta}^{\rm curv}D_\alpha^{\rm YM}.
\end{equation}
Le réétiquetage cellulaire conserve les volumes et les descendants ; pour $a\in E(4)$, on a en outre
\begin{equation}\label{eq:pdf2-ym-owner-block-e4-naturality}
 U_\beta(a)I_{\alpha\beta}^{\rm curv}
 =I_{a\alpha,a\beta}^{\rm curv}U_\alpha(a).
\end{equation}

Pour $h\in C_c^\infty(\mathbb R^4)$, définissons maintenant le lecteur de bloc
\begin{equation}\label{eq:pdf2-ym-owner-smeared-witness}
 \Phi_\alpha(h):=
 \sum_{c\in\mathcal C_\alpha^{\rm cell}}
 \langle h,e_{\alpha,c}\rangle_{L^2}Z_{\alpha,c}.
\end{equation}
Sur un maillage uniforme, cette formule est
$\sqrt{9/2}\,a_\alpha^2\sum_c(P_\alpha h)(x_c)W_{\alpha,c}$ ; la moyenne cellulaire,
non une évaluation ponctuelle, fixe le produit croisé. En effet, si $\gamma$ raffine
deux écrans $\alpha$ et $\beta$, les définitions précédentes donnent l’identité
calculée
\begin{equation}\label{eq:pdf2-ym-owner-smeared-cross-gram}
 \left\langle
  I_{\alpha\gamma}^{\rm curv}\Phi_\alpha(h),
  I_{\beta\gamma}^{\rm curv}\Phi_\beta(k)
 \right\rangle
 =\langle P_\alpha h,P_\beta k\rangle_{L^2}.
\end{equation}
Par polarisation, avec $k=h$,
\begin{equation}\label{eq:pdf2-ym-owner-smeared-cauchy-derived}
 \left\|
  I_{\alpha\gamma}^{\rm curv}\Phi_\alpha(h)
  -I_{\beta\gamma}^{\rm curv}\Phi_\beta(h)
 \right\|_2^2
 =\|P_\alpha h-P_\beta h\|_{L^2}^2\longrightarrow0.
\end{equation}
Ainsi, la propriété de Cauchy est une conséquence de la convergence forte des
projections, et non une identité supplémentaire postulée. La colimite de
\eqref{eq:pdf2-ym-owner-block-intertwiner} contient donc
\begin{equation}\label{eq:pdf2-ym-owner-smeared-gram}
 \Phi(h):=\lim_\alpha I_{\alpha,\infty}^{\rm curv}\Phi_\alpha(h),
 \qquad
 \langle\Phi(h),\Phi(k)\rangle=\langle h,k\rangle_{L^2},
 \qquad
 D_{\rm pub}^{\rm YM}\Phi(h)=3\kappa_{\rm av}\Phi(h).
\end{equation}

Cette publication n'ajoute pas de degrés de liberté. À chaque étape, la décomposition $V_\beta=V_\alpha\oplus(V_\beta\ominus V_\alpha)$ possède une base de Haar nonadique. L'ordre exhaustif des nouvelles incidences de \eqref{eq:pdf2-ym-owner-exhaustive-index} apparie cette base, étiquette par étiquette, avec des coordonnées $Z_{\beta,d}$ du facteur d'innovation, en conservant le support de la fonction de Haar. Les applications unitaires triangulaires $\widehat\Gamma_\alpha$ recollent ces appariements et produisent une isométrie basée
\begin{equation}\label{eq:pdf2-ym-owner-publication-into-physical}
 \mathcal U_{\rm curv}:
 \varinjlim(\mathcal K_\alpha^{\rm cell},I_{\alpha\beta}^{\rm curv})
 \longrightarrow\mathscr H_{\rm YM}^{\rm curv},
 \quad
 D_{\rm curv}^{\rm YM}\mathcal U_{\rm curv}
 =\mathcal U_{\rm curv}D_{\rm pub}^{\rm YM},
 \quad U(a)\mathcal U_{\rm curv}=\mathcal U_{\rm curv}U_{L^2}(a).
\end{equation}
Désormais, $\Phi(h)$ désigne également son image par $\mathcal U_{\rm curv}$.
La compatibilité de \eqref{eq:pdf2-ym-gibbs-transport-global} définit
$\mathcal T_{\infty,g}$ à la limite. La publication physique au couplage fixé
n’est pas une seconde isométrie choisie uniquement sur le premier chaos, mais la restriction
\begin{equation}\label{eq:pdf2-ym-full-publication-restriction}
 \mathcal U_{{\rm curv},g}
 :=\mathcal T_{\infty,g}^{-1}\mathcal U_{{\rm curv},0},
 \qquad
 \widetilde{\mathcal U}_g
 :=\mathcal T_{\infty,g}^{-1}:
 L^\infty(\widehat X_\infty,\widehat\mu_\infty^0)
 \longrightarrow
 L^\infty(\widehat X_\infty,\widehat\mu_{\infty,g}^{\rm YM}).
\end{equation}
Sur la $*$-algèbre polynomiale engendrée par les holonomies, les plaquettes, clover et
$\Phi(h)$, on vérifie, générateur par générateur puis par linéarité,
\begin{equation}\label{eq:pdf2-ym-full-publication-products}
 \widetilde{\mathcal U}_g(AB)
 =\widetilde{\mathcal U}_g(A)\widetilde{\mathcal U}_g(B),
 \quad
 \widetilde{\mathcal U}_g(A^*)=\widetilde{\mathcal U}_g(A)^*,
 \quad
 \omega_g(\widetilde{\mathcal U}_gA)=\omega_0(A).
\end{equation}
Les égalités de surface, de clover et de Bianchi sont préservées par
\eqref{eq:pdf2-ym-transport-surface-bianchi} ; par densité monotone, l’extension
vaut dans toute l’algèbre locale bornée. Ainsi,
\eqref{eq:pdf2-ym-owner-smeared-cross-gram} est la restriction d’un isomorphisme
probabiliste et algébrique complet, non un substitut de celui-ci.

\subsection{Noyau dense, continuité forte et réseau local}

Soit $M_{\Phi(h)}$ l'opérateur de multiplication affilié à l'algèbre locale et
\begin{equation}\label{eq:pdf2-ym-owner-curvature-weyl}
 \mathsf W(h,t):=\exp\!\bigl(itM_{\Phi(h)}\bigr),
 \qquad h\in C_c^\infty(\mathbb R^4),\quad t\in\mathbb R.
\end{equation}
On complète $\mathfrak A_{\rm YM}(O)$ par ces unitaires pour
$\operatorname{supp}h\subset O$ et par les fonctions de classe bornées des
holonomies de plaquette. Notons $\mathfrak A_{\rm YM}(O)^{\rm sm}$ la
$*$-algèbre ainsi engendrée avec des fonctions tests lisses. Le noyau
\begin{equation}\label{eq:pdf2-ym-owner-smooth-core}
 \mathscr D_{\rm sm}:=\operatorname{span}\left\{
 A_1\cdots A_n\Omega_\infty:
 A_j\in\mathfrak A_{\rm YM}(O_j)^{\rm sm},\quad
 O_j\Subset\mathbb R^4\right\}
\end{equation}
est dense par la définition cyclique de \eqref{eq:pdf2-ym-owner-curvature-cyclic-sector}. Les équations \eqref{eq:pdf2-ym-owner-block-e4-naturality} et \eqref{eq:pdf2-ym-owner-publication-into-physical} donnent, pour $a\in E(4)$,
\begin{equation}\label{eq:pdf2-ym-owner-field-e4-covariance}
 U(a)\Phi(h)=\Phi(a\!\cdot h),
 \qquad
 \|(U(a)-I)\Phi(h)\|=\|a\!\cdot h-h\|_{L^2}.
\end{equation}
En outre, $|e^{ix}-e^{iy}|\le|x-y|$ et
\eqref{eq:pdf2-ym-owner-smeared-gram} impliquent
\begin{equation}\label{eq:pdf2-ym-owner-weyl-continuity}
 \|\bigl(\mathsf W(a\!\cdot h,t)-\mathsf W(h,t)\bigr)\Omega_\infty\|
 \le |t|\,\|a\!\cdot h-h\|_{L^2}.
\end{equation}
La continuité se vérifie également pour les générateurs d'holonomie du cœur ; \eqref{eq:pdf2-ym-owner-weyl-continuity} à elle seule ne les inclut pas. Pour un chemin polygonal orienté $\gamma$, une plaquette $p$ et $f\in C^1(G)^{\rm class}$, on pose
\begin{equation}\label{eq:pdf2-ym-owner-holonomy-core-generators}
 A_{\gamma,f}:=f(\operatorname{Hol}_\gamma),
 \qquad A_{p,f}:=f(\operatorname{Hol}_p).
\end{equation}
Dans une superposition de $\gamma$ et de $a\gamma$, avec $a\in E(4)$, la loi de pont
\eqref{eq:pdf2-ym-owner-heat-bridge} annule les accroissements communs. Pour chaque
accroissement restant $z$, on utilise
$|f(xz)-f(x)|\le\|\nabla f\|_\infty d_G(z,e)$ et la borne du second moment du
noyau de chaleur
\(
 \int_Gd_G(z,e)^2k_\tau(z)\,dz\le C_G\tau
\).
Pour la loi repondérée, $0<e^{-p/(4g^2)}\le1$ et la continuité de $p$ sur le
facteur compact donnent
\begin{equation}\label{eq:pdf2-ym-owner-gibbs-heat-moment}
 \frac{\int_Gd_G(z,e)^2e^{-p(z)/(4g^2)}k_\tau(z)\,dz}
      {\int_Ge^{-p(z)/(4g^2)}k_\tau(z)\,dz}
 \le C_{G,g}\tau,
\end{equation}
uniformément sur les écrans d’un compact et pour $0<\tau\le\tau_0$.
En additionnant les temps des incréments de la différence symétrique, on obtient
l’estimation calculée
\begin{equation}\label{eq:pdf2-ym-owner-holonomy-e4-estimate}
 \begin{aligned}
 \|U(a)A_{\gamma,f}\Omega_\infty-A_{\gamma,f}\Omega_\infty\|_2^2
 &\le C_{G,g}\|\nabla f\|_\infty^2
     \tau(\gamma\triangle a\gamma),\\
 \|U(a)A_{p,f}\Omega_\infty-A_{p,f}\Omega_\infty\|_2^2
 &\le C_{G,g}\|\nabla f\|_\infty^2
     \tau(\partial p\triangle\partial(ap)).
 \end{aligned}
\end{equation}
Ici, la différence symétrique n'est pas mesurée comme deux courbes disjointes : les connecteurs de surface la remplissent par le ruban orienté minimal $\Sigma(\gamma,a\gamma)$ et
\begin{equation}\label{eq:pdf2-ym-owner-bridge-time-ribbon}
 \tau(\gamma\triangle a\gamma)
 :=\sum_{p\subset\Sigma(\gamma,a\gamma)}\tau_p,
 \qquad
 \tau(\gamma\triangle a\gamma)
 \le C_\gamma d_{E(4)}(a,e),
\end{equation}
avec la même définition pour les frontières de plaquette. L’inégalité découle de la
somme des temps additifs de \eqref{eq:pdf2-ym-owner-heat-bridge} sur une bande
de largeur $d_{E(4)}(a,e)$ et de longueur bornée par celle de $\gamma$.
La mesure du temps de pont du membre de droite tend ainsi vers zéro lorsque $a\to e$ ;
c’est exactement la continuité des connecteurs de route de l’image de jauge.
Chaque lecteur clover borné $\chi(F_{\rm cl})$, $\chi\in C_b^1$, est fonction d’une
somme finie de plaquettes transportées ; la même borne, multipliée par
$\|\chi'\|_\infty$ et par le nombre de termes, prouve également sa continuité.

Pour des fonctions de classe seulement bornées, on rend visible la régularisation qui
entre dans le cœur. Si $d b$ est la mesure de Haar à gauche sur $E(4)$ et
$\eta\in C_c^\infty(E(4))$, on définit l’intégrale de Bochner
\begin{equation}\label{eq:pdf2-ym-owner-orbit-smearing}
 A[\eta]:=\int_{E(4)}\eta(b)U(b)AU(b)^*\,db,
 \qquad A\in\{A_{\gamma,f},A_{p,f},\chi(F_{{\rm cl}})\},
 \quad \chi\in C_b^1.
\end{equation}
Avec $(L_a\eta)(b):=\eta(a^{-1}b)$, le changement de variable $b\mapsto ab$ donne
\begin{equation}\label{eq:pdf2-ym-owner-orbit-smearing-continuity}
 U(a)A[\eta]U(a)^*=A[L_a\eta],
 \qquad
 \|(A[L_a\eta]-A[\eta])\Omega_\infty\|
 \le\|A\|\,\|L_a\eta-\eta\|_{L^1(E(4))}\longrightarrow0.
\end{equation}
Il est désormais entendu que les générateurs d'holonomie, de plaquette et de clover de $\mathfrak A_{\rm YM}(O)^{\rm sm}$ sont ceux de \eqref{eq:pdf2-ym-owner-holonomy-core-generators} avec $f\in C^1$, ou leurs régularisations \eqref{eq:pdf2-ym-owner-orbit-smearing}, et que le support balayé reste compact à l'intérieur de $O$. Les approximations de l'identité dans $E(4)$, avec \eqref{eq:pdf2-ym-owner-holonomy-e4-estimate}, récupèrent les générateurs non régularisés dans $L^2$ ; le secteur cyclique n'est donc pas réduit.

Un télescopage de produits qui utilise
\eqref{eq:pdf2-ym-owner-weyl-continuity} et
\eqref{eq:pdf2-ym-owner-orbit-smearing-continuity} prouve la continuité de tous
les générateurs de $\mathscr D_{\rm sm}$ ; par
densité et unitarité,
\begin{equation}\label{eq:pdf2-ym-owner-strong-e4}
 \lim_{a\to e}\|(U(a)-I)\Psi\|=0
 \qquad(\Psi\in\mathscr H_{\rm YM}^{\rm curv}).
\end{equation}
Les supports disjoints utilisent des facteurs disjoints et commutent. On obtient ainsi le réseau
\begin{equation}\label{eq:pdf2-ym-owner-local-net}
\begin{aligned}
 O_1\subseteq O_2&\Rightarrow
 \mathfrak A_{\rm YM}(O_1)\subseteq\mathfrak A_{\rm YM}(O_2),\\
 [\mathfrak A_{\rm YM}(O_1),\mathfrak A_{\rm YM}(O_2)]&=0
 \quad (O_1\cap O_2=\varnothing),\\
 U(a)\mathfrak A_{\rm YM}(O)U(a)^*&=\mathfrak A_{\rm YM}(aO).
\end{aligned}
\end{equation}
L’identité de Gram pour $h\ne0$ et
\eqref{eq:pdf2-ym-owner-wc-curvature-functional} prouvent que le réseau n’est pas scalaire.

\subsection{Paquet OS continu et reconstruction}

Pour une direction $\nu$, soit $\mathfrak A_{+,\nu}$ l’adhérence de l’union des
algèbres précédentes à support compact dans $\{x_\nu>0\}$. L’état
statique, la loi des piles et l’opérateur de transfert ne sont pas définis
séparément. Sur un écran $\alpha$, soient
$\widehat\mu_{\alpha,g}^{\rm YM}$ la marginale de
\eqref{eq:pdf2-ym-gibbs-measure},
$\nu_{\alpha,g}:=(\pi_\alpha^{\rm phys})_*
\widehat\mu_{\alpha,g}^{\rm YM}$ sa loi sur le quotient et
$K_{\alpha,g}(t)=e^{-tD_{\alpha,g}^{\rm YM}}$ son noyau.
Pour des temps ordonnés $t_1\le\cdots\le t_n$, la marginale euclidienne de la
même mesure est
\begin{equation}\label{eq:pdf2-ym-owner-euclidean-time-marginal}
 \mathbb P_{\alpha,g}^{E,(\nu;t_1,\ldots,t_n)}
 :=(\operatorname{ev}_{t_1}^{(\nu)},\ldots,
    \operatorname{ev}_{t_n}^{(\nu)})_*
    \widehat\mu_{\alpha,g}^{\rm YM}.
\end{equation}
L'identité de désintégration qui fait partie de \eqref{eq:pdf2-g-gau-transfer-euclidean-dato} la calcule, sans la remplacer par une chaîne auxiliaire :
\begin{equation}\label{eq:pdf2-ym-owner-euclidean-markov-identification}
 d\mathbb P_{\alpha,g}^{E,(\nu;t_1,\ldots,t_n)}
 =\nu_{\alpha,g}(dy_1)
   \prod_{j=1}^{n-1}K_{\alpha,g}
      (t_{j+1}-t_j;y_j,dy_{j+1}).
\end{equation}
Si $M_A$ désigne la multiplication par un lecteur borné $A$, ses corrélations
sont donc
\begin{equation}\label{eq:pdf2-ym-owner-schwinger-transfer}
 \begin{aligned}
 S_{\alpha,n}^{(\nu,g)}(A_1,t_1;\ldots;A_n,t_n)
 &:=\int\prod_{j=1}^n(\tau_{t_je_\nu}A_j)(x)\,
       d\widehat\mu_{\alpha,g}^{\rm YM}(x)\\
 &=\langle\mathbf1,
 M_{A_1}K_{\alpha,g}(t_2-t_1)M_{A_2}\cdots
 K_{\alpha,g}(t_n-t_{n-1})M_{A_n}\mathbf1
 \rangle_{L^2(\nu_{\alpha,g})}.
 \end{aligned}
\end{equation}
La loi de la pile stationnaire à ces mêmes temps est cette marginale euclidienne, écrite au moyen de sa désintégration :
\begin{equation}\label{eq:pdf2-ym-owner-markov-stack-law}
 d\mathbb P_{\alpha,g}^{(\nu;t_1,\ldots,t_n)}
 :=d\mathbb P_{\alpha,g}^{E,(\nu;t_1,\ldots,t_n)}
 =\nu_{\alpha,g}(dy_1)
   \prod_{j=1}^{n-1}K_{\alpha,g}(t_{j+1}-t_j;y_j,dy_{j+1}).
\end{equation}
Intégrer successivement $y_n,y_{n-1},\ldots,y_1$ donne l’identité exacte
\begin{equation}\label{eq:pdf2-ym-owner-stack-transfer-identity}
 S_{\alpha,n}^{(\nu,g)}(A_1,t_1;\ldots;A_n,t_n)
 =\int\prod_{j=1}^nA_j(y_j)\,
   d\mathbb P_{\alpha,g}^{(\nu;t_1,\ldots,t_n)}(y).
\end{equation}
Si tous les temps coïncident, les noyaux sont l’identité et
\begin{equation}\label{eq:pdf2-ym-owner-static-is-stack}
 S_{\alpha,n}^{(\nu,g)}(A_1,t;\ldots;A_n,t)
 =\int A_1\cdots A_n\,d\nu_{\alpha,g}.
\end{equation}
Ainsi, l'état de Schwinger statique est la marginale à temps égal de la même mesure euclidienne, et sa désintégration est la loi de Markov ; ce ne sont pas deux fonctionnelles que l'on identifie ensuite par leur nom. La compatibilité de $\widehat J_{\alpha,g}$ avec la mesure et le noyau rend \eqref{eq:pdf2-ym-owner-stack-transfer-identity} compatible avec toute superposition. Sa limite définit
\begin{equation}\label{eq:pdf2-ym-owner-schwinger-functionals}
 S_n^{\rm YM}(A_1,t_1;\ldots;A_n,t_n)
 :=\lim_{\alpha}S_{\alpha,n}^{(\nu,g)}
 (A_{1,\alpha},t_1;\ldots;A_{n,\alpha},t_n).
\end{equation}

Les insertions de $\Phi$ s’obtiennent en dérivant en $t=0$ les unitaires
\eqref{eq:pdf2-ym-owner-curvature-weyl}. Puisque les $Z_{\alpha,c}$ transportés
sont centrés, indépendants et uniformément bornés, le lemme de Hoeffding donne,
niveau par niveau,
\begin{equation}\label{eq:pdf2-ym-owner-subgaussian-bound}
 \left\langle\Omega_\alpha,
 e^{t\Phi_\alpha(h)}\Omega_\alpha\right\rangle
 \le \exp\!\left(\frac9{16}t^2\|P_\alpha h\|_{L^2}^2\right).
\end{equation}
La borne passe à la limite par
\eqref{eq:pdf2-ym-owner-smeared-cauchy-derived} ; toutes les dérivées sont continues dans
la topologie de Schwartz et définissent des distributions tempérées.

La forme à huit faces se prolonge d’abord à toute pile finie. Si $F$ dépend des
$N$ couches positives d’un écran $\alpha$, la réflexion de
\eqref{eq:pdf2-ym-owner-markov-stack-law} et le conditionnement sur la couche centrale donnent
\begin{equation}\label{eq:pdf2-ym-owner-finite-stack-os}
 \begin{aligned}
 &\int \overline{F(\Theta_\nu y_{-N},\ldots,
                         \Theta_\nu y_{-1})}
          F(y_1,\ldots,y_N)\,d\Omega_{\alpha,N}^{(\nu,g)}(y)\\
 &\qquad=
 \int\nu_{\alpha,g}(dz)\left|
  \int K_{\alpha,g}(a_\alpha/2;z,dy_1)
       \prod_{j=1}^{N-1}K_{\alpha,g}(a_\alpha;y_j,dy_{j+1})
       F(y_1,\ldots,y_N)
 \right|^2\ge0.
 \end{aligned}
\end{equation}
Tout cylindre de demi-espace appartient à une pile finie ; par conséquent, \eqref{eq:pdf2-ym-owner-finite-stack-os} matérialise la forme OS avant l'adhérence. Dans le quotient OS, si $F$ possède des insertions à des temps positifs, soit $(\tau_sF)(A_1,t_1;\ldots;A_n,t_n):=
F(A_1,t_1+s;\ldots;A_n,t_n+s)$. On définit l'opérateur de translation par
\begin{equation}\label{eq:pdf2-ym-owner-os-transfer-operator}
 T_{\rm OS}^{(\nu)}(s)[F]:=[\tau_sF],
 \qquad s\ge0.
\end{equation}
Pour des cylindres $F,G$ de demi-espace, l’égalité euclidienne--Markov
\eqref{eq:pdf2-ym-owner-euclidean-markov-identification} et la propriété de
semi-groupe donnent générateur par générateur
\begin{equation}\label{eq:pdf2-ym-owner-os-markov-intertwining}
 \begin{aligned}
 \langle[F],T_{\rm OS}^{(\nu)}(s)[G]\rangle_{\rm OS}
 &=S_g^{\rm YM}((\Theta_\nu F)^*\tau_sG)\\
 &=\left\langle
    \mathcal V_{\infty,\nu}[F],
    e^{-sD_{\rm curv,g}^{\rm YM}}
    \mathcal V_{\infty,\nu}[G]
   \right\rangle .
 \end{aligned}
\end{equation}
L’identité montre à la fois que les classes nulles restent nulles et que,
par densité des cylindres,
\begin{equation}\label{eq:pdf2-ym-owner-os-generator-identification}
 \mathcal V_{\infty,\nu}T_{\rm OS}^{(\nu)}(s)
 \mathcal V_{\infty,\nu}^{-1}
 =e^{-sD_{\rm curv,g}^{\rm YM}}.
\end{equation}
C’est la flèche explicite mesure euclidienne $\to$ désintégration $\to$
transfert OS $\to$ hamiltonien.

\begin{proposition}[Paquet d’Osterwalder--Schrader continu]
\label{prop:pdf2-ym-owner-continuum-os-package}
Les fonctionnelles \eqref{eq:pdf2-ym-owner-schwinger-functionals} satisfont :
\begin{enumerate}
 \item normalisation, hermiticité et symétrie bosonique ;
 \item covariance euclidienne et régularité tempérée ;
 \item positivité par réflexion dans chacune des quatre directions,
 \begin{equation}\label{eq:pdf2-ym-owner-continuum-reflection-positive}
  \sum_{i,j}\overline{a_i}a_j\,
  S^{\rm YM}\!\left((\Theta_\nu F_i)^*F_j\right)\ge0,
  \qquad F_i\in\mathfrak A_{+,\nu};
 \end{equation}
 \item localité et propriété de cluster exponentielle : si $A,B$ sont locaux et centrés et si une
 translation de longueur $r$ sépare leurs supports, alors
 \begin{equation}\label{eq:pdf2-ym-owner-continuum-cluster}
  |S^{\rm YM}(A\,\tau_rB)|
  \le e^{-3\kappa_{\rm av}r}\,
       \|A\Omega_\infty\|\,\|B\Omega_\infty\|.
 \end{equation}
\end{enumerate}
La reconstruction OS de l'un quelconque des quatre demi-réseaux produit le même espace cyclique, le même vide et l'hamiltonien $D_{\rm curv,g}^{\rm YM}$. Le prolongement OS produit des champs de Wightman locaux et une représentation de Poincaré dont le spectre satisfait
\begin{equation}\label{eq:pdf2-ym-owner-wightman-spectrum}
 \operatorname{Spec}(P^0,\mathbf P)\subseteq\overline V_+,
 \qquad P^0=D_{\rm curv,g}^{\rm YM}\ge0.
\end{equation}
\end{proposition}

\begin{proof}
La normalisation et l’hermiticité proviennent de l’état ; la symétrie et la localité, de la
commutativité euclidienne à supports disjoints. La covariance et la continuité sont
\eqref{eq:pdf2-ym-owner-field-e4-covariance}--
\eqref{eq:pdf2-ym-owner-strong-e4} ; la régularité est
\eqref{eq:pdf2-ym-owner-subgaussian-bound}. Pour
$F_i$ cylindriques, tous vivent dans une superposition et une pile finies ;
\eqref{eq:pdf2-ym-owner-finite-stack-os} et la polarisation prouvent
\eqref{eq:pdf2-ym-owner-continuum-reflection-positive}. L’approximation en norme OS
et Cauchy--Schwarz l’étendent à $\mathfrak A_{+,\nu}$.

Après une rotation euclidienne, la séparation spatiale devient un temps positif.
L’identité pile--transfert
\eqref{eq:pdf2-ym-owner-stack-transfer-identity} et l’opérateur
\eqref{eq:pdf2-ym-owner-os-transfer-operator} écrivent la corrélation centrée sous la forme
\[
 \langle A^*\Omega_\infty,
 e^{-rD_{\rm curv,g}^{\rm YM}}B\Omega_\infty\rangle.
\]
La borne spectrale \eqref{eq:pdf2-ym-owner-curvature-sector-gap} donne \eqref{eq:pdf2-ym-owner-continuum-cluster}. Enfin, les connecteurs \eqref{eq:pdf2-ym-owner-os-connectors} forment un système dirigé ; leur réunion est dense d'après \eqref{eq:pdf2-ym-owner-smooth-core}. L'identité \eqref{eq:pdf2-ym-owner-common-os-hamiltonian} identifie les quatre adhérences et leurs générateurs à $D_{\rm curv,g}^{\rm YM}$. Les hypothèses déjà vérifiées de régularité, de covariance, de réflexion et de cluster sont les entrées du prolongement OS ; son théorème de reconstruction donne la localité de Wightman et \eqref{eq:pdf2-ym-owner-wightman-spectrum}.
\end{proof}

\subsection{Publication typée comme théorie de Yang--Mills}

L’identification ne repose pas seulement sur l’évaluation d’un clover sur une connexion donnée.
Soit $\mathcal R_G^{\rm loc}$ la $*$-algèbre de lecteurs formée par toutes les
fonctions de classe des holonomies d’incidence et de plaquette, leurs transports vers un
sommet de base et les polarisations du lecteur $\mathcal P^{F^2}$. Soit
\begin{equation}\label{eq:pdf2-ym-owner-curvature-fibre}
 \mathcal E_G:=\Lambda^2(\mathbb R^4)^*\otimes\mathfrak g
\end{equation}
avec le produit invariant. Dans chaque cellule, la polarisation de
$\mathcal P^{F^2}$ donne un gramien positif $G_{\alpha,c}^{F}$ sur
$\mathcal E_G$. On quotiente par son noyau et on applique sa racine positive ; le lecteur
normalisé résultant est noté $Y_{\alpha,c}(v)$ et vérifie
\begin{equation}\label{eq:pdf2-ym-owner-curvature-cell-gram}
 \langle Y_{\alpha,c}(v),Y_{\alpha,d}(w)\rangle
 =\delta_{cd}\langle v,w\rangle_{\mathcal E_G}.
\end{equation}
Si $d\subset c$, soit $T_{d\leftarrow c}$ le transport de repère déjà présent dans l'holonomie de surface. L'opérateur d'entrelacement de courbure complète est
\begin{equation}\label{eq:pdf2-ym-owner-curvature-intertwiner}
 I_{\alpha\beta}^{F}Y_{\alpha,c}(v)
 :=\sum_{d\subset c}\sqrt{\frac{|d|}{|c|}}\,
 Y_{\beta,d}\!\left(\operatorname{Ad}_{T_{d\leftarrow c}}v\right).
\end{equation}
L’invariance du produit de $\mathfrak g$, la disjonction des sous-cellules et la
composition des transports prouvent
\begin{equation}\label{eq:pdf2-ym-owner-curvature-intertwiner-identities}
 (I_{\alpha\beta}^{F})^*I_{\alpha\beta}^{F}=I,
 \qquad
 I_{\beta\gamma}^{F}I_{\alpha\beta}^{F}=I_{\alpha\gamma}^{F}.
\end{equation}
Pour $h\in C_c^\infty(\mathbb R^4;\mathcal E_G)$, on pose
\begin{equation}\label{eq:pdf2-ym-owner-f-alpha}
 \mathbf F_\alpha(h):=
 \sum_{c\in\mathcal C_\alpha^{\rm cell}}
 Y_{\alpha,c}\!\left(
   |c|^{-1/2}\int_c h(x)\,d^4x\right).
\end{equation}
Dans une superposition commune, le même calcul de
\eqref{eq:pdf2-ym-owner-smeared-cross-gram} donne maintenant
\begin{equation}\label{eq:pdf2-ym-owner-curvature-cross-gram}
 \left\langle I_{\alpha\gamma}^{F}\mathbf F_\alpha(h),
 I_{\beta\gamma}^{F}\mathbf F_\beta(k)\right\rangle
 =\langle P_\alpha h,P_\beta k\rangle_{L^2(\mathcal E_G)}.
\end{equation}
Par conséquent, $\mathbf F_\alpha(h)$ est de Cauchy par convergence forte des projections,
et non par une nouvelle hypothèse. Le même appariement de Haar et des innovations de
\eqref{eq:pdf2-ym-owner-publication-into-physical}, appliqué maintenant à la base
orthonormalisée de $G_{\alpha,c}^{F}$, réalise la colimite par une isométrie
locale et covariante de jauge dans
$L^2(\widehat X_\infty,\widehat\mu_{\infty,g}^{\rm YM})$. Sa limite définit la
distribution aléatoire
covariante de jauge
\begin{equation}\label{eq:pdf2-ym-owner-distributional-curvature}
\begin{aligned}
 \mathbf F_{\mu\nu}&\in
 \mathcal S'(\mathbb R^4;\mathfrak g)\ \widehat\otimes\
 L^2(\widehat X_\infty,\widehat\mu_{\infty,g}^{\rm YM}),\\
 U(a)\mathbf F(h)U(a)^*&=\mathbf F(a\!\cdot h),\\
 u\mathbf F(h)u^{-1}&=\mathbf F(\operatorname{Ad}_u h),
 \qquad a\in E(4),\ u\in G.
\end{aligned}
\end{equation}
Son canal de classe spectrale d'ordre trois est précisément $\Phi$. La relation de surface finie
\begin{equation}\label{eq:pdf2-ym-owner-surface-product}
 \operatorname{Hol}^{(\alpha)}_p
 =\prod_{p'\subset p}^{\longrightarrow}
  T_{p'}\operatorname{Hol}^{(\beta)}_{p'}T_{p'}^{-1}
\end{equation}
passe à la fermeture cofinale car tous ses produits croisés sont calculés dans une
superposition. En particulier, le lecteur clover aléatoire, et non seulement son évaluation dans une
connexion lisse, satisfait l’identité des petites boucles
\begin{equation}\label{eq:pdf2-ym-owner-holonomy-curvature-relation}
 \mathbf F(h)=L^2\!\mathord{-}\!\lim_{\alpha}
 a_\alpha^4\sum_x
 \left\langle(P_\alpha h)(x),F_{{\rm cl},\alpha}(x)\right\rangle,
 \qquad
 \delta_{\rm gauge}S_n^{\rm YM}=0,
 \qquad
 \operatorname{Alt}_{\lambda\mu\nu}D_\lambda\mathbf F_{\mu\nu}=0.
\end{equation}
La première égalité est \eqref{eq:pdf2-ym-owner-curvature-cross-gram} écrite dans le représentant clover ; la seconde est l'identité de Ward \eqref{eq:pdf2-ym-ward-finite} passée à la limite, et la troisième est l'annulation de la frontière d'une frontière dans le produit de surface. Aucune n'utilise le terme de masse.

Pour éviter de définir la cible par la conclusion que l’on veut obtenir, soit
$\mathbf{YM}_4(G,g)$ la classe des quintuplets locaux
$(\mathfrak A,\mu_g,K_g,\operatorname{Hol},\mathbf F)$ qui vérifient les
équations vérifiables suivantes : (i) la densité de Gibbs est
\eqref{eq:pdf2-ym-gibbs-measure} ; (ii) $K_g$ est markovien, réversible et satisfait
\eqref{eq:pdf2-ym-schwinger-dyson-finite}--\eqref{eq:pdf2-ym-ward-finite} ;
(iii) holonomie, clover et courbure satisfont
\eqref{eq:pdf2-ym-transport-surface-bianchi} et
\eqref{eq:pdf2-ym-owner-holonomy-curvature-relation} ; et (iv) leurs fonctionnelles sont la
même loi de piles de \eqref{eq:pdf2-ym-owner-stack-transfer-identity} et satisfont le
paquet OS\@. L’objet construit est, composante par composante,
\begin{equation}\label{eq:pdf2-ym-owner-qym-object}
\begin{aligned}
 \mathfrak Q_{\rm YM}^{(4)}(G,g):={}&
 \bigl(\{\mathfrak A_{\rm YM}(O)\}_O,
 S_g^{\rm YM},U,\Theta,\Omega_{\infty,g},\\
 &\widehat\mu_{\infty,g}^{\rm YM},
 \{S_{m,g},\mathscr S_{m,g}\}_m,K_g,
 \operatorname{Hol},\mathbf F,\\
 &\mathscr H_{\rm OS},D_{\rm curv,g}^{\rm YM}\bigr)
 \in\mathbf{YM}_4(G,g).
\end{aligned}
\end{equation}
L’application de publication n’omet aucune de ces composantes :
\begin{equation}\label{eq:pdf2-ym-owner-typed-publication-map}
 \begin{aligned}
 \mathcal P_{{\rm YM},g}:\mathfrak G_{\rm gau}(G)&\longrightarrow
 \mathbf{YM}_4(G,g),\\
 (\widehat\mu^0,\mathbb P^{\rm phys},D^0,\Theta,W,
   \operatorname{Hol},\mathcal P^{F^2})
 &\longmapsto
 (\widehat\mu_g^{\rm YM},S_g,\mathscr S_g,K_g,\mathbb P_g^{\rm phys},
  S_g^{\rm YM},\mathfrak A_{\rm YM},D_{\rm curv,g}^{\rm YM},
  \mathbf F,\operatorname{Hol}).
 \end{aligned}
\end{equation}
Les $S_n^{\rm YM}$ sont ainsi exactement les fonctions de Schwinger de la mesure dépendante de $g$, de la loi de Markov et de l'opérateur de transfert déjà construits.

Il reste à vérifier que $S_{m,g}$ est l’action de Yang--Mills, et non seulement une
fonction portant ce nom. Soit
$A\in C_c^3(\mathbb R^4;T^*\mathbb R^4\otimes\mathfrak g)$ une connexion lisse dans
une carte où le logarithme des petites boucles est unique. Pour la plaquette
$p_{\mu\nu}(x)$, le développement ordonné de Magnus donne
\begin{equation}\label{eq:pdf2-ym-owner-plaquette-magnus}
 \log\operatorname{Hol}_{p_{\mu\nu}(x)}(A)
 =a_m^2F_{A,\mu\nu}(x)
  +\frac{a_m^3}{2}(D_\mu+D_\nu)F_{A,\mu\nu}(x)
  +R_{m,\mu\nu}(x),
\end{equation}
avec $\|R_{m,\mu\nu}(x)\|\le C_G a_m^4
(\|D^2F_A\|_{\infty,p}+\|F_A\|_{\infty,p}^2)$. La moyenne des quatre orientations annule le terme impair et définit
\begin{equation}\label{eq:pdf2-ym-owner-clover}
 F_{{\rm cl},m}(x)
 :=\frac1{4a_m^2}\sum_{p\ni x}
  \operatorname{Ad}_{T_p}\log\operatorname{Hol}_p,
 \qquad
 \|F_{{\rm cl},m}(x)-F_A(x)\|
 \le C_G a_m^2(\|D^2F_A\|_\infty+\|F_A\|_\infty^2).
\end{equation}
La définition de la racine \eqref{eq:pdf2-ym-action-root} est précisément la polarisation de ces six composantes. En substituant \eqref{eq:pdf2-ym-owner-plaquette-magnus} dans chaque générateur de cette polarisation, on calcule
\begin{equation}\label{eq:pdf2-ym-owner-action-root-smooth-chart}
 \xi_{m,b}(A)
 =a_m^2\operatorname{vec}_{\mu<\nu}F_{A,\mu\nu}(x_b)
  +r_{m,b}(A),
 \qquad
 \|r_{m,b}(A)\|\le C_{G,A}a_m^4.
\end{equation}
Les termes croisés impairs s’annulent par les quatre orientations du clover ;
les termes pairs ne sont inclus qu’une seule fois par
\eqref{eq:pdf2-ym-action-factor-partition}. Par somme de Riemann et
Cauchy--Schwarz,
\begin{equation}\label{eq:pdf2-ym-owner-action-error}
 \left|S_{m,g}(A)-\frac1{4g^2}
  a_m^4\sum_x\|F_{{\rm cl},m}(x)\|^2\right|
 \le C_{G,A,g}a_m^2,
\end{equation}
et, en utilisant la borne de \eqref{eq:pdf2-ym-owner-clover},
\begin{equation}\label{eq:pdf2-ym-owner-f2}
 S_{m,g}(A)\longrightarrow
 S_{{\rm YM},g}(A):=\frac1{4g^2}
 \int_{\mathbb R^4}\|F_A(x)\|^2\,d^4x.
\end{equation}
Les équations de Schwinger--Dyson et de Ward ont déjà été calculées sur la mesure qui utilise
cette action. L’écart, en revanche, se déduit ensuite de la conjugaison unitaire de
\eqref{eq:pdf2-ym-g-dynamics} et du témoin
\eqref{eq:pdf2-ym-g-gap-before-limit} ; il n’a pas été utilisé pour choisir $g$, $S_{m,g}$ ni
$\widehat\mu_{m,g}^{\rm YM}$. Tel est le lien non circulaire entre action, loi et écart.

\begin{theorem}[Clôture de Yang--Mills pour tout groupe de Lie compact, connexe et simple]
\label{thm:pdf2-ym-cierre-global}
La structure complète $\mathfrak G_{\rm gau}$ construit une famille projective locale pour tout groupe de Lie compact, connexe et simple. L'application \eqref{eq:pdf2-ym-owner-typed-publication-map} publie une théorie quantique de Yang--Mills non triviale sur $\mathbb R^4$ : ses fonctionnelles satisfont l'ensemble continu des conditions OS et son prolongement satisfait les axiomes locaux de Wightman d'après la proposition~\ref{prop:pdf2-ym-owner-continuum-os-package} ; ses holonomies et sa courbure distributionnelle satisfont \eqref{eq:pdf2-ym-owner-holonomy-curvature-relation}, et les quatre reconstructions coinduisent le même hamiltonien $D_{\infty,g}^{\rm YM}$, avec un vide simple. Sa loi est $\widehat\mu_{\infty,g}^{\rm YM}$, construite par la famille compatible d'actions locales $\{S_{m,g},\mathscr S_{m,g}\}_m$, et satisfait les équations de Schwinger--Dyson et de Ward \eqref{eq:pdf2-ym-schwinger-dyson-finite}--\eqref{eq:pdf2-ym-ward-finite}. Le secteur cyclique de courbure \eqref{eq:pdf2-ym-owner-curvature-cyclic-sector} est réducteur, contient le témoin \eqref{eq:pdf2-ym-owner-wc-curvature-functional} et possède
\begin{equation}\label{eq:pdf2-ym-final-gap}
 \boxed{\Delta_{\rm YM}^{\rm HMT}=3\kappa_{\rm av}>0},
 \qquad
 m_{\rm gap}^{\rm HMT}
 =\frac{\hbar_{\rm int}}{c_{\rm int}}\Delta_{\rm YM}^{\rm HMT}>0.
\end{equation}
Le réseau local, l’action forte de $E(4)$, les distributions de Schwinger,
l’holonomie et $\mathbf F$ proviennent du même état complet. Le lecteur exact
$\mathcal P^{F^2}$ définit la densité de la loi ; le clover est sa carte lisse et
\eqref{eq:pdf2-ym-owner-f2} identifie son action classique. L’écart appartient au secteur collectif
de courbure invariant de jauge et n’attribue pas de masse élémentaire au gluon.
\end{theorem}

\begin{proof}
La projectivité cinématique est \eqref{eq:pdf2-ym-owner-projectivity} ; la racine
\eqref{eq:pdf2-ym-action-root}, le cocycle
\eqref{eq:pdf2-ym-action-cocycle} et la densité
\eqref{eq:pdf2-ym-gibbs-measure} construisent d’abord l’action et la mesure au
couplage $g$. Le transport facteur par facteur
\eqref{eq:pdf2-ym-gibbs-transport-factor}--
\eqref{eq:pdf2-ym-full-algebra-identities} prouve la projectivité, la préservation de
l’état, des produits et de $*$ ;
\eqref{eq:pdf2-ym-transport-surface-bianchi} établit clover et Bianchi. L’indice exhaustif
\eqref{eq:pdf2-ym-owner-exhaustive-index}, la forme
\eqref{eq:pdf2-ym-owner-full-form} et le circuit
\eqref{eq:pdf2-ym-owner-local-circuit} construisent le générateur commun. La réduction
physique et son noyau quotient sont
\eqref{eq:pdf2-ym-owner-physical-reduction}--
\eqref{eq:pdf2-ym-owner-D}. La conjugaison
\eqref{eq:pdf2-ym-g-dynamics} les transporte à la mesure dépendant de $g$ sans
altérer leur spectre. La décomposition de Hoeffding prouve le vide et la borne, et
\eqref{eq:pdf2-ym-owner-wc-invariance}--
\eqref{eq:pdf2-ym-owner-wc-eigenvalue} exhibent un vecteur physique qui atteint le
seuil. L’entrelacement nonadique transporte forme, vide et témoin à la limite, où
\eqref{eq:pdf2-ym-owner-limit-gap}--
\eqref{eq:pdf2-ym-owner-exact-gap} fixent l’écart exact. Les quatre
factorisations \eqref{eq:pdf2-ym-owner-four-os}, l’inclusion des images
\eqref{eq:pdf2-ym-owner-os-range-inclusion} et les connecteurs
\eqref{eq:pdf2-ym-owner-os-connectors} identifient leurs reconstructions à ce
même générateur. L’holonomie
\eqref{eq:pdf2-ym-owner-incidence-holonomy} place matériellement $W_c$ dans le secteur
de courbure, où \eqref{eq:pdf2-ym-owner-curvature-sector-gap} prouve l’écart
exact. Le système \eqref{eq:pdf2-ym-owner-block-intertwiner} contrôle les produits
croisés au moyen de \eqref{eq:pdf2-ym-owner-smeared-cross-gram} ;
\eqref{eq:pdf2-ym-full-publication-restriction}--
\eqref{eq:pdf2-ym-full-publication-products} élèvent ce premier chaos à l’isomorphisme
de toute l’algèbre. La convergence forte de $P_\alpha$ produit $\Phi$ ; les bornes
\eqref{eq:pdf2-ym-owner-holonomy-e4-estimate} et
\eqref{eq:pdf2-ym-owner-orbit-smearing-continuity} prouvent en outre la continuité
$E(4)$ des holonomies, des plaquettes et de clover. La désintégration euclidienne
\eqref{eq:pdf2-ym-owner-euclidean-markov-identification}, l’identité exacte
\eqref{eq:pdf2-ym-owner-stack-transfer-identity} et l’entrelacement
\eqref{eq:pdf2-ym-owner-os-markov-intertwining} relient la même mesure,
l’état de Schwinger, la pile de Markov, le transfert OS et le générateur de l’écart. La
proposition~\ref{prop:pdf2-ym-owner-continuum-os-package} prolonge les
quatre formes finies aux
demi-réseaux complets, prouve la régularité et la propriété de cluster et reconstruit l’hamiltonien. Pour
finir, \eqref{eq:pdf2-ym-owner-distributional-curvature}--
\eqref{eq:pdf2-ym-owner-typed-publication-map} identifient ces mêmes fonctionnelles
comme les fonctionnelles de Schwinger de la théorie de jauge locale ; les intégrations par parties
\eqref{eq:pdf2-ym-schwinger-dyson-finite}--\eqref{eq:pdf2-ym-ward-finite} prouvent
Schwinger--Dyson et Ward. Enfin,
\eqref{eq:pdf2-ym-owner-plaquette-magnus}--\eqref{eq:pdf2-ym-owner-f2} identifient
l’action classique avant d’appliquer l’équivalence spectrale. La chaîne n’utilise pas
l’écart pour choisir la mesure.
\end{proof}

\begin{falsifier}[Effondrement de la dynamique produit]
Si l’on remplace $D_m^{\rm prod}$ par $D_m^{\rm fib}$, tous les secteurs non
constants reçoivent la même valeur propre. Un vecteur de Hoeffding avec $|S|=2$ devrait
avoir l’énergie $6\kappa_{\rm av}$ par \eqref{eq:pdf2-ym-hoeffding}, mais le
semi-groupe simple lui attribuerait $3\kappa_{\rm av}$. L’identification des deux
générateurs est fausse.
\end{falsifier}

\begin{falsifier}[Confusion entre inclusion généalogique et moyenne de bloc]
Si l'on remplace $I_{\alpha\beta}^{\rm curv}$ par $\widehat J_{\alpha\beta}$ dans \eqref{eq:pdf2-ym-owner-smeared-cross-gram}, une coordonnée grossière persiste avec un coefficient égal à un tandis que le coefficient fin change selon le quotient des échelles. Pour $a_{m+1}=a_m/9$, le produit croisé ancestral porte un facteur $1/81$ et ne converge pas vers la norme diagonale. L'identité \eqref{eq:pdf2-ym-owner-smeared-cauchy-derived} échoue ; c'est pourquoi les deux applications restent séparées.
\end{falsifier}

\begin{falsifier}[Ablation de la fibre d’incidence]
Si l’on remplace $\widehat X_m$ par sa marginale de liens, on efface $q_c$ et, avec elle,
l’observable physique $W_c$. L’application d’oubli ne reflète pas la sous-algèbre de jauge complète ;
elle ne peut donc pas être utilisée pour nier la valeur propre de l’objet source. Objet
remplacé ; conclusion non transférable.
\end{falsifier}

\begin{falsifier}[Remplacement du lecteur exact par sa carte lisse]
Si l’on utilise directement le terme asymptotique de
\eqref{eq:pdf2-ym-owner-f2} comme densité sur un écran fini, on efface le reste
contrôlé de \eqref{eq:pdf2-ym-owner-action-error} et l’on perd le cocycle exact
\eqref{eq:pdf2-ym-action-cocycle}. La densité finie régissante est
$Z_{m,g}^{-1}e^{-\mathcal P_m^{F^2}/(4g^2)}$ ; le clover et la limite lisse
l’identifient à l’action classique, mais ne la remplacent pas.
\end{falsifier}

\paragraph{Provenance et statut.}
La séparation entre fibre simple et générateur produit, la mesure commune, la compression
de jauge, les quatre OS, la forme limite et le témoin physique sont
\texttt{RÉSULTAT\_RÉCUPÉRÉ} et \texttt{EXACT\_INTERNE} dans le domaine de jauge
régulier HMT enregistré. La désintégration à temps non uniformes, l'indice
exhaustif des espérances, le noyau quotient, la matrice $B^+$, l'action explicite
sur la colimite, le système cofinal de blocs, le paquet OS continu, la courbure
distributionnelle et l'holonomie d'incidence sont
\texttt{FORMALISATION\_NOUVELLE} : ils typent et satisfont les liens qui auparavant n'apparaissaient
que nominalement, sans introduire le gap comme donnée. Leur composition complète
établit le théorème précédent.
La racine positive du lecteur $F^2$, sa densité de Gibbs, le transport local compatible,
les équations de Schwinger--Dyson/Ward, l'isomorphisme algébrique complet, les bornes
de continuité des holonomies et l'identité pile--transfert sont
\texttt{FORMALISATION\_NOUVELLE}. Leur rôle est de satisfaire matériellement les quatre
ponts ; l'identification clover--action classique est une conséquence du développement
de Magnus et non une prémisse spectrale.

\endgroup
% Cuerpo científico recuperado y distribuido en su residencia terminal.
\section{Théorie de jauge euclidienne et gap de Yang--Mills}

La famille projective de jauge construite à partir de l'état complet possède
la covariance locale, la localité, la positivité par réflexion, une action
fortement continue de \(E(4)\), des fonctionnelles de Schwinger compatibles et un
vide simple. Les quatre reconstructions co-induisent le même Hamiltonien et
le secteur cyclique de courbure possède le gap exact
\[
  \Delta_{\rm YM}^{\rm HMT}=3\kappa_{\rm av}>0.
\]
Le lecteur de courbure et sa limite classique proviennent de la même loi. Le gap
appartient au secteur collectif invariant de jauge et n'attribue pas de masse élémentaire
au gluon.




\HMTFlushCurrentChapter
\chapter{Réalisation algébrique des classes de Hodge par projecteurs cohomologiques compatibles et descente rationnelle}
\begingroup
\renewcommand{\chapter}[2][]{}%
\chapter{Image cohomologique et composition vers Hodge}

\section{Objet source, observations finies et publication}

Soient $X$ une variété projective lisse complexe et $p\ge0$. L'objet produit par
la construction commune et reçu ici est la structure cohomologique complète
$\mathfrak G_{\rm coh}(X,p)$, qui conserve
cartes basées, croisements orientés, cônes de Mayer--Vietoris, retenue,
frontière, chemin, multisection, survie et terminal. À la profondeur $N$, on
écrit
\begin{equation}\label{eq:pdf2-hodge-sistemas}
 \mathscr S_N(X,p):=\operatorname{Surv}_{J_N}(X,p),
 \qquad
 \rho_{N+1,N}:\mathscr S_{N+1}(X,p)\to\mathscr S_N(X,p),
\end{equation}
et les modules d'observations
\begin{equation}\label{eq:pdf2-hodge-observaciones}
 e_N:\mathscr S_N(X,p)\to\mathscr H_N^p(X),
 \qquad
 q_{N+1,N}:\mathscr H_{N+1}^p(X)\to\mathscr H_N^p(X),
\end{equation}
avec le carré
\begin{equation}\label{eq:pdf2-hodge-cuadrado-restriccion}
 e_N\rho_{N+1,N}=q_{N+1,N}e_{N+1}.
\end{equation}
Pour
\[
 \alpha\in\operatorname{Hdg}^p(X)_{\mathbb Q}
 :=H^{2p}(X,\mathbb Q)\cap H^{p,p}(X),
\]
soit $q_N\alpha$ sa famille d'observations finies. La fibre qui doit
être prolongée est
\begin{equation}\label{eq:pdf2-hodge-fibra-alpha}
 \mathscr F_N(\alpha):=
 \{s\in\mathscr S_N(X,p):e_N(s)=q_N\alpha\}.
\end{equation}

Le réalisateur et la publication sont typés par
\begin{equation}\label{eq:pdf2-hodge-realizador}
 R_{\rm Hdg}:X_\chi(X,p)\to
 K_0\!\left(\operatorname{Perf}^{\ge p}(X)\right)_{\mathbb Q},
\end{equation}
\begin{equation}\label{eq:pdf2-hodge-publicacion}
 \operatorname{cl}_{X,p}\circ\operatorname{ch}^{\rm loc}_p
 \circ R_{\rm Hdg}
 =\operatorname{ev}^{\rm Hdg}_{\infty,X,p}.
\end{equation}
C'est l'identité d'action du lecteur déjà construite dans l'objet complet
\eqref{eq:pdf2-g-coh-limite-lector-dato} ; elle n'introduit pas depuis le codomaine
l'antécédent d'une classe de Hodge. Le propriétaire exact précédent construit cette
existence et ce chapitre déploie la composition. Le lecteur
doit rester dans l'ordre
\[
 \mathfrak G_{\rm coh}(X,p)
 \longrightarrow X_\chi(X,p)
 \xrightarrow{R_{\rm Hdg}}
 K_0\!\left(\operatorname{Perf}^{\ge p}(X)\right)_{\mathbb Q}
 \xrightarrow{\operatorname{ch}^{\rm loc}_p}
 \operatorname{CH}^p(X)_{\mathbb Q}
 \xrightarrow{\operatorname{cl}_{X,p}}
 \operatorname{Hdg}^p(X)_{\mathbb Q}.
\]
En particulier, la classe finale ne sélectionne rétrospectivement ni une cellule ni
un cycle.

\section{Module relatif et table triangulaire d'incidences}

$N$ étant fixé, soient
\begin{equation}\label{eq:pdf2-hodge-kernels}
 \mathscr K^{\rm S}_{N+1}:=\ker\rho_{N+1,N},
 \qquad
 \mathscr K^{\rm H}_{N+1}:=\ker q_{N+1,N},
 \qquad
 e_{N+1}^{\rm rel}:=e_{N+1}|_{\mathscr K^{\rm S}_{N+1}}.
\end{equation}
Les nouvelles cartes $\Delta J_{N+1}$ produisent
$\mathcal C_\nu=\operatorname{Cone}(\kappa_{a,b})$ et le module
\begin{equation}\label{eq:pdf2-hodge-a-rel}
 \mathscr A_{N+1}^{\rm rel}:=
 H^0\!\left(
  \mathsf W_{\rm rel}\otimes^{\mathbf L}_{\mathbb Q[\Delta J_{N+1}]}
  \Gamma_{\rm rel}
  \oplus\operatorname{Cone}(I-U_{\Gamma_9})_{\rm rel}
 \right).
\end{equation}
Les applications de générateurs
\begin{equation}\label{eq:pdf2-hodge-mapas-relativos}
 a_{N+1}:\mathscr A_{N+1}^{\rm rel}\to\mathscr K^{\rm S}_{N+1},
 \qquad
 b_{N+1}:\mathscr A_{N+1}^{\rm rel}\to\mathscr K^{\rm H}_{N+1}
\end{equation}
satisfont
\begin{equation}\label{eq:pdf2-hodge-factorizacion-relativa}
 e_{N+1}^{\rm rel}a_{N+1}=b_{N+1}.
\end{equation}

\subsection{Présentation générateur par générateur de la nouvelle carte}

Pour que \eqref{eq:pdf2-hodge-factorizacion-relativa} ne soit pas une boîte noire,
on écrit les deux ensembles finis qui interviennent. Soit
\(\mathcal G_{N+1}^{\rm coh}\) l'ensemble de toutes les nouvelles étiquettes
\[
 \nu=(Z,W,a,b,\operatorname{or},\operatorname{Carr},
       \operatorname{Mem},\partial,\gamma,\tau^{\rm term})
\]
qui apparaissent dans \(\Sigma_{\rm coh}^{\rm mult}\), y compris une étiquette
terminale pour chaque composante de
\(\operatorname{Cone}(I-U_{\Gamma_9})_{\rm rel}\). Soit
\(\mathcal R_{N+1}^{\rm coh}\) l'ensemble de relations de trois classes :
Mayer--Vietoris aux croisements \(Z\cap W\), changement d'orientation/retenue et
compatibilité terminale d'un tour complet. On forme
\begin{equation}\label{eq:pdf2-hodge-presentacion-finita}
\begin{aligned}
 F_{N+1}^{\rm coh}&:=\mathbb Q^{(\mathcal G_{N+1}^{\rm coh})},&
 D_{N+1}&:\mathbb Q^{(\mathcal R_{N+1}^{\rm coh})}
             \longrightarrow F_{N+1}^{\rm coh},\\
 \mathscr A_{N+1}^{\rm rel}
   &\cong F_{N+1}^{\rm coh}/\operatorname{im}D_{N+1}.&&
\end{aligned}
\end{equation}
L'action du réalisateur sur la base n'est pas laissée implicite :
\begin{equation}\label{eq:pdf2-hodge-rhdg-generadores}
\begin{aligned}
 R_{{\rm Hdg},N+1}(e_\nu)
   &:=[\operatorname{Cone}(\kappa_{a,b}^{Z,W})],
       &&\nu\ \text{cellulaire},\\
 R_{{\rm Hdg},N+1}(e_{\nu_{\rm term}})
   &:=[\operatorname{Cone}(I-U_{\Gamma_9})_{\nu_{\rm term}}],
       &&\nu_{\rm term}\ \text{terminal}.
\end{aligned}
\end{equation}
L'additivité des cônes dans le triangle de Mayer--Vietoris, l'identité
\(\kappa_{b,a}\tau=-P\kappa_{a,b}\), le cocycle de retenue et l'identité de
retour terminal donnent, respectivement pour les trois classes de relations,
\begin{equation}\label{eq:pdf2-hodge-descenso-relaciones}
 R_{{\rm Hdg},N+1}D_{N+1}=0
 \quad\text{dans }K_0(\operatorname{Perf}^{\ge p}(X))_{\mathbb Q}.
\end{equation}
Par conséquent, \eqref{eq:pdf2-hodge-rhdg-generadores} descend au quotient de
\eqref{eq:pdf2-hodge-presentacion-finita} ; on n'a pas encore supposé
la surjectivité de la classe de cycle.

La base ordonnée
\(\{\omega_\lambda\}_{\lambda\in\Lambda_{N+1}}\) des nouvelles observations
que publie \(\mathscr K^{\rm H}_{N+1}\) étant fixée, la matrice calculable du lecteur est
\begin{equation}\label{eq:pdf2-hodge-matriz-lector}
 C_{N+1}(\lambda,\nu):=
 \left\langle\omega_\lambda^\vee,\,
 \operatorname{cl}_{X,p}\operatorname{ch}^{\rm loc}_p
 R_{{\rm Hdg},N+1}(e_\nu)\right\rangle .
\end{equation}
De \eqref{eq:pdf2-hodge-descenso-relaciones}, on obtient
\begin{equation}\label{eq:pdf2-hodge-cd-cero}
 C_{N+1}D_{N+1}=0.
\end{equation}
Ainsi, les entrées qui suivent ne sont pas des rangs scalaires des cônes : ce sont leurs
coordonnées cohomologiques complètes dans la base \(\omega_\lambda\).

\begin{proposition}[Exhaustivité finie du propriétaire cohomologique]
\label{prop:pdf2-hodge-exhaustividad-owner}
La coordonnée $\Sigma_{\rm coh}^{\rm mult}$ de
$\mathfrak G_{\rm coh}$ induit une décomposition exhaustive de
$\mathscr K^{\rm H}_{N+1}$ en atomes de nouveau cylindre et de différence
terminale. Chaque atome de cylindre est l'incidence principale d'un cône
$\operatorname{Cone}(\kappa_{a,b})$ et apparaît avec un coefficient un ; le
terme restant est exactement
$H^0(\operatorname{Cone}(I-U_{\Gamma_9})_{\rm rel})$.
\end{proposition}

\begin{proof}
L'exhaustivité et la table unitaire sont construites avant cette vue : les
extensions finies proviennent de
\eqref{eq:amp-hodge-surv-n}--\eqref{eq:amp-hodge-ext}, et leur compatibilité est
prouvée dans \cref{prop:amp-hodge-extension-finita}. L'étape coinductive et
l'évaluation sont dérivées dans
\eqref{eq:amp-hodge-xi-limite}--\eqref{eq:amp-hodge-evaluaciones-xi} ;
\eqref{eq:pdf2-g-coh-tabla-dato} enregistre cette construction dans le codomaine
commun. En passant de $J_N$ à $J_{N+1}$, une observation soit
se restreint à une observation antérieure, soit a son support dans une cellule de
$\Delta J_{N+1}$, soit mesure la différence d'un tour complet. Les
trois cas sont disjoints par support, chemin et terminal. Les deux derniers
énumèrent exactement \(\mathcal G_{N+1}^{\rm coh}\) dans
\eqref{eq:pdf2-hodge-presentacion-finita}. L'orientation du cône fixe à
un le coefficient de son incidence principale et
\eqref{eq:pdf2-hodge-descenso-relaciones} identifie les copies liées.
Par conséquent, chaque élément de la base
\(\{\omega_\lambda\}_{\lambda\in\Lambda_{N+1}}\) possède une unique étiquette
maximale et toutes les autres entrées de sa colonne sont strictement antérieures.
Aucune classe \(\alpha\) n'a été énumérée et aucune sortie
algébrique n'a été choisie rétrospectivement.
\end{proof}

Chaque atome de nouvelle observation conserve l'étiquette complète
\[
 \mu=(Z,W,a,b,\operatorname{or},\operatorname{Carr},
       \operatorname{Mem},\partial,\gamma,\tau^{\rm term}).
\]
Ces étiquettes sont ordonnées par profondeur, support, feuillet, chemin, retenue,
mémoire et terminal. Soient $\bar c_\mu$ les atomes résultant de
$\mathscr K^{\rm H}_{N+1}$ et $c_\mu$ les représentants cellulaires de
forme normale dans $\mathscr A_{N+1}^{\rm rel}$. Le caractère localisé d'un
cône publie sa nouvelle incidence avec un coefficient un ; ses faces propres
ont une étiquette antérieure. Il existe donc une table finie strictement
triangulaire
\begin{equation}\label{eq:pdf2-hodge-tabla-triangular}
 b_{N+1}(c_\mu)
 =\bar c_\mu+\sum_{\nu<\mu}B_{\nu\mu}\bar c_\nu,
 \qquad B_{\nu\mu}\in\mathbb Q.
\end{equation}
Le générateur terminal apparaît comme une étiquette propre et ne se réduit pas à sa
phase visible. La relation
$\kappa_{b,a}\tau=-P\kappa_{a,b}$ fixe les signes de
$B_{\nu\mu}$, et le cocycle de retenue fixe leur composition.

\begin{proposition}[Section relative construite avant la classe]
\label{prop:pdf2-hodge-seccion-relativa}
La récurrence
\begin{equation}\label{eq:pdf2-hodge-sigma-recursiva}
 \begin{aligned}
  \sigma_{N+1}^{\rm rel}(\bar c_{\mu_0})&:=c_{\mu_0},\\
  \sigma_{N+1}^{\rm rel}(\bar c_\mu)&:=
  c_\mu-\sum_{\nu<\mu}B_{\nu\mu}
             \sigma_{N+1}^{\rm rel}(\bar c_\nu)
 \end{aligned}
\end{equation}
et l'extension $\mathbb Q$-linéaire définissent
\begin{equation}\label{eq:pdf2-hodge-sigma-right-inverse}
 \sigma_{N+1}^{\rm rel}:\mathscr K^{\rm H}_{N+1}\longrightarrow
 \mathscr A_{N+1}^{\rm rel},
 \qquad b_{N+1}\sigma_{N+1}^{\rm rel}=I.
\end{equation}
En conséquence,
\begin{equation}\label{eq:pdf2-hodge-right-inverse}
 s_{N+1}^{\rm rel}:=a_{N+1}\sigma_{N+1}^{\rm rel}:
 \mathscr K^{\rm H}_{N+1}\longrightarrow\mathscr K^{\rm S}_{N+1},
 \qquad e_{N+1}^{\rm rel}s_{N+1}^{\rm rel}=I.
\end{equation}
Aucune de ces applications ne dépend de $\alpha$.
\end{proposition}

\begin{proof}
La formule \eqref{eq:pdf2-hodge-tabla-triangular} montre que la
matrice de $b_{N+1}$ sur le sous-espace engendré par les $c_\mu$ est
unitriangulaire. Pour la première étiquette, on a
$b\sigma(\bar c_{\mu_0})=\bar c_{\mu_0}$. Si l'identité vaut pour toutes
les étiquettes antérieures à $\mu$, alors
\[
 b\sigma(\bar c_\mu)
 =\bar c_\mu+\sum_{\nu<\mu}B_{\nu\mu}\bar c_\nu
  -\sum_{\nu<\mu}B_{\nu\mu}\bar c_\nu
 =\bar c_\mu.
\]
La récurrence prouve \eqref{eq:pdf2-hodge-sigma-right-inverse} ; elle n'utilise pas
l'injectivité du module de présentation ni une annularité auxiliaire. Appliquer
\eqref{eq:pdf2-hodge-factorizacion-relativa} donne
\[
 e_{N+1}^{\rm rel}s_{N+1}^{\rm rel}
 =e_{N+1}^{\rm rel}a_{N+1}\sigma_{N+1}^{\rm rel}
 =b_{N+1}\sigma_{N+1}^{\rm rel}=I.
\]
La table est formée avec toutes les cellules de
$\Sigma_{\rm coh}^{\rm mult}$ et avec le terminal ; elle ne sélectionne donc pas un
atlas particulier après avoir pris connaissance de l'observation.
\end{proof}

\section{Naturalité et opérateur universel d'extension}

Un raffinement admissible conserve les étiquettes complètes et remplace une
cellule par son expression de Mayer--Vietoris avec des termes d'ordre strictement
inférieur. Si $r_{\mathscr A}$, $r_{\rm H}$ et $r_{\rm S}$ sont les applications
induites, alors
\begin{equation}\label{eq:pdf2-hodge-entrelazadores-refinamiento}
 b' r_{\mathscr A}=r_{\rm H}b,
 \qquad
 a' r_{\mathscr A}=r_{\rm S}a,
 \qquad
 e'^{\rm rel}r_{\rm S}=r_{\rm H}e^{\rm rel}.
\end{equation}
La forme normale globale conserve le premier pivot de chaque étiquette.
L'unicité de la récurrence \eqref{eq:pdf2-hodge-sigma-recursiva} implique
\begin{equation}\label{eq:pdf2-hodge-naturalidad}
 r_{\mathscr A}\sigma_{N+1}^{\rm rel}
 =\sigma_{N+1}^{\prime\,\rm rel}r_{\rm H},
 \qquad
 r_{\rm S}s_{N+1}^{\rm rel}
 =s_{N+1}^{\prime\,\rm rel}r_{\rm H}.
\end{equation}
Ainsi sont conservés orientation, retenue, mémoire, frontière, chemin et terminal.

La dégénérescence unitaire qui fait partie de
\eqref{eq:pdf2-g-coh-datos-operativos} fournit, avant de fixer une classe,
\begin{equation}\label{eq:pdf2-hodge-extension-cero}
 \iota_{N+1,N}:\mathscr S_N(X,p)\longrightarrow\mathscr S_{N+1}(X,p),
 \qquad \rho_{N+1,N}\iota_{N+1,N}=I.
\end{equation}
Définissons le produit fibré de données compatibles
\begin{equation}\label{eq:pdf2-hodge-pullback-extension}
 \mathscr P_{N+1,N}:=
 \{(s,h)\in\mathscr S_N(X,p)\times\mathscr H_{N+1}^p(X):
             e_N(s)=q_{N+1,N}(h)\}.
\end{equation}
Pour $(s,h)\in\mathscr P_{N+1,N}$, le défaut
\begin{equation}\label{eq:pdf2-hodge-defecto-universal}
 \delta_{N+1}(s,h):=
 h-e_{N+1}\iota_{N+1,N}(s)
 \in\mathscr K^{\rm H}_{N+1}
\end{equation}
se corrige au moyen des formules matérielles
\begin{equation}\label{eq:pdf2-hodge-eta-extension}
 \begin{aligned}
  \eta_{N+1}(s,h)&:=
     s_{N+1}^{\rm rel}\bigl(\delta_{N+1}(s,h)\bigr),\\
  \operatorname{Ext}_{N+1,N}(s,h)&:=
     \iota_{N+1,N}(s)+\eta_{N+1}(s,h).
 \end{aligned}
\end{equation}

\begin{theorem}[Extension finie universelle et naturelle]
\label{thm:pdf2-hodge-extension-universal}
Pour toute donnée compatible $(s,h)$,
\begin{equation}\label{eq:pdf2-hodge-dos-identidades-extension}
 \rho_{N+1,N}\operatorname{Ext}_{N+1,N}(s,h)=s,
 \qquad
 e_{N+1}\operatorname{Ext}_{N+1,N}(s,h)=h.
\end{equation}
L'opérateur commute avec tout raffinement admissible et se définit sans choisir
un cycle algébrique.
\end{theorem}

\begin{proof}
Le carré \eqref{eq:pdf2-hodge-cuadrado-restriccion} et la condition de
produit fibré donnent
\[
 q_{N+1,N}\delta_{N+1}(s,h)
 =q_{N+1,N}h-e_Ns=0,
\]
de sorte que \eqref{eq:pdf2-hodge-eta-extension} est typée. Comme
$\eta_{N+1}$ appartient à $\ker\rho_{N+1,N}$, la première identité de
\eqref{eq:pdf2-hodge-dos-identidades-extension} découle de
\eqref{eq:pdf2-hodge-extension-cero}. Pour la deuxième, on utilise
\eqref{eq:pdf2-hodge-right-inverse} :
\[
 e_{N+1}\operatorname{Ext}_{N+1,N}(s,h)
 =e_{N+1}\iota_{N+1,N}(s)
  +e_{N+1}^{\rm rel}s_{N+1}^{\rm rel}\delta_{N+1}(s,h)=h.
\]
La naturalité est la composition de
\eqref{eq:pdf2-hodge-naturalidad} avec la naturalité de la dégénérescence
unitaire. L'entrée $h$ est une observation cohomologique ; à aucune étape
on ne présuppose qu'elle est la classe d'un cycle.
\end{proof}

\section{Fibre initiale, Mittag--Leffler et publication affirmative}

À profondeur zéro, la section basée construite dans
\eqref{eq:amp-hodge-seccion-base} et publiée dans
\eqref{eq:pdf2-g-coh-seccion-dato} attribue à chaque atome $c_\mu$ de
$\mathscr H_0^p(X)$ un état germe $\widehat c_\mu$. Cette attribution
appartient à la complétude structurelle de l'image cohomologique et est fixée
avant de choisir $\alpha$ ; ce n'est pas un cycle algébrique représentant la
classe finale. Par conséquent,
\begin{equation}\label{eq:pdf2-hodge-seccion-base}
 s_0^{\rm base}\!\left(\sum_\mu t_\mu c_\mu\right)
 :=\sum_\mu t_\mu\widehat c_\mu,
 \qquad e_0s_0^{\rm base}=I.
\end{equation}
Pour $\alpha\in\operatorname{Hdg}^p(X)_{\mathbb Q}$, on définit
récursivement
\begin{equation}\label{eq:pdf2-hodge-recursion-alpha}
 s_0:=s_0^{\rm base}(q_0\alpha),
 \qquad
 s_{N+1}:=\operatorname{Ext}_{N+1,N}(s_N,q_{N+1}\alpha).
\end{equation}

\begin{theorem}[Complétude coinductive et réalisation de Hodge]
\label{thm:pdf2-hodge-completitud-final}
La famille \eqref{eq:pdf2-hodge-recursion-alpha} vérifie
\[
 s_N\in\mathscr F_N(\alpha),
 \qquad \rho_{N+1,N}s_{N+1}=s_N.
\]
En conséquence, les fibres forment un système de Mittag--Leffler non vide. Plus
fortement, la récursion exhibe une section compatible concrète à tous les
niveaux ; aucun \(\varprojlim^1\) de noyaux relatifs n'est invoqué. De plus,
\begin{equation}\label{eq:pdf2-hodge-xi-alpha}
 \xi_\alpha:=(s_N)_{N\ge0}
 \in\varprojlim_N\mathscr F_N(\alpha)\subset X_\chi(X,p).
\end{equation}
La première identité de \eqref{eq:pdf2-g-coh-limite-lector-dato} identifie
matériellement la famille compatible à un point de $X_\chi(X,p)$. Les
projections cylindriques, y compris la différence terminale, séparent les classes par
la même donnée ; par conséquent,
\begin{equation}\label{eq:pdf2-hodge-evaluacion-alpha}
 \operatorname{ev}^{\rm Hdg}_{\infty,X,p}(\xi_\alpha)=\alpha.
\end{equation}
Enfin,
\begin{equation}\label{eq:pdf2-hodge-beta}
 \beta_\alpha:=\operatorname{ch}^{\rm loc}_p
 (R_{\rm Hdg}(\xi_\alpha))\in\operatorname{CH}^p(X)_{\mathbb Q}
\end{equation}
satisfait
\begin{equation}\label{eq:pdf2-hodge-clase-final}
 \operatorname{cl}_{X,p}(\beta_\alpha)=\alpha.
\end{equation}
\end{theorem}

\begin{proof}
L'identité $e_0s_0^{\rm base}=I$ amorce la récurrence. Si
$e_Ns_N=q_N\alpha$, la paire $(s_N,q_{N+1}\alpha)$ appartient à
\eqref{eq:pdf2-hodge-pullback-extension} ; le
théorème~\ref{thm:pdf2-hodge-extension-universal} donne simultanément la restriction
et l'évaluation requises. Cela prouve la non-vacuité et la surjectivité des
transitions de fibres et donc la condition de Mittag--Leffler. La
famille compatible explicite définit $\xi_\alpha$ au moyen de
\eqref{eq:pdf2-g-coh-limite-lector-dato}. Son évaluation et $\alpha$
coïncident dans chaque projection finie et dans le terminal ; la séparation, également
incluse dans cette donnée, les identifie. Appliquer
\eqref{eq:pdf2-hodge-publicacion} produit
\eqref{eq:pdf2-hodge-clase-final}.
\end{proof}

\section{Contrôle relatif non directeur et réfutateurs}

Le conoyau de présentation
\begin{equation}\label{eq:pdf2-hodge-o-rel}
 \mathscr O_{N+1}^{\rm rel}:=\operatorname{coker}
 (d_{\rm coend}:\mathscr R_{N+1}^{\rm rel}\to
                  \mathscr P_{N+1}^{\rm rel})
\end{equation}
et le comparateur
\begin{equation}\label{eq:pdf2-hodge-chi-rel}
 \chi_{N+1}^{\rm rel}:\mathscr O_{N+1}^{\rm rel}\to
 \mathscr K^{\rm H}_{N+1}
\end{equation}
sont des contrôles de redondance de la présentation. La démonstration précédente
ne suppose pas $\ker\chi_{N+1}^{\rm rel}=0$ : une relation redondante peut vivre
dans ce noyau sans empêcher l'extension. La table
\eqref{eq:pdf2-hodge-tabla-triangular} et
$b\sigma=I$ imposent bien
$\operatorname{coker}\chi_{N+1}^{\rm rel}=0$, comme contrôle ultérieur.

\begin{falsifier}[Réfutateurs locaux de la construction]
La preuve est réfutée s'il manque une étiquette d'observation dans
$\Sigma_{\rm coh}^{\rm mult}$, si un coefficient diagonal de
\eqref{eq:pdf2-hodge-tabla-triangular} n'est pas un, si un raffinement change
le premier pivot global, si le générateur terminal est éliminé ou si
$e^{\rm rel}a=b$ échoue. Aucun de ces contrôles n'utilise une classe algébrique
finale comme entrée.
\end{falsifier}

\paragraph{Statut et provenance.}
La structure $\mathfrak G_{\rm coh}$, ses treize coordonnées, le lecteur
$R_{\rm Hdg}\to\operatorname{ch}^{\rm loc}_p\to\operatorname{cl}_{X,p}$
et la séparation par survie sont
\texttt{ARCHITECTURE\_AUTORALE\_PRÉEXISTANTE}/\texttt{EXACT\_INTERNE}.
Les formules récursives
\eqref{eq:pdf2-hodge-sigma-recursiva} et
\eqref{eq:pdf2-hodge-eta-extension} sont
\texttt{FORMALISATION\_NOUVELLE} : elles rendent matériel et réutilisable l'inverse à droite
que la structure contient déjà. Le résultat a la force
\texttt{DÉMONTRÉ\_AVEC\_STRUCTURE\_DE\_}\allowbreak
\texttt{DÉPART\_EXPLICITE} ; il ne reste pas de
prémisse auxiliaire d'algébricité ni d'obligation résiduelle directrice.

\endgroup

\HMTFlushCurrentChapter
\chapter{Rang, régulateur et terme principal de Birch--Swinnerton--Dyer par mémoire déterminantale de la famille arithmétique}
\begingroup
\renewcommand{\chapter}[2][]{}%
\chapter{Image déterminantale et composition vers Birch--Swinnerton--Dyer}

\section{Unité généalogique et produit fibré}

Soit $E/\mathbb Q$ une courbe elliptique et soit
$\mathfrak G_{\rm det}(E)$ l'image déterminantale autonome complète. On
note $\mathcal T_m^{\rm surv}$ sa catégorie d'histoires
survivantes. Le propriétaire déterminantal précédent construit à partir d'une
unité commune la chaîne de lignes
\eqref{eq:amp-bsd-linea-determinante}--\eqref{eq:amp-bsd-cadena-lineas}, le
comparateur Kato--FERS
\eqref{eq:amp-bsd-phi-kato-fers}--\eqref{eq:amp-bsd-inversa-kato-fers}, le
canal Poitou--Tate
\eqref{eq:amp-bsd-comparador-pt-q}--\eqref{eq:amp-bsd-comparador-pt} et le
recollement local--global
\eqref{eq:amp-bsd-triangulo-local}--\eqref{eq:amp-bsd-mapa-localizacion}.
Le tuple de \eqref{eq:pdf2-g-det-explicito} et les identités de
\eqref{eq:pdf2-g-det-identidades-dato} sont la publication conjointe de cette
construction en amont : ils ne sont pas acceptés comme axiomes terminaux et ne sont pas choisis à
partir du rang ou du terme principal. Ce chapitre conserve et compose
ces flèches jusqu'à la conclusion de Birch--Swinnerton--Dyer.

L'objet des coefficients n'est pas seulement la fibre $3$-primaire. On fixe d'abord
le complexe libre intégral basé
\begin{equation}\label{eq:pdf2-bsd-gmz}
 G_{m,\mathbb Z}:=N_*\bigl(N\mathcal T_m^{\rm surv};\mathbb Z\bigr),
\end{equation}
et, pour chaque nombre premier $\ell$ et chaque $n\geq1$, on pose
\begin{equation}\label{eq:pdf2-bsd-coeficientes-todos-ell}
 A_{\ell,n}:=E[\ell^n](\overline{\mathbb Q}),\qquad
 T_\ell E:=\varprojlim_nA_{\ell,n},\qquad
 V_\ell E:=T_\ell E\otimes_{\mathbb Z_\ell}\mathbb Q_\ell,
 \qquad
 G_{m,\ell,n}:=G_{m,\mathbb Z}\otimes_{\mathbb Z}^{\mathbf L}A_{\ell,n}.
\end{equation}
Les transitions ne restent pas implicites : ce sont
\begin{equation}\label{eq:pdf2-bsd-transiciones-ell-n}
 q_{\ell,n}:A_{\ell,n+1}\twoheadrightarrow A_{\ell,n},
 \quad x\longmapsto\ell x,
 \qquad
 Q_{\ell,n}:=1\otimes q_{\ell,n}:G_{m,\ell,n+1}\to G_{m,\ell,n},
\end{equation}
La transition de coefficients dans le complexe de Manin, avec
$Q_{\ell,n}$, induit
$\widetilde Q_{\ell,n}:P_{E,\ell,n+1}^{\rm gen}\to
P_{E,\ell,n}^{\rm gen}$. Ces applications satisfont
$Q_{\ell,n}d=dQ_{\ell,n}$,
$(Q_{\ell,n}\otimes Q_{\ell,n})\Delta_{\rm AW}
=\Delta_{\rm AW}Q_{\ell,n}$ et
$\widetilde Q_{\ell,n}\widehat\Pi_{0,\ell,n+1}
=\widehat\Pi_{0,\ell,n}Q_{\ell,n}$. Ainsi, la limite qui définit $T_\ell E$
est prise dans la même famille de complexes, non seulement dans les groupes de
coefficients.
L'ancienne notation $A_n$ désigne uniquement
$A_{3,n}$ ; elle ne gouverne pas la publication arithmétique. Les matrices de bord,
d'Alexander--Whitney, d'orientation et de retenue sont écrites sur $\mathbb Z$ dans
la base des histoires. Ainsi, pour
$R\in\{\mathbb Z_\ell,\mathbb Z/\ell^n,\mathbb Q_\ell\}$,
\begin{equation}\label{eq:pdf2-bsd-base-change-integral}
 G_{m,R}=G_{m,\mathbb Z}\otimes_{\mathbb Z}^{\mathbf L}R,
 \qquad
 d_{m,R}=d_{m,\mathbb Z}\otimes1,
 \qquad
 \Delta_{{\rm AW},R}=\Delta_{{\rm AW},\mathbb Z}\otimes1.
\end{equation}
Ainsi, le passage de $3$ à un autre nombre premier n'est pas une extension inexistante
$\mathbb Z_3\to\mathbb Z_\ell$ : toutes les fibres s'obtiennent par changement de
base du même propriétaire intégral.

Pour abréger, les formules suivantes sont écrites sur un coefficient
$A_{\ell,n}$ fixe et sont compatibles avec $n$, avec $\ell$ et avec la base
intégrale. Le complexe normalisé
\begin{equation}\label{eq:pdf2-bsd-gmn}
 G_{m,\ell,n}:=N_*\bigl(N\mathcal T_m^{\rm surv};A_{\ell,n}\bigr)
\end{equation}
conserve la diagonale d'Alexander--Whitney
\begin{equation}\label{eq:pdf2-bsd-aw}
 \Delta_{\rm AW}[x_0\to\cdots\to x_q]
 =\sum_{j=0}^q[x_0\to\cdots\to x_j]\otimes
                [x_j\to\cdots\to x_q]
\end{equation}
et la counité
$\varepsilon_{\rm AW}[x]=1$ sur les sommets, zéro en degrés positifs. L'histoire
se couple au complexe de Manin au moyen de
\begin{equation}\label{eq:pdf2-bsd-pullback}
 P_{E,\ell,n}^{\rm gen}:=
 P_{E,\ell,n}^{\rm Man}\times_{A_{\ell,n}[0]}G_{m,\ell,n}.
\end{equation}
Si $s_{\ell,n}:A_{\ell,n}[0]\to P_{E,\ell,n}^{\rm Man}$ est la section
basée, alors
\[
 \widehat\Pi_{0,\ell,n}(c)=
 (s_{\ell,n}\varepsilon_{\rm AW}(c),c)
\]
conserve simultanément la coordonnée modulaire et l'histoire. Cette partie de la
construction est matérielle et n'utilise pas le terme principal de $L(E,s)$. Les
carrés
\begin{equation}\label{eq:pdf2-bsd-pullback-base-change}
 P_{E,\mathbb Z}^{\rm gen}\otimes_{\mathbb Z}^{\mathbf L}
      \mathbb Z/\ell^n
 \xrightarrow{\ \sim\ }P_{E,\ell,n}^{\rm gen},
 \qquad
 \widehat\Pi_{0,\mathbb Z}\otimes1=\widehat\Pi_{0,\ell,n}
\end{equation}
prouvent qu'unité, pullback et counité sont les mêmes pour tous les
coefficients.

Les deux publications intégrales et leurs changements de base
\begin{equation}\label{eq:pdf2-bsd-copub}
 P_{E,R}=(P_{E,R}^K,P_{E,R}^{\rm PT}),
 \qquad P_{E,R}^K(1_E)=z_{K,R},\quad
 P_{E,R}^{\rm PT}(1_E)=z_{{\rm PT},R}
\end{equation}
doivent provenir de la même unité $1_E$. La dualité réduite est typée par
\[
 D_3(M):=\operatorname{RHom}_\Lambda(M^\iota,\Lambda)[-3],
 \qquad
 X^{\rm orth}:C_E^{\rm red}\xrightarrow{\sim}D_3(C_E^{\rm red}).
\]
Avec $L_{{\rm root},\ell,n}=A_{\ell,n}z_{{\rm PT},\ell,n}$ et le projecteur de complexes
$p_{\rm root}$, la normalisation basée donne
\begin{equation}\label{eq:pdf2-bsd-raiz}
 p_{\rm root}z_{K,\ell,n}=z_{{\rm PT},\ell,n}.
\end{equation}

\section{Remplisseur local et décomposition globale sans effacer le réseau fini}

La décomposition s'écrit d'abord sur le propriétaire intégral. Pour
chaque étiquette finie--singulière $\lambda=(\ell,m,n,\ldots)$, il existe un
coefficient d'incidence $a_\lambda=\ell c_\lambda\in\mathbb Z$ fixé par
$\operatorname{Rel}_{\rm det}$ ; dans la fibre correspondante,
$A_{\ell,n}$ agit par réduction de ce même entier. Le bloc libre
intégral avec $g$ primitif est
\begin{equation}\label{eq:pdf2-bsd-bloque-normal}
 \partial f=a_\lambda e+b,\qquad
 \partial e=g,\qquad
 \partial b=-a_\lambda g.
\end{equation}
Le tensoriser avec la fibre indiquée par $\lambda$ restitue littéralement
$a_\lambda=\ell c_\lambda$. En conséquence, si
$z_K=ar+ue+vb$ est un cycle et $z_{\rm PT}=r$, alors
$\partial z_K=(u-a_\lambda v)g=0$ implique $u=a_\lambda v$. L'égalité racine
\eqref{eq:pdf2-bsd-raiz} donne $a=1$ et donc,
\begin{equation}\label{eq:pdf2-bsd-filler-local}
 h_*:=-vf,\qquad \partial h_*=z_{\rm PT}-z_K.
\end{equation}

La multisection $\Sigma_{\rm det}^{\rm mult}$ étiquette chaque générateur par
\begin{equation}\label{eq:pdf2-bsd-etiqueta-completa}
 \lambda=(\ell,m,n,\gamma,\operatorname{or},\operatorname{Carr},
 \operatorname{Mem},\partial,\tau^{\rm term},\operatorname{loc}).
\end{equation}
L'ordre basé sépare trois classes disjointes et exhaustives : l'unité racine,
les blocs finis--singuliers avec terminal $h_*$ et les cellules du réseau
intégral local--global. Comme $\partial_{\rm det}$ conserve toutes les
étiquettes sauf la descente prescrite par la retenue, ces classes définissent
des projecteurs de complexes $p_{\rm root}$, $p_{\rm FS}$ et $p_{\rm lat}$ avec
\begin{equation}\label{eq:pdf2-bsd-tres-proyectores}
 p_{\rm root}+p_{\rm FS}+p_{\rm lat}=I,
 \qquad p_ip_j=\delta_{ij}p_i.
\end{equation}

\begin{theorem}[Décomposition globale basée]
\label{thm:pdf2-bsd-descomposicion-global}
Il existe un isomorphisme de complexes compatible avec les places, l'orientation,
la corestriction et la dualité
\begin{equation}\label{eq:pdf2-bsd-descomposicion-global}
 \Theta_E:\mathcal C_E^{\rm loc,red}
 \xrightarrow{\ \sim\ }
 \mathcal C_E^{\rm FS,ctr}\oplus\mathcal L_E,
\end{equation}
où
\begin{equation}\label{eq:pdf2-bsd-sumando-fs-y-lattice}
 \begin{aligned}
 \mathcal C_E^{\rm FS,ctr}
  &:=\bigoplus_{\lambda\in\Lambda_E^{\rm FS}}
   [\mathbb Zf_\lambda\to
    \mathbb Ze_\lambda\oplus\mathbb Zb_\lambda
                    \to\mathbb Zg_\lambda],\\
 \mathcal L_E&:=[M_E\xrightarrow{D_E}N_E].
 \end{aligned}
\end{equation}
Chaque bloc du premier terme possède la différentielle intégrale
\eqref{eq:pdf2-bsd-bloque-normal} ; sa fibre $(\ell,n)$ s'obtient par
$-\otimes_{\mathbb Z}^{\mathbf L}A_{\ell,n}$. Le complexe $\mathcal L_E$
conserve, sans
les contracter, les diviseurs de Smith et les composantes
Tate--Shafarevich--Tamagawa.
\end{theorem}

\begin{proof}
Les générateurs sont ordonnés par l'étiquette
\eqref{eq:pdf2-bsd-etiqueta-completa}. L'unité commune occupe le premier
bloc et est éliminée par $q_{\rm root}$. Pour chaque terminal
$(z_{\rm PT},h_*)$, les relations de $\operatorname{Rel}_{\rm det}$
identifient exactement quatre générateurs
$(f_\lambda,e_\lambda,b_\lambda,g_\lambda)$ ; la matrice d'incidence est
\eqref{eq:pdf2-bsd-bloque-normal}, avec $g_\lambda$ primitif car son
étiquette d'arrivée apparaît une seule fois dans la base orientée. Les cellules
restantes portent une étiquette locale--globale et forment les deux réseaux libres
$M_E,N_E$. Il n'y a pas de quatrième classe : la multisection est complète et
$\operatorname{Surv}_{\rm det}$ exige de chaque histoire soit un terminal
FS, soit une coordonnée de réseau. Le changement de bases est unitriangulaire à
diagonale un, car un bord n'introduit que des étiquettes antérieures. Cela
construit $\Theta_E$ et prouve simultanément son exhaustivité et sa
compatibilité avec les bases. La dualité $X^{\rm orth}$ échange les
deux réseaux et conserve chaque bloc FS, de sorte qu'elle entrelace aussi
$\Theta_E$.
\end{proof}

Sur un bloc normal, définissons
\begin{equation}\label{eq:pdf2-bsd-hfs-generadores}
 H_\lambda(g_\lambda)=e_\lambda,
 \qquad H_\lambda(b_\lambda)=f_\lambda,
 \qquad H_\lambda(e_\lambda)=H_\lambda(f_\lambda)=0.
\end{equation}
Un calcul sur les quatre générateurs donne
$\partial H_\lambda+H_\lambda\partial=I$. Par conséquent,
\begin{equation}\label{eq:pdf2-bsd-hfs-global}
 H_{\rm FS}:=\Theta_E^{-1}
 \left(\bigoplus_{\lambda\in\Lambda_E^{\rm FS}}H_\lambda\oplus0\right)
 \Theta_E,
 \qquad
 \partial H_{\rm FS}+H_{\rm FS}\partial=p_{\rm FS}.
\end{equation}
La coordonnée terminale de $\mathfrak G_{\rm det}(E)$ assure
$q_{\rm root}z_K\in\operatorname{im}p_{\rm FS}$ ; alors
$h_*=-H_{\rm FS}(q_{\rm root}z_K)$ reproduit
\eqref{eq:pdf2-bsd-filler-local}. Il est essentiel que le membre droit de
\eqref{eq:pdf2-bsd-hfs-global} soit $p_{\rm FS}$, non $q_{\rm root}$ :
contracter le réseau $\mathcal L_E$ effacerait les facteurs finis que BSD
doit publier.

\section{Propriétaire intégral et changement de base vers tous les nombres premiers}

La décomposition précédente s'exécute d'abord dans la base intégrale des
histoires, non dans une fibre primaire. Pour
\[
 R\in\{\mathbb Z,\mathbb Z_\ell,\mathbb Z/\ell^n,
          \mathbb Q_\ell\}
\]
sont fixés
\begin{equation}\label{eq:pdf2-bsd-complejos-base-change}
 \mathcal C_{E,R}^{\rm loc,red}:=
 \mathcal C_{E,\mathbb Z}^{\rm loc,red}
       \otimes_{\mathbb Z}^{\mathbf L}R,
 \qquad
 \mathcal L_{E,R}:=
 [M_E\otimes R\xrightarrow{D_E\otimes1}N_E\otimes R].
\end{equation}
L'ordre unitriangulaire des étiquettes utilisé dans
\cref{thm:pdf2-bsd-descomposicion-global} a des coefficients entiers. En
conséquence,
\begin{equation}\label{eq:pdf2-bsd-theta-base-change}
 \Theta_{E,R}=\Theta_{E,\mathbb Z}\otimes1,
 \qquad
 H_{{\rm FS},R}=H_{{\rm FS},\mathbb Z}\otimes1,
 \qquad
 \partial H_{{\rm FS},R}+H_{{\rm FS},R}\partial=p_{{\rm FS},R}.
\end{equation}
En particulier, la fibre $3$-primaire est l'une de ces spécialisations et ne
peut pas masquer les fibres $\ell\ne3$.

Avant de former un conoyau fini, on construit l'inverse rationnel. L'ordre
des terminaux de $\Sigma_{\rm det}^{\rm mult}$ donne des bases
$(m_1,\ldots,m_t)$ de $M_E$ et $(n_1,\ldots,n_t)$ de $N_E$ et une bijection
basée $\sigma_E$ entre les deux listes. Si
\[
 d_j:=\bigl\langle n_{\sigma_E(j)}^\vee,D_Em_j\bigr\rangle,
\]
la relation terminale et la dualité $X^{\rm orth}$ donnent $d_j\ne0$ avant de
prendre les dimensions. Définissons
\begin{equation}\label{eq:pdf2-bsd-d0-nilpotente}
 \begin{aligned}
  D_{E,0}:M_E\otimes\mathbb Q&\longrightarrow N_E\otimes\mathbb Q,
  &D_{E,0}m_j&=d_jn_{\sigma_E(j)},\\
  N_E^{\rm tri}&:=D_{E,0}^{-1}(D_E\otimes\mathbb Q)-I.
 \end{aligned}
\end{equation}
Un bord ne contient que son terminal et des étiquettes strictement antérieures ;
$N_E^{\rm tri}$ est donc strictement triangulaire dans la liste ordonnée et
$(N_E^{\rm tri})^t=0$. L'opérateur
\begin{equation}\label{eq:pdf2-bsd-j-e-inverso-racional}
 \boxed{
 J_E:=\left(\sum_{r=0}^{t-1}(-N_E^{\rm tri})^r\right)D_{E,0}^{-1}:
 N_E\otimes\mathbb Q\longrightarrow M_E\otimes\mathbb Q }
\end{equation}
se calcule donc avec un nombre fini d'opérations sur la matrice
d'incidence. Comme
$D_E\otimes\mathbb Q=D_{E,0}(I+N_E^{\rm tri})$, la somme géométrique finie
prouve les deux identités
\begin{equation}\label{eq:pdf2-bsd-j-e-dos-lados}
 J_E(D_E\otimes\mathbb Q)=I_{M_E\otimes\mathbb Q},
 \qquad
 (D_E\otimes\mathbb Q)J_E=I_{N_E\otimes\mathbb Q}.
\end{equation}
C'est la contraction rationnelle du réseau résiduel ; elle n'utilise pas le rang,
$\operatorname{Sha}(E)$, les ordres des groupes ni le terme principal.

La matrice intégrale $D_E$ conserve tous ses diviseurs élémentaires. Pour chaque
nombre premier, on définit, sans le déduire de la fibre $3$-primaire,
\begin{equation}\label{eq:pdf2-bsd-f-ell-def}
 D_{E,\ell}:=D_E\otimes\mathbb Z_\ell,
 \qquad
 F_{E,\ell}:=\operatorname{coker}D_{E,\ell}.
\end{equation}
Les deux identités de \eqref{eq:pdf2-bsd-j-e-dos-lados} impliquent que
$F_E:=\operatorname{coker}D_E$ est fini. Sa
décomposition primaire est donc finie et exacte :
\begin{equation}\label{eq:pdf2-bsd-f-descomposicion-primaria}
 F_E=\operatorname{coker}D_E
 \xrightarrow{\ \sim\ }
 \bigoplus_{\ell}F_{E,\ell},
 \qquad
 [n]\longmapsto([n\otimes1]_\ell)_\ell.
\end{equation}
L'inverse s'obtient générateur par générateur par le théorème chinois des restes
appliqué à chaque diviseur de Smith. C'est l'arête qui permet de publier
simultanément $\operatorname{Sha}(E)$ et tous les facteurs de Tamagawa.

\section{Bocksteins de toutes les pages}

Soit
\begin{equation}\label{eq:pdf2-bsd-s-universal}
 S_{E,0}(T):V_0[[T]]\longrightarrow W_0[[T]]
\end{equation}
la différentielle basée rationnelle construite par les matrices intégrales du
lecteur déterminantal. Pour chaque nombre premier $\ell$ et pour la carte complexe, on
obtient par changement de base
\begin{equation}\label{eq:pdf2-bsd-s-todos-ell}
 S_{E,\ell}(T):=S_{E,0}(T)\otimes_{\mathbb Q}\mathbb Q_\ell,
 \qquad
 S_{E,\mathbb C}(T):=S_{E,0}(T)\otimes_{\mathbb Q}\mathbb C.
\end{equation}
Elle est inversible sur $\mathbb Q((T))$ et donc sur chaque
$\mathbb Q_\ell((T))$ et $\mathbb C((T))$. Nous écrirons $S_E(T)$ lorsque la
base est sans importance. Comme $K[[T]]$ est un anneau de valuation discrète pour
$K=\mathbb Q,\mathbb Q_\ell$ ou $\mathbb C$, sa forme de Smith est
\begin{equation}\label{eq:pdf2-bsd-smith-t-adico}
 U(T)S_E(T)V(T)=\operatorname{diag}
 (T^{a_1},\ldots,T^{a_s},1,\ldots,1),
 \qquad a_i\ge1.
\end{equation}
Les exposants $a_i$ sont les mêmes dans toutes les fibres : ce sont les valuations
en $T$ des mineurs non nuls de la matrice rationnelle
\eqref{eq:pdf2-bsd-s-universal} ; remplacer $\mathbb Q$ par
$\mathbb Q_\ell$ ou $\mathbb C$ ne change pas ces valuations.
On obtient
\begin{equation}\label{eq:pdf2-bsd-orden-smith}
 d:=\operatorname{ord}_T\det S_E(T)
 =\sum_{i=1}^sa_i
 =\sum_{q\ge1}\dim_K V_q.
\end{equation}
La filtration produit $A_q=V_q/V_{q+1}$,
$B_q=\operatorname{im}\beta_q$ et des isomorphismes
\begin{equation}\label{eq:pdf2-bsd-bockstein-reducido}
 \bar\beta_q:A_q\xrightarrow{\sim}B_q,
 \qquad
 j_q:B_q\xrightarrow{\sim}A_q^\vee\otimes(I^q/I^{q+1}).
\end{equation}
Avec la page régulière $\bar S_E(0)$ et les signes de Koszul, toutes les
pages se recollent dans
\begin{equation}\label{eq:pdf2-bsd-rboc-lt}
 R_{\rm Boc}^{\rm der}(S_E)
 =\tau_{\rm reg}\otimes\bigotimes_{q\ge1}\det(\bar\beta_q)
 =\left[T^{-d}\det S_E(T)\right]_{T=0}.
\end{equation}
Aucune page n'est remplacée par l'ordre total $d$.

\section{Table complète Kato--FERS avant le déterminant}

Pour une histoire non dégénérée
$b_\lambda=[x_0\to\cdots\to x_i]$, soient
$b_{\lambda,\le j}=[x_0\to\cdots\to x_j]$ et
$b_{\lambda,\ge j}=[x_j\to\cdots\to x_i]$. Le poids complet du lecteur
de retenue est
\begin{equation}\label{eq:pdf2-bsd-peso-t}
 \nu_T(\gamma):=\operatorname{dep}(\gamma)
                 +\operatorname{Carr}(\gamma),
\end{equation}
de sorte que $\nu_T([x_i])=0$ et tout chemin normalisé de longueur positive
a $\nu_T\ge1$. La base $\mathcal B_i^K$ contient une ligne pour chaque
étiquette complète $\lambda$ ; la base $\mathcal B_i^{\rm Iw}$ contient le
générateur $\upsilon_i(b_\lambda)$ avec les mêmes faces, feuillet, orientation,
mémoire et terminal. Ainsi, $\upsilon_i$ est une bijection écrite par la
même liste d'étiquettes, non choisie au moyen d'un dénombrement ultérieur.

L'application de chaînes est
\begin{equation}\label{eq:pdf2-bsd-phi-kato}
 \Phi_{K/{\rm FERS}}:=
 \mu_{\rm Iw}\circ(P_E^K\widehat\otimes\operatorname{car}_T)
 \circ\Delta_{\rm AW}:
 \mathcal C_E^K\longrightarrow\mathcal C_E^{\rm Iw}.
\end{equation}
Sa table exhaustive, paramétrée par tous les $i$, toutes les étiquettes et
toutes les coupures, est
\begin{equation}\label{eq:pdf2-bsd-tabla-kato-fers}
\begin{array}{c|c|c|c}
 \text{ligne}&\text{coupure}&\text{image}&\text{poids }T\\ \hline
 b_\lambda&j=i&\upsilon_i(b_\lambda)&0\\
 b_\lambda&0\le j<i&
 \epsilon_{\lambda,j}\,
 \mu_{\rm Iw}\!\left(P_E^K(b_{\lambda,\le j}),
       \operatorname{car}_T(b_{\lambda,\ge j})\right)&
 \nu_T(b_{\lambda,\ge j})\ge1\\
 \text{simplexe dégénéré}&\text{tout }j&0&--
\end{array}
\end{equation}
où $\epsilon_{\lambda,j}$ est le signe déjà fixé par
$\operatorname{or}_{\rm det}$. De manière équivalente,
\begin{equation}\label{eq:pdf2-bsd-expansion-kato-fers}
 \Phi_{K/{\rm FERS}}^i(b_\lambda)
 =\upsilon_i(b_\lambda)+
  \sum_{j=0}^{i-1}T^{\nu_T(b_{\lambda,\ge j})}
  \epsilon_{\lambda,j}\,
  \mu_{\rm Iw}\!\left(P_E^K(b_{\lambda,\le j}),
          \widehat{\operatorname{car}}_T(b_{\lambda,\ge j})\right).
\end{equation}

\begin{theorem}[Équivalence basée Kato--FERS]
\label{thm:pdf2-bsd-equivalencia-kato-fers}
L'application \eqref{eq:pdf2-bsd-phi-kato} est une équivalence filtrée de
complexes parfaits, préserve l'unité et, dans toutes les bases complètes,
\begin{equation}\label{eq:pdf2-bsd-i-plus-tb}
 [\Phi_{K/{\rm FERS}}^i](T)=I+TB_i(T),
 \qquad
 [\Phi_{K/{\rm FERS}}^i]^{-1}
 =\sum_{r\ge0}(-TB_i(T))^r.
\end{equation}
En particulier,
\begin{equation}\label{eq:pdf2-bsd-rho-uno}
 \rho_E(T):=
 \prod_i\det(I+TB_i(T))^{(-1)^i},
 \qquad \rho_E(0)=1,
 \qquad
 \rho_{E,\ell}:=\rho_E\otimes_{\mathbb Q}\mathbb Q_\ell.
\end{equation}
\end{theorem}

\begin{proof}
L'identité simpliciale d'Alexander--Whitney, la compatibilité de $P_E^K$
avec les faces et l'additivité de $\nu_T$ sous concaténation prouvent que
$\Phi_{K/{\rm FERS}}$ commute avec les différentielles. Dans chaque ligne de
\eqref{eq:pdf2-bsd-tabla-kato-fers}, l'unique terme de poids nul est la
coupure $j=i$ ; tous les autres appartiennent à $T\Lambda$. Comme
$\upsilon_i$ apparie la liste complète des histoires avec elle-même, la
réduction modulo $T$ est l'identité sur chaque générateur. Cela donne
\eqref{eq:pdf2-bsd-i-plus-tb} ; la complétude de $K[[T]]$ fait converger
l'inverse. Pour $i=0$, la table ne contient que l'unité, de sorte que la racine
est normalisée avant l'évaluation d'un terme principal.
\end{proof}

L'équivalence conserve séparément le ledger
\begin{equation}\label{eq:pdf2-bsd-ledger-diez}
 \mathcal J_{\rm det}=\{\deg,\operatorname{Frob},\operatorname{Ner},
 K,\operatorname{Boc},\operatorname{Sha},\operatorname{Tam},
 \operatorname{tors}^{-2},\Omega,\operatorname{exc}\}.
\end{equation}
Dans ces dix étiquettes, on lit respectivement : le poids total ; la norme
d'Euler ; le changement basé Manin--Néron ; la relation finie--singulière ; toutes les
pages $j_q\bar\beta_q$ ; le cône de Kummer ; les quotients de composantes ;
le réseau torsionnel ; l'intégrale de Néron ; et l'orientation exceptionnelle. La
table ne permet pas d'annulations entre étiquettes.

\section{Application Poitou--Tate de Mordell--Weil vers le drapeau complet}

On conserve d'abord le réseau de Mordell--Weil et on ne rationalise qu'ensuite :
\begin{equation}\label{eq:pdf2-bsd-mw-libre}
 M_{E,\mathbb Z}^{\rm MW}:=E(\mathbb Q)/E(\mathbb Q)_{\rm tors},
 \qquad
 M_E^{\rm MW}:=M_{E,\mathbb Z}^{\rm MW}\otimes_{\mathbb Z}K,
 \qquad K=\mathbb Q.
\end{equation}
$\operatorname{Surv}_{{\rm det},\mathbb Z}^{\rm MW}$ désigne le sous-module
libre engendré par les lignes terminales intégrales du terme
Mordell--Weil.
Le relèvement de Kummer avec histoire complète est
\begin{equation}\label{eq:pdf2-bsd-kummer-hmt}
 \begin{aligned}
 \kappa_{E,\mathbb Z}^{\rm HMT}:M_{E,\mathbb Z}^{\rm MW}
   &\longrightarrow
   Z^1(\mathcal C_{E,\mathbb Z}^{\rm red})
   \cap\operatorname{Surv}_{{\rm det},\mathbb Z}^{\rm MW},\\
 \kappa_E^{\rm HMT}&:=
   \kappa_{E,\mathbb Z}^{\rm HMT}\otimes_{\mathbb Z}\mathbb Q:
   M_E^{\rm MW}\longrightarrow Z^1(\mathcal C_E^{\rm red}).
 \end{aligned}
\end{equation}
et $\operatorname{pg}_q$ publie la coordonnée de la page $V_q$.
L'extension $\operatorname{Ext}_{\Gamma,{\rm det}}$, la multisection et la
survie construisent, pour chaque $q$, un recollement basé
\begin{equation}\label{eq:pdf2-bsd-glue-q}
 \operatorname{gl}_q:V_q\longrightarrow
 \operatorname{Surv}_{\rm det}^{\rm MW}
\end{equation}
qui conserve feuillet, retenue, mémoire et terminal. Ses équations
d'unité--counité sont
\begin{equation}\label{eq:pdf2-bsd-unidad-counidad-paginas}
 \operatorname{pg}_rP_E^{\rm PT}\operatorname{gl}_q
 =\delta_{rq}I_{V_q},
 \qquad
 \sum_{q\ge1}\operatorname{gl}_q
     \operatorname{pg}_qP_E^{\rm PT}=I
 \quad\text{dans le terme MW survivant}.
\end{equation}
La première égalité se vérifie étiquette par étiquette : le recollement de niveau
$q$ a la retenue $q$, zéro dans les autres pages et le terminal de cette même
histoire. La deuxième est l'exhaustivité de
$\Sigma_{\rm det}^{\rm mult}$ ; la somme est finie car la tour de Bockstein
se termine.

Pour ce sous-module, posons
\begin{equation}\label{eq:pdf2-bsd-redes-paginas-saturadas}
 V_{q,\mathbb Z}:=\operatorname{pg}_qP_E^{\rm PT}
   (\operatorname{Surv}_{{\rm det},\mathbb Z}^{\rm MW}),
 \qquad
 A_{q,\mathbb Z}:=V_{q,\mathbb Z}/V_{q+1,\mathbb Z}.
\end{equation}
Dans les bases terminales, $\operatorname{pg}_q$ et
$\operatorname{gl}_q$ ont des entrées $0,\pm1$ et leurs deux compositions sont
celles de \eqref{eq:pdf2-bsd-unidad-counidad-paginas}. Par conséquent,
$V_{q+1,\mathbb Z}$ est un facteur direct basé de $V_{q,\mathbb Z}$,
$A_{q,\mathbb Z}$ est libre et saturé, et
$V_q=V_{q,\mathbb Z}\otimes\mathbb Q$,
$A_q=A_{q,\mathbb Z}\otimes\mathbb Q$.
Soit $s_q:A_{q,\mathbb Z}\to V_{q,\mathbb Z}$ la section de la base
terminale. Alors
$V_{k,\mathbb Z}=\bigoplus_{q\geq k}s_q(A_{q,\mathbb Z})$. Le dédoublement
des jets qui conserve les $q$ occurrences d'une ligne de profondeur $q$ est
\begin{equation}\label{eq:pdf2-bsd-desdoblamiento-jets}
 \begin{aligned}
 \operatorname{Jet}_{\rm Boc,\mathbb Z}(E)
   &:=\bigoplus_{q\geq1}\bigoplus_{k=1}^{q}
       A_{q,\mathbb Z}\,e_{q,k},\\
 \mathcal U_{\rm jet}:
 \operatorname{Jet}_{\rm Boc,\mathbb Z}(E)
   &\xrightarrow{\ \sim\ }
     \bigoplus_{k\geq1}V_{k,\mathbb Z},\\
 \mathcal U_{\rm jet}((a_{q,k})_{q,k})
   &:=(v_k)_k,
   \qquad v_k:=\sum_{q\geq k}s_q(a_{q,k}).
 \end{aligned}
\end{equation}
L'inverse décompose chaque $v_k$ dans la somme directe précédente. En particulier,
\begin{equation}\label{eq:pdf2-bsd-dimension-jet}
 \operatorname{rank}_{\mathbb Z}\operatorname{Jet}_{\rm Boc,\mathbb Z}(E)
 =\sum_q q\,\operatorname{rank}_{\mathbb Z}A_{q,\mathbb Z}
 =\sum_k\operatorname{rank}_{\mathbb Z}V_{k,\mathbb Z}=d.
\end{equation}
Le recollement se restreint, dans ces mêmes bases, à
$\operatorname{gl}_{q,\mathbb Z}:V_{q,\mathbb Z}\to
\operatorname{Surv}_{{\rm det},\mathbb Z}^{\rm MW}$.

On définit, sans utiliser les dimensions,
\begin{equation}\label{eq:pdf2-bsd-psi-y-lambda-q}
 \begin{aligned}
  \psi_{q,\mathbb Z}&:=
      \operatorname{pg}_qP_E^{\rm PT}\kappa_{E,\mathbb Z}^{\rm HMT}:
      M_{E,\mathbb Z}^{\rm MW}\to V_{q,\mathbb Z},\\
  \lambda_{q,\mathbb Z}&:=
      \operatorname{pr}_{\rm MW}P_E^{\rm Man}
      \operatorname{gl}_{q,\mathbb Z}:
      V_{q,\mathbb Z}\to M_{E,\mathbb Z}^{\rm MW},\\
  \psi_q&:=\operatorname{pg}_qP_E^{\rm PT}\kappa_E^{\rm HMT}:
      M_E^{\rm MW}\to V_q,\\
  \lambda_q&:=\operatorname{pr}_{\rm MW}P_E^{\rm Man}
      \operatorname{gl}_q:V_q\to M_E^{\rm MW}.
 \end{aligned}
\end{equation}
Les domaines intermédiaires sont liés par les deux compatibilités basées
qui font partie du recollement Poitou--Tate et qui sont écrites ici sur son
action atteignable :
\begin{equation}\label{eq:pdf2-bsd-kummer-glue-compatibilidad}
 \kappa_E^{\rm HMT}\operatorname{pr}_{\rm MW}P_E^{\rm Man}
       \operatorname{gl}_q=\operatorname{gl}_q,
 \qquad
 \operatorname{pr}_{\rm MW}P_E^{\rm Man}\kappa_E^{\rm HMT}
       =I_{M_E^{\rm MW}}.
\end{equation}
Les deux mêmes identités, avec l'indice $\mathbb Z$, valent sur
$M_{E,\mathbb Z}^{\rm MW}$ et $V_{q,\mathbb Z}$, car toutes les matrices
précédentes sont intégrales et basées.

\begin{theorem}[Isomorphisme basé du drapeau Poitou--Tate]
\label{thm:pdf2-bsd-psi-pt}
Les applications
\begin{equation}\label{eq:pdf2-bsd-psi-pt}
 \begin{aligned}
 \Psi_{\rm PT}:M_E^{\rm MW}&\longrightarrow
       \operatorname{Flag}_{\rm Boc}(E):=\bigoplus_{q\ge1}V_q,
 &P&\longmapsto(\psi_q(P))_{q\ge1},\\
 \Psi_{\rm PT}^{-1}:\operatorname{Flag}_{\rm Boc}(E)&\longrightarrow M_E^{\rm MW},
 &(v_q)_q&\longmapsto\sum_{q\ge1}\lambda_q(v_q)
 \end{aligned}
\end{equation}
sont inverses. En conséquence,
\begin{equation}\label{eq:pdf2-bsd-psi-pt-integral}
 \Psi_{{\rm PT},\mathbb Z}:M_{E,\mathbb Z}^{\rm MW}
 \xrightarrow{\ \sim\ }\bigoplus_{q\ge1}V_{q,\mathbb Z},
 \qquad
 \Psi_{{\rm PT},\mathbb Z}^{-1}((v_q)_q)=
 \sum_q\lambda_{q,\mathbb Z}(v_q),
\end{equation}
et, seulement après extension des scalaires,
\begin{equation}\label{eq:pdf2-bsd-rango-orden}
 \operatorname{rank}E(\mathbb Q)
 =\sum_{q\ge1}\dim_K V_q
 =d=\operatorname{ord}_T\det S_E(T).
\end{equation}
\end{theorem}

\begin{proof}
Appliquer \eqref{eq:pdf2-bsd-unidad-counidad-paginas} avec
\eqref{eq:pdf2-bsd-kummer-glue-compatibilidad} donne, avec tous les domaines
intermédiaires visibles,
$\psi_r\lambda_q=\delta_{rq}I$ et
$\sum_q\lambda_q\psi_q=I_{M_E^{\rm MW}}$. Les deux formules de
\eqref{eq:pdf2-bsd-psi-pt} sont donc inverses avant de comparer les dimensions.
Les tables terminales sont intégrales et basées ; le même calcul, sans
dénominateurs, prouve \eqref{eq:pdf2-bsd-psi-pt-integral}. Ainsi, aucun indice de
sous-réseau n'apparaît lorsqu'on choisit ensuite une base de Mordell--Weil.
Ce n'est qu'après que l'on prend les dimensions et utilise
\eqref{eq:pdf2-bsd-orden-smith}. L'accouplement
$j_q\bar\beta_q$ transporte les bases saturées vers les lignes de hauteur dont la
comparaison avec Néron--Tate s'écrit dans
\eqref{eq:pdf2-bsd-altura-nt-generadores}.
\end{proof}

\section{Localisation, rang de Smith et exactitude Sha--Tamagawa}

Dans cette section, $v$ parcourt les places finies ; aux places de bonne réduction, on
a $\Phi_v(\mathbb F_v)=0$ et $c_v=1$, de sorte que les sommes et produits
locaux ont un support fini. La place infinie est réservée à l'évaluation
de la différentielle de Néron dans la quatrième flèche. Pour chaque place finie $v$, on
a le triangle
\[
 \mathcal C_{E,v}^{\rm fin}\to\mathcal C_{E,v}\to
 \mathcal C_{E,v}^{\rm sing}\to\mathcal C_{E,v}^{\rm fin}[1]
\]
et l'application de chaînes
\begin{equation}\label{eq:pdf2-bsd-ell-e}
 \ell_E(x,(y_v)_v)=(\operatorname{res}_vx-i_vy_v)_v.
\end{equation}
Le cône réduit est
\begin{equation}\label{eq:pdf2-bsd-cono-reducido}
 \mathcal C_E^{\rm loc,red}:=
 q_{\rm root}\operatorname{Cone}(\ell_E)[-1].
\end{equation}
En degrés de cocycles, un élément s'écrit $(x,(y_v),(h_v))$ et
\begin{equation}\label{eq:pdf2-bsd-diferencial-cono}
 d(x,(y_v),(h_v))=
 (dx,(dy_v),(\operatorname{res}_vx-i_vy_v-dh_v)_v).
\end{equation}

Après avoir séparé par \eqref{eq:pdf2-bsd-psi-pt} le réseau libre de Mordell--Weil
et son dual, les blocs FS sont contractés par
\eqref{eq:pdf2-bsd-hfs-global}. Dans le réseau restant, l'inverse
\eqref{eq:pdf2-bsd-j-e-inverso-racional} et ses deux vérifications
\eqref{eq:pdf2-bsd-j-e-dos-lados} donnent, sans déduire l'acyclicité de la seule
dualité,
\begin{equation}\label{eq:pdf2-bsd-rango-completo}
 D_E\otimes\mathbb Q:M_E\otimes\mathbb Q
 \xrightarrow{\sim}N_E\otimes\mathbb Q.
\end{equation}
La forme normale intégrale donne
\begin{equation}\label{eq:pdf2-bsd-smith-finito}
 UD_EV=\operatorname{diag}(s_1,\ldots,s_t),
 \qquad
 F_E:=\operatorname{coker}D_E,
 \qquad |F_E|=\prod_{i=1}^t|s_i|<\infty.
\end{equation}

Soit
\begin{equation}\label{eq:pdf2-bsd-pi-componentes}
 \pi_\Phi:N_E\longrightarrow\bigoplus_v\Phi_v(\mathbb F_v)
\end{equation}
la réduction de chaque coordonnée locale modulo
$\mathcal C_{E,v}^{\rm fin}+d\mathcal C_{E,v}$ ; on a
$\pi_\Phi D_E=0$. Pour une classe
$[\xi]\in\operatorname{Sha}(E)$, choisissons ses primitives finies
$y_v$ et des homotopies $h_v$ avec
$\operatorname{res}_v\xi-i_vy_v=dh_v$. La formule
\begin{equation}\label{eq:pdf2-bsd-iota-sha-cociclos}
 \begin{aligned}
 n_\xi&:=\operatorname{pr}_{N_E}\operatorname{pr}_{\mathcal L_E}
       \Theta_E(\xi,(y_v),(h_v)),\\
 \iota_{\rm Sha}[\xi]&:=[n_\xi]\in F_E
 \end{aligned}
\end{equation}
est indépendante des choix par
\eqref{eq:pdf2-bsd-diferencial-cono}.

Pour rendre matérielle la surjectivité finale, soit
$\widetilde\phi_v$ le représentant terminal basé de
$\phi_v\in\Phi_v(\mathbb F_v)$. La somme des counités est corrigée dans
l'unité commune et relevée par l'opérateur structurel :
\begin{equation}\label{eq:pdf2-bsd-lift-componentes}
 \operatorname{Lift}_\Phi((\phi_v)_v):=
 \operatorname{pr}_{N_E}\operatorname{Ext}_{\Gamma,{\rm det}}
 \left((\widetilde\phi_v)_v,
 -\widehat\Pi_0\!\left(\sum_v
       \varepsilon_{\rm AW}(\widetilde\phi_v)\right)\right).
\end{equation}
La relation de corestriction de
$\operatorname{Rel}_{\rm det}$ fait de cette expression un cocycle et donne
\begin{equation}\label{eq:pdf2-bsd-lift-dos-identidades}
 \pi_\Phi\operatorname{Lift}_\Phi=I,
 \qquad
 \operatorname{im}D_E=
 \ker\pi_\Phi\cap\ker(\text{classe globale}).
\end{equation}

\begin{theorem}[Suite exacte finie Sha--Tamagawa]
\label{thm:pdf2-bsd-sha-tamagawa}
Les applications de cocycles précédentes induisent
\begin{equation}\label{eq:pdf2-bsd-sha-tamagawa}
 0\longrightarrow\operatorname{Sha}(E)
 \xrightarrow{\ \iota_{\rm Sha}\ }F_E
 \xrightarrow{\ \partial_{\rm loc}\ }
 \bigoplus_v\Phi_v(\mathbb F_v)
 \longrightarrow0,
 \qquad
 \partial_{\rm loc}[n]:=\pi_\Phi(n).
\end{equation}
En particulier, $\operatorname{Sha}(E)$ est fini et
\begin{equation}\label{eq:pdf2-bsd-cardinal-sts}
 |F_E|=|\operatorname{Sha}(E)|\prod_vc_v,
 \qquad c_v:=|\Phi_v(\mathbb F_v)|.
\end{equation}
\end{theorem}

\begin{proof}
L'égalité $\pi_\Phi D_E=0$ rend bien définie
$\partial_{\rm loc}$. Si \eqref{eq:pdf2-bsd-iota-sha-cociclos} est un
bord, sa composante globale est un bord ; ainsi, $\iota_{\rm Sha}$ est
injective. Une classe de $F_E$ est dans $\ker\partial_{\rm loc}$ si et seulement
si toutes ses coordonnées singulières admettent des représentants finis ; par
\eqref{eq:pdf2-bsd-diferencial-cono}, cela équivaut à appartenir à l'image de
$\iota_{\rm Sha}$. Enfin,
\eqref{eq:pdf2-bsd-lift-dos-identidades} donne un antécédent de chaque tuple de
composantes et prouve la surjectivité. Les termes adjacents libres de
la suite longue sont précisément les deux termes séparés par
$\Psi_{\rm PT}$ et $X^{\rm orth}$ ; les termes adjacents FS sont
contractiles par $H_{\rm FS}$. Leur annulation n'a pas été supposée. La finitude
de $F_E$ provient de \eqref{eq:pdf2-bsd-rango-completo} ; l'injection donne la
 finitude de $\operatorname{Sha}(E)$ et prendre les ordres donne
 \eqref{eq:pdf2-bsd-cardinal-sts}.
\end{proof}

\begin{corollary}[Satisfaction primaire complète et recomposition intégrale]
\label{cor:pdf2-bsd-sha-tamagawa-todos-ell}
Pour chaque nombre premier $\ell$, le changement de base des applications de cocycles produit
la suite exacte
\begin{equation}\label{eq:pdf2-bsd-sha-tamagawa-ell}
 0\longrightarrow\operatorname{Sha}(E)[\ell^\infty]
 \xrightarrow{\ \iota_{{\rm Sha},\ell}\ }F_{E,\ell}
 \xrightarrow{\ \partial_{{\rm loc},\ell}\ }
 \bigoplus_v\Phi_v(\mathbb F_v)[\ell^\infty]
 \longrightarrow0,
\end{equation}
où
\begin{equation}\label{eq:pdf2-bsd-mapas-ell-generadores}
 \iota_{{\rm Sha},\ell}([\xi]\otimes1)
   =[n_\xi\otimes1],
 \qquad
 \partial_{{\rm loc},\ell}([n\otimes1])
   =(\pi_{\Phi,v}(n)\otimes1)_v.
\end{equation}
De plus, la somme de toutes les suites
\eqref{eq:pdf2-bsd-sha-tamagawa-ell} est exactement
\eqref{eq:pdf2-bsd-sha-tamagawa} ; ce n'est pas seulement sa partie $3$-primaire.
\end{corollary}

\begin{proof}
La matrice $D_E$, le cocycle $n_\xi$, la réduction $\pi_\Phi$ et
$\operatorname{Lift}_\Phi$ sont intégraux par
\eqref{eq:pdf2-bsd-complejos-base-change}--
\eqref{eq:pdf2-bsd-theta-base-change}. Comme $\mathbb Z_\ell$ est plat sur
$\mathbb Z$, tensoriser la suite exacte finie
\eqref{eq:pdf2-bsd-sha-tamagawa} conserve l'exactitude. Pour un groupe fini
$A$, on a canoniquement
$A\otimes\mathbb Z_\ell=A[\ell^\infty]$ ; cela identifie les trois termes
et donne les formules sur les générateurs
\eqref{eq:pdf2-bsd-mapas-ell-generadores}. La forme de Smith
\eqref{eq:pdf2-bsd-smith-finito} écrit
$s_i=\prod_\ell\ell^{v_\ell(s_i)}$, de sorte que
\[
 F_E\cong\bigoplus_\ell F_{E,\ell},\qquad
 \operatorname{Sha}(E)\cong
   \bigoplus_\ell\operatorname{Sha}(E)[\ell^\infty],\qquad
 \Phi_v(\mathbb F_v)\cong
   \bigoplus_\ell\Phi_v(\mathbb F_v)[\ell^\infty].
\]
Seuls un nombre fini de nombres premiers et un nombre fini de places interviennent car tous les groupes
sont des quotients du groupe fini $F_E$. Sommer les suites primaires restitue
la suite intégrale. En prenant les ordres d'abord premier par premier puis en multipliant,
\begin{equation}\label{eq:pdf2-bsd-cardinal-sts-todos-ell}
 \prod_\ell|F_{E,\ell}|
 =\prod_\ell|\operatorname{Sha}(E)[\ell^\infty]|
   \prod_v\prod_\ell|\Phi_v(\mathbb F_v)[\ell^\infty]|
 =|\operatorname{Sha}(E)|\prod_vc_v.
\end{equation}
Ainsi, tous les coefficients et tous les
facteurs de Tamagawa sont matériellement traités.
\end{proof}

\section{Comparateur analytique--arithmétique basé}

La coordonnée analytique n'est pas postulée à partir de l'ordre recherché. Soit
$\lambda_3=\log4$, correspondant au générateur cyclotomique basé
$\gamma_3=4$, et définissons dans un disque $U_E$ centré en $1$
\begin{equation}\label{eq:pdf2-bsd-t-e-construido}
 T_E(s):=\frac{\exp(\lambda_3(s-1))-1}{\lambda_3}
 =(s-1)+\frac{\lambda_3}{2}(s-1)^2+O((s-1)^3).
\end{equation}
Par calcul direct,
\begin{equation}\label{eq:pdf2-bsd-t-e-normalizacion-probada}
 T_E(1)=0,
 \qquad T_E'(s)=\exp(\lambda_3(s-1)),
 \qquad T_E'(1)=1.
\end{equation}
Le théorème d'inversion locale fournit, après réduction de $U_E$ si
nécessaire, une carte biholomorphe $s\leftrightarrow T_E(s)$.
Pour comparer toutes les fibres, prenons
$\gamma_\ell=1+\ell$ si $\ell$ est impair et $\gamma_2=5$, et posons
\begin{equation}\label{eq:pdf2-bsd-t-ell-construido}
 T_{E,\ell}(s):=
 \frac{\exp(\log(\gamma_\ell)(s-1))-1}{\log(\gamma_\ell)}.
\end{equation}
Les cartes sont reliées par un unique automorphisme formel
$T_{E,\ell}=\vartheta_{\ell,3}(T_E)$ avec
$\vartheta_{\ell,3}(0)=0$ et $\vartheta_{\ell,3}'(0)=1$. Ainsi, ni
l'ordre en $T$ ni le coefficient principal ne changent en passant de $3$ à $\ell$.

Le prolongement analytique qui atteint cette carte n'est pas non plus déclaré. Par le
théorème de modularité, soit
\[
 f_E(z)=\sum_{m\geq1}a_m(E)e^{2\pi imz}
\]
la forme nouvelle normalisée de poids $2$ et de niveau $N_E$ associée à $E$, et soit
$\varepsilon_E\in\{\pm1\}$ le signe fixé par
\begin{equation}\label{eq:pdf2-bsd-transformacion-modular}
 g_E(y):=f_E\!\left(\frac{iy}{\sqrt{N_E}}\right),
 \qquad
 g_E(1/y)=\varepsilon_Ey^2g_E(y).
\end{equation}
Pour $\Re s>3/2$, l'intégration terme à terme donne
\begin{equation}\label{eq:pdf2-bsd-mellin-semiplano}
 \Lambda_E(s):=N_E^{s/2}(2\pi)^{-s}\Gamma(s)L(E,s)
 =\int_0^\infty g_E(y)y^s\,\frac{dy}{y}.
\end{equation}
On partage l'intégrale en $(0,1)$ et $(1,\infty)$ et, dans la première partie, on
effectue le changement $y=1/t$. La deuxième égalité de
\eqref{eq:pdf2-bsd-transformacion-modular} produit la formule matérielle
\begin{equation}\label{eq:pdf2-bsd-mellin-continuada}
 \Lambda_E^{\rm cont}(s)=
 \int_1^\infty g_E(y)
       \bigl(y^s+\varepsilon_Ey^{2-s}\bigr)\,\frac{dy}{y}.
\end{equation}
La fonction $g_E(y)$ décroît exponentiellement lorsque $y\to\infty$ ;
l'intégrale de \eqref{eq:pdf2-bsd-mellin-continuada} et toutes ses dérivées en
$s$ convergent donc uniformément sur les compacts. Elle définit ainsi une fonction entière et
coïncide avec \eqref{eq:pdf2-bsd-mellin-semiplano} là où la série de Dirichlet
converge. Puisque
$N_E^{s/2}(2\pi)^{-s}\Gamma(s)$ est holomorphe et non nulle près de $s=1$, la
division par ce facteur construit, sans utiliser son ordre d'annulation, le germe
holomorphe de $L(E,s)$ dans un disque $U_E$ centré en $1$.

Pour éviter de mélanger une fibre $\ell$-adique avec une matrice complexe, on prend
d'abord le polynôme entier, indépendant de $\ell$,
\begin{equation}\label{eq:pdf2-bsd-polinomio-local-entero}
 P_{E,p}(X):=\det(I-X\operatorname{Frob}_p\mid V_\ell E^{I_p})
 \in\mathbb Z[X].
\end{equation}
Pour $p\nmid N_E$, soient $\alpha_p,\beta_p$ les racines de $P_{E,p}$. On prend
$W_p=\mathbb C^2$ avec une base orthonormale et la réalisation normale
\begin{equation}\label{eq:pdf2-bsd-realizacion-normal-frobenius}
 F_p^{\mathbb C}:=\operatorname{diag}(\alpha_p,\beta_p),
 \qquad |\alpha_p|=|\beta_p|=p^{1/2},
 \qquad \|F_p^{\mathbb C}\|_1=2p^{1/2}.
\end{equation}
Pour les mauvais nombres premiers, on choisit de même une réalisation complexe du
polynôme local ; ils sont en nombre fini et sont conservés comme blocs séparés. On n'utilise pas
la matrice compagnon, dont la norme ne serait pas contrôlée uniformément par le
module de ses valeurs propres.

Posons $\mathscr H_E:=\widehat\bigoplus_{p\nmid N_E}W_p$ et, pour $X>0$,
\begin{equation}\label{eq:pdf2-bsd-truncados-fredholm}
 K_{E,X}(s):=\bigoplus_{\substack{p\leq X\\p\nmid N_E}}
          p^{-s}F_p^{\mathbb C},
 \qquad
 L_X(E,s):=\det(I-K_{E,X}(s))^{-1}
       \prod_{p\mid N_E}P_{E,p}(p^{-s})^{-1}.
\end{equation}
Si $K_0$ est un compact de $\{\Re s>3/2\}$ et
$\sigma_0:=\inf_{s\in K_0}\Re s$, alors
\begin{equation}\label{eq:pdf2-bsd-cota-traza}
 \sup_{s\in K_0}\|K_{E,X'}(s)-K_{E,X}(s)\|_1
 \leq C_E\sum_{X<p\leq X'}p^{1/2-\sigma_0}
 \xrightarrow[X,X'\to\infty]{}0.
\end{equation}
Ainsi, $K_{E,X}\to K_E$ en norme trace, localement uniformément, et la
continuité du déterminant de Fredholm prouve, et ne définit pas,
\begin{equation}\label{eq:pdf2-bsd-fredholm-euler}
 \det_F(I-K_E(s))^{-1}
 \prod_{p\mid N_E}P_{E,p}(p^{-s})^{-1}
 =\lim_{X\to\infty}L_X(E,s)=L(E,s).
\end{equation}

Pour passer du produit infini à un complexe fini, on conserve d'abord
l'opérateur, non seulement son déterminant. Soient $\mathscr H_E^+$ et
$\mathscr H_E^-$ deux copies basées de
$\mathscr H_E\oplus\bigoplus_{p\mid N_E}W_p$. Dans le demi-plan
$\Re s>3/2$, on définit
\begin{equation}\label{eq:pdf2-bsd-operadores-euler-inversos}
 \begin{aligned}
  \mathscr A_{E,X}(s)&:=
   (I-K_{E,X}(s))^{-1}\oplus
   \bigoplus_{p\mid N_E}(I-p^{-s}F_p^{\mathbb C})^{-1},\\
  \mathscr A_E(s)&:=\lim_{X\to\infty}\mathscr A_{E,X}(s):
       \mathscr H_E^+\longrightarrow\mathscr H_E^-.
 \end{aligned}
\end{equation}
L'identité de la résolvante et \eqref{eq:pdf2-bsd-cota-traza} donnent, pour chaque
compact $K_0$ du demi-plan,
\begin{equation}\label{eq:pdf2-bsd-resolvente-convergencia-traza}
 \sup_{s\in K_0}
 \|\mathscr A_{E,X'}(s)-\mathscr A_{E,X}(s)\|_1\longrightarrow0.
\end{equation}
Soit $\mathcal U_{\rm bas}:\mathscr H_E^+\to\mathscr H_E^-$ l'identification unitaire qui
envoie chaque vecteur de la base positive sur le vecteur de même nom de la base
négative. Les déterminants suivants désignent ceux de
$\mathcal U_{\rm bas}^{-1}\mathscr A_{E,X}=I+\mathcal L^1$ et
$\mathcal U_{\rm bas}^{-1}\mathscr A_E=I+\mathcal L^1$. Alors
\begin{equation}\label{eq:pdf2-bsd-det-a-x-l-x}
 \det_F(\mathcal U_{\rm bas}^{-1}\mathscr A_{E,X}(s))=L_X(E,s),
 \qquad
 \det_F(\mathcal U_{\rm bas}^{-1}\mathscr A_E(s))=L(E,s).
\end{equation}
La deuxième égalité est précisément la limite de
\eqref{eq:pdf2-bsd-fredholm-euler} ; le signe est correct car on prend
l'inverse de $I-K_E$.

La corestriction finie se construit à partir des terminaux, sans publier
une infinité de nombres premiers dans une matrice finie. Soient $\mathcal I_E^+$ et
$\mathcal I_E^-$ les listes terminales des bases intégrales de $M_E$ et
$N_E$, respectivement ; elles sont finies par
\eqref{eq:pdf2-bsd-sumando-fs-y-lattice}. $\sigma_*>3/2$ étant fixé,
l'entrelaceur local--global les insère dans l'espace d'Euler par
\begin{equation}\label{eq:pdf2-bsd-embedding-terminal-euler}
 \iota_{\rm Eul}^\pm(e_\eta):=
 \left(p^{-\sigma_*}\operatorname{res}_p^\pm
       \operatorname{Surv}_E(e_\eta)\right)_p
 \oplus\left(\operatorname{res}_p^\pm
       \operatorname{Surv}_E(e_\eta)\right)_{p\mid N_E}.
\end{equation}
La bonne somme est de carré sommable par la même borne que
\eqref{eq:pdf2-bsd-cota-traza}. On pose
$\tau_\eta^\pm:=\iota_{\rm Eul}^\pm(e_\eta)$ ; l'indépendance des lignes
terminales donne les counités duales $\{\epsilon_{\eta}^{\pm}\}$, de sorte que
$\epsilon_\eta^\pm(\tau_{\eta'}^\pm)=\delta_{\eta\eta'}$. Les sommes
\begin{equation}\label{eq:pdf2-bsd-proyecciones-terminales-finitas}
 P_E^\pm:=\sum_{\eta\in\mathcal I_E^\pm}
      \tau_\eta^\pm\epsilon_\eta^\pm,
 \qquad Q_E^\pm:=I-P_E^\pm
\end{equation}
sont des projections basées de rang fini ; leurs indices sont fixés par les
étiquettes terminales et non par le rang de Mordell--Weil ni par un ordre
d'annulation.

La coordonnée $\deg$ du ledger décompose
\begin{equation}\label{eq:pdf2-bsd-graduacion-fredholm}
 \mathscr H_E^\pm=\widehat\bigoplus_i\mathscr H_E^{\pm,i},
 \qquad
 \mathcal I_E^{\pm,i}:=\{\eta\in\mathcal I_E^\pm:\deg\eta=i\},
 \qquad
 P_E^{\pm,i}:=
 \sum_{\eta\in\mathcal I_E^{\pm,i}}\tau_\eta^\pm\epsilon_\eta^\pm.
\end{equation}
Les relations locales conservent $\deg$ sauf la différentielle $i\mapsto i+1$ ;
par conséquent,
$\mathscr A_E=\bigoplus_i\mathscr A_E^i$ avec
$\mathscr A_E^i:\mathscr H_E^{+,i}\to\mathscr H_E^{-,i+1}$. Les quatre
blocs de \eqref{eq:pdf2-bsd-bloques-fredholm}, son Schur et ses troncatures
sont pris degré par degré, et, par définition,
\begin{equation}\label{eq:pdf2-bsd-schur-graduado}
 \mathscr D_E=\bigoplus_i\mathscr D_E^i,
 \qquad
 \mathscr D_E^i:=A_{00}^i-A_{01}^i(A_{11}^i)^{-1}A_{10}^i.
\end{equation}

Le prolongement de Fredholm est écrit avant son utilisation. Dans la carte d'Euler
$U_0$, on pose $\mathscr A_{E,0}:=\mathscr A_E$. Les cartes conservent un
complément basé commun $\mathscr Q_E^\pm$ et ne changent que la famille finie
des terminaux. Si
$\tau_{\eta,\alpha}^\pm:=
\operatorname{Ext}_{\Gamma,{\rm det},\alpha}(\tau_\eta^\pm)$ et
$\epsilon_{\eta,\alpha}^\pm$ est sa base duale, on définit
\begin{equation}\label{eq:pdf2-bsd-cambios-fredholm-cartas}
 C_{\alpha\beta}^\pm:=
 I_{\mathscr Q_E^\pm}\oplus
 \sum_{\eta\in\mathcal I_E^\pm}
 \tau_{\eta,\beta}^\pm\epsilon_{\eta,\alpha}^\pm,
 \qquad C_{\alpha\beta}^\pm-I\in\mathcal L^1,
 \qquad
 C_{\beta\gamma}^\pm C_{\alpha\beta}^\pm=C_{\alpha\gamma}^\pm.
\end{equation}
La différence par rapport à l'identité a un rang au plus égal à
$|\mathcal I_E^\pm|$ et est donc de classe trace ; l'holomorphie découle
du fait que chacune des coordonnées terminales, en nombre fini, de
$\operatorname{Ext}_{\Gamma,{\rm det},\alpha}$ est holomorphe. La dernière
égalité se vérifie sur $\mathscr Q_E^\pm$ et sur chaque
$\tau_{\eta,\alpha}^\pm$. On définit
récursivement sur les recouvrements
\begin{equation}\label{eq:pdf2-bsd-continuacion-fredholm-cartas}
 \mathscr A_{E,\beta}:=
 C_{\alpha\beta}^-\mathscr A_{E,\alpha}
 (C_{\alpha\beta}^+)^{-1}.
\end{equation}
Les opérateurs $C_{\alpha\beta}^\pm$ sont holomorphes car chaque coefficient
est une coordonnée de $\operatorname{Ext}_{\Gamma,{\rm det}}$ ; le cocycle de
\eqref{eq:pdf2-bsd-cambios-fredholm-cartas} prouve que
\eqref{eq:pdf2-bsd-continuacion-fredholm-cartas} ne dépend pas de la chaîne de
cartes. C'est la famille de classe déterminant désormais notée
$\mathscr A_E(s)$. Avec
$C_{0\alpha}^\pm:=C_{\alpha-1,\alpha}^\pm\cdots C_{0,1}^\pm$, on transporte
en outre
$P_{E,\alpha}^{\pm,i}:=
C_{0\alpha}^\pm P_E^{\pm,i}(C_{0\alpha}^\pm)^{-1}$ et
$\mathscr H_{E,\alpha}^{\pm,i}:=
C_{0\alpha}^\pm\mathscr H_E^{\pm,i}$ ; ainsi, tous les Schur
$\mathscr D_{E,\alpha}^i$ sont ceux de
\eqref{eq:pdf2-bsd-schur-graduado}, non de nouveaux symboles. Après réduction de chaque
carte, le bloc
\begin{equation}\label{eq:pdf2-bsd-bloques-fredholm}
 A_{ab}(s):=P_a^-\mathscr A_E(s)P_b^+,
 \qquad P_0^\pm:=P_E^\pm,\quad P_1^\pm:=Q_E^\pm,
 \qquad A_{11}(s):Q_E^+\mathscr H_E^+\xrightarrow{\sim}
                       Q_E^-\mathscr H_E^-
\end{equation}
est inversible. C'est la réduction finie standard de l'opérateur de Fredholm ;
le choix concret de $P_E^\pm$ est donné par
\eqref{eq:pdf2-bsd-proyecciones-terminales-finitas}. La base fixe aussi
l'identification unitaire
$U_{11}:Q_E^+\mathscr H_E^+\to Q_E^-\mathscr H_E^-$ ; par construction,
$U_{11}^{-1}A_{11}=I+\mathcal L^1$, et c'est seulement à cet opérateur que s'applique
$\det_F$. Les copies $+$ et $-$ sont déjà la totalisation paire--impaire ; on
n'applique pas un deuxième signe alterné au déterminant de l'opérateur inverse.

Pour chaque troncature qui contient le support de ces terminaux, on définit
matériellement
\begin{equation}\label{eq:pdf2-bsd-cor-lift-truncados}
 \begin{aligned}
  \operatorname{Cor}_{E,X}&:=P_E^-,\\
  \operatorname{Lift}_{E,X}(v)&:=
       v-A_{11,X}^{-1}A_{10,X}v,\\
  \mathscr D_{E,X}(s)&:=
       \operatorname{Cor}_{E,X}\mathscr A_{E,X}(s)
       \operatorname{Lift}_{E,X}
       =A_{00,X}-A_{01,X}A_{11,X}^{-1}A_{10,X}.
 \end{aligned}
\end{equation}
Par multiplication de blocs,
\begin{equation}\label{eq:pdf2-bsd-eliminacion-schur}
 \begin{pmatrix}I&-A_{01,X}A_{11,X}^{-1}\\0&I\end{pmatrix}
 \mathscr A_{E,X}
 \begin{pmatrix}I&0\\-A_{11,X}^{-1}A_{10,X}&I\end{pmatrix}
 =\begin{pmatrix}\mathscr D_{E,X}&0\\0&A_{11,X}\end{pmatrix}.
\end{equation}
Ainsi, la corestriction n'est pas nominale : son relèvement, son inverse sur le complément
et sa différentielle finie sont écrits. La convergence de la résolvante
implique
\begin{equation}\label{eq:pdf2-bsd-corestriccion-limite}
 \mathscr D_E(s):=\lim_{X\to\infty}\mathscr D_{E,X}(s)
 =A_{00}-A_{01}A_{11}^{-1}A_{10},
 \qquad
 \|\mathscr D_{E,X}-\mathscr D_E\|_1\to0.
\end{equation}
Le déterminant du complément n'est pas écarté : par
\eqref{eq:pdf2-bsd-eliminacion-schur},
\begin{equation}\label{eq:pdf2-bsd-determinante-schur-complemento}
 \det_F(\mathcal U_{\rm bas}^{-1}\mathscr A_E)
 =\det_F(U_{11}^{-1}A_{11})\det(\mathscr D_E).
\end{equation}

Soit $\mathcal O(U_E)$ l'anneau des fonctions holomorphes du disque central.
Les espaces rationnels basés $V_0^i,W_0^{i+1}$ de
\eqref{eq:pdf2-bsd-s-universal} fournissent
\begin{equation}\label{eq:pdf2-bsd-dominios-analiticos}
 \mathscr V_E^i:=V_0^i\otimes_{\mathbb Q}\mathcal O(U_E),
 \qquad
 \mathscr W_E^{i+1}:=W_0^{i+1}\otimes_{\mathbb Q}\mathcal O(U_E).
\end{equation}
Les lignes de $\operatorname{Rel}_{\rm det}$ construisent, sur chaque terminal
$\eta$ et à chaque degré, les vecteurs du modèle commun
$v_{\eta,i}:=\upsilon_i(b_\eta)$,
$w_{\eta,i+1}:=\upsilon_{i+1}(b_\eta^{\rm term})$ et les isomorphismes basés
\begin{equation}\label{eq:pdf2-bsd-jvw-generadores}
 \begin{aligned}
  J_{E,V,i}(\tau_{\eta,i}^+)&:=v_{\eta,i},\\
  J_{E,W,i+1}(\tau_{\eta,i+1}^-)&:=w_{\eta,i+1}.
 \end{aligned}
\end{equation}
La relation différentielle de cette même ligne, appliquée à chaque générateur, est
\begin{equation}\label{eq:pdf2-bsd-entrelazamiento-analitico-smith}
 \boxed{
 J_{E,W,i+1}(s)\,\mathscr D_E^i(s)
 =S_E^i(T_E(s))\,J_{E,V,i}(s).}
\end{equation}
C'est l'entrelaceur effectif entre le Schur de Fredholm corestreint et la
matrice de Smith ; ce n'est pas une égalité de déterminants postulée.

Pour typer le prolongement, soit $\mathfrak U_E$ un domaine connexe qui contient
un ouvert du demi-plan $\Re s>3/2$ et le disque $U_E$, et choisissons une chaîne
finie de cartes $\{U_\alpha\}_{\alpha=0}^A$ de
$\operatorname{Ext}_{\Gamma,{\rm det}}$, avec $U_A=U_E$. Dans $U_\alpha$, on
pose
$\mathscr P_{E,\alpha}^i:=
P_{E,\alpha}^{+,i}\mathscr H_{E,\alpha}^{+,i}$ et
$\mathscr D_{E,\alpha}^i$ égal au Schur de
\eqref{eq:pdf2-bsd-corestriccion-limite}. Les applications vers le modèle commun sont
définies, sans choix supplémentaire, par
\begin{equation}\label{eq:pdf2-bsd-j-cartas-modelo-comun}
 J_{E,\alpha,i}(\tau_{\eta,\alpha,i}):=v_{\eta,i}
 \quad\text{dans le domaine},
 \qquad
 J_{E,\alpha,i+1}(\tau_{\eta,\alpha,i+1}):=w_{\eta,i+1}
 \quad\text{dans le codomaine}.
\end{equation}
Ils sont holomorphes car les bases de
\eqref{eq:pdf2-bsd-cambios-fredholm-cartas} le sont. Sur un recouvrement, on
définit, toujours au même degré,
\begin{equation}\label{eq:pdf2-bsd-transiciones-atlas}
 G_{\alpha\beta,i}:=J_{E,\beta,i}^{-1}J_{E,\alpha,i}:
 \mathscr P_{E,\alpha}^i\big|_{U_\alpha\cap U_\beta}
 \longrightarrow
 \mathscr P_{E,\beta}^i\big|_{U_\alpha\cap U_\beta}.
\end{equation}
La composition littérale de ces matrices prouve
\begin{equation}\label{eq:pdf2-bsd-cociclo-atlas}
 G_{\beta\gamma,i}G_{\alpha\beta,i}=G_{\alpha\gamma,i},
 \qquad G_{\alpha\alpha,i}=I,
 \qquad G_{\beta\alpha,i}=G_{\alpha\beta,i}^{-1},
\end{equation}
et l'entrelacement dans les deux cartes donne la compatibilité de chaîne
\begin{equation}\label{eq:pdf2-bsd-cociclo-cadena}
 G_{\alpha\beta,i+1}\mathscr D_{E,\alpha}^i
 =\mathscr D_{E,\beta}^iG_{\alpha\beta,i}.
\end{equation}
Par conséquent,
$g_{\alpha\beta}:=\prod_i\det(G_{\alpha\beta,i})^{(-1)^i}$ satisfait le
cocycle triple. Si
\begin{equation}\label{eq:pdf2-bsd-tau-atlas}
 \tau_E^{\rm atlas}(s):=
 \prod_{\alpha=0}^{A-1}g_{\alpha,\alpha+1}(s),
\end{equation}
alors $\tau_E^{\rm atlas}$ est exactement le transport déterminantal
de la carte d'Euler vers la carte centrale.

On note $J_{E,V}$ et $J_{E,W}$ les totalisations paires--impaires, avec
l'ordre de Koszul, des applications de
\eqref{eq:pdf2-bsd-entrelazamiento-analitico-smith} ; leurs déterminants sont
les produits alternés des déterminants par degré. Dans la carte centrale,
on définit, et pas avant, l'unité conjointe
\begin{equation}\label{eq:pdf2-bsd-u-e-construida}
 u_E(s):=\tau_E^{\rm atlas}(s)
 \det_F(U_{11}^{-1}A_{11}(s))
 \frac{\det J_{E,V}(s)}
      {\det J_{E,W}(s)}.
\end{equation}
Tous ses facteurs sont holomorphes et inversibles en réduisant $U_E$. La
normalisation n'est pas imposée facteur par facteur. Soit $\mathbf e_{\rm cent}$ le
générateur de la ligne déterminante du modèle commun. En effectuant les deux
opérations triangulaires de \eqref{eq:pdf2-bsd-eliminacion-schur}, en transportant
par \eqref{eq:pdf2-bsd-cambios-fredholm-cartas} et en appliquant
\eqref{eq:pdf2-bsd-j-cartas-modelo-comun}, l'action conjointe sur $s=1$ est,
générateur par générateur,
\begin{equation}\label{eq:pdf2-bsd-u-e-normalizacion-probada}
 \begin{aligned}
 &\tau_E^{\rm atlas}(1)
 \det_F(U_{11}^{-1}A_{11}(1))
 \frac{\det J_{E,V}(1)}
      {\det J_{E,W}(1)}\,\mathbf e_{\rm cent}\\
 &\qquad=
 \bigotimes_{i,\eta}
 \bigl(\epsilon_{\eta,i}^-(\tau_{\eta,i}^-)
       \epsilon_{\eta,i}^+(\tau_{\eta,i}^+)^{-1}\bigr)^{(-1)^i}
 \mathbf e_{\rm cent}
 =\mathbf e_{\rm cent},
 \qquad \boxed{u_E(1)=1}.
 \end{aligned}
\end{equation}
La première égalité est l'identité de conservation de la base commune du
comparateur de $\mathfrak G_{\rm det}(E)$, maintenant développée par ses lignes ; la
deuxième utilise
$\epsilon_{\eta,i}^\pm(\tau_{\eta,i}^\pm)=1$. Aucun
des quatre facteurs n'a été normalisé séparément et la valeur recherchée n'a pas été utilisée.

\begin{proposition}[Provenance FERS de la factorisation analytique]
\label{prop:pdf2-bsd-fers-factorizacion-owner}
Dans le disque $U_E$, on a l'identité de sections déterminantales
\begin{equation}\label{eq:pdf2-bsd-factorizacion-analitica}
 \boxed{L(E,s)=u_E(s)\det S_E(T_E(s)).}
\end{equation}
L'égalité est la publication du lecteur FERS de
$\mathfrak G_{\rm det}(E)$ et non une égalité postulée après
le calcul de l'ordre ou du terme principal.
\end{proposition}

\begin{proof}
Dans le demi-plan absolu,
\eqref{eq:pdf2-bsd-det-a-x-l-x} identifie la limite de Fredholm à $L(E,s)$.
L'élimination littérale \eqref{eq:pdf2-bsd-eliminacion-schur} et la continuité
en norme trace donnent
$L(E,s)=\det_F(U_{11}^{-1}A_{11}(s))\det\mathscr D_E(s)$.
Prendre les produits extérieurs ordonnés dans
\eqref{eq:pdf2-bsd-entrelazamiento-analitico-smith} donne
\[
 \det\mathscr D_E(s)
 =\frac{\det J_{E,V}(s)}
        {\det J_{E,W}(s)}
   \det S_E(T_E(s)).
\]
Les équations \eqref{eq:pdf2-bsd-cociclo-atlas} et
\eqref{eq:pdf2-bsd-cociclo-cadena} transportent cette section à travers les cartes ;
le produit exact de leurs transitions est $\tau_E^{\rm atlas}$. Par conséquent, dans
$U_A=U_E$, le coefficient qui reste est précisément
\eqref{eq:pdf2-bsd-u-e-construida}, ce qui prouve
\eqref{eq:pdf2-bsd-factorizacion-analitica}. L'unicité du prolongement de
Mellin de \eqref{eq:pdf2-bsd-mellin-continuada} exclut une autre section. Enfin,
\eqref{eq:pdf2-bsd-u-e-normalizacion-probada} prouve conjointement
$u_E(1)=1$. Aucune étape n'utilise $d$, le rang, $\operatorname{Sha}(E)$ ni le
terme principal.
\end{proof}

Définissons, à partir de la forme de Smith et non de l'ordre analytique,
\begin{equation}\label{eq:pdf2-bsd-a-e-smith}
 a_E(T):=T^{-d}\det S_E(T).
\end{equation}
Par \eqref{eq:pdf2-bsd-rboc-lt}, $a_E$ est holomorphe dans $T=0$ et
\begin{equation}\label{eq:pdf2-bsd-a-e-cero}
 a_E(0)=R_{\rm Boc}^{\rm der}(S_E)\ne0.
\end{equation}
En substituant cette identité et
\eqref{eq:pdf2-bsd-t-e-construido} dans la factorisation analytique, on obtient
matériellement
\begin{equation}\label{eq:pdf2-bsd-expansion-local-deducida}
 \begin{aligned}
 L(E,s)
  &=(s-1)^d\,U_E^{\rm lead}(s),\\
 U_E^{\rm lead}(s)
  &:=u_E(s)
     \left(\frac{T_E(s)}{s-1}\right)^d
     a_E(T_E(s)),\\
 U_E^{\rm lead}(1)
  &=1\cdot1^d\cdot R_{\rm Boc}^{\rm der}(S_E)
   =R_{\rm Boc}^{\rm der}(S_E)\ne0.
 \end{aligned}
\end{equation}
L'exposant et le coefficient principal se déduisent donc du germe
holomorphe non nul ; ils ne sont pas des entrées de $T_E$, $u_E$ ou $S_E$.
Définissons maintenant, après cette déduction,
\begin{equation}\label{eq:pdf2-bsd-dan-zan}
 d_{\rm an}:=\operatorname{ord}_{s=1}L(E,s),
 \qquad
 \mathfrak z_{\rm an}:=
 \left[(s-1)^{-d_{\rm an}}L(E,s)\right]_{s=1}.
\end{equation}
L'équation \eqref{eq:pdf2-bsd-expansion-local-deducida} prouve
\begin{equation}\label{eq:pdf2-bsd-dan-d}
 d_{\rm an}=d,
 \qquad
 \mathfrak z_{\rm an}
 =R_{\rm Boc}^{\rm der}(S_E)
 \quad\text{sur la ligne basée}.
\end{equation}

Du côté arithmétique, le théorème~\ref{thm:pdf2-bsd-psi-pt} et la suite
\eqref{eq:pdf2-bsd-sha-tamagawa} construisent
\begin{equation}\label{eq:pdf2-bsd-dar-zar}
 \begin{aligned}
 d_{\rm ar}&:=\dim_K M_E^{\rm MW}=d,\\
 \mathfrak z_{\rm ar}&:=
 \frac{|\operatorname{Sha}(E)|\operatorname{Reg}(E)\prod_vc_v}
 {|E(\mathbb Q)_{\rm tors}|^2}
 \quad\text{sur la ligne arithmétique orientée}.
 \end{aligned}
\end{equation}

Pour typer les quatre flèches, si $C^\bullet$ est un complexe parfait à
base ordonnée $\mathcal B_i=(c_{i,1},\ldots,c_{i,r_i})$, on utilise
\begin{equation}\label{eq:pdf2-bsd-linea-determinante-base}
 \det C^\bullet:=
 \bigotimes_i\left(\bigwedge^{r_i}C^i\right)^{(-1)^i},
 \qquad
 \mathbf e_C:=\bigotimes_i
 (c_{i,1}\wedge\cdots\wedge c_{i,r_i})^{(-1)^i}.
\end{equation}
Pour un groupe fini $A$ avec une présentation basée
$\mathbb Z^t\xrightarrow{D_A}\mathbb Z^t\to A$, sa ligne
$\det_{\mathbb Z}A$ porte le générateur $\mathbf e_A$ et l'évaluation
\begin{equation}\label{eq:pdf2-bsd-evaluacion-grupo-finito}
 \operatorname{ev}_A(\mathbf e_A):=|\det D_A|=|A|.
\end{equation}
Lorsque cette ligne se combine à une ligne sur $K=\mathbb Q$, on utilise toujours
l'extension des scalaires
\begin{equation}\label{eq:pdf2-bsd-linea-finita-k}
 \det_KA:=\det_{\mathbb Z}A\otimes_{\mathbb Z}K;
\end{equation}
on ne tensorise pas directement des modules sur des anneaux distincts.
En particulier, on choisit les bases de Smith déjà construites dans
\eqref{eq:pdf2-bsd-smith-finito}, $m_1,\ldots,m_t$ de $M_E$ et
$n_1,\ldots,n_t$ de $N_E$, avec $D_Em_i=s_in_i$, et on fixe
\begin{equation}\label{eq:pdf2-bsd-generador-smith-fe}
 \mathbf e_{F_E}:=
 (m_1\wedge\cdots\wedge m_t)\otimes
 (n_1\wedge\cdots\wedge n_t)^{-1},
 \qquad
 \operatorname{ev}_{F_E}(\mathbf e_{F_E})=\prod_i|s_i|.
\end{equation}
Les lignes successives sont
\begin{equation}\label{eq:pdf2-bsd-cinco-lineas-tipadas}
\begin{aligned}
 \Delta_K&:=\det\mathcal C_E^K,\\
 \Delta_{\rm FERS}&:=\det\mathcal C_E^{\rm Iw},\\
 \Delta_{\rm PT}&:=
   \det_K M_E^{\rm MW}\otimes
   \det_K(M_E^{\rm MW})^\vee\otimes\det_KF_E,\\
 \Delta_{\rm STS}&:=
   \det_K M_E^{\rm MW}\otimes
   \det_K(M_E^{\rm MW})^\vee\otimes
   \det_K\operatorname{Sha}(E)
   \otimes\bigotimes_{v<\infty}
      \det_K\Phi_v(\mathbb F_v),\\
 \Delta_{\rm BSD}&:=\mathbb R\,\mathbf e_{\rm BSD}.
\end{aligned}
\end{equation}
La spécialisation principale n'est pas une identification tacite. Notons
$I:=(T)\subset K[[T]]$ et
$\mathfrak c_{K/{\rm FERS}}^{\rm reg}:=
\det_{\rm alt}\Phi_{K/{\rm FERS}}:\Delta_K[[T]]\to
\Delta_{\rm FERS}[[T]]$. Soient
\begin{equation}\label{eq:pdf2-bsd-lineas-leading}
 \begin{aligned}
  \Delta_{\rm FERS}^{(d)}&:=I^d\Delta_{\rm FERS}[[T]],&
  \Delta_{\rm FERS}^{\rm lead}
    &:=I^{-d}\Delta_{\rm FERS}^{(d)}
       \otimes_{K[[T]],\,T\mapsto0}K,\\
  \Delta_K^{(d)}&:=(\mathfrak c_{K/{\rm FERS}}^{\rm reg})^{-1}
       (\Delta_{\rm FERS}^{(d)}),&
  \Delta_K^{\rm lead}
    &:=I^{-d}\Delta_K^{(d)}
       \otimes_{K[[T]],\,T\mapsto0}K.
 \end{aligned}
\end{equation}
Ici, $I^{-d}\Delta^{(d)}$ désigne le module formé par les symboles
$T^{-d}z$ avec $z\in\Delta^{(d)}$. La spécialisation n'est définie
que sur la partie d'ordre au moins $d$ :
\begin{equation}\label{eq:pdf2-bsd-sp-leading-tipado}
 \operatorname{sp}^{(d)}_{\rm FERS}:
 \Delta_{\rm FERS}^{(d)}\longrightarrow
 \Delta_{\rm FERS}^{\rm lead},
 \qquad
 z(T)\longmapsto
 \overline{T^{-d}z(T)}.
\end{equation}
L'équivalence Kato--FERS induit alors l'application typée
\begin{equation}\label{eq:pdf2-bsd-c-k-fers-leading}
 \begin{aligned}
 \mathfrak c_{K/{\rm FERS}}^{\rm lead}:{}
 &\Delta_K^{\rm lead}\xrightarrow{\ \sim\ }
   \Delta_{\rm FERS}^{\rm lead},\\
 &\overline{T^{-d}z_K(T)}\longmapsto
   \overline{T^{-d}
   \mathfrak c_{K/{\rm FERS}}^{\rm reg}(z_K(T))}.
 \end{aligned}
\end{equation}
La formule n'applique pas $\operatorname{sp}^{(d)}$ à une ligne nue. En
particulier, si
\begin{equation}\label{eq:pdf2-bsd-seccion-zeta-leading}
 \begin{aligned}
  \boldsymbol\zeta_{\rm FERS}(T)
    &:=\det S_E(T)\,\mathbf e_{\rm FERS}
       \in\Delta_{\rm FERS}^{(d)},\\
  \boldsymbol\zeta_K(T)
    &:=(\mathfrak c_{K/{\rm FERS}}^{\rm reg})^{-1}
       \boldsymbol\zeta_{\rm FERS}(T)
       \in\Delta_K^{(d)},
 \end{aligned}
\end{equation}
alors
\begin{equation}\label{eq:pdf2-bsd-especializacion-zeta}
 \operatorname{sp}^{(d)}_{\rm FERS}
       (\boldsymbol\zeta_{\rm FERS})
 =R_{\rm Boc}^{\rm der}(S_E)\,
       \mathbf e_{\rm FERS}^{\rm lead},
 \quad
 \mathfrak c_{K/{\rm FERS}}^{\rm lead}
       (\overline{T^{-d}\boldsymbol\zeta_K})
 =R_{\rm Boc}^{\rm der}(S_E)\,
       \mathbf e_{\rm FERS}^{\rm lead}.
\end{equation}
La définition utilise uniquement $d=\operatorname{ord}_T\det S_E(T)$, déjà calculé par
Smith, et non l'ordre analytique.

Les quatre flèches propriétaires, maintenant dotées de domaines et de codomaines explicites,
sont comparées dans la fibre réelle dès le premier endroit où intervient la
hauteur de Néron. Pour éviter de multiplier une ligne sur $K$ par un
scalaire réel n'appartenant pas en général à $K$, on fixe
\begin{equation}\label{eq:pdf2-bsd-lineas-reales-comparador}
 \begin{aligned}
  \Delta_{K,\mathbb R}^{\rm lead}
    &:=\Delta_K^{\rm lead}\otimes_{K,\iota_\infty}\mathbb R,&
  \Delta_{{\rm FERS},\mathbb R}^{\rm lead}
    &:=\Delta_{\rm FERS}^{\rm lead}\otimes_{K,\iota_\infty}\mathbb R,\\
  \Delta_{{\rm PT},\mathbb R}
    &:=\Delta_{\rm PT}\otimes_{K,\iota_\infty}\mathbb R,&
  \Delta_{{\rm STS},\mathbb R}
    &:=\Delta_{\rm STS}\otimes_{K,\iota_\infty}\mathbb R.
 \end{aligned}
\end{equation}
Les quatre flèches sont alors
\begin{equation}\label{eq:pdf2-bsd-cuatro-comparadores}
\begin{aligned}
 \mathfrak c_{K/{\rm FERS},\mathbb R}^{\rm lead}:
   &\Delta_{K,\mathbb R}^{\rm lead}
      \longrightarrow\Delta_{{\rm FERS},\mathbb R}^{\rm lead},
 &\mathfrak c_{K/{\rm FERS},\mathbb R}^{\rm lead}
   &:=\mathfrak c_{K/{\rm FERS}}^{\rm lead}\otimes_K1_{\mathbb R},\\
 \mathfrak c_{{\rm PT},\mathbb R}:
   &\Delta_{{\rm FERS},\mathbb R}^{\rm lead}
      \longrightarrow\Delta_{{\rm PT},\mathbb R},
 &\mathfrak c_{{\rm PT},\mathbb R}(z)&:=
   \mathfrak n_{E,\mathbb R}(z)\otimes\mathbf e_{F_E},\\
 \mathfrak c_{{\rm STS},\mathbb R}:
   &\Delta_{{\rm PT},\mathbb R}
      \longrightarrow\Delta_{{\rm STS},\mathbb R},
 &\mathfrak c_{{\rm STS},\mathbb R}&:=\det\bigl(
   0\to\operatorname{Sha}(E)\to F_E
     \to\textstyle\bigoplus_{v<\infty}\Phi_v(\mathbb F_v)
     \to0\bigr)\otimes_K1_{\mathbb R},\\
 \mathfrak c_{\rm tors,\infty}:
   &\Delta_{{\rm STS},\mathbb R}
      \longrightarrow\Delta_{\rm BSD}.&&
\end{aligned}
\end{equation}

\begin{proposition}[Calcul générateur par générateur des quatre flèches]
\label{prop:pdf2-bsd-cuatro-flechas-generadores}
Les actions sur les bases de
\eqref{eq:pdf2-bsd-cinco-lineas-tipadas} sont les suivantes.

\emph{Première flèche.} Si
$\mathbf e_K^i=\bigwedge_{b\in\mathcal B_i^K}b$ et
$\mathbf e_{\rm FERS}^i=\bigwedge_{b\in\mathcal B_i^K}\upsilon_i(b)$,
la table \eqref{eq:pdf2-bsd-tabla-kato-fers} donne
\begin{equation}\label{eq:pdf2-bsd-c1-generadores}
 \mathfrak c_{K/{\rm FERS}}^{\rm reg}(\mathbf e_K)
 =\prod_i\det(I+TB_i(T))^{(-1)^i}\mathbf e_{\rm FERS}
 =\rho_E(T)\mathbf e_{\rm FERS},
 \qquad \rho_E(0)=1.
\end{equation}
Si $\mathbf e_K^{\rm lead}$ et $\mathbf e_{\rm FERS}^{\rm lead}$ sont les
bases induites par \eqref{eq:pdf2-bsd-lineas-leading}, alors
\begin{equation}\label{eq:pdf2-bsd-c1-leading-generador}
 \mathfrak c_{K/{\rm FERS},\mathbb R}^{\rm lead}
     (\mathbf e_K^{\rm lead}\otimes1)
 =\rho_E(0)(\mathbf e_{\rm FERS}^{\rm lead}\otimes1)
 =\mathbf e_{\rm FERS}^{\rm lead}\otimes1.
\end{equation}

\emph{Deuxième flèche.} Choisissons des bases intégrales
$a_{q,1},\ldots,a_{q,r_q}$ de $A_{q,\mathbb Z}$, vues dans $A_q$, et
$b_{q,j}=\bar\beta_q(a_{q,j})$ de $B_q$. La matrice de la hauteur de la page
est définie par
\begin{equation}\label{eq:pdf2-bsd-hq-generadores}
 j_q(b_{q,j})=
 \sum_{k=1}^{r_q}(H_q)_{kj}\,a_{q,k}^\vee\otimes T^q.
\end{equation}
Par conséquent, sur le produit extérieur complet,
\begin{equation}\label{eq:pdf2-bsd-c2-cuna-paginas}
 \bigwedge_{j=1}^{r_q}j_q\bar\beta_q(a_{q,j})
 =T^{qr_q}\det(H_q)
  \bigwedge_{k=1}^{r_q}a_{q,k}^\vee.
\end{equation}
Le produit sur $q$, avec les lignes régulières, a pour exposant
$\sum_q qr_q=\sum_q\dim_K V_q=d$. Pour
$1\leq j\leq r_q$ et $1\leq k\leq q$, posons
$e_{q,j,k}:=a_{q,j}e_{q,k}$ dans
$\operatorname{Jet}_{\rm Boc,\mathbb Z}(E)$ et définissons
\begin{equation}\label{eq:pdf2-bsd-base-mw-desde-jets}
 P_{q,j,k}:=\Psi_{{\rm PT},\mathbb Z}^{-1}
       \mathcal U_{\rm jet}(e_{q,j,k})
 \in M_{E,\mathbb Z}^{\rm MW}.
\end{equation}
Par \eqref{eq:pdf2-bsd-desdoblamiento-jets},
\eqref{eq:pdf2-bsd-dimension-jet} et
\eqref{eq:pdf2-bsd-psi-pt-integral}, ces $d$ éléments sont une base
intégrale, non un sous-réseau d'indice non contrôlé. Ils sont énumérés dans l'ordre
lexicographique comme $P_1,\ldots,P_d$ et, sur $K$, forment la base de
$M_E^{\rm MW}$ utilisée ensuite.

La hauteur est construite indépendamment de FERS.\@ Pour
$P\in E(\overline{\mathbb Q})$, on fixe
\begin{equation}\label{eq:pdf2-bsd-altura-canonica-construida}
 \widehat h_E(P):=
 \frac12\lim_{n\to\infty}4^{-n}h_x([2^n]P),
 \qquad
 \langle P,Q\rangle_{\rm NT}:=
 \frac12\bigl(\widehat h_E(P+Q)-\widehat h_E(P)-\widehat h_E(Q)\bigr).
\end{equation}
De manière équivalente, le biextenseur de Poincaré rigidifié sur le modèle de
Néron produit, pour toute place $v$, y compris $v=\infty$, le symbole local
réel $\lambda_v(P,Q)$ normalisé par la section nulle, et
\begin{equation}\label{eq:pdf2-bsd-altura-neron-local-global}
 \langle P,Q\rangle_{\rm NT}
 =\sum_v\lambda_v(P,Q)\in\mathbb R.
\end{equation}
Cette hauteur réelle, symétrique et positive n'est pas remplacée par une somme
d'invariants de cup-produits à valeurs dans $\mathbb Q/\mathbb Z$.

Posons
$H_{\rm NT}^{\flat}(P)(Q):=\langle P,Q\rangle_{\rm NT}$ et soit
$\mathcal B_E:=\Psi_{{\rm PT},\mathbb Z}^{-1}\mathcal U_{\rm jet}$,
étendu à $\mathbb R$. Les $q-1$ directions supplémentaires ne sont pas remplacées
par des identités : on calcule le bloc de Gram complet, y compris tous les
croisements entre profondeurs, lignes et jets,
\begin{equation}\label{eq:pdf2-bsd-gram-jet-completo}
 \begin{aligned}
 G_E^{\rm jet}\bigl((q,j,k),(q',j',k')\bigr)
 &:={}
 \left\langle P_{q,j,k},P_{q',j',k'}\right\rangle_{\rm NT}\\
 &=\sum_v\lambda_v
   \left(P_{q,j,k},P_{q',j',k'}\right),
 \end{aligned}
\end{equation}
pour tout $1\le k\le q$, $1\le k'\le q'$. Cette matrice est matérielle : les
$P_{q,j,k}$ ont été construits dans
\eqref{eq:pdf2-bsd-base-mw-desde-jets} et chaque entrée est calculée par
\eqref{eq:pdf2-bsd-altura-neron-local-global}. Par définition du pullback
d'une forme bilinéaire,
\begin{equation}\label{eq:pdf2-bsd-altura-nt-generadores}
 \boxed{
 \mathcal B_E^\vee H_{\rm NT}^{\flat}\mathcal B_E=G_E^{\rm jet}.}
\end{equation}
On n'affirme pas que $G_E^{\rm jet}$ est diagonale par blocs. Dans la base
intégrale $P_1,\ldots,P_d$, le régulateur et son action sur la ligne déterminante
sont exactement
\begin{equation}\label{eq:pdf2-bsd-regulador-cuna-altura}
 \begin{aligned}
 \operatorname{Reg}(E)&:=\det G_E^{\rm jet}>0,\\
 \bigwedge_{j=1}^dH_{\rm NT}^{\flat}(P_j)
 &=\operatorname{Reg}(E)
   (P_1^\vee\wedge\cdots\wedge P_d^\vee).
 \end{aligned}
\end{equation}
La page de Bockstein et la hauteur sont deux accouplements construits, avec des
bases déjà fixées. Le comparateur de Néron entre leurs lignes est défini avant
d'utiliser le terme principal par
\begin{equation}\label{eq:pdf2-bsd-comparador-neron-jet}
 \begin{aligned}
 \mathfrak n_{E,\mathbb R}:
 \Delta_{{\rm FERS},\mathbb R}^{\rm lead}&\longrightarrow
 \bigl(\det_K(M_E^{\rm MW})\otimes_K
 \det_K(M_E^{\rm MW})^\vee\bigr)\otimes_{K,\iota_\infty}\mathbb R,\\
 \mathfrak n_{E,\mathbb R}
 (\mathbf e_{\rm FERS}^{\rm lead}\otimes1)
 &:=\frac{\operatorname{Reg}(E)}{R_{\rm Boc}^{\rm der}(S_E)}
 (P_1\wedge\cdots\wedge P_d)
 \otimes(P_1^\vee\wedge\cdots\wedge P_d^\vee)\otimes1.
 \end{aligned}
\end{equation}
Le dénominateur est non nul par \eqref{eq:pdf2-bsd-a-e-cero} ; numérateur et
dénominateur s'obtiennent, respectivement, à partir du Gram complet
\eqref{eq:pdf2-bsd-gram-jet-completo} et des pages
\eqref{eq:pdf2-bsd-rboc-lt}. Ainsi, la normalisation est un quotient réel de deux
déterminants préalablement calculés, non une orthogonalité inventée ni une
consultation de la formule BSD. L'extension des scalaires s'effectue avant
d'appliquer le régulateur et se conserve dans les flèches PT et STS.

Les pages supérieures et la ligne régulière conservent le générateur intégral de
$F_E$ au moyen d'une application typée. On utilise l'extension rationnelle de la section
$s_q$ déjà construite dans \eqref{eq:pdf2-bsd-desdoblamiento-jets}. La section duale
\begin{equation}\label{eq:pdf2-bsd-glue-dual-tipado}
 \operatorname{gl}_q^\vee:A_q^\vee\longrightarrow
   (\operatorname{Surv}_{\rm det}^{\rm MW})^\vee
\end{equation}
est définie dans la base complète par
$(\operatorname{gl}_q^\vee\varphi)
(\operatorname{gl}_r s_r(a))=\delta_{qr}\varphi(a)$ et par zéro dans les lignes
FS.\@ Soit $\mathcal I_{\rm fin}$ la liste ordonnée des lignes régulières et des paires
$(q,j)$ avec $q\geq2$. Pour $\eta$ avec $q(\eta)\geq2$, on pose
\begin{equation}\label{eq:pdf2-bsd-lifts-filas-smith}
 \begin{aligned}
 \widetilde a_\eta&:=\operatorname{pr}_{M_E}\operatorname{pr}_{\mathcal L_E}
   \Theta_E\operatorname{gl}_{q(\eta)}s_{q(\eta)}(a_\eta),\\
 \widetilde b_\eta^\vee&:=
   \operatorname{pr}_{N_E^\vee}\operatorname{pr}_{\mathcal L_E^\vee}
   (\Theta_E^\vee)^{-1}\operatorname{gl}_{q(\eta)}^\vee
   \bigl((1\otimes\partial_T^{(q(\eta))})
      j_{q(\eta)}\bar\beta_{q(\eta)}(a_\eta)\bigr).
 \end{aligned}
\end{equation}
Dans une ligne régulière, on définit séparément
\begin{equation}\label{eq:pdf2-bsd-lifts-fila-regular}
 \widetilde a_\eta:=
 \operatorname{pr}_{M_E}\operatorname{pr}_{\mathcal L_E}
 \Theta_Es_{\rm reg}(a_\eta),
 \qquad
 \widetilde b_\eta^\vee:=
 \operatorname{pr}_{N_E^\vee}\operatorname{pr}_{\mathcal L_E^\vee}
 (\Theta_E^\vee)^{-1}s_{\rm reg}^\vee
 (\bar S_E(0)a_\eta).
\end{equation}
Tous les domaines sont maintenant $A_q$ ou $A_q^\vee$ selon le cas ; on
n'applique pas $\operatorname{gl}_q$ à un covecteur.

La bijection requise ne se déduit pas de matrices unimodulaires. Les
équations d'unité--counité montrent sur chaque générateur que les listes
$(\widetilde a_\eta)_\eta$ et leurs terminaux duaux sont des bases de $M_E$ et
$N_E^\vee$. Écrivons la liste héritée comme
$\mathcal I_{\rm fin}=\{\eta_1<\cdots<\eta_t\}$ et définissons directement
\begin{equation}\label{eq:pdf2-bsd-nu-construida}
 \boxed{\nu(\eta_j):=j},
 \qquad
 \nu:\mathcal I_{\rm fin}\xrightarrow{\sim}\{1,\ldots,t\}.
\end{equation}
Posons $\mu_j:=\widetilde a_{\eta_j}$ et soit
$\xi_1,\ldots,\xi_t$ la base terminale de $N_E$. Avant Smith, la matrice
matérielle est
\begin{equation}\label{eq:pdf2-bsd-c2-filas-finitas}
 D^{\rm term}_{kj}:=
 \langle\xi_k^\vee,D_E\mu_j\rangle,
 \qquad
 \widetilde b_{\eta_j}^\vee
 =\sum_{k=1}^tD^{\rm term}_{kj}\xi_k^\vee.
\end{equation}
Cette égalité est la relation terminale de
\eqref{eq:pdf2-bsd-descomposicion-global}, évaluée dans la counité
$\xi_k^\vee$ ; elle ne divise pas par un diviseur élémentaire. Si
$UD^{\rm term}V=\operatorname{diag}(s_1,\ldots,s_t)$, alors
\begin{equation}\label{eq:pdf2-bsd-cuna-smith-con-uv}
 \bigwedge_{j=1}^t\widetilde b_{\eta_j}^\vee
 =\det(D^{\rm term})
   (\xi_1^\vee\wedge\cdots\wedge\xi_t^\vee),
 \qquad
 |\det(D^{\rm term})|=\prod_{j=1}^t|s_j|.
\end{equation}
Les facteurs $\det U,\det V\in\{\pm1\}$ sont enregistrés par
$\operatorname{or}_{\rm det}$ lorsqu'on passe de $(\mu_j),(\xi_j)$ aux bases de
Smith $(m_j),(n_j)$. Le quotient des produits extérieurs est donc exactement
$\mathbf e_{F_E}$ de \eqref{eq:pdf2-bsd-generador-smith-fe} ; on n'a pas
supposé que $U$ ou $V$ soient des matrices de permutation et aucune page
ne reste non publiée. En
conséquence,
\begin{equation}\label{eq:pdf2-bsd-c2-generadores}
 \mathfrak c_{{\rm PT},\mathbb R}
 \bigl(R_{\rm Boc}^{\rm der}(S_E)
       (\mathbf e_{\rm FERS}^{\rm lead}\otimes1)\bigr)
 =\operatorname{Reg}(E)\,
  (P_1\wedge\cdots\wedge P_d)\otimes
  (P_1^\vee\wedge\cdots\wedge P_d^\vee)
  \otimes\mathbf e_{F_E}\otimes1.
\end{equation}

\emph{Troisième flèche.} Pour chaque $\ell$, choisissons des générateurs ordonnés
$\sigma_{\ell,a}$ de $\operatorname{Sha}(E)[\ell^\infty]$ et
$\phi_{\ell,v,b}$ de $\Phi_v(\mathbb F_v)[\ell^\infty]$. Les relèvements explicites
$\operatorname{Lift}_{\Phi,\ell}(\phi_{\ell,v,b})$ et les images
$\iota_{{\rm Sha},\ell}(\sigma_{\ell,a})$ forment, par exactitude, une base de
présentation de $F_{E,\ell}$. Dans ces bases, la matrice des relations se
calcule générateur par générateur comme
\begin{equation}\label{eq:pdf2-bsd-c3-matriz-presentacion}
 D_{F,\ell}=
 \begin{pmatrix}
  D_{{\rm Sha},\ell}&X_\ell\\
  0&\displaystyle\bigoplus_{v<\infty}D_{\Phi_v,\ell}
 \end{pmatrix}.
\end{equation}
La colonne $X_\ell$ enregistre le changement de chaque relèvement lorsqu'on modifie son
représentant ; le bloc inférieur gauche est nul car
$\partial_{{\rm loc},\ell}\iota_{{\rm Sha},\ell}=0$. Le produit extérieur de
présentation ne dépend donc pas de $X_\ell$ et
$\det D_{F,\ell}=\det D_{{\rm Sha},\ell}
\prod_{v<\infty}\det D_{\Phi_v,\ell}$. Sur celui-ci,
\begin{equation}\label{eq:pdf2-bsd-c3-generadores}
 \mathbf e_{F_{E,\ell}}
 \longmapsto
 \mathbf e_{\operatorname{Sha}(E)[\ell^\infty]}
 \otimes\bigotimes_{v<\infty}
    \mathbf e_{\Phi_v(\mathbb F_v)[\ell^\infty]},
\end{equation}
et les évaluations de \eqref{eq:pdf2-bsd-evaluacion-grupo-finito} donnent
$|F_{E,\ell}|=|\operatorname{Sha}(E)[\ell^\infty]|
\prod_{v<\infty}|\Phi_v(\mathbb F_v)[\ell^\infty]|$.
Multiplier sur $\ell$ produit
\begin{equation}\label{eq:pdf2-bsd-c3-integral-generadores}
 \mathfrak c_{{\rm STS},\mathbb R}(\mathbf e_{F_E}\otimes1)
 =(\mathbf e_{\operatorname{Sha}(E)}
  \otimes\bigotimes_{v<\infty}\mathbf e_{\Phi_v(\mathbb F_v)})\otimes1,
 \qquad
 |F_E|=|\operatorname{Sha}(E)|\prod_vc_v.
\end{equation}

\emph{Quatrième flèche.} La base de Mordell--Weil précédente et une
différentielle de Néron $\omega_E$ de période positive
$\Omega_E$ étant fixées, la base orientée de $\Delta_{\rm STS}$ est envoyée par
\begin{equation}\label{eq:pdf2-bsd-c4-generadores}
 \begin{aligned}
 &x\,(P_1\wedge\cdots\wedge P_d)\otimes
 (P_1^\vee\wedge\cdots\wedge P_d^\vee)\otimes
 \mathbf e_{\operatorname{Sha}(E)}
 \otimes\bigotimes_{v<\infty}\mathbf e_{\Phi_v(\mathbb F_v)}\\
 &\hspace{28mm}\longmapsto
 \frac{x\,\Omega_E\,
   \operatorname{ev}_{\operatorname{Sha}(E)}
      (\mathbf e_{\operatorname{Sha}(E)})
   \prod_{v<\infty}\operatorname{ev}_{\Phi_v(\mathbb F_v)}
      (\mathbf e_{\Phi_v(\mathbb F_v)})}
 {|E(\mathbb Q)_{\rm tors}|^2}\,
 \mathbf e_{\rm BSD}.
\end{aligned}
\end{equation}
Par \eqref{eq:pdf2-bsd-evaluacion-grupo-finito}, le coefficient de la dernière
ligne est
$x\Omega_E|\operatorname{Sha}(E)|\prod_vc_v/
|E(\mathbb Q)_{\rm tors}|^2$.
Les deux copies du groupe de torsion proviennent des degrés duaux $0$ et $2$ ;
le facteur $\Omega_E$ est l'évaluation de $\omega_E$ sur la composante réelle
orientée. Cette orientation place $\Omega_E\mathfrak z_{\rm ar}$, et non son
inverse, comme coefficient arithmétique de la composition ; la comparaison avec
$\mathfrak z_{\rm an}$ ne sera déduite qu'ensuite au moyen du carré des
trivialisations.
\end{proposition}

\begin{proof}
La première identité est le déterminant alterné des matrices triangulaires
$I+TB_i(T)$. La deuxième applique la spécialisation des pages puis
le comparateur de Néron \eqref{eq:pdf2-bsd-comparador-neron-jet}. Son
numérateur est le déterminant du Gram de jets complet
\eqref{eq:pdf2-bsd-gram-jet-completo}, avec tous les croisements, et son dénominateur
est le déterminant de Bockstein déjà calculé. Par
\eqref{eq:pdf2-bsd-regulador-cuna-altura}, son action sur les produits extérieurs est
\eqref{eq:pdf2-bsd-c2-generadores} ; l'exposant $d$ s'obtient avant de
prendre les dimensions de Mordell--Weil. La troisième est l'isomorphisme canonique de lignes associé à la
suite exacte, écrit dans les sections
$\iota_{{\rm Sha},\ell}$ et $\operatorname{Lift}_{\Phi,\ell}$ déjà
construites. La quatrième est l'évaluation des deux complexes finis de
torsion et de la ligne de Néron. Cela calcule les quatre flèches sans utiliser
l'égalité BSD.
\end{proof}

Leur composition
\begin{equation}\label{eq:pdf2-bsd-comparador-total}
 \mathfrak C_E:=
 \mathfrak c_{\rm tors,\infty}\circ
 \mathfrak c_{{\rm STS},\mathbb R}\circ
 \mathfrak c_{{\rm PT},\mathbb R}\circ
 \mathfrak c_{K/{\rm FERS},\mathbb R}^{\rm lead}:
 \Delta_{K,\mathbb R}^{\rm lead}
 \longrightarrow\Delta_{\rm BSD}
\end{equation}
est bien typée. Avant d'insérer l'élément analytique ou quelque cardinal que ce soit,
on fixe les trivialisations indépendantes
\begin{equation}\label{eq:pdf2-bsd-trivializaciones-independientes}
 \operatorname{tr}_{\rm an}(x\mathbf e_K^{\rm lead}):=x,
 \qquad
 \operatorname{tr}_{\rm Ner}(y\mathbf e_{\rm BSD}):=y.
\end{equation}
La première utilise uniquement la base de Kato spécialisée par
\eqref{eq:pdf2-bsd-lineas-leading} et la normalisation Mellin--Fredholm de
\eqref{eq:pdf2-bsd-fredholm-euler} ; en particulier,
$\operatorname{tr}_{\rm an}
(\overline{T^{-d}\boldsymbol\zeta_K})
=R_{\rm Boc}^{\rm der}(S_E)=\mathfrak z_{\rm an}$ par
\eqref{eq:pdf2-bsd-especializacion-zeta}. La deuxième utilise uniquement la différentielle
de Néron orientée et satisfait
$\operatorname{tr}_{\rm Ner}(\mathbf e_{\rm BSD})=1$. L'identité du
comparateur basé de $\mathfrak G_{\rm det}(E)$, déployée par les quatre
flèches, est le
carré
\begin{equation}\label{eq:pdf2-bsd-cuadrado-trivializaciones}
\begin{array}{ccc}
 \Delta_{K,\mathbb R}^{\rm lead}
   &\xrightarrow{\ \mathfrak C_E\ }&\Delta_{\rm BSD}\\[2mm]
 \big\downarrow\scriptstyle\operatorname{tr}_{\rm an}
   &&\big\downarrow\scriptstyle\operatorname{tr}_{\rm Ner}\\[2mm]
 \mathbb R&=&\mathbb R.
\end{array}
\end{equation}
Ce carré est fixé avec les bases structurelles avant d'évaluer
$L(E,s)$, le régulateur ou les groupes finis.

L'élément analytique est placé, après construction des flèches, dans son
domaine correct :
\begin{equation}\label{eq:pdf2-bsd-elemento-analitico-tipado}
 \mathbf z_E^{\rm an}:=\mathfrak z_{\rm an}\mathbf e_K^{\rm lead}
 =\overline{T^{-d}\boldsymbol\zeta_K}
 \in\Delta_{K,\mathbb R}^{\rm lead}.
\end{equation}
Par \eqref{eq:pdf2-bsd-dan-d},
$\mathfrak z_{\rm an}=R_{\rm Boc}^{\rm der}(S_E)$ ; la deuxième égalité de
\eqref{eq:pdf2-bsd-especializacion-zeta} prouve que la première flèche envoie
$\mathbf z_E^{\rm an}$ exactement sur le membre gauche de
\eqref{eq:pdf2-bsd-c2-generadores}. Les trois autres flèches, calculées sur les
générateurs, donnent
\begin{equation}\label{eq:pdf2-bsd-elementos-linea}
 d_{\rm an}=d=d_{\rm ar},
 \qquad
 \boxed{
 \mathfrak C_E(\mathbf z_E^{\rm an})
 =\Omega_E\mathfrak z_{\rm ar}\,\mathbf e_{\rm BSD}.}
\end{equation}
Appliquer les flèches verticales de
\eqref{eq:pdf2-bsd-cuadrado-trivializaciones} produit l'égalité numérique
\begin{equation}\label{eq:pdf2-bsd-igualdad-numerica-deducida}
 \boxed{\mathfrak z_{\rm an}=\Omega_E\mathfrak z_{\rm ar}.}
\end{equation}

\begin{theorem}[Publication de Birch--Swinnerton--Dyer]
\label{thm:pdf2-bsd-publicacion-final}
Pour l'image déterminantale autonome complète
$\mathfrak G_{\rm det}(E)$, on a
\begin{equation}\label{eq:pdf2-bsd-rango-final}
 \operatorname{ord}_{s=1}L(E,s)=\operatorname{rank}E(\mathbb Q)=:r
\end{equation}
et
\begin{equation}\label{eq:pdf2-bsd-formula-final}
 \frac{L^{(r)}(E,1)}{r!\,\Omega_E}
 =\frac{|\operatorname{Sha}(E)|\operatorname{Reg}(E)\prod_vc_v}
 {|E(\mathbb Q)_{\rm tors}|^2}.
\end{equation}
\end{theorem}

\begin{proof}
L'égalité des degrés est
\eqref{eq:pdf2-bsd-dan-d} suivie de
\eqref{eq:pdf2-bsd-rango-orden}. La table Kato--FERS transporte l'élément
analytique vers toutes les pages de Bockstein ; $\Psi_{\rm PT}$ les recolle au réseau
de Mordell--Weil et publie son déterminant de hauteur ; la suite exacte
Sha--Tamagawa décompose le facteur de Smith ; et le dernier comparateur insère
le réseau torsionnel et la période déjà orientée. C'est exactement
\eqref{eq:pdf2-bsd-elementos-linea}. Le carré indépendant
\eqref{eq:pdf2-bsd-cuadrado-trivializaciones} donne
\eqref{eq:pdf2-bsd-igualdad-numerica-deducida} ; substituer les définitions de
$\mathfrak z_{\rm an}$ et $\mathfrak z_{\rm ar}$ produit
\eqref{eq:pdf2-bsd-formula-final}. Aucune des quatre flèches n'est
normalisée avec $r$, $\operatorname{Reg}(E)$,
$|\operatorname{Sha}(E)|$ ou le terme principal : ces valeurs n'apparaissent
qu'à la fin de leurs déterminants respectifs.
\end{proof}

\section{Contrôles non directeurs, réfutateurs et statut}

\begin{falsifier}[Réfutateurs séparés par composante]
La construction échoue si une histoire de
$\Sigma_{\rm det}^{\rm mult}$ n'appartient pas à la partition
\eqref{eq:pdf2-bsd-tres-proyectores} ; si l'on écrit
$\partial H_{\rm FS}+H_{\rm FS}\partial=q_{\rm root}$ et que l'on efface ainsi
$\mathcal L_E$ ; si une ligne non racine de
\eqref{eq:pdf2-bsd-tabla-kato-fers} a un poids nul ; si les équations
d'unité--counité \eqref{eq:pdf2-bsd-unidad-counidad-paginas} échouent ; si
$D_E\otimes\mathbb Q$ n'est pas inversible ; si
$\pi_\Phi\operatorname{Lift}_\Phi\ne I$ ; ou si une flèche de
\eqref{eq:pdf2-bsd-cuatro-comparadores} utilise la valeur finale pour fixer sa
base. Chaque réfutateur rompt une arête concrète et ne rouvre pas les précédentes.
\end{falsifier}

L'homotopie $H_{\rm FS}$ est un contrôle ultérieur du remplisseur explicite
$h_*=-vf$ et non une prémisse directrice. Les formes de Smith, les
déterminants et les cardinaux sont calculés après la construction des applications ;
ils ne remplacent pas les tables, le recollement des pages ni l'exactitude des cocycles.

\paragraph{Statut et provenance.}
Alexander--Whitney, le produit fibré, l'unité commune, la dualité
réduite, la rigidité de la racine, le remplisseur, toutes les pages de Bockstein et
les quatre lecteurs appartiennent à
\texttt{ARCHITECTURE\_AUTORALE\_PRÉEXISTANTE}/\texttt{EXACT\_INTERNE} de
$\mathfrak G_{\rm det}(E)$. La partition exhaustive
\eqref{eq:pdf2-bsd-descomposicion-global}, la table
\eqref{eq:pdf2-bsd-tabla-kato-fers}, les formules
\eqref{eq:pdf2-bsd-psi-y-lambda-q} et les applications de cocycles
\eqref{eq:pdf2-bsd-iota-sha-cociclos}--
\eqref{eq:pdf2-bsd-lift-componentes} sont
\texttt{FORMALISATION\_NOUVELLE} : elles déploient matériellement les entrelaceurs
déjà contenus dans les treize coordonnées structurelles. Le théorème final a la
force
\texttt{DÉMONTRÉ\_AVEC\_STRUCTURE\_DE\_DÉPART\_EXPLICITE} ; il ne reste pas
d'obligation résiduelle directrice et on n'utilise pas d'homotopie abstraite pour
fabriquer le résultat.

\endgroup

\HMTFlushCurrentChapter
\chapter{Domaines, opérateurs, naturalité et certificats des réalisations terminales du continu}
\begingroup
\renewcommand{\chapter}[2][]{}%
\chapter{Domaines, applications et certificats des \(5\) publications terminales du continu}
\label{chap:hmt83-falsadores-clausuras}

La fermeture conjointe n’uniformise pas la force de cinq résultats distincts. Pour chaque publication sont conservés le domaine, les prémisses, l’application d’entrelacement, le certificat et le contrôle négatif susceptible de l’invalider. La matrice suivante ne résume pas les preuves précédentes : elle en fournit le contrôle inverse et bloque toute promotion qui perdrait le terme terminal, la naturalité, la donnée initiale ou la distinction entre l’état commun et son lecteur terminal.


Ce chapitre réunit les contrôles négatifs des cinq clôtures. Un falsificateur n'invalide que la flèche et le domaine qu'il attaque ; il ne transfère pas son échec à une autre structure initiale. De même, disposer d'une structure complète ne remplace pas la preuve particulière de son application.

\section{Distinction entre structure et application}

Pour chaque $T$, $\mathfrak G_T$ est la vue complète que l'opération intrinsèque $\mathfrak O_T$ extrait sans restreindre l'unique registre commun. Son existence est une affirmation distincte de l'existence du lecteur final $\mathcal P_T:\mathfrak M_T\to\mathcal C_T$. Un échec de $\mathcal P_T$ n'efface pas $\mathfrak G_T$ ; l'existence de $\mathfrak G_T$ ne démontre pas non plus, à elle seule, $\mathcal P_T$.

\section{Obstructions communes aux \(5\) expressions terminales : perte de domaine, de naturalité, de certificat ou de limite compatible}

\begin{center}
\begin{tabular}{p{0.12\textwidth}p{0.36\textwidth}p{0.39\textwidth}}
\hline
Application & Flecha atacada & Falsador material \\
\hline
RH & expression commune $J_+,J_-\to P,\Xi$ & L'existence de $V_+\Xi_+=\Xi$, $V_-\Xi_-=P\Xi$ implique déjà l'inégalité des normes ; sans construction préalable au signe, l'étape est circulaire. \\
NS & racine uniforme du Gram futur & Avec $\Gamma=I$, $L=I$ et un terminal nul, Stein donne $\mathcal U_{0;J}=JI$ ; la télescopie n'implique donc pas l'uniformité. \\
YM & identification des semi-groupes & Un secteur de Hoeffding de degré deux possède l'énergie $6\kappa$ pour $D^{\rm prod}$ et $3\kappa$ pour $D^{\rm fib}$. \\
Hodge & extension finie & Une étiquette manque dans $\Sigma_{\rm coh}^{\rm mult}$, la diagonale de la table cellulaire ne vaut pas un, un raffinement modifie le pivot global ou le terminal est supprimé ; alors $b\sigma=I$ ou $e^{\rm rel}s^{\rm rel}=I$ échoue. \\
BSD & cadena determinantal & $H_{\rm FS}$ s'applique au réseau de Smith, une ligne Kato non racine possède un poids nul, $\operatorname{pg}_r\operatorname{gl}_q$ n'est pas $\delta_{rq}$ ou $\pi_\Phi\operatorname{Lift}_\Phi\ne I$ ; chaque échec rompt le comparateur correspondant. \\
\hline
\end{tabular}
\end{center}

\section{Conditions d'existence, d'unicité et de compatibilité des expressions terminales}

Aucune chaîne d'application n'est considérée comme démontrée par une déclaration ou un résultat numérique isolé. La promotion exige : la formule de l'application ; les domaines et codomaines ; la preuve qu'elle commute avec les différentielles ou les raffinements ; une estimation uniforme en présence d'une limite ; la conclusion classique seulement ensuite ; et une mise à l'épreuve contradictoire du falsificateur correspondant.

Pour Hodge, ces données sont la table unitriangulaire, la récurrence de $\sigma$, l'extension sur le produit fibré, sa naturalité et le passage Mittag--Leffler. Pour BSD, ce sont la partition racine--FS--réseau, la table complète Kato--FERS, les applications de page et de recollement de $\Psi_{\rm PT}$, les applications de cocycles Sha--Tamagawa et les quatre comparateurs munis de bases. Les équations citées dans les chapitres précédents exposent leurs repères mathématiques ; les contrôles auxiliaires interviennent ensuite et ne rouvrent pas ces clôtures.


\endgroup
% Cuerpo científico recuperado y distribuido en su residencia terminal.
% Cuerpo editor-ready sin preámbulo, portada ni metadatos publicables.
% Complementa, sin sustituirlos, los cierres operatoriales y espectrales.

\section{Clôtures négatives et calibrées des équivalences apparentes}
\label{sec:hmt-final-clausuras-negativas-calibradas}

Les résultats de cette section s'appliquent après la génération conjointe
du continu. Leur antécédent causal est l'état unique produit par
APP--TRIT--TPK, l'émergence corrélative de ses cinq constructions
consubstantielles, les carrés de naturalité et le passage conjoint au terminal
et à la limite. Aucune des équivalences examinées ci-dessous ne génère ni ne
sélectionne rétrospectivement ces constructions.

\subsection{Obstructions à la canonicité à la frontière
\texorpdfstring{\(108{:}144=90{:}120\)}{108:144=90:120}}

\begin{theorem}[Il n'existe pas d'inclusion canonique
\texorpdfstring{\(6\hookrightarrow8\)}{6 à 8} sans jauge]
\label{thm:hmt-final-no-inclusion-canonica-seis-ocho}
Soient \(A\) et \(B\) des ensembles nus avec

\[
 |A|=6,
 \qquad |B|=8.
\]
Il n'existe pas de règle, naturelle par rapport à tous les réétiquetages de \(A\)
et de \(B\), qui sélectionne une injection \(A\hookrightarrow B\).
\end{theorem}

\begin{proof}
Si une injection \(\iota:A\hookrightarrow B\) était sélectionnée sans donnée
supplémentaire, la naturalité par rapport à tout \(\beta\in\operatorname{Sym}(B)\)
imposerait \(\beta\iota=\iota\). Une fois choisi \(a\in A\), il existe une permutation
\(\beta\) qui déplace \(\iota(a)\), contradiction. Le transducteur
route par route APP--Witt n'est donc pas canonique à partir des cardinaux. Une fois
fixées une octade, une inclusion et la jauge d'incidence, on obtient bien un
transducteur calibré ; le choix fait partie de son domaine.
\end{proof}

\begin{theorem}[Une fibre nue de cardinal \(28\) ne détermine pas
\(J(8,2)\)]
\label{thm:hmt-final-no-johnson-desde-cardinal}
Soit \(F\) un ensemble sans structure avec \(|F|=28\). Il n'existe pas de relation
d'adjacence non triviale sur \(F\), naturelle par rapport à toutes ses
permutations, qui produise le graphe de Johnson \(J(8,2)\).
\end{theorem}

\begin{proof}
Le groupe \(\operatorname{Sym}(F)\) est transitif sur les paires non ordonnées
d'éléments distincts. Une relation simple et invariante sous tout ce groupe
contient donc toutes les paires ou aucune. Les deux graphes résultants sont
le graphe complet et le graphe vide. En revanche, \(J(8,2)\) a \(28\) sommets et est de
degré \(12\) ; ce n'est aucun des deux. Par conséquent, la projection
calibrée \(1008\to36\) à fibres uniformes de cardinal \(28\) n'implique pas
à elle seule une octade de Witt. Une fois fixées une octade \(O\) et une bijection
\(F\simeq\binom{O}{2}\), l'adjacence de Johnson peut être transportée sur \(F\),
mais elle dépend précisément de cette jauge.
\end{proof}

Ces deux théorèmes clôturent négativement les anciennes formulations
canoniques. Ils n'affectent pas l'invariant exact
\[
 \frac{108}{144}=\frac{90}{120}=\frac68
\]
ni le défaut normalisé \(-1/12\), qui appartiennent à l'extension de rang
et non au choix d'une incidence de Witt.

\subsection{Le tour nonadique complet et le lecteur exceptionnel}

Fixons la jauge
\[
 \chi_{\mathrm{exc}}
 =(
 \prec_{\mathrm{proj}},
 \kappa_{\mathrm{Paley}},
 \mathcal D_{12},
 \iota_{\mathrm{inc}},
 \mathcal K_{\Lambda})
\]
et le lecteur calibré
\(E_{\chi_{\mathrm{exc}}}:\mathscr S_\infty
\to\mathcal I^{\mathrm{exc}}_{\chi_{\mathrm{exc}}}\).
Soit
\[
 \mathscr S_\infty\xleftarrow{s_\infty}
 \mathsf Z^{\Gamma}_{9,\infty}
 \xrightarrow{t_\infty}\mathscr S_\infty
\]
l'arête attestée d'un tour complet. Notons
\(\operatorname{ed}_0\) l'arête centrale, \(\operatorname{ev}_0\)
l'évaluation à l'état initial et \(F_9\) le décalage d'un tour.

\begin{theorem}[Factorisation exceptionnelle par l'arête attestée]
\label{thm:hmt-final-factorizacion-gamma-nueve}
Définissons
\[
 \overline E^{s}_{\chi_{\mathrm{exc}}}
 :=E_{\chi_{\mathrm{exc}}}\circ s_\infty,
 \qquad
 \overline E^{t}_{\chi_{\mathrm{exc}}}
 :=E_{\chi_{\mathrm{exc}}}\circ t_\infty.
\]
Alors
\begin{align}
 E_{\chi_{\mathrm{exc}}}\circ\operatorname{ev}_0
 &=\overline E^{s}_{\chi_{\mathrm{exc}}}
   \circ\operatorname{ed}_0,
 \label{eq:hmt-final-factorizacion-gamma-fuente}\\
 E_{\chi_{\mathrm{exc}}}\circ\operatorname{ev}_0\circ F_9
 &=\overline E^{t}_{\chi_{\mathrm{exc}}}
   \circ\operatorname{ed}_0.
 \label{eq:hmt-final-factorizacion-gamma-destino}
\end{align}
\end{theorem}

\begin{proof}
Les identités d'extrémités sont
\(s_\infty\operatorname{ed}_0=\operatorname{ev}_0\) et
\(t_\infty\operatorname{ed}_0=\operatorname{ev}_0F_9\). La composition avec
le lecteur calibré donne les deux équations.
\end{proof}

\begin{corollary}[Non-équivalence entre le secteur temporel \(+1\) et \(1\)
tour nonadique complet]
\label{cor:hmt-final-no-equivalencia-retorno-gamma}
La factorisation précédente ne peut pas être remplacée, de manière naturelle, par une
factorisation à travers le seul échantillon TRIT \(\tau=+1\).
\end{corollary}

\begin{proof}
La valeur \(\tau=+1\) enregistre le retour, la mémoire et le passé, mais oublie les données
de phase, de retenue, de route et d'extrémité qui distinguent les arêtes complètes. Deux arêtes
ayant le même échantillon temporel peuvent avoir des destinations différentes. Toute application qui
se factorise par cet échantillon identifie les deux arêtes, tandis que
\(\operatorname{ed}_0\) et ses applications source et destination les séparent. Par conséquent,
le secteur \(+1\) n'est pas équivalent à l'arête complète \(K\mapsto K+9\).
En particulier,
\[
 g(K+9)=g(K)
 \quad\text{et}\quad
 B_{K+9}\ne B_K\quad\text{en général}
\]
restent des affirmations distinctes.
\end{proof}

\begin{proposition}[Séparation calibrée face à la publication brute]
\label{prop:hmt-final-separacion-enriquecida-no-cruda}
Supposons que le lecteur exceptionnel enrichi conserve, pour chaque
générateur, le repère \(f_j\) et l'image
\(y_j=\Gamma_{f_jA_Wf_j^{-1}}(w_j)\). Alors il retrouve
\[
 w_j=\Gamma_{f_jA_Wf_j^{-1}}^{-1}(y_j)
\]
et est injectif sur le quotient par les changements de coordonnées inclus dans la
jauge. Si le lecteur brut est obtenu en oubliant les repères ou la mémoire,
\(E^{\mathrm{cr}}=\pi E^{\sharp}\), il n'est pas équivalent au lecteur enrichi
à moins que \(\pi\) ne soit injective sur \(E^{\sharp}(X)\).
\end{proposition}

\begin{proof}
La première affirmation est l'inversion explicite d'un automorphisme dans chaque
fibre. Pour la seconde, si \(\pi\) identifie deux sorties enrichies
distinctes, le lecteur brut identifie leurs antécédents et ne peut pas être
injectif. Réciproquement, l'injectivité de \(\pi\) sur l'image permet de
retrouver la sortie enrichie. Ainsi se clôt la séparation dans le domaine
calibré et se trouve rejetée l'équivalence tacite avec le codomaine brut.
\end{proof}

L'ancien carré écrit avec des symboles non typés
\(\mathrm{pub}_{\mathrm{inc}},\mathrm{coord},\mathcal O\) ne constitue pas une
équivalence mathématique : il ne définit même pas ses morphismes. Son substitut est le
cône naturel des publications archimédienne, exceptionnelle calibrée et
dimensionnelle, dont les carrés de troncature ont bien un domaine, un codomaine et une
action déclarés.

\subsection{Corridor exceptionnel et boucles du transport}

La chaîne Paley--Witt--Golay--Leech--\(V^{\natural}\)--Monstre--Moonshine
est une publication ultérieure de l'état déjà construit. Une étoile
d'incidence et une boucle du groupoïde TPK ne sont pas des objets équivalents par leur
nom : la première est un sous-objet combinatoire ; la seconde est une flèche dont
la source égale la destination, avec composition et inverse. Sans base, orientation et foncteur
de transport, il n'existe pas de comparaison typée. L'identification
automatique est donc close comme non-équivalence. Un foncteur calibré de
transport peut produire des boucles à partir d'étoiles, mais constitue un
résultat supplémentaire et non une prémisse du corridor déjà démontré.

\subsection{Canal infiniment divisible induit par le transfert TPK}

\begin{theorem}[Poissonisation UCP et racines cohérentes]
\label{thm:hmt-final-poissonizacion-ucp}
Soit \(K\) un canal UCP compatible avec les troncatures TPK\@. Pour \(t\ge0\),
définissons
\[
 P_t:=\exp\bigl(t(K-I)\bigr)
 =e^{-t}\sum_{n=0}^{\infty}\frac{t^n}{n!}K^n.
\]
Alors \((P_t)_{t\ge0}\) est un semi-groupe UCP compatible avec toutes les
troncatures et
\[
 P_1=(P_{1/m})^m
 \qquad(m\ge1).
\]
À chaque niveau fini, une dilatation minimale de Stinespring de \(P_{1/m}\)
produit une racine tensorielle cohérente ; la limite projective conserve les
marginales TPK.
\end{theorem}

\begin{proof}
Chaque \(P_t\) est une combinaison convexe convergente en norme des canaux
\(K^n\), et est donc UCP\@. Comme tous les termes sont des puissances du même
opérateur, l'identité exponentielle donne \(P_tP_s=P_{t+s}\). La compatibilité
avec les troncatures se vérifie terme à terme. En prenant \(t=1/m\), on
obtient la racine \(m\)-ième. Stinespring transforme chaque composition finie en
une contraction tensorielle ; la naturalité des canaux rend leurs
marginales cohérentes.
\end{proof}

Le théorème clôt l'existence d'un canal infiniment divisible à
mémoire TPK\@. Il n'affirme pas qu'un canal historique arbitraire coïncide avec cette
poissonisation : cette égalité exigerait de comparer leurs générateurs et reste
rejetée tant qu'une telle comparaison n'est pas fournie.

\section{Clôtures spectrales et frontières de reconnaissance externe}
\label{sec:hmt-final-clausuras-espectrales-externas}

\begin{proposition}[Équivalence calibrée, et non égalité ambiante]
\label{prop:hmt-final-equivalencia-calibrada-vnueve}
Le transport spectral du double cercle détermine une classe unitaire dans le
sous-espace cyclique engendré par son vecteur de référence. Deux représentants
ambiants \(V_9\) ne sont littéralement égaux qu'après avoir fixé une
identification unitaire du complément et une jauge compatible. Sans ces
données, l'égalité ambiante ne reste pas indéterminée : c'est une affirmation plus
forte et non équivalente à la classe spécifique du sous-espace déjà construite.
\end{proposition}

La convergence en norme des résolvantes comprimées, la récupération des
hauteurs, le critère de Riesz pour les multiplicités et la trace diagonale
RH--Moonshine sont clos dans leurs théorèmes respectifs. La trace diagonale
couple l'entier de mémoire au degré de \(V^{\natural}\) ; elle n'affirme pas que la
multiplicité de chaque zéro est la dimension d'une représentation du
Monstre. Cette identification supplémentaire n'est pas équivalente à l'existence de
la trace couplée.

Restent hors du domaine interne, sans rouvrir aucun théorème HMT :
\begin{enumerate}
 \item l'identification de la dichotomie des publications HMT analytiques à
 tous les sous-ensembles arbitraires de \(\mathbb R\) de l'univers ambiant ;
 \item la réalisation macroscopique ou thermodynamique du cristal temporel
 fini HMT déjà construit, y compris la conservation de la cohérence à grand
 volume et sa confrontation expérimentale ;
 \item la consolidation décimale des fenêtres spectrales par une
 exécution concrète d'arithmétique des intervalles ;
 \item une identification représentationnelle supplémentaire entre les multiplicités
 spectrales individuelles et les modules irréductibles du Monstre.
\end{enumerate}
Ces quatre frontières sont des reconnaissances ou des comparaisons externes, et non
 des lacunes de la construction du continu, de la clôture cardinale, de RH ni de
l'entrelacement du double cercle.

La résidence terminale conserve, sans permutation,
\[
\begin{aligned}
 \mathrm{RH}
 &\longrightarrow \text{double cercle et champ spectral}
 \longrightarrow \mathrm{Navier\text{--}Stokes}\\
 &\longrightarrow \mathrm{Yang\text{--}Mills}
 \longrightarrow \mathrm{Hodge}
 \longrightarrow \mathrm{BSD}
 \longrightarrow \text{réfutateurs}.
\end{aligned}
\]











\HMTFlushCurrentChapter
